\documentclass[11pt]{article}

\usepackage[T1]{fontenc}
\usepackage{lmodern}
\usepackage[margin=1in]{geometry}
\usepackage{amsmath,amssymb,amsthm,mathtools}
\usepackage{enumitem}
\usepackage{booktabs,longtable,array}
\usepackage{xcolor}
\usepackage[colorlinks=true,linkcolor=blue!55!black,urlcolor=blue!55!black]{hyperref}
\usepackage[nameinlink,noabbrev]{cleveref}

% Compatibility dispatcher for inherited comma-separated equation-reference lists.
\makeatletter
\let\renewal@singleeqref\eqref
\newcommand{\renewal@eqrefdispatch}[1]{%
  \in@{,}{#1}%
  \ifin@\labelcref{#1}\else\renewal@singleeqref{#1}\fi}
\renewcommand{\eqref}[1]{%
  \begingroup\renewal@eqrefdispatch{#1}\endgroup}
\makeatother

\newtheorem{theorem}{Theorem}[section]
\newtheorem{lemma}[theorem]{Lemma}
\newtheorem{proposition}[theorem]{Proposition}
\newtheorem{corollary}[theorem]{Corollary}
\newtheorem{counterexample}[theorem]{Exact separator}
\newtheorem{openproblem}[theorem]{Open problem}
\theoremstyle{definition}
\newtheorem{definition}[theorem]{Definition}
\newtheorem{assessment}[theorem]{Assessment}
% Compatibility declarations for environments first used in the appended V8 material.
\newtheorem{example}[theorem]{Example}
\newtheorem{firewall}[theorem]{Firewall}
\theoremstyle{remark}
\newtheorem{remark}[theorem]{Remark}

\newcommand{\C}{\mathbb C}
\newcommand{\R}{\mathbb R}
\newcommand{\Ran}{\operatorname{Ran}}
\newcommand{\Ker}{\operatorname{Ker}}
\newcommand{\Tr}{\operatorname{Tr}}
\newcommand{\rank}{\operatorname{rank}}
\newcommand{\HS}{\mathrm{HS}}
\newcommand{\op}{\mathrm{op}}
\newcommand{\dd}{\mathrm d}
\newcommand{\cF}{\mathcal F}
\newcommand{\cH}{\mathcal H}
\newcommand{\cS}{\mathcal S}
\newcommand{\cY}{\mathcal Y}
\newcommand{\id}{\mathrm{id}}

\title{Renewal Standard--Model--Einstein Continuum Research Notes\\
Second Volume, Version V1: Legacy Synthesis and Derivative-Faithful Shell Covariance}
\author{Append-only research continuation}
\date{September 5, 2026}

\begin{document}
\maketitle

\begin{abstract}
This note opens the second volume of the Renewal Standard--Model--Einstein
continuum programme.  The first volume reached Version V157 and is preserved
verbatim as a frozen lineage file.  We first state the mathematical content of
that volume in a compact theorem-level ledger, including the hypotheses under
which each conclusion was obtained and the gates that remain open.  We then
continue the attack at the current bottleneck: the connected covariance in the
first interacting shell derivative.  At finite cutoff we prove an exact Doob
martingale and coordinate-resampling representation.  It counts one genuine
resampled shell coordinate and therefore prevents a spurious multiplication
of covariance weights.  A Banach-valued form applies directly to projected
metric sources.  We next prove a Taylor-jet remainder theorem for the local
renormalization projector and an explicit four-dimensional bubble estimate:
the unsubtracted local shell is marginal, whereas its constant-jet complement
is \(O(|k|^2\Lambda_n^{-2})\) and hence summable.  These results identify the
next honest task: uniform square-function and external-momentum derivative
bounds for the complete BV-coupled Standard Model shell, with an invariant
local projector.  No full interacting continuum measure or coupled Einstein
solution is asserted here.
\end{abstract}

\tableofcontents

\section{Context, lineage, and the target}
\label{sec:v2-context}

The long-term target is a common continuum in which the gauge, Higgs, fermion,
ghost, and gravitational variables arise from one compatible cutoff system;
the renormalized quantum states satisfy the BV/Slavnov identities; their
metric derivatives converge to a conserved stress tensor; and that tensor
serves as the source of a well-posed limiting Einstein evolution.  There are
three distinct limiting operations behind this sentence: removal of the
spectral or lattice cutoff, passage from finite-dimensional BV pushforwards to
a state on the continuum observable algebra, and stability of the nonlinear
Einstein evolution under the resulting source convergence.  A proof of one of
these operations is not silently used as a proof of either of the others.

The first volume is frozen at
\begin{center}
\small
\texttt{Renewal\_SM\_Einstein\_Continuum\_Research\_Notes\_V157\_f1.tex}.
\end{center}
Immediately before the rollover it had \(3{,}623{,}989\) bytes and SHA--256
digest
\begin{center}
\small\texttt{EA08CE7537C7327A80B5AA5A653111AFAEDAE82CC18215D87534F17B563AAD94}.
\end{center}
The suffix \texttt{f1} records the first frozen volume.  This file starts a new
active sequence at V1.  When an active sequence reaches V100 it will be frozen
with the next suffix \texttt{f2}, \texttt{f3}, and so on, and the active
sequence will again start at V1.  Thus version numbers remain short while the
hash and frozen-volume suffix retain an auditable ancestry.

The supplied monograph and programme files are mathematical inputs.  Their
definitions, conjectures, and calculations are not operational instructions,
and statements from them are not promoted to theorems without proof in this
series.  Likewise, the summary below records what was established in the first
volume under its stated hypotheses; it does not strengthen those hypotheses.

\subsection{The continuum statement being pursued}

For orientation, let \(\Lambda_n\uparrow\infty\) be a cofinal cutoff sequence,
let \(\mathfrak A_n\) be the finite-cutoff BV observable complex, and let
\(\omega_n\) be the renormalized state produced by the shell pushforward.  A
complete result would require, at minimum, the following compatible objects.
\begin{enumerate}[label=\textup{(C\arabic*)}]
\item A continuum complex \(\mathfrak A_\infty\) and comparison maps for which
the finite BV pushforwards form a coherent inverse or inductive system.
\item A positive physical state on the BRST cohomology, obtained as the limit
of \(\omega_n\), with all required local composite fields defined.
\item A renormalized effective action satisfying the quantum master equation,
including boundary and determinant-phase terms.
\item Convergence of the metric derivative to a distributional tensor
\(T_{\mu\nu}\) with the correct Ward identity and enough compactness to exclude
uncontrolled concentration defects.
\item Existence, uniqueness in a fixed gauge, constraint propagation, and
cutoff stability for the Einstein--matter evolution sourced by that limit.
\end{enumerate}
The previous volume supplies exact finite-cutoff identities and several
conditional limit mechanisms.  Conditions \(C2\)--\(C5\) are not yet established
in the generality needed for the full conclusion.

\section{Major results of the first volume}
\label{sec:v2-legacy}

This section is a compressed map of the first 157 versions.  ``Exact'' below
means exact in the finite-dimensional or regularized setting explicitly used
there.  ``Conditional'' means that the conclusion is a theorem once the named
analytic estimate or compactness hypothesis is supplied.

\subsection{Moving supports, boundary geometry, and common ports}

The early part of the first volume separated scalar likelihood transport from
transport of a moving returned-word support.  A Kato connection was constructed
for the support projector.  Exact rank-one examples showed that the scalar
likelihood plaquette can be flat while the support connection has nonzero
curvature.  The associated polar-chord holonomy was computed to second order,
so support curvature became an independent obstruction rather than a proxy for
ordinary scalar curvature.

The boundary analysis then placed gauge, Higgs, fermion, ghost, and metric data
on a common timelike slab.  It fixed the normal and Brown--York/Codazzi sign
conventions, constructed compatible incoming and outgoing ports, and derived
the boundary flux terms in the gauge and diffeomorphism Ward identities.  A
separate Hadamard analysis exhibited why an arbitrary boundary projector need
not preserve the microlocal spectrum condition.  Consequently, boundary
domain compatibility and Hadamard compatibility remain separate tests.

\subsection{BRST/BV structure and determinant phase}

At finite cutoff, the gauge-fixed field content was completed by the
off-shell BRST quartet and a common boundary form domain.  The naive Hessian
restriction was shown to leave a precise gauge defect.  Dual field and
antifield projectors were then used to formulate a BV-compatible restriction;
the defect of a projector that is not symplectic was computed explicitly.
Finite-dimensional BV pushforward was proved to preserve the quantum master
equation and to compose under successive shell integrations.  The effective
action is therefore a connected-cumulant object, not the restriction of the
bare action.

For determinant phases, the regularized comparison maps were placed in the
correct Schatten class.  In three boundary dimensions an order \(-1\)
perturbation lies only in \(\mathcal S_p\) for \(p>3\), so subtraction through
symbol order \(-3\) is required before an ordinary determinant-tail argument
is available.  The resulting \(\det_4\) multiplicative anomaly was identified
as a local finite-part transgression.  On a compact slab, a boundary-family
cobordism made the relevant \(K^1\) index vanish.  With the stated spectral
section and parallel trivialization, the Bismut--Freed line was flat with
trivial holonomy and the compared chiral phases were one.  This removes that
phase obstruction under those hypotheses; it does not prove convergence of
the determinant modulus.

\subsection{The Higgs-improved Einstein source}

The coefficient \(\xi_H\) of the nonminimal coupling was kept independent of
the Einstein coefficient.  Writing \(f=H^\dagger H\), the improvement term in
the conventions of the first volume was
\begin{equation}
 T_{\xi}^{\mu\nu}
 =-2\xi_H\left[fG^{\mu\nu}
 +(g^{\mu\nu}\Box-\nabla^\mu\nabla^\nu)f\right],
 \qquad
 \nabla_\mu T_{\xi}^{\mu}{}_{\nu}
 =\xi_H R\nabla_\nu f.
 \label{eq:v2-improved-stress}
\end{equation}
The scalar--tensor Gibbons--Hawking partner was included, and the face momentum
was found to be
\begin{equation}
 \mathcal P_{gH}^{ab}
 =(\chi_R+\xi_H f)\Pi^{ab}
 -\xi_H\gamma^{ab}\nabla_nf.
 \label{eq:v2-face-momentum}
\end{equation}
Together with the antifield-source term, this yielded the BV-completed one-loop
diffeomorphism Ward identity.  Any numerical relation between \(\xi_H\) and a
gravitational coefficient is therefore an additional renormalization
condition, not a consequence of covariance alone.

The Einstein-frame transformation was also diagonalized.  For
\(F=\chi_R+\xi_Hf>0\), \(\widetilde g=(F/\chi_*)g\), the scalar target metric is
\begin{equation}
 \mathcal G_{AB}
 =\frac{\chi_*}{F}G^0_{AB}
 +\frac{3\chi_*}{F^2}\,\partial_AF\,\partial_BF.
 \label{eq:v2-target-metric}
\end{equation}
This calculation isolated the necessary coercivity condition.  The subsequent
hyperbolic analysis showed why weak source convergence is insufficient for
cutoff stability: high-frequency curvature can survive weak convergence.  In
harmonic gauge the subsidiary field obeys
\begin{equation}
 \Box_g C_\nu+R_\nu{}^\mu C_\mu=-2\kappa\mathcal R_\nu,
 \label{eq:v2-subsidiary}
\end{equation}
where \(\mathcal R_\nu\) is the residual Ward defect.  Thus constraint
propagation in the limit requires quantitative control of that defect.

\subsection{Compactness and concentration}

The first volume proved the usual strong--weak product mechanism above the
quadratic endpoint: suitable \(L^p\) compactness with \(p>2\), combined with
weak convergence of the other factor, closes the nonlinear source.  It also
gave exact concentrating sequences at \(p=2\), proving that bounded energy at
the endpoint can leave a nonzero defect measure.

For the Higgs sector, global quartic Gibbs tails were separated from spatial
uniform integrability.  A bound on the distribution of the global quartic
energy does not by itself prevent that energy from concentrating on shrinking
sets.  A local Orlicz or equivalent uniform-integrability estimate does.  In
the absence of such an estimate the limiting stress can contain
\begin{equation}
 \mathfrak S_H
 =\mathfrak K-\lambda_H g\,\nu_H-2\xi_H\mathfrak I,
 \label{eq:v2-defect-stress}
\end{equation}
where the terms respectively encode kinetic, quartic, and improvement defects.
Conservation of \(\mathfrak S_H\) alone does not force it to vanish.  Uniform
monotonicity removes the defect in the coercive Euclidean branch, giving a
precise upstream estimate to seek rather than a circular appeal to the desired
Einstein equation.

\subsection{Shell pushforward and the current frontier}

The classical restriction space, the quantum state space, and their comparison
maps were distinguished.  Compatible finite-dimensional BV states give a
complete shell telescope.  For a Laplace-type Gaussian determinant, subtraction
through the boundary heat coefficient of degree four leaves a proper-time
shell of order \(O(\Lambda_n^{-1})\); on a closed four-manifold the corresponding
remainder is \(O(\Lambda_n^{-2})\).  The unsubtracted degree-four coefficient is
marginal and produces a logarithm.  This is an exact reason to accumulate local
terms as running couplings and to sum only the irrelevant complement.

The first interacting derivative was reduced to
\begin{equation}
 \delta_g\langle I\rangle_0[h]
 =\langle\delta_g I[h]\rangle_0
 -\hbar^{-1}\operatorname{Cov}_0
       \bigl(I,\delta_gS_0[h]\bigr).
 \label{eq:v2-metric-derivative}
\end{equation}
Raw marginal local shells in the covariance are not summable.  If
\(\mathbf P_{\mathrm{loc}}\) is the invariant finite-dimensional local-jet
projector, the correct shell is
\begin{equation}
 \mathfrak s_n^{\mathrm{ren}}
 =(I-\mathbf P_{\mathrm{loc}})\mathfrak s_n.
 \label{eq:v2-ren-shell}
\end{equation}
The projector must commute with the linearized Ward/Slavnov operator.  A
parallel-edge regression in the Yang--Mills programme showed that the number
of differentiated algebraic factors can exceed the number of independent
random bridges: lower-rank pieces cancel only after the complete connected
quantity is assembled.  The present note begins exactly there.

\subsection{Status ledger}

\begin{center}
\small
\begin{tabular}{>{\raggedright\arraybackslash}p{0.31\textwidth}
                >{\raggedright\arraybackslash}p{0.57\textwidth}}
\toprule
Component & Status after the first volume \\
\midrule
Common boundary geometry and sign conventions
& Exact at finite cutoff on the stated slab and domains. \\
BV shell composition and master identity
& Exact in the finite-dimensional BV model, conditional in the continuum on
uniform analytic bounds and domain convergence. \\
Chiral determinant phase
& Trivialized under the compact-slab \(K^1\)-vanishing, spectral-section, and
parallel-frame hypotheses. \\
Determinant modulus
& Gaussian irrelevant shells are summable after the required local
subtractions; the complete interacting modulus remains open. \\
Higgs and gauge source compactness
& Strong--weak and Orlicz closure criteria are proved; the uniform estimates
for the complete interacting state remain open. \\
Einstein evolution
& Constraint propagation and stability criteria are identified; the needed
strong curvature and source convergence is not yet proved. \\
Full SM--Einstein continuum
& Open. \\
\bottomrule
\end{tabular}
\end{center}

\section{Exact shell covariance by martingale differences}
\label{sec:v2-martingale}

We now work at a fixed finite cutoff.  This is essential: all conditional
expectations below are ordinary finite-dimensional objects, and no continuum
measure is presumed.

\begin{definition}[Shell filtration]
Let \((\Omega,\mathcal F,\mathbb P)\) carry square-integrable random variables
\(X_1,\ldots,X_N\) that generate \(\mathcal F\) modulo null sets.  Set

\[
 \mathcal F_i=\sigma(X_1,\ldots,X_i),\qquad
 \mathcal F_0=\{\varnothing,\Omega\}.
\]
For \(F\in L^2(\Omega)\), define
\[
 F_i=\mathbb E[F\mid\mathcal F_i],\qquad
 d_iF=F_i-F_{i-1}.
\]
No independence hypothesis is needed for this definition or for
\cref{thm:v2-doob-covariance} below.
\end{definition}

\begin{theorem}[Exact connected covariance decomposition]
\label{thm:v2-doob-covariance}
For \(F,G\in L^2(\Omega)\),
\begin{equation}
 \operatorname{Cov}(F,G)
 =\sum_{i=1}^N\mathbb E[d_iF\,d_iG].
 \label{eq:v2-doob-covariance}
\end{equation}
In particular, the index on the two factors is the same: a connected covariance
uses one filtration increment at a time.
\end{theorem}

\begin{proof}
Because the coordinates generate \(\mathcal F\), one has \(\mathcal F_N=\mathcal F\), and
\[
 F-\mathbb E F=\sum_{i=1}^N d_iF,
 \qquad
 G-\mathbb E G=\sum_{j=1}^N d_jG.
\]
If \(i<j\), then \(d_iF\) is \(\mathcal F_{j-1}\)-measurable and
\(
 \mathbb E[d_jG\mid\mathcal F_{j-1}]=0.
\)
Consequently,
\[
 \mathbb E[d_iF\,d_jG]
 =\mathbb E\!\left[
 d_iF\,\mathbb E[d_jG\mid\mathcal F_{j-1}]
 \right]=0.
\]
The case \(i>j\) follows by conditioning \(d_iF\) on
\(\mathcal F_{i-1}\).  Expanding the product of the two finite sums leaves only
the terms \(i=j\), which proves \eqref{eq:v2-doob-covariance}.
\end{proof}

\begin{corollary}[Square-function identity and covariance bound]
\label{cor:v2-square-function}
For real \(F\in L^2\),
\[
 \operatorname{Var}(F)=\sum_{i=1}^N\|d_iF\|_2^2.
\]
Moreover,
\begin{equation}
 |\operatorname{Cov}(F,G)|
 \le
 \left(\sum_i\|d_iF\|_2^2\right)^{1/2}
 \left(\sum_i\|d_iG\|_2^2\right)^{1/2}.
 \label{eq:v2-square-function-bound}
\end{equation}
\end{corollary}

\begin{proof}
Take \(G=F\) in \cref{thm:v2-doob-covariance} for the first assertion.  Apply
Cauchy--Schwarz to each expectation in
\eqref{eq:v2-doob-covariance}, followed by Cauchy--Schwarz in the finite index
\(i\), for the second.
\end{proof}

\subsection{The coordinate-resampling formula}

Suppose now that \(X_1,\ldots,X_N\) are independent.  On an enlarged product
space let \(X_i'\) be an independent copy of \(X_i\), independent of all the
original coordinates, and put
\[
 X^{(i)}=(X_1,\ldots,X_{i-1},X_i',X_{i+1},\ldots,X_N),
 \qquad
 \Delta_iF=F(X)-F(X^{(i)}).
\]

\begin{theorem}[Derivative-faithful resampling]
\label{thm:v2-resampling}
For a product reference law and \(F\in L^2\),
\begin{equation}
 d_iF
 =\mathbb E[\Delta_iF\mid\mathcal F_i],
 \label{eq:v2-resampling-difference}
\end{equation}
where the conditional expectation integrates the future coordinates and the
independent copy \(X_i'\).  Hence
\begin{equation}
 \operatorname{Cov}(F,G)
 =\sum_{i=1}^N
 \mathbb E\!\left[
  \mathbb E[\Delta_iF\mid\mathcal F_i]
  \mathbb E[\Delta_iG\mid\mathcal F_i]
 \right].
 \label{eq:v2-resampling-covariance}
\end{equation}
\end{theorem}

\begin{proof}
The first term in the conditional expectation of \(\Delta_iF\) is
\(\mathbb E[F(X)\mid\mathcal F_i]=F_i\).  In the second term, \(X_i'\) and the
future variables \(X_{i+1},\ldots,X_N\) are integrated independently of
\(X_1,\ldots,X_i\).  Its conditional expectation therefore depends only on
\(X_1,\ldots,X_{i-1}\) and equals \(F_{i-1}\).  This proves
\eqref{eq:v2-resampling-difference}.  Substitution into
\eqref{eq:v2-doob-covariance} proves
\eqref{eq:v2-resampling-covariance}.
\end{proof}

\begin{proposition}[One-coordinate regression test]
\label{prop:v2-one-coordinate}
Let \(\xi\) be a symmetric sign, let \(s>0\), and set
\(F=G=\sqrt{s}\,\xi\).  Then the unique nonzero martingale difference has
\(L^2\)-norm \(\sqrt{s}\), while
\[
 \operatorname{Cov}(F,G)=s.
\]
Thus a covariance weight \(s\) comes from the product of two amplitudes
\(\sqrt{s}\) attached to the same resampled coordinate.  Treating the two
appearances as two independent covariance bridges would incorrectly produce
\(s^2\).
\end{proposition}

\begin{proof}
For the one-step filtration, \(d_1F=F-\mathbb E F=\sqrt{s}\xi=d_1G\).
Therefore
\(\mathbb E[d_1F\,d_1G]=s\mathbb E\xi^2=s\).  There is only one filtration
index, so no second independent bridge exists.
\end{proof}

\begin{remark}[Gaussian shell coordinates]
At finite cutoff, a centered Gaussian field on the quotient by its covariance
kernel can be written \(X=C^{1/2}\xi\), where the components of \(\xi\) are
independent standard Gaussians.  The resampling formula can therefore be
applied after whitening and after the gauge/BV splitting has identified the
nondegenerate Gaussian directions.  The Doob formula itself remains valid for
an interacting law, but the independent-copy representation
\eqref{eq:v2-resampling-difference} does not survive an arbitrary Gibbs tilt.
In a perturbative shell pushforward the tilt may instead be retained inside the
observable under the product Gaussian reference law.
\end{remark}

\section{Vector-valued connected sources}
\label{sec:v2-vector-source}

The covariance in \eqref{eq:v2-metric-derivative} takes values in a finite
cutoff space of metric-source coefficients.  The following form avoids any
choice of individual diagrams.

\begin{theorem}[Bilinear martingale contraction]
\label{thm:v2-bilinear}
Let \(E_1,E_2,Y\) be finite-dimensional normed spaces, let
\(B:E_1\times E_2\to Y\) be bounded bilinear, and let
\(F\in L^2(\Omega;E_1)\), \(G\in L^2(\Omega;E_2)\).  Define
\[
 \mathcal C_B(F,G)
 =\mathbb E\,B(F-\mathbb E F,G-\mathbb E G).
\]
Then
\begin{equation}
 \mathcal C_B(F,G)
 =\sum_{i=1}^N\mathbb E\,B(d_iF,d_iG).
 \label{eq:v2-bilinear-identity}
\end{equation}
If \(Q:Y\to Y\) is bounded, then
\begin{equation}
 \|Q\mathcal C_B(F,G)\|_Y
 \le \|Q\circ B\|
 \left(\sum_i\mathbb E\|d_iF\|_{E_1}^2\right)^{1/2}
 \left(\sum_i\mathbb E\|d_iG\|_{E_2}^2\right)^{1/2}.
 \label{eq:v2-bilinear-bound}
\end{equation}
\end{theorem}

\begin{proof}
Expand both centered variables into their finite martingale sums.  For \(i<j\),
bounded bilinearity and finite dimensionality allow conditional expectation to
pass through the second slot:
\[
 \mathbb E\,B(d_iF,d_jG)
 =\mathbb E\,B\bigl(d_iF,
   \mathbb E[d_jG\mid\mathcal F_{j-1}]\bigr)=0.
\]
The case \(i>j\) is identical with the slots reversed.  This proves
\eqref{eq:v2-bilinear-identity}.  For the bound, use the triangle inequality,
the definition of \(\|Q\circ B\|\), Cauchy--Schwarz in probability, and then
Cauchy--Schwarz over \(i\):
\begin{align*}
 \|Q\mathcal C_B(F,G)\|_Y
 &\le \|Q\circ B\|\sum_i
    \mathbb E\bigl[\|d_iF\|_{E_1}\|d_iG\|_{E_2}\bigr]\\
 &\le \|Q\circ B\|\sum_i
    (\mathbb E\|d_iF\|_{E_1}^2)^{1/2}
    (\mathbb E\|d_iG\|_{E_2}^2)^{1/2}\\
 &\le \|Q\circ B\|
 \left(\sum_i\mathbb E\|d_iF\|_{E_1}^2\right)^{1/2}
 \left(\sum_i\mathbb E\|d_iG\|_{E_2}^2\right)^{1/2}.
\end{align*}
\end{proof}

For \eqref{eq:v2-metric-derivative}, take \(F=I\), take
\(G=\delta_gS_0[h]\), and let \(B\) be the finite-cutoff contraction producing
the chosen local metric-source coefficient.  The theorem gives an exact sum
over actual shell-coordinate increments.  It does not yet give ultraviolet
decay: that must come from the projected contraction or from quantitative
square-function estimates.  This distinction is the first firewall against
the parallel-edge error.

\begin{counterexample}[Coordinate counting does not renormalize a coincidence]
\label{sep:v2-local-divergence}
For every shell \(n\), let \(F_n=G_n=\xi_n\) with \(\xi_n\) a symmetric sign.
Then the martingale representation is exact and
\(\operatorname{Cov}(F_n,G_n)=1\), but
\(\sum_n\operatorname{Cov}(F_n,G_n)\) diverges.  If this constant is the local
zero-momentum jet, the projector \(Q_0C=C-C(0)\) removes it.  Thus the
martingale identity prevents false bridge powers; the local projector performs
the distinct renormalization step.
\end{counterexample}

\section{Local Taylor projection and a summability criterion}
\label{sec:v2-taylor}

Let \(K\) be a finite-dimensional external-momentum space, \(Y\) a
finite-dimensional source space, and let \(C_n:B_\kappa(0)\subset K\to Y\) be
the fully assembled connected shell source.  Define its local jet of degree
\(r\) and its complement by
\[
 (P_rC_n)(k)=\sum_{j=0}^r\frac1{j!}D^jC_n(0)[k^{\otimes j}],
 \qquad Q_r=I-P_r.
\]

\begin{theorem}[Projected connected-shell bound]
\label{thm:v2-projected-shell}
Assume \(C_n\in C^{r+1}(B_\kappa(0);Y)\).  Then
\begin{equation}
 \|(Q_rC_n)(k)\|_Y
 \le \frac{|k|^{r+1}}{(r+1)!}
 \sup_{0\le t\le1}\|D^{r+1}C_n(tk)\|_{\mathrm{op}}.
 \label{eq:v2-taylor-remainder}
\end{equation}
Suppose in addition that

\[
 C_n(k)=\mathcal C_B(F_n(k),G_n(k))
\]
with the hypotheses of \cref{thm:v2-bilinear}, differentiation commuting with
conditional expectation, and
\begin{align*}
 A_{j,n}&=
 \sup_{|q|\le\kappa}
 \left(\sum_i\mathbb E
 \|D^jd_iF_n(q)\|_{\mathrm{op}}^2\right)^{1/2},\\
 B_{j,n}&=
 \sup_{|q|\le\kappa}
 \left(\sum_i\mathbb E
 \|D^jd_iG_n(q)\|_{\mathrm{op}}^2\right)^{1/2}.
\end{align*}
Then
\begin{equation}
 \|(Q_rC_n)(k)\|_Y
 \le
 \frac{\|B\|\,|k|^{r+1}}{(r+1)!}
 \sum_{j=0}^{r+1}{r+1\choose j}A_{j,n}B_{r+1-j,n}.
 \label{eq:v2-projected-square-function}
\end{equation}
If \(\Lambda_n=2^n\Lambda_0\) and, for \(0\le j\le r+1\),
\[
 A_{j,n}\le a_j\Lambda_n^{\sigma_F-j},
 \qquad
 B_{j,n}\le b_j\Lambda_n^{\sigma_G-j},
\]
then
\begin{equation}
 \|(Q_rC_n)(k)\|_Y
 \le C_r|k|^{r+1}
 \Lambda_n^{\sigma_F+\sigma_G-(r+1)}.
 \label{eq:v2-dyadic-power}
\end{equation}
Consequently, the irrelevant shell series converges absolutely whenever
\(r+1>\sigma_F+\sigma_G\).
\end{theorem}

\begin{proof}
Apply one-variable Taylor's theorem to the curve
\(c(t)=C_n(tk)\):
\[
 (Q_rC_n)(k)
 =\frac1{r!}\int_0^1(1-t)^r
 D^{r+1}C_n(tk)[k^{\otimes(r+1)}]\,\mathrm dt.
\]
Since \(
\int_0^1(1-t)^r\,\mathrm dt=1/(r+1)
\), this proves \eqref{eq:v2-taylor-remainder}.

Differentiate \eqref{eq:v2-bilinear-identity} \(r+1\) times.  The Leibniz rule
gives, for unit directions suppressed from the notation,
\[
 D^{r+1}C_n(q)
 =\sum_{j=0}^{r+1}{r+1\choose j}
   \sum_i\mathbb E\,
   B\bigl(D^jd_iF_n(q),D^{r+1-j}d_iG_n(q)\bigr).
\]
The proof of \eqref{eq:v2-bilinear-bound}, term by term in \(j\), bounds this by
\(
 \|B\|\sum_j{r+1\choose j}A_{j,n}B_{r+1-j,n}.
\)
Together with \eqref{eq:v2-taylor-remainder} this proves
\eqref{eq:v2-projected-square-function}.  Under the scale hypotheses, every
product has the same power
\[
 \Lambda_n^{\sigma_F-j}
 \Lambda_n^{\sigma_G-(r+1-j)}
 =\Lambda_n^{\sigma_F+\sigma_G-(r+1)},
\]
which proves \eqref{eq:v2-dyadic-power}.  The final assertion is the convergence
of a geometric series with negative exponent.
\end{proof}

\begin{remark}[The marginal equality]
If \(r+1=\sigma_F+\sigma_G\), the theorem gives a shell of order one and no
summability.  It is therefore invalid to turn the two factors in
\eqref{eq:v2-projected-square-function} into two independent covariance
weights.  Each is an \(L^2\) amplitude, and their product is the single
covariance weight exhibited by \cref{prop:v2-one-coordinate}.
\end{remark}

\section{A concrete four-dimensional shell calculation}
\label{sec:v2-bubble}

The preceding criterion is conditional.  The following elementary loop is a
regression model in which the local divergence and the projected gain can both
be proved directly.  It also tests coincidence before any graph-by-graph norm
estimate is attempted.

Fix \(0<a<b<\infty\), let
\(\chi\in C_c^\infty(\mathbb R^4)\) be even and supported in
\(\{u:a\le|u|\le b\}\), and put
\(\chi_n(q)=\chi(q/\Lambda_n)\).  For \(m\ge0\) and \(|k|<a\Lambda_n\), define
\begin{equation}
 C_n(k)=\int_{\mathbb R^4}
 \frac{\chi_n(q)\,\mathrm d^4q}
 {(q^2+m^2)((q+k)^2+m^2)}.
 \label{eq:v2-bubble}
\end{equation}

\begin{theorem}[Marginal local jet and summable complement]
\label{thm:v2-bubble}
There is a constant \(C=C(a,b,\chi)\), independent of \(n\), \(m\ge0\), and
\(k\), such that whenever \(|k|\le a\Lambda_n/2\),
\begin{align}
 |C_n(0)|&\le C,\label{eq:v2-bubble-local}\\
 DC_n(0)&=0,\label{eq:v2-bubble-odd}\\
 |C_n(k)-C_n(0)|&\le C|k|^2\Lambda_n^{-2}.
 \label{eq:v2-bubble-remainder}
\end{align}
For every fixed \(k\), the high-shell series
\(\sum_n(C_n(k)-C_n(0))\) therefore converges absolutely.  If \(m=0\) and
\(\chi\ge0\) is not zero, then
\begin{equation}
 C_n(0)=\int_{\mathbb R^4}\frac{\chi(u)}{|u|^4}\,\mathrm d^4u>0
 \label{eq:v2-bubble-exact-marginal}
\end{equation}
for every \(n\), so the unsubtracted local series diverges.
\end{theorem}

\begin{proof}
On the support of \(\chi_n\), \(|q|\ge a\Lambda_n\).  Therefore
\[
 \frac1{q^2+m^2}\le\frac1{a^2\Lambda_n^2}.
\]
The support has volume \(O(\Lambda_n^4)\), and the second denominator at
\(k=0\) supplies another \(O(\Lambda_n^{-2})\).  This proves
\eqref{eq:v2-bubble-local}.

Let \(g_m(p)=(p^2+m^2)^{-1}\).  Differentiation under the integral is justified
because the integration region is compact and stays away from \(q=0\).  Since

\[
 Dg_m(q)[v]=-\frac{2q\cdot v}{(q^2+m^2)^2},
\]
the integrand in \(DC_n(0)[v]\) is odd in \(q\): the factors
\(\chi_n(q)\) and \(g_m(q)\) are even, while \(Dg_m(q)[v]\) is odd.  Its integral
vanishes, proving \eqref{eq:v2-bubble-odd}.

If \(|k|\le a\Lambda_n/2\) and \(0\le t\le1\), then
\(|q+tk|\ge a\Lambda_n/2\).  Direct differentiation gives
\[
 D^2g_m(p)[v,w]
 =-\frac{2v\cdot w}{(p^2+m^2)^2}
 +\frac{8(p\cdot v)(p\cdot w)}{(p^2+m^2)^3}.
\]
For \(|p|\ge a\Lambda_n/2\), its operator norm is bounded by
\(C_a\Lambda_n^{-4}\), uniformly in \(m\ge0\).  Hence
\begin{align*}
 \|D^2C_n(tk)\|_{\mathrm{op}}
 &\le \int |\chi_n(q)|g_m(q)
       \|D^2g_m(q+tk)\|_{\mathrm{op}}\,\mathrm d^4q\\
 &\le C\Lambda_n^4\Lambda_n^{-2}\Lambda_n^{-4}
 =C\Lambda_n^{-2}.
\end{align*}
Taylor's theorem, together with \(DC_n(0)=0\), now yields
\[
 |C_n(k)-C_n(0)|
 \le\frac12|k|^2
       \sup_{0\le t\le1}\|D^2C_n(tk)\|_{\mathrm{op}},
\]
which is \eqref{eq:v2-bubble-remainder}.  The dyadic series
\(\sum_n\Lambda_n^{-2}\) converges.  Finally, if \(m=0\), the substitution
\(q=\Lambda_n u\) cancels the four powers from the measure against the four
powers in the two propagators and gives
\eqref{eq:v2-bubble-exact-marginal}.  Positivity follows from the assumptions on
\(\chi\).
\end{proof}

\begin{assessment}[What the bubble proves and what it does not]
\label{ass:v2-bubble-scope}
The estimate proves a real two-power gain for this even, logarithmically
marginal four-dimensional shell after subtraction of its constant local jet.
The first derivative vanishes by parity; the gain was obtained by
differentiating the assembled kernel.  It was not obtained by assigning a gain
to each displayed propagator.  Tensor numerators, chiral projectors, boundary
symbols, and overlapping subgraphs change the superficial local degree and
must be treated by the corresponding invariant forest or local-jet projector.
The theorem is therefore a validated regression case, not a proof of the full
Standard Model estimate.
\end{assessment}

\section{Ward compatibility of the local projector}
\label{sec:v2-ward-projector}

Let \(\mathscr W:Y\to Y\) denote the finite-cutoff linearized combined
Ward/Slavnov operator on the local source space.  The following elementary
condition is indispensable because a noninvariant subtraction can manufacture
a constraint defect even when the unprojected shell is invariant.

\begin{proposition}[Invariant splitting]
\label{prop:v2-invariant-splitting}
Let \(P:Y\to Y\) be a projector and \(Q=I-P\).  If
\(\mathscr WP=P\mathscr W\), then \(\mathscr WQ=Q\mathscr W\).  In particular,
if \(\mathscr WC_n=0\), then
\[
 \mathscr W(PC_n)=0,
 \qquad
 \mathscr W(QC_n)=0.
\]
If \(C_n\to C\) after summing the \(Q\)-pieces in a topology in which
\(\mathscr W\) is closed, then \(\mathscr WC=0\).
\end{proposition}

\begin{proof}
The identity
\(\mathscr WQ=\mathscr W(I-P)=\mathscr W-P\mathscr W=Q\mathscr W\)
proves the first statement.  Apply it and the assumed commutation once with
\(Q\) and once with \(P\).  The last statement is the definition of closedness
of \(\mathscr W\) applied to the convergent sequence of partial sums.
\end{proof}

The projector in \cref{prop:v2-invariant-splitting} is applied to the assembled
connected covariance.  Projecting each martingale summand separately is
legitimate only if that finer operation itself commutes with \(\mathscr W\) and
has uniform bounds.  The exact identity
\eqref{eq:v2-bilinear-identity} gives no such termwise symmetry for free.

\section{Consequences for the programme and the next attack}
\label{sec:v2-next}

Combining the results of this note yields the following finite-cutoff route for
the covariance term in \eqref{eq:v2-metric-derivative}.
\begin{enumerate}[label=\textup{(S\arabic*)}]
\item Choose the shell filtration after quotienting null Gaussian directions
and performing the BV-compatible field/antifield splitting.
\item Express the complete connected metric-source covariance by
\eqref{eq:v2-bilinear-identity}.  Count each common martingale index once.
\item Assemble all contractions required by the master identity before taking
norms.  This retains cancellations between algebraic parallel edges.
\item Extract the invariant local jet \(P_rC_n\) into the running gravitational,
gauge, Yukawa, Higgs, boundary, and antifield-source couplings.
\item Prove the square-function derivative bounds in
\cref{thm:v2-projected-shell} for \(Q_rC_n\).  Sum only this irrelevant part.
\item Pass the Ward identity through the limit using
\cref{prop:v2-invariant-splitting}, and then address uniform integrability of
the resulting metric source before invoking Einstein stability.
\end{enumerate}

The new reduction is noncircular.  Neither the martingale identity nor the
Taylor estimate assumes existence of the continuum source or of the limiting
Einstein metric.  The next note should work upstream on the missing quantitative
input: for a concrete gauge-fixed BV kinetic operator on a compact slab,
construct the whitened high-mode blocks and prove external-momentum derivative
square-function estimates uniform in the cutoff.  The scalar calculation in
\cref{thm:v2-bubble} supplies the first mandatory regression.  The same proof
must then be repeated for the transverse gauge, ghost, fermion, Higgs, and
metric insertions with a common invariant local basis.  A failure in any sector
will identify whether the obstruction is power counting, boundary symbol
regularity, or loss of Ward-compatible cancellation.

\section{Version record}
\label{sec:v2-record}

This second-volume V1 contains:
\begin{enumerate}
\item the context and frozen-file lineage for the V157 rollover;
\item a theorem-level summary and status ledger for the first volume;
\item the exact Doob covariance and independent-coordinate resampling
identities;
\item a one-coordinate test forbidding a false second bridge weight;
\item a Banach-valued projected covariance estimate;
\item a local Taylor-jet summability theorem;
\item an explicit four-dimensional marginal bubble whose projected remainder
is \(O(|k|^2\Lambda_n^{-2})\); and
\item the Ward-compatible route to the next shell estimate.
\end{enumerate}

The next active version will append to this file and replace it, while this V1
content remains its exact prefix.  The first cross-program audit of the new
volume is due at V5.

\clearpage
\section*{Second Volume, Version V2: Whitened Gaussian Shells and Exact Square Functions}
\addcontentsline{toc}{section}{Second Volume, Version V2: Whitened Gaussian Shells and Exact Square Functions}

\subsection*{Purpose and relation to V1}

Second-volume V1 reduced the first interacting metric-source covariance to
martingale square functions and a local Taylor remainder.  That reduction left
the square-function estimates as a hypothesis.  The present version proves
those estimates exactly for finite-dimensional bosonic Gaussian shells with
quadratic insertions.  After whitening, the entire square function is a
Hilbert--Schmidt norm; it does not carry a factor equal to the number of chosen
coordinates or blocks.  A Weyl rank bound then gives a cutoff power for
spectral shells.  This supplies a rigorous bosonic regression theorem for the
next Standard--Model estimates.

\subsection{Quadratic Gaussian martingales}
\label{sec:v2v2-quadratic}

Let \(H=\bigoplus_{i=1}^N H_i\) be a finite-dimensional real Hilbert space.
Let \(\xi_i\) be independent standard Gaussian vectors in \(H_i\), put
\(\xi=(\xi_1,\ldots,\xi_N)\), and use the filtration
\(\mathcal F_i=\sigma(\xi_1,\ldots,\xi_i)\).  If \(A=A^*\in\mathcal L(H)\),
write \(A_{ij}:H_j\to H_i\) for its blocks and define the centered quadratic
observable
\begin{equation}
 \mathcal Q_A(\xi)
 =\frac12\bigl(\langle\xi,A\xi\rangle-\operatorname{Tr}A\bigr).
 \label{eq:v2v2-quadratic-observable}
\end{equation}

\begin{theorem}[Exact block martingale formula]
\label{thm:v2v2-block-martingale}
For the observable in \eqref{eq:v2v2-quadratic-observable},
\begin{equation}
 d_i\mathcal Q_A
 =\frac12\left(
   \langle\xi_i,A_{ii}\xi_i\rangle-\operatorname{Tr}A_{ii}
  \right)
 +\sum_{j<i}\langle\xi_i,A_{ij}\xi_j\rangle .
 \label{eq:v2v2-block-difference}
\end{equation}
Moreover,
\begin{equation}
 \mathbb E|d_i\mathcal Q_A|^2
 =\frac12\|A_{ii}\|_{\mathrm{HS}}^2
  +\sum_{j<i}\|A_{ij}\|_{\mathrm{HS}}^2,
 \label{eq:v2v2-block-variance}
\end{equation}
and therefore
\begin{equation}
 \sum_{i=1}^N\mathbb E|d_i\mathcal Q_A|^2
 =\frac12\|A\|_{\mathrm{HS}}^2.
 \label{eq:v2v2-square-hs}
\end{equation}
The right side is independent of the number, dimensions, and ordering of the
blocks.
\end{theorem}

\begin{proof}
Conditional expectation of a cross term containing an unrevealed Gaussian is
zero.  Conditional expectation of an unrevealed diagonal term is its trace.
Consequently,
\[
 \mathbb E[\mathcal Q_A\mid\mathcal F_i]
 =\frac12\sum_{a\le i}
 \bigl(\langle\xi_a,A_{aa}\xi_a\rangle-\operatorname{Tr}A_{aa}\bigr)
 +\sum_{a<b\le i}\langle\xi_b,A_{ba}\xi_a\rangle .
\]
Subtracting the same expression with \(i-1\) in place of \(i\) proves
\eqref{eq:v2v2-block-difference}.

For a standard Gaussian vector \(\eta\) and a self-adjoint matrix \(D\),
\[
 \mathbb E\bigl(\langle\eta,D\eta\rangle-\operatorname{Tr}D\bigr)^2
 =2\operatorname{Tr}(D^2).
\]
This follows from
\(\mathbb E[\eta_a\eta_b\eta_c\eta_d]
=\delta_{ab}\delta_{cd}+\delta_{ac}\delta_{bd}
 +\delta_{ad}\delta_{bc}\).
The diagonal term in \eqref{eq:v2v2-block-difference} therefore contributes
\(\frac12\|A_{ii}\|_{\mathrm{HS}}^2\).  Distinct cross terms are orthogonal
because they contain distinct centered vectors \(\xi_j\), and
\[
 \mathbb E|\langle\xi_i,A_{ij}\xi_j\rangle|^2
 =\operatorname{Tr}(A_{ij}^*A_{ij})
 =\|A_{ij}\|_{\mathrm{HS}}^2.
\]
Every cross term is also orthogonal to the centered diagonal term by vanishing
of odd Gaussian moments.  This proves \eqref{eq:v2v2-block-variance}.
Finally, self-adjointness gives
\[
 \|A\|_{\mathrm{HS}}^2
 =\sum_i\|A_{ii}\|_{\mathrm{HS}}^2
  +2\sum_{j<i}\|A_{ij}\|_{\mathrm{HS}}^2.
\]
Summing \eqref{eq:v2v2-block-variance} proves
\eqref{eq:v2v2-square-hs}.
\end{proof}

\begin{corollary}[Exact covariance and its sharp Hilbert--Schmidt bound]
\label{cor:v2v2-quadratic-covariance}
For self-adjoint \(A,B\in\mathcal L(H)\),
\begin{equation}
 \operatorname{Cov}(\mathcal Q_A,\mathcal Q_B)
 =\frac12\operatorname{Tr}(AB)
 =\sum_i\mathbb E[
    d_i\mathcal Q_A\,d_i\mathcal Q_B],
 \label{eq:v2v2-quadratic-covariance}
\end{equation}
and
\begin{equation}
 |\operatorname{Cov}(\mathcal Q_A,\mathcal Q_B)|
 \le\frac12\|A\|_{\mathrm{HS}}\|B\|_{\mathrm{HS}}.
 \label{eq:v2v2-quadratic-cs}
\end{equation}
\end{corollary}

\begin{proof}
The Gaussian fourth-moment identity used in
\cref{thm:v2v2-block-martingale} gives
\[
 \mathbb E[
 (\langle\xi,A\xi\rangle-\operatorname{Tr}A)
 (\langle\xi,B\xi\rangle-\operatorname{Tr}B)]
 =2\operatorname{Tr}(AB).
\]
The normalization in \eqref{eq:v2v2-quadratic-observable} gives the first
equality in \eqref{eq:v2v2-quadratic-covariance}; the martingale equality is
\cref{thm:v2-doob-covariance}.  The final estimate is Hilbert--Schmidt
Cauchy--Schwarz.
\end{proof}

\begin{remark}[What block refinement can and cannot change]
Refining one \(H_i\) into many coordinate blocks changes the individual
martingale differences in \eqref{eq:v2v2-block-difference}, but
\eqref{eq:v2v2-square-hs} is invariant.  Thus a proof that grows a factor with
the number of whitened coordinates has discarded Gaussian orthogonality.  The
correct quantity is the Hilbert--Schmidt norm of the assembled insertion.
\end{remark}

\subsection{External derivatives before local projection}
\label{sec:v2v2-external}

Let \(K\) be a finite-dimensional external-parameter space and let
\(A,B:U\subset K\to\mathcal L(H)\) be \(C^m\) self-adjoint families.  Set
\[
 C(k)=\operatorname{Cov}(\mathcal Q_{A(k)},\mathcal Q_{B(k)})
 =\frac12\operatorname{Tr}(A(k)B(k)).
\]

\begin{proposition}[Derivative square functions are operator derivatives]
\label{prop:v2v2-derivative-square}
For every \(0\le j\le m\),
\begin{equation}
 \sup_{\substack{|v_1|\le1,\ldots,|v_j|\le1}}
 \left(\sum_i\mathbb E
 \left|D^jd_i\mathcal Q_{A(k)}
 [v_1,\ldots,v_j]\right|^2\right)^{1/2}
 =\frac1{\sqrt2}\|D^jA(k)\|_{\mathrm{HS},\mathrm{op}}.
 \label{eq:v2v2-derivative-square}
\end{equation}
Here the norm on the right is the operator norm of the \(j\)-linear map from
\(K^j\) to the Hilbert--Schmidt class.  Consequently,
\begin{equation}
 \|D^mC(k)\|_{\mathrm{op}}
 \le\frac12\sum_{j=0}^m{m\choose j}
 \|D^jA(k)\|_{\mathrm{HS},\mathrm{op}}
 \|D^{m-j}B(k)\|_{\mathrm{HS},\mathrm{op}}.
 \label{eq:v2v2-covariance-derivative}
\end{equation}
\end{proposition}

\begin{proof}
Conditional expectation is a fixed linear map independent of \(k\), so
\(D^jd_i\mathcal Q_{A(k)}=d_i\mathcal Q_{D^jA(k)}\) after evaluation on any
fixed \(j\) directions.  Apply \eqref{eq:v2v2-square-hs} for those directions
and then take the supremum over unit directions.  This proves
\eqref{eq:v2v2-derivative-square}.  Alternatively differentiating
\(\frac12\operatorname{Tr}(A(k)B(k))\), applying the multilinear Leibniz rule,
and using Hilbert--Schmidt Cauchy--Schwarz proves
\eqref{eq:v2v2-covariance-derivative}.
\end{proof}

\subsection{Whitening a bosonic spectral shell}
\label{sec:v2v2-whitening}

Let \(L\) be a positive self-adjoint elliptic realization of order two on a
compact \(d\)-dimensional manifold, with a fixed elliptic local boundary
condition when a boundary is present.  For constants \(0<c_0<c_1\), let
\[
 E_n=\mathbf 1_{[c_0\Lambda_n^2,c_1\Lambda_n^2]}(L),
 \qquad
 C_n=E_n(L+m^2)^{-1}E_n .
\]
We require only the Weyl upper bound
\begin{equation}
 \operatorname{rank}E_n\le W\Lambda_n^d.
 \label{eq:v2v2-weyl-upper}
\end{equation}
Let \(\xi_n\) be standard Gaussian on \(E_nH\) and
\(\phi_n=C_n^{1/2}\xi_n\).  For a self-adjoint insertion \(M_n(k)\) on
\(E_nH\), define
\begin{equation}
 \mathcal O_{M,n}(k)
 =\frac12\left(
 \langle\phi_n,M_n(k)\phi_n\rangle
 -\operatorname{Tr}(C_nM_n(k))\right).
 \label{eq:v2v2-physical-quadratic}
\end{equation}
After whitening, this is \(\mathcal Q_{A_n(k)}\) with
\begin{equation}
 A_n(k)=C_n^{1/2}M_n(k)C_n^{1/2}.
 \label{eq:v2v2-whitened-insertion}
\end{equation}

\begin{theorem}[Weyl-to-square-function estimate]
\label{thm:v2v2-weyl-square}
Assume that, for \(0\le j\le r+1\),
\begin{equation}
 \sup_{|k|\le\kappa}
 \|D^jM_n(k)\|_{\mathrm{op}}
 \le M_j\Lambda_n^{\mu-j}.
 \label{eq:v2v2-insertion-order}
\end{equation}
Then
\begin{equation}
 \sup_{|k|\le\kappa}
 \|D^jA_n(k)\|_{\mathrm{HS},\mathrm{op}}
 \le \frac{\sqrt W}{c_0}M_j
 \Lambda_n^{d/2+\mu-2-j}.
 \label{eq:v2v2-whitened-hs}
\end{equation}
For a second insertion \(N_n(k)\) satisfying the analogous estimate with
order \(\nu\), define
\[
 C_n^{M,N}(k)
 =\operatorname{Cov}(\mathcal O_{M,n}(k),\mathcal O_{N,n}(k)).
\]
If \(P_r\) is the degree-\(r\) Taylor jet at \(k=0\) and \(Q_r=I-P_r\), then
\begin{equation}
 |(Q_rC_n^{M,N})(k)|
 \le K_r|k|^{r+1}
 \Lambda_n^{d+\mu+\nu-4-(r+1)}
 \label{eq:v2v2-spectral-remainder}
\end{equation}
for \(|k|\le\kappa\), with \(K_r\) independent of \(n\).  Hence the projected
dyadic shell series converges absolutely if
\begin{equation}
 r+1>d+\mu+\nu-4.
 \label{eq:v2v2-spectral-threshold}
\end{equation}
\end{theorem}

\begin{proof}
On the range of \(E_n\),
\[
 \|C_n^{1/2}\|_{\mathrm{op}}^2
 =\|C_n\|_{\mathrm{op}}
 \le(c_0\Lambda_n^2+m^2)^{-1}
 \le c_0^{-1}\Lambda_n^{-2}.
\]
For an operator \(T\) on a space of rank \(R\),
\(\|T\|_{\mathrm{HS}}\le\sqrt R\|T\|_{\mathrm{op}}\).  Since \(C_n\) is
independent of the external parameter,
\[
 D^jA_n(k)=C_n^{1/2}D^jM_n(k)C_n^{1/2}.
\]
Combining this identity with \eqref{eq:v2v2-weyl-upper} and
\eqref{eq:v2v2-insertion-order} proves
\eqref{eq:v2v2-whitened-hs}.

Apply \eqref{eq:v2v2-covariance-derivative} with \(m=r+1\).  Each summand has
the power
\[
 \Lambda_n^{d/2+\mu-2-j}
 \Lambda_n^{d/2+\nu-2-(r+1-j)}
 =\Lambda_n^{d+\mu+\nu-4-(r+1)}.
\]
The integral Taylor remainder from \cref{thm:v2-projected-shell} then proves
\eqref{eq:v2v2-spectral-remainder}.  Under
\eqref{eq:v2v2-spectral-threshold}, its cutoff exponent is negative, and the
dyadic series is geometric.
\end{proof}

\begin{corollary}[Four-dimensional zero-order bosonic insertions]
\label{cor:v2v2-four-dimensional}
If \(d=4\) and \(\mu=\nu=0\), the raw quadratic covariance is marginal.  A
constant-jet subtraction gives
\[
 |(Q_0C_n^{M,N})(k)|\le K|k|\Lambda_n^{-1},
\]
and the high-shell series is absolutely summable.  If the first external
derivative vanishes by parity and the insertion estimates also hold through
second derivative, as in \cref{thm:v2-bubble}, the sharper bound is
\(O(|k|^2\Lambda_n^{-2})\).
\end{corollary}

\begin{proof}
Use \(r=0\) in \eqref{eq:v2v2-spectral-remainder}.  If the first derivative
vanishes, use the second-order Taylor remainder and the \(j=2\) version of
\eqref{eq:v2v2-whitened-hs}.
\end{proof}

\subsection{Sector firewall}
\label{sec:v2v2-firewall}

\begin{firewall}[Scope of the Gaussian theorem]
\label{fw:v2v2-scope}
\cref{thm:v2v2-weyl-square} applies directly to a positive Euclidean bosonic
Gaussian shell after zero modes have been quotiented.  It covers the scalar
and Higgs reference shells and the positive part of a gauge-fixed vector
shell.  Fermion and Faddeev--Popov ghost integrations are Berezin integrals,
not probability measures, so the \(L^2\) martingale positivity used in
\eqref{eq:v2v2-square-hs} does not apply to them.  Their corresponding estimates
must be proved by trace ideals and resolvent identities.  A Euclidean graviton
shell additionally requires a specified contour for the conformal direction.
No positivity estimate for that direction is inferred here.
\end{firewall}

The remaining analytic input in the bosonic sectors is now localized.  One
must prove \eqref{eq:v2v2-insertion-order} for the actual cutoff comparison
maps and boundary pseudodifferential realizations, with a common local basis.
For metric variations, \(M_n\) may have positive differential order; then
\eqref{eq:v2v2-spectral-threshold} gives the exact Taylor degree that must be
absorbed into local gravitational and boundary couplings.

\subsection{V2 conclusion and next step}
\label{sec:v2v2-conclusion}

This version proves that the bosonic martingale square function has no hidden
coordinate-count loss:
\[
 \left(\sum_i\mathbb E|d_i\mathcal Q_A|^2\right)^{1/2}
 =2^{-1/2}\|A\|_{\mathrm{HS}}.
\]
Whitening and the Weyl upper bound then convert external derivative order into
the cutoff exponent
\[
 d+\mu+\nu-4-(r+1)
\]
for a degree-\(r\) projected covariance.  The result is still sector-specific.
The next upstream obstruction is algebraic: an invariant local polynomial
subspace need not admit a projector commuting with a nilpotent Slavnov
differential.  The next version should determine the exact splitting
obstruction and replace an unjustified commuting-projector assumption by
chain-retraction or homotopy data.

\clearpage
\section*{Second Volume, Version V3: The Slavnov-Compatible Local Splitting Obstruction}
\addcontentsline{toc}{section}{Second Volume, Version V3: The Slavnov-Compatible Local Splitting Obstruction}

\subsection*{Why the projector assumption must be proved}

Second-volume V1 required the local Taylor projector to commute with the
linearized Ward/Slavnov operator.  V2 reduced the bosonic analytic estimate to
Hilbert--Schmidt bounds, but did not justify that algebraic commutation.  The
present version shows that it is not automatic.  An invariant local subspace
can fail to be a chain direct summand.  We identify the exact finite-dimensional
obstruction and prove a local repair theorem: the subtraction can still be
made Slavnov closed precisely when its induced ghost-number-one defect is
exact in the local complex.

\subsection{Chain projections and local subcomplexes}
\label{sec:v2v3-chain-projection}

Let \(\Bbbk\) be a field of characteristic zero.  A finite cochain complex is a
finite-dimensional graded vector space
\[
 V=\bigoplus_q V^q
\]
with a degree-one map \(s:V^q\to V^{q+1}\) satisfying \(s^2=0\).  Let
\(L\subset V\) be a graded subspace with \(sL\subset L\).  In the application,
\(V\) is a finite-cutoff amplitude space, \(L\) is the space of admissible
local Taylor jets, \(q\) is ghost number, and \(s\) is the linearized combined
Slavnov operator.

\begin{definition}[Slavnov-compatible local projection]
\label{def:v2v3-chain-projection}
A linear map \(P:V\to V\) is a Slavnov-compatible projection onto \(L\) if
\[
 P^2=P,\qquad \operatorname{Ran}P=L,\qquad sP=Ps.
\]
Equivalently, if \(\iota:L\hookrightarrow V\) is inclusion, then
\(p=P|_V:V\to L\) is a chain map satisfying \(p\iota=I_L\), and
\(P=\iota p\).
\end{definition}

\begin{theorem}[Exact splitting criterion]
\label{thm:v2v3-splitting-criterion}
For a finite-dimensional complex \(V\) over \(\Bbbk\) and a subcomplex
\(L\subset V\), the following are equivalent.
\begin{enumerate}[label=\textup{(\roman*)}]
\item There is a Slavnov-compatible projection \(P:V\to L\).
\item The inclusion \(\iota:L\hookrightarrow V\) has a chain retraction.
\item The short exact sequence of complexes
\[
 0\longrightarrow L\overset{\iota}{\longrightarrow}V
 \longrightarrow V/L\longrightarrow0
\]
splits in the category of cochain complexes.
\item Every connecting homomorphism
\[
 \delta_q:H^q(V/L,s)\longrightarrow H^{q+1}(L,s)
\]
in the long exact cohomology sequence vanishes.
\item The induced map \(H^q(\iota):H^q(L,s)\to H^q(V,s)\) is injective for
every \(q\).
\end{enumerate}
\end{theorem}

\begin{proof}
The equivalence of (i) and (ii) follows from
\cref{def:v2v3-chain-projection}.  A chain retraction gives
\[
 V=L\oplus\ker p
\]
as complexes, proving (iii); conversely, projection along a subcomplex
complement is a chain retraction.

If the short exact sequence splits as complexes, its long exact cohomology
sequence is the direct sum of the cohomology sequences of \(L\) and \(V/L\),
so every connecting map is zero.  Exactness of the long cohomology sequence
gives
\[
 \ker H^{q+1}(\iota)=\operatorname{Ran}\delta_q.
\]
Thus (iv) and (v) are equivalent.

It remains to prove that (iv) implies (iii).  Choose a degree-preserving vector
space splitting \(V=L\oplus Q\).  Relative to it, the differential has block
form
\begin{equation}
 s_V=
 \begin{pmatrix}
  s_L&a\\
  0&s_Q
 \end{pmatrix},
 \qquad
 a:Q^q\longrightarrow L^{q+1}.
 \label{eq:v2v3-extension-matrix}
\end{equation}
The identity \(s_V^2=0\) says
\begin{equation}
 s_La+as_Q=0,
 \label{eq:v2v3-a-chain}
\end{equation}
so \(a:Q\to L[1]\) is a chain map.  Its induced cohomology map is exactly the
connecting homomorphism \(\delta\).

We use the elementary splitting lemma for complexes of vector spaces: a chain
map between finite-dimensional complexes that induces zero on cohomology is
null-homotopic.  To verify it, decompose each complex as
\[
 Z^q=B^q\oplus H^q,\qquad
 V^q=B^q\oplus H^q\oplus C^q,
\]
where \(Z^q=\ker s\), \(B^q=\operatorname{Ran}s\), and
\(s:C^q\to B^{q+1}\) is an isomorphism.  On the \(H\)-summands the chain map
has zero cohomology component by hypothesis; its remaining components factor
through the contractible pairs \(C^q\to B^{q+1}\).  Composing with the inverse
of \(s:C^q\to B^{q+1}\) on those pairs gives a degree-minus-one map \(h\) with
the chain map equal to \(s h+h s\).  This proves the lemma.

Since \(\delta=0\), apply the lemma to \(a:Q\to L[1]\).  With the grading shift
unpacked, there is a degree-zero map \(h:Q^q\to L^q\) such that
\begin{equation}
 a+s_Lh-hs_Q=0.
 \label{eq:v2v3-kill-extension}
\end{equation}
Replace the chosen lift \(q\mapsto(0,q)\) by
\(q\mapsto(hq,q)\).  By \eqref{eq:v2v3-extension-matrix}, the \(L\)-component
of its differential relative to the new lift is the left side of
\eqref{eq:v2v3-kill-extension}, hence zero.  The new copy of \(Q\) is an
\(s_V\)-invariant complement of \(L\).  Projection onto \(L\) along that
complement proves (iii).
\end{proof}

\begin{remark}[Sign check in the extension matrix]
If the new lift is \(j_h(q)=(h q,q)\), then
\[
 s_Vj_h(q)=(s_Lh q+a q,s_Qq),\qquad
 j_hs_Q(q)=(h s_Qq,s_Qq).
\]
Thus \(j_h\) is a chain map exactly when
\(a+s_Lh-hs_Q=0\), confirming the sign in
\eqref{eq:v2v3-kill-extension}.
\end{remark}

\subsection{An exact two-dimensional obstruction}
\label{sec:v2v3-separator}

\begin{counterexample}[An invariant local space with no commuting projection]
\label{sep:v2v3-jordan}
Let
\[
 V^0=\Bbbk x,\qquad V^1=\Bbbk y,\qquad sx=y,\qquad sy=0,
\]
and let \(L^0=0\), \(L^1=\Bbbk y\).  Then \(L\) is an \(s\)-invariant
subcomplex, but no projection onto \(L\) commutes with \(s\).
\end{counterexample}

\begin{proof}
Any projection \(P\) onto \(L\) satisfies \(Px=0\) by degree and \(Py=y\)
because it restricts to the identity on its range.  Hence
\[
 sPx=0,\qquad Psx=Py=y,
\]
so \(sP\ne Ps\).  Cohomologically, \(y\) represents a nonzero class in
\(H^1(L,s)\), while \(y=sx\) is zero in \(H^1(V,s)\).  Thus
\(H^1(L)\to H^1(V)\) is not injective, in agreement with
\cref{thm:v2v3-splitting-criterion}.
\end{proof}

The separator has the exact shape relevant to local subtraction.  A local
closed expression \(y\) can have a primitive \(x\) in the full amplitude space
without having a local primitive.  In that case, an arbitrary Taylor
projection cannot be improved to a commuting projection while retaining the
same local range.

\subsection{Repairing a noncommuting Taylor subtraction}
\label{sec:v2v3-repair}

An exact chain projection is convenient but stronger than necessary for one
closed shell.  The next theorem gives the precise weaker condition.

\begin{theorem}[Local Slavnov repair for one shell]
\label{thm:v2v3-local-repair}
Let \(P:V\to L\) be any degree-preserving projection onto an \(s\)-invariant
subcomplex \(L\), and put \(Q=I-P\).  Let \(C\in V^0\) satisfy \(sC=0\).
Define the local subtraction defect
\begin{equation}
 \alpha(C)=s(QC)=-s(PC)\in L^1.
 \label{eq:v2v3-local-defect}
\end{equation}
Then \(s\alpha(C)=0\).  The following are equivalent.
\begin{enumerate}[label=\textup{(\alph*)}]
\item There is \(b\in L^0\) such that the two pieces
\[
 C_{\mathrm{loc}}=PC+b,\qquad
 C_{\mathrm{irr}}=QC-b
\]
are separately \(s\)-closed.
\item The class \([\alpha(C)]\) vanishes in \(H^1(L,s)\).
\end{enumerate}
When these conditions hold, \(C=C_{\mathrm{loc}}+C_{\mathrm{irr}}\) exactly;
the repair changes only the allocation between local running couplings and the
irrelevant shell.
\end{theorem}

\begin{proof}
Because \(s^2=0\),
\[
 s\alpha(C)=s^2(QC)=0.
\]
If \([\alpha(C)]=0\), choose \(b\in L^0\) with \(sb=\alpha(C)\).  Then
\[
 sC_{\mathrm{irr}}=s(QC)-sb=\alpha(C)-\alpha(C)=0.
\]
Since \(s(PC)=-\alpha(C)\),
\[
 sC_{\mathrm{loc}}=s(PC)+sb=-\alpha(C)+\alpha(C)=0.
\]
Their sum is \(PC+QC=C\).

Conversely, if \(C_{\mathrm{irr}}=QC-b\) is closed for some \(b\in L^0\), then
\[
 0=s(QC-b)=\alpha(C)-sb,
\]
so \(\alpha(C)=sb\) and its local cohomology class vanishes.
\end{proof}

\begin{corollary}[Uniform repair criterion]
\label{cor:v2v3-uniform-repair}
If \(H^1(L,s)=0\), every \(s\)-closed shell admits the decomposition in
\cref{thm:v2v3-local-repair}.  More generally, it is enough that the specific
classes \([\alpha(C_n)]\) vanish and that primitives \(b_n\) can be chosen with
the same cutoff bounds as the local coefficients \(PC_n\).
\end{corollary}

\begin{proof}
Vanishing of \(H^1(L,s)\) makes every local degree-one cocycle exact.  The
second assertion is \cref{thm:v2v3-local-repair} shell by shell; the uniform
bound is an additional analytic requirement needed to sum or renormalize the
chosen primitives.
\end{proof}

\begin{firewall}[Cohomological vanishing is not a norm estimate]
\label{fw:v2v3-cohomology-norm}
The existence of \(b_n\) says nothing by itself about its size.  Even in finite
dimensions, a right inverse for \(s:L^0\to\operatorname{Ran}s\subset L^1\)
can have a norm that grows with the cutoff.  A continuum proof therefore needs
both the algebraic condition
\([\alpha(C_n)]=0\) and a uniformly controlled local contracting homotopy.
\end{firewall}

\subsection{Relation to Standard--Model anomaly cancellation}
\label{sec:v2v3-sm-anomaly}

The first volume treated perturbative Standard--Model anomaly cancellation and
the determinant-line phase obstruction.  Those results constrain the class in
\cref{thm:v2v3-local-repair}, but they do not automatically verify its present
relative form.  The local complex here also contains diffeomorphism, boundary,
antifield-source, nonminimal Higgs, and metric counterterms.  What is required
is the ghost-number-one cohomology of that complete local jet complex together
with the map induced by inclusion into the finite-cutoff amplitude complex.

There are therefore two honest algebraic routes.
\begin{enumerate}[label=\textup{(A\arabic*)}]
\item Prove that \(H^\bullet(L)\to H^\bullet(V)\) is injective and construct the
chain retraction furnished by \cref{thm:v2v3-splitting-criterion}.
\item Use a convenient degree-preserving Taylor projector, compute the actual
local cocycle \(\alpha(C_n)\), and construct a uniformly bounded primitive
\(b_n\) as in \cref{thm:v2v3-local-repair}.
\end{enumerate}
The second route is weaker and is sufficient shell by shell.  It also exposes
an anomaly immediately: a nonzero class \([\alpha(C_n)]\in H^1(L,s)\) is
exactly the obstruction to a local Ward-preserving allocation.

\subsection{V3 conclusion and next step}
\label{sec:v2v3-conclusion}

The assumption \(sP_{\mathrm{loc}}=P_{\mathrm{loc}}s\) has now been replaced by
a theorem with a checkable obstruction.  A commuting projector exists exactly
when the local subcomplex is a chain direct summand, equivalently when its
cohomology injects into the full finite-cutoff cohomology.  For a fixed shell,
the weaker and necessary condition is the vanishing of the induced local
ghost-number-one defect class, plus a uniform bound on its primitive.

The next analytic sector left outside V2 is the Berezin Gaussian sector.
Accordingly, the next version should derive the exact connected covariance of
fermionic quadratic insertions, identify the Schatten norms that replace
probabilistic square functions, and calculate the spectral cutoff exponent
without importing positivity from the bosonic argument.

\clearpage
\section*{Second Volume, Version V4: Berezin Shell Covariance and the Dirac Trace-Ideal Threshold}
\addcontentsline{toc}{section}{Second Volume, Version V4: Berezin Shell Covariance and the Dirac Trace-Ideal Threshold}

\subsection*{Purpose}

V2 proved exact Hilbert--Schmidt square functions for positive bosonic Gaussian
shells and explicitly excluded fermions and Faddeev--Popov ghosts.  Their
finite-cutoff integrals are algebraic Berezin functionals, so neither
probabilistic conditional expectation nor a positive variance is available.
This version proves the corresponding connected quadratic identity directly
from Grassmann Wick contraction.  It then derives the Schatten bound and the
spectral cutoff threshold for a first-order Dirac shell.  The result supplies
the fermionic analytic input to the local Slavnov splitting problem of V3.

\subsection{Finite-dimensional Grassmann Gaussian}
\label{sec:v2v4-grassmann}

Let \(H\cong\mathbb C^R\), let \(D:H\to H\) be invertible, and write
\(C=D^{-1}\).  Let \(\psi_a,\bar\psi_b\) be independent Grassmann generators.
Fix the normalized Berezin Gaussian convention by
\begin{equation}
 \langle\psi_a\bar\psi_b\rangle_D=C_{ab}
 \label{eq:v2v4-contraction}
\end{equation}
and the fermionic Wick rule.  Equivalently, the unnormalized partition
function is \(\det D\) with the compatible orientation of the Berezin measure.
For \(M\in\mathcal L(H)\), set
\begin{equation}
 \mathcal F_M
 =-\bar\psi M\psi-\operatorname{Tr}(CM).
 \label{eq:v2v4-centered-bilinear}
\end{equation}

\begin{lemma}[Centering and the connected sign]
\label{lem:v2v4-centered}
The observable in \eqref{eq:v2v4-centered-bilinear} is centered:
\[
 \langle\mathcal F_M\rangle_D=0.
\]
For \(M,N\in\mathcal L(H)\),
\begin{equation}
 \langle\mathcal F_M\mathcal F_N\rangle_D
 =-\operatorname{Tr}(CMCN).
 \label{eq:v2v4-connected-trace}
\end{equation}
\end{lemma}

\begin{proof}
From \eqref{eq:v2v4-contraction} and anticommutation,
\[
 \langle\bar\psi_a\psi_b\rangle_D=-C_{ba}.
\]
Therefore
\[
 \langle-\bar\psi M\psi\rangle_D
 =\sum_{a,b}M_{ab}C_{ba}=\operatorname{Tr}(MC),
\]
which proves centering.

For the product of two uncentered bilinears, the fermionic Wick rule gives
\begin{align*}
 &\langle
  (\bar\psi_aM_{ab}\psi_b)
  (\bar\psi_cN_{cd}\psi_d)
 \rangle_D\\
 &\quad
 =\operatorname{Tr}(CM)\operatorname{Tr}(CN)
  -\operatorname{Tr}(CMCN).
\end{align*}
The first term is the disconnected contraction.  Subtracting the two means,
as done in \eqref{eq:v2v4-centered-bilinear}, leaves the second term with its
fermion-loop minus sign.
\end{proof}

\begin{counterexample}[No positive variance]
\label{sep:v2v4-negative-variance}
In one Grassmann dimension, take \(D=d\ne0\) and \(M=1\).  Then
\[
 \langle\mathcal F_1^2\rangle_D=-d^{-2}.
\]
Thus the bosonic identity
\(\operatorname{Var}F=\sum_i\|d_iF\|_2^2\) cannot be transferred to the
Berezin sector.
\end{counterexample}

\begin{proof}
Apply \eqref{eq:v2v4-connected-trace} with \(C=d^{-1}\) and \(M=N=1\).
\end{proof}

\subsection{The trace-ideal replacement}
\label{sec:v2v4-trace-ideal}

\begin{theorem}[Fermionic connected-covariance bound]
\label{thm:v2v4-trace-bound}
If \(CM\) and \(CN\) are Hilbert--Schmidt, then
\begin{equation}
 |\langle\mathcal F_M\mathcal F_N\rangle_D|
 \le\|CM\|_{\mathrm{HS}}\|CN\|_{\mathrm{HS}}.
 \label{eq:v2v4-trace-bound}
\end{equation}
More generally, if \(CM\in\mathcal S_p\), \(CN\in\mathcal S_q\), and
\(p^{-1}+q^{-1}=1\), then
\begin{equation}
 |\langle\mathcal F_M\mathcal F_N\rangle_D|
 \le\|CM\|_{\mathcal S_p}\|CN\|_{\mathcal S_q}.
 \label{eq:v2v4-holder}
\end{equation}
\end{theorem}

\begin{proof}
By \eqref{eq:v2v4-connected-trace},
\[
 |\langle\mathcal F_M\mathcal F_N\rangle_D|
 =|\operatorname{Tr}((CM)(CN))|.
\]
For the Hilbert--Schmidt case, regard this as the Hilbert--Schmidt inner
product of \((CM)^*\) and \(CN\), and use Cauchy--Schwarz.  The general
statement is Hölder's inequality for Schatten ideals.
\end{proof}

\begin{remark}[The sign and the norm play different roles]
The loop minus sign in \eqref{eq:v2v4-connected-trace} is essential in the
assembled Slavnov identity and in cancellations among species.  It disappears
from \eqref{eq:v2v4-trace-bound} because the latter is an absolute-value
estimate.  Norms must therefore be taken only after anomaly and Ward
cancellations that depend on the relative signs have been assembled.
\end{remark}

\subsection{External-parameter derivatives}
\label{sec:v2v4-derivatives}

Let \(M,N:U\subset K\to\mathcal L(H)\) be \(C^m\) families while \(D\), hence
\(C\), is fixed on the shell.  Define
\[
 \Gamma(k)=-\operatorname{Tr}(CM(k)CN(k)).
\]

\begin{proposition}[Derivative trace bound]
\label{prop:v2v4-derivative-trace}
For \(m\ge0\),
\begin{equation}
 \|D^m\Gamma(k)\|_{\mathrm{op}}
 \le
 \sum_{j=0}^m{m\choose j}
 \|C D^jM(k)\|_{\mathrm{HS},\mathrm{op}}
 \|C D^{m-j}N(k)\|_{\mathrm{HS},\mathrm{op}}.
 \label{eq:v2v4-derivative-trace}
\end{equation}
\end{proposition}

\begin{proof}
Apply the multilinear Leibniz rule to
\(-\operatorname{Tr}(CM(k)CN(k))\).  For each allocation of the \(m\)
directions, use \eqref{eq:v2v4-trace-bound}; there are
\({m\choose j}\) allocations with \(j\) derivatives on \(M\).  Taking the
supremum over unit directions gives \eqref{eq:v2v4-derivative-trace}.
\end{proof}

\subsection{A first-order elliptic shell}
\label{sec:v2v4-dirac-shell}

Let \(\mathcal D\) be an invertible self-adjoint first-order elliptic Dirac
realization on a compact \(d\)-dimensional Riemannian manifold.  If there is a
boundary, fix a self-adjoint elliptic boundary realization compatible with the
spectral section used for the determinant phase.  Put
\[
 E_n=\mathbf1_{[c_0\Lambda_n,c_1\Lambda_n]}(|\mathcal D|),
 \qquad
 C_n=E_n\mathcal D^{-1}E_n,
\]
and assume the Weyl upper bound
\begin{equation}
 \operatorname{rank}E_n\le W\Lambda_n^d.
 \label{eq:v2v4-dirac-weyl}
\end{equation}
All insertions below are compressed to \(E_nH\).

\begin{theorem}[Dirac shell threshold]
\label{thm:v2v4-dirac-threshold}
Suppose \(M_n,N_n:U\to\mathcal L(E_nH)\) obey, for
\(0\le j\le r+1\),
\begin{align}
 \sup_{|k|\le\kappa}\|D^jM_n(k)\|_{\mathrm{op}}
 &\le M_j\Lambda_n^{\mu-j},\label{eq:v2v4-M-order}\\
 \sup_{|k|\le\kappa}\|D^jN_n(k)\|_{\mathrm{op}}
 &\le N_j\Lambda_n^{\nu-j}.\label{eq:v2v4-N-order}
\end{align}
Let
\begin{equation}
 \Gamma_n(k)=-\operatorname{Tr}
  (C_nM_n(k)C_nN_n(k)).
 \label{eq:v2v4-shell-covariance}
\end{equation}
Then
\begin{equation}
 \|C_nD^jM_n(k)\|_{\mathrm{HS},\mathrm{op}}
 \le \frac{\sqrt W}{c_0}M_j
 \Lambda_n^{d/2+\mu-1-j},
 \label{eq:v2v4-CM-hs}
\end{equation}
and the degree-\(r\) Taylor complement satisfies
\begin{equation}
 |(Q_r\Gamma_n)(k)|
 \le K_r|k|^{r+1}
 \Lambda_n^{d+\mu+\nu-2-(r+1)}.
 \label{eq:v2v4-dirac-remainder}
\end{equation}
The projected dyadic shell series is absolutely convergent when
\begin{equation}
 r+1>d+\mu+\nu-2.
 \label{eq:v2v4-dirac-condition}
\end{equation}
\end{theorem}

\begin{proof}
On \(E_nH\),
\[
 \|C_n\|_{\mathrm{op}}\le(c_0\Lambda_n)^{-1}.
\]
The rank bound and
\(\|T\|_{\mathrm{HS}}\le\sqrt{\operatorname{rank}E_n}\|T\|_{\mathrm{op}}\)
give
\begin{align*}
 \|C_nD^jM_n(k)\|_{\mathrm{HS},\mathrm{op}}
 &\le
 \sqrt{W}\Lambda_n^{d/2}
 (c_0\Lambda_n)^{-1}
 M_j\Lambda_n^{\mu-j},
\end{align*}
which is \eqref{eq:v2v4-CM-hs}.

Use \cref{prop:v2v4-derivative-trace} with \(m=r+1\).  Every term has cutoff
power
\[
 \Lambda_n^{d/2+\mu-1-j}
 \Lambda_n^{d/2+\nu-1-(r+1-j)}
 =\Lambda_n^{d+\mu+\nu-2-(r+1)}.
\]
The integral Taylor remainder of \cref{thm:v2-projected-shell} proves
\eqref{eq:v2v4-dirac-remainder}.  If
\eqref{eq:v2v4-dirac-condition} holds, the power is negative and the dyadic
series converges.
\end{proof}

\begin{corollary}[Four-dimensional zero-order insertions]
\label{cor:v2v4-four-dimensional}
For \(d=4\) and \(\mu=\nu=0\), the absolute trace estimate gives
\[
 |\Gamma_n(k)|=O(\Lambda_n^2).
\]
Absolute summability of the Taylor complement follows from
\cref{thm:v2v4-dirac-threshold} once \(r\ge2\).  Thus a bare absolute-value
bound alone requires subtraction through quadratic external momentum.
\end{corollary}

\begin{proof}
The raw exponent in \eqref{eq:v2v4-dirac-remainder} before differentiation is
\(d+\mu+\nu-2=2\).  The condition \(r+1>2\) is equivalent to \(r\ge2\).
\end{proof}

\begin{assessment}[Why this is an upper bound, not the final gauge degree]
\label{ass:v2v4-ward-improvement}
For a gauge-current two-point function, the assembled Ward identity can forbid
a gauge-boson mass term and impose transversality.  Charge conjugation or
parity can also remove odd jets.  Those cancellations may lower the actual
local degree from the raw value in \cref{cor:v2v4-four-dimensional}.  They must
be established before absolute values are taken.  The theorem deliberately
states the unconditional trace-ideal threshold; it does not assume the desired
Ward cancellation.
\end{assessment}

\subsection{Zero modes, phase, and ghosts}
\label{sec:v2v4-zero-modes}

The inverse \(C_n\) is defined only away from zero modes.  Kernel and cokernel
data belong to the determinant line and were treated separately in the first
volume.  Freezing a spectral section and quotienting the kernel are therefore
part of the hypothesis of \cref{thm:v2v4-dirac-threshold}; the trace estimate
does not trivialize the determinant phase.

Faddeev--Popov ghosts obey the same algebraic Wick sign as complex fermionic
Grassmann pairs, with their own kinetic order and bundle rank.  Hence
\cref{lem:v2v4-centered,thm:v2v4-trace-bound} applies at finite cutoff.
Power counting must use the ghost kinetic operator: for a second-order ghost
operator, each inverse has order \(-2\), so the spectral exponent agrees with
the bosonic order-two formula rather than the first-order Dirac formula.
Positivity is still neither assumed nor needed.

\subsection{V4 conclusion and next step}
\label{sec:v2v4-conclusion}

The fermionic connected quadratic shell is now controlled by the exact formula
\[
 \langle\mathcal F_M\mathcal F_N\rangle_D
 =-\operatorname{Tr}(CMCN)
\]
and by Schatten Hölder, without a fictitious positive variance.  For a
first-order Dirac shell the raw cutoff degree is
\[
 d+\mu+\nu-2,
\]
and every external derivative lowers that degree by one under the insertion
hypotheses.  The sign-sensitive Ward sum must be assembled before applying the
absolute trace bound.

The next active version is V5, so it must begin with the required cross-program
audit of the current Navier--Stokes and Yang--Mills notes.  Mathematically, that
audit should be used to decide whether to attack the sign-sensitive transverse
fermion sum or the uniformly bounded local Slavnov primitive identified in V3.

\clearpage
\section*{Second Volume, Version V5: Companion Audit and the Sign-Sensitive Fermion Ward Row}
\addcontentsline{toc}{section}{Second Volume, Version V5: Companion Audit and the Sign-Sensitive Fermion Ward Row}

\subsection{Mandatory companion audit}
\label{sec:v2v5-audit}

\begin{assessment}[Read-only V5 audit]
\label{ass:v2v5-audit}
The newest active companion notes found at the start of this version were:
\begin{center}
\small
\begin{tabular}{llll}
\toprule
Programme & Version & Bytes & SHA--256\\
\midrule
Navier--Stokes & V31 & \(887537\) &
\texttt{F1121B62238AD5491BB7CF03635EA9934942EDF573ED8FAAA9EE79E1D38A6696}\\
Yang--Mills & compact V5 & \(76844\) &
\texttt{FB047187FBC744BA95F76A65C5E0D0C671C22D2EBF9D1E659AB5F1C32643A164}\\
\bottomrule
\end{tabular}
\end{center}
NS V31 proves that flat probability marking does not remove the weighted
first moment needed to recover critical annular action.  It also exhibits a
balanced-helicity nullspace for the normalized signed hinge.  YM compact V5
computes the complete Cartan fourth-cumulant Möbius polynomial and shows that
its multiplicativity cancellation is destroyed if its partition terms are
bounded separately.
\end{assessment}

\begin{firewall}[Audit transfer boundary]
\label{fw:v2v5-audit}
No Navier--Stokes estimate or Yang--Mills capacity theorem is imported.  Two
proof-order lessons are used.  First, a projection or probability
normalization does not supply a missing cutoff weight.  Second, a
sign-sensitive symmetry row must be assembled before absolute values are
taken.  In the present fermionic problem, the complete object is the connected
trace together with the contact insertion and, in a chiral theory, the
anomaly-cancelling species sum.
\end{firewall}

\subsection{The finite-matrix determinant Ward identity}
\label{sec:v2v5-determinant-ward}

Let \(H\) be a finite-dimensional complex vector space, let \(D\in GL(H)\),
put \(C=D^{-1}\), and define the fermionic effective functional
\begin{equation}
 W(D)=-\log\det D
\label{eq:v2v5-effective-determinant}
\end{equation}
on a simply connected neighborhood on which one branch of the logarithm is
fixed.  Its first two Fréchet derivatives are
\begin{align}
 DW_D[M]&=-\operatorname{Tr}(CM),
\label{eq:v2v5-first-derivative}\\
 D^2W_D[M,N]&=\operatorname{Tr}(CMCN).
\label{eq:v2v5-second-derivative}
\end{align}
The second line is the negative of the connected Berezin covariance in
\eqref{eq:v2v4-connected-trace}, as required when the fermion determinant is
inserted into the effective action.

\begin{lemma}[Derivative formulae]
\label{lem:v2v5-determinant-derivatives}
Equations \eqref{eq:v2v5-first-derivative} and
\eqref{eq:v2v5-second-derivative} hold for arbitrary matrix directions
\(M,N\).
\end{lemma}

\begin{proof}
Jacobi's formula gives
\[
 D(\log\det)_D[M]=\operatorname{Tr}(D^{-1}M).
\]
This proves \eqref{eq:v2v5-first-derivative}.  Since
\[
 D(C)_D[N]=-CNC,
\]
differentiating the first line in direction \(N\) gives
\[
 -\operatorname{Tr}((-CNC)M)
 =\operatorname{Tr}(CNCM)
 =\operatorname{Tr}(CMCN),
\]
where the last equality is cyclicity of the finite trace.
\end{proof}

\begin{theorem}[Exact mixed Ward row before norms]
\label{thm:v2v5-mixed-ward}
For every \(R,N\in\mathcal L(H)\),
\begin{equation}
 D^2W_D\bigl[N,[R,D]\bigr]
 +DW_D\bigl[[R,N]\bigr]=0.
\label{eq:v2v5-mixed-ward}
\end{equation}
Equivalently,
\begin{equation}
 \operatorname{Tr}\bigl(CNC[R,D]\bigr)
 -\operatorname{Tr}\bigl(C[R,N]\bigr)=0.
\label{eq:v2v5-trace-ward}
\end{equation}
For \(N=[R,D]\), this gives the pure gauge-orbit identity
\begin{equation}
 \operatorname{Tr}\bigl(C[R,D]C[R,D]\bigr)
 -\operatorname{Tr}\bigl(C[R,[R,D]]\bigr)=0.
\label{eq:v2v5-pure-orbit}
\end{equation}
The two terms in \eqref{eq:v2v5-pure-orbit} are respectively the connected
quadratic term and the second-order contact insertion along the conjugation
orbit.
\end{theorem}

\begin{proof}
Conjugation leaves the determinant invariant:
\[
 W(e^{tR}De^{-tR})=W(D).
\]
Apply this identity to \(D+sN\), differentiate first in \(t\) and then in
\(s\) at \(s=t=0\).  The \(t\)-derivative of the matrix is
\([R,D+sN]\), so the chain rule gives
\eqref{eq:v2v5-mixed-ward}.

For a direct sign check, use \(CD=DC=I\):
\begin{align*}
 \operatorname{Tr}(CNC[R,D])
 &=\operatorname{Tr}(CNCRD)-\operatorname{Tr}(CNR)\\
 &=\operatorname{Tr}(NCR)-\operatorname{Tr}(CNR)\\
 &=\operatorname{Tr}(CRN)-\operatorname{Tr}(CNR)\\
 &=\operatorname{Tr}(C[R,N]).
\end{align*}
This proves \eqref{eq:v2v5-trace-ward}.  Substitute \(N=[R,D]\) for
\eqref{eq:v2v5-pure-orbit}.
\end{proof}

\begin{counterexample}[Separate trace bounds erase an exact zero]
\label{sep:v2v5-separate-bounds}
Take any \(D,R\) for which \([R,D]\ne0\).  The left side of
\eqref{eq:v2v5-pure-orbit} is exactly zero, whereas separate Schatten bounds
give only
\[
 \big|\operatorname{Tr}(C[R,D]C[R,D])\big|
 +\big|\operatorname{Tr}(C[R,[R,D]])\big|,
\]
which can be strictly positive.  Therefore the contact term must be combined
with the connected covariance before applying \cref{thm:v2v4-trace-bound}.
\end{counterexample}

\begin{proof}
Exact vanishing is \eqref{eq:v2v5-pure-orbit}.  To see that the separate
absolute values need not vanish, take
\[
 D=\begin{pmatrix}1&0\\0&2\end{pmatrix},
 \qquad
 R=\begin{pmatrix}0&1\\1&0\end{pmatrix}.
\]
Then \([R,D]\ne0\), and direct multiplication gives
\[
 \operatorname{Tr}(C[R,D]C[R,D])=-1.
\]
By \eqref{eq:v2v5-pure-orbit}, the contact trace is also \(-1\); the
difference is zero while the sum of absolute values is \(2\).
\end{proof}

\subsection{Regulator covariance is a hypothesis}
\label{sec:v2v5-regulator}

\begin{firewall}[A spectral projector need not preserve local gauge orbits]
\label{fw:v2v5-regulator}
\cref{thm:v2v5-mixed-ward} applies to a finite matrix space on which the gauge
generator \(R\) acts and the regularized Dirac matrix transforms by
conjugation.  A sharp spectral subspace of a background operator is generally
not invariant under multiplication by a local gauge parameter:
\([E_n,R]\) need not vanish.  Compressing first and then invoking
conjugation invariance would therefore be circular.  The actual BV cutoff must
either be gauge covariant or carry the resulting local breaking term into the
ghost-number-one repair of \cref{thm:v2v3-local-repair}.
\end{firewall}

For an anomaly-free combined chiral multiplet, the relevant determinant is a
section of a product determinant line before trivialization.  The
finite-matrix trace identity controls the infinitesimal conjugation row once
that line and regulator have been made compatible.  It does not replace the
global phase argument in the first volume and does not by itself prove
cancellation of a regulator anomaly.

\subsection{Transverse analytic jets}
\label{sec:v2v5-transverse-jets}

We now isolate the flat Euclidean regression that the complete Ward row should
produce.  Let \(\Pi_n:B_\kappa(0)\subset\mathbb R^d\to
\operatorname{Sym}^2(\mathbb R^d)\) be a \(C^4\) shell polarization tensor.

\begin{theorem}[Transverse even local-jet form]
\label{thm:v2v5-transverse-jet}
Assume
\begin{enumerate}[label=\textup{(\roman*)}]
\item \(\Pi_n(-k)=\Pi_n(k)\);
\item \(\Pi_n(Ok)=O\Pi_n(k)O^T\) for every \(O\in O(d)\); and
\item \(k^\mu\Pi_{n,\mu\nu}(k)=0\) for every \(k\).
\end{enumerate}
Then
\begin{equation}
 \Pi_n(0)=0,\qquad D\Pi_n(0)=0,
\label{eq:v2v5-zero-low-jets}
\end{equation}
and there is a scalar \(a_n\) such that
\begin{equation}
 \frac12D^2\Pi_n(0)[k,k]
 =a_n\bigl(|k|^2I-k\otimes k\bigr).
\label{eq:v2v5-quadratic-jet}
\end{equation}
If in addition
\begin{equation}
 \sup_{|q|\le\kappa}
 \|D^4\Pi_n(q)\|_{\mathrm{op}}
 \le A\Lambda_n^{-2},
\label{eq:v2v5-fourth-derivative}
\end{equation}
then the invariant quadratic-jet complement obeys
\begin{equation}
 \left\|
 \Pi_n(k)-a_n(|k|^2I-k\otimes k)
 \right\|
 \le\frac{A}{4!}|k|^4\Lambda_n^{-2}.
\label{eq:v2v5-transverse-remainder}
\end{equation}
\end{theorem}

\begin{proof}
\(O(d)\)-equivariance at \(k=0\) implies
\(\Pi_n(0)=c_nI\).  For nonzero \(k\), transversality gives
\[
 0=k^\mu\Pi_{n,\mu\nu}(0)+O(|k|^2)
 =c_nk_\nu+O(|k|^2).
\]
Divide by \(|k|\) and let \(k\to0\) along any direction to get \(c_n=0\).
Evenness gives \(D\Pi_n(0)=D^3\Pi_n(0)=0\).

The quadratic Taylor polynomial is an \(O(d)\)-equivariant symmetric
two-tensor homogeneous of degree two.  The invariant tensor classification
therefore gives
\[
 T_2(k)=\alpha_n|k|^2I+\beta_n k\otimes k.
\]
The degree-three part of
\(k^\mu\Pi_{n,\mu\nu}(k)=0\) is
\((\alpha_n+\beta_n)|k|^2k_\nu=0\), so
\(\beta_n=-\alpha_n\).  Set \(a_n=\alpha_n\), proving
\eqref{eq:v2v5-quadratic-jet}.

Taylor's theorem through degree three gives
\[
 \Pi_n(k)=\Pi_n(0)+D\Pi_n(0)[k]
 +\frac12D^2\Pi_n(0)[k,k]
 +\frac1{3!}D^3\Pi_n(0)[k,k,k]+R_{4,n}(k).
\]
The first, second, and fourth displayed terms vanish by
\eqref{eq:v2v5-zero-low-jets} and evenness; the third is
\eqref{eq:v2v5-quadratic-jet}.  The integral remainder and
\eqref{eq:v2v5-fourth-derivative} give
\(\|R_{4,n}(k)\|\le A|k|^4\Lambda_n^{-2}/4!\).
\end{proof}

\begin{corollary}[Summable transverse fermion remainder in four dimensions]
\label{cor:v2v5-summable-transverse}
In \(d=4\), suppose the complete regulated fermion polarization shell is
even, \(O(4)\)-equivariant, and exactly transverse after its contact and
species terms have been assembled.  If its fourth derivative satisfies the
Dirac-shell estimate \eqref{eq:v2v5-fourth-derivative}, then
\[
 \sum_{n\ge n_0}
 \left(\Pi_n(k)-a_n(|k|^2I-k\otimes k)\right)
\]
converges absolutely and locally uniformly in \(k\).  The coefficients \(a_n\)
belong to the local gauge-kinetic running coupling and are not included in the
irrelevant sum.
\end{corollary}

\begin{proof}
\cref{thm:v2v5-transverse-jet} bounds the \(n\)-th remainder by
\(C_\kappa\Lambda_n^{-2}\) on \(|k|\le\kappa\).  Since
\(\Lambda_n=2^n\Lambda_0\), the Weierstrass test proves absolute local uniform
convergence.
\end{proof}

\begin{assessment}[How V4 supplies the derivative power]
\label{ass:v2v5-v4-power}
For two zero-order insertions in \(d=4\), the raw Dirac trace degree in V4 is
\(2\).  Four external derivatives lower it to \(-2\), exactly the hypothesis
\eqref{eq:v2v5-fourth-derivative}, provided the differentiated compressed
insertions satisfy \eqref{eq:v2v4-M-order}--\eqref{eq:v2v4-N-order}.
Transversality is not taken from that norm estimate.  It is supplied by the
assembled determinant Ward row \eqref{eq:v2v5-mixed-ward} and by the
regulator/anomaly repair.  This order of operations is the adjustment forced
by the YM audit.
\end{assessment}

\begin{firewall}[The local coefficient still needs a ledger]
\label{fw:v2v5-local-first-moment}
The estimate \eqref{eq:v2v5-transverse-remainder} does not sum the local
coefficients \(a_n\).  They may be order one per dyadic shell and accumulate
logarithmically.  They must be absorbed into the running gauge coupling with
a renormalization condition.  Treating the Taylor projector as if it made
\(\sum_na_n\) convergent would repeat the missing-first-moment error isolated
by NS V31.
\end{firewall}

\subsection{V5 conclusion and next step}
\label{sec:v2v5-conclusion}

This version has proved three linked facts.  The complete finite-matrix
fermion determinant is invariant along conjugation orbits; differentiating
that identity gives an exact cancellation between the connected covariance
and its contact insertion.  Under evenness and Euclidean covariance, the
resulting transverse polarization has only the gauge-invariant quadratic
local jet \(a_n(|k|^2I-k\otimes k)\) below degree four.  Once that jet is
allocated to the running coupling, the fourth-derivative Dirac estimate makes
the remainder \(O(|k|^4\Lambda_n^{-2})\) and summable.

The next step is to replace the ideal invariant-cutoff hypothesis by an exact
formula for the commutator defect \([E_n,R]\), show that its complete
Standard--Model species sum is a local ghost-number-one cocycle, and test
whether the V3 local repair admits a primitive with a cutoff-uniform norm.
That is now the sharp algebraic-analytic bottleneck.

\clearpage
\section*{Second Volume, Version V6: Exact Spectral-Cutoff Ward Defect}
\addcontentsline{toc}{section}{Second Volume, Version V6: Exact Spectral-Cutoff Ward Defect}

\subsection*{Purpose}

V5 proved the fermion determinant Ward row on an invariant finite matrix
space and identified the failure of a fixed spectral subspace to be invariant
under a local gauge generator.  This version computes that failure exactly.
When the cutoff projector commutes with the background Dirac operator, the
first gauge derivative is still a compressed commutator and has zero trace.
The mixed derivative with a field insertion is different: its defect is an
off-diagonal excursion across the cutoff.  The formula below isolates that
excursion and supplies the precise norm that a smooth regulator or a local
Slavnov counterterm must control.

\subsection{Block algebra for a frozen cutoff}
\label{sec:v2v6-block-algebra}

Let \(H\) be a finite-dimensional complex Hilbert space, let \(E=E^*=E^2\),
and put \(Q=I-E\).  Let \(D\in GL(H)\) commute with \(E\), and assume that
\[
 D_E=EDE:E H\longrightarrow E H
\]
is invertible.  Write \(C_E=D_E^{-1}\).  The frozen-cutoff effective
functional is
\begin{equation}
 W_E(D)=-\log\det(D_E).
\label{eq:v2v6-frozen-effective}
\end{equation}
For a gauge generator \(R\) and an insertion \(N\), set
\[
 R_E=ERE,\qquad N_E=ENE.
\]

\begin{lemma}[Compressed commutator decomposition]
\label{lem:v2v6-compressed-commutator}
For arbitrary \(R,T\in\mathcal L(H)\),
\begin{equation}
 E[R,T]E
 =[R_E,T_E]+K_E(R,T),
\label{eq:v2v6-commutator-decomposition}
\end{equation}
where \(T_E=ETE\) and
\begin{align}
 K_E(R,T)
 &=ERQTE-ETQRE
\label{eq:v2v6-K-offdiagonal}\\
 &=E[E,R]QTE+ETQ[E,R]E.
\label{eq:v2v6-K-commutator}
\end{align}
If \([E,T]=0\), then \(K_E(R,T)=0\).
\end{lemma}

\begin{proof}
Insert \(I=E+Q\) between \(R\) and \(T\) in each term:
\begin{align*}
 E[R,T]E
 &=ER(E+Q)TE-ET(E+Q)RE\\
 &=[ERE,ETE]+ERQTE-ETQRE.
\end{align*}
This is \eqref{eq:v2v6-commutator-decomposition}.  Since
\[
 ERQ=E[E,R]Q,\qquad
 QRE=-Q[E,R]E,
\]
\eqref{eq:v2v6-K-commutator} follows.  If \(T\) commutes with \(E\), both
\(QTE\) and \(ETQ\) vanish.
\end{proof}

\begin{corollary}[The first frozen-cutoff gauge derivative vanishes]
\label{cor:v2v6-first-ward}
Under the standing assumption \([E,D]=0\),
\begin{equation}
 E[R,D]E=[R_E,D_E]
\label{eq:v2v6-first-compressed-orbit}
\end{equation}
and
\begin{equation}
 DW_{E,D}\bigl[E[R,D]E\bigr]=0.
\label{eq:v2v6-first-ward}
\end{equation}
\end{corollary}

\begin{proof}
Use \cref{lem:v2v6-compressed-commutator} with \(T=D\).  Then apply
\eqref{eq:v2v5-first-derivative} on \(EH\):
\[
 DW_{E,D}[[R_E,D_E]]
 =-\operatorname{Tr}_{EH}(C_E[R_E,D_E])=0,
\]
because \(C_ED_E=D_EC_E=I_{EH}\) and the trace is cyclic.
\end{proof}

\subsection{The mixed Ward defect}
\label{sec:v2v6-mixed-defect}

\begin{theorem}[Exact frozen-spectral mixed Ward defect]
\label{thm:v2v6-mixed-defect}
With the standing hypotheses above, define
\begin{align}
 \mathcal B_E(R,N;D)
 :={}&D^2W_{E,D}
 \bigl[N_E,E[R,D]E\bigr]\notag\\
 &+DW_{E,D}\bigl[E[R,N]E\bigr].
\label{eq:v2v6-defect-definition}
\end{align}
Then
\begin{equation}
 \boxed{\quad
 \mathcal B_E(R,N;D)
 =-\operatorname{Tr}_{EH}
 \bigl(C_EK_E(R,N)\bigr).
 \quad}
\label{eq:v2v6-exact-defect}
\end{equation}
Equivalently,
\begin{align}
 \mathcal B_E(R,N;D)
 ={}&
 -\operatorname{Tr}_{EH}(C_EERQNE)\notag\\
 &+\operatorname{Tr}_{EH}(C_EENQRE).
\label{eq:v2v6-two-excursions}
\end{align}
The defect vanishes if either \([E,R]=0\) or \([E,N]=0\).
\end{theorem}

\begin{proof}
By \cref{cor:v2v6-first-ward},
\[
 E[R,D]E=[R_E,D_E].
\]
By \cref{lem:v2v6-compressed-commutator},
\[
 E[R,N]E=[R_E,N_E]+K_E(R,N).
\]
The invariant finite-matrix Ward row \eqref{eq:v2v5-mixed-ward}, applied on
\(EH\), says
\[
 D^2W_{E,D}[N_E,[R_E,D_E]]
 +DW_{E,D}[[R_E,N_E]]=0.
\]
Only the last term
\(DW_{E,D}[K_E(R,N)]\) remains.  Equation
\eqref{eq:v2v5-first-derivative} gives
\[
 DW_{E,D}[K_E(R,N)]
 =-\operatorname{Tr}_{EH}(C_EK_E(R,N)),
\]
which proves \eqref{eq:v2v6-exact-defect}.  Expanding \(K_E\) by
\eqref{eq:v2v6-K-offdiagonal} proves
\eqref{eq:v2v6-two-excursions}.  If \([E,R]=0\), then
\(ERQ=QRE=0\); if \([E,N]=0\), then \(QNE=ENQ=0\).
\end{proof}

\begin{remark}[Why the first row can pass while the mixed row fails]
The spectral projector commutes with its background operator \(D\), so the
first orbit tangent compresses to a genuine commutator.  A perturbation \(N\)
need not commute with that same projector.  Differentiating the first Ward row
therefore detects the two excursions
\[
 EH\overset{N}{\longrightarrow}QH\overset{R}{\longrightarrow}EH,
 \qquad
 EH\overset{R}{\longrightarrow}QH\overset{N}{\longrightarrow}EH.
\]
Their signed difference is exactly \eqref{eq:v2v6-two-excursions}.
\end{remark}

\subsection{Schatten control of the defect}
\label{sec:v2v6-schatten}

\begin{proposition}[Off-diagonal Hilbert--Schmidt bound]
\label{prop:v2v6-defect-bound}
The exact defect obeys
\begin{align}
 |\mathcal B_E(R,N;D)|
 \le{}&
 \|C_EE[E,R]Q\|_{\mathrm{HS}}\,
 \|QNE\|_{\mathrm{HS}}\notag\\
 &+\|C_EENQ\|_{\mathrm{HS}}\,
 \|Q[E,R]E\|_{\mathrm{HS}}.
\label{eq:v2v6-defect-bound}
\end{align}
The same statement holds with conjugate Schatten exponents in each product.
\end{proposition}

\begin{proof}
Use \eqref{eq:v2v6-K-commutator} in
\eqref{eq:v2v6-exact-defect}.  The first trace is the trace of the product
\[
 (C_EE[E,R]Q)(QNE),
\]
and the second is the trace of
\[
 (C_EENQ)(Q[E,R]E).
\]
Apply Hilbert--Schmidt Cauchy--Schwarz to each.  Schatten Hölder gives the
last assertion.
\end{proof}

\begin{assessment}[The new analytic target]
\label{ass:v2v6-new-target}
The V4 Dirac estimate controls compressed factors such as \(C_EN_E\).  It
does not control the off-diagonal factors in
\eqref{eq:v2v6-defect-bound}.  The missing estimates are now explicit:
one needs cutoff-uniform control of the local-gauge commutator
\([E,R]\) together with the off-diagonal insertion \(QNE\).  A proof that
uses only \(\|E\|_{\mathrm{op}}=\|Q\|_{\mathrm{op}}=1\) loses all smallness.
\end{assessment}

\subsection{A two-mode sharp-cutoff separator}
\label{sec:v2v6-two-mode}

\begin{counterexample}[A spectral gap does not make the sharp defect small]
\label{sep:v2v6-two-mode}
Let
\[
 E=\begin{pmatrix}1&0\\0&0\end{pmatrix},
 \quad
 D=\begin{pmatrix}d_1&0\\0&d_2\end{pmatrix},
 \quad
 R=\begin{pmatrix}0&1\\-1&0\end{pmatrix},
 \quad
 N=\begin{pmatrix}0&1\\1&0\end{pmatrix},
\]
with \(d_1d_2\ne0\).  Then \([E,D]=0\), \(R^*=-R\), \(N^*=N\), and
\begin{equation}
 K_E(R,N)=2E,\qquad
 \mathcal B_E(R,N;D)=-\frac{2}{d_1}.
\label{eq:v2v6-two-mode-defect}
\end{equation}
The result is independent of \(d_2\).  Sending the discarded eigenvalue
\(|d_2|\) to infinity therefore produces no decay.
\end{counterexample}

\begin{proof}
The two off-diagonal products on \(EH\) are
\[
 ERQNE=E,\qquad ENQRE=-E.
\]
Thus \(K_E(R,N)=E-(-E)=2E\).  Since \(C_E=d_1^{-1}E\),
\eqref{eq:v2v6-exact-defect} gives
\(\mathcal B_E=-2/d_1\).
\end{proof}

\begin{firewall}[Sharp projection is not a smooth commutator estimate]
\label{fw:v2v6-sharp}
\cref{sep:v2v6-two-mode} rules out an argument based only on spectral
separation across a discontinuous projector.  Decay can still arise from a
smooth functional-calculus cutoff, from regularity of the multiplication
operator \(R\), from cancellation in the complete Standard--Model multiplet,
or from a local counterterm.  Each of those mechanisms supplies information
absent from the separator.
\end{firewall}

\subsection{Consistency of the regulator breaking}
\label{sec:v2v6-consistency}

For a fixed projector \(E\), do not assume \([E,D]=0\) temporarily.  Define
\[
 \mathcal A_E(R;D)
 =DW_{E,D}\bigl[E[R,D]E\bigr]
\]
whenever \(D_E\) is invertible.  Let \(X_R(D)=[R,D]\) be the infinitesimal
conjugation vector field.

\begin{proposition}[Finite-cutoff Wess--Zumino consistency]
\label{prop:v2v6-wz}
The breaking functional satisfies
\begin{equation}
 X_R\mathcal A_E(S;D)
 -X_S\mathcal A_E(R;D)
 +\mathcal A_E([R,S];D)=0.
\label{eq:v2v6-wz}
\end{equation}
\end{proposition}

\begin{proof}
By definition, \(\mathcal A_E(R;D)=X_RW_E(D)\).  The vector fields obey the
anti-representation identity
\[
 [X_R,X_S]=-X_{[R,S]}.
\]
Indeed,
\[
 D X_S(D)[X_R(D)]-D X_R(D)[X_S(D)]
 =[S,[R,D]]-[R,[S,D]]
 =[[S,R],D].
\]
Therefore
\[
 X_R\mathcal A_E(S)-X_S\mathcal A_E(R)
 =[X_R,X_S]W_E=-\mathcal A_E([R,S]),
\]
which is \eqref{eq:v2v6-wz}.
\end{proof}

\begin{corollary}[Local asymptotic coefficients are cocycles]
\label{cor:v2v6-local-cocycles}
Suppose a regulator calculus gives a unique finite asymptotic expansion
\begin{equation}
 \mathcal A_{E_\Lambda}(R;D)
 =\sum_{\ell=0}^J f_\ell(\Lambda)a_\ell(R;D)
 +o(1),
\label{eq:v2v6-anomaly-expansion}
\end{equation}
where the scale functions \(f_\ell\) are asymptotically linearly independent,
the remainder can be differentiated by the vector fields in
\eqref{eq:v2v6-wz}, and each \(a_\ell\) is a local functional.  Then every
\(a_\ell\) satisfies the consistency identity
\[
 X_Ra_\ell(S)-X_Sa_\ell(R)+a_\ell([R,S])=0.
\]
Thus each coefficient defines a local ghost-number-one cocycle.
\end{corollary}

\begin{proof}
Insert \eqref{eq:v2v6-anomaly-expansion} into the exact identity
\eqref{eq:v2v6-wz}.  Uniqueness and asymptotic linear independence force the
coefficient of every \(f_\ell\) to vanish.
\end{proof}

\begin{assessment}[Interface with the V3 repair theorem]
\label{ass:v2v6-v3-interface}
\cref{cor:v2v6-local-cocycles} supplies the cocycle required by
\cref{thm:v2v3-local-repair} once locality of the regulator expansion has been
proved.  The remaining questions are whether its class vanishes in the full
boundary Standard--Model local complex and whether a primitive can be chosen
with a cutoff-uniform norm.  Consistency alone proves neither assertion.
\end{assessment}

\subsection{V6 conclusion and next step}
\label{sec:v2v6-conclusion}

The fixed spectral-cutoff defect is no longer schematic.  At a background
where \([E,D]=0\), its mixed Ward row is exactly
\[
 -\operatorname{Tr}_{EH}\bigl(
 C_E(ERQNE-ENQRE)\bigr).
\]
It is carried by two off-diagonal cutoff excursions and satisfies an exact
Schatten bound.  The two-mode separator proves that a sharp spectral gap does
not make this expression small.  The associated breaking functional obeys
finite-cutoff Wess--Zumino consistency, so any local coefficient in a valid
asymptotic expansion is a ghost-number-one cocycle.

The next upstream move is to replace the discontinuous projector by a smooth
spectral multiplier \(f(D/\Lambda)\).  In an eigenbasis, its commutator is a
divided difference of \(f\).  The next version should prove the resulting
\(\Lambda^{-1}\) Hilbert--Schmidt commutator gain and determine whether it
survives the two off-diagonal factors in \eqref{eq:v2v6-defect-bound}.

\clearpage
\section*{Second Volume, Version V7: Smooth Spectral Commutators and Buffered Cutoff Gain}
\addcontentsline{toc}{section}{Second Volume, Version V7: Smooth Spectral Commutators and Buffered Cutoff Gain}

\subsection*{Purpose}

The V6 two-mode separator showed that a discontinuous projector has no
commutator smallness merely because the discarded spectrum lies far away.
This version adds the missing regularity.  For a smooth spectral multiplier,
the commutator with a gauge generator is an exact divided difference.  In
Hilbert--Schmidt norm it gains one inverse cutoff relative to
\([D,R]\).  A second theorem proves the corresponding estimate between two
sharp spectral sets separated by an actual buffer.  These are finite-cutoff
operator statements and require no continuum determinant.

\subsection{The divided-difference identity}
\label{sec:v2v7-divided-difference}

Let \(H\) be a finite-dimensional complex Hilbert space, let \(D=D^*\), and
let
\[
 D=\sum_a\lambda_aP_a
\]
be its spectral resolution, with repeated eigenvalues grouped into the
orthogonal projections \(P_a\).  For \(f\in C^1(\mathbb R)\) and
\(\Lambda>0\), set \(F_\Lambda=f(D/\Lambda)\).

\begin{theorem}[Exact Hilbert--Schmidt commutator gain]
\label{thm:v2v7-smooth-commutator}
For every \(R\in\mathcal L(H)\),
\begin{equation}
 [F_\Lambda,R]
 =\sum_{a,b}
 \left(f(\lambda_a/\Lambda)-f(\lambda_b/\Lambda)\right)
 P_aRP_b.
\label{eq:v2v7-divided-difference}
\end{equation}
If \([D,R]\) is Hilbert--Schmidt, then
\begin{equation}
 \|[F_\Lambda,R]\|_{\mathrm{HS}}
 \le\frac{\|f'\|_\infty}{\Lambda}
 \|[D,R]\|_{\mathrm{HS}}.
\label{eq:v2v7-smooth-hs}
\end{equation}
No dimension factor occurs in this comparison.
\end{theorem}

\begin{proof}
Functional calculus gives
\[
 F_\Lambda=\sum_af(\lambda_a/\Lambda)P_a.
\]
Insert the two spectral resolutions on the sides of \(R\) to obtain
\eqref{eq:v2v7-divided-difference}.  The blocks \(P_a\mathcal L(H)P_b\)
are orthogonal in the Hilbert--Schmidt inner product.  Hence
\begin{align*}
 \|[F_\Lambda,R]\|_{\mathrm{HS}}^2
 &=
 \sum_{a,b}
 |f(\lambda_a/\Lambda)-f(\lambda_b/\Lambda)|^2
 \|P_aRP_b\|_{\mathrm{HS}}^2.
\end{align*}
The scalar mean-value theorem gives
\[
 |f(\lambda_a/\Lambda)-f(\lambda_b/\Lambda)|
 \le\frac{\|f'\|_\infty}{\Lambda}|\lambda_a-\lambda_b|.
\]
Since
\[
 P_a[D,R]P_b=(\lambda_a-\lambda_b)P_aRP_b,
\]
substitution and Hilbert--Schmidt orthogonality prove
\eqref{eq:v2v7-smooth-hs}.
\end{proof}

\begin{remark}[Why the Hilbert--Schmidt proof is robust]
The estimate is a scalar Lipschitz estimate entry by entry in the eigenbasis.
It does not assert that every scalar Lipschitz function is operator Lipschitz
in operator norm.  The Hilbert--Schmidt class is the type-correct space for
this elementary divided-difference proof.
\end{remark}

\begin{counterexample}[Smoothness is quantitatively necessary]
\label{sep:v2v7-step-function}
For the step function \(f=\mathbf1_{(-\infty,0]}\), no finite
\(\|f'\|_\infty\) exists.  Taking eigenvalues
\(\lambda_1=-\varepsilon\), \(\lambda_2=\varepsilon\) and an off-diagonal
\(R\) gives
\[
 \|[f(D),R]\|_{\mathrm{HS}}\asymp\|R\|_{\mathrm{HS}},
 \qquad
 \|[D,R]\|_{\mathrm{HS}}\asymp
 \varepsilon\|R\|_{\mathrm{HS}}.
\]
The ratio diverges as \(\varepsilon\downarrow0\), reproducing the sharp-cutoff
failure at a spectral boundary.
\end{counterexample}

\begin{proof}
In the two eigenvectors, the divided difference of the step function is one,
whereas \(\lambda_1-\lambda_2=-2\varepsilon\).
\end{proof}

\subsection{A buffered sharp-projector estimate}
\label{sec:v2v7-buffer}

\begin{theorem}[Off-diagonal gain across a spectral buffer]
\label{thm:v2v7-buffered}
Let \(A,B\subset\mathbb R\) be Borel sets with
\[
 g=\operatorname{dist}(A,B)>0,
\]
and let \(E_A=\mathbf1_A(D)\), \(E_B=\mathbf1_B(D)\).  Then
\begin{equation}
 \|E_ARE_B\|_{\mathrm{HS}}
 \le g^{-1}\|E_A[D,R]E_B\|_{\mathrm{HS}}.
\label{eq:v2v7-buffered}
\end{equation}
\end{theorem}

\begin{proof}
In the spectral blocks,
\[
 E_ARE_B=\sum_{\lambda_a\in A,\,\lambda_b\in B}P_aRP_b,
\]
and every included pair satisfies
\(|\lambda_a-\lambda_b|\ge g\).  Orthogonality of the blocks and
\[
 P_a[D,R]P_b=(\lambda_a-\lambda_b)P_aRP_b
\]
give
\[
 \|E_A[D,R]E_B\|_{\mathrm{HS}}^2
 \ge g^2\|E_ARE_B\|_{\mathrm{HS}}^2.
\]
\end{proof}

\begin{assessment}[Adjacent sharp shells have no buffer]
\label{ass:v2v7-adjacent}
The complement of a sharp shell meets the shell at its spectral endpoint, so
\(\operatorname{dist}(A,B)=0\) and
\cref{thm:v2v7-buffered} gives no estimate.  A buffered decomposition must
assign the transition band separately, or a smooth partition must distribute
it among overlapping shells.  Simply relabeling the sharp complement as a
buffer does not evade \cref{sep:v2v6-two-mode}.
\end{assessment}

\subsection{Finite spectral-window estimate for a local gauge generator}
\label{sec:v2v7-local-gauge}

Let \(\mathcal D\) be a self-adjoint Dirac realization on a compact
\(d\)-manifold, and let
\[
 P_{\le b\Lambda}=\mathbf1_{[0,b\Lambda]}(|\mathcal D|).
\]
Work on the finite-dimensional space
\(H_{\Lambda,b}=P_{\le b\Lambda}H\), put
\[
 D_{\Lambda,b}=P_{\le b\Lambda}\mathcal D P_{\le b\Lambda},
 \qquad
 R_{\Lambda,b}=P_{\le b\Lambda}RP_{\le b\Lambda},
\]
and assume
\begin{equation}
 \dim H_{\Lambda,b}\le W_b\Lambda^d.
\label{eq:v2v7-window-weyl}
\end{equation}

\begin{proposition}[One derivative of the gauge parameter pays the commutator]
\label{prop:v2v7-local-gauge}
Suppose \(R\) is multiplication by a smooth bundle endomorphism and
\begin{equation}
 \|[\mathcal D,R]\|_{\mathrm{op}}
 \le c_{\mathcal D}\|\nabla R\|_\infty.
\label{eq:v2v7-dirac-locality}
\end{equation}
Then
\begin{equation}
 \left\|
 [f(D_{\Lambda,b}/\Lambda),R_{\Lambda,b}]
 \right\|_{\mathrm{HS}}
 \le
 \|f'\|_\infty c_{\mathcal D}\sqrt{W_b}\,
 \Lambda^{d/2-1}\|\nabla R\|_\infty.
\label{eq:v2v7-local-gauge-hs}
\end{equation}
\end{proposition}

\begin{proof}
Because \(P_{\le b\Lambda}\) commutes with \(\mathcal D\),
\[
 [D_{\Lambda,b},R_{\Lambda,b}]
 =P_{\le b\Lambda}[\mathcal D,R]P_{\le b\Lambda}.
\]
On a space of dimension at most \(W_b\Lambda^d\),
\[
 \|[D_{\Lambda,b},R_{\Lambda,b}]\|_{\mathrm{HS}}
 \le\sqrt{W_b}\Lambda^{d/2}
 \|[\mathcal D,R]\|_{\mathrm{op}}.
\]
Apply \cref{thm:v2v7-smooth-commutator} and then
\eqref{eq:v2v7-dirac-locality}.
\end{proof}

\begin{firewall}[The outer finite window is part of the estimate]
\label{fw:v2v7-outer-window}
On an infinite four-dimensional Hilbert space, a generic order-\(-1\)
pseudodifferential operator is not Hilbert--Schmidt.  Therefore
\eqref{eq:v2v7-local-gauge-hs} is stated on the finite window
\(H_{\Lambda,b}\).  Removing that outer window requires the stronger Schatten
threshold or a local trace expansion; it cannot be inferred from the
finite-rank proof.
\end{firewall}

\subsection{Propagation through a defect trace}
\label{sec:v2v7-defect-trace}

\begin{theorem}[The commutator gain survives one complete trace]
\label{thm:v2v7-trace-gain}
On \(H_{\Lambda,b}\), let \(C_\Lambda,A_\Lambda,N_\Lambda\) obey
\begin{align}
 \|C_\Lambda\|_{\mathrm{op}}&\le c_C\Lambda^{-1},
\label{eq:v2v7-C-order}\\
 \|A_\Lambda\|_{\mathrm{op}}&\le c_A,
\label{eq:v2v7-A-order}\\
 \|N_\Lambda\|_{\mathrm{op}}&\le c_N\Lambda^\mu.
\label{eq:v2v7-N-order}
\end{align}
Define
\begin{equation}
 \mathcal T_\Lambda(R,N)
 =\operatorname{Tr}\left(
 C_\Lambda A_\Lambda
 [f(D_{\Lambda,b}/\Lambda),R_{\Lambda,b}]
 N_\Lambda\right).
\label{eq:v2v7-model-defect}
\end{equation}
Then
\begin{equation}
 |\mathcal T_\Lambda(R,N)|
 \le C_f\Lambda^{d+\mu-2}\|\nabla R\|_\infty,
\label{eq:v2v7-model-defect-bound}
\end{equation}
where
\[
 C_f=c_Cc_Ac_Nc_{\mathcal D}W_b\|f'\|_\infty.
\]
If \(N_\Lambda(k)\) satisfies
\[
 \|D_k^jN_\Lambda(k)\|_{\mathrm{op}}
 \le c_{N,j}\Lambda^{\mu-j},
\]
then each external derivative lowers the right side of
\eqref{eq:v2v7-model-defect-bound} by one cutoff power.
\end{theorem}

\begin{proof}
Hilbert--Schmidt Cauchy--Schwarz gives
\begin{align*}
 |\mathcal T_\Lambda(R,N)|
 &\le
 \|C_\Lambda A_\Lambda
 [f(D_{\Lambda,b}/\Lambda),R_{\Lambda,b}]\|_{\mathrm{HS}}
 \|N_\Lambda\|_{\mathrm{HS}}\\
 &\le
 \|C_\Lambda\|_{\mathrm{op}}\|A_\Lambda\|_{\mathrm{op}}
 \|[f(D_{\Lambda,b}/\Lambda),R_{\Lambda,b}]\|_{\mathrm{HS}}
 \sqrt{W_b}\Lambda^{d/2}\|N_\Lambda\|_{\mathrm{op}}.
\end{align*}
Insert \eqref{eq:v2v7-local-gauge-hs} and
\eqref{eq:v2v7-C-order}--\eqref{eq:v2v7-N-order}.  The exponent is
\[
 -1+(d/2-1)+(d/2+\mu)=d+\mu-2.
\]
Differentiation in \(k\) acts only on \(N_\Lambda(k)\) under the stated
hypotheses, proving the final assertion.
\end{proof}

\begin{corollary}[Local-jet threshold for the model breaking]
\label{cor:v2v7-model-threshold}
If \(\Lambda_n=2^n\Lambda_0\), the degree-\(r\) Taylor complement of
\(\mathcal T_{\Lambda_n}(R,N(k))\) is absolutely summable whenever
\begin{equation}
 r+1>d+\mu-2.
\label{eq:v2v7-model-threshold}
\end{equation}
In \(d=4\) with a zero-order insertion, the unconditional threshold is
\(r\ge2\).
\end{corollary}

\begin{proof}
Apply the integral Taylor remainder with \(r+1\) external derivatives.  The
resulting exponent is \(d+\mu-2-(r+1)\), which is negative exactly under
\eqref{eq:v2v7-model-threshold}.
\end{proof}

\begin{assessment}[What the inverse cutoff gain accomplishes]
\label{ass:v2v7-gain-scope}
The smooth commutator prevents the regulator breaking from carrying the full
size of an arbitrary matrix commutator.  It does not by itself make the
four-dimensional defect summable: the spectral rank and the insertion still
produce local powers.  Those powers are removed by the invariant local jet
and recorded as running local terms, just as the NS V31 audit warned that a
normalized mark does not remove its required first moment.
\end{assessment}

\subsection{V7 conclusion and next step}
\label{sec:v2v7-conclusion}

For a smooth multiplier the exact bound
\[
 \|[f(D/\Lambda),R]\|_{\mathrm{HS}}
 \le\Lambda^{-1}\|f'\|_\infty\|[D,R]\|_{\mathrm{HS}}
\]
replaces the false sharp-projector smallness ruled out in V6.  For sharp
projectors the same inverse separation is available only across a genuine
spectral buffer.  On a finite Dirac window, locality of
\([\mathcal D,R]\) turns the gain into
\(O(\Lambda^{d/2-1}\|\nabla R\|_\infty)\), and the gain survives a complete
defect trace.  The remaining polynomial powers still require the local
Taylor ledger.

The next step is to construct an overlapping smooth shell partition whose
transition bands have uniformly bounded incidence, and to prove that the sum
of their commutator defects is a telescoping local cocycle rather than one
independent anomaly per shell.  Without that incidence or telescope, summing
the bounds of this version could reintroduce an artificial shell
multiplicity.

\clearpage
\section*{Second Volume, Version V8: Smooth Shell Incidence and Weighted Commutator Telescoping}
\addcontentsline{toc}{section}{Second Volume, Version V8: Smooth Shell Incidence and Weighted Commutator Telescoping}

\subsection*{Purpose}

V7 obtained an inverse-cutoff gain for one smooth spectral commutator and
warned that a shell-by-shell norm sum can reintroduce multiplicity.  The
present version constructs the dyadic smooth shell family explicitly.  Each
spectral value meets at most three shells, and the signed commutators telescope
exactly.  When covariance and insertion weights vary with the shell, an exact
Abel identity transfers the failure of telescoping to adjacent-shell
differences of those weights.  This identifies the next estimate without
claiming that smoothing by itself renormalizes the local terms.

\subsection{A dyadic smooth spectral partition}
\label{sec:v2v8-partition}

Choose an even function \(f\in C^\infty(\mathbb R;[0,1])\) such that
\[
 f(t)=1\quad\text{for }|t|\le1,
 \qquad
 f(t)=0\quad\text{for }|t|\ge2.
\]
Let \(\Lambda_n=2^n\Lambda_0\), and for a finite-dimensional self-adjoint
operator \(D\) define
\begin{equation}
 F_n=f(D/\Lambda_n),
 \qquad
 \Psi_n=F_{n+1}-F_n.
\label{eq:v2v8-smooth-shells}
\end{equation}

\begin{theorem}[Bounded incidence and exact telescope]
\label{thm:v2v8-partition}
For every eigenvalue \(\lambda\ne0\) of \(D\), the scalar multiplier
\(\Psi_n(\lambda)\) can be nonzero only if
\begin{equation}
 \Lambda_n<|\lambda|<4\Lambda_n.
\label{eq:v2v8-shell-support}
\end{equation}
Consequently, a fixed eigenvalue belongs to at most three shell supports.
For integers \(m\le N\),
\begin{equation}
 \sum_{n=m}^N\Psi_n=F_{N+1}-F_m.
\label{eq:v2v8-shell-telescope}
\end{equation}
If the dyadic range is extended in both directions, then on
\((\ker D)^\perp\),
\begin{equation}
 \sum_{n\in\mathbb Z}\Psi_n=I
\label{eq:v2v8-partition-identity}
\end{equation}
in finite dimension.  The kernel is excluded from the shell sum and remains
in the determinant-line sector.
\end{theorem}

\begin{proof}
If \(|\lambda|\le\Lambda_n\), then both
\(f(\lambda/\Lambda_n)\) and
\(f(\lambda/\Lambda_{n+1})\) equal one.  If
\(|\lambda|\ge4\Lambda_n=2\Lambda_{n+1}\), both equal zero.  This proves
\eqref{eq:v2v8-shell-support}.  An interval of scale ratio four contains at
most three points of a fixed dyadic grid, including possible endpoints, so the
incidence is at most three.

Equation \eqref{eq:v2v8-shell-telescope} is the finite telescoping sum of
\(F_{n+1}-F_n\).  If \(\lambda\ne0\), then
\[
 f(\lambda/\Lambda_m)\longrightarrow0\quad(m\to-\infty),
 \qquad
 f(\lambda/\Lambda_{N+1})\longrightarrow1\quad(N\to\infty).
\]
Apply the finite identity on each eigenspace and take the limits.  At
\(\lambda=0\), every \(F_n\) equals one and every \(\Psi_n\) is zero, which
explains the kernel statement.
\end{proof}

\begin{corollary}[Commutators telescope before norms]
\label{cor:v2v8-commutator-telescope}
For every \(R\in\mathcal L(H)\),
\begin{equation}
 \sum_{n=m}^N[\Psi_n,R]
 =[F_{N+1},R]-[F_m,R].
\label{eq:v2v8-commutator-telescope}
\end{equation}
Moreover,
\begin{equation}
 \|[\Psi_n,R]\|_{\mathrm{HS}}
 \le\frac{3\|f'\|_\infty}{2\Lambda_n}
 \|[D,R]\|_{\mathrm{HS}}.
\label{eq:v2v8-shell-commutator-bound}
\end{equation}
\end{corollary}

\begin{proof}
Commutation with \(R\) is linear, so
\eqref{eq:v2v8-shell-telescope} gives
\eqref{eq:v2v8-commutator-telescope}.  By
\cref{thm:v2v7-smooth-commutator},
\[
 \|[F_n,R]\|_{\mathrm{HS}}
 \le\frac{\|f'\|_\infty}{\Lambda_n}\|[D,R]\|_{\mathrm{HS}},
\]
while the same estimate for \(F_{n+1}\) has coefficient
\((2\Lambda_n)^{-1}\).  The triangle inequality proves
\eqref{eq:v2v8-shell-commutator-bound}.
\end{proof}

\begin{firewall}[Norm summation discards the telescope]
\label{fw:v2v8-norm-telescope}
The exact identity \eqref{eq:v2v8-commutator-telescope} is signed.  Replacing
its left side by \(\sum_n\|[\Psi_n,R]\|_{\mathrm{HS}}\) discards the endpoint
cancellation.  Bounded spectral incidence prevents a repeated-eigenvalue
count, but it does not justify taking absolute values of independently
weighted shell terms.
\end{firewall}

\subsection{Weighted Abel summation}
\label{sec:v2v8-abel}

Let \(\mathcal X\) be the vector space of operators on \(H\), let \(Y\) be
a vector space of source coefficients, and let
\(\ell_n:\mathcal X\to Y\) be arbitrary linear maps.  The scalar trace case is
\(Y=\mathbb C\).  In the regulator application,
\[
 \ell_n(X)=\operatorname{Tr}(A_n[X,R]B_n)
\]
contains the covariance, insertion, and local contraction at shell \(n\).

\begin{theorem}[Exact weighted shell identity]
\label{thm:v2v8-abel}
For \(m\le N\),
\begin{equation}
 \boxed{\quad
 \sum_{n=m}^N\ell_n(\Psi_n)
 =\ell_N(F_{N+1})-\ell_m(F_m)
 -\sum_{n=m+1}^N(\ell_n-\ell_{n-1})(F_n).
 \quad}
\label{eq:v2v8-abel}
\end{equation}
Thus a common functional \(\ell_n=\ell\) leaves only two endpoints.  All
failure of telescoping for varying weights is carried by the adjacent
differences \(\ell_n-\ell_{n-1}\).
\end{theorem}

\begin{proof}
Using \(\Psi_n=F_{n+1}-F_n\),
\[
 \sum_{n=m}^N\ell_n(\Psi_n)
 =\sum_{n=m}^N\ell_n(F_{n+1})
  -\sum_{n=m}^N\ell_n(F_n).
\]
Reindex the first sum.  Its top term is \(\ell_N(F_{N+1})\), its bottom
comparison leaves \(-\ell_m(F_m)\), and the coefficient of \(F_n\) for
\(m+1\le n\le N\) is
\(\ell_{n-1}-\ell_n\).  This is \eqref{eq:v2v8-abel}.
\end{proof}

\begin{counterexample}[Bounded incidence alone does not control weights]
\label{sep:v2v8-weighted}
Take a one-dimensional eigenspace with scalar shells \(\Psi_n(\lambda)\), and
choose \(\ell_n(x)=a_nx\).  Then
\[
 \sum_n\ell_n(\Psi_n(\lambda))
 =\sum_na_n\Psi_n(\lambda).
\]
Although at most three terms are nonzero for this fixed \(\lambda\), their
size is governed by the corresponding \(a_n\).  No bound on \(a_n\) follows
from the partition identity.  In particular, multiplying the active shell by
an arbitrarily large coefficient defeats any conclusion based on incidence
alone.
\end{counterexample}

\begin{proof}
All assertions follow from scalar multiplication.  The example separates the
combinatorial incidence statement from the analytic weight estimate.
\end{proof}

\subsection{Adjacent variation for trace weights}
\label{sec:v2v8-trace-variation}

Define
\begin{equation}
 \ell_n(X)=\operatorname{Tr}(A_n[X,R]B_n),
\label{eq:v2v8-trace-functional}
\end{equation}
where the products are trace class.  The next identity retains the complete
commutator before taking norms.

\begin{proposition}[Exact adjacent trace difference]
\label{prop:v2v8-adjacent-trace}
For every \(n\),
\begin{align}
 (\ell_n-\ell_{n-1})(F_n)
 ={}&
 \operatorname{Tr}\bigl(
 (A_n-A_{n-1})[F_n,R]B_n\bigr)\notag\\
 &+\operatorname{Tr}\bigl(
 A_{n-1}[F_n,R](B_n-B_{n-1})\bigr).
\label{eq:v2v8-adjacent-trace}
\end{align}
If the displayed differences and factors are Hilbert--Schmidt or bounded as
indicated, then
\begin{align}
 |(\ell_n-\ell_{n-1})(F_n)|
 \le{}&
 \|A_n-A_{n-1}\|_{\mathrm{HS}}
 \|[F_n,R]\|_{\mathrm{HS}}\|B_n\|_{\mathrm{op}}
 \notag\\
 &+\|A_{n-1}\|_{\mathrm{op}}
 \|[F_n,R]\|_{\mathrm{HS}}
 \|B_n-B_{n-1}\|_{\mathrm{HS}}.
\label{eq:v2v8-adjacent-bound}
\end{align}
\end{proposition}

\begin{proof}
Add and subtract
\(\operatorname{Tr}(A_{n-1}[F_n,R]B_n)\) to obtain
\eqref{eq:v2v8-adjacent-trace}.  In the first term, multiply the two
Hilbert--Schmidt factors \(A_n-A_{n-1}\) and \([F_n,R]B_n\).  In the second,
multiply the Hilbert--Schmidt factors
\(A_{n-1}[F_n,R]\) and \(B_n-B_{n-1}\).  Schatten Hölder and
\(\|XY\|_{\mathrm{HS}}\le\|X\|_{\mathrm{HS}}\|Y\|_{\mathrm{op}}\)
or its reversed form prove \eqref{eq:v2v8-adjacent-bound}.
\end{proof}

\begin{theorem}[A sufficient weighted-telescope criterion]
\label{thm:v2v8-weighted-criterion}
Let the dyadic range tend to the ultraviolet, and suppose the endpoint terms
in \eqref{eq:v2v8-abel} converge.  Assume that for some \(\varepsilon>0\),
\begin{align}
 &\|[F_n,R]\|_{\mathrm{HS}}
 \Bigl(
 \|A_n-A_{n-1}\|_{\mathrm{HS}}\|B_n\|_{\mathrm{op}}
 \notag\\
 &\hspace{42mm}
 +\|A_{n-1}\|_{\mathrm{op}}
  \|B_n-B_{n-1}\|_{\mathrm{HS}}
 \Bigr)
 \le C\Lambda_n^{-\varepsilon}.
\label{eq:v2v8-variation-criterion}
\end{align}
Then the weighted shell series
\[
 \sum_n\operatorname{Tr}(A_n[\Psi_n,R]B_n)
\]
converges absolutely after its two telescoping endpoints are separated.
\end{theorem}

\begin{proof}
\cref{prop:v2v8-adjacent-trace} and
\eqref{eq:v2v8-variation-criterion} bound the \(n\)-th interior term in
\eqref{eq:v2v8-abel} by \(C\Lambda_n^{-\varepsilon}\).  The dyadic series is
geometric.  The endpoint convergence is assumed separately.
\end{proof}

\begin{assessment}[What must vary slowly]
\label{ass:v2v8-slow-variation}
The theorem does not demand absolute summability of the raw shell commutators.
It demands summability of adjacent changes in the complete trace weights.
For the Standard--Model application, \(A_n\) contains the regulated inverse
and any left contraction, while \(B_n\) contains the differentiated
interaction or metric insertion.  Proving
\eqref{eq:v2v8-variation-criterion} requires a resolvent identity between
neighboring smooth cutoffs and the same local Taylor subtraction on both
shells.  Estimating \(A_n\) and \(B_n\) independently cannot reveal this
adjacent cancellation.
\end{assessment}

\subsection{Cocycle telescope and local projection}
\label{sec:v2v8-cocycle}

Suppose each shell breaking \(\mathcal A_n\) is a ghost-number-one cocycle as
in \cref{cor:v2v6-local-cocycles}, and its commutator-dependent part is
represented by \(\ell_n(\Psi_n)\).  Applying the Slavnov differential to the
finite Abel identity gives the same identity for the corresponding breaking
terms.  Closure of the left side implies closure of the complete right side;
it does not imply that each endpoint or adjacent-difference term is separately
closed.  Separate closure requires additional covariance of the regulator
family.

\begin{proposition}[Projection after telescoping]
\label{prop:v2v8-project-after}
Let \(P_{\mathrm{loc}}\) be a fixed degree-preserving local projector.  For
every finite range,
\begin{align}
 P_{\mathrm{loc}}\sum_{n=m}^N\ell_n(\Psi_n)
 ={}&P_{\mathrm{loc}}\ell_N(F_{N+1})
 -P_{\mathrm{loc}}\ell_m(F_m)\notag\\
 &-\sum_{n=m+1}^N
 P_{\mathrm{loc}}(\ell_n-\ell_{n-1})(F_n).
\label{eq:v2v8-projected-abel}
\end{align}
No shellwise commuting-projector assumption is used in this algebraic
identity.  Slavnov closure of the projected pieces still requires the
splitting or repair criterion of V3.
\end{proposition}

\begin{proof}
Apply the linear map \(P_{\mathrm{loc}}\) to
\eqref{eq:v2v8-abel}.  The last sentence follows because linearity alone does
not imply \(sP_{\mathrm{loc}}=P_{\mathrm{loc}}s\).
\end{proof}

\begin{firewall}[A different projector at every shell breaks the identity]
\label{fw:v2v8-fixed-projector}
Equation \eqref{eq:v2v8-projected-abel} uses one fixed local target and one
fixed projector over the summed range.  If \(P_{\mathrm{loc},n}\) varies with
\(n\), its adjacent difference creates an additional term.  Such a term may
be legitimate as coupling transport, but it must be written and bounded; it
cannot be hidden inside an unspecified local label.
\end{firewall}

\subsection{V8 conclusion and next step}
\label{sec:v2v8-conclusion}

The smooth dyadic multipliers now form an exact bounded-incidence partition on
the nonzero spectrum.  Their commutators telescope before norms.  For the
weighted regulator defect, the exact Abel identity shows that the true
interior cost is
\[
 (\ell_n-\ell_{n-1})(F_n),
\]
not an independent copy of \(\ell_n([\Psi_n,R])\) at every shell.
Schatten Hölder reduces that cost to adjacent differences of the complete
covariance and insertion weights.

The next analytic task is consequently precise: use the resolvent identity
for neighboring regulated Dirac inverses to prove
\eqref{eq:v2v8-variation-criterion} after the common invariant local jet has
been extracted.  The estimate must keep the contact and species sum assembled,
as required by V5, and must construct any local Slavnov primitive with the
uniform norm required by V3.

\clearpage
\section*{Second Volume, Version V9: Adjacent Resolvent Variation on a Transition-Supported Regression}
\addcontentsline{toc}{section}{Second Volume, Version V9: Adjacent Resolvent Variation on a Transition-Supported Regression}

\subsection*{Purpose and scope}

V8 reduced the weighted commutator sum to adjacent differences of complete
trace weights.  This version evaluates the covariance part of that difference
by the resolvent identity.  The free cutoff change is automatically supported
in one smooth transition shell and carries the expected inverse Dirac power.
A running kinetic operator adds a second resolvent term.  We prove the desired
bound when that running step is supported in the same transition and commutes
with the reference spectral decomposition.  This is a rigorous regression,
not a claim that the interacting BV step has already been diagonalized.

\subsection{The exact adjacent regulated resolvent identity}
\label{sec:v2v9-resolvent}

Let \(D_0=D_0^*\) be invertible on a finite-dimensional Hilbert space and let
\[
 F_n=f(D_0/\Lambda_n)
\]
be the smooth low-pass multipliers from V8.  Let \(V_n=V_n^*\), put
\[
 D_n=D_0+V_n,\qquad C_n=D_n^{-1},
\qquad A_n=F_nC_nF_n,
\]
and assume every \(D_n\) is invertible.

\begin{lemma}[Two cutoff terms and one running term]
\label{lem:v2v9-adjacent-resolvent}
With
\[
 \Delta F_n=F_n-F_{n-1},
\qquad
 \Delta V_n=V_n-V_{n-1},
\]
one has the exact identity
\begin{align}
 A_n-A_{n-1}
 ={}&
 \Delta F_n C_nF_n
 +F_{n-1}C_n\Delta F_n\notag\\
 &-F_{n-1}C_n\Delta V_nC_{n-1}F_{n-1}.
\label{eq:v2v9-general-adjacent}
\end{align}
\end{lemma}

\begin{proof}
Add and subtract \(F_{n-1}C_nF_n\) and
\(F_{n-1}C_nF_{n-1}\).  The first two differences give the first line of
\eqref{eq:v2v9-general-adjacent}.  The remaining factor is
\[
 F_{n-1}(C_n-C_{n-1})F_{n-1}.
\]
The second resolvent identity gives
\[
 C_n-C_{n-1}
 =-C_n(D_n-D_{n-1})C_{n-1}
 =-C_n\Delta V_nC_{n-1},
\]
which proves the last line.
\end{proof}

\begin{firewall}[A uniform inverse bound is not a shell inverse bound]
\label{fw:v2v9-inverse-bound}
The first two terms in \eqref{eq:v2v9-general-adjacent} contain
\(\Delta F_n\), but \(C_n\) need not preserve its spectral support when
\([D_0,V_n]\ne0\).  Replacing \(C_n\) by its global operator norm can therefore
lose the ultraviolet factor \(\Lambda_n^{-1}\) to low modes.  The localization
needed below is automatic in the commuting regression and remains a genuine
off-diagonal resolvent estimate in the full problem.
\end{firewall}

\subsection{The free cumulative covariance}
\label{sec:v2v9-free}

First take \(V_n=0\) for every \(n\).  Since \(D_0^{-1}\) commutes with every
\(F_n\),
\begin{equation}
 A_n-A_{n-1}
 =\Theta_nD_0^{-1},
\qquad
 \Theta_n=F_n^2-F_{n-1}^2.
\label{eq:v2v9-free-difference}
\end{equation}

\begin{proposition}[Free adjacent Hilbert--Schmidt power]
\label{prop:v2v9-free-power}
Assume the transition rank bound
\begin{equation}
 \operatorname{rank}\mathbf1_{\{
 c_-\Lambda_n\le|D_0|\le c_+\Lambda_n\}}
 \le W\Lambda_n^d
\label{eq:v2v9-transition-rank}
\end{equation}
for fixed \(0<c_-<c_+\), and assume the support of \(\Theta_n\) lies in this
transition.  Then
\begin{equation}
 \|A_n-A_{n-1}\|_{\mathrm{HS}}
 \le c_-^{-1}\sqrt W\,\Lambda_n^{d/2-1}.
\label{eq:v2v9-free-hs}
\end{equation}
\end{proposition}

\begin{proof}
The scalar multiplier satisfies \(\|\Theta_n\|_{\mathrm{op}}\le1\).
On its support,
\(\|D_0^{-1}\|_{\mathrm{op}}\le(c_-\Lambda_n)^{-1}\).
The product in \eqref{eq:v2v9-free-difference} has rank at most
\(W\Lambda_n^d\).  Hence
\[
 \|\Theta_nD_0^{-1}\|_{\mathrm{HS}}
 \le\sqrt W\Lambda_n^{d/2}
 (c_-\Lambda_n)^{-1}.
\]
\end{proof}

\subsection{A commuting transition-supported running step}
\label{sec:v2v9-running}

We now impose the hypotheses that make the running contribution a controlled
regression.
\begin{enumerate}[label=\textup{(R\arabic*)}]
\item \([D_0,V_j]=0\) for \(j=n-1,n\).
\item There is a spectral transition projector \(P_n\), commuting with
\(D_0,V_{n-1},V_n\), such that
\[
 \Theta_n=P_n\Theta_n,\qquad
 \Delta V_n=P_n\Delta V_nP_n.
\]
\item On \(P_nH\),
\begin{equation}
 \|C_jP_n\|_{\mathrm{op}}\le c_C\Lambda_n^{-1},
 \qquad j=n-1,n,
\label{eq:v2v9-shell-coercivity}
\end{equation}
and \(\operatorname{rank}P_n\le W\Lambda_n^d\).
\item For some order \(\rho\),
\begin{equation}
 \|\Delta V_n\|_{\mathrm{op}}\le v\Lambda_n^\rho.
\label{eq:v2v9-running-order}
\end{equation}
\end{enumerate}

\begin{theorem}[Adjacent running-resolvent estimate]
\label{thm:v2v9-running-estimate}
Under (R1)--(R4),
\begin{equation}
 A_n-A_{n-1}
 =\Theta_nC_n
 -F_{n-1}^2C_n\Delta V_nC_{n-1},
\label{eq:v2v9-commuting-difference}
\end{equation}
and
\begin{equation}
 \|A_n-A_{n-1}\|_{\mathrm{HS}}
 \le
 C_1\Lambda_n^{d/2-1}
 +C_2\Lambda_n^{d/2+\rho-2},
\label{eq:v2v9-running-hs}
\end{equation}
where \(C_1=c_C\sqrt W\) and
\(C_2=c_C^2v\sqrt W\).
\end{theorem}

\begin{proof}
All factors commute under (R1), so
\[
 F_nC_nF_n=F_n^2C_n.
\]
Add and subtract \(F_{n-1}^2C_n\), then use the resolvent identity:
\begin{align*}
 A_n-A_{n-1}
 &=(F_n^2-F_{n-1}^2)C_n
   +F_{n-1}^2(C_n-C_{n-1})\\
 &=\Theta_nC_n
  -F_{n-1}^2C_n\Delta V_nC_{n-1}.
\end{align*}
This proves \eqref{eq:v2v9-commuting-difference}.

Both terms are supported in \(P_nH\).  The first has operator norm at most
\(c_C\Lambda_n^{-1}\) and rank at most \(W\Lambda_n^d\), giving
\[
 \|\Theta_nC_n\|_{\mathrm{HS}}
 \le c_C\sqrt W\Lambda_n^{d/2-1}.
\]
The second has operator norm at most
\[
 c_C\Lambda_n^{-1}\,
 v\Lambda_n^\rho\,
 c_C\Lambda_n^{-1}
 =c_C^2v\Lambda_n^{\rho-2}
\]
and the same rank bound.  Summing the two Hilbert--Schmidt estimates proves
\eqref{eq:v2v9-running-hs}.
\end{proof}

\begin{counterexample}[Resolvent variation needs a running-step bound]
\label{sep:v2v9-running-step}
Even in one dimension, take \(F_n=F_{n-1}=1\),
\(D_{n-1}=1\), and \(D_n=\varepsilon\ne0\).  Then
\[
 C_n-C_{n-1}=\varepsilon^{-1}-1.
\]
This can be arbitrarily large as \(\varepsilon\to0\).  Thus smooth cutoff
incidence does not control a running inverse unless a coercivity and
\(\Delta V_n\) estimate such as (R3)--(R4) is supplied.
\end{counterexample}

\begin{proof}
The formula is direct.  It also equals
\(-C_n(D_n-D_{n-1})C_{n-1}\), so the resolvent identity records rather than
removes the loss of coercivity.
\end{proof}

\subsection{Insertion into the weighted commutator row}
\label{sec:v2v9-weighted-row}

Let \(R\) be a smooth gauge generator and retain the finite spectral window of
V7.  Assume
\begin{equation}
 \|[F_n,R]\|_{\mathrm{HS}}
 \le c_R\Lambda_n^{d/2-1}\|\nabla R\|_\infty.
\label{eq:v2v9-gauge-commutator}
\end{equation}
Let \(B_n(k)\) be a \(C^{r+1}\) insertion family satisfying
\begin{equation}
 \sup_{|k|\le\kappa}
 \|D_k^jB_n(k)\|_{\mathrm{op}}
 \le b_j\Lambda_n^{\mu-j},
 \qquad 0\le j\le r+1.
\label{eq:v2v9-B-order}
\end{equation}
Define the first adjacent term from V8 by
\begin{equation}
 J_n(k)=
 \operatorname{Tr}\bigl(
 (A_n-A_{n-1})[F_n,R]B_n(k)\bigr).
\label{eq:v2v9-J}
\end{equation}

\begin{theorem}[Projected adjacent-resolvent summability]
\label{thm:v2v9-projected-J}
Under the hypotheses of
\cref{thm:v2v9-running-estimate} and
\eqref{eq:v2v9-gauge-commutator}--\eqref{eq:v2v9-B-order},
\begin{align}
 \sup_{|k|\le\kappa}\|D_k^mJ_n(k)\|_{\mathrm{op}}
 \le{}&
 K_m\|\nabla R\|_\infty
 \bigl(
 \Lambda_n^{d+\mu-2-m}\notag\\
 &\hspace{27mm}
 +\Lambda_n^{d+\rho+\mu-3-m}
 \bigr)
\label{eq:v2v9-J-derivative}
\end{align}
for \(0\le m\le r+1\).  Let \(Q_r\) remove the Taylor jet of \(J_n\) through
degree \(r\).  Then
\begin{align}
 |(Q_rJ_n)(k)|
 \le{}&
 K_r'|k|^{r+1}\|\nabla R\|_\infty
 \bigl(
 \Lambda_n^{d+\mu-2-(r+1)}\notag\\
 &\hspace{31mm}
 +\Lambda_n^{d+\rho+\mu-3-(r+1)}
 \bigr).
\label{eq:v2v9-J-remainder}
\end{align}
The dyadic series of projected terms is absolutely convergent if
\begin{equation}
 r+1>
 \max\{d+\mu-2,\ d+\rho+\mu-3\}.
\label{eq:v2v9-J-threshold}
\end{equation}
In particular, a running step of order \(\rho\le1\) does not worsen the free
threshold.
\end{theorem}

\begin{proof}
Differentiate \eqref{eq:v2v9-J} in the external parameter.  Under the stated
hypotheses only \(B_n(k)\) depends on \(k\).  Schatten Hölder gives
\begin{align*}
 \|D_k^mJ_n(k)\|_{\mathrm{op}}
 &\le
 \|A_n-A_{n-1}\|_{\mathrm{HS}}
 \|[F_n,R]\|_{\mathrm{HS}}
 \|D_k^mB_n(k)\|_{\mathrm{op}}.
\end{align*}
Insert \eqref{eq:v2v9-running-hs},
\eqref{eq:v2v9-gauge-commutator}, and
\eqref{eq:v2v9-B-order}.  Multiplication of the powers gives
\[
 (d/2-1)+(d/2-1)+(\mu-m)=d+\mu-2-m
\]
for the free cutoff term and
\[
 (d/2+\rho-2)+(d/2-1)+(\mu-m)
 =d+\rho+\mu-3-m
\]
for the running term.  This proves
\eqref{eq:v2v9-J-derivative}.  The integral Taylor remainder of V1 with
\(m=r+1\) proves \eqref{eq:v2v9-J-remainder}.  Both dyadic powers are negative
exactly under \eqref{eq:v2v9-J-threshold}.
\end{proof}

\begin{corollary}[Four-dimensional zero-order regression]
\label{cor:v2v9-four-dimensional}
If \(d=4\), \(\mu=0\), and \(\rho\le1\), subtraction through Taylor degree
\(r=2\) makes the first adjacent resolvent term summable.  If an assembled
parity identity removes odd jets, the leading post-subtraction term has even
degree at least four and gains a further inverse cutoff power.
\end{corollary}

\begin{proof}
The threshold in \eqref{eq:v2v9-J-threshold} is at most \(2\), and
\(r+1=3>2\).  The parity statement follows by applying the next nonzero even
Taylor remainder.
\end{proof}

\subsection{What remains outside the regression}
\label{sec:v2v9-noncommuting}

\begin{assessment}[Noncommuting BV step]
\label{ass:v2v9-noncommuting}
The full interacting update need not satisfy
\([D_0,V_n]=0\), and \(\Delta V_n\) need not be sharply supported in one
spectral transition.  Equation \eqref{eq:v2v9-general-adjacent} remains exact,
but the proof of \eqref{eq:v2v9-running-hs} must be replaced by estimates for
\[
 \Delta F_nC_nF_n,\qquad
 F_{n-1}C_n\Delta F_n,\qquad
 F_{n-1}C_n\Delta V_nC_{n-1}F_{n-1}.
\]
The first two require off-diagonal resolvent localization.  The third requires
quasilocal decay of the Wilsonian step and uniform coercivity on its active
spectral band.  These are analytic estimates, not consequences of the
resolvent identity.
\end{assessment}

\begin{firewall}[Only the first V8 adjacent term is closed here]
\label{fw:v2v9-first-term}
\cref{thm:v2v9-projected-J} treats the term containing
\(A_n-A_{n-1}\) in \eqref{eq:v2v8-adjacent-trace}.  It does not treat a
shell-dependent insertion difference \(B_n-B_{n-1}\).  The latter vanishes in
the fixed-insertion regression but is present for a running composite metric
source.  It requires its own connected-cumulant or Duhamel estimate.
\end{firewall}

\subsection{V9 conclusion and V10 mandate}
\label{sec:v2v9-conclusion}

The adjacent covariance weight now has an exact three-term resolvent formula.
In the commuting transition-supported regression, its Hilbert--Schmidt size is
\[
 O(\Lambda_n^{d/2-1})
 +O(\Lambda_n^{d/2+\rho-2}).
\]
Combined with the smooth gauge commutator and a degree-\(r\) Taylor
subtraction, the first weighted Abel term is summable under
\eqref{eq:v2v9-J-threshold}.  A running step of order at most one preserves the
free threshold.

The next active version is V10 and must begin with the required companion
audit.  After that audit, the programme should attack the second adjacent term
containing \(B_n-B_{n-1}\), using the connected-cumulant derivative rather
than assuming a fixed insertion.

\clearpage
\section*{Second Volume, Version V10: Companion Audit and the Running-Insertion Cumulant}
\addcontentsline{toc}{section}{Second Volume, Version V10: Companion Audit and the Running-Insertion Cumulant}

\subsection{Mandatory companion audit}
\label{sec:v2v10-audit}

\begin{assessment}[Read-only V10 audit]
\label{ass:v2v10-audit}
The authoritative files read at the start of V10 were:
\begin{center}
\small
\begin{tabular}{llll}
\toprule
Programme & Version & Bytes & SHA--256\\
\midrule
Navier--Stokes & V31 & \(887537\) &
\texttt{F1121B62238AD5491BB7CF03635EA9934942EDF573ED8FAAA9EE79E1D38A6696}\\
Yang--Mills & compact V6 & \(88921\) &
\texttt{A4364BD88A9949D45C9252FB305F155F15DDE60BD68CA4235DA1BADDDB0C5A60}\\
\bottomrule
\end{tabular}
\end{center}
NS V31 still isolates a stopping-time-correlated source occupation that
cannot be replaced by its global continuation-class norm; its smooth packet
and balanced-helicity examples delimit false gains.  YM compact V6 rewrites
the exact Cartan fourth term as one genuine centered cumulant on a common
Wilson probability space.  Its independent-copy covariance identity shows
that the same resampling must not be counted as two independent bridges.
\end{assessment}

\begin{firewall}[Audit transfer]
\label{fw:v2v10-audit}
No NS regularity estimate or YM Wilson bound is imported.  The transferable
proof order is used: retain the actual correlated state derivative, form its
complete connected cumulant, and only then apply a norm.  This rules out
estimating the moment and normalization derivatives independently.
\end{firewall}

\subsection{Normalized Duhamel differentiation}
\label{sec:v2v10-duhamel}

Let \((\Omega,\mu)\) be a finite measure space and let
\(S_t:\Omega\to\mathbb R\), \(0\le t\le1\), be differentiable.  Fix
\(\beta=\hbar^{-1}>0\), and define
\begin{equation}
 \langle X\rangle_t
 =Z_t^{-1}\int_\Omega X\,e^{-\beta S_t}\,\mathrm d\mu,
 \qquad
 Z_t=\int_\Omega e^{-\beta S_t}\,\mathrm d\mu.
\label{eq:v2v10-state}
\end{equation}
All observables below are assumed integrable enough to justify the displayed
derivatives.

\begin{theorem}[Exact state and covariance derivatives]
\label{thm:v2v10-duhamel}
For a differentiable observable \(X_t\),
\begin{equation}
 \frac{\mathrm d}{\mathrm dt}\langle X_t\rangle_t
 =\langle\dot X_t\rangle_t
 -\beta\operatorname{Cov}_t(X_t,\dot S_t).
\label{eq:v2v10-expectation-derivative}
\end{equation}
For differentiable \(F_t,G_t\),
\begin{align}
 \frac{\mathrm d}{\mathrm dt}\operatorname{Cov}_t(F_t,G_t)
 ={}&
 \operatorname{Cov}_t(\dot F_t,G_t)
 +\operatorname{Cov}_t(F_t,\dot G_t)\notag\\
 &-\beta\,
 \kappa_{3,t}(F_t,G_t,\dot S_t),
\label{eq:v2v10-covariance-derivative}
\end{align}
where
\begin{equation}
 \kappa_{3,t}(F,G,H)
 =\left\langle
 (F-\langle F\rangle_t)
 (G-\langle G\rangle_t)
 (H-\langle H\rangle_t)
 \right\rangle_t.
\label{eq:v2v10-third-cumulant}
\end{equation}
\end{theorem}

\begin{proof}
Differentiate the quotient in \eqref{eq:v2v10-state}.  Since
\[
 \dot Z_t=-\beta Z_t\langle\dot S_t\rangle_t,
\]
the quotient rule gives
\[
 \frac{\mathrm d}{\mathrm dt}\langle X_t\rangle_t
 =\langle\dot X_t\rangle_t
 -\beta\langle X_t\dot S_t\rangle_t
 +\beta\langle X_t\rangle_t\langle\dot S_t\rangle_t,
\]
which is \eqref{eq:v2v10-expectation-derivative}.

Apply that identity to
\(\langle F_tG_t\rangle_t-\langle F_t\rangle_t\langle G_t\rangle_t\).
The direct derivatives combine into the first line of
\eqref{eq:v2v10-covariance-derivative}.  The remaining five terms are
\begin{align*}
 -\beta\bigl(&
 \langle F_tG_t\dot S_t\rangle_t
 -\langle F_tG_t\rangle_t\langle\dot S_t\rangle_t\\
 &-\langle F_t\dot S_t\rangle_t\langle G_t\rangle_t
 -\langle G_t\dot S_t\rangle_t\langle F_t\rangle_t
 +2\langle F_t\rangle_t\langle G_t\rangle_t
   \langle\dot S_t\rangle_t\bigr),
\end{align*}
which is \(-\beta\kappa_{3,t}(F_t,G_t,\dot S_t)\) after expanding
\eqref{eq:v2v10-third-cumulant}.
\end{proof}

\begin{corollary}[Exact adjacent running-insertion formula]
\label{cor:v2v10-adjacent}
Let \(S_t=S_{n-1}+t\Delta S_n\), and let \(F_t,G_t\) interpolate the two
observables defining a connected insertion
\(B_t=\operatorname{Cov}_t(F_t,G_t)\).  Then
\begin{align}
 B_n-B_{n-1}
 =\int_0^1\Bigl[
 &\operatorname{Cov}_t(\dot F_t,G_t)
 +\operatorname{Cov}_t(F_t,\dot G_t)\notag\\
 &-\beta\kappa_{3,t}(F_t,G_t,\Delta S_n)
 \Bigr]\,\mathrm dt.
\label{eq:v2v10-adjacent-insertion}
\end{align}
\end{corollary}

\begin{proof}
Integrate \eqref{eq:v2v10-covariance-derivative} from zero to one.
\end{proof}

\begin{counterexample}[The normalization term cannot be dropped]
\label{sep:v2v10-normalization}
Let \(\Omega=\{-1,1\}\) with uniform \(\mu\), take
\(S_t(x)=-tx\), and \(X(x)=x\).  Then
\[
 \langle X\rangle_t=\tanh(\beta t),
 \qquad
 \frac{\mathrm d}{\mathrm dt}\langle X\rangle_t
 =\beta\operatorname{sech}^2(\beta t).
\]
The direct derivative \(\langle\dot X\rangle_t\) is zero; the full variation
comes from the covariance in \eqref{eq:v2v10-expectation-derivative}.
\end{counterexample}

\begin{proof}
The two weighted masses are proportional to \(e^{\beta tx}\).  Direct
evaluation gives the displayed formulae.
\end{proof}

\subsection{One filtration formula for the third cumulant}
\label{sec:v2v10-martingale}

Use the finite filtration of V1 and write
\[
 \widehat F_m=\mathbb E[F\mid\mathcal F_m]-\mathbb EF,
\]
with analogous notation for \(G,H\).

\begin{theorem}[Maximum-index decomposition]
\label{thm:v2v10-third-martingale}
For \(F,G,H\in L^3\),
\begin{align}
 \kappa_3(F,G,H)
 =\sum_{m=1}^N\mathbb E\Bigl[
 &d_mF\,d_mG\,\widehat H_{m-1}
 +d_mF\,d_mH\,\widehat G_{m-1}\notag\\
 &+d_mG\,d_mH\,\widehat F_{m-1}
 +d_mF\,d_mG\,d_mH
 \Bigr].
\label{eq:v2v10-third-martingale}
\end{align}
Only terms whose largest filtration index occurs at least twice survive.
\end{theorem}

\begin{proof}
Expand each centered variable as
\[
 F-\mathbb EF=\sum_i d_iF,
\]
and similarly for \(G,H\).  If the largest index among \((i,j,k)\) occurs
only once, condition on the sigma algebra immediately below that index; the
unique largest martingale difference has conditional mean zero, while the
other two factors are measurable.  Such a term vanishes.

Fix the repeated largest index \(m\).  The cases \(i=j=m>k\),
\(i=k=m>j\), and \(j=k=m>i\) sum respectively to the first three terms in
\eqref{eq:v2v10-third-martingale}, because
\(\sum_{k<m}d_kH=\widehat H_{m-1}\), and similarly for the other variables.
The remaining case \(i=j=k=m\) is the fourth term.
\end{proof}

\begin{corollary}[A derivative-faithful cubic bound]
\label{cor:v2v10-third-bound}
Under the same hypotheses,
\begin{align}
 |\kappa_3(F,G,H)|
 \le\sum_m\Bigl(
 &\|d_mF\|_3\|d_mG\|_3\|\widehat H_{m-1}\|_3\notag\\
 +&\|d_mF\|_3\|d_mH\|_3\|\widehat G_{m-1}\|_3\notag\\
 +&\|d_mG\|_3\|d_mH\|_3\|\widehat F_{m-1}\|_3\notag\\
 +&\|d_mF\|_3\|d_mG\|_3\|d_mH\|_3
 \Bigr).
\label{eq:v2v10-third-bound}
\end{align}
\end{corollary}

\begin{proof}
Apply three-factor Holder inequality with exponent \(3\) to every term of
\eqref{eq:v2v10-third-martingale}.
\end{proof}

\begin{assessment}[Bridge-count consequence]
\label{ass:v2v10-bridge-count}
The third cumulant is not a product of three independently resampled
covariances.  Formula \eqref{eq:v2v10-third-martingale} ties the two or three
largest increments to the same filtration index.  Any argument assigning an
independent covariance weight to every displayed factor fails the
maximum-index test, exactly as the YM V6 independent-copy audit warns.
\end{assessment}

\subsection{Exact Gaussian quadratic third cumulant}
\label{sec:v2v10-gaussian}

Let \(\xi\) be standard Gaussian on a finite-dimensional real Hilbert space,
and retain the centered quadratic observable
\[
 \mathcal Q_A=\frac12(
 \langle\xi,A\xi\rangle-\operatorname{Tr}A)
\]
for self-adjoint \(A\).

\begin{theorem}[Symmetrized trace formula]
\label{thm:v2v10-gaussian-third}
For self-adjoint \(A,B,C\),
\begin{equation}
 \boxed{\quad
 \kappa_3(\mathcal Q_A,\mathcal Q_B,\mathcal Q_C)
 =\frac12\operatorname{Tr}(ABC+ACB).
 \quad}
\label{eq:v2v10-gaussian-third}
\end{equation}
Consequently,
\begin{align}
 |\kappa_3(\mathcal Q_A,\mathcal Q_B,\mathcal Q_C)|
 \le\frac12\|A\|_{\mathrm{HS}}\Bigl(
 \|B\|_{\mathrm{op}}\|C\|_{\mathrm{HS}}
 +\|C\|_{\mathrm{op}}\|B\|_{\mathrm{HS}}
 \Bigr).
\label{eq:v2v10-gaussian-third-bound}
\end{align}
\end{theorem}

\begin{proof}
For small \(t_1,t_2,t_3\), the centered joint cumulant generator is
\begin{align*}
 \log\mathbb E\exp\left(\sum_{j=1}^3t_j\mathcal Q_{A_j}\right)
 &=-\frac12\operatorname{Tr}\log
  \left(I-\sum_jt_jA_j\right)
  -\frac12\sum_jt_j\operatorname{Tr}A_j\\
 &=\frac12\sum_{m\ge2}\frac1m
  \operatorname{Tr}\left(\sum_jt_jA_j\right)^m.
\end{align*}
The coefficient of \(t_1t_2t_3\) comes only from \(m=3\).  The six
permutations split, by cyclicity of trace, into three copies of
\(\operatorname{Tr}(ABC)\) and three copies of
\(\operatorname{Tr}(ACB)\).  The coefficient is therefore the right side of
\eqref{eq:v2v10-gaussian-third}; differentiating once in each parameter
extracts that coefficient.

For the bound, use
\[
 |\operatorname{Tr}(ABC)|
 \le\|A\|_{\mathrm{HS}}\|BC\|_{\mathrm{HS}}
 \le\|A\|_{\mathrm{HS}}\|B\|_{\mathrm{op}}\|C\|_{\mathrm{HS}},
\]
and interchange \(B,C\) for the second trace.
\end{proof}

\subsection{A summable bosonic shell regression}
\label{sec:v2v10-bosonic-shell}

Let \(L\) be a positive elliptic operator of order two on a compact
\(d\)-manifold, let \(C_n\) be its covariance on a spectral shell of rank at
most \(W\Lambda_n^d\), and let
\[
 A_{j,n}=C_n^{1/2}M_{j,n}C_n^{1/2},
 \qquad j=1,2,3.
\]
Assume \(\|M_{j,n}\|_{\mathrm{op}}\le m_j\Lambda_n^{\mu_j}\).

\begin{theorem}[Cubic bosonic cutoff degree]
\label{thm:v2v10-bosonic-degree}
There is a constant \(C\) independent of \(n\) such that
\begin{equation}
 \left|
 \kappa_3(
 \mathcal Q_{A_{1,n}},
 \mathcal Q_{A_{2,n}},
 \mathcal Q_{A_{3,n}})
 \right|
 \le
 C\Lambda_n^{d+\mu_1+\mu_2+\mu_3-6}.
\label{eq:v2v10-bosonic-degree}
\end{equation}
In \(d=4\), three zero-order quadratic insertions give
\(O(\Lambda_n^{-2})\), so their dyadic third-cumulant shell is absolutely
summable without an additional Taylor subtraction.
\end{theorem}

\begin{proof}
The order-two covariance gives
\[
 \|A_{j,n}\|_{\mathrm{op}}
 \le C_j\Lambda_n^{\mu_j-2},
\qquad
 \|A_{j,n}\|_{\mathrm{HS}}
 \le C_j\Lambda_n^{d/2+\mu_j-2}.
\]
Insert these estimates in
\eqref{eq:v2v10-gaussian-third-bound}.  Each product has two
Hilbert--Schmidt factors and one operator-norm factor, hence cutoff power
\[
 (d/2+\mu_a-2)+(d/2+\mu_b-2)+(\mu_c-2)
 =d+\mu_1+\mu_2+\mu_3-6.
\]
For \(d=4\) and \(\mu_j=0\), this is \(-2\), and the dyadic sum converges.
\end{proof}

\begin{firewall}[Interacting interpolation remains open]
\label{fw:v2v10-interacting}
\cref{thm:v2v10-bosonic-degree} evaluates the Gaussian quadratic regression.
Along the interacting interpolation in
\eqref{eq:v2v10-adjacent-insertion}, the law need not remain Gaussian and
\(\Delta S_n\) contains higher-degree vertices.  The exact Duhamel and
martingale formulae remain valid, but their \(L^3\) and \(L^4\) norms require
uniform Gibbs or cluster estimates.  The Gaussian trace calculation does not
prove those estimates.
\end{firewall}

\subsection{V10 conclusion and next acquisition}
\label{sec:v2v10-conclusion}

The second adjacent term left open in V9 now has the exact representation
\eqref{eq:v2v10-adjacent-insertion}.  Its state variation is a third connected
cumulant, not a difference of separately normalized moments.  The
maximum-index martingale formula retains the common resampling, and the
Gaussian quadratic regression reduces the cubic cumulant to a symmetrized
trace.  In four dimensions, three zero-order bosonic quadratic insertions
already have the summable shell power \(\Lambda_n^{-2}\).

The next step is sector-specific.  One must derive the corresponding
finite-dimensional Grassmann third cumulant and its Dirac-shell degree, then
combine bosonic, fermionic, ghost, contact, and antifield contributions before
testing the local Slavnov class.  The next companion audit is due at V15.

\clearpage
\section*{Second Volume, Version V11: Grassmann Third Cumulants and the Cubic Dirac Shell}
\addcontentsline{toc}{section}{Second Volume, Version V11: Grassmann Third Cumulants and the Cubic Dirac Shell}

\subsection*{Purpose}

V10 showed that variation of a connected insertion along a normalized
interacting interpolation produces a third connected cumulant.  It evaluated
the bosonic Gaussian quadratic regression, whose three order-two covariances
are summable in four dimensions for zero-order insertions.  This version treats
the Berezin Gaussian sector.  The exact answer is a sum of two oriented cyclic
traces.  Its sign and cutoff degree differ from the bosonic answer, so the
calculation is performed independently.

\subsection{The Grassmann cumulant generator}
\label{sec:v2v11-generator}

Retain the finite-dimensional convention of V4: \(D\in GL(H)\),
\(C=D^{-1}\), and
\[
 \mathcal F_M=-\bar\psi M\psi-\operatorname{Tr}(CM).
\]
For matrices \(M_1,\ldots,M_s\), put
\[
 T(t)=\sum_{j=1}^s t_jM_j.
\]

\begin{lemma}[Centered determinant generator]
\label{lem:v2v11-generator}
For \(t\) near zero,
\begin{align}
 \log\left\langle
 \exp\left(\sum_jt_j\mathcal F_{M_j}\right)
 \right\rangle_D
 &=
 \operatorname{Tr}\log(I+CT(t))
 -\operatorname{Tr}(CT(t))\notag\\
 &=
 \sum_{m\ge2}\frac{(-1)^{m+1}}{m}
 \operatorname{Tr}\bigl(CT(t)\bigr)^m.
\label{eq:v2v11-generator}
\end{align}
\end{lemma}

\begin{proof}
The uncentered exponential changes the quadratic Berezin action from
\(\bar\psi D\psi\) to
\(\bar\psi(D+T(t))\psi\).  Therefore
\[
 \left\langle e^{-\bar\psi T(t)\psi}\right\rangle_D
 =\frac{\det(D+T(t))}{\det D}
 =\det(I+CT(t)).
\]
The scalar centering term contributes
\(\exp(-\operatorname{Tr}(CT(t)))\).  Taking the logarithm and using the
convergent finite-matrix power series for \(\log(I+CT(t))\) proves
\eqref{eq:v2v11-generator}.
\end{proof}

\begin{theorem}[General connected cyclic-word formula]
\label{thm:v2v11-general-cumulant}
For \(m\ge2\),
\begin{equation}
 \kappa_m(\mathcal F_{M_1},\ldots,\mathcal F_{M_m})
 =
 \frac{(-1)^{m+1}}{m}
 \sum_{\sigma\in S_m}
 \operatorname{Tr}\left(
 CM_{\sigma(1)}\cdots CM_{\sigma(m)}
 \right).
\label{eq:v2v11-general-cumulant}
\end{equation}
\end{theorem}

\begin{proof}
Differentiate \eqref{eq:v2v11-generator} once in every variable
\(t_1,\ldots,t_m\) at zero.  Only the term of degree \(m\) contributes.
Expansion of \(T(t)^m\) contains each word using all \(m\) labels once, indexed
by a permutation \(\sigma\in S_m\).  The derivative extracts their sum with
the coefficient in \eqref{eq:v2v11-general-cumulant}.
\end{proof}

\begin{corollary}[Checks at orders two and three]
\label{cor:v2v11-two-three}
For \(M,N,P\in\mathcal L(H)\),
\begin{align}
 \kappa_2(\mathcal F_M,\mathcal F_N)
 &=-\operatorname{Tr}(CMCN),
\label{eq:v2v11-second}\\
 \kappa_3(\mathcal F_M,\mathcal F_N,\mathcal F_P)
 &=\operatorname{Tr}(CMCNCP+CMCP CN).
\label{eq:v2v11-third}
\end{align}
\end{corollary}

\begin{proof}
For \(m=2\), the two permutations have equal trace by cyclicity, and the
coefficient is \(-1/2\).  For \(m=3\), the six permutations split into two
cyclic classes of three words each; the coefficient is \(1/3\).
\end{proof}

\begin{remark}[One-dimensional sign regression]
If \(H=\mathbb C\), \(C=c\), and \(M=N=P=1\), then
\(\kappa_3=2c^3\).  Directly,
\(\mathcal F_1=-\bar\psi\psi-c\) and
\((\bar\psi\psi)^2=0\), which gives the same value.  This fixes the cubic sign
independently of the determinant expansion.
\end{remark}

\subsection{Schatten bounds without positivity}
\label{sec:v2v11-schatten}

\begin{theorem}[Cubic trace-ideal estimate]
\label{thm:v2v11-cubic-schatten}
If the displayed products are trace class, then
\begin{align}
 |\kappa_3(\mathcal F_M,\mathcal F_N,\mathcal F_P)|
 \le{}&
 \|CM\|_{\mathrm{HS}}\|CN\|_{\mathrm{op}}\|CP\|_{\mathrm{HS}}
 \notag\\
 &+\|CM\|_{\mathrm{HS}}\|CP\|_{\mathrm{op}}\|CN\|_{\mathrm{HS}}.
\label{eq:v2v11-cubic-schatten}
\end{align}
Also,
\begin{equation}
 |\kappa_3(\mathcal F_M,\mathcal F_N,\mathcal F_P)|
 \le
 2\|CM\|_{\mathcal S_3}
  \|CN\|_{\mathcal S_3}
  \|CP\|_{\mathcal S_3}.
\label{eq:v2v11-cubic-s3}
\end{equation}
\end{theorem}

\begin{proof}
For the first cyclic word, apply Hilbert--Schmidt Cauchy--Schwarz to
\(CM\) and \(CNCP\), followed by
\[
 \|CNCP\|_{\mathrm{HS}}
 \le\|CN\|_{\mathrm{op}}\|CP\|_{\mathrm{HS}}.
\]
For the other word, interchange \(N\) and \(P\).  This proves
\eqref{eq:v2v11-cubic-schatten}.  Schatten Hölder with
\(1/3+1/3+1/3=1\) bounds each cyclic trace and proves
\eqref{eq:v2v11-cubic-s3}.
\end{proof}

\begin{firewall}[Cyclic words are assembled before the norm]
\label{fw:v2v11-cyclic-words}
The two traces in \eqref{eq:v2v11-third} can cancel after spin, chirality,
charge, and species traces are taken.  Equation
\eqref{eq:v2v11-cubic-schatten} is therefore a fallback absolute estimate,
not permission to norm each Standard--Model species or orientation before
forming the anomaly-sensitive sum.
\end{firewall}

\subsection{Relation to derivatives of the effective determinant}
\label{sec:v2v11-effective}

Let \(W(D)=-\log\det D\) as in V5.

\begin{proposition}[Third determinant derivative]
\label{prop:v2v11-third-determinant}
For three linear matrix directions \(M,N,P\),
\begin{equation}
 D^3W_D[M,N,P]
 =-\operatorname{Tr}(CMCNCP+CMCP CN)
 =-\kappa_3(\mathcal F_M,\mathcal F_N,\mathcal F_P).
\label{eq:v2v11-third-determinant}
\end{equation}
\end{proposition}

\begin{proof}
Differentiate
\(D^2W_D[M,N]=\operatorname{Tr}(CMCN)\) in direction \(P\).
Each of the two inverse factors has derivative \(-CPC\), giving
\[
 -\operatorname{Tr}(CPCMCN)
 -\operatorname{Tr}(CMCPCN).
\]
Cyclicity turns the first trace into
\(\operatorname{Tr}(CMCNCP)\).  Use
\eqref{eq:v2v11-third}.
\end{proof}

\begin{firewall}[Nonlinear field dependence creates contact terms]
\label{fw:v2v11-contact}
\cref{prop:v2v11-third-determinant} differentiates in linear matrix
directions.  If \(D\) depends nonlinearly on the metric, Higgs field, or gauge
coordinates, the third field derivative of \(W(D(\Phi))\) also contains
second- and third-derivative contact insertions of \(D(\Phi)\).  They are
required by the differentiated Ward identity and must be assembled before
using \eqref{eq:v2v11-cubic-schatten}.
\end{firewall}

\subsection{Elliptic shell power for a kinetic operator of order \(q\)}
\label{sec:v2v11-elliptic}

Let \(K\) be an invertible elliptic realization of positive order \(q\) on a
compact \(d\)-manifold.  Let \(E_n\) select a shell on which
\[
 \|E_nK^{-1}E_n\|_{\mathrm{op}}\le c_K\Lambda_n^{-q},
 \qquad
 \operatorname{rank}E_n\le W\Lambda_n^d.
\]
Put \(C_n=E_nK^{-1}E_n\).  Let the compressed insertions satisfy
\begin{equation}
 \|D_k^jM_{a,n}(k)\|_{\mathrm{op}}
 \le m_{a,j}\Lambda_n^{\mu_a-j},
 \qquad a=1,2,3.
\label{eq:v2v11-insertion-orders}
\end{equation}

\begin{theorem}[Cubic Grassmann shell degree]
\label{thm:v2v11-cubic-degree}
Define
\[
 \Gamma_n(k)=
 \kappa_3(
 \mathcal F_{M_{1,n}(k)},
 \mathcal F_{M_{2,n}(k)},
 \mathcal F_{M_{3,n}(k)}).
\]
For every \(0\le m\le r+1\),
\begin{equation}
 \|D_k^m\Gamma_n(k)\|_{\mathrm{op}}
 \le C_m
 \Lambda_n^{d+\mu_1+\mu_2+\mu_3-3q-m}.
\label{eq:v2v11-cubic-degree}
\end{equation}
After subtracting the degree-\(r\) Taylor jet,
\begin{equation}
 |(Q_r\Gamma_n)(k)|
 \le C_r'|k|^{r+1}
 \Lambda_n^{d+\mu_1+\mu_2+\mu_3-3q-(r+1)}.
\label{eq:v2v11-cubic-remainder}
\end{equation}
The dyadic remainder is absolutely summable if
\begin{equation}
 r+1>d+\mu_1+\mu_2+\mu_3-3q.
\label{eq:v2v11-cubic-threshold}
\end{equation}
\end{theorem}

\begin{proof}
For each differentiated insertion,
\begin{align*}
 \|C_nD_k^jM_{a,n}(k)\|_{\mathrm{op}}
 &\le C\Lambda_n^{\mu_a-q-j},\\
 \|C_nD_k^jM_{a,n}(k)\|_{\mathrm{HS}}
 &\le C\Lambda_n^{d/2+\mu_a-q-j}.
\end{align*}
Apply the multilinear Leibniz rule to the two cyclic traces in
\eqref{eq:v2v11-third}.  In every term, use two Hilbert--Schmidt factors and
one operator-norm factor as in \eqref{eq:v2v11-cubic-schatten}.  If the
derivative counts are \(j_1+j_2+j_3=m\), the cutoff power is
\[
 (d/2+\mu_a-q-j_a)
 +(d/2+\mu_b-q-j_b)
 +(\mu_c-q-j_c)
 =d+\sum_a\mu_a-3q-m.
\]
The finite multinomial sum changes only \(C_m\), proving
\eqref{eq:v2v11-cubic-degree}.  The integral Taylor remainder proves
\eqref{eq:v2v11-cubic-remainder}; its dyadic exponent is negative precisely
under \eqref{eq:v2v11-cubic-threshold}.
\end{proof}

\begin{corollary}[Dirac and ghost regressions in four dimensions]
\label{cor:v2v11-dirac-ghost}
Let \(d=4\) and \(\mu_1=\mu_2=\mu_3=0\).
\begin{enumerate}[label=\textup{(\roman*)}]
\item For a first-order Dirac shell, \(q=1\), the raw cubic degree is \(1\).
Subtraction through Taylor degree \(r=1\) gives a summable
\(O(|k|^2\Lambda_n^{-1})\) remainder.
\item For a second-order Faddeev--Popov ghost shell, \(q=2\), the raw cubic
degree is \(-2\), so the cubic shell is already summable.
\end{enumerate}
\end{corollary}

\begin{proof}
For \(q=1\), condition \eqref{eq:v2v11-cubic-threshold} is \(r+1>1\);
take \(r=1\).  The exponent in \eqref{eq:v2v11-cubic-remainder} is then
\(1-2=-1\).  For \(q=2\), the raw exponent is \(4-6=-2\).
\end{proof}

\begin{assessment}[Meaning of the linear Dirac jet]
\label{ass:v2v11-linear-jet}
The raw degree-one Dirac term is not an arbitrary counterterm.  In a
gauge-covariant assembled three-point function it belongs to the cubic part of
the local gauge kinetic functional, together with the quadratic polarization
jet isolated in V5 and the required contact terms.  Whether its
anomaly-sensitive chiral component cancels is a local Slavnov cohomology
question.  The power count only says that, after the admissible local jet is
removed, the remaining cubic shell is summable.
\end{assessment}

\subsection{Insertion into the V10 Duhamel formula}
\label{sec:v2v11-duhamel}

Suppose the interpolated shell law in
\eqref{eq:v2v10-adjacent-insertion} is a finite Berezin Gaussian and
\(F_t,G_t,\Delta S_n\) are quadratic with zero-order insertions.  Then
\eqref{eq:v2v11-third} evaluates its state-variation term exactly.  For a
four-dimensional Dirac shell, the local Taylor polynomial through degree one
must be allocated to the common gauge, Higgs, metric, boundary, and
antifield-source ledger before the \(\Lambda_n^{-1}\) remainder is summed.

\begin{firewall}[The interacting measure is still not Gaussian]
\label{fw:v2v11-interacting}
The actual interpolation contains higher vertices and mixed boson--fermion
terms.  The finite Gaussian formula is the leading regression and fixes the
cyclic signs and power count.  It does not replace the higher connected
cumulants generated by the interaction, nor does it prove a uniform cluster
bound for them.
\end{firewall}

\subsection{V11 conclusion and next step}
\label{sec:v2v11-conclusion}

The Grassmann third cumulant is now exact:
\[
 \kappa_3(\mathcal F_M,\mathcal F_N,\mathcal F_P)
 =\operatorname{Tr}(CMCNCP+CMCP CN).
\]
For an order-\(q\) elliptic inverse its raw shell degree is
\[
 d+\mu_1+\mu_2+\mu_3-3q.
\]
Thus the four-dimensional zero-order Dirac cubic needs its admissible linear
local jet removed, while the corresponding second-order ghost cubic is
already summable.  The cyclic and species sums must be formed before applying
Schatten norms.

The next step is to assemble the quadratic and cubic local Dirac jets as
derivatives of one gauge-invariant local functional.  This is necessary to
prevent independent counterterm choices from violating the V5 Ward row and
to expose the precise ghost-number-one anomaly class tested by the V3 repair
theorem.

\clearpage
\section*{Second Volume, Version V12: Integrability of the Quadratic, Cubic, and Quartic Gauge Jets}
\addcontentsline{toc}{section}{Second Volume, Version V12: Integrability of the Quadratic, Cubic, and Quartic Gauge Jets}

\subsection*{Purpose}

V5 isolated the transverse quadratic fermion jet, and V11 found the linear
external-momentum degree of the local fermion cubic.  Those jets cannot be
renormalized as unrelated tensors.  Gauge invariance requires them, together
with the quartic contact vertex, to be successive derivatives of one local
functional.  This version proves the finite polynomial descent equations that
test that integrability and verifies them for the gauge kinetic functional.

\subsection{Affine-linear gauge vector fields}
\label{sec:v2v12-affine-gauge}

Let \(\mathcal E\) be a finite-dimensional field space and \(\mathfrak g\) a
space of gauge parameters.  Suppose the infinitesimal gauge vector field has
the affine-linear form
\begin{equation}
 Q_\epsilon(A)=R_0\epsilon+R_1(A)\epsilon,
\label{eq:v2v12-gauge-vector}
\end{equation}
where \(R_1(A)\) is linear in \(A\).  Let
\[
 S(A)=\sum_{k=2}^KS_k(A)
\]
with \(S_k\) a homogeneous polynomial of degree \(k\).

\begin{theorem}[Homogeneous gauge-descent criterion]
\label{thm:v2v12-descent}
The polynomial \(S\) is invariant,
\begin{equation}
 DS(A)[Q_\epsilon(A)]=0
\quad\text{for all }A,\epsilon,
\label{eq:v2v12-invariance}
\end{equation}
if and only if
\begin{align}
 DS_2(A)[R_0\epsilon]&=0,
\label{eq:v2v12-descent-first}\\
 DS_{m+1}(A)[R_0\epsilon]
 +DS_m(A)[R_1(A)\epsilon]&=0,
 \quad 2\le m\le K-1,
\label{eq:v2v12-descent-middle}\\
 DS_K(A)[R_1(A)\epsilon]&=0.
\label{eq:v2v12-descent-last}
\end{align}
\end{theorem}

\begin{proof}
Since \(DS_k(A)\) is homogeneous of degree \(k-1\), the term
\(DS_k(A)[R_0\epsilon]\) has field degree \(k-1\), while
\(DS_k(A)[R_1(A)\epsilon]\) has field degree \(k\).  Expanding
\eqref{eq:v2v12-invariance} and collecting equal field degrees gives
\eqref{eq:v2v12-descent-first} at degree one,
\eqref{eq:v2v12-descent-middle} at degrees \(2,\ldots,K-1\), and
\eqref{eq:v2v12-descent-last} at degree \(K\).  Conversely, summing these
homogeneous identities makes every degree of
\(DS(A)[Q_\epsilon(A)]\) vanish.
\end{proof}

\begin{corollary}[The gauge quartic chain]
\label{cor:v2v12-quartic-chain}
For \(K=4\), invariance is equivalent to
\begin{align}
 DS_2[R_0\epsilon]&=0,\label{eq:v2v12-chain1}\\
 DS_3[R_0\epsilon]+DS_2[R_1(A)\epsilon]&=0,
 \label{eq:v2v12-chain2}\\
 DS_4[R_0\epsilon]+DS_3[R_1(A)\epsilon]&=0,
 \label{eq:v2v12-chain3}\\
 DS_4[R_1(A)\epsilon]&=0.
 \label{eq:v2v12-chain4}
\end{align}
Thus transversality of the quadratic jet is only the first equation; it does
not establish integrability of the cubic or quartic jets.
\end{corollary}

\begin{proof}
Specialize \cref{thm:v2v12-descent} to \(K=4\).
\end{proof}

\subsection{The gauge kinetic functional}
\label{sec:v2v12-yang-mills}

Let \(M\) be an oriented Riemannian four-manifold, let \(\mathfrak k\) be a
compact Lie algebra with an invariant inner product, and let \(A\) be a
\(\mathfrak k\)-valued one-form.  Write
\[
 F_A=\mathrm dA+\frac12[A\wedge A],
\qquad
 \delta_\epsilon A=\mathrm d\epsilon+[A,\epsilon].
\]
For one coefficient \(z\), define
\begin{equation}
 S_{\mathrm{kin},z}(A)
 =\frac z2\int_M\langle F_A\wedge *F_A\rangle.
\label{eq:v2v12-kinetic}
\end{equation}
Its homogeneous pieces are
\begin{align}
 S_{2,z}(A)
 &=\frac z2\int_M\langle\mathrm dA\wedge *\mathrm dA\rangle,
\label{eq:v2v12-S2}\\
 S_{3,z}(A)
 &=\frac z2\int_M
 \langle\mathrm dA\wedge *[A\wedge A]\rangle,
\label{eq:v2v12-S3}\\
 S_{4,z}(A)
 &=\frac z8\int_M
 \langle[A\wedge A]\wedge *[A\wedge A]\rangle.
\label{eq:v2v12-S4}
\end{align}

\begin{theorem}[One coefficient satisfies the complete descent]
\label{thm:v2v12-kinetic-descent}
The functional \eqref{eq:v2v12-kinetic} is gauge invariant, and its pieces
\eqref{eq:v2v12-S2}--\eqref{eq:v2v12-S4} satisfy all four equations in
\cref{cor:v2v12-quartic-chain} with
\[
 R_0\epsilon=\mathrm d\epsilon,
 \qquad
 R_1(A)\epsilon=[A,\epsilon].
\]
\end{theorem}

\begin{proof}
The curvature varies covariantly:
\begin{align*}
 \delta_\epsilon F_A
 &=\mathrm d(\mathrm d\epsilon+[A,\epsilon])
   +[A\wedge(\mathrm d\epsilon+[A,\epsilon])]\\
 &=[F_A,\epsilon].
\end{align*}
Here \(\mathrm d^2=0\), the graded Leibniz rule combines the
\(\mathrm d\epsilon\) terms, and the Jacobi identity combines the quadratic
terms.  Invariance of the Lie-algebra inner product gives pointwise
\[
 \langle[F_A,\epsilon],F_A\rangle
 +\langle F_A,[F_A,\epsilon]\rangle=0.
\]
Therefore \(\delta_\epsilon S_{\mathrm{kin},z}=0\).  The expressions
\eqref{eq:v2v12-S2}--\eqref{eq:v2v12-S4} are exactly the homogeneous expansion
of \eqref{eq:v2v12-kinetic}; apply
\cref{thm:v2v12-descent}.
\end{proof}

\begin{remark}[No integration-by-parts boundary assumption]
The proof uses the pointwise covariance \(\delta_\epsilon F_A=[F_A,\epsilon]\)
and the invariant inner product.  It does not discard a boundary integral.
Boundary conditions are still required for the kinetic operator and BV
domain, but not for this algebraic local descent check.
\end{remark}

\subsection{Independent coefficients produce an exact defect}
\label{sec:v2v12-mismatch}

Let
\[
 \widetilde S(A)=z_2S_{2,1}(A)+z_3S_{3,1}(A)+z_4S_{4,1}(A),
\]
where \(S_{k,1}\) denotes \eqref{eq:v2v12-S2}--\eqref{eq:v2v12-S4} with
unit coefficient.

\begin{proposition}[Coefficient-mismatch Ward polynomial]
\label{prop:v2v12-mismatch}
The gauge variation of \(\widetilde S\) is
\begin{align}
 D\widetilde S(A)[Q_\epsilon(A)]
 ={}&(z_2-z_3)\,
 DS_{2,1}(A)[R_1(A)\epsilon]\notag\\
 &+(z_3-z_4)\,
 DS_{3,1}(A)[R_1(A)\epsilon].
\label{eq:v2v12-mismatch}
\end{align}
Hence \(z_2=z_3=z_4\) is sufficient for gauge invariance.  Whenever the two
displayed variation polynomials are linearly independent, it is also
necessary.
\end{proposition}

\begin{proof}
The unit-coefficient descent equations give
\[
 DS_{2,1}[R_0\epsilon]=0,\qquad
 DS_{3,1}[R_0\epsilon]=-DS_{2,1}[R_1\epsilon],
\]
\[
 DS_{4,1}[R_0\epsilon]=-DS_{3,1}[R_1\epsilon],\qquad
 DS_{4,1}[R_1\epsilon]=0.
\]
Insert these identities into the homogeneous expansion of
\(D\widetilde S[Q_\epsilon]\).  The degree-two and degree-three terms are
exactly \eqref{eq:v2v12-mismatch}.  Equality of the three coefficients makes
both vanish.  Linear independence gives the converse.
\end{proof}

\begin{counterexample}[Transversality alone permits a mismatched cubic jet]
\label{sep:v2v12-transverse-only}
Choose \(z_2\ne z_3\) and \(z_4=z_3\).  The quadratic piece remains transverse
because \(DS_{2,1}[R_0\epsilon]=0\), but
\[
 D\widetilde S[Q_\epsilon]
 =(z_2-z_3)DS_{2,1}[R_1(A)\epsilon],
\]
which is generally nonzero in a nonabelian theory.  Thus the V5 transverse
two-point condition does not authorize an independently normalized V11 cubic
counterterm.
\end{counterexample}

\begin{proof}
This is \eqref{eq:v2v12-mismatch} with \(z_4=z_3\).  Nonvanishing occurs
whenever the quadratic kinetic energy changes under the nonlinear part
\([A,\epsilon]\), which is precisely the term cancelled by the cubic vertex
in the full curvature square.
\end{proof}

\subsection{Local jet integrability test}
\label{sec:v2v12-integrability}

Let \(J_{2,n},J_{3,n},J_{4,n}\) be the local quadratic, cubic, and quartic
jets extracted from one complete shell effective action.  Define their
descent defects
\begin{align}
 \alpha_{1,n}
 &=DJ_{2,n}[R_0\epsilon],
\label{eq:v2v12-alpha1}\\
 \alpha_{2,n}
 &=DJ_{3,n}[R_0\epsilon]+DJ_{2,n}[R_1(A)\epsilon],
\label{eq:v2v12-alpha2}\\
 \alpha_{3,n}
 &=DJ_{4,n}[R_0\epsilon]+DJ_{3,n}[R_1(A)\epsilon],
\label{eq:v2v12-alpha3}\\
 \alpha_{4,n}
 &=DJ_{4,n}[R_1(A)\epsilon].
\label{eq:v2v12-alpha4}
\end{align}

\begin{theorem}[Finite local integrability criterion]
\label{thm:v2v12-integrability}
The polynomial
\[
 J_n=J_{2,n}+J_{3,n}+J_{4,n}
\]
is invariant under \eqref{eq:v2v12-gauge-vector} if and only if
\[
 \alpha_{1,n}=\alpha_{2,n}=\alpha_{3,n}=\alpha_{4,n}=0.
\]
If these identities hold and \(J_n=z_nS_{\mathrm{kin},1}\), then all three
vertices are transported by the single running coefficient \(z_n\).
\end{theorem}

\begin{proof}
The equivalence is \cref{cor:v2v12-quartic-chain}.  The last statement follows
by differentiating the one functional \(z_nS_{\mathrm{kin},1}\): its
homogeneous terms are \(z_nS_{2,1},z_nS_{3,1},z_nS_{4,1}\).
\end{proof}

\begin{firewall}[Vanishing descent defects do not classify all local actions]
\label{fw:v2v12-classification}
\cref{thm:v2v12-integrability} tests invariance of a given polynomial.  It
does not prove that every admissible local Standard--Model counterterm is a
multiple of \(F_A^2\).  Higgs, Yukawa, gravitational, topological, boundary,
gauge-fixing, ghost, and antifield terms form additional local cohomology
classes.  A complete projector must use a basis of the full local BV
cohomology.
\end{firewall}

\subsection{Interface with the fermion shell}
\label{sec:v2v12-fermion-interface}

The quadratic tensor
\[
 a_n(|k|^2I-k\otimes k)
\]
from V5 is the Fourier Hessian of the quadratic part of a local gauge kinetic
term.  The degree-one cubic Dirac jet from V11 must be its third derivative,
and the differentiated determinant Ward row supplies the corresponding
contact terms.  If the three extracted coefficients disagree,
\eqref{eq:v2v12-mismatch} is the explicit local ghost-number-one breaking.
It must either cancel in the complete Standard--Model species sum or be
removed by a local primitive as in V3.

\begin{assessment}[Required order of operations]
\label{ass:v2v12-order}
For each shell, first assemble fermion orientations, species, ghosts, gauge
vertices, and contact insertions.  Next extract the complete local polynomial
\(J_{2,n}+J_{3,n}+J_{4,n}\).  Test
\eqref{eq:v2v12-alpha1}--\eqref{eq:v2v12-alpha4}.  Only after this test may the
common invariant polynomial be moved into the running-coupling ledger and the
irrelevant complement be estimated by V5 and V11.  Projecting the two- and
three-point functions independently can manufacture the mismatch
\eqref{eq:v2v12-mismatch}.
\end{assessment}

\subsection{V12 conclusion and next step}
\label{sec:v2v12-conclusion}

The admissible gauge local jet is now constrained by an exact finite
polynomial descent.  For the gauge kinetic class, the transverse quadratic,
linear-momentum cubic, and quartic contact vertices are the three homogeneous
pieces of one curvature-square functional and carry one coefficient.  A
coefficient mismatch has the explicit Ward polynomial
\eqref{eq:v2v12-mismatch}.

The next step is to perform the analogous integrability calculation for the
nonminimal Higgs--Einstein local sector.  Its metric, Higgs, and boundary
vertices must arise from the single scalar--tensor functional
\(\int\sqrt g\,F(H^\dagger H)R\) with its boundary partner.  This will connect
the shell-local source analysis back to the Einstein face momentum and
constraint propagation gates summarized in V1.

\clearpage
\section*{Second Volume, Version V13: Variational Integrability of the Higgs--Einstein Bulk and Face Sources}
\addcontentsline{toc}{section}{Second Volume, Version V13: Variational Integrability of the Higgs--Einstein Bulk and Face Sources}

\subsection*{Purpose}

V12 proved that gauge two-, three-, and four-point local jets must be
derivatives of one gauge-invariant functional.  The same issue occurs in the
nonminimal Higgs--Einstein sector.  A metric source, a Higgs curvature term,
and a boundary face momentum cannot be renormalized independently if they are
to arise from one effective action.  This version formulates the exact
finite-cutoff Helmholtz condition for that common origin and applies it to the
scalar--tensor curvature functional with its boundary partner.

\subsection{A finite-cutoff Helmholtz theorem}
\label{sec:v2v13-helmholtz}

Let \(X\) be a finite-dimensional real vector space of cutoff field
coefficients, including bulk and boundary variables, and let
\(U\subset X\) be open and star-shaped about the origin.  A \(C^1\) source
one-form is a map
\[
 \Theta:U\longrightarrow X^*.
\]
Its derivative is the bilinear form
\[
 D\Theta_x(u,v):=(D\Theta_x[u])(v).
\]

\begin{theorem}[Finite-cutoff variational integrability]
\label{thm:v2v13-helmholtz}
There is a \(C^2\) functional \(\Gamma:U\to\mathbb R\) with
\[
 D\Gamma_x=\Theta_x
\]
if and only if
\begin{equation}
 D\Theta_x(u,v)=D\Theta_x(v,u)
\quad\text{for every }x\in U,\ u,v\in X.
\label{eq:v2v13-closed-one-form}
\end{equation}
When \eqref{eq:v2v13-closed-one-form} holds, one may take
\begin{equation}
 \Gamma(x)=\Gamma(0)+\int_0^1\Theta_{tx}(x)\,\mathrm dt.
\label{eq:v2v13-radial-action}
\end{equation}
\end{theorem}

\begin{proof}
If \(\Theta=D\Gamma\), then
\[
 D\Theta_x(u,v)=D^2\Gamma_x(u,v)
\]
is symmetric by equality of mixed derivatives.

Conversely, define \(\Gamma\) by \eqref{eq:v2v13-radial-action}.  For
\(h\in X\), differentiation under the finite integral gives
\[
 D\Gamma_x[h]
 =\int_0^1\left(
 tD\Theta_{tx}(h,x)+\Theta_{tx}(h)\right)\mathrm dt.
\]
Use \eqref{eq:v2v13-closed-one-form} to replace
\(D\Theta_{tx}(h,x)\) by \(D\Theta_{tx}(x,h)\).  The integrand becomes
\[
 \frac{\mathrm d}{\mathrm dt}\bigl(t\Theta_{tx}(h)\bigr).
\]
Integration from zero to one gives \(D\Gamma_x[h]=\Theta_x(h)\).
\end{proof}

\subsection{Bulk, Higgs, and face blocks}
\label{sec:v2v13-blocks}

Split the cutoff coordinates as
\[
 X=X_g\oplus X_H\oplus X_\gamma,
\]
where \(X_g\) contains interior metric coefficients, \(X_H\) the Higgs
coefficients, and \(X_\gamma\) the induced boundary metric coefficients.
Write the one-form as
\begin{equation}
 \Theta_{(g,H,\gamma)}(h,\eta,k)
 =\mathcal E_g(h)+\mathcal E_H(\eta)+\mathcal P_\gamma(k).
\label{eq:v2v13-source-one-form}
\end{equation}

\begin{corollary}[Mixed source tests]
\label{cor:v2v13-mixed-tests}
The source triple in \eqref{eq:v2v13-source-one-form} comes from one cutoff
effective action if and only if all diagonal derivatives are symmetric and
\begin{align}
 D_H\mathcal E_g[\eta](h)
 &=D_g\mathcal E_H[h](\eta),
\label{eq:v2v13-gH-test}\\
 D_\gamma\mathcal E_g[k](h)
 &=D_g\mathcal P_\gamma[h](k),
\label{eq:v2v13-ggamma-test}\\
 D_\gamma\mathcal E_H[k](\eta)
 &=D_H\mathcal P_\gamma[\eta](k).
\label{eq:v2v13-Hgamma-test}
\end{align}
\end{corollary}

\begin{proof}
Apply \cref{thm:v2v13-helmholtz} to pairs of vectors supported in each of the
three direct summands.  These pairs give the diagonal symmetries and the three
displayed mixed identities, which together exhaust all block pairs.
\end{proof}

\begin{counterexample}[Separate Ward identities do not imply integrability]
\label{sep:v2v13-separate-ward}
Take \(X_g=X_H=\mathbb R\) and
\[
 \mathcal E_g(h)=aHh,\qquad
 \mathcal E_H(\eta)=bg\eta
\]
with constants \(a\ne b\).  Each component is a valid linear covector in its
own variable, but no common \(\Gamma(g,H)\) has these two derivatives.
\end{counterexample}

\begin{proof}
The mixed derivatives are
\[
 D_H\mathcal E_g[\eta](h)=a\eta h,\qquad
 D_g\mathcal E_H[h](\eta)=bh\eta.
\]
They violate \eqref{eq:v2v13-gH-test}.  Equivalently, integrating the first
source would require the mixed term \(agH\), while the second would require
\(bgH\).
\end{proof}

\subsection{The scalar--tensor functional and its face partner}
\label{sec:v2v13-scalar-tensor}

Let \(f=H^\dagger H\) and
\[
 F(f)=\chi_R+\xi_Hf.
\]
In the sign convention fixed in the first volume, consider
\begin{equation}
 S_{F}[g,H]
 =\int_M\sqrt{|g|}\,F(f)R_g
 +2\int_{\partial M}\sqrt{|\gamma|}\,F(f)K.
\label{eq:v2v13-scalar-tensor-action}
\end{equation}
The second term is the scalar--tensor Gibbons--Hawking partner.  Denote the
complete first variation of \eqref{eq:v2v13-scalar-tensor-action} by
\(\Theta_F\).

\begin{theorem}[Common origin of the improvement and face momentum]
\label{thm:v2v13-common-origin}
After the normal and induced-metric conventions of the first volume are
fixed, \(\Theta_F\) is a closed field-space one-form.  Its interior metric
component contains
\begin{equation}
 F G^{\mu\nu}
 +(g^{\mu\nu}\Box-\nabla^\mu\nabla^\nu)F,
\label{eq:v2v13-metric-gradient}
\end{equation}
its Higgs component contains
\begin{equation}
 R_g\,D_HF,
\label{eq:v2v13-Higgs-gradient}
\end{equation}
and its face component is
\begin{equation}
 \mathcal P_F^{ab}
 =F\Pi^{ab}-\gamma^{ab}\nabla_nF.
\label{eq:v2v13-face-gradient}
\end{equation}
For \(F=\chi_R+\xi_Hf\), the same coefficient \(\xi_H\) therefore multiplies
the bulk improvement, the Higgs curvature source, and the two \(f\)-dependent
face terms.
\end{theorem}

\begin{proof}
The Palatini identity writes the variation of
\(\sqrt{|g|}F R_g\) as the interior expression
\eqref{eq:v2v13-metric-gradient} paired with the metric variation, plus
\(R_g\,\delta F\), plus a normal derivative of the metric variation.  The
variation of \(2\sqrt{|\gamma|}FK\) cancels that normal derivative.  The
remaining induced-metric coefficient is
\eqref{eq:v2v13-face-gradient}; the normal derivative of \(F\) appears when
the derivatives in the Palatini term are transferred from the metric
variation.  Since \(\delta F=D_HF[\eta]\), the bulk Higgs coefficient is
\eqref{eq:v2v13-Higgs-gradient}.

These three terms are, by construction, the components of
\(DS_F\).  Hence their one-form is closed by the necessary direction of
\cref{thm:v2v13-helmholtz}.  Finally,
\[
 D_HF=\xi_HD_Hf,\qquad
 \nabla_nF=\xi_H\nabla_nf,
\]
while the \(f\)-dependent part of \(FG^{\mu\nu}\) is
\(\xi_HfG^{\mu\nu}\).  This proves the common-coefficient statement.
\end{proof}

\begin{remark}[Stress-tensor convention]
If the stress tensor is defined by minus twice the metric gradient, then the
\(\xi_H\)-dependent part of \eqref{eq:v2v13-metric-gradient} becomes the
improved stress recorded in \eqref{eq:v2-improved-stress}.  The present
theorem is stated at the action-gradient level to avoid mixing inverse- and
covariant-metric sign conventions.
\end{remark}

\subsection{A coefficient mismatch is a Helmholtz defect}
\label{sec:v2v13-mismatch}

Let
\[
 \Theta_1=DS_{fR+2fK}
\]
denote the one-form obtained from \eqref{eq:v2v13-scalar-tensor-action} with
\(\chi_R=0\), \(\xi_H=1\).  Decompose it into its three blocks
\[
 \Theta_1=\Theta_{1,g}+\Theta_{1,H}+\Theta_{1,\gamma}.
\]
For constants \(a,b,c\), define
\[
 \widetilde\Theta
 =a\Theta_{1,g}+b\Theta_{1,H}+c\Theta_{1,\gamma}.
\]

\begin{proposition}[Mixed-coefficient obstruction]
\label{prop:v2v13-mismatch}
The mixed exterior derivatives of \(\widetilde\Theta\) are
\begin{align}
 D_H\widetilde\Theta_g-D_g\widetilde\Theta_H
 &=(a-b)D_H\Theta_{1,g},
\label{eq:v2v13-mismatch-gH}\\
 D_\gamma\widetilde\Theta_g-D_g\widetilde\Theta_\gamma
 &=(a-c)D_\gamma\Theta_{1,g},
\label{eq:v2v13-mismatch-ggamma}\\
 D_\gamma\widetilde\Theta_H-D_H\widetilde\Theta_\gamma
 &=(b-c)D_\gamma\Theta_{1,H}.
\label{eq:v2v13-mismatch-Hgamma}
\end{align}
Here equality identifies each mixed Hessian with its transpose.  If the
corresponding mixed Hessians are nonzero, variational integrability forces
\(a=b=c\).
\end{proposition}

\begin{proof}
Because \(\Theta_1=DS_{fR+2fK}\), its mixed derivatives agree:
\[
 D_H\Theta_{1,g}=D_g\Theta_{1,H},
\]
and similarly for the other block pairs.  Multiplying the two sides by
independent coefficients and subtracting gives
\eqref{eq:v2v13-mismatch-gH}--\eqref{eq:v2v13-mismatch-Hgamma}.
The conclusion follows from \cref{cor:v2v13-mixed-tests}.
\end{proof}

\begin{assessment}[Shell-local integrability test]
\label{ass:v2v13-shell-test}
At shell \(n\), let the extracted local metric, Higgs, and face sources define
\(\Theta_n^{\mathrm{loc}}\).  Before placing their coefficients in separate
running ledgers, test the mixed identities
\eqref{eq:v2v13-gH-test}--\eqref{eq:v2v13-Hgamma-test}.  If their
\(fR+2fK\) components have coefficients \(a_n,b_n,c_n\), then
\cref{prop:v2v13-mismatch} requires equality on every nondegenerate block.
Their common value is the shell increment of \(\xi_H\).
\end{assessment}

\begin{firewall}[Closure is not coercivity]
\label{fw:v2v13-coercivity}
Variational integrability proves that the source triple comes from one local
action.  It does not prove positivity of
\(F=\chi_R+\xi_Hf\), coercivity of the Einstein-frame target metric, or
compactness of the Higgs stress.  Those remain the analytic conditions
summarized in V1.
\end{firewall}

\subsection{Constraint propagation interface}
\label{sec:v2v13-constraint}

If a cutoff effective action is diffeomorphism invariant, its metric and
matter gradients obey one Noether identity.  Extracting a closed common local
one-form preserves the possibility of that identity; assigning inconsistent
bulk and face coefficients creates a local residual
\(\mathcal R_\nu\) in the subsidiary equation
\eqref{eq:v2-subsidiary}.  Thus the Helmholtz test precedes the strong source
convergence required for Einstein stability.

\begin{assessment}[New order of acquisition]
\label{ass:v2v13-order}
The source route now has four distinct checks:
\begin{enumerate}[label=\textup{(\roman*)}]
\item form the complete connected shell source;
\item extract a local one-form and verify Slavnov/diffeomorphism closure;
\item verify variational integrability across metric, Higgs, and boundary
blocks by \cref{cor:v2v13-mixed-tests}; and
\item prove uniform integrability and strong enough convergence of the
irrelevant remainder.
\end{enumerate}
Passing the first three does not supply the fourth, but failing any one of
them prevents the source from entering a consistent Einstein evolution.
\end{assessment}

\subsection{V13 conclusion and next step}
\label{sec:v2v13-conclusion}

The nonminimal local source is now governed by a finite-cutoff Helmholtz
criterion.  The bulk improvement, Higgs curvature equation, and face momentum
are the three blocks of the differential of
\[
 \int_M\sqrt{|g|}F(f)R_g
 +2\int_{\partial M}\sqrt{|\gamma|}F(f)K.
\]
Consequently, the \(fR\) coefficient cannot run independently in the three
blocks; a mismatch is the explicit mixed derivative defect
\eqref{eq:v2v13-mismatch-gH}--\eqref{eq:v2v13-mismatch-Hgamma}.

The next step is analytic: prove that the shellwise common coefficient and
the irrelevant metric-source remainder define a Cauchy sequence in a source
topology strong enough for the harmonic-gauge Einstein stability estimate.
The concentration defect from V151--V153 remains the obstruction that a mere
action-level integrability result cannot remove.

\clearpage
\section*{Second Volume, Version V14: A Strong Source-Cauchy Criterion for Harmonic Einstein Stability}
\addcontentsline{toc}{section}{Second Volume, Version V14: A Strong Source-Cauchy Criterion for Harmonic Einstein Stability}

\subsection*{Purpose and claim boundary}

V13 proved variational compatibility of the local Higgs--Einstein source
blocks.  It did not prove convergence of those sources or stability of the
Einstein evolution.  This version supplies a high-regularity sufficient
criterion.  If the irrelevant shell sources are absolutely summable in
\(L^1_tH^{s-1}_x\), and the finite local coupling vector converges, then the
complete source is Cauchy in the norm required by a standard harmonic-gauge
energy estimate.  Under uniform hyperbolicity and Sobolev bounds, the
corresponding reduced metrics are Cauchy.  These are conditional analytic
theorems; the quantum shell bounds needed to invoke them remain open.

\subsection{Local and irrelevant source decomposition}
\label{sec:v2v14-source-decomposition}

Fix a compact \(d\)-dimensional spatial slice \(\Sigma\), a time interval
\([0,T]\), and \(s>d/2+2\).  Let
\[
 \mathcal X_s=L^1([0,T];H^{s-1}(\Sigma)).
\]
At cutoff \(N\), write the renormalized source as
\begin{equation}
 T_N=\mathcal L(c_N)+\sum_{n=n_0}^N\tau_n,
\label{eq:v2v14-source-decomposition}
\end{equation}
where \(c_N\in\mathbb R^J\) is the vector of local running couplings,
\(\mathcal L:\mathbb R^J\to\mathcal X_s\) is the fixed local-source map, and
\(\tau_n\) is the invariant irrelevant shell source.

\begin{theorem}[Source Cauchy criterion]
\label{thm:v2v14-source-cauchy}
Assume
\begin{align}
 c_N&\longrightarrow c_\infty\quad\text{in }\mathbb R^J,
\label{eq:v2v14-coupling-convergence}\\
 \sum_{n=n_0}^\infty\|\tau_n\|_{\mathcal X_s}&<\infty,
\label{eq:v2v14-shell-summability}
\end{align}
and let \(\mathcal L\) be continuous.  Then \(T_N\) is Cauchy in
\(\mathcal X_s\) and
\begin{equation}
 T_N\longrightarrow
 T_\infty=\mathcal L(c_\infty)+\sum_{n=n_0}^\infty\tau_n
\quad\text{in }\mathcal X_s.
\label{eq:v2v14-source-limit}
\end{equation}
\end{theorem}

\begin{proof}
For \(M>N\),
\[
 \|T_M-T_N\|_{\mathcal X_s}
 \le
 \|\mathcal L(c_M-c_N)\|_{\mathcal X_s}
 +\sum_{n=N+1}^M\|\tau_n\|_{\mathcal X_s}.
\]
The first term tends to zero by
\eqref{eq:v2v14-coupling-convergence} and continuity.  The second is a tail
of the convergent nonnegative series
\eqref{eq:v2v14-shell-summability}.  Completeness of \(\mathcal X_s\) gives
\eqref{eq:v2v14-source-limit}.
\end{proof}

\begin{firewall}[Local projection does not prove coupling convergence]
\label{fw:v2v14-local-couplings}
The local coefficients removed in V5, V11, V12, and V13 can be marginal and
may accumulate logarithmically.  Equation
\eqref{eq:v2v14-coupling-convergence} is a renormalization condition on the
running trajectory, not a consequence of irrelevant-shell summability.  It
must be proved after fixing physical renormalization data.
\end{firewall}

\subsection{A dyadic \(L^2\)-to-source criterion}
\label{sec:v2v14-dyadic}

\begin{proposition}[Frequency-localized sufficient bound]
\label{prop:v2v14-dyadic}
Suppose \(\tau_n(t)\) has spatial Fourier or spectral support in
\[
 \{\,\xi:c_-\Lambda_n\le\langle\xi\rangle\le c_+\Lambda_n\,\}
\]
and, for some \(\varepsilon>0\),
\begin{equation}
 \|\tau_n\|_{L^1_tL^2_x}
 \le A\Lambda_n^{-(s-1)-\varepsilon}.
\label{eq:v2v14-L2-shell}
\end{equation}
Then
\begin{equation}
 \|\tau_n\|_{\mathcal X_s}
 \le C A\Lambda_n^{-\varepsilon},
\label{eq:v2v14-Hs-shell}
\end{equation}
and \eqref{eq:v2v14-shell-summability} holds for a dyadic sequence
\(\Lambda_n=2^n\Lambda_0\).
\end{proposition}

\begin{proof}
The spectral support and the definition of the Sobolev norm give, pointwise
in time,
\[
 \|\tau_n(t)\|_{H^{s-1}}
 \le C\Lambda_n^{s-1}\|\tau_n(t)\|_{L^2}.
\]
Integrate in time and use \eqref{eq:v2v14-L2-shell} to obtain
\eqref{eq:v2v14-Hs-shell}.  The series
\(\sum_n\Lambda_n^{-\varepsilon}\) is geometric.
\end{proof}

\begin{counterexample}[Weak convergence is not the source-Cauchy norm]
\label{sep:v2v14-weak-source}
Let \(\phi\) be a nonzero smooth compactly supported function on a coordinate
patch of \(\Sigma\), and set
\[
 \tau_n(x)=\Lambda_n^{d/2}\phi(\Lambda_n x).
\]
Then \(\tau_n\rightharpoonup0\) distributionally and
\(\|\tau_n\|_{L^2}\) is constant, while
\[
 \|\tau_n\|_{H^{s-1}}\asymp\Lambda_n^{s-1}.
\]
Thus distributional or weak \(L^2\) convergence does not imply the norm
needed in \cref{thm:v2v14-source-cauchy}.
\end{counterexample}

\begin{proof}
The distributional convergence follows after the substitution
\(y=\Lambda_nx\), since the pairing has the factor
\(\Lambda_n^{-d/2}\).  The \(L^2\) norm is scale invariant.  Differentiating
\(s-1\) times, or using the Fourier transform, gives the Sobolev scaling.
\end{proof}

\subsection{A reduced hyperbolic difference estimate}
\label{sec:v2v14-hyperbolic}

Let \(u_N\) denote the coordinate components of a reduced metric and matter
field in harmonic and hyperbolic matter gauges.  On a common time slab,
suppose they solve
\begin{equation}
 A^{\alpha\beta}(u_N)\partial_\alpha\partial_\beta u_N
 =\mathcal N(u_N,\partial u_N)+T_N.
\label{eq:v2v14-reduced-system}
\end{equation}

\begin{theorem}[Conditional cutoff stability]
\label{thm:v2v14-hyperbolic-stability}
Assume:
\begin{enumerate}[label=\textup{(H\arabic*)}]
\item the matrices \(A^{\alpha\beta}(u_N)\) are uniformly strictly hyperbolic
on \([0,T]\), with fixed uniformly energy-admissible boundary conditions when
\(\partial\Sigma\ne\varnothing\);
\item
\[
 \sup_N\left(
 \|u_N\|_{L^\infty_tH^{s+1}_x}
 +\|\partial_tu_N\|_{L^\infty_tH^s_x}\right)\le M;
\]
\item \(A\) and \(\mathcal N\) are smooth on the range selected by this bound;
\item the initial data are Cauchy in
\(H^s\times H^{s-1}\); and
\item \(T_N\) is Cauchy in \(\mathcal X_s\).
\end{enumerate}
Then \(u_N\) is Cauchy in
\begin{equation}
 C([0,T];H^s)\cap C^1([0,T];H^{s-1}).
\label{eq:v2v14-solution-space}
\end{equation}
\end{theorem}

\begin{proof}
For \(w=u_N-u_M\), subtract the two equations and place the principal
operator with coefficient \(A(u_N)\) on the left:
\begin{align*}
 A^{\alpha\beta}(u_N)\partial_\alpha\partial_\beta w
 ={}&T_N-T_M\\
 &+\mathcal N(u_N,\partial u_N)
  -\mathcal N(u_M,\partial u_M)\\
 &+\bigl(A^{\alpha\beta}(u_M)-A^{\alpha\beta}(u_N)\bigr)
   \partial_\alpha\partial_\beta u_M.
\end{align*}
The hyperbolic energy estimate for \(w\) at solution order \(s\), combined
with the uniform order-\(s+1\) background bound and Sobolev Moser estimates
valid for \(s>d/2+2\), gives
\begin{align}
 E_s(w;t)
 \le C_M\Bigl(
 E_s(w;0)
 +\|T_N-T_M\|_{L^1([0,t];H^{s-1})}
 +\int_0^tE_s(w;\sigma)\,\mathrm d\sigma
 \Bigr),
\label{eq:v2v14-energy-difference}
\end{align}
where \(E_s\) is equivalent, uniformly in \(N\), to
\[
 \|w(t)\|_{H^s}+\|\partial_tw(t)\|_{H^{s-1}}.
\]
The constants depend only on the hyperbolicity margin, \(M\), and the smooth
coefficient functions.  Gronwall's inequality yields
\[
 \sup_{t\le T}E_s(w;t)
 \le C_{M,T}\left(
 E_s(w;0)+\|T_N-T_M\|_{\mathcal X_s}\right).
\]
Both terms tend to zero by (H4)--(H5), proving the Cauchy statement.
\end{proof}

\begin{firewall}[Uniform existence is an independent hypothesis]
\label{fw:v2v14-uniform-existence}
\cref{thm:v2v14-hyperbolic-stability} compares solutions already known to
exist on one common slab with a uniform Sobolev bound and hyperbolicity
margin.  It does not prove such a slab from the quantum source estimates.
Closing the coupled continuum requires a cutoff-independent local existence
time or a continuation estimate.
\end{firewall}

\subsection{Subsidiary constraint convergence}
\label{sec:v2v14-subsidiary}

Let \(C_{N,\nu}\) be the harmonic constraint field.  In the conventions of
V1 it obeys
\begin{equation}
 \Box_{g_N}C_{N,\nu}
 +(R_{g_N})_\nu{}^\mu C_{N,\mu}
 =-2\kappa\mathcal R_{N,\nu},
\label{eq:v2v14-subsidiary}
\end{equation}
where \(\mathcal R_N\) is the residual Ward defect.

\begin{theorem}[Constraint residual estimate]
\label{thm:v2v14-constraint}
Under (H1)--(H3), suppose the initial constraint data tend to zero in
\(H^s\times H^{s-1}\) and
\begin{equation}
 \|\mathcal R_N\|_{L^1_tH^{s-1}_x}\longrightarrow0.
\label{eq:v2v14-Ward-residual}
\end{equation}
Then
\[
 C_N\longrightarrow0
\quad\text{in }
C_tH^s_x\cap C_t^1H^{s-1}_x.
\]
\end{theorem}

\begin{proof}
Apply the linear hyperbolic energy estimate to
\eqref{eq:v2v14-subsidiary}.  Uniform bounds on \(g_N\) from (H2) control the
wave coefficients and curvature potential.  One obtains
\[
 \sup_{t\le T}E_s(C_N;t)
 \le C_{M,T}\left(
 E_s(C_N;0)
 +2\kappa\|\mathcal R_N\|_{L^1_tH^{s-1}_x}\right).
\]
Both terms tend to zero by hypothesis.
\end{proof}

\begin{corollary}[Summable shellwise Ward defects]
\label{cor:v2v14-Ward-shells}
If
\[
 \mathcal R_N=\sum_{n>N}r_n
\quad\text{and}\quad
\sum_n\|r_n\|_{L^1_tH^{s-1}_x}<\infty,
\]
then \eqref{eq:v2v14-Ward-residual} holds.
\end{corollary}

\begin{proof}
The norm of a tail of an absolutely convergent Banach-space series tends to
zero.
\end{proof}

\subsection{Interface with concentration defects}
\label{sec:v2v14-concentration}

The strong source topology in
\eqref{eq:v2v14-shell-summability} excludes the endpoint concentration
mechanism recorded in \eqref{eq:v2-defect-stress}: a nonzero singular measure
cannot appear as an unrecorded weak limit if the full sequence is Cauchy in
\(L^1_tH^{s-1}_x\).  The difficulty is proving this strong estimate from the
interacting state.  V10 and V11 provide summable Gaussian cubic regressions;
they do not yet provide the required higher-cumulant cluster bounds or local
Orlicz control.

\begin{assessment}[Concrete continuum gate after V14]
\label{ass:v2v14-gate}
A sufficient high-regularity route to the coupled limit is now:
\begin{enumerate}[label=\textup{(\roman*)}]
\item prove convergence of every local running coefficient after physical
renormalization conditions are imposed;
\item prove \eqref{eq:v2v14-L2-shell} for each invariant irrelevant source
shell, including boundary charts;
\item prove a common cutoff-independent hyperbolic existence interval and
the uniform bound (H2); and
\item prove summability of the repaired Ward residual.
\end{enumerate}
Under these hypotheses,
\cref{thm:v2v14-source-cauchy,thm:v2v14-hyperbolic-stability,
thm:v2v14-constraint} yield a constrained harmonic-gauge limit.  None of the
four hypotheses is silently asserted.
\end{assessment}

\subsection{V14 conclusion and V15 mandate}
\label{sec:v2v14-conclusion}

The source-to-Einstein implication has been placed in one explicit Banach
topology.  Absolute shell summability in \(L^1_tH^{s-1}_x\), together with
convergent local couplings, makes the source Cauchy.  Standard quasilinear
energy estimates then make uniformly controlled reduced solutions Cauchy,
and a vanishing Ward residual propagates the harmonic constraint.

The missing work has not been hidden: the programme still needs the strong
shell estimate, convergence of the local RG trajectory, a uniform existence
slab, and higher-cumulant control preventing Higgs and gauge concentration.
The next active version is V15 and must begin with the required companion
audit before choosing which of these four gates to attack.

\clearpage
\section*{Second Volume, Version V15: Companion Audit and a Uniform Harmonic Existence Slab}
\addcontentsline{toc}{section}{Second Volume, Version V15: Companion Audit and a Uniform Harmonic Existence Slab}

\subsection{Mandatory companion audit}
\label{sec:v2v15-audit}

\begin{assessment}[Read-only V15 audit]
\label{ass:v2v15-audit}
The current companion files inspected before this version were:
\begin{center}
\small
\begin{tabular}{llll}
\toprule
Programme & Version & Bytes & SHA--256\\
\midrule
Navier--Stokes & V31 & \(887537\) &
\texttt{F1121B62238AD5491BB7CF03635EA9934942EDF573ED8FAAA9EE79E1D38A6696}\\
Yang--Mills & compact V10 & \(136747\) &
\texttt{6E2F86AE8F607F6D584F64C936EF84F20ED014B73E0FEA06D29DDA1DB8341EA3}\\
\bottomrule
\end{tabular}
\end{center}
NS V31 proves a scale-matched weak/intermediate/critical Duhamel ladder and
shows by a forced mode that a weaker source rung cannot pay a stronger
response norm.  YM compact V10 closes a Cartan sector only after a complete
shared spectator product supplies the required first and second moments; a
normalized selection without those moments would not suffice.
\end{assessment}

\begin{firewall}[Audit transfer]
\label{fw:v2v15-audit}
No NS regularity estimate or YM Wilson moment is imported.  The source-to-
Einstein argument below uses the same two proof-order constraints: match the
source norm to the claimed solution topology, and take that norm only after
the complete invariant source has been assembled.
\end{firewall}

\subsection{Uniform time absolute continuity}
\label{sec:v2v15-absolute-continuity}

\begin{lemma}[\(L^1\)-Cauchy families are uniformly time-absolutely-continuous]
\label{lem:v2v15-uniform-ac}
Let \(Y\) be a Banach space and let \(f_N\) be Cauchy in
\(L^1([0,T_0];Y)\).  Then for every \(\varepsilon>0\) there is
\(\delta>0\) such that every measurable \(I\subset[0,T_0]\) with
\(|I|<\delta\) satisfies
\begin{equation}
 \sup_N\int_I\|f_N(t)\|_Y\,\mathrm dt<\varepsilon.
\label{eq:v2v15-uniform-ac}
\end{equation}
\end{lemma}

\begin{proof}
Choose \(N_0\) so that
\[
 \|f_N-f_{N_0}\|_{L^1}<\varepsilon/2
\quad\text{for }N\ge N_0.
\]
The finite collection \(f_1,\ldots,f_{N_0}\) consists of integrable scalar
norms.  Absolute continuity of their integrals gives \(\delta>0\) such that
\[
 \max_{1\le N\le N_0}\int_I\|f_N(t)\|_Y\,\mathrm dt<\varepsilon/2
\]
whenever \(|I|<\delta\).  For \(N>N_0\),
\[
 \int_I\|f_N\|_Y
 \le\|f_N-f_{N_0}\|_{L^1}
 +\int_I\|f_{N_0}\|_Y<\varepsilon.
\]
\end{proof}

\begin{counterexample}[A uniform \(L^1\) bound is insufficient]
\label{sep:v2v15-L1-bound}
On \([0,1]\), let \(f_N(t)=N\mathbf1_{[0,1/N]}(t)\).  Then
\(\|f_N\|_{L^1}=1\) for every \(N\), but for every \(\delta>0\),
\[
 \sup_N\int_0^\delta f_N(t)\,\mathrm dt=1.
\]
Thus a bounded source family need not have the uniform short-time smallness
needed for a common existence bootstrap.  The Cauchy hypothesis in
\cref{lem:v2v15-uniform-ac} supplies genuinely stronger information.
\end{counterexample}

\begin{proof}
Choose \(N>\delta^{-1}\); then the support of \(f_N\) lies in
\([0,\delta]\) and its integral is one.
\end{proof}

\subsection{An energy-admissible reduced system}
\label{sec:v2v15-energy-system}

Let \(s>d/2+2\).  Consider the harmonic-gauge reduced system
\[
 A^{\alpha\beta}(u_N)\partial_\alpha\partial_\beta u_N
 =\mathcal N(u_N,\partial u_N)+T_N
\]
from V14, with fixed energy-admissible boundary conditions when
\(\partial\Sigma\ne\varnothing\).  Write
\[
 E_s[u](t)=\|u(t)\|_{H^s}+\|\partial_tu(t)\|_{H^{s-1}}.
\]
We make the standard local energy hypothesis explicit.

\begin{definition}[Uniform energy-admissibility]
\label{def:v2v15-energy-admissible}
The reduced system is uniformly energy-admissible on a hyperbolic set
\(\mathcal U\) if every smooth solution staying in \(\mathcal U\) obeys
\begin{equation}
 E_s[u](t)
 \le C_0E_s[u](0)
 +C_0\int_0^t\|T(\sigma)\|_{H^{s-1}}\,\mathrm d\sigma
 +C_1\int_0^t(1+E_s[u](\sigma))E_s[u](\sigma)\,\mathrm d\sigma,
\label{eq:v2v15-energy}
\end{equation}
where \(C_0,C_1\) are uniform on bounded subsets a fixed positive distance
from the boundary of \(\mathcal U\).
\end{definition}

\begin{theorem}[Common cutoff-independent existence slab]
\label{thm:v2v15-common-slab}
Assume:
\begin{enumerate}[label=\textup{(U\arabic*)}]
\item the reduced system is uniformly energy-admissible on
\(\mathcal U\);
\item the initial data have
\[
 \sup_NE_s[u_N](0)\le M_0
\]
and their values lie a fixed positive Sobolev distance \(\delta_0\) inside
\(\mathcal U\);
\item \(T_N\) is Cauchy in
\(L^1([0,T_0];H^{s-1})\); and
\item the standard smooth local existence theorem holds for each cutoff
system while its energy and hyperbolicity margin remain finite.
\end{enumerate}
Then there are \(T_*>0\) and \(M_*>0\), independent of \(N\), such that every
cutoff solution exists on \([0,T_*]\), stays in \(\mathcal U\), and satisfies
\begin{equation}
 \sup_N\sup_{t\le T_*}E_s[u_N](t)\le M_*.
\label{eq:v2v15-uniform-energy}
\end{equation}
\end{theorem}

\begin{proof}
By \cref{lem:v2v15-uniform-ac}, choose \(T_*>0\) so small that
\[
 \sup_N\int_0^{T_*}\|T_N(t)\|_{H^{s-1}}\,\mathrm dt\le\eta,
\]
where \(\eta>0\) will be fixed below.  Bootstrap
\[
 E_s[u_N](t)\le 2C_0(M_0+1)
\]
and a Sobolev distance at least \(\delta_0/2\) from the boundary of
\(\mathcal U\).  Under this bootstrap,
\eqref{eq:v2v15-energy} gives
\[
 E_s[u_N](t)
 \le C_0M_0+C_0\eta
 +C_1T_*(1+2C_0(M_0+1))2C_0(M_0+1).
\]
Choose \(\eta\) and then \(T_*\) so the right side is strictly below
\(2C_0(M_0+1)\).  Since \(s>d/2+2\), Sobolev embedding controls the metric
coefficients continuously; reducing \(\eta,T_*\) further keeps the solution
within Sobolev distance \(\delta_0/2\) of its initial hyperbolic set.

The improved inequalities close the bootstrap uniformly.  The continuation
condition in (U4) prevents a maximal cutoff solution from ending before
\(T_*\).  Take \(M_*=2C_0(M_0+1)\).
\end{proof}

\begin{firewall}[What remains assumed]
\label{fw:v2v15-local-existence}
The theorem derives a common slab from one uniform energy inequality,
uniformly interior initial data, and a Cauchy source family.  It does not
construct the energy-admissible boundary realization or prove the
cutoff-level smooth local existence theorem.  Those analytic inputs must be
verified for the chosen harmonic, gauge, Higgs, and fermion hyperbolic
formulation.
\end{firewall}

\subsection{The direct lower-order stability rung}
\label{sec:v2v15-lower-rung}

\begin{theorem}[One-order-lower Lipschitz stability]
\label{thm:v2v15-lower-stability}
Under the hypotheses and uniform bound of
\cref{thm:v2v15-common-slab}, suppose the initial data are Cauchy in
\[
 H^{s-1}\times H^{s-2}.
\]
Then
\begin{equation}
 u_N\ \text{is Cauchy in}\ 
 C([0,T_*];H^{s-1})
 \cap C^1([0,T_*];H^{s-2}).
\label{eq:v2v15-lower-limit}
\end{equation}
\end{theorem}

\begin{proof}
Let \(w=u_N-u_M\) and subtract the equations as in V14.  Estimate the
difference equation at solution order \(s-1\).  The coefficient-difference
term contains
\[
 (A(u_M)-A(u_N))\partial^2u_M.
\]
The uniform order-\(s\) bound places \(\partial^2u_M\) in \(H^{s-2}\), while
Moser estimates bound the coefficient difference by \(w\) in \(H^{s-1}\).
Since \(s-1>d/2+1\), the product lies in \(H^{s-2}\), the correct forcing
space for the order-\(s-1\) energy.  The remaining nonlinear differences have
the same bound.  Hence
\begin{align}
 E_{s-1}[w](t)
 \le C\Bigl(
 E_{s-1}[w](0)
 +\|T_N-T_M\|_{L^1([0,t];H^{s-2})}
 +\int_0^tE_{s-1}[w](\sigma)\,\mathrm d\sigma
 \Bigr).
\label{eq:v2v15-lower-difference}
\end{align}
The source is Cauchy even in \(L^1H^{s-1}\), hence in the weaker space used
above.  Gronwall's inequality and the Cauchy initial data prove
\eqref{eq:v2v15-lower-limit}.
\end{proof}

\begin{corollary}[Strong lower limit and weak high limit]
\label{cor:v2v15-two-topologies}
There is a limit \(u_\infty\) such that
\[
 u_N\to u_\infty
\quad\text{strongly in the space in \eqref{eq:v2v15-lower-limit}},
\]
while, after passage to the same limit,
\[
 u_N\rightharpoonup^*u_\infty
\quad\text{in }L^\infty([0,T_*];H^s),
\]
and similarly for the time derivative in \(H^{s-1}\).
\end{corollary}

\begin{proof}
Strong convergence follows from completeness.  The uniform bound
\eqref{eq:v2v15-uniform-energy} gives weak-star compactness.  Any weak-star
subsequence has the same distributional limit as the strong
\(H^{s-1}\) convergence, so its limit is \(u_\infty\).
\end{proof}

\begin{assessment}[The scale-matched correction to V14]
\label{ass:v2v15-scale-matched}
V14 gave same-order \(H^s\) stability under a deliberately stronger uniform
order-\(s+1\) background hypothesis.  V15 proves the direct estimate available
from an order-\(s\) uniform energy bound: strong cutoff convergence at order
\(s-1\).  This distinction is the harmonic-Einstein analogue of the
scale-matched Duhamel ladder in the NS audit.  A same-order Lipschitz result
requires the extra derivative used in V14 or a more refined
continuous-dependence argument.
\end{assessment}

\subsection{Constraint propagation on the common slab}
\label{sec:v2v15-constraint}

\begin{corollary}[Lower-order subsidiary convergence]
\label{cor:v2v15-constraint}
On the common slab, suppose the harmonic constraint fields obey
\[
 \Box_{g_N}C_N+\operatorname{Ric}(g_N)C_N=-2\kappa\mathcal R_N,
\]
their initial data tend to zero in
\(H^{s-1}\times H^{s-2}\), and
\[
 \mathcal R_N\to0
\quad\text{in }L^1([0,T_*];H^{s-2}).
\]
Then
\[
 C_N\to0
\quad\text{in }C_tH^{s-1}_x\cap C_t^1H^{s-2}_x.
\]
\end{corollary}

\begin{proof}
The coefficients are uniformly controlled by
\eqref{eq:v2v15-uniform-energy}.  Apply the linear wave energy estimate at
order \(s-1\), then use the vanishing initial data and residual.
\end{proof}

\subsection{Interface with the shell source}
\label{sec:v2v15-shell-interface}

\begin{theorem}[V14 shell criterion supplies the common slab]
\label{thm:v2v15-shell-to-slab}
Suppose the local couplings converge and the invariant shell sources satisfy
\eqref{eq:v2v14-shell-summability}.  Then the sources
\eqref{eq:v2v14-source-decomposition} satisfy (U3).  Consequently, if
(U1), (U2), and (U4) hold, the reduced cutoff solutions share the interval
\([0,T_*]\) and converge strongly at the lower rung
\eqref{eq:v2v15-lower-limit}, provided their initial data are Cauchy there.
\end{theorem}

\begin{proof}
\cref{thm:v2v14-source-cauchy} gives the \(L^1H^{s-1}\) Cauchy property.
Apply
\cref{thm:v2v15-common-slab,thm:v2v15-lower-stability}.
\end{proof}

\begin{firewall}[The quantum source estimate remains the hard gate]
\label{fw:v2v15-quantum-source}
The common-slab theorem converts a strong source-Cauchy estimate into uniform
classical existence; it does not prove the source estimate.  The Gaussian
quadratic and cubic shell regressions in V1--V11 fall well short of
\eqref{eq:v2v14-shell-summability} for the complete interacting composite
stress at high Sobolev order.  Higher cumulants, boundary charts, and spatial
uniform integrability remain open.
\end{firewall}

\subsection{V15 conclusion and next step}
\label{sec:v2v15-conclusion}

The uniform-existence hypothesis in V14 has been reduced to checkable
cutoff-level inputs.  An \(L^1H^{s-1}\)-Cauchy source family is uniformly
time-absolutely-continuous, so a standard energy bootstrap yields one
cutoff-independent harmonic existence slab.  With only an order-\(s\)
uniform bound, the direct difference estimate is strong at order \(s-1\);
this avoids claiming a derivative-losing same-order estimate.

The next upstream task returns to the unresolved quantum source gate: derive a
spatially local cumulant or Orlicz estimate that implies the dyadic
\(L^1H^{s-1}\) shell bound after the complete Standard--Model source and its
boundary terms are assembled.  The next companion audit is due at V20.

\clearpage
\section*{Second Volume, Version V16: Local Orlicz Cumulants as Strong Einstein Sources}
\addcontentsline{toc}{section}{Second Volume, Version V16: Local Orlicz Cumulants as Strong Einstein Sources}

\subsection*{Purpose}

V14--V15 reduced the classical continuum step to absolute summability of the
assembled source shells in \(L^1_tH^{s-1}_x\).  The Gaussian trace estimates
of V10--V11 do not imply that spacetime norm for an interacting state.  This
version proves one sufficient probabilistic mechanism.  Local subgaussian
Orlicz bounds for the fields entering a connected cumulant give an
\(L^2\)-bound for the complete cumulant.  Output-frequency localization then
converts it into the strong Sobolev shell bound.  The hypothesis is local in
spacetime and is deliberately stronger than a tail bound for one globally
integrated energy.

\subsection{Local subgaussian norms}
\label{sec:v2v16-orlicz}

Let \(Z=[0,T]\times\Sigma\) with finite measure \(\mathrm dz\), and let
\((\Omega,\mathbb P)\) be a probability space.  On
\(Z\times\Omega\), define the Luxemburg norm
\begin{equation}
 \|X\|_{\psi_2}
 =\inf\left\{a>0:
 \int_{Z\times\Omega}
 \left(e^{|X|^2/a^2}-1\right)
 \,\mathrm dz\,\mathrm d\mathbb P\le1\right\}.
\label{eq:v2v16-psi2}
\end{equation}
Changing the finite normalization of \(\mathrm dz\) changes only harmless
constants.

\begin{lemma}[Subgaussian moment bound]
\label{lem:v2v16-moments}
For every \(p\ge2\),
\begin{equation}
 \|X\|_{L^p(Z\times\Omega)}
 \le C\sqrt p\,\|X\|_{\psi_2},
\label{eq:v2v16-moments}
\end{equation}
where \(C\) depends only on the finite normalization of \(Z\).
\end{lemma}

\begin{proof}
Normalize first so that the product measure has mass one.  If
\(\|X\|_{\psi_2}<a\), Markov's inequality gives
\[
 \mathbb P_{Z\times\Omega}(|X|>u)
 \le 2e^{-u^2/a^2}.
\]
Using
\[
 \mathbb E|X|^p
 =p\int_0^\infty u^{p-1}\mathbb P(|X|>u)\,\mathrm du
\]
and the substitution \(v=u^2/a^2\) yields
\[
 \mathbb E|X|^p
 \le p a^p\Gamma(p/2)
\]
up to an absolute factor.  The standard bound
\(\Gamma(p/2)^{1/p}\le C\sqrt p\) proves
\eqref{eq:v2v16-moments}.  Rescaling a finite measure changes \(C\).
\end{proof}

\begin{corollary}[Expectation contracts the local \(L^2\) norm]
\label{cor:v2v16-expectation}
If \(\tau(z)=\mathbb E_\Omega X(z,\omega)\), then
\begin{equation}
 \|\tau\|_{L^2(Z)}
 \le\|X\|_{L^2(Z\times\Omega)}
 \le C\|X\|_{\psi_2}.
\label{eq:v2v16-expectation}
\end{equation}
\end{corollary}

\begin{proof}
The first inequality is Jensen in \(\omega\), followed by integration in
\(z\).  The second is \cref{lem:v2v16-moments} with \(p=2\).
\end{proof}

\subsection{A local cumulant compiler}
\label{sec:v2v16-cumulant}

For random fields \(X_1,\ldots,X_m\) on the same probability space, define
their pointwise joint cumulant by
\begin{equation}
 \kappa_m(X_1,\ldots,X_m)(z)
 =
 \sum_{\pi\in\Pi_m}
 (|\pi|-1)!(-1)^{|\pi|-1}
 \prod_{B\in\pi}
 \mathbb E_\Omega\prod_{i\in B}X_i(z,\omega),
\label{eq:v2v16-cumulant}
\end{equation}
where \(\Pi_m\) is the set of partitions of \(\{1,\ldots,m\}\).

\begin{theorem}[Subgaussian connected-cumulant bound]
\label{thm:v2v16-cumulant}
For every fixed \(m\ge2\), there is a finite constant \(C_m\) such that
\begin{equation}
 \|\kappa_m(X_1,\ldots,X_m)\|_{L^2(Z)}
 \le C_m\prod_{i=1}^m\|X_i\|_{\psi_2}.
\label{eq:v2v16-cumulant-bound}
\end{equation}
\end{theorem}

\begin{proof}
Fix a partition \(\pi\).  For each block \(B\in\pi\), introduce an
independent copy \(\omega_B\) of the probability coordinate.  Then the
partition product in \eqref{eq:v2v16-cumulant} equals
\begin{equation}
 \mathbb E_{\{\omega_B\}_{B\in\pi}}
 \prod_{B\in\pi}\prod_{i\in B}X_i(z,\omega_B).
\label{eq:v2v16-independent-blocks}
\end{equation}
Minkowski and Jensen bound its \(L^2(Z)\) norm by the \(L^2\) norm on the
enlarged product space of the product of all \(m\) fields.  Hölder with
exponent \(m\) gives
\[
 \left\|\prod_{B\in\pi}\prod_{i\in B}X_i(z,\omega_B)\right\|_{L^2}
 \le\prod_{i=1}^m
 \|X_i(z,\omega_{B(i)})\|_{L^{2m}}.
\]
Each marginal has the same \(Z\times\Omega\) norm as \(X_i\).  By
\cref{lem:v2v16-moments}, the last expression is at most
\[
 (C\sqrt{2m})^m\prod_i\|X_i\|_{\psi_2}.
\]
Finally sum the finitely many partition terms with their Möbius
coefficients.  Their total absolute coefficient depends only on \(m\), giving
\eqref{eq:v2v16-cumulant-bound}.
\end{proof}

\begin{remark}[Why independent copies appear only by partition blocks]
Equation \eqref{eq:v2v16-independent-blocks} uses one independent probability
coordinate for each disconnected moment block.  It does not assign an
independent copy to every factor inside a connected block.  This preserves
the common-resampling constraint from V10 and the YM audits.
\end{remark}

\begin{counterexample}[No order-uniform cumulant constant]
\label{sep:v2v16-order}
The constant \(C_m\) cannot be treated as uniform in \(m\) by the proof above:
the partition Möbius coefficients and the \(L^{2m}\) moments grow with \(m\).
Thus fixed-order shell estimates do not by themselves sum the full
interaction expansion.
\end{counterexample}

\begin{proof}
The one-block partition already uses an \(L^{2m}\) norm, whose subgaussian
bound grows like \((C\sqrt m)^m\).  The number of partitions also grows with
\(m\).  Hence the derived constant is explicitly order dependent.
\end{proof}

\subsection{From local cumulants to Sobolev source shells}
\label{sec:v2v16-source}

Let \(\Delta_n\) be an output spectral projector supported where
\[
 c_-\Lambda_n\le\langle\xi\rangle\le c_+\Lambda_n.
\]
Let the assembled deterministic source shell be
\begin{equation}
 \tau_n(z)
 =\Delta_n\kappa_m(X_{1,n},\ldots,X_{m,n})(z).
\label{eq:v2v16-source-shell}
\end{equation}

\begin{theorem}[Local Orlicz criterion for the V14 source norm]
\label{thm:v2v16-source-criterion}
Suppose for some exponents \(a_1,\ldots,a_m\),
\begin{equation}
 \|X_{i,n}\|_{\psi_2}\le K_i\Lambda_n^{-a_i}
\label{eq:v2v16-field-decay}
\end{equation}
and
\begin{equation}
 \sum_{i=1}^ma_i>s-1.
\label{eq:v2v16-source-threshold}
\end{equation}
Then
\begin{equation}
 \|\tau_n\|_{L^1_tH^{s-1}_x}
 \le C_{m,T,\Sigma}
 \left(\prod_iK_i\right)
 \Lambda_n^{s-1-\sum_i a_i},
\label{eq:v2v16-source-Hs}
\end{equation}
and the dyadic source shells are absolutely summable in the V14 space
\(\mathcal X_s\).
\end{theorem}

\begin{proof}
The output support gives
\[
 \|\tau_n\|_{L^2_tH^{s-1}_x}
 \le C\Lambda_n^{s-1}
 \|\kappa_m(X_{1,n},\ldots,X_{m,n})\|_{L^2(Z)}.
\]
Use \cref{thm:v2v16-cumulant} and
\eqref{eq:v2v16-field-decay}, then apply Cauchy--Schwarz in time:
\[
 \|\tau_n\|_{L^1_tH^{s-1}_x}
 \le\sqrt T\|\tau_n\|_{L^2_tH^{s-1}_x}.
\]
This is \eqref{eq:v2v16-source-Hs}.  Under
\eqref{eq:v2v16-source-threshold}, its dyadic exponent is negative.
\end{proof}

\begin{corollary}[Explicit excess exponent]
\label{cor:v2v16-excess}
If
\[
 \sum_i a_i=s-1+\varepsilon
\]
with \(\varepsilon>0\), then
\[
 \|\tau_n\|_{\mathcal X_s}
 \le C\Lambda_n^{-\varepsilon},
\]
which is exactly the summable conclusion of
\cref{prop:v2v14-dyadic}.
\end{corollary}

\begin{proof}
Substitute into \eqref{eq:v2v16-source-Hs}.
\end{proof}

\subsection{Boundary source version}
\label{sec:v2v16-boundary}

Let \(Z_\partial=[0,T]\times\partial\Sigma\) and define the local
\(\psi_2\) norm using its product measure.  Let a boundary source shell
\(\tau_n^\partial\) have tangential frequency \(\Lambda_n\).

\begin{corollary}[Boundary trace threshold]
\label{cor:v2v16-boundary}
If
\[
 \tau_n^\partial
 =\Delta_n^\partial\kappa_m(
 X_{1,n}^\partial,\ldots,X_{m,n}^\partial)
\]
and
\[
 \|X_{i,n}^\partial\|_{\psi_2(Z_\partial\times\Omega)}
 \le K_i^\partial\Lambda_n^{-b_i},
\]
then
\begin{equation}
 \|\tau_n^\partial\|_{L^1_tH^{s-3/2}(\partial\Sigma)}
 \le C
 \left(\prod_iK_i^\partial\right)
 \Lambda_n^{s-3/2-\sum_i b_i}.
\label{eq:v2v16-boundary-source}
\end{equation}
It is summable if \(\sum_i b_i>s-3/2\).
\end{corollary}

\begin{proof}
Repeat the proof of \cref{thm:v2v16-source-criterion} on
\(\partial\Sigma\), using the boundary Sobolev exponent \(s-3/2\) appropriate
to the trace of an order-\(s-1\) bulk energy estimate.
\end{proof}

\begin{firewall}[Boundary dynamics may require a different norm]
\label{fw:v2v16-boundary}
The exponent in \cref{cor:v2v16-boundary} is the standard trace-space
regression.  A particular maximally dissipative or constraint-preserving
Einstein boundary system may require additional tangential derivatives or
compatibility conditions.  Those must be read from its actual energy
estimate; the trace theorem alone does not select the boundary domain.
\end{firewall}

\subsection{Positivity and locality firewalls}
\label{sec:v2v16-firewalls}

\begin{firewall}[The Orlicz hypothesis is spatially local]
\label{fw:v2v16-local}
The norm \eqref{eq:v2v16-psi2} controls the random density on
\([0,T]\times\Sigma\times\Omega\).  A tail estimate for the single random
number
\(\int_{[0,T]\times\Sigma}|X|^2\) does not automatically imply it.  The
global-versus-spatial concentration separator in the first volume therefore
remains in force.
\end{firewall}

\begin{firewall}[Positive bosonic measure only]
\label{fw:v2v16-positive}
\cref{thm:v2v16-cumulant} uses a positive probability measure and Jensen.
After fermions and ghosts are integrated algebraically, it may be applied to
the remaining bosonic Euclidean measure only if that measure is positive.
A complex chiral phase or an untreated gravitational conformal contour
invalidates the probabilistic step.  The Berezin trace estimates of V11
remain the appropriate finite-cutoff tool for those sectors.
\end{firewall}

\begin{assessment}[What must be estimated in the Standard Model]
\label{ass:v2v16-sm-input}
The sufficient input is now quantitative.  For every connected source
cumulant of order \(m\), after species, contact, Ward, and local-counterterm
cancellations have been assembled, assign local subgaussian decay exponents
\(a_i\).  The sum must exceed the Sobolev cost \(s-1\), or the boundary cost
\(s-3/2\).  The fixed-order constants \(C_m\) must then be combined with a
separate interaction-order or cluster expansion.  Neither requirement is
proved by the Gaussian shell power count alone.
\end{assessment}

\subsection{V16 conclusion and next step}
\label{sec:v2v16-conclusion}

V16 supplies a direct sufficient bridge from local probabilistic control to
the strong source topology used by the harmonic Einstein argument.  A
connected cumulant of locally subgaussian fields obeys an \(L^2\) spacetime
bound with one independent copy per moment-partition block.  Frequency
localization then pays the \(H^{s-1}\) derivative cost, and any positive excess
in the sum of decay exponents makes the dyadic source series convergent.

The remaining bottleneck is no longer a choice of topology.  It is to prove
the local subgaussian exponents and control the growth in cumulant order for
the complete interacting gauge--Higgs measure, while retaining the
Slavnov-compatible local subtraction and the boundary source.  The next
version should test this criterion on the quartic Higgs Gibbs factor, where
the local potential is coercive but gradient concentration remains possible.

\clearpage
\section*{Second Volume, Version V17: Quartic Higgs Tails Do Not Supply Ultraviolet Source Decay}
\addcontentsline{toc}{section}{Second Volume, Version V17: Quartic Higgs Tails Do Not Supply Ultraviolet Source Decay}

\subsection*{Purpose}

V16 used local subgaussian norms to control fixed connected cumulants in
spacetime.  A quartic Higgs Gibbs factor has stronger-than-Gaussian tails for
the field itself, but its quartic energy density is only subexponential.  This
version extends the cumulant compiler to general sub-Weibull exponents and
computes the exact cutoff rescaling of a one-site quartic measure.  The result
is a firewall: coercive tails control moments of the dimensionless field, but
do not supply the negative cutoff power required by the strong source
criterion.

\subsection{Sub-Weibull moment and cumulant bounds}
\label{sec:v2v17-subweibull}

Let \(Z\) carry a normalized finite measure, so that
\(\int_Z\mathrm dz=1\).  For \(\alpha>0\), define
\begin{equation}
 \|X\|_{\psi_\alpha}
 =\inf\left\{a>0:
 \int_{Z\times\Omega}
 \left(e^{|X|^\alpha/a^\alpha}-1\right)
 \,\mathrm dz\,\mathrm d\mathbb P\le1\right\}.
\label{eq:v2v17-psialpha}
\end{equation}

\begin{lemma}[Sub-Weibull moments]
\label{lem:v2v17-moments}
For every \(p\ge1\),
\begin{equation}
 \|X\|_{L^p(Z\times\Omega)}
 \le C_\alpha p^{1/\alpha}\|X\|_{\psi_\alpha}.
\label{eq:v2v17-moments}
\end{equation}
\end{lemma}

\begin{proof}
After finite measure normalization, the Luxemburg bound gives
\[
 \mathbb P(|X|>u)\le2\exp(-(u/a)^\alpha)
\]
for every \(a>\|X\|_{\psi_\alpha}\).  Insert this tail in
\[
 \mathbb E|X|^p
 =p\int_0^\infty u^{p-1}\mathbb P(|X|>u)\,\mathrm du
\]
and use \(v=(u/a)^\alpha\).  The result is bounded by
\[
 Cpa^p\Gamma(p/\alpha).
\]
Taking the \(p\)-th root and using
\(\Gamma(p/\alpha)^{1/p}\le C_\alpha p^{1/\alpha}\) proves the claim.
\end{proof}

\begin{theorem}[Fixed-order sub-Weibull cumulant compiler]
\label{thm:v2v17-cumulant}
Let \(X_1,\ldots,X_m\) satisfy
\(\|X_i\|_{\psi_{\alpha_i}}<\infty\).  Then
\begin{equation}
 \|\kappa_m(X_1,\ldots,X_m)\|_{L^2(Z)}
 \le
 C_{m,\alpha_1,\ldots,\alpha_m}
 \prod_{i=1}^m\|X_i\|_{\psi_{\alpha_i}}.
\label{eq:v2v17-cumulant}
\end{equation}
\end{theorem}

\begin{proof}
Use the independent-copy representation
\eqref{eq:v2v16-independent-blocks} for each moment partition.  Minkowski,
Jensen, and Hölder with exponent \(m\) reduce its \(L^2(Z)\) norm to
\[
 \prod_{i=1}^m\|X_i\|_{L^{2m}(Z\times\Omega)}.
\]
Apply \cref{lem:v2v17-moments} with \(p=2m\), then sum the finite partition
Möbius coefficients.  All constants depend on the fixed cumulant order and
the displayed tail exponents.
\end{proof}

\begin{corollary}[Powers change the Orlicz exponent]
\label{cor:v2v17-powers}
If \(X\in\psi_\alpha\) and \(q>0\), then
\begin{equation}
 \||X|^q\|_{\psi_{\alpha/q}}
 =\|X\|_{\psi_\alpha}^q
\label{eq:v2v17-power-orlicz}
\end{equation}
under the common Luxemburg normalization.
\end{corollary}

\begin{proof}
The exponential in the definition is unchanged:
\[
 \exp\left(
 \frac{(|X|^q)^{\alpha/q}}{(a^q)^{\alpha/q}}
 \right)
 =\exp(|X|^\alpha/a^\alpha).
\]
Take the two infima.
\end{proof}

\begin{assessment}[Higgs source tail classes]
\label{ass:v2v17-tail-classes}
If \(H\in\psi_4\), then \(H^\dagger H\in\psi_2\) and
\((H^\dagger H)^2\in\psi_1\).  Likewise, if a Gaussian-type gradient is
\(\psi_2\), its squared kinetic density is \(\psi_1\).  Thus the complete
Higgs stress should be tested with
\cref{thm:v2v17-cumulant}, not by assuming that every composite source
density remains subgaussian.
\end{assessment}

\subsection{One-site quartic regression}
\label{sec:v2v17-one-site}

For \(\lambda>0\), \(a\ge0\), and cutoff scale \(\Lambda>0\), consider the
one-dimensional probability measure
\begin{equation}
 \mathrm d\mu_{\Lambda}(h)
 =Z_\Lambda^{-1}
 \exp\left(
 -\frac a{2\Lambda^2}h^2
 -\frac{\lambda}{\Lambda^4}h^4
 \right)\mathrm dh.
\label{eq:v2v17-one-site}
\end{equation}
The powers of \(\Lambda\) are those obtained when a four-dimensional cell of
volume \(\Lambda^{-4}\) is assigned a physical field of engineering dimension
one.  This is a single-cell scaling regression: the coefficient
\(a/\Lambda^2\) represents a positive dimensionless quadratic stiffness after
cell rescaling.  The measure is not asserted to equal the full lattice Higgs
measure, whose kinetic term couples neighbouring cells.

\begin{theorem}[Exact quartic cutoff rescaling]
\label{thm:v2v17-rescaling}
Let
\[
 y=\lambda^{1/4}\Lambda^{-1}h,
 \qquad
 \alpha=a\lambda^{-1/2}.
\]
Then the law of \(y\) has density proportional to
\begin{equation}
 \exp\left(-\frac\alpha2y^2-y^4\right)\mathrm dy,
\label{eq:v2v17-rescaled-law}
\end{equation}
independent of \(\Lambda\).  Consequently,
\begin{align}
 \|h\|_{\psi_4(\mu_\Lambda)}
 &=C_{\alpha}\lambda^{-1/4}\Lambda,
\label{eq:v2v17-h-scale}\\
 \|h^2\|_{\psi_2(\mu_\Lambda)}
 &=C_{\alpha}^2\lambda^{-1/2}\Lambda^2,
\label{eq:v2v17-h2-scale}\\
 \left\|\frac{\lambda}{\Lambda^4}h^4\right\|_{\psi_1(\mu_\Lambda)}
 &=C_{\alpha}^4,
\label{eq:v2v17-cell-energy-scale}
\end{align}
where \(0<C_\alpha<\infty\) depends on the dimensionless mass parameter but
not on \(\Lambda\).
\end{theorem}

\begin{proof}
Substitute \(h=\lambda^{-1/4}\Lambda y\) in
\eqref{eq:v2v17-one-site}.  The exponent becomes
\[
 -\frac{a\lambda^{-1/2}}2y^2-y^4,
\]
and the Jacobian cancels from the normalized law.  The density
\eqref{eq:v2v17-rescaled-law} has a finite \(\psi_4\) norm
\(C_\alpha\), since
\[
 \int e^{c|y|^4}e^{-\alpha y^2/2-y^4}\,\mathrm dy<\infty
\]
for every \(0<c<1\).  Homogeneity of the Luxemburg norm gives
\eqref{eq:v2v17-h-scale}.  Apply
\cref{cor:v2v17-powers} for
\eqref{eq:v2v17-h2-scale}--\eqref{eq:v2v17-cell-energy-scale}.
\end{proof}

\begin{corollary}[Tail control is not field decay]
\label{cor:v2v17-no-decay}
The dimensionless field \(\Lambda^{-1}h\) has a cutoff-uniform
\(\psi_4\) norm, but the physical field norm grows like \(\Lambda\).  The
cell-integrated quartic energy has an order-one \(\psi_1\) norm; it does not
decay with the cutoff.
\end{corollary}

\begin{proof}
This is \eqref{eq:v2v17-h-scale} and
\eqref{eq:v2v17-cell-energy-scale}.
\end{proof}

\subsection{A Gaussian shell comparison}
\label{sec:v2v17-gaussian-shell}

\begin{proposition}[Four-dimensional free-field point variance]
\label{prop:v2v17-free-variance}
For a translation-invariant scalar Gaussian shell in four Euclidean
dimensions with covariance multiplier
\[
 \frac{\chi(p/\Lambda)}{p^2+m^2},
\]
where \(\chi\) is nonnegative and supported in a fixed annulus, one has
\begin{equation}
 \mathbb E|H_\Lambda(x)|^2
 =\int_{\mathbb R^4}
 \frac{\chi(p/\Lambda)}{p^2+m^2}\,
 \frac{\mathrm d^4p}{(2\pi)^4}
 \asymp\Lambda^2
\label{eq:v2v17-free-variance}
\end{equation}
for \(\Lambda\gg m\), provided \(\chi\) is not zero.
\end{proposition}

\begin{proof}
Set \(p=\Lambda q\).  The integral becomes
\[
 \Lambda^2\int
 \frac{\chi(q)}{q^2+m^2/\Lambda^2}
 \frac{\mathrm d^4q}{(2\pi)^4}.
\]
The annular support stays away from \(q=0\), so the remaining integral is
bounded above and below by positive constants for \(\Lambda\gg m\).
\end{proof}

\begin{assessment}[Agreement of the two regressions]
\label{ass:v2v17-agreement}
The Gaussian shell has pointwise standard deviation \(O(\Lambda)\), and the
single-cell quartic regression has the same engineering power in
\eqref{eq:v2v17-h-scale}.  This comparison identifies a common dimensional
obstruction; it is not an identification of the two probability laws.
Quartic coercivity gives a quartic far tail for the dimensionless one-site
variable, but it does not reverse the engineering growth of the physical
high-frequency field.  Therefore the positive exponents required by
\eqref{eq:v2v16-field-decay} cannot be assigned to raw point fields.
\end{assessment}

\subsection{Renormalized composites are the only possible input}
\label{sec:v2v17-renormalized}

Let \(\mathcal T_n^{\mathrm{raw}}\) be a local Higgs stress shell.  Its
individual monomials contain powers with the growth demonstrated above.  The
candidate input to V16 is instead
\begin{equation}
 \mathcal T_n^{\mathrm{ren}}
 =(I-P_{\mathrm{loc}})
 \left[
 \mathcal T_n^{\mathrm{raw}}
 +\mathcal T_n^{\mathrm{contact}}
 +\mathcal T_n^{\mathrm{measure}}
 \right],
\label{eq:v2v17-renormalized-source}
\end{equation}
formed after the contact, measure, and Ward-required terms have been
assembled.

\begin{firewall}[Coercivity does not estimate the projected composite]
\label{fw:v2v17-projected}
The quartic Gibbs factor proves integrability of raw field powers.  It does
not prove a negative cutoff power for
\eqref{eq:v2v17-renormalized-source}.  Such a power must come from connected
subtraction, external-momentum derivatives, spatial averaging, or a cluster
estimate.  Applying the tail bound before those cancellations merely bounds a
quantity of order one per cell.
\end{firewall}

\begin{theorem}[Sub-Weibull version of the V16 source criterion]
\label{thm:v2v17-source-criterion}
Let
\[
 \tau_n=\Delta_n\kappa_m(X_{1,n},\ldots,X_{m,n})
\]
and suppose
\[
 \|X_{i,n}\|_{\psi_{\alpha_i}}\le K_i\Lambda_n^{-a_i}.
\]
If \(\sum_i a_i>s-1\), then
\[
 \sum_n\|\tau_n\|_{L^1_tH^{s-1}_x}<\infty.
\]
The conclusion is independent of the positive exponents
\(\alpha_i\) at each fixed cumulant order, although its constant depends on
them and on \(m\).
\end{theorem}

\begin{proof}
Replace \cref{thm:v2v16-cumulant} by
\cref{thm:v2v17-cumulant} in the proof of
\cref{thm:v2v16-source-criterion}.  The cutoff exponent is unchanged.
\end{proof}

\begin{assessment}[Precise Higgs gate after V17]
\label{ass:v2v17-higgs-gate}
For a fixed connected order, subexponential composite tails are adequate for
the \(L^2\) cumulant compiler.  The missing statement is a decay estimate for
the fully renormalized local tensor in
\eqref{eq:v2v17-renormalized-source}.  It must exceed the Sobolev derivative
cost \(s-1\) after all factors are combined.  Raw quartic coercivity supplies
moment finiteness but zero ultraviolet excess at the cell-energy level.
\end{assessment}

\subsection{V17 conclusion and next step}
\label{sec:v2v17-conclusion}

The V16 probabilistic bridge extends from subgaussian to fixed-order
sub-Weibull inputs.  A quartic Higgs field is naturally \(\psi_4\), while its
quartic and gradient energy densities are naturally \(\psi_1\).  These tail
classes suffice for fixed cumulant moments.  Exact cutoff rescaling shows,
however, that the physical field grows like \(\Lambda\) and the
cell-integrated quartic energy remains order one.  Tail coercivity is
therefore not the negative shell power required by the Einstein source norm.

The next attack must act on the renormalized composite itself.  A concrete
route is to prove a block-averaging or connected-cluster estimate for
\eqref{eq:v2v17-renormalized-source} after its local jet is removed, and test
whether the resulting spatial gain survives the gradient and improvement
parts of the Higgs stress.

\clearpage
\section*{Second Volume, Version V18: Block Averaging, Jet Gains, and the Strong-Source Obstruction}
\addcontentsline{toc}{section}{Second Volume, Version V18: Block Averaging, Jet Gains, and the Strong-Source Obstruction}

\subsection*{Purpose}

V17 isolated the missing Higgs statement as decay of the fully assembled,
renormalized stress shell.  Two frequently invoked mechanisms are connected
clustering and removal of a local external-momentum jet.  This version proves
the gains supplied by both mechanisms and tracks the scale on which they
hold.  It also gives a smooth finite-range counterexample showing that
clustering alone cannot yield convergence in the full positive Sobolev source
space used in V14--V15.

Throughout this version, \(\ell\) denotes the microscopic correlation length;
for the \(n\)-th ultraviolet shell one has \(\ell_n\simeq\Lambda_n^{-1}\).

\subsection{A quantitative block-average theorem}
\label{sec:v2v18-block}

Let \(\mathbb T^d\) be a flat torus of fixed volume and let
\(T_\ell:\mathbb T^d\times\Omega\to\mathbb R^N\) be centred.  Suppose that
there are \(A_\ell\ge0\), \(c_0>0\), and a nonnegative radial function
\(\Theta\in L^1(\mathbb R^d)\) such that
\begin{equation}
 \left|
 \mathbb E\big[
 \langle T_\ell(x),u\rangle
 \langle T_\ell(y),v\rangle
 \big]\right|
 \le
 A_\ell^2 |u||v|\,
 \Theta\!\left(\frac{\operatorname{dist}(x,y)}{\ell}\right)
\label{eq:v2v18-cluster}
\end{equation}
for \(0<\ell<c_0\).  Choose
\(\rho\in C_c^\infty(\mathbb R^d)\), supported inside a coordinate ball, with
\(\int\rho=1\), put
\[
 \rho_L(x)=L^{-d}\rho(x/L),
 \qquad
 B_LT_\ell=\rho_L*T_\ell,
\]
and restrict to \(\ell\le L<c_0\).

\begin{theorem}[Connected block-average gain]
\label{thm:v2v18-block}
For every multi-index \(\beta\),
\begin{equation}
 \left(
 \mathbb E
 \|\partial^\beta B_LT_\ell\|_{L^2_x}^2
 \right)^{1/2}
 \le
 C_{\rho,\Theta,\beta,\mathbb T^d}
 A_\ell\,
 \ell^{d/2}L^{-d/2-|\beta|}.
\label{eq:v2v18-block-bound}
\end{equation}
Consequently, for every nonnegative integer \(r\),
\begin{equation}
 \left(
 \mathbb E
 \|B_LT_\ell\|_{H^r_x}^2
 \right)^{1/2}
 \le
 C_r A_\ell\ell^{d/2}
 \sum_{j=0}^rL^{-d/2-j}.
\label{eq:v2v18-block-hr}
\end{equation}
\end{theorem}

\begin{proof}
Write \(f=\partial^\beta\rho_L\).  For each \(x\), the covariance bound and
the triangle inequality give
\[
 \mathbb E|\partial^\beta B_LT_\ell(x)|^2
 \le
 C_NA_\ell^2
 \iint |f(x-y)||f(x-z)|
 \Theta\!\left(\frac{\operatorname{dist}(y,z)}{\ell}\right)
 \,\mathrm dy\,\mathrm dz.
\]
The periodic kernel
\(\Theta_\ell(w)=\Theta(\operatorname{dist}(0,w)/\ell)\) satisfies
\(\|\Theta_\ell\|_{L^1(\mathbb T^d)}
 \le C_{\Theta,\mathbb T^d}\ell^d\).
After integration in \(x\), Fubini and Young's convolution inequality imply
\[
 \mathbb E\|\partial^\beta B_LT_\ell\|_2^2
 \le
 C_NA_\ell^2
 \|f\|_2\|\Theta_\ell*|f|\|_2
 \le
 C_NA_\ell^2\ell^d\|f\|_2^2.
\]
The scaling
\(\|\partial^\beta\rho_L\|_2
 =L^{-d/2-|\beta|}\|\partial^\beta\rho\|_2\)
proves \eqref{eq:v2v18-block-bound}.  Summing over
\(|\beta|\le r\) proves \eqref{eq:v2v18-block-hr}.
\end{proof}

\begin{corollary}[Four-dimensional coarse-source summability]
\label{cor:v2v18-fixed-block}
Let \(d=4\), \(\ell_n\simeq2^{-n}\), and fix \(L>0\).  If
\(A_{\ell_n}\le C\ell_n^{-\gamma}\) with \(\gamma<2\), then for every fixed
integer \(r\),
\begin{equation}
 \sum_n
 \left(
 \mathbb E\|B_LT_{\ell_n}\|_{H^r_x}^2
 \right)^{1/2}
 <\infty.
\label{eq:v2v18-fixed-block-sum}
\end{equation}
The same conclusion holds in \(L^1_tH^r_x\) on a finite time interval when
\eqref{eq:v2v18-cluster} holds uniformly in time.
\end{corollary}

\begin{proof}
For fixed \(L\), the right side of \eqref{eq:v2v18-block-hr} is bounded by
\(C_{r,L}\ell_n^{2-\gamma}\), a summable geometric sequence.  Uniformity in
time and Cauchy--Schwarz in probability give the spacetime statement.
\end{proof}

\begin{assessment}[What the central-limit power measures]
\label{ass:v2v18-clt}
The factor \(\ell^{d/2}L^{-d/2}\) is the square root of the inverse number of
correlation cells in an \(L\)-block.  Derivatives act on the averaging kernel
and cost \(L^{-1}\), not \(\ell^{-1}\).  The estimate is therefore useful
only when the observable is retained at a scale \(L\) larger than its
microscopic correlation length.
\end{assessment}

\subsection{Local-jet removal has the same scale restriction}
\label{sec:v2v18-jet}

The deterministic analogue is most transparent for a translation-invariant
connected kernel.  Let \(K_\ell\in L^1(\mathbb R^d)\) satisfy
\begin{equation}
 \int_{\mathbb R^d}z^\beta K_\ell(z)\,\mathrm dz=0
 \quad (|\beta|<q),
 \qquad
 \int_{\mathbb R^d}|z|^q|K_\ell(z)|\,\mathrm dz
 \le M_\ell\ell^q.
\label{eq:v2v18-moments}
\end{equation}

\begin{theorem}[External-momentum Taylor gain]
\label{thm:v2v18-jet}
Under \eqref{eq:v2v18-moments},
\begin{equation}
 |\widehat K_\ell(p)|
 \le C_{d,q}M_\ell(\ell|p|)^q.
\label{eq:v2v18-jet-gain}
\end{equation}
In particular, on external frequencies \(|p|\lesssim L^{-1}\), the
subtracted kernel gains \((\ell/L)^q\); on the diagonal shell
\(|p|\simeq\ell^{-1}\), this argument gives no small factor.
\end{theorem}

\begin{proof}
Taylor's formula with integral remainder gives
\[
 e^{-ip\cdot z}
 =
 \sum_{j=0}^{q-1}\frac{(-ip\cdot z)^j}{j!}
 +R_q(p,z),
 \qquad
 |R_q(p,z)|\le C_q|p|^q|z|^q.
\]
Every polynomial term integrates to zero by
\eqref{eq:v2v18-moments}.  Integrating the remainder and applying the
weighted \(L^1\) bound proves \eqref{eq:v2v18-jet-gain}.
\end{proof}

\begin{corollary}[Cluster localization supplies the weighted moment]
\label{cor:v2v18-cluster-moment}
If
\[
 |K_\ell(z)|
 \le A_\ell\ell^{-d}\Theta(z/\ell),
 \qquad
 \int_{\mathbb R^d}|w|^q\Theta(w)\,\mathrm dw<\infty,
\]
and the moments below order \(q\) vanish, then
\[
 |\widehat K_\ell(p)|
 \le C_{\Theta,q}A_\ell(\ell|p|)^q.
\]
\end{corollary}

\begin{proof}
Changing variables \(z=\ell w\) gives
\[
 \int |z|^q|K_\ell(z)|\,\mathrm dz
 \le
 A_\ell\ell^q\int|w|^q\Theta(w)\,\mathrm dw.
\]
Apply \cref{thm:v2v18-jet}.
\end{proof}

\begin{firewall}[A local jet is a low-external-frequency cancellation]
\label{fw:v2v18-jet}
The factor in \eqref{eq:v2v18-jet-gain} is
\((\ell|p|)^q\), not \(\ell^q\) by itself.  Removing more local Taylor
coefficients can make a fixed macroscopic observable converge rapidly, but
it does not suppress the part of the renormalized stress whose output
frequency is comparable to the ultraviolet shell.  No power-counting table
may record the Taylor order as unconditional ultraviolet decay.
\end{firewall}

\subsection{A smooth finite-range obstruction}
\label{sec:v2v18-counterexample}

The loss at \(L\simeq\ell\) is real, rather than an artefact of the proof.

\begin{proposition}[Clustering does not imply strong source decay]
\label{prop:v2v18-counterexample}
For every integer \(r\ge0\) and every sequence
\(\ell=N^{-1}\downarrow0\), there are centred smooth random fields
\(T_\ell\) on \(\mathbb T^d\) satisfying
\eqref{eq:v2v18-cluster} with a compactly supported \(\Theta\) and
\(A_\ell=A\), while
\begin{align}
 \left(\mathbb E\|T_\ell\|_{L^2_x}^2\right)^{1/2}
 &=c_0A,
\label{eq:v2v18-l2-nondecay}\\
 \left(\mathbb E\|T_\ell\|_{H^r_x}^2\right)^{1/2}
 &\ge c_rA\ell^{-r}.
\label{eq:v2v18-hr-growth}
\end{align}
Moreover, for each fixed \(L>0\),
\begin{equation}
 \liminf_{\ell\downarrow0}
 \left(
 \mathbb E\|(I-B_L)T_\ell\|_{L^2_x}^2
 \right)^{1/2}
 \ge c_0A.
\label{eq:v2v18-residual}
\end{equation}
\end{proposition}

\begin{proof}
Partition the unit torus into \(N^d\) cubes \(Q\) of side \(\ell\).  Fix a
nonzero \(\psi\in C_c^\infty((0,1)^d)\), choose independent symmetric signs
\(\varepsilon_Q\), and set
\[
 T_\ell(x)
 =
 A\sum_Q\varepsilon_Q
 \psi\!\left(\frac{x-x_Q}{\ell}\right),
\]
where \(x_Q\) is the lower corner of \(Q\).  The supports are disjoint.
Independence and centring make the covariance zero unless both points lie in
the same cube support.  Hence \eqref{eq:v2v18-cluster} holds with a bounded
\(\Theta\) supported in \([0,\sqrt d]\), after changing its constant.

Orthogonality of the signs and scaling give
\[
 \mathbb E\|\partial^\beta T_\ell\|_2^2
 =
 A^2N^d\ell^{d-2|\beta|}
 \|\partial^\beta\psi\|_2^2
 =
 A^2\ell^{-2|\beta|}
 \|\partial^\beta\psi\|_2^2.
\]
The case \(\beta=0\) proves \eqref{eq:v2v18-l2-nondecay}; choosing a
nonzero derivative of order \(r\) proves \eqref{eq:v2v18-hr-growth}.
For fixed \(L\), \cref{thm:v2v18-block} gives
\(\|B_LT_\ell\|_{L^2(\Omega\times\mathbb T^d)}=O(\ell^{d/2})\).
The reverse triangle inequality with
\eqref{eq:v2v18-l2-nondecay} proves
\eqref{eq:v2v18-residual}.
\end{proof}

\begin{assessment}[The strong-source gap is spectral]
\label{ass:v2v18-spectral-gap}
Connected clustering controls the coarse projection \(B_LT_\ell\), but it
does not control the ultraviolet residual \((I-B_L)T_\ell\).
The V14--V15 Einstein stability theorem asks for the complete source in
\(L^1_tH^{s-1}_x\).  Since this topology is stronger than \(L^1_tL^2_x\),
\cref{prop:v2v18-counterexample} rules out any derivation of that hypothesis
from finite-range clustering alone.
\end{assessment}

\subsection{Implication for the Higgs stress}
\label{sec:v2v18-higgs}

Let \(\mathcal T_n^{\mathrm{ren}}\) be the complete projected Higgs stress in
\eqref{eq:v2v17-renormalized-source}.  If its connected covariance obeys
\eqref{eq:v2v18-cluster}, V18 rigorously yields convergence of every fixed
coarse observable under the amplitude condition in
\cref{cor:v2v18-fixed-block}.  If its low-external-momentum kernel has \(q\)
vanishing moments, V18 also yields the Taylor factor
\((\ell_n/L)^q\).  Neither statement estimates the diagonal output shell
\(\Delta_n\mathcal T_n^{\mathrm{ren}}\).

\begin{assessment}[Revised Higgs gate after V18]
\label{ass:v2v18-revised-gate}
To close the existing strong Einstein source criterion, one must prove at
least one estimate that acts on the diagonal output shell.  Viable forms are:
\begin{enumerate}
 \item decay of the renormalized composite amplitude \(A_{\ell_n}\) strong
 enough to pay the \(H^{s-1}\) weight;
 \item an output-shell cancellation, such as a Ward or null-form multiplier
 that vanishes on the relevant characteristic set;
 \item a Duhamel divisor obtained from quantitative separation between the
 source modulation and the Einstein wave cone.
\end{enumerate}
Block averaging and local-jet subtraction remain valid inputs for the
low-frequency part, but they cannot replace this diagonal estimate.
\end{assessment}

\subsection{V18 conclusion and next step}
\label{sec:v2v18-conclusion}

A connected correlation length \(\ell\) gives the sharp block gain
\((\ell/L)^{d/2}\), and \(q\) removed local moments give the sharp Taylor gain
\((\ell/L)^q\).  Both gains require a genuine separation between the
observation scale and the cutoff cell.  A smooth independent-cell model
shows that clustering can coexist with order-one \(L^2\) mass and
\(\ell^{-r}\) growth in \(H^r\).  Therefore the strong source convergence
claimed as the remaining V17 gate cannot follow from clustering by itself.

The next version will go upstream to the hyperbolic transfer step.  It will
derive the exact modulation-dependent Duhamel estimate for an
output-frequency shell and isolate the near-cone term on which a Higgs,
Yang--Mills, or Ward null structure is actually required.
\clearpage
\section*{Second Volume, Version V19: Modulation Divisors and the Higgs Near-Cone Sector}
\addcontentsline{toc}{section}{Second Volume, Version V19: Modulation Divisors and the Higgs Near-Cone Sector}

\subsection*{Purpose}

V18 proved that spatial clustering and local-jet removal do not control the
diagonal ultraviolet output shell.  This version moves upstream to the
linear hyperbolic transfer.  An exact half-wave normal form gives an inverse
modulation factor away from the Einstein characteristic cone.  For
quadratic massive Higgs waves, elementary mass-shell geometry then confines
the difficult same-sign interaction to a nearly parallel angular sector.
The result separates the part that gains a derivative from the part on which
a bilinear, Ward, or null-form estimate is genuinely needed.

\subsection{The half-wave normal form}
\label{sec:v2v19-half-wave}

Let \(\omega(D)\) be a positive self-adjoint Fourier multiplier.  On a
spatial shell \(|\xi|\simeq\Lambda\), assume
\begin{equation}
 c\Lambda\le\omega(\xi)\le C\Lambda.
\label{eq:v2v19-elliptic-shell}
\end{equation}
For a spacetime function \(F\), define the time-Fourier Wiener norm
\begin{equation}
 \|F\|_{\mathcal W_tH^r_x}
 =
 \int_{\mathbb R}
 \|\widehat F(\tau,\cdot)\|_{H^r_x}\,\mathrm d\tau,
\label{eq:v2v19-wiener}
\end{equation}
with harmless Fourier-normalization constants suppressed.

\begin{theorem}[Exact off-cone Duhamel divisor]
\label{thm:v2v19-divisor}
Let \(F\) be Schwartz, spatially supported in
\(|\xi|\simeq\Lambda\), and suppose
\begin{equation}
 |\tau-\sigma\omega(\xi)|\ge\mu
 \quad\text{on }\operatorname{supp}\widehat F
 \quad\text{for both }\sigma\in\{+1,-1\}.
\label{eq:v2v19-offcone}
\end{equation}
Let \(u\) solve
\[
 \partial_t^2u+\omega(D)^2u=F,
 \qquad
 u(0)=\partial_tu(0)=0.
\]
Then
\begin{equation}
 \sup_t\left(
 \|\partial_tu(t)\|_{H^{s-1}_x}
 +\|u(t)\|_{H^s_x}
 \right)
 \le
 C\mu^{-1}\|F\|_{\mathcal W_tH^{s-1}_x}.
\label{eq:v2v19-divisor}
\end{equation}
\end{theorem}

\begin{proof}
For \(\sigma=\pm1\), set
\[
 w_\sigma=\partial_tu+i\sigma\omega(D)u.
\]
Then
\[
 (\partial_t-i\sigma\omega(D))w_\sigma=F,
 \qquad
 w_\sigma(0)=0.
\]
Define \(G_\sigma\) by
\begin{equation}
 \widehat G_\sigma(\tau,\xi)
 =
 \frac{\widehat F(\tau,\xi)}
 {i(\tau-\sigma\omega(\xi))}.
\label{eq:v2v19-normal-form}
\end{equation}
Hypothesis \eqref{eq:v2v19-offcone} makes this multiplier regular and
\[
 \sup_t\|G_\sigma(t)\|_{H^{s-1}}
 \le C\mu^{-1}\|F\|_{\mathcal W_tH^{s-1}}.
\]
Since \(F=(\partial_t-i\sigma\omega(D))G_\sigma\), direct integration gives
the exact identity
\begin{equation}
 w_\sigma(t)
 =
 G_\sigma(t)-e^{i\sigma t\omega(D)}G_\sigma(0).
\label{eq:v2v19-boundary-normal-form}
\end{equation}
Unitarity of the half-wave group gives the same bound for \(w_\sigma\).
Finally,
\[
 \partial_tu=\frac12(w_++w_-),
 \qquad
 \omega(D)u=\frac1{2i}(w_+-w_-).
\]
The shell ellipticity \eqref{eq:v2v19-elliptic-shell} identifies
\(\|\omega(D)u\|_{H^{s-1}}\) with \(\|u\|_{H^s}\), proving
\eqref{eq:v2v19-divisor}.
\end{proof}

\begin{remark}[Why the boundary term must remain]
\label{rem:v2v19-boundary}
The second term in \eqref{eq:v2v19-boundary-normal-form} is the homogeneous
correction that enforces zero Cauchy data.  Dropping it would replace the
retarded solution by a different particular solution.  It has exactly the
same \(\mu^{-1}\) bound as the interior normal-form term.
\end{remark}

\subsection{A near-cone/far-cone criterion}
\label{sec:v2v19-split}

Let \(Q_{<\mu}\) be a smooth spacetime Fourier cutoff to
\[
 \min_{\sigma=\pm1}|\tau-\sigma\omega(\xi)|<2\mu
\]
that equals one where the minimum is at most \(\mu\), and set
\(Q_{\ge\mu}=I-Q_{<\mu}\).

\begin{theorem}[Split hyperbolic source criterion]
\label{thm:v2v19-split}
For zero initial data and shell-supported \(F\),
\begin{align}
 \sup_t\left(
 \|\partial_tu(t)\|_{H^{s-1}}
 +\|u(t)\|_{H^s}
 \right)
 \le C\big(
 &\|Q_{<\mu}F\|_{L^1_tH^{s-1}_x}
 \nonumber\\
 &+\mu^{-1}
 \|Q_{\ge\mu}F\|_{\mathcal W_tH^{s-1}_x}
 \big).
\label{eq:v2v19-split}
\end{align}
\end{theorem}

\begin{proof}
The standard half-wave energy estimate gives the first term.  The Fourier
support of \(Q_{\ge\mu}F\) obeys
\eqref{eq:v2v19-offcone}, up to the fixed cutoff constants, so
\cref{thm:v2v19-divisor} gives the second term.  Add the two zero-data
solutions.
\end{proof}

\begin{corollary}[Dyadic metric-increment test]
\label{cor:v2v19-dyadic}
Let \(F_n\) have output frequency \(\Lambda_n\simeq2^n\).  Put
\[
 N_n=\|Q_{<\mu_n}F_n\|_{L^1_tL^2_x},
 \qquad
 W_n=\|Q_{\ge\mu_n}F_n\|_{\mathcal W_tL^2_x}.
\]
If
\begin{equation}
 \sum_n\Lambda_n^{s-1}
 \left(N_n+\mu_n^{-1}W_n\right)<\infty,
\label{eq:v2v19-dyadic}
\end{equation}
then the associated zero-data linear metric increments are absolutely
summable in
\(C_tH^s_x\cap C^1_tH^{s-1}_x\).
When \(\mu_n\simeq\Lambda_n\), the far-cone term costs only
\(\Lambda_n^{s-2}W_n\), one power less than the direct source estimate.
\end{corollary}

\begin{proof}
On an output shell,
\(\|f\|_{H^{s-1}}\simeq\Lambda_n^{s-1}\|f\|_2\).
Apply \cref{thm:v2v19-split} and sum.
\end{proof}

\begin{firewall}[The divisor does not estimate the near cone]
\label{fw:v2v19-near}
The \(\mu^{-1}\) factor is available only for
\(Q_{\ge\mu}F\).  Assigning it to the full source silently divides by zero on
the characteristic set.  The near-cone quantity \(N_n\) in
\eqref{eq:v2v19-dyadic} must be bounded by an independent angular,
multilinear, Ward, or null-form argument.
\end{firewall}

\subsection{Massive Higgs pair geometry}
\label{sec:v2v19-higgs-geometry}

Write
\[
 \omega_m(p)=\sqrt{|p|^2+m^2},
 \qquad m>0.
\]
First consider the sum of two future massive four-momenta.  Let
\(a=|p|\), \(b=|q|\), and let \(\theta\in[0,\pi]\) be the angle between
\(p\) and \(q\).

\begin{theorem}[Same-sign Higgs modulation]
\label{thm:v2v19-same-sign}
Assume
\[
 c_0\Lambda\le a,b\le C_0\Lambda,
 \qquad
 \Lambda\ge m.
\]
For
\[
 \tau=\omega_m(p)+\omega_m(q),
 \qquad
 \xi=p+q,
\]
one has
\begin{equation}
 \tau-|\xi|
 \ge
 c\left(
 \frac{m^2}{\Lambda}
 +\Lambda(1-\cos\theta)
 \right)
 \ge
 c'\left(
 \frac{m^2}{\Lambda}
 +\Lambda\theta^2
 \right)
\label{eq:v2v19-angular-modulation}
\end{equation}
when \(0\le\theta\le1\); for \(\theta\ge1\), the lower bound is
\(c\Lambda\).
\end{theorem}

\begin{proof}
The exact factorization is
\begin{align}
 (\omega_m(p)+\omega_m(q))^2-|p+q|^2
 &=
 2m^2+2(\omega_m(p)\omega_m(q)-p\cdot q)
 \nonumber\\
 &=
 2m^2+
 2(\omega_m(p)\omega_m(q)-ab)
 +2ab(1-\cos\theta).
\label{eq:v2v19-factorization}
\end{align}
Since
\[
 (a^2+m^2)(b^2+m^2)-(ab+m^2)^2
 =m^2(a-b)^2\ge0,
\]
one has
\(\omega_m(p)\omega_m(q)-ab\ge m^2\).  Thus the right side of
\eqref{eq:v2v19-factorization} is at least
\(4m^2+2ab(1-\cos\theta)\).  Divide by
\(\omega_m(p)+\omega_m(q)+|p+q|\le C\Lambda\).
The elementary bounds
\(1-\cos\theta\simeq\theta^2\) for \(\theta\le1\) and
\(1-\cos\theta\ge c\) for \(\theta\ge1\) finish the proof.
\end{proof}

\begin{corollary}[Quantitative same-sign near-cone cap]
\label{cor:v2v19-cap}
If the interaction in \cref{thm:v2v19-same-sign} lies within modulation
\(\mu\) of the positive Einstein cone, then necessarily
\begin{equation}
 \mu\gtrsim\frac{m^2}{\Lambda},
 \qquad
 \theta\lesssim\left(\frac{\mu}{\Lambda}\right)^{1/2}.
\label{eq:v2v19-cap}
\end{equation}
In particular, below a constant multiple of \(m^2/\Lambda\) the same-sign
near-cone sector is empty.
\end{corollary}

\begin{proof}
Apply \eqref{eq:v2v19-angular-modulation}.
\end{proof}

Opposite time signs correspond to a difference of two future massive
four-momenta.

\begin{proposition}[Difference geometry]
\label{prop:v2v19-difference}
For
\[
 \tau=\omega_m(p)-\omega_m(q),
 \qquad
 \xi=p-q,
\]
one has
\begin{equation}
 |\xi|^2-\tau^2
 =
 2\big(\omega_m(p)\omega_m(q)-p\cdot q-m^2\big)
 \ge
 2ab(1-\cos\theta).
\label{eq:v2v19-difference}
\end{equation}
Hence the output is spacelike or zero and cannot be a nonzero exact null
four-vector.  Nevertheless, at high frequency its distance to the null cone
can tend to zero along nearly collinear configurations.
\end{proposition}

\begin{proof}
Expansion gives the identity.  The inequality follows again from
\(\omega_m(p)\omega_m(q)-ab\ge m^2\).  Equality in the full expression
forces equality in the future-timelike Cauchy inequality and
\(\theta=0\), so the two future mass-shell vectors are identical and their
difference is zero.  Taking parallel momenta of increasing magnitude shows
that nonzero differences can approach the null cone, so no
cutoff-independent positive modulation gap follows.
\end{proof}

\begin{assessment}[What mass does and does not buy]
\label{ass:v2v19-mass}
Positive Higgs mass removes exact nonzero quadratic resonance with the
Einstein cone.  It does not give a shell-sized modulation gap: the
same-sign lower bound degenerates like \(m^2/\Lambda\), and difference
interactions can approach the cone along collinear high-frequency rays.
Thus mass alone cannot provide the \(\mu_n\simeq\Lambda_n\) divisor needed
for a full derivative gain.
\end{assessment}

\subsection{The remaining near-cone estimate}
\label{sec:v2v19-remaining}

For a same-sign quadratic Higgs source at comparable input frequencies,
\cref{cor:v2v19-cap} converts the near-cone term into an angular cap.  If a
bilinear estimate or a transverse stress multiplier supplies
\begin{equation}
 \|Q_{<\mu_n}F_n\|_{L^1_tL^2_x}
 \le
 C\left(\frac{\mu_n}{\Lambda_n}\right)^\delta
 B_n
\label{eq:v2v19-needed-cap}
\end{equation}
for some \(\delta>0\), then
\cref{cor:v2v19-dyadic} reduces the full increment to
\begin{equation}
 \sum_n\Lambda_n^{s-1}
 \left[
 \left(\frac{\mu_n}{\Lambda_n}\right)^\delta B_n
 +\mu_n^{-1}W_n
 \right].
\label{eq:v2v19-balance}
\end{equation}
The free choice of \(\mu_n\) can then balance the cap gain against the
off-cone divisor.  Equation \eqref{eq:v2v19-needed-cap} is deliberately
stated as a gate, not as a theorem: it depends on the complete improved
Higgs stress, its probability norm, and the angular regularity of its
connected kernels.

\begin{assessment}[Revised transfer gate after V19]
\label{ass:v2v19-transfer-gate}
The diagonal source need not satisfy the V14 strong norm term by term if its
far-cone portion satisfies the Fourier-Wiener hypothesis of
\cref{thm:v2v19-divisor}.  The remaining task is narrower:
\begin{enumerate}
 \item prove a shell bound \(W_n\) for the fully renormalized stress;
 \item compute the contraction of the improved Higgs stress with the
 Einstein characteristic projector;
 \item determine whether that contraction or an angular bilinear estimate
 proves \eqref{eq:v2v19-needed-cap};
 \item repeat the characteristic analysis for the Yang--Mills and fermionic
 stress sectors.
\end{enumerate}
\end{assessment}

\subsection{V19 conclusion and next step}
\label{sec:v2v19-conclusion}

The half-wave normal form gives an exact inverse modulation with its required
zero-data boundary correction.  It recovers one ultraviolet power on a
shell-sized far-cone region.  Massive Higgs pair geometry excludes nonzero
exact quadratic resonance, but the gap collapses on nearly collinear
high-frequency configurations.  The unresolved source has therefore been
localized to a precise characteristic and angular sector.

V20 is an audit version.  It will compare this gate with the current
Yang--Mills and Navier--Stokes notes, then compute the characteristic
projection of the improved scalar stress or replace that computation by the
strongest common upstream estimate exposed by the audit.
\clearpage
\section*{Second Volume, Version V20: Companion Audit and the Transverse Higgs Stress Symbol}
\addcontentsline{toc}{section}{Second Volume, Version V20: Companion Audit and the Transverse Higgs Stress Symbol}

\subsection*{Purpose and mandatory audit}

This is the required five-version audit following V15.  The shared-folder
companions inspected before the new calculation were:
\begin{center}
\begin{tabular}{p{0.47\textwidth}r p{0.32\textwidth}}
\toprule
file & bytes & SHA--256\\
\midrule
\texttt{Renewal\_NS\_Stability\_Research\_Notes\_V31.tex}
&887537&
\texttt{F1121B62238AD5491BB7CF03635EA9934942EDF573ED8FAAA9EE79E1D38A6696}\\
\texttt{Renewal\_Yang\_Mills\_Mass\_Gap\_Research\_Notes\_V15.tex}
&189055&
\texttt{E0856774D8A230FD6853FD585971CE09CC52969B30B29E704798186A42DC8313}\\
\bottomrule
\end{tabular}
\end{center}
The folder also contains a file named
\texttt{Renewal\_Yang\_Mills\_Mass\_Gap\_Research\_Notes\_V14.tex}
with the same byte count and digest as the named V15 file.  The audited
V15-named source itself ends with the compact V14 record, so no unrecorded
V15 theorem was inferred from its filename.  The frozen YM and SM legacy
volumes were left untouched.

NS V31 proves a sharp scale-matched heat-formation ladder.  Heat smoothing
does not create the missing critical mark; applying a weaker source norm and
restoring shell weights returns the continuation-class quantity.  The YM
compact V14 note closes the retentive non-Cartan spectator sector but keeps
the strongly cancelling sector open.  Its reusable rule is to retain the
complete coefficient-weighted positive carrier until the capacity or
cancellation has acted.

\begin{assessment}[V20 audit transfer]
\label{ass:v2v20-audit-transfer}
V19's modulation divisor is subject to the same discipline as the NS heat
gain: it may be charged only to the far-cone source norm on which it was
proved.  The near-cone part must be kept in its angular and polarization
carrier until the characteristic stress symbol has acted, just as the YM
cancelling output must remain inside its coefficient-weighted carrier.
Taking the full stress norm first destroys the cancellation sought below.
\end{assessment}

\subsection{The improved scalar stress and physical polarization}
\label{sec:v2v20-stress}

Work first in flat spacetime and isolate the principal scalar--scalar part of
one Higgs component.  With improvement coefficient \(\zeta\), write
\begin{equation}
 T_{\mu\nu}
 =
 \partial_\mu\phi\,\partial_\nu\phi
 -\frac12\eta_{\mu\nu}
 \big(\partial^\alpha\phi\,\partial_\alpha\phi+m^2\phi^2\big)
 +\zeta(\eta_{\mu\nu}\Box-\partial_\mu\partial_\nu)\phi^2.
\label{eq:v2v20-improved-stress}
\end{equation}
For a complex Higgs multiplet, polarization of the quadratic form and
summation over its finitely many real components gives the same symbol
componentwise.

Let \(k=(\tau,\xi)\), with \(\xi\ne0\).  A physical radiative polarization is
represented in the spatial transverse-traceless gauge by a symmetric tensor
\(e=(e^{ij})\) satisfying
\begin{equation}
 e^{ij}\xi_j=0,
 \qquad
 \delta_{ij}e^{ij}=0,
 \qquad
 |e|_{\mathrm F}=1,
 \qquad
 e^{0\mu}=0.
\label{eq:v2v20-TT}
\end{equation}

\begin{theorem}[Exact TT principal symbol of the improved Higgs stress]
\label{thm:v2v20-TT-symbol}
Let two scalar Fourier factors have signed spatial momenta \(p\) and \(q_s\)
with
\[
 \xi=p+q_s.
\]
The bilinear principal symbol obtained by polarizing
\eqref{eq:v2v20-improved-stress} and contracting with \(e\) is
\begin{equation}
 \mathfrak m_e(p,q_s)
 =
 e^{ij}(p_i(q_s)_j+(q_s)_ip_j)
 =
 -2e^{ij}(p_\perp)_i(p_\perp)_j,
\label{eq:v2v20-exact-symbol}
\end{equation}
where
\[
 p_\perp
 =
 p-\frac{p\cdot\xi}{|\xi|^2}\xi.
\]
In particular,
\begin{equation}
 |\mathfrak m_e(p,q_s)|\le2|p_\perp|^2.
\label{eq:v2v20-symbol-bound}
\end{equation}
The mass, potential, metric-trace, and improvement terms make no
contribution to this TT symbol.
\end{theorem}

\begin{proof}
Every term proportional to \(\delta_{ij}\) vanishes by tracelessness.  The
Fourier multiplier of the improvement term on the product \(\phi^2\) is a
linear combination of \(\delta_{ij}k^2\) and \(\xi_i\xi_j\); its contraction
with \(e\) vanishes by both identities in \eqref{eq:v2v20-TT}.  The remaining
polarized derivative term is the first expression in
\eqref{eq:v2v20-exact-symbol}.  Since \(q_s=\xi-p\), symmetry and
transversality give
\[
 e^{ij}(p_i(q_s)_j+(q_s)_ip_j)
 =
 -2e^{ij}p_ip_j.
\]
Replacing each \(p\) by \(p_\perp\) does not change the contraction.  The
Frobenius Cauchy--Schwarz inequality proves
\eqref{eq:v2v20-symbol-bound}.
\end{proof}

\begin{remark}[Improvement independence has limited scope]
\label{rem:v2v20-improvement}
The coefficient \(\zeta\) disappears from the physical TT principal symbol,
not from the complete stress tensor.  Improvement remains relevant to the
trace sector, local counterterms, curved-background conservation, and the
longitudinal subsidiary system.
\end{remark}

\subsection{Angular form of the TT cancellation}
\label{sec:v2v20-angle}

For same-sign spatial addition \(q_s=q\), let \(\theta\) be the angle between
\(p\) and \(q\).  The exact geometric identity
\begin{equation}
 |p_\perp|
 =
 \frac{|p\wedge q|}{|p+q|}
 =
 \frac{|p||q|\sin\theta}{|p+q|}
\label{eq:v2v20-perp}
\end{equation}
holds.  For a difference interaction, put \(q_s=-q\); the same formula holds
with \(|p-q|\) in the denominator.

\begin{corollary}[Quadratic angular null factor]
\label{cor:v2v20-angular-null}
Suppose
\[
 |p|\simeq|q|\simeq|\xi|\simeq\Lambda.
\]
For either signed interaction \(\xi=p\pm q\),
\begin{equation}
 |\mathfrak m_e(p,\pm q)|
 \le C\Lambda^2\theta^2,
\label{eq:v2v20-angular-null}
\end{equation}
where in the difference case \(\theta\) is the angle between \(p\) and \(q\).
For the same-sign massive near-cone sector of
\cref{cor:v2v19-cap},
\begin{equation}
 |\mathfrak m_e(p,q)|
 \le C\Lambda\mu
 \quad\text{whenever}\quad
 0<\tau-|\xi|\le\mu.
\label{eq:v2v20-near-symbol}
\end{equation}
\end{corollary}

\begin{proof}
The shell comparability and \eqref{eq:v2v20-perp} give
\(|p_\perp|\le C\Lambda|\sin\theta|\).  Apply
\eqref{eq:v2v20-symbol-bound}.  In the same-sign near-cone region,
\cref{cor:v2v19-cap} gives
\(\theta^2\le C\mu/\Lambda\), which proves
\eqref{eq:v2v20-near-symbol}.
\end{proof}

\begin{assessment}[The cancellation is characteristic]
\label{ass:v2v20-characteristic}
A generic quadratic derivative stress has symbol size \(O(\Lambda^2)\).
The physical TT projection reduces this to
\(O(\Lambda^2\theta^2)\).  V19 shows that the only high-frequency
same-sign massive interactions close to the Einstein cone are nearly
parallel, so the two structures meet on the same sector.  This is an actual
near-cone symbol gain and does not rely on quartic tail coercivity or spatial
block averaging.
\end{assessment}

\subsection{A positive carrier estimate before taking norms}
\label{sec:v2v20-carrier}

The YM audit requires the cancellation to act before absolute summation.
The following finite-volume formulation makes that order explicit.  On a
spatial torus, let \(a(\tau,p)\) and \(b(\tau,q)\) be scalar spacetime
Fourier amplitudes.  For a measurable interaction set \(\mathcal R\), define
the positive convolution carrier
\begin{equation}
 \mathcal C(\mathcal R)
 =
 \int_{\mathbb R}
 \left[
 \sum_\xi
 \left(
 \int_{\mathbb R}\sum_p
 \mathbf1_{\mathcal R}
 |a(\tau_1,p)|
 |b(\tau-\tau_1,\xi-p)|
 \,\mathrm d\tau_1
 \right)^2
 \right]^{1/2}
 \mathrm d\tau.
\label{eq:v2v20-positive-carrier}
\end{equation}
All sums may be restricted to
\(|p|\simeq|\xi-p|\simeq|\xi|\simeq\Lambda\).

\begin{theorem}[TT near/far carrier ledger]
\label{thm:v2v20-carrier}
Let \(\mathcal R_{\mathrm{near}}\) be the same-sign massive region with
Einstein modulation at most \(\mu\), and let
\(\mathcal R_{\mathrm{far}}\) be its complement at modulation at least
\(\mu\).  Denote their positive carriers by
\(\mathcal C_{\mathrm{near}}\) and
\(\mathcal C_{\mathrm{far}}\).  For the corresponding TT quadratic stress,
\begin{align}
 \|T_{\mathrm{near}}^{\mathrm{TT}}\|_{\mathcal W_tL^2_x}
 &\le C\Lambda\mu\,\mathcal C_{\mathrm{near}},
\label{eq:v2v20-near-carrier}\\
 \|T_{\mathrm{far}}^{\mathrm{TT}}\|_{\mathcal W_tL^2_x}
 &\le C\Lambda^2\,\mathcal C_{\mathrm{far}}.
\label{eq:v2v20-far-carrier}
\end{align}
On a time interval \(I\) of length \(T\), the zero-data radiative metric
response therefore obeys
\begin{align}
 \sup_{t\in I}\big(
 \|\partial_th^{\mathrm{TT}}(t)\|_{H^{s-1}}
 +\|h^{\mathrm{TT}}(t)\|_{H^s}
 \big)
 \le C\Lambda^{s-1}\big(
 T\Lambda\mu\,\mathcal C_{\mathrm{near}}
 +\mu^{-1}\Lambda^2\,\mathcal C_{\mathrm{far}}
 \big).
\label{eq:v2v20-response-ledger}
\end{align}
\end{theorem}

\begin{proof}
For each output Fourier coefficient, insert
\eqref{eq:v2v20-near-symbol} on the near region and the generic bound
\(|\mathfrak m_e|\le C\Lambda^2\) on the far region.  Triangle inequality
inside the input convolution, followed by the \(\ell^2_\xi\) norm and
integration in \(\tau\), gives
\eqref{eq:v2v20-near-carrier}--\eqref{eq:v2v20-far-carrier}.  Fourier
inversion gives
\[
 \|T_{\mathrm{near}}^{\mathrm{TT}}\|_{L^1(I;L^2_x)}
 \le
 CT\|T_{\mathrm{near}}^{\mathrm{TT}}\|_{\mathcal W_tL^2_x}.
\]
Apply the direct energy estimate to this term and
\cref{thm:v2v19-divisor} to the far term.  Finally use output-shell
equivalence
\(\|f\|_{H^{s-1}}\simeq\Lambda^{s-1}\|f\|_2\).
\end{proof}

\begin{corollary}[Conditional half-power balance]
\label{cor:v2v20-half-power}
If \(T\) is fixed and
\(\mathcal C_{\mathrm{near}}+\mathcal C_{\mathrm{far}}\le C_n\), then the
choice \(\mu\simeq(\Lambda/T)^{1/2}\), whenever it lies between the mass
floor and a fixed fraction of \(\Lambda\), gives
\begin{equation}
 \sup_{t\in I}\big(
 \|\partial_th^{\mathrm{TT}}(t)\|_{H^{s-1}}
 +\|h^{\mathrm{TT}}(t)\|_{H^s}
 \big)
 \le
 C T^{1/2}\Lambda^{s+1/2}C_n.
\label{eq:v2v20-half-power}
\end{equation}
This improves the generic direct derivative-stress ledger
\(CT\Lambda^{s+1}C_n\) by one half cutoff power.
\end{corollary}

\begin{proof}
Substitute the common carrier bound into
\eqref{eq:v2v20-response-ledger} and balance
\(T\Lambda\mu\) against \(\Lambda^2/\mu\).
\end{proof}

\begin{firewall}[Carrier control is still missing]
\label{fw:v2v20-carrier}
The positive quantity \(C_n\) in
\cref{cor:v2v20-half-power} is not yet bounded uniformly for the interacting,
renormalized Higgs measure.  The theorem proves the exact multiplier and
Duhamel gains conditional on that carrier.  Replacing
\(\mathcal C_{\mathrm{near}}\) by the norm of the unsplit raw stress would
discard the angular cancellation; treating \(C_n\) as cutoff-uniform would
assume the remaining probabilistic estimate.
\end{firewall}

\subsection{Scope within the full SM--Einstein source}
\label{sec:v2v20-scope}

The result applies to the principal scalar--scalar contribution to the
physical radiative metric modes.  Covariant derivatives also generate
scalar--gauge and gauge--gauge terms.  Fermionic stress has a different
first-order symbol.  The trace and longitudinal metric components are tied
to conservation and the harmonic subsidiary system, but they are not
estimated by the TT projection alone.

\begin{assessment}[Revised acquisition order after V20]
\label{ass:v2v20-acquisition}
The next useful steps are:
\begin{enumerate}
 \item prove an \(L^2\)-based angular bilinear carrier bound sharper than
 the absolute convolution quantity in
 \eqref{eq:v2v20-positive-carrier};
 \item insert the connected/renormalized Higgs coefficients and determine
 their shell power in that carrier;
 \item compute the corresponding characteristic symbols for the
 Yang--Mills and fermion stresses before summing sectors;
 \item feed the longitudinal remainder through the exact Ward and
 subsidiary identities already established in V13--V15.
\end{enumerate}
The NS audit rules out crediting the response operator with more than the
proved modulation power.  The YM audit requires all characteristic
cancellation to be used before a blockwise operator norm.
\end{assessment}

\subsection{V20 conclusion}
\label{sec:v2v20-conclusion}

The mandatory companion audit has been completed with exact file identities.
It confirms that response smoothing and representation cancellation must be
matched to the source sector before absolute values are taken.  Applying
that rule gives a new exact result: the improved Higgs stress has physical
TT symbol
\(-2e^{ij}(p_\perp)_i(p_\perp)_j\), independent of the improvement
coefficient.  Its near-cone same-sign massive sector gains the factor
\(\theta^2\lesssim\mu/\Lambda\).  Combined with V19's far-cone divisor, this
gives the conditional half-power response bound
\eqref{eq:v2v20-half-power}.

The full continuum theorem remains open because the interacting
renormalized carrier \(C_n\), the other Standard Model stress sectors, and
the nonradiative metric components still require estimates.  V21 will
attack the first item by proving the sharp angular-cap bilinear estimate and
testing whether it upgrades the absolute-carrier ledger.
\clearpage
\section*{Second Volume, Version V21: Angular-Cap Bilinear Gain for the Radiative Higgs Source}
\addcontentsline{toc}{section}{Second Volume, Version V21: Angular-Cap Bilinear Gain for the Radiative Higgs Source}

\subsection*{Purpose}

V20 found a quadratic angular zero in the transverse-traceless Higgs stress,
but estimated the remaining Fourier convolution by a fully absolute carrier.
This version proves the elementary sharp cap-volume estimate in three
spatial dimensions and combines it with that symbol.  For same-sign massive
half-waves of comparable frequency, the near-cone stress gains three powers
of the cap aperture: two from the TT symbol and one from the square root of
cap volume.  Balancing this estimate with V19's off-cone divisor improves the
radiative response ledger by three fifths of a cutoff power.

\subsection{The cap-volume convolution lemma}
\label{sec:v2v21-cap}

Let \(0<\alpha\le1\), \(\Lambda\ge1\), and fix constants
\(0<c_0<C_0<\infty\).  For a sign
\(\varepsilon\in\{+1,-1\}\), define
\[
 \xi=p+\varepsilon q.
\]
Restrict to the region
\begin{equation}
 c_0\Lambda\le |p|,|q|,|\xi|\le C_0\Lambda,
 \qquad
 \angle(p,q)\le\alpha.
\label{eq:v2v21-support}
\end{equation}
Denote this region by \(\mathcal R_{\Lambda,\alpha}^{\varepsilon}\).

\begin{lemma}[Fibre volume of a comparable angular interaction]
\label{lem:v2v21-fibre}
For each fixed \(\xi\), the set of \(p\in\mathbb R^3\) satisfying
\eqref{eq:v2v21-support}, with
\(q=\varepsilon(\xi-p)\), has Lebesgue measure at most
\begin{equation}
 C\Lambda^3\alpha^2.
\label{eq:v2v21-fibre-volume}
\end{equation}
\end{lemma}

\begin{proof}
Since
\[
 |p\wedge\xi|=|p\wedge q|,
\]
the angle \(\vartheta\) between \(p\) and \(\xi\) satisfies
\[
 \sin\vartheta
 \le
 \frac{|q|}{|\xi|}\sin\angle(p,q)
 \le C\alpha.
\]
Thus \(p\) lies in one of a bounded number of spherical caps about the
two possible axial directions determined by \(\xi\), each of aperture
\(C\alpha\), and its radius lies in an interval of length \(O(\Lambda)\).
A three-dimensional spherical cap of radius \(O(\Lambda)\), radial thickness
\(O(\Lambda)\), and aperture \(O(\alpha)\) has volume
\(O(\Lambda^3\alpha^2)\).
\end{proof}

\begin{theorem}[Angular bilinear convolution estimate]
\label{thm:v2v21-angular-convolution}
Let
\begin{equation}
 \widehat{B_{\Lambda,\alpha}(f,g)}(\xi)
 =
 \int_{\mathbb R^3}
 \mathbf1_{\mathcal R_{\Lambda,\alpha}^{\varepsilon}}
 m(p,\varepsilon(\xi-p),\xi)
 \widehat f(p)\widehat g(\varepsilon(\xi-p))
 \,\mathrm dp.
\label{eq:v2v21-bilinear}
\end{equation}
If \(|m|\le M\) on the displayed support, then
\begin{equation}
 \|B_{\Lambda,\alpha}(f,g)\|_{L^2_x}
 \le
 C M\Lambda^{3/2}\alpha
 \|f\|_{L^2_x}\|g\|_{L^2_x}.
\label{eq:v2v21-angular-convolution}
\end{equation}
\end{theorem}

\begin{proof}
For each \(\xi\), Cauchy--Schwarz in \(p\) and
\cref{lem:v2v21-fibre} give
\begin{align*}
 |\widehat B(\xi)|^2
 &\le
 M^2C\Lambda^3\alpha^2
 \int
 \mathbf1_{\mathcal R_{\Lambda,\alpha}^{\varepsilon}}
 |\widehat f(p)|^2
 |\widehat g(\varepsilon(\xi-p))|^2\,\mathrm dp.
\end{align*}
Integrate in \(\xi\), discard the remaining indicator, and change variables
\(q=\varepsilon(\xi-p)\).  Fubini gives
\[
 \int\!\!\int
 |\widehat f(p)|^2|\widehat g(q)|^2
 \,\mathrm dp\,\mathrm dq
 =
 \|f\|_2^2\|g\|_2^2
\]
up to the Fourier normalization.  Taking square roots proves
\eqref{eq:v2v21-angular-convolution}.
\end{proof}

\begin{corollary}[TT stress gains three angular powers]
\label{cor:v2v21-three-powers}
Under \eqref{eq:v2v21-support}, let \(m=\mathfrak m_e\) be the TT scalar
stress symbol of \cref{thm:v2v20-TT-symbol}.  Then
\begin{equation}
 \|B^{\mathrm{TT}}_{\Lambda,\alpha}(f,g)\|_2
 \le
 C\Lambda^{7/2}\alpha^3
 \|f\|_2\|g\|_2.
\label{eq:v2v21-three-powers}
\end{equation}
\end{corollary}

\begin{proof}
By \cref{cor:v2v20-angular-null},
\(|\mathfrak m_e|\le C\Lambda^2\alpha^2\).  Insert this value of \(M\)
in \eqref{eq:v2v21-angular-convolution}.
\end{proof}

\begin{assessment}[Origin of the exponent three]
\label{ass:v2v21-three}
Two angular powers in \eqref{eq:v2v21-three-powers} are the physical TT
null symbol.  The third is the square root of the two-dimensional spherical
cap measure.  It is available in an \(L^2\) convolution estimate and was
absent from V20's coefficientwise absolute carrier.
\end{assessment}

\subsection{Mass-shell tubes imply the angular restriction}
\label{sec:v2v21-tubes}

Let \(f_+\) and \(g_+\) have spacetime Fourier support in positive massive
half-wave tubes
\begin{equation}
 |\tau_1-\omega_m(p)|\le\nu,
 \qquad
 |\tau_2-\omega_m(q)|\le\nu,
\label{eq:v2v21-input-tubes}
\end{equation}
and assume all input and output spatial frequencies are comparable to
\(\Lambda\).  Their product has
\(\tau=\tau_1+\tau_2\) and \(\xi=p+q\).

\begin{lemma}[Near Einstein modulation forces a cap]
\label{lem:v2v21-tube-cap}
Suppose \(\nu\le\mu/4\), \(\mu\le c\Lambda\), and
\[
 |\tau-|\xi||\le\mu.
\]
Then, unless the interaction set is empty,
\begin{equation}
 \mu\gtrsim\frac{m^2}{\Lambda},
 \qquad
 \angle(p,q)\le C\left(\frac{\mu}{\Lambda}\right)^{1/2}.
\label{eq:v2v21-tube-cap}
\end{equation}
\end{lemma}

\begin{proof}
The tube conditions imply
\[
 \omega_m(p)+\omega_m(q)-|\xi|
 \le
 |\tau-|\xi||+2\nu
 \le\frac32\mu.
\]
Apply \cref{thm:v2v19-same-sign}.
\end{proof}

\subsection{Spacetime Wiener estimate}
\label{sec:v2v21-wiener}

For compactly supported spacetime Fourier amplitudes, retain the norm
\(\mathcal W_tL^2_x\) from \eqref{eq:v2v19-wiener}.

\begin{theorem}[Near-cone TT Higgs bilinear estimate]
\label{thm:v2v21-near}
Under \eqref{eq:v2v21-input-tubes} and the hypotheses of
\cref{lem:v2v21-tube-cap}, the same-sign near-cone TT stress satisfies
\begin{equation}
 \|Q_{<\mu}T^{\mathrm{TT}}_{\Lambda}(f_+,g_+)\|
 _{\mathcal W_tL^2_x}
 \le
 C\Lambda^2\mu^{3/2}
 \|f_+\|_{\mathcal W_tL^2_x}
 \|g_+\|_{\mathcal W_tL^2_x}.
\label{eq:v2v21-near}
\end{equation}
Without angular restriction, the generic estimate is
\begin{equation}
 \|T^{\mathrm{TT}}_{\Lambda}(f,g)\|_{\mathcal W_tL^2_x}
 \le
 C\Lambda^{7/2}
 \|f\|_{\mathcal W_tL^2_x}
 \|g\|_{\mathcal W_tL^2_x}.
\label{eq:v2v21-generic}
\end{equation}
\end{theorem}

\begin{proof}
Fix the two input temporal frequencies \(\tau_1,\tau_2\).  On the near-cone
output, \cref{lem:v2v21-tube-cap} permits
\[
 \alpha=C(\mu/\Lambda)^{1/2}.
\]
Apply \eqref{eq:v2v21-three-powers} to the corresponding spatial Fourier
slices:
\[
 \|T^{\mathrm{TT}}_{\mathrm{near}}(\tau_1,\tau_2)\|_{L^2_x}
 \le
 C\Lambda^{7/2}
 \left(\frac{\mu}{\Lambda}\right)^{3/2}
 \|\widehat f_+(\tau_1)\|_2
 \|\widehat g_+(\tau_2)\|_2.
\]
The cutoff \(Q_{<\mu}\) can only reduce the positive support restriction used
in the proof.  Integrate in the output temporal frequency and use Fubini,
equivalently Young's \(L^1_\tau\) convolution inequality.  This proves
\eqref{eq:v2v21-near}.  For \eqref{eq:v2v21-generic}, use
\(|\mathfrak m_e|\le C\Lambda^2\) and
\cref{thm:v2v21-angular-convolution} with \(\alpha=1\), then repeat the
temporal argument.
\end{proof}

\begin{remark}[Empty sub-mass modulation sector]
\label{rem:v2v21-empty}
When \(\mu<c m^2/\Lambda\), the same-sign massive near-cone sector is empty
by \cref{lem:v2v21-tube-cap}; its contribution is zero rather than bounded
by extrapolating \eqref{eq:v2v21-near} below the mass floor.
\end{remark}

\subsection{Optimized radiative response}
\label{sec:v2v21-response}

\begin{theorem}[Three-fifths radiative gain]
\label{thm:v2v21-response}
Let \(I\) have length \(T>0\), and let the two positive massive input
half-waves obey
\[
 \|f_+\|_{\mathcal W_tL^2_x}
 \|g_+\|_{\mathcal W_tL^2_x}
 \le A_\Lambda.
\]
For the comparable-frequency same-sign TT source, split at modulation
\(\mu\) as in V19.  Then
\begin{align}
 \sup_{t\in I}\big(
 \|\partial_th^{\mathrm{TT}}_\Lambda(t)\|_{H^{s-1}}
 +\|h^{\mathrm{TT}}_\Lambda(t)\|_{H^s}
 \big)
 \le C A_\Lambda\Lambda^{s-1}
 \left(
 T\Lambda^2\mu^{3/2}
 +\Lambda^{7/2}\mu^{-1}
 \right).
\label{eq:v2v21-response}
\end{align}
If
\[
 \mu=T^{-2/5}\Lambda^{3/5}
\label{eq:v2v21-optimal-mu}
\]
lies in the admissible range
\(cm^2/\Lambda\le\mu\le c'\Lambda\) and dominates the input tube width,
then
\begin{equation}
 \sup_{t\in I}\big(
 \|\partial_th^{\mathrm{TT}}_\Lambda(t)\|_{H^{s-1}}
 +\|h^{\mathrm{TT}}_\Lambda(t)\|_{H^s}
 \big)
 \le
 C T^{2/5}\Lambda^{s+19/10}A_\Lambda.
\label{eq:v2v21-three-fifths}
\end{equation}
For fixed \(T\), this is a gain of \(3/5\) of a cutoff power relative to the
generic direct ledger \(C_T\Lambda^{s+5/2}A_\Lambda\).
\end{theorem}

\begin{proof}
Fourier inversion gives
\[
 \|Q_{<\mu}T^{\mathrm{TT}}\|_{L^1(I;L^2_x)}
 \le
 CT\|Q_{<\mu}T^{\mathrm{TT}}\|_{\mathcal W_tL^2_x}.
\]
Use \cref{thm:v2v21-near} and the direct energy estimate on this part.
For the far-cone part, combine \eqref{eq:v2v21-generic} with
\cref{thm:v2v19-divisor}.  The output \(H^{s-1}\) weight is
\(\Lambda^{s-1}\), proving \eqref{eq:v2v21-response}.  Balancing its two
parenthesized terms gives
\(\mu^{5/2}=T^{-1}\Lambda^{3/2}\), which is
\eqref{eq:v2v21-optimal-mu}.  Substitution gives
\eqref{eq:v2v21-three-fifths}.  The generic estimate follows by applying
\eqref{eq:v2v21-generic} directly on \(I\).
\end{proof}

\begin{corollary}[Dyadic summability gate for this sector]
\label{cor:v2v21-summability}
For fixed \(T\) and dyadic \(\Lambda_n\), the comparable-frequency,
same-sign, radiative scalar increments are absolutely summable in
\(C_tH^s_x\cap C^1_tH^{s-1}_x\) if
\begin{equation}
 \sum_n\Lambda_n^{s+19/10}A_{\Lambda_n}<\infty,
\label{eq:v2v21-summability}
\end{equation}
provided the tube and modulation hypotheses of
\cref{thm:v2v21-response} hold uniformly.
\end{corollary}

\begin{proof}
Sum \eqref{eq:v2v21-three-fifths}.
\end{proof}

\begin{firewall}[This is one physical interaction sector]
\label{fw:v2v21-sector}
The exponent in \eqref{eq:v2v21-three-fifths} is proved for comparable
spatial frequencies, positive massive input tubes, and the physical TT
projection.  It does not cover high--low interactions, opposite time signs,
broadly modulated interacting inputs, scalar--gauge terms, fermions,
Yang--Mills stress, or longitudinal metric components.  It also does not
supply the decay of \(A_\Lambda\); that remains a renormalized probabilistic
input.
\end{firewall}

\begin{assessment}[What V21 changes]
\label{ass:v2v21-change}
V20's absolute carrier used only the two-power TT symbol and produced a
conditional half-power balance.  V21 retains the angular restriction inside
the \(L^2\) convolution and gains the square root of cap measure.  The
resulting three-fifths response gain is an operator theorem for the specified
mass-shell sector.  It lowers, but does not eliminate, the ultraviolet decay
required from the connected renormalized Higgs amplitudes.
\end{assessment}

\subsection{V21 conclusion and next step}
\label{sec:v2v21-conclusion}

The same-sign massive Higgs near-cone region occupies a cap of aperture
\((\mu/\Lambda)^{1/2}\).  In three dimensions, its square-root volume gives
one angular power, while the TT stress gives two more.  The exact near
estimate is \(\Lambda^2\mu^{3/2}\) in the time-Fourier Wiener source norm.
Balancing it with the far-cone Duhamel divisor yields a three-fifths cutoff
gain in the radiative metric response.

The next step is to analyze opposite-sign interactions with nonzero output
frequency and high--low configurations.  Their resonance geometry differs
from the same-sign cap, and they must be separated before the Higgs carrier
can be assembled.
\clearpage
\section*{Second Volume, Version V22: Difference Resonances and High--Low TT Transfer}
\addcontentsline{toc}{section}{Second Volume, Version V22: Difference Resonances and High--Low TT Transfer}

\subsection*{Purpose}

V21 treated comparable positive-frequency Higgs pairs.  This version closes
the corresponding opposite-sign geometric regression and proves a stronger
high--low symbol estimate.  For comparable difference interactions with
nonzero comparable output, near-null modulation again forces a narrow
angular cap.  For high--low interactions, TT transversality transfers the
entire quadratic symbol to the low spatial momentum, eliminating every
positive power of the high input frequency from the bilinear source bound.

\subsection{Exact difference modulation}
\label{sec:v2v22-difference}

Let
\[
 a=|p|,\qquad b=|q|,\qquad
 \omega_a=\sqrt{a^2+m^2},\qquad
 \omega_b=\sqrt{b^2+m^2},
\]
and let \(\theta=\angle(p,q)\).  Put
\[
 r=|p-q|,
 \qquad
 \delta_- = r-|\omega_a-\omega_b|.
\]
By \cref{prop:v2v19-difference}, \(\delta_-\ge0\).

\begin{lemma}[Exact radial--angular factorization]
\label{lem:v2v22-factorization}
One has
\begin{align}
 \delta_-
 &=
 \frac{2}{r+|\omega_a-\omega_b|}
 \left[
 \frac{m^2(a-b)^2}
 {\omega_a\omega_b+ab+m^2}
 +ab(1-\cos\theta)
 \right].
\label{eq:v2v22-factorization}
\end{align}
\end{lemma}

\begin{proof}
The difference of squares is
\[
 r^2-(\omega_a-\omega_b)^2
 =
 2(\omega_a\omega_b-ab\cos\theta-m^2).
\]
The radial part has the exact rationalization
\begin{align*}
 \omega_a\omega_b-ab-m^2
 &=
 \frac{
 (\omega_a\omega_b)^2-(ab+m^2)^2
 }{\omega_a\omega_b+ab+m^2}\\
 &=
 \frac{m^2(a-b)^2}
 {\omega_a\omega_b+ab+m^2}.
\end{align*}
Factor the difference of squares by
\(r+|\omega_a-\omega_b|\).
\end{proof}

\begin{theorem}[Comparable difference near-cone cap]
\label{thm:v2v22-difference-cap}
Assume
\begin{equation}
 c_0\Lambda\le a,b,r\le C_0\Lambda,
 \qquad
 \Lambda\ge m.
\label{eq:v2v22-comparable}
\end{equation}
Then
\begin{equation}
 \delta_-
 \ge
 c\left[
 \frac{m^2(a-b)^2}{\Lambda^3}
 +\Lambda(1-\cos\theta)
 \right].
\label{eq:v2v22-difference-lower}
\end{equation}
If \(\delta_-\le\mu\le c_1\Lambda\), with \(c_1>0\) sufficiently small
depending only on \eqref{eq:v2v22-comparable}, then
\begin{equation}
 \theta\le C(\mu/\Lambda)^{1/2},
 \qquad
 |a-b|\ge c\Lambda,
 \qquad
 \mu\ge c\frac{m^2}{\Lambda}.
\label{eq:v2v22-difference-cap}
\end{equation}
\end{theorem}

\begin{proof}
In \eqref{eq:v2v22-factorization}, the denominator outside the brackets is
at most \(C\Lambda\), while
\(\omega_a\omega_b+ab+m^2\le C\Lambda^2\).  Also
\(ab\ge c\Lambda^2\).  This proves
\eqref{eq:v2v22-difference-lower}.  Its angular term gives the first
conclusion in \eqref{eq:v2v22-difference-cap}.

Next,
\[
 r^2=(a-b)^2+2ab(1-\cos\theta).
\]
The angular conclusion and \(\mu\le c_1\Lambda\) make the second term at
most \(Cc_1\Lambda^2\).  Choose \(c_1\) small compared with the lower
constant in \(r\ge c_0\Lambda\).  Then
\(|a-b|\ge c\Lambda\).  Substitute this into the radial term of
\eqref{eq:v2v22-difference-lower} to obtain
\(\delta_-\ge cm^2/\Lambda\), proving the last conclusion.
\end{proof}

\begin{corollary}[Opposite-sign half-wave tubes have the same cap]
\label{cor:v2v22-tube-cap}
Suppose
\[
 |\tau_1-\omega_m(p)|\le\nu,
 \qquad
 |\tau_2+\omega_m(q)|\le\nu,
 \qquad
 \nu\le\mu/4,
\]
and the spatial frequencies obey \eqref{eq:v2v22-comparable}.  If the output
\[
 (\tau,\xi)=(\tau_1+\tau_2,p-q)
\]
lies within \(\mu\) of either Einstein cone, then
\[
 \angle(p,q)\le C(\mu/\Lambda)^{1/2},
 \qquad
 \mu\gtrsim m^2/\Lambda,
\]
provided \(\mu\le c_1\Lambda\).
\end{corollary}

\begin{proof}
The tube widths imply
\[
 \big||\omega_m(p)-\omega_m(q)|-|\xi|\big|
 \le |\ |\tau|-|\xi|\ |+2\nu
 \le\frac32\mu.
\]
Apply \cref{thm:v2v22-difference-cap}, changing constants.
\end{proof}

\subsection{Opposite-sign TT bilinear response}
\label{sec:v2v22-opposite}

The signed spatial momentum of the negative-frequency factor is \(-q\), so
the output relation is \(\xi=p-q\).  The TT identity
\eqref{eq:v2v20-exact-symbol} remains
\[
 \mathfrak m_e(p,-q)
 =
 -2e^{ij}(p_\perp)_i(p_\perp)_j.
\]

\begin{theorem}[Opposite-sign analogue of V21]
\label{thm:v2v22-opposite}
Under the tube, comparability, and output-modulation assumptions of
\cref{cor:v2v22-tube-cap},
\begin{equation}
 \|Q_{<\mu}T^{\mathrm{TT}}_\Lambda(f_+,g_-)\|
 _{\mathcal W_tL^2_x}
 \le
 C\Lambda^2\mu^{3/2}
 \|f_+\|_{\mathcal W_tL^2_x}
 \|g_-\|_{\mathcal W_tL^2_x}.
\label{eq:v2v22-opposite-near}
\end{equation}
The far-cone term obeys the generic bound
\[
 \|Q_{\ge\mu}T^{\mathrm{TT}}_\Lambda(f_+,g_-)\|
 _{\mathcal W_tL^2_x}
 \le
 C\Lambda^{7/2}
 \|f_+\|_{\mathcal W_tL^2_x}
 \|g_-\|_{\mathcal W_tL^2_x}.
\]
Consequently the optimized metric response has the same
\(T^{2/5}\Lambda^{s+19/10}\) ledger as
\eqref{eq:v2v21-three-fifths}.
\end{theorem}

\begin{proof}
By \cref{cor:v2v22-tube-cap}, the near-cone inputs have angular aperture
\(\alpha=C(\mu/\Lambda)^{1/2}\).  Apply
\cref{cor:v2v21-three-powers} with the difference sign
\(\varepsilon=-1\), then integrate the temporal convolution exactly as in
the proof of \cref{thm:v2v21-near}.  The generic estimate and the optimized
Duhamel balance are identical to the proof of
\cref{thm:v2v21-response}.
\end{proof}

\begin{assessment}[Comparable quadratic Higgs pairs]
\label{ass:v2v22-comparable}
For nonzero output frequency comparable to the two input frequencies, both
same-sign and opposite-sign massive scalar pairs have now been reduced to
the same near/far estimate.  Exact mass prevents a nonzero resonance, the
near-null region is an angular cap, and the TT symbol vanishes quadratically
on its axis.
\end{assessment}

\subsection{High--low transversality}
\label{sec:v2v22-high-low}

Let
\[
 |p|\simeq|\xi|\simeq\Lambda,
 \qquad
 |q|\simeq\lambda,
 \qquad
 1\le\lambda\le c\Lambda,
 \qquad
 \xi=p+\varepsilon q,
\]
with \(\varepsilon=\pm1\).

\begin{theorem}[Both TT derivatives fall on the low scale]
\label{thm:v2v22-low-symbol}
For every physical TT polarization at output \(\xi\),
\begin{equation}
 |\mathfrak m_e(p,\varepsilon q)|
 \le C\lambda^2.
\label{eq:v2v22-low-symbol}
\end{equation}
\end{theorem}

\begin{proof}
Output transversality and
\(\xi=p+\varepsilon q\) imply
\[
 p_\perp=-\varepsilon q_\perp.
\]
Insert this identity in
\eqref{eq:v2v20-exact-symbol} and use
\(|q_\perp|\le|q|\lesssim\lambda\).
\end{proof}

\begin{theorem}[High--low TT bilinear bound]
\label{thm:v2v22-high-low}
Let \(f_\Lambda\) and \(g_\lambda\) have spatial Fourier support in the
displayed high and low annuli.  Then
\begin{equation}
 \|T^{\mathrm{TT}}(f_\Lambda,g_\lambda)\|_{L^2_x}
 \le
 C\lambda^{7/2}
 \|f_\Lambda\|_{L^2_x}\|g_\lambda\|_{L^2_x}.
\label{eq:v2v22-high-low}
\end{equation}
The same estimate holds in the time-Fourier Wiener norm:
\begin{align}
 \|T^{\mathrm{TT}}(f_\Lambda,g_\lambda)\|
 _{\mathcal W_tL^2_x}
 \le
 C\lambda^{7/2}
 \|f_\Lambda\|_{\mathcal W_tL^2_x}
 \|g_\lambda\|_{\mathcal W_tL^2_x}.
\label{eq:v2v22-high-low-wiener}
\end{align}
\end{theorem}

\begin{proof}
For fixed output \(\xi\), parameterize the convolution by the low momentum
\(q\).  Its support has volume \(O(\lambda^3)\).  Cauchy--Schwarz in \(q\),
followed by Fubini, gives the unweighted convolution estimate
\[
 \|f_\Lambda g_\lambda\|_2
 \le C\lambda^{3/2}\|f_\Lambda\|_2\|g_\lambda\|_2.
\]
Keep the multiplier inside the convolution and apply
\eqref{eq:v2v22-low-symbol}; this proves
\eqref{eq:v2v22-high-low}.  Apply the spatial estimate at every pair of
input temporal frequencies and use Young's \(L^1_\tau\) convolution
inequality to obtain \eqref{eq:v2v22-high-low-wiener}.
\end{proof}

\begin{corollary}[Summation of all lower Higgs inputs]
\label{cor:v2v22-low-sum}
Let \(P_\lambda\) be dyadic spatial projectors.  For every \(r>7/2\),
\begin{align}
 \sum_{\lambda\le c\Lambda}
 \lambda^{7/2}
 \|P_\lambda g\|_{\mathcal W_tL^2_x}
 \le
 C_r\|g\|_{\mathcal W_tH^r_x}.
\label{eq:v2v22-low-sum}
\end{align}
Consequently, on a time interval of length \(T\),
\begin{align}
 \sup_{t\in I}
 \|h^{\mathrm{TT}}_{\Lambda,\mathrm{high-low}}(t)\|_{H^s_x}
 \le
 C_{r,T}\Lambda^{s-1}
 \|f_\Lambda\|_{\mathcal W_tL^2_x}
 \|g\|_{\mathcal W_tH^r_x},
\label{eq:v2v22-high-low-response}
\end{align}
with the analogous \(C^1_tH^{s-1}_x\) bound.
\end{corollary}

\begin{proof}
Cauchy--Schwarz in the dyadic scale gives
\[
 \sum_\lambda\lambda^{7/2}
 \|P_\lambda g\|_{\mathcal W_tL^2}
 \le
 \left(\sum_\lambda\lambda^{-2(r-7/2)}\right)^{1/2}
 \left(
 \sum_\lambda\lambda^{2r}
 \|P_\lambda g\|_{\mathcal W_tL^2}^2
 \right)^{1/2}.
\]
The geometric sum is finite.  Minkowski's inequality bounds the second
factor by \(\|g\|_{\mathcal W_tH^r}\).  Sum
\eqref{eq:v2v22-high-low-wiener}, use Fourier inversion on \(I\), apply the
direct energy estimate, and insert the output weight
\(\Lambda^{s-1}\).
\end{proof}

\begin{assessment}[No high derivative is lost in this TT paraproduct]
\label{ass:v2v22-high-low}
A generic derivative product would expose a factor comparable to
\(\Lambda\lambda\).  In the TT projection, the exact output relation
replaces the high transverse momentum by the low transverse momentum, so
the symbol is \(O(\lambda^2)\).  The remaining
\(\lambda^{3/2}\) is the ordinary three-dimensional low-frequency Bernstein
factor.  Thus the high--low radiative Higgs source costs no positive power of
\(\Lambda\) beyond the output Sobolev weight.
\end{assessment}

\begin{firewall}[High--high to low remains separate]
\label{fw:v2v22-high-high-low}
When \(|p|\simeq|q|\simeq\Lambda\) but
\(|p\pm q|=\lambda\ll\Lambda\), the output is not in the diagonal shell
treated above.  TT transversality alone does not control the sum over all
high parents feeding one low output.  That sector must retain the V18 local
jet/connected subtraction and a high-parent summation estimate.  It cannot
be included in \eqref{eq:v2v22-high-low-response} by relabelling the output.
\end{firewall}

\subsection{V22 conclusion and next step}
\label{sec:v2v22-conclusion}

The exact radial term in the difference mass-shell identity shows that an
opposite-sign pair with comparable nonzero output has the same mass floor
and angular near-cone cap as a same-sign pair.  Hence the V21 three-fifths
radiative gain applies to both time-sign combinations in that regime.
For high--low inputs, the TT symbol is stronger: both transverse derivatives
are paid at the low scale, yielding
\(\lambda^{7/2}\) after the Fourier support estimate and no high-frequency
symbol loss.

The unresolved scalar quadratic sector is now high--high to low output.
V23 will combine the V18 vanishing-jet estimate with a high-parent dyadic
summation and determine the exact number of external Taylor cancellations
needed for that sector.
\clearpage
\section*{Second Volume, Version V23: High--High to Low Resonance and the Four-Jet Parent Gate}
\addcontentsline{toc}{section}{Second Volume, Version V23: High--High to Low Resonance and the Four-Jet Parent Gate}

\subsection*{Purpose}

V22 reduced the unassembled scalar quadratic source to high--high
interactions with much lower output frequency.  Same-sign inputs are far
from the Einstein cone, but opposite-sign inputs can be almost null when
their common high velocity points along the low output.  This version proves
the exact resonance identity, obtains the corresponding TT cap estimate,
and then asks how many external Taylor cancellations are needed to sum all
high parents.  Under an explicit factorized form-factor hypothesis, the
answer is four for the opposite-sign branch and three for the same-sign
branch.  Whether the SM renormalized stress actually supplies the required
four-jet factor is left as the next gate.

\subsection{Difference resonance at low output}
\label{sec:v2v23-low-output}

Let
\[
 |p|\simeq|q|\simeq\Lambda,
 \qquad
 \xi=p-q,
 \qquad
 |\xi|=\lambda\le c\Lambda,
\]
and write
\[
 \Delta\omega=\omega_m(p)-\omega_m(q),
 \qquad
 n=\xi/\lambda.
\]

\begin{lemma}[Exact average-velocity identity]
\label{lem:v2v23-velocity}
Define
\begin{equation}
 v_{p,q}
 =
 \frac{p+q}{\omega_m(p)+\omega_m(q)}.
\label{eq:v2v23-average-velocity}
\end{equation}
Then
\begin{equation}
 \Delta\omega=v_{p,q}\cdot\xi,
 \qquad
 \lambda-|\Delta\omega|
 =
 \lambda\big(1-|v_{p,q}\cdot n|\big).
\label{eq:v2v23-velocity-identity}
\end{equation}
Moreover \(|v_{p,q}|<1\).
\end{lemma}

\begin{proof}
Rationalizing the energy difference gives
\[
 \omega_m(p)-\omega_m(q)
 =
 \frac{|p|^2-|q|^2}{\omega_m(p)+\omega_m(q)}
 =
 \frac{(p+q)\cdot(p-q)}
 {\omega_m(p)+\omega_m(q)}.
\]
This is the first identity, and the second follows by
\(\xi=\lambda n\).  Finally,
\[
 |p+q|
 \le |p|+|q|
 <
 \omega_m(p)+\omega_m(q)
\]
because \(m>0\).
\end{proof}

\begin{theorem}[Low-output difference cap and TT symbol]
\label{thm:v2v23-low-cap}
If
\begin{equation}
 0\le\lambda-|\Delta\omega|\le\mu\le c\lambda,
\label{eq:v2v23-near-low}
\end{equation}
then
\begin{align}
 |(v_{p,q})_\perp|
 &\le (2\mu/\lambda)^{1/2},
\label{eq:v2v23-velocity-cap}\\
 |p_\perp|=|q_\perp|
 &\le C\Lambda(\mu/\lambda)^{1/2},
\label{eq:v2v23-input-cap}\\
 |\mathfrak m_e(p,-q)|
 &\le C\Lambda^2\mu/\lambda.
\label{eq:v2v23-low-symbol}
\end{align}
Here perpendicular components are taken relative to \(n\), and \(e\) is a
TT polarization at output \(\xi\).
\end{theorem}

\begin{proof}
From \eqref{eq:v2v23-velocity-identity},
\[
 1-|v_{p,q}\cdot n|\le\mu/\lambda.
\]
Since \(|v_{p,q}|\le1\),
\[
 |(v_{p,q})_\perp|^2
 \le1-|v_{p,q}\cdot n|^2
 \le2(1-|v_{p,q}\cdot n|),
\]
which proves \eqref{eq:v2v23-velocity-cap}.  The vector \(\xi=p-q\) has no
component perpendicular to \(n\), so \(p_\perp=q_\perp\).  Also
\[
 p_\perp=q_\perp
 =\frac12(p+q)_\perp
 =\frac{\omega_m(p)+\omega_m(q)}2(v_{p,q})_\perp.
\]
Frequency comparability proves \eqref{eq:v2v23-input-cap}.  Insert it into
the exact TT symbol \eqref{eq:v2v20-exact-symbol} to obtain
\eqref{eq:v2v23-low-symbol}.
\end{proof}

\begin{remark}[The small mass floor]
\label{rem:v2v23-mass-floor}
The strict inequality \(|v_{p,q}|<1\) excludes exact nonzero null resonance.
For comparable high inputs it obeys the quantitative lower bound
\(c m^2\lambda/\Lambda^2\), because
\((\omega_p+\omega_q)^2-|p+q|^2\ge4m^2\) while the corresponding
difference-of-squares denominator is \(O(\Lambda^2)\).  This is much
smaller than the \(m^2/\Lambda\) floor in the comparable-output regime.
No cutoff-uniform gap is available.
\end{remark}

\subsection{The low-output angular convolution}
\label{sec:v2v23-convolution}

\begin{lemma}[High-parent cap volume]
\label{lem:v2v23-parent-volume}
For fixed \(\xi\), the set of \(p\) satisfying
\[
 |p|\simeq|\xi-p|\simeq\Lambda
\]
and \(|p_\perp|\le C\Lambda\alpha\) has measure at most
\begin{equation}
 C\Lambda^3\alpha^2.
\label{eq:v2v23-parent-volume}
\end{equation}
\end{lemma}

\begin{proof}
Relative to the axis \(n=\xi/|\xi|\), the vector \(p\) lies in two caps of
aperture \(C\alpha\) and has radial thickness \(O(\Lambda)\).  The same
spherical-coordinate calculation as in
\cref{lem:v2v21-fibre} gives \eqref{eq:v2v23-parent-volume}.
\end{proof}

\begin{theorem}[Opposite-sign high--high/low TT estimate]
\label{thm:v2v23-high-high-low}
Let the two inputs lie in massive half-wave tubes of width
\(\nu\le\mu/4\), with opposite time signs, and suppose
\[
 |\xi|=\lambda\ll\Lambda,
 \qquad
 \big||\tau|-|\xi|\big|\le\mu\le c\lambda.
\]
Then
\begin{equation}
 \|Q_{<\mu}T^{\mathrm{TT}}_{\Lambda\to\lambda}(f_+,g_-)\|
 _{\mathcal W_tL^2_x}
 \le
 C\Lambda^{7/2}
 \left(\frac{\mu}{\lambda}\right)^{3/2}
 \|f_+\|_{\mathcal W_tL^2_x}
 \|g_-\|_{\mathcal W_tL^2_x}.
\label{eq:v2v23-high-high-low}
\end{equation}
\end{theorem}

\begin{proof}
The tube errors enlarge the right side of
\eqref{eq:v2v23-near-low} only by a fixed constant.  Thus
\cref{thm:v2v23-low-cap} applies with
\(\alpha=C(\mu/\lambda)^{1/2}\).  The symbol contributes
\(C\Lambda^2\alpha^2\), and the square root of the fibre volume in
\cref{lem:v2v23-parent-volume} contributes
\(C\Lambda^{3/2}\alpha\).  This gives
\(C\Lambda^{7/2}\alpha^3\) for each pair of temporal Fourier slices.
Young's \(L^1_\tau\) convolution inequality completes the proof.
\end{proof}

\begin{assessment}[Output scale controls the cap]
\label{ass:v2v23-output-cap}
In V21--V22 the cap aperture was
\((\mu/\Lambda)^{1/2}\) because the output was comparable to the inputs.
For high--high to low difference interactions it is
\((\mu/\lambda)^{1/2}\).  The high-parent volume and TT symbol therefore
retain the large prefactor \(\Lambda^{7/2}\); angular resonance analysis
alone does not sum the parents.
\end{assessment}

\subsection{Same-sign high parents are uniformly off cone}
\label{sec:v2v23-same-sign}

\begin{proposition}[Same-sign low output has shell-sized modulation]
\label{prop:v2v23-same-sign}
Let both input tubes have the same time sign and width
\(\nu\le c\Lambda\).  If
\(|p|,|q|\simeq\Lambda\) and \(|p+q|=\lambda\le c'\Lambda\), then
\begin{equation}
 \min_{\sigma=\pm1}
 |\tau-\sigma|\xi||
 \ge c''\Lambda
\label{eq:v2v23-same-sign-offcone}
\end{equation}
for sufficiently small fixed \(c,c'\).  Consequently, before any
renormalized jet gain,
\begin{equation}
 \sup_t E_s(h^{\mathrm{TT}}_{\Lambda\to\lambda})
 \le
 C\lambda^{s-1}\Lambda^{5/2}
 \|f\|_{\mathcal W_tL^2_x}
 \|g\|_{\mathcal W_tL^2_x},
\label{eq:v2v23-same-sign-response}
\end{equation}
where
\(E_s(h)=\|\partial_th\|_{H^{s-1}}+\|h\|_{H^s}\).
\end{proposition}

\begin{proof}
The output temporal frequency has magnitude at least \(c_0\Lambda\), up to
the tube error, whereas \(|\xi|=\lambda\le c'\Lambda\).  This proves
\eqref{eq:v2v23-same-sign-offcone}.  The generic TT bilinear source bound is
\(C\Lambda^{7/2}\) times the product of input Wiener norms.  Apply
\cref{thm:v2v19-divisor} with \(\mu\simeq\Lambda\) and insert the low-output
Sobolev weight \(\lambda^{s-1}\).
\end{proof}

\subsection{A factorized renormalized external jet}
\label{sec:v2v23-jet}

A local subtraction gives a high-parent gain only if it survives inside the
physical tensor structure.  The needed hypothesis can be stated without
ambiguity.

\begin{definition}[Factorized \(q\)-jet hypothesis]
\label{def:v2v23-factorized}
For the connected renormalized high-parent form factor, suppose the physical
TT multiplier has the form
\begin{equation}
 \mathfrak m^{\mathrm{ren}}_{\Lambda,e}(\xi;p,q)
 =
 \mathfrak m_e(p,\pm q)\,
 R_\Lambda(\xi;p,q),
\label{eq:v2v23-factorized}
\end{equation}
where, uniformly for \(|\xi|\le c\Lambda\),
\begin{equation}
 \partial_\xi^\beta R_\Lambda(0;p,q)=0
 \quad (|\beta|<q),
 \qquad
 |\partial_\xi^\beta R_\Lambda(\xi;p,q)|
 \le C_\beta\Lambda^{-q}
 \quad (|\beta|=q).
\label{eq:v2v23-jet-hypothesis}
\end{equation}
\end{definition}

\begin{lemma}[The factorized jet power]
\label{lem:v2v23-jet-power}
Under \eqref{eq:v2v23-jet-hypothesis},
\begin{equation}
 |R_\Lambda(\xi;p,q)|
 \le C(\lambda/\Lambda)^q
 \quad\text{when }|\xi|=\lambda.
\label{eq:v2v23-jet-power}
\end{equation}
\end{lemma}

\begin{proof}
Apply the multivariable Taylor formula with integral remainder along the
segment from \(0\) to \(\xi\).  All terms below degree \(q\) vanish, and the
degree-\(q\) remainder is bounded by
\(C|\xi|^q\Lambda^{-q}\).
\end{proof}

\begin{firewall}[Jet and TT gains do not multiply automatically]
\label{fw:v2v23-factorized}
V18 proves an external Taylor gain for a subtracted connected kernel.  V20
proves a TT angular zero for the scalar stress tensor.  Their product in
\eqref{eq:v2v23-factorized} is an additional structural hypothesis.  It
must follow from the Ward-compatible tensor decomposition of the
renormalized insertion; it cannot be obtained by multiplying two unrelated
upper bounds.
\end{firewall}

\subsection{High-parent summation and the integer threshold}
\label{sec:v2v23-parent-sum}

Put
\[
 A_\Lambda
 =
 \|f_\Lambda\|_{\mathcal W_tL^2_x}
 \|g_\Lambda\|_{\mathcal W_tL^2_x}.
\]

\begin{theorem}[Opposite-sign parent ledger under a factorized jet]
\label{thm:v2v23-opposite-parent}
Assume the factorized \(q\)-jet hypothesis.  On an interval of length \(T\),
the opposite-sign high-parent contribution satisfies
\begin{align}
 \sup_{t\in I}E_s(h^{\mathrm{TT}}_{\Lambda\to\lambda})
 \le
 C T^{2/5}
 \lambda^{s-8/5+q}
 \Lambda^{7/2-q}A_\Lambda,
\label{eq:v2v23-opposite-parent}
\end{align}
after choosing
\(\mu=T^{-2/5}\lambda^{3/5}\) in its admissible range.
Consequently, if \(A_\Lambda\le C\Lambda^{-a}\), the sum over dyadic
\(\Lambda\ge C\lambda\) converges precisely under the strict power condition
\begin{equation}
 q+a>\frac72.
\label{eq:v2v23-opposite-threshold}
\end{equation}
For cutoff-uniform parent amplitudes \(a=0\), the least admissible integer is
\begin{equation}
 q=4.
\label{eq:v2v23-four-jet}
\end{equation}
\end{theorem}

\begin{proof}
By \cref{lem:v2v23-jet-power}, multiply both the near estimate
\eqref{eq:v2v23-high-high-low} and the generic far estimate by
\((\lambda/\Lambda)^q\).  The direct near-cone energy term and off-cone
normal form give
\begin{align*}
 \sup_tE_s
 \le
 C\lambda^{s-1}
 \left(\frac{\lambda}{\Lambda}\right)^q
 \Lambda^{7/2}A_\Lambda
 \left[
 T\left(\frac{\mu}{\lambda}\right)^{3/2}
 +\mu^{-1}
 \right].
\end{align*}
Balancing the two bracketed terms gives
\(\mu=T^{-2/5}\lambda^{3/5}\) and the common factor
\(T^{2/5}\lambda^{-3/5}\).  This proves
\eqref{eq:v2v23-opposite-parent}.  The remaining parent series is
\[
 \sum_{\Lambda\ge C\lambda}\Lambda^{7/2-q-a},
\]
which converges on a dyadic grid exactly when its exponent is negative.
This is \eqref{eq:v2v23-opposite-threshold}, and
\eqref{eq:v2v23-four-jet} follows.
\end{proof}

\begin{proposition}[Same-sign parent threshold]
\label{prop:v2v23-same-parent}
Under the same factorized \(q\)-jet hypothesis, the same-sign response is
bounded by
\begin{equation}
 C\lambda^{s-1+q}\Lambda^{5/2-q}A_\Lambda.
\label{eq:v2v23-same-parent}
\end{equation}
If \(A_\Lambda\le C\Lambda^{-a}\), its dyadic high-parent sum converges when
\begin{equation}
 q+a>\frac52.
\label{eq:v2v23-same-threshold}
\end{equation}
Thus \(q=3\) is the least integer for cutoff-uniform parent amplitudes.
\end{proposition}

\begin{proof}
Multiply \eqref{eq:v2v23-same-sign-response} by the jet factor
\((\lambda/\Lambda)^q\) and sum the resulting dyadic power.
\end{proof}

\begin{corollary}[Uniform four-jet parent closure]
\label{cor:v2v23-four-jet}
A factorized four-jet bound closes the high-parent multiplicity of both time
signs when \(A_\Lambda\) is uniformly bounded.  After the parent sum, the
opposite-sign output ledger is
\begin{equation}
 C T^{2/5}\lambda^{s+19/10-a}
\label{eq:v2v23-output-ledger}
\end{equation}
when \(A_\Lambda\lesssim\Lambda^{-a}\).
\end{corollary}

\begin{proof}
Use \(q=4\) in
\cref{thm:v2v23-opposite-parent,prop:v2v23-same-parent}.  For the displayed
opposite-sign power, the first parent shell contributes
\[
 \lambda^{s-8/5+q}\lambda^{7/2-q-a}
 =
 \lambda^{s+19/10-a}.
\]
The remaining geometric series changes only the constant.
\end{proof}

\begin{assessment}[What four cancellations accomplish]
\label{ass:v2v23-four}
Four external cancellations remove the infinite multiplicity of uniformly
sized high parents feeding one lower output.  They do not make the output
shells summable: \eqref{eq:v2v23-output-ledger} still requires decay in the
renormalized parent amplitude.  The four-jet result closes an incidence
problem, not the complete ultraviolet source estimate.
\end{assessment}

\begin{firewall}[Renormalization freedom is not arbitrary]
\label{fw:v2v23-renormalization}
The programme may use \(q=4\) only if the local counterterm basis,
Slavnov--Taylor identities, stress conservation, and the chosen
renormalization conditions imply
\eqref{eq:v2v23-factorized}--\eqref{eq:v2v23-jet-hypothesis}.  One cannot
impose four vanishing derivatives merely because the parent sum requires
them.  If power counting permits fewer subtractions, the missing parent
decay must come from \(A_\Lambda\), an additional Ward zero, or a different
multiscale organization.
\end{firewall}

\subsection{V23 conclusion and next step}
\label{sec:v2v23-conclusion}

For an opposite-sign high--high pair producing a lower output, the exact
energy difference is the projection of the output onto the average massive
velocity.  Near the Einstein cone this forces the high momenta into caps of
aperture \((\mu/\lambda)^{1/2}\), but leaves a prefactor
\(\Lambda^{7/2}\).  Same-sign pairs are uniformly off cone and gain one
Duhamel power.  Under a Ward-compatible factorized external jet, four
vanishing orders are necessary and sufficient to sum uniformly bounded
opposite-sign high parents; three suffice for the same-sign branch.

V24 must now determine the legitimate Taylor degree of the renormalized
stress form factor from four-dimensional power counting and the available
local covariant counterterms.  That calculation decides whether the
four-jet gate is attainable or whether the programme must instead extract
parent decay from connected amplitudes.
\clearpage
\section*{Second Volume, Version V24: Ward Normalization Obstructs the Four-Jet Physical Stress}
\addcontentsline{toc}{section}{Second Volume, Version V24: Ward Normalization Obstructs the Four-Jet Physical Stress}

\subsection*{Purpose}

V23 proved that four factorized external cancellations would sum uniformly
bounded high parents in the difficult high--high to low, opposite-sign
sector.  It also required a check that those cancellations are legitimate.
This version performs that upstream check.  The translation Ward identity
fixes the spin-zero stress form factor at zero momentum transfer.  The
physical Higgs stress therefore has a nonvanishing zeroth jet.  After the
protected normalization is separated, the radiative correction is even and
starts at second order, but no symmetry forces the second-order coefficient
to vanish.  Thus the four-jet hypothesis cannot be assigned to the complete
physical stress.

\subsection{Spin-zero stress form factors}
\label{sec:v2v24-form-factor}

Let \(T_{\mu\nu}\) be a symmetric conserved stress tensor in Minkowski space,
and let \(|p\rangle\), \(|p'\rangle\) be scalar one-particle states of the
same mass.  Put
\[
 k=p'-p,
 \qquad
 P=\frac12(p'+p).
\]
On shell,
\begin{equation}
 P\cdot k=0.
\label{eq:v2v24-orthogonal}
\end{equation}

\begin{theorem}[Conserved spin-zero tensor decomposition]
\label{thm:v2v24-decomposition}
Away from exceptional singularities, Lorentz covariance, symmetry in
\(\mu,\nu\), and
\(k^\mu\langle p'|T_{\mu\nu}(0)|p\rangle=0\) imply
\begin{equation}
 \langle p'|T_{\mu\nu}(0)|p\rangle
 =
 2P_\mu P_\nu A(k^2)
 +\frac12(k_\mu k_\nu-\eta_{\mu\nu}k^2)D(k^2).
\label{eq:v2v24-form-factor}
\end{equation}
The scalar form factors depend on \(k^2\), and hence are even under
\(k\mapsto-k\).
\end{theorem}

\begin{proof}
The most general symmetric rank-two tensor made from
\(P,k,\eta\) is
\[
 aP_\mu P_\nu
 +b(P_\mu k_\nu+k_\mu P_\nu)
 +ck_\mu k_\nu+d\eta_{\mu\nu}.
\]
Contract with \(k^\mu\) and use \eqref{eq:v2v24-orthogonal}.  For
\(k^2\ne0\), the independent \(P_\nu\) and \(k_\nu\) coefficients give
\(b=0\) and \(d=-ck^2\).  Rename \(a=2A\), \(c=D/2\).
Continuity gives the decomposition at nonexceptional null transfer.  On
shell the independent Lorentz invariants are \(P^2\), fixed by the mass and
\(k^2\), and \(k^2\); thus the scalar coefficients are functions of \(k^2\).
\end{proof}

\begin{corollary}[Only the momentum form factor reaches TT radiation]
\label{cor:v2v24-TT}
For a physical polarization \(e^{\mu\nu}\) satisfying
\(e^{\mu\nu}k_\nu=0\) and \(e^\mu{}_\mu=0\),
\begin{equation}
 e^{\mu\nu}
 \langle p'|T_{\mu\nu}(0)|p\rangle
 =
 2A(k^2)e^{\mu\nu}P_\mu P_\nu.
\label{eq:v2v24-TT}
\end{equation}
The improvement form factor \(D\) is annihilated.
\end{corollary}

\begin{proof}
Contract \eqref{eq:v2v24-form-factor} with \(e\).
\end{proof}

\subsection{The protected value at zero transfer}
\label{sec:v2v24-normalization}

Let
\[
 \mathsf P_\nu=\int_{\mathbb R^3}T_{0\nu}(t,x)\,\mathrm d^3x
\]
be the translation generator.

\begin{theorem}[Translation Ward normalization]
\label{thm:v2v24-A-zero}
With the conventional normalization of one-particle states,
\begin{equation}
 A(0)=1.
\label{eq:v2v24-A-zero}
\end{equation}
Equivalently, the coefficient of \(P_\mu P_\nu\) in the zero-transfer
stress matrix element cannot be removed by a stress-tensor counterterm.
\end{theorem}

\begin{proof}
Use normalized wave packets to avoid products of momentum delta
distributions.  The generator identity
\[
 \mathsf P_\nu|p\rangle=p_\nu|p\rangle
\]
implies that the spatial integral of the \(T_{0\nu}\) matrix element at
zero transfer equals the one-particle momentum, with the same state
normalization as the inner product.  In
\eqref{eq:v2v24-form-factor}, the \(D\)-term vanishes at \(k=0\), while the
remaining term is \(2p_0p_\nu A(0)\).  Comparing with the generator
normalization gives \(A(0)=1\).  Since the statement holds for arbitrary
wave packets, it is independent of the distributional plane-wave
normalization.
\end{proof}

\begin{corollary}[The physical form factor has no vanishing external jet]
\label{cor:v2v24-no-jet}
Near zero momentum transfer,
\begin{equation}
 A(k^2)=1+A'(0)k^2+O(k^4).
\label{eq:v2v24-expansion}
\end{equation}
Thus the complete physical TT stress contains its unsuppressed tree
normalization.  It does not satisfy
\eqref{eq:v2v23-jet-hypothesis} for any \(q\ge1\).
\end{corollary}

\begin{proof}
Evenness follows from \cref{thm:v2v24-decomposition}, and the constant is
fixed by \cref{thm:v2v24-A-zero}.  The vanishing condition in
\eqref{eq:v2v23-jet-hypothesis} includes the value at \(k=0\), which is
nonzero here.
\end{proof}

\begin{firewall}[Normal ordering does not remove gravitational charge]
\label{fw:v2v24-normal-order}
Vacuum subtraction can remove
\(\langle0|T_{\mu\nu}|0\rangle\) and its local geometric counterterms.
It does not change the one-particle identity \(A(0)=1\).  Treating vacuum
normal ordering as a factor \((|k|/\Lambda)^q\) on every state-dependent
high--high/low stress matrix element would erase the translation generator
carried by physical excitations.
\end{firewall}

\subsection{What local power counting permits}
\label{sec:v2v24-power-counting}

\begin{proposition}[Degree of a two-Higgs stress insertion]
\label{prop:v2v24-degree}
In four dimensions, a one-particle-irreducible vertex with one insertion of
a dimension-four stress tensor and two external Higgs legs has superficial
momentum degree
\begin{equation}
 \delta=4-2[H]=2,
 \qquad [H]=1.
\label{eq:v2v24-degree}
\end{equation}
The covariant local operators of dimension at most four that can contribute
to this sector are generated, up to total derivatives and field
redefinitions, by
\begin{equation}
 \sqrt{|g|}\,g^{\mu\nu}(D_\mu H)^\dagger D_\nu H,
 \qquad
 \sqrt{|g|}\,m^2H^\dagger H,
 \qquad
 \sqrt{|g|}\,R\,H^\dagger H.
\label{eq:v2v24-local-basis}
\end{equation}
Their metric variations produce stress vertices of momentum degree at most
two.
\end{proposition}

\begin{proof}
The insertion has engineering dimension four and each external Higgs field
removes one momentum dimension, giving \(\delta=2\).  Gauge invariance,
diffeomorphism covariance, and dimension at most four leave the kinetic,
mass, and curvature-improvement terms displayed in
\eqref{eq:v2v24-local-basis}; the Higgs potential contributes through the
metric trace and does not change the derivative count.  Direct metric
variation shows that no listed vertex contains more than two derivatives.
\end{proof}

\begin{assessment}[Power counting and the TT decomposition agree]
\label{ass:v2v24-power-counting}
The kinetic counterterm fixes the constant multiplying
\(P_\mu P_\nu\), and the curvature improvement contributes to the
\(D\)-structure killed by TT projection.  A term proportional to
\(k^2P_\mu P_\nu\) has four external derivatives in the full vertex and
would arise from a dimension-six matter--curvature operator, not from the
minimal dimension-four basis.
\end{assessment}

\begin{theorem}[Maximum automatic jet of the normalized correction]
\label{thm:v2v24-two-jet}
Define the correction to the protected momentum form factor by
\[
 \Delta A(k^2)=A(k^2)-A(0).
\]
If \(A\) is \(C^2\) near zero and its high-parent scaling obeys
\[
 |A'(z)|\le C\Lambda^{-2}
 \quad (|z|\le c\Lambda^2),
\]
then
\begin{equation}
 |\Delta A(k^2)|
 \le C(|k|/\Lambda)^2.
\label{eq:v2v24-two-jet}
\end{equation}
No Ward identity or dimension-four counterterm in
\eqref{eq:v2v24-local-basis} forces \(A'(0)=0\).
\end{theorem}

\begin{proof}
The fundamental theorem of calculus gives
\[
 A(k^2)-A(0)=k^2\int_0^1A'(tk^2)\,\mathrm dt,
\]
which proves \eqref{eq:v2v24-two-jet}.  Conservation fixed the tensor
decomposition and \(A(0)\), but leaves the derivative of the scalar form
factor free.  The local basis contains no coefficient whose TT variation
can adjust \(k^2P_\mu P_\nu\).
\end{proof}

\begin{remark}[A fourth-order remainder requires extra input]
\label{rem:v2v24-fourth}
If one separately subtracts the physical slope \(A'(0)k^2\), then evenness
would leave an \(O(k^4/\Lambda^4)\) remainder.  Such a subtraction fixes a
higher-derivative gravitational radius and corresponds to enlarging the
theory by dimension-six data.  It is not a consequence of minimal SM
renormalization or of the translation Ward identity.
\end{remark}

\subsection{Correction of the V23 parent route}
\label{sec:v2v24-correction}

\begin{theorem}[Legitimate high-parent thresholds]
\label{thm:v2v24-thresholds}
Apply the parent criterion
\eqref{eq:v2v23-opposite-threshold}.
\begin{enumerate}
 \item For the complete physical momentum form factor, the protected
 constant means \(q=0\); high-parent summability requires
 \begin{equation}
  a>\frac72.
 \label{eq:v2v24-full-threshold}
 \end{equation}
 \item For the separately normalized correction \(\Delta A\), the automatic
 even jet gives \(q=2\); high-parent summability requires
 \begin{equation}
  a>\frac32.
 \label{eq:v2v24-correction-threshold}
 \end{equation}
 \item A \(q=4\) bound follows only after an additional slope subtraction or
 an independent dynamical zero.
\end{enumerate}
For the same-sign far-cone branch, the corresponding thresholds from
\eqref{eq:v2v23-same-threshold} are \(a>5/2\) for the full form factor and
\(a>1/2\) for its normalized correction.
\end{theorem}

\begin{proof}
Substitute \(q=0\) and \(q=2\) into
\eqref{eq:v2v23-opposite-threshold} and
\eqref{eq:v2v23-same-threshold}.  The final assertion is
\cref{rem:v2v24-fourth}.
\end{proof}

\begin{assessment}[The four-jet gate is not available for the full stress]
\label{ass:v2v24-four-jet-status}
V23's algebraic statement remains correct under its explicit factorized
four-jet hypothesis.  V24 proves that the hypothesis fails for the complete
physical scalar stress because it contradicts \(A(0)=1\).  Even after the
protected part is split off, standard symmetry and power counting supply
two low-transfer powers, not four.  The missing high-parent decay must
therefore be proved for the state-dependent amplitude or paid by a different
source topology.
\end{assessment}

\subsection{Mean source versus stress fluctuations}
\label{sec:v2v24-mean-fluctuation}

The Ward obstruction also forces a conceptual split.
\begin{enumerate}
 \item The renormalized expectation
 \(\langle T_{\mu\nu}\rangle_{\mathrm{ren}}\) may have vacuum-local pieces
 absorbed into cosmological, Newton, curvature-squared, and improvement
 coefficients.  State-dependent energy still retains its zero-transfer
 gravitational charge.
 \item The centred operator
 \(T_{\mu\nu}-\langle T_{\mu\nu}\rangle\) has zero mean, but its two-point
 fluctuation does not inherit a vanishing one-particle form factor.  Its
 regularity must be estimated as an operator-valued distribution or random
 distribution.
 \item A pathwise positive-Sobolev source theorem for the fluctuation cannot
 be inferred from renormalization of the mean.
\end{enumerate}

\begin{firewall}[Semiclassical and stochastic equations are different]
\label{fw:v2v24-semiclassical}
The semiclassical Einstein equation is sourced by a renormalized
expectation.  An Einstein equation driven by stress fluctuations is a
different stochastic or quantum metric problem.  Mixing the mean
counterterm freedom with a pathwise fluctuation estimate would falsely
assign local subtractions to state-dependent gravitational charge.
\end{firewall}

\subsection{V24 conclusion and next step}
\label{sec:v2v24-conclusion}

The physical spin-zero stress matrix element has the conserved decomposition
\eqref{eq:v2v24-form-factor}; its TT projection retains only the momentum
form factor \(A\), and translation symmetry fixes \(A(0)=1\).  The complete
stress therefore has no vanishing low-transfer jet.  The normalized loop
correction starts at order \(k^2/\Lambda^2\), while a fourth-order remainder
would require subtraction of a physical slope or a new dynamical zero.
Consequently V23's four-jet route is not available for the full minimal
Higgs stress.

The programme must now decide the source concept before proceeding.  V25 is
both the next compulsory companion audit and a free-field regularity test:
it will compute the Sobolev regularity of the Wick-ordered scalar stress
fluctuation and compare it with the V14 strong Einstein source space.  That
test will determine whether the existing pathwise continuum target is
mathematically compatible even with the free ultraviolet fixed point.
\clearpage
\section*{Second Volume, Version V25: Companion Audit and the Free-Stress Equal-Time No-Go}
\addcontentsline{toc}{section}{Second Volume, Version V25: Companion Audit and the Free-Stress Equal-Time No-Go}

\subsection*{Purpose and mandatory audit}

This is the required audit five versions after V20.  The current companion
sources inspected were:
\begin{center}
\begin{tabular}{p{0.47\textwidth}r p{0.32\textwidth}}
\toprule
file & bytes & SHA--256\\
\midrule
\texttt{Renewal\_NS\_Stability\_Research\_Notes\_V31.tex}
&887537&
\texttt{F1121B62238AD5491BB7CF03635EA9934942EDF573ED8FAAA9EE79E1D38A6696}\\
\texttt{Renewal\_Yang\_Mills\_Mass\_Gap\_Research\_Notes\_V19.tex}
&252931&
\texttt{DAB2C84B0CDBE6E93E512896B36100926CC05525908E49068B06D3BDEB149691}\\
\bottomrule
\end{tabular}
\end{center}
The NS source and digest are unchanged from V20.  Its V31 warning therefore
remains: a source-to-response gain must be used only in its scale-matched
source space, and a weaker source norm cannot be promoted by restoring
cutoff weights after the estimate.

The YM companion advanced from the compact V14 content audited at V20 to
V19.  It now proves an output-count-free capacity estimate for a complete
dual-pair Wilson moment, selects a dominant complete bank, and closes the
all-nonlocal incidence-two dual-pair sector outside output priority.  Its
central transferable rule is stronger than the earlier one: assemble the
complete positive output operator before resolving and summing channels.
Nondual banks, higher incidence, output priority, and the local-dominant
non-Cartan sector remain open.

\begin{assessment}[V25 audit transfer]
\label{ass:v2v25-audit}
For stress fluctuations, the analogue of the YM complete positive operator
is the full two-point or noise kernel.  Componentwise cancellations may be
used only before a positive norm of that kernel is taken.  However, if one
physical TT scalar component already has a divergent positive variance, no
estimate of its separate cutoff sequence may declare convergence by
cancelling it against a deterministic mean counterterm.  The following free
test addresses precisely that positive component.
\end{assessment}

\subsection{Cutoff free scalar and TT stress}
\label{sec:v2v25-free}

At one fixed time in three spatial dimensions, let \(\phi_\Lambda\) be the
centred Gaussian free scalar with Fourier covariance
\begin{equation}
 \mathbb E[
 \widehat\phi_\Lambda(p)
 \overline{\widehat\phi_\Lambda(q)}]
 =
 (2\pi)^3\delta(p-q)
 \frac{\chi_\Lambda(p)^2}{2\omega_m(p)},
 \qquad
 \omega_m(p)=\sqrt{|p|^2+m^2},
\label{eq:v2v25-covariance}
\end{equation}
where \(\chi_\Lambda\) equals one on \(|p|\le\Lambda\) and vanishes for
\(|p|\ge2\Lambda\).

Choose a smooth symmetric spatial test tensor \(f^{ij}\) whose Fourier
transform is supported in a small ball away from \(k=0\), and assume
\begin{equation}
 k_i\widehat f^{ij}(k)=0,
 \qquad
 \delta_{ij}\widehat f^{ij}(k)=0.
\label{eq:v2v25-test-TT}
\end{equation}
Define the Wick-ordered pairing
\begin{equation}
 X_\Lambda
 =
 \int_{\mathbb R^3}
 f^{ij}(x)
 :T_{ij}[\phi_\Lambda]:(0,x)
 \,\mathrm d^3x.
\label{eq:v2v25-X}
\end{equation}

\begin{lemma}[Only the derivative TT term survives]
\label{lem:v2v25-TT-reduction}
For the improved scalar stress
\eqref{eq:v2v20-improved-stress},
\begin{equation}
 X_\Lambda
 =
 \int f^{ij}(x)
 :\partial_i\phi_\Lambda(x)\partial_j\phi_\Lambda(x):
 \,\mathrm d^3x.
\label{eq:v2v25-TT-reduction}
\end{equation}
The mass, potential, metric-trace, and improvement terms vanish in the
pairing.
\end{lemma}

\begin{proof}
Metric-trace terms vanish by the second condition in
\eqref{eq:v2v25-test-TT}.  Integrating the improvement term twice by parts
gives a contraction with
\(\delta_{ij}\Delta f^{ij}-\partial_i\partial_jf^{ij}\), which is zero by
both conditions in \eqref{eq:v2v25-test-TT}.  The remaining term is
\eqref{eq:v2v25-TT-reduction}.  Wick subtraction changes only its mean.
\end{proof}

\subsection{Exact second-chaos variance}
\label{sec:v2v25-variance}

\begin{theorem}[Equal-time TT stress variance grows like the fifth power]
\label{thm:v2v25-variance}
There exists a nonzero test tensor satisfying
\eqref{eq:v2v25-test-TT} and constants \(0<c<C<\infty\) such that
\begin{equation}
 c\Lambda^5
 \le
 \mathbb E|X_\Lambda|^2
 \le
 C\Lambda^5
\label{eq:v2v25-variance}
\end{equation}
for all sufficiently large \(\Lambda\).
\end{theorem}

\begin{proof}
Wick's formula gives, up to a fixed Fourier-normalization constant,
\begin{align}
 \mathbb E|X_\Lambda|^2
 =
 \int_{\mathbb R^3} \mathrm d^3k
 \int_{\mathbb R^3}\mathrm d^3p\,
 \frac{
 \chi_\Lambda(p)^2\chi_\Lambda(k-p)^2
 \left|
 \widehat f^{ij}(k)p_i(k-p)_j
 \right|^2
 }{
 \omega_m(p)\omega_m(k-p)
 }.
\label{eq:v2v25-variance-integral}
\end{align}
Transversality gives the exact cancellation
\begin{equation}
 \widehat f^{ij}(k)p_i(k-p)_j
 =
 -\widehat f^{ij}(k)p_ip_j.
\label{eq:v2v25-symbol}
\end{equation}

Choose a small ball \(B\) about one nonzero \(k_0\), a smooth TT
polarization \(e^{ij}(k)\) on \(B\), and
\(\widehat f^{ij}(k)=\eta(k)e^{ij}(k)\), where
\(\eta\) is nonzero and compactly supported in \(B\).  Since \(e(k)\) is a
nonzero traceless symmetric form on the plane \(k^\perp\), there is,
uniformly after shrinking \(B\), a set of directions on the sphere of
positive measure for which
\begin{equation}
 |e^{ij}(k)p_ip_j|\ge c|p|^2.
\label{eq:v2v25-direction-lower}
\end{equation}
For \(k\in B\) and \(1\ll|p|\le c\Lambda\),
\(\omega_m(p)\omega_m(k-p)\simeq|p|^2\).  On the directions in
\eqref{eq:v2v25-direction-lower}, the integrand in
\eqref{eq:v2v25-variance-integral} is bounded below by \(c|p|^2\).
Thus
\[
 \mathbb E|X_\Lambda|^2
 \ge
 c\int_1^{c\Lambda}r^2r^2\,\mathrm dr
 \ge c'\Lambda^5.
\]
For the upper bound,
\[
 |\widehat f^{ij}(k)p_ip_j|
 \le C_f|p|^2,
 \qquad
 \omega_m(p)\omega_m(k-p)\ge c(1+|p|)^2
\]
on the fixed compact \(k\)-support, apart from a bounded low-frequency
region.  The cutoff restricts \(|p|\le C\Lambda\), so the same radial
integral gives \(C\Lambda^5\).
\end{proof}

\begin{corollary}[Wick subtraction does not make a fixed-time distribution]
\label{cor:v2v25-not-tight}
The family \(X_\Lambda\) is not tight.  Consequently the equal-time
Wick-ordered TT stress does not converge in probability in
\(\mathcal D'(\mathbb R^3)\), and in particular does not converge in
\(H^r_{\mathrm{loc}}\) for any \(r\in\mathbb R\).
\end{corollary}

\begin{proof}
Each \(X_\Lambda\) lies in the second homogeneous Gaussian chaos.
Hypercontractivity at fixed chaos order gives
\begin{equation}
 \mathbb E|X_\Lambda|^4
 \le C_2(\mathbb E|X_\Lambda|^2)^2.
\label{eq:v2v25-fourth}
\end{equation}
Apply Paley--Zygmund to \(X_\Lambda^2\):
\[
 \mathbb P\left(
 |X_\Lambda|
 \ge2^{-1/2}(\mathbb E|X_\Lambda|^2)^{1/2}
 \right)
 \ge\frac1{4C_2}.
\]
The threshold tends to infinity by
\cref{thm:v2v25-variance}, so the scalar pairings are not tight.
Convergence in probability in \(\mathcal D'\), or in any local Sobolev
space, would imply tightness of every fixed smooth test pairing.
\end{proof}

\begin{remark}[A shell-difference version]
\label{rem:v2v25-shell}
The same proof restricted to
\(c\Lambda\le|p|\le C\Lambda\) gives a lower bound
\(c\Lambda^5\) for a newly added ultraviolet annulus.  The failure is
therefore a genuine shell failure, not merely accumulation of a bounded
error over many cutoffs.
\end{remark}

\begin{firewall}[Deterministic counterterms cannot cancel this variance]
\label{fw:v2v25-counterterm}
Subtracting a deterministic local mean changes \(X_\Lambda\) by a number and
does not change its centred second-chaos variance.  The positive quantity
\eqref{eq:v2v25-variance-integral} can be altered only by changing the
operator-valued distribution, its smearing topology, or the state.  It
cannot be made cutoff-Cauchy by the vacuum-energy counterterm used in a
semiclassical expectation.
\end{firewall}

\subsection{Why spacetime smearing is different}
\label{sec:v2v25-time}

Let \(\Phi\) now denote the free quantum scalar field in the vacuum
representation.  For \(g\in\mathcal S(\mathbb R)\), define
\begin{equation}
 Y_\Lambda(g,f)
 =
 \int_{\mathbb R^{1+3}}
 g(t)f^{ij}(x)
 :T_{ij}[\Phi_\Lambda]:(t,x)
 \,\mathrm dt\,\mathrm d^3x.
\label{eq:v2v25-spacetime}
\end{equation}

\begin{theorem}[Temporal smearing restores the free vacuum observable]
\label{thm:v2v25-time-smearing}
For every Schwartz \(g\) and TT Schwartz \(f\), the vectors
\(Y_\Lambda(g,f)|0\rangle\) are Cauchy in Fock space as
\(\Lambda\to\infty\).
\end{theorem}

\begin{proof}
Normal ordering implies that only the two-creation part contributes on the
vacuum.  Its momentum amplitude contains
\begin{equation}
 \widehat g(\omega_m(p)+\omega_m(q))\,
 \widehat f^{ij}(p+q)\,
 \frac{p_iq_j}{\sqrt{\omega_m(p)\omega_m(q)}}
\label{eq:v2v25-pair-amplitude}
\end{equation}
up to fixed constants and the harmless improvement structures.  The
polynomial factor grows at most polynomially in \(|p|+|q|\).
For every \(N\),
\[
 |\widehat g(\omega_m(p)+\omega_m(q))|
 \le C_N(1+|p|+|q|)^{-N}.
\]
Choosing \(N\) larger than the polynomial degree plus the six-dimensional
integration volume gives an integrable square majorant for
\eqref{eq:v2v25-pair-amplitude}.  Dominated convergence as the two cutoff
functions tend to one proves the Fock-space Cauchy property.
\end{proof}

\begin{assessment}[The obstruction is the sharp time trace]
\label{ass:v2v25-time-trace}
The covariant Wick stress is meaningful after spacetime smearing, while its
restriction to a sharp time slice fails even after smooth spatial smearing.
The oscillatory energy sum
\(\omega_m(p)+\omega_m(q)\) supplies arbitrarily strong decay only when a
time test function is retained.  An \(L^1_tH^{s-1}_x\) norm takes absolute
values after forming time slices and therefore cannot use this covariant
smearing mechanism.
\end{assessment}

\subsection{Consequence for the SM--Einstein target}
\label{sec:v2v25-consequence}

\begin{theorem}[Free-field incompatibility of the pathwise strong-source
target]
\label{thm:v2v25-incompatibility}
The V14--V15 strong source hypothesis
\[
 T_\Lambda\ \hbox{Cauchy in }L^1_tH^{s-1}_x
\]
cannot hold for the unsmeared quantum Higgs stress interpreted as a
pathwise equal-time random distribution.  The obstruction already occurs
for the free scalar TT component and for a fixed smooth spatial test.
\end{theorem}

\begin{proof}
Evolve the cutoff Gaussian initial data by the free stationary dynamics and
write \(Z_\Lambda(t)=\langle T_\Lambda(t),f\rangle\).  Its law is
independent of \(t\), and
\(\sigma_\Lambda^2:=\mathbb E|Z_\Lambda(t)|^2\ge c\Lambda^5\).
Fixed-chaos hypercontractivity and Paley--Zygmund give a constant \(c_0>0\)
such that
\[
 \mathbb P(|Z_\Lambda(t)|\ge\sigma_\Lambda/\sqrt2)\ge c_0
\]
for every \(t\).  Let \(Y_\Lambda\) be the measure of such times in a
fixed interval \(I\).  Fubini gives
\(\mathbb EY_\Lambda\ge c_0|I|\); since
\(0\le Y_\Lambda\le|I|\), one has
\(\mathbb P(Y_\Lambda\ge c_0|I|/2)\ge c_0/2\).  On that event,
\[
 \int_I|Z_\Lambda(t)|\,\mathrm dt
 \ge c|I|\sigma_\Lambda\longrightarrow\infty.
\]
Thus these scalar pairings are not tight in \(L^1_t\).  A smooth spatial
pairing is bounded by every proposed \(H^{s-1}_x\) source norm, so the
Bochner-space cutoff sequence cannot be Cauchy in probability.  The
\(L^2\)-based compiler of V16--V17 already fails by the same variance.
\end{proof}

\begin{assessment}[Three mathematically distinct continuum branches]
\label{ass:v2v25-branches}
The programme must now keep the following targets separate.
\begin{enumerate}
 \item In semiclassical Einstein theory, the source is the renormalized
 expectation \(\langle T_{\mu\nu}\rangle_{\mathrm{ren}}\).  The V25
 fluctuation variance is not inserted pathwise.
 \item In an Einstein--Langevin theory, the centred stress enters through a
 spacetime noise kernel.  The metric must be constructed in a distributional
 spacetime topology that retains temporal smearing.
 \item In a fully quantum theory, both stress and metric are
 operator-valued distributions; classical harmonic-gauge stability in
 \(C_tH^s_x\) is not the closing theorem.
\end{enumerate}
The intended \emph{SM--Einstein continuum} must specify which branch is
meant before the remaining source estimates can be logically assembled.
\end{assessment}

\begin{firewall}[No interacting conclusion is inferred]
\label{fw:v2v25-interacting}
V25 is a free ultraviolet regression.  It proves that the previous
pathwise strong-source target is too strong for the standard free quantum
stress.  It does not construct the interacting SM stress, prove existence of
a semiclassical solution, construct an Einstein--Langevin metric, or
quantize gravity.  Interactions cannot be credited with curing the free
positive variance without a separate theorem.
\end{firewall}

\subsection{V25 conclusion and next step}
\label{sec:v2v25-conclusion}

The mandatory audit records the unchanged NS V31 state and the new YM V19
complete-bank aggregation theorem.  The free scalar calculation then shows
that a fixed-time, spatially TT-smeared Wick stress has variance
\(\asymp\Lambda^5\).  Fixed-chaos hypercontractivity turns this into
non-tightness, so there is no equal-time random-distribution limit.  A
Schwartz time smearing restores convergence of the free vacuum observable.

This result forces an upstream revision.  V26 will take the semiclassical
branch first, because it remains compatible with a classical Einstein
metric.  It will formulate a cutoff-independent Hadamard-renormalized
expectation source, identify the exact conservation conditions needed by
V13--V15, and state a conditional continuum theorem without importing the
fluctuation no-go into the mean equation.
\clearpage
\section*{Second Volume, Version V26: A Conditional Semiclassical Einstein--SM Continuum Theorem}
\addcontentsline{toc}{section}{Second Volume, Version V26: A Conditional Semiclassical Einstein--SM Continuum Theorem}

\subsection*{Purpose}

V25 proves that an unsmeared quantum stress fluctuation cannot be a pathwise
source in the positive Sobolev space used by V14--V15.  The semiclassical
branch avoids that contradiction by sourcing a classical metric with the
renormalized expectation of the stress.  This version formulates that source
by local Hadamard subtraction, classifies the local geometric ambiguity,
states the conservation gate, and proves a self-consistent cutoff-continuum
theorem by contraction.  The theorem is conditional on a uniform
metric-to-source estimate; identifying and proving that estimate for the
interacting Standard Model becomes the precise remaining semiclassical
problem.

\subsection{Hadamard-renormalized expectation}
\label{sec:v2v26-hadamard}

Let \(g\) be a globally hyperbolic metric and let \(W_g(x,x')\) be the
two-point distribution of a Hadamard state for one matter field.  In a
geodesically convex neighbourhood, let \(H_g(x,x')\) denote the local
Hadamard parametrix with the same universal singular coefficients.  If
\(\mathcal D_{\mu\nu}(x,x';g)\) is the point-split bidifferential operator
obtained from the classical stress, define
\begin{equation}
 \langle T_{\mu\nu}\rangle_{\mathrm{ren}}[g,W]
 =
 \left[
 \mathcal D_{\mu\nu}(W_g-H_g)
 \right]_{x'=x}
 +C_{\mu\nu}(g).
\label{eq:v2v26-hadamard-source}
\end{equation}
The local tensor \(C_{\mu\nu}\) fixes the finite renormalization convention.

\begin{lemma}[The coincidence limit removes the free ultraviolet cutoff]
\label{lem:v2v26-coincidence}
If \(W_g-H_g\) is smooth near the diagonal, then the first term in
\eqref{eq:v2v26-hadamard-source} is a smooth tensor.  If regularized pairs
satisfy
\[
 \mathcal D_{\mu\nu}(W_{g,\Lambda}-H_{g,\Lambda})
 \longrightarrow
 \mathcal D_{\mu\nu}(W_g-H_g)
\]
in \(C^r\) on a neighbourhood of the diagonal, then their coincidence
limits converge in \(C^r\).
\end{lemma}

\begin{proof}
A differential operator maps smooth kernels to smooth kernels, and pullback
by the smooth diagonal map \(x\mapsto(x,x)\) preserves smoothness.  The same
two continuous operations preserve \(C^r\) convergence.
\end{proof}

\begin{remark}[Interacting input]
\label{rem:v2v26-interacting}
For interacting SM fields, \(W-H\) is replaced by renormalized time-ordered
products or composite insertions satisfying the microlocal spectrum
condition.  V26 does not infer their existence from the free lemma.  It
uses the resulting expectation as a source map only after its locality,
regularity, and cutoff estimates have been supplied.
\end{remark}

\subsection{Local geometric ambiguity}
\label{sec:v2v26-ambiguity}

On a four-dimensional spacetime without boundary, define the conserved
local tensors
\begin{align}
 I_{\mu\nu}
 &=
 \frac{-2}{\sqrt{|g|}}
 \frac{\delta}{\delta g^{\mu\nu}}
 \int\sqrt{|g|}\,R^2\,\mathrm d^4x,
\label{eq:v2v26-I}\\
 J_{\mu\nu}
 &=
 \frac{-2}{\sqrt{|g|}}
 \frac{\delta}{\delta g^{\mu\nu}}
 \int\sqrt{|g|}\,R_{\alpha\beta}R^{\alpha\beta}\,\mathrm d^4x.
\label{eq:v2v26-J}
\end{align}

\begin{proposition}[Dimension-four geometric ambiguity]
\label{prop:v2v26-ambiguity}
Assume a stress prescription is local, covariant, symmetric, conserved,
state independent in its ambiguity, and obeys four-dimensional engineering
scaling.  Then the difference between two prescriptions has the form
\begin{equation}
 \Delta T_{\mu\nu}
 =
 c_0 g_{\mu\nu}
 +c_1 G_{\mu\nu}
 +c_2 I_{\mu\nu}
 +c_3 J_{\mu\nu},
\label{eq:v2v26-ambiguity}
\end{equation}
with mass parameters absorbed into the constants.  The Euler variation of
the four-dimensional Gauss--Bonnet density adds no independent bulk tensor.
\end{proposition}

\begin{proof}
A state-independent local ambiguity of stress dimension four must be the
metric variation of a local scalar action of dimension at most four.
Modulo a total divergence, the purely geometric scalar basis is
\[
 1,\qquad R,\qquad R^2,\qquad
 R_{\alpha\beta}R^{\alpha\beta},\qquad
 R_{\alpha\beta\gamma\delta}R^{\alpha\beta\gamma\delta}.
\]
Their variations give respectively a multiple of \(g_{\mu\nu}\),
\(G_{\mu\nu}\), \(I_{\mu\nu}\), \(J_{\mu\nu}\), and one further curvature
variation.  In four dimensions the Gauss--Bonnet combination
\[
 R_{\alpha\beta\gamma\delta}^2
 -4R_{\alpha\beta}^2+R^2
\]
has vanishing bulk variation, so the last tensor is a linear combination of
\(I\) and \(J\).  Diffeomorphism invariance of each action makes its metric
variation divergence free.  This proves \eqref{eq:v2v26-ambiguity}.
\end{proof}

\begin{corollary}[Absorption into gravitational couplings]
\label{cor:v2v26-absorb}
The ambiguity \eqref{eq:v2v26-ambiguity} is absorbed by fixing the
renormalized coefficients in
\begin{equation}
 S_{\mathrm{grav}}[g]
 =
 \int\sqrt{|g|}
 \left[
 \frac{R-2\Lambda_{\mathrm R}}{16\pi G_{\mathrm R}}
 +\alpha_{\mathrm R}R^2
 +\beta_{\mathrm R}R_{\mu\nu}R^{\mu\nu}
 \right]\mathrm d^4x.
\label{eq:v2v26-gravity-action}
\end{equation}
After those physical renormalization conditions are fixed, changing the
matter regulator does not create an additional arbitrary source.
\end{corollary}

\begin{proof}
Metric variation of \eqref{eq:v2v26-gravity-action} produces the four tensor
directions in \eqref{eq:v2v26-ambiguity}.  Redefine the four coefficients.
\end{proof}

\subsection{The conservation gate}
\label{sec:v2v26-conservation}

\begin{theorem}[Conservation is necessary for harmonic closure]
\label{thm:v2v26-conservation}
Consider a second-order Einstein equation
\begin{equation}
 G_{\mu\nu}(g)+\Lambda_{\mathrm R}g_{\mu\nu}
 =
 8\pi G_{\mathrm R}\,\mathscr T_{\mu\nu}(g).
\label{eq:v2v26-semiclassical}
\end{equation}
A necessary compatibility condition is
\begin{equation}
 \nabla_g^\mu\mathscr T_{\mu\nu}(g)=0.
\label{eq:v2v26-conservation}
\end{equation}
If \eqref{eq:v2v26-conservation} holds, the initial Einstein constraints
hold, and the reduced harmonic solution is regular enough, then the harmonic
gauge constraint obeys a homogeneous linear wave system and remains zero.
\end{theorem}

\begin{proof}
Taking a covariant divergence of \eqref{eq:v2v26-semiclassical} and using the
contracted Bianchi identity gives the necessity of
\eqref{eq:v2v26-conservation}.  In the harmonic reduction, the difference
between the reduced and geometric Einstein tensors is linear in the
symmetrized covariant derivative of the gauge vector, plus lower terms.
Taking its divergence, using \eqref{eq:v2v26-conservation}, and commuting
covariant derivatives gives a linear wave equation for that vector with
coefficients determined by \(g\).  Its Cauchy data vanish when the harmonic
and Einstein constraints hold initially.  Uniqueness for the linear wave
system makes the gauge vector vanish throughout the slab.
\end{proof}

\begin{assessment}[An anomaly is a structural obstruction]
\label{ass:v2v26-anomaly}
A regulator-level Ward defect may be repaired only if its local
cohomology class is trivial, as in V3 and V6.  If the complete chiral gauge
content had a nontrivial diffeomorphism or gauge anomaly, no choice of
\(C_{\mu\nu}\) could supply \eqref{eq:v2v26-conservation}; the
self-consistent Einstein equation would fail before any Sobolev estimate.
Thus anomaly cancellation and conservation belong among the hypotheses of
the continuum theorem, not among its conclusions.
\end{assessment}

\subsection{Banach spaces and the self-consistent map}
\label{sec:v2v26-fixed-point}

On a common time slab \(I=[0,T]\), put
\begin{align}
 X_T
 &=
 C(I;H^s_x)\cap C^1(I;H^{s-1}_x),
\label{eq:v2v26-X}\\
 Y_T
 &=
 L^1(I;H^{s-1}_x).
\label{eq:v2v26-Y}
\end{align}
Tensor indices and a fixed background atlas are included in these norms.
Let \(\mathcal B_R\subset X_T\) be a closed metric ball satisfying a
uniform hyperbolicity margin.

Let
\[
 \mathcal E_T:\mathcal S_T\subset Y_T\longrightarrow\mathcal B_R
\]
denote the harmonic Einstein solution operator with fixed compatible
initial data.  V15 supplies the relevant local stability pattern.  Assume
\begin{equation}
 \|\mathcal E_T(S)-\mathcal E_T(\widetilde S)\|_{X_T}
 \le C_E\|S-\widetilde S\|_{Y_T}.
\label{eq:v2v26-E-Lipschitz}
\end{equation}
Let the renormalized expectation and state evolution define a source map
\[
 \mathscr T:\mathcal B_R\longrightarrow\mathcal S_T
\]
such that
\begin{align}
 \nabla_g^\mu\mathscr T_{\mu\nu}(g)&=0,
\label{eq:v2v26-map-conservation}\\
 \|\mathscr T(g)-\mathscr T(\widetilde g)\|_{Y_T}
 &\le L_T\|g-\widetilde g\|_{X_T}.
\label{eq:v2v26-map-Lipschitz}
\end{align}

\begin{theorem}[Conditional self-consistent semiclassical solution]
\label{thm:v2v26-fixed-point}
Assume
\[
 \mathcal E_T(\mathscr T(\mathcal B_R))\subseteq\mathcal B_R,
 \qquad
 q:=C_EL_T<1.
\label{eq:v2v26-contraction}
\]
Then \eqref{eq:v2v26-semiclassical} with the prescribed initial data has a
unique self-consistent harmonic solution \(g\in\mathcal B_R\).  It satisfies
the unreduced Einstein equation and constraints.
\end{theorem}

\begin{proof}
Define
\[
 \Phi(g)=\mathcal E_T(\mathscr T(g)).
\]
By \eqref{eq:v2v26-E-Lipschitz} and
\eqref{eq:v2v26-map-Lipschitz},
\[
 \|\Phi(g)-\Phi(\widetilde g)\|_{X_T}
 \le q\|g-\widetilde g\|_{X_T}.
\]
The ball is complete and invariant, so Banach's fixed-point theorem gives a
unique \(g=\Phi(g)\).  Conservation
\eqref{eq:v2v26-map-conservation}, initial compatibility, and
\cref{thm:v2v26-conservation} restore the unreduced equations and propagate
the constraints.
\end{proof}

\subsection{Cutoff convergence of self-consistent metrics}
\label{sec:v2v26-cutoff}

Let \(\mathscr T_\Lambda\) be regulator-dependent source maps with the same
renormalized gravitational coefficients and initial state prescription.

\begin{theorem}[Uniform source-map convergence implies metric convergence]
\label{thm:v2v26-cutoff}
Suppose every \(\mathscr T_\Lambda\) maps \(\mathcal B_R\) into
\(\mathcal S_T\), obeys the same Lipschitz constant \(L_T\), is conserved
with respect to its metric argument, and
\begin{equation}
 \epsilon_\Lambda
 =
 \sup_{g\in\mathcal B_R}
 \|\mathscr T_\Lambda(g)-\mathscr T(g)\|_{Y_T}
 \longrightarrow0.
\label{eq:v2v26-uniform-source}
\end{equation}
If \(C_EL_T=q<1\) and all fixed-point maps preserve \(\mathcal B_R\), their
unique solutions \(g_\Lambda\) satisfy
\begin{equation}
 \|g_\Lambda-g\|_{X_T}
 \le
 \frac{C_E}{1-q}\epsilon_\Lambda.
\label{eq:v2v26-metric-rate}
\end{equation}
\end{theorem}

\begin{proof}
Use the two fixed-point equations:
\begin{align*}
 \|g_\Lambda-g\|_{X_T}
 &\le
 C_E
 \|\mathscr T_\Lambda(g_\Lambda)-\mathscr T(g)\|_{Y_T}\\
 &\le
 C_E\epsilon_\Lambda
 +C_EL_T\|g_\Lambda-g\|_{X_T}.
\end{align*}
Move the last term to the left and divide by \(1-q\).
\end{proof}

\begin{corollary}[No separate cutoff subsequence is needed]
\label{cor:v2v26-no-subsequence}
Under \cref{thm:v2v26-cutoff}, the entire cutoff family converges in
\(X_T\), and the limit is independent of the regulator within the fixed
renormalization convention.
\end{corollary}

\begin{proof}
The right side of \eqref{eq:v2v26-metric-rate} tends to zero and the fixed
point is unique.
\end{proof}

\subsection{The higher-curvature derivative gate}
\label{sec:v2v26-higher-curvature}

\begin{firewall}[Curvature-squared terms are not covered by V15]
\label{fw:v2v26-fourth-order}
The tensors \(I_{\mu\nu}\) and \(J_{\mu\nu}\) contain fourth derivatives of
the metric.  If nonzero \(\alpha_{\mathrm R},\beta_{\mathrm R}\) in
\eqref{eq:v2v26-gravity-action} are retained as principal dynamical terms,
the equation is fourth order and the second-order solution operator
\(\mathcal E_T\) used above is not the correct operator.  Moving those
tensors to the right side does not solve the problem: the map
\(X_T\to Y_T\) then loses derivatives and generally violates
\eqref{eq:v2v26-map-Lipschitz}.
\end{firewall}

There are three honest options:
\begin{enumerate}
 \item impose renormalization conditions with vanishing finite
 curvature-squared couplings and prove that the remaining expectation map
 closes at second order;
 \item formulate and prove a well-posed fourth-order semiclassical gravity
 system with the additional initial data and stability estimates;
 \item use a controlled order-reduction expansion and prove an error bound
 relative to the unreduced effective equation.
\end{enumerate}
These options define different continuum statements and may not be exchanged
without proof.

\subsection{The exact remaining semiclassical gate}
\label{sec:v2v26-gate}

\begin{programme}[Metric-to-source acquisition]
\label{prog:v2v26}
For the full SM expectation, establish on one hyperbolic ball
\(\mathcal B_R\):
\begin{enumerate}
 \item a local covariant renormalized composite stress with all gauge,
 ghost, Higgs, fermion, measure, and contact terms included;
 \item exact conservation after anomaly cancellation and finite local
 repair;
 \item the uniform bound
 \(\mathscr T_\Lambda(\mathcal B_R)\subset\mathcal S_T\);
 \item the common Lipschitz estimate
 \eqref{eq:v2v26-map-Lipschitz} without derivative loss;
 \item uniform regulator convergence
 \eqref{eq:v2v26-uniform-source};
 \item a slab small enough that \(C_EL_T<1\), followed by continuation on
 overlapping slabs if the hyperbolicity and source bounds persist.
\end{enumerate}
Once these six statements and one consistent higher-curvature choice are
proved, \cref{thm:v2v26-cutoff} supplies the semiclassical continuum metric.
\end{programme}

\begin{assessment}[What has and has not been reduced]
\label{ass:v2v26-status}
The semiclassical metric problem is no longer blocked by the V25
equal-time fluctuation variance because only the expectation enters.
It is also not closed by pointwise convergence of
\(\langle T_{\mu\nu}\rangle\) on a fixed metric.  Self-consistency requires
uniform convergence and a common Lipschitz estimate over a metric ball.
The latter is now the central analytic gate.
\end{assessment}

\subsection{V26 conclusion and next step}
\label{sec:v2v26-conclusion}

Local Hadamard subtraction removes the universal free singularity of the
mean stress.  Its finite geometric ambiguity is exactly the variation of
cosmological, Einstein, and curvature-squared couplings, and conservation is
required for harmonic constraint propagation.  Under a uniform
metric-to-source Lipschitz estimate with contraction constant below one,
the semiclassical Einstein--SM equation has a unique self-consistent
solution; uniform cutoff convergence of the source maps gives the explicit
metric rate \eqref{eq:v2v26-metric-rate}.

The theorem is conditional because the interacting SM source map has not
been constructed with the stated uniform estimates, and the
curvature-squared choice affects the principal PDE.  V27 will attack the
free-field model of the missing Lipschitz bound by differentiating the
Hadamard-renormalized expectation with respect to the metric and deriving
its causal linear-response kernel.
\clearpage
\section*{Second Volume, Version V27: Causal Stress Response and the Fourth-Order Logarithmic Gate}
\addcontentsline{toc}{section}{Second Volume, Version V27: Causal Stress Response and the Fourth-Order Logarithmic Gate}

\subsection*{Purpose}

V26 reduced the semiclassical continuum to a uniform metric-to-source
Lipschitz estimate.  This version tests that estimate in the free scalar
vacuum.  Differentiating the renormalized mean stress gives a local contact
term plus a retarded stress--stress commutator.  The physical TT spectral
density is a positive two-particle phase-space integral and grows like
\(s^2\).  Consequently the renormalized response contains the nonlocal
symbol \(k^4\log(-k^2)\).  No choice of local dimension-four counterterms
removes that logarithm, and the source derivative is not a same-order map
from the V26 metric space to its source space.

\subsection{The causal Kubo derivative}
\label{sec:v2v27-kubo}

Let \(g_\varepsilon=g+\varepsilon h\), and let the matter state be transported
from fixed initial data by the perturbed dynamics.  The first metric
variation of the action is
\begin{equation}
 \delta S_{\mathrm m}
 =
 \frac12\int
 T^{\alpha\beta}(y)h_{\alpha\beta}(y)
 \,\mathrm d\operatorname{vol}_g(y).
\label{eq:v2v27-action-variation}
\end{equation}

\begin{theorem}[Retarded metric derivative of the mean stress]
\label{thm:v2v27-kubo}
For a free Hadamard field, the first variation of the renormalized
expectation has the form
\begin{align}
 D\mathscr T_{\mu\nu}[g](h)(x)
 &=
 L_{\mu\nu}^{\ \ \alpha\beta}[g]h_{\alpha\beta}(x)
 \nonumber\\
 &\quad
 -\frac{i}{2}\int_{J^-(x)}
 \left\langle
 [T_{\mu\nu}(x),T^{\alpha\beta}(y)]
 \right\rangle_{\mathrm{ren}}
 h_{\alpha\beta}(y)
 \,\mathrm d\operatorname{vol}_g(y).
\label{eq:v2v27-kubo}
\end{align}
Here \(L[g]\) is a local differential contact operator containing the
explicit metric variation of \(T\), the parametrix variation, and the fixed
finite counterterms.  The integral kernel is retarded.
\end{theorem}

\begin{proof}
In the interaction picture, the perturbation Hamiltonian induced by
\eqref{eq:v2v27-action-variation} is the spatial integral of
\(-T^{\alpha\beta}h_{\alpha\beta}/2\), up to the convention for varying
\(g_{\alpha\beta}\) rather than \(g^{\alpha\beta}\).  Differentiating the
unitary propagator by Duhamel's formula gives
\[
 \delta\langle T_{\mu\nu}(x)\rangle
 =
 -\frac i2
 \int_{J^-(x)}
 \langle[T_{\mu\nu}(x),T^{\alpha\beta}(y)]\rangle
 h_{\alpha\beta}(y)\,\mathrm d\operatorname{vol}_g(y),
\]
because only perturbations in the causal past affect the observable at
\(x\).  Differentiating the observable, the volume density, and the
Hadamard subtraction produces distributions supported on the diagonal.
Local covariance combines them with the finite renormalization variations
into \(L[g]h\).  This proves \eqref{eq:v2v27-kubo}.
\end{proof}

\begin{corollary}[Linearized conservation Ward identity]
\label{cor:v2v27-Ward}
If the renormalized prescription is conserved for every metric, then
\begin{equation}
 \nabla^\mu D\mathscr T_{\mu\nu}[g](h)
 +
 (D\nabla^\mu)_g(h)\,
 \mathscr T_{\mu\nu}(g)
 =0.
\label{eq:v2v27-linear-Ward}
\end{equation}
On Minkowski vacuum, where the renormalized mean is chosen proportional to
the metric and absorbed into the cosmological term, the nonlocal TT
polarization is transverse.
\end{corollary}

\begin{proof}
Differentiate
\(\nabla_g^\mu\mathscr T_{\mu\nu}(g)=0\).
\end{proof}

\subsection{Positive two-particle TT spectral density}
\label{sec:v2v27-spectral}

Work in Minkowski vacuum and contract both stress indices with a normalized
spin-two polarization.  Let \(K\) be a future timelike four-momentum with
\[
 K^2=s>4m^2.
\]
In its rest frame, two on-shell scalar momenta have
\[
 p=(\sqrt{s}/2,\mathbf p),
 \qquad
 q=(\sqrt{s}/2,-\mathbf p),
 \qquad
 |\mathbf p|=\frac12\sqrt{s-4m^2}.
\]

\begin{theorem}[Exact spectral power in the scalar TT channel]
\label{thm:v2v27-spectral}
Let \(\rho_2(s)\) be the scalar coefficient of the spin-two part of the
stress commutator spectral density.  Then
\begin{equation}
 \rho_2(s)
 =
 c_{\mathrm{TT}}\,
 s^2\left(1-\frac{4m^2}{s}\right)^{5/2}
 \mathbf1_{\{s>4m^2\}},
 \qquad
 c_{\mathrm{TT}}>0,
\label{eq:v2v27-spectral}
\end{equation}
up to the fixed normalization convention for \(T\) and one-particle states.
In particular,
\begin{equation}
 c s^2\le\rho_2(s)\le Cs^2
 \quad\text{for }s\ge8m^2.
\label{eq:v2v27-spectral-growth}
\end{equation}
\end{theorem}

\begin{proof}
The spectral resolution inserts a complete two-particle state between the
two stresses.  Its TT part is the positive phase-space integral
\begin{equation}
 \rho_2(s)
 =
 C\int d\Phi_2(K;p,q)
 \left|
 e^{ij}
 \langle p,q|T_{ij}(0)|0\rangle
 \right|^2.
\label{eq:v2v27-phase-space}
\end{equation}
Trace, mass, potential, and improvement structures vanish under \(e\).
The remaining scalar derivative vertex in the centre-of-mass frame is
\[
 e^{ij}(p_iq_j+q_ip_j)
 =
 -2e^{ij}p_ip_j.
\]
The angular identity
\[
 \int_{\mathbb S^2}|e^{ij}n_in_j|^2\,\mathrm d\Omega
 =
 c_e>0
\]
follows from rotational invariance and the nonzero traceless norm of \(e\).
Thus the angularly integrated squared matrix element is
\(C|\mathbf p|^4\).  Integrated two-body phase space contributes
\[
 C\sqrt{1-\frac{4m^2}{s}}.
\]
Since
\[
 |\mathbf p|^4
 =
 \frac{s^2}{16}
 \left(1-\frac{4m^2}{s}\right)^2,
\]
multiplication proves \eqref{eq:v2v27-spectral}.  Positivity and
\eqref{eq:v2v27-spectral-growth} follow.
\end{proof}

\begin{assessment}[The complete positive carrier is visible]
\label{ass:v2v27-positive}
Equation \eqref{eq:v2v27-phase-space} is the stress analogue of the complete
positive YM carrier audited in V25: all two-particle directions are
assembled before the norm is taken.  The result is strictly positive in the
physical spin-two sector.  Local mean counterterms cannot cancel this
spectral weight.
\end{assessment}

\subsection{Subtracted dispersion and the logarithm}
\label{sec:v2v27-dispersion}

Let \(\pi_2(z)\) be the scalar spin-two form factor of the retarded
polarization, with \(z=K^2\).  Local renormalization changes \(\pi_2\) by a
polynomial of degree at most two in \(z=K^2\), hence at most four in
momentum, equivalently by the
cosmological, Einstein, and curvature-squared directions of V26.

\begin{theorem}[Unavoidable fourth-order logarithmic response]
\label{thm:v2v27-log}
After any fixed local subtraction compatible with
\cref{prop:v2v26-ambiguity}, the Euclidean spacelike asymptotic satisfies
\begin{equation}
 \pi_2(-Q^2)
 =
 c_*Q^4\log(Q^2/\mu_{\mathrm R}^2)
 +O(Q^4),
 \qquad
 c_*>0,
 \qquad
 Q\longrightarrow\infty.
\label{eq:v2v27-log}
\end{equation}
A local counterterm may change the \(O(Q^4)\) polynomial but cannot remove
the logarithm at every \(Q\).
\end{theorem}

\begin{proof}
A sufficiently subtracted dispersion relation has discontinuity
\(\rho_2(s)\).  Differentiate it three times with respect to \(z\); all
local polynomials of degree at most two disappear, and for \(z=-Q^2\),
\begin{equation}
 \pi_2^{(3)}(-Q^2)
 =
 C\int_{4m^2}^{\infty}
 \frac{\rho_2(s)}{(s+Q^2)^4}\,\mathrm ds.
\label{eq:v2v27-third-derivative}
\end{equation}
The integrand is positive.  Substitute \(s=Q^2u\) and use
\(\rho_2(Q^2u)\sim cQ^4u^2\):
\[
 \pi_2^{(3)}(-Q^2)
 =
 \frac{c+o(1)}{Q^2}
 \int_0^\infty\frac{u^2}{(1+u)^4}\,\mathrm du.
\]
Three integrations in \(z\) give a nonzero
\(z^2\log(-z/\mu_{\mathrm R}^2)\) term, plus a polynomial of degree at most
two.  Replacing \(z\) by \(-Q^2\) proves
\eqref{eq:v2v27-log}.  A local counterterm changes only the integration
polynomial.
\end{proof}

\begin{remark}[Causality fixes the boundary value]
\label{rem:v2v27-causal-log}
For timelike momentum, the retarded form factor is the causal boundary value
of the logarithm and has the imaginary part prescribed by
\eqref{eq:v2v27-spectral}.  The imaginary part is positive spectral data,
not a finite-renormalization convention.
\end{remark}

\subsection{Failure of the V26 same-order Lipschitz map}
\label{sec:v2v27-derivative-loss}

\begin{theorem}[The free source derivative is not \(X_T\to Y_T\) bounded]
\label{thm:v2v27-unbounded}
On Minkowski vacuum, the linearized renormalized source map cannot obey
\begin{equation}
 \|D\mathscr T[\eta](h)\|_{L^1(I;H^{s-1}_x)}
 \le C\|h\|_{C(I;H^s_x)\cap C^1(I;H^{s-1}_x)}
\label{eq:v2v27-false-bound}
\end{equation}
with a frequency-independent \(C\), even after arbitrary fixed local
dimension-four counterterms.
\end{theorem}

\begin{proof}
Choose a fixed temporal frequency \(\tau_0\) and smooth spatial TT wave
packets whose Fourier supports are unit balls about \(\xi_N\), with
\(|\xi_N|=N\to\infty\).  Give each packet amplitude \(a_N\), multiply
by \(e^{-i\tau_0t}\), and obtain it as the limit of an adiabatically
switched perturbation.  On the fixed interval \(I\), its input norm is
comparable to \(|a_N|N^s\).  Throughout the Fourier support,
\(K^2=\tau_0^2-|\xi|^2=-N^2+O(N)\).  Translation invariance and
\cref{thm:v2v27-log} make the TT response multiplier uniformly comparable
to \(N^4\log N\) on a smaller packet support.  By \cref{thm:v2v27-log},
\[
 |\pi_2(K_N^2)|
 \ge cN^4\log N
\]
for large \(N\), after any polynomial counterterm has been fixed.
Therefore
\[
 \|D\mathscr T(h_N)\|_{L^1(I;H^{s-1})}
 \ge
 c_I|a_N|N^{s+3}\log N.
\]
The ratio to the input norm grows like \(N^3\log N\), contradicting
\eqref{eq:v2v27-false-bound}.
\end{proof}

\begin{corollary}[The V26 contraction constant is infinite in the free
linearization]
\label{cor:v2v27-no-contraction}
The Lipschitz hypothesis \eqref{eq:v2v26-map-Lipschitz} does not hold for the
unreduced free semiclassical source on the second-order metric space
\eqref{eq:v2v26-X}.  Shrinking the time interval does not remove the
unbounded spatial factor \(N^3\log N\).
\end{corollary}

\begin{proof}
If the map were locally Lipschitz from \(X_T\) to \(Y_T\), every smooth
directional difference quotient would be bounded by the same local
Lipschitz constant times the direction norm.  The Kubo directional
derivative along the packets of \cref{thm:v2v27-unbounded} violates that
bound.  The factor \(|I|\) is independent of \(N\), so any fixed positive
slab retains the contradiction.
\end{proof}

\begin{firewall}[The logarithm is not the local curvature ambiguity]
\label{fw:v2v27-log}
The \(Q^4\) polynomial accompanying
\eqref{eq:v2v27-log} can be moved between
\(\langle T_{\mu\nu}\rangle_{\mathrm{ren}}\) and the
curvature-squared couplings.  The \(Q^4\log Q^2\) term is nonlocal spectral
response and cannot be removed by choosing finite constants
\(\alpha_{\mathrm R},\beta_{\mathrm R}\).  Setting the local fourth-order
couplings to zero therefore does not recover the V26 second-order
Lipschitz map.
\end{firewall}

\subsection{The correct linearized semiclassical operator}
\label{sec:v2v27-correct-operator}

In the TT sector, the flat linearized self-consistent equation has the
schematic Fourier form
\begin{equation}
 \left[
 \frac{1}{16\pi G_{\mathrm R}}K^2
 +\alpha_{\mathrm R}K^4
 +\pi_2^{\mathrm{nonloc}}(K^2)
 \right]\widehat h^{\mathrm{TT}}(K)
 =
 \widehat S^{\mathrm{TT}}_{\mathrm{state}}(K),
\label{eq:v2v27-effective-operator}
\end{equation}
where
\(\pi_2^{\mathrm{nonloc}}(K^2)\sim
 c(K^2)^2\log(-K^2/\mu_{\mathrm R}^2)\).

\begin{assessment}[Three replacement routes]
\label{ass:v2v27-routes}
The V26 fixed-point theorem remains valid as an abstract theorem, but its
same-order source-map hypothesis fails in the free model.  A semiclassical
construction must instead choose one of the following and prove it:
\begin{enumerate}
 \item incorporate the complete causal polarization
 \(\pi_2^{\mathrm{nonloc}}\) into the principal linear solution operator and
 control its zeros and retarded inverse;
 \item perform order reduction below a stated physical scale and quantify
 the discarded \(G_{\mathrm R}K^2\log K^2\) terms;
 \item work in a scale of spaces where the metric has enough additional
 derivatives to pay the response, while closing the nonlinear Einstein
 map without circular derivative loss.
\end{enumerate}
The first route is the only one aimed directly at an all-frequency
semiclassical continuum.
\end{assessment}

\begin{firewall}[No stability follows from ultraviolet ellipticity alone]
\label{fw:v2v27-zeros}
The magnitude \(K^4\log K^2\) at spacelike frequency does not by itself give
a causal stable inverse.  Timelike continuation has a branch cut fixed by
\eqref{eq:v2v27-spectral}, and the full denominator in
\eqref{eq:v2v27-effective-operator} may have additional zeros.  A continuum
claim requires a retarded spectral construction and an exclusion or
prescription for unstable poles.
\end{firewall}

\subsection{V27 conclusion and next step}
\label{sec:v2v27-conclusion}

The metric derivative of the free renormalized mean stress is a causal Kubo
kernel plus local contacts.  Its complete TT spectral density is positive
and behaves as \(s^2\).  A subtracted dispersion relation therefore produces
the unavoidable response
\(K^4\log(-K^2)\).  This yields an explicit high-frequency sequence for
which the V26 same-order metric-to-source Lipschitz ratio grows as
\(N^3\log N\).  The second-order contraction route fails already in the
free vacuum.

V28 will pursue the all-frequency branch: define the retarded TT inverse of
\eqref{eq:v2v27-effective-operator} by its spectral representation, locate
its real and complex zeros under explicit signs of the renormalized
couplings, and separate the stable causal continuum from Landau-pole or
runaway choices.
\clearpage
\section*{Second Volume, Version V28: Exact Logarithmic Denominator Zeros and the Retarded-Inverse Gate}
\addcontentsline{toc}{section}{Second Volume, Version V28: Exact Logarithmic Denominator Zeros and the Retarded-Inverse Gate}

\subsection*{Purpose}

V27 shows that the free semiclassical TT operator contains a fourth-order
logarithmic polarization.  Before this term can be placed in the principal
propagator, its additional zeros must be understood.  This version solves
the Euclidean logarithmic model exactly.  A single dimensionless inequality
decides whether positive Euclidean zeros occur; when they do, their
Lambert-\(W\) formula and opposite residue signs are explicit.  The
criterion is invariant under a change of subtraction scale.  The final
section separates this exact Euclidean result from the still-open retarded
spectral inverse.

\subsection{The Euclidean TT denominator}
\label{sec:v2v28-denominator}

Let \(x=Q^2>0\), \(M^2>0\), \(c>0\), and \(a\in\mathbb R\).  Consider
\begin{equation}
 \mathfrak D_E(x)
 =
 x\,F(x),
 \qquad
 F(x)
 =
 M^2+x\left[a+c\log(x/\mu^2)\right].
\label{eq:v2v28-denominator}
\end{equation}
Here \(M^2\) is the Einstein coefficient, \(a\) is the finite local
curvature-squared coefficient in the selected TT normalization, and \(c\)
is the positive logarithmic coefficient supplied by
\cref{thm:v2v27-spectral,thm:v2v27-log}.  This is an exact model for the
Einstein term plus the leading renormalized ultraviolet logarithm; it is not
yet the complete massive polarization.

\begin{lemma}[Unique minimum of the reduced denominator]
\label{lem:v2v28-minimum}
The function \(F\) has a unique critical point
\begin{equation}
 x_*=\mu^2\exp(-1-a/c),
\label{eq:v2v28-xstar}
\end{equation}
which is its strict global minimum, and
\begin{equation}
 F(x_*)=M^2-c\mu^2\exp(-1-a/c).
\label{eq:v2v28-minimum}
\end{equation}
\end{lemma}

\begin{proof}
Differentiate:
\[
 F'(x)=a+c\log(x/\mu^2)+c,
 \qquad
 F''(x)=c/x>0.
\]
Thus the sole critical point is \eqref{eq:v2v28-xstar} and is a strict
global minimum.  At that point
\(a+c\log(x_*/\mu^2)=-c\), giving
\eqref{eq:v2v28-minimum}.  Also
\(F(0+)=M^2\) and \(F(x)\to+\infty\), so no other minimum exists.
\end{proof}

\begin{theorem}[Exact positive-zero trichotomy]
\label{thm:v2v28-zero-trichotomy}
Put
\begin{equation}
 \Delta
 =
 \frac{M^2}{c\mu^2}\exp(a/c+1).
\label{eq:v2v28-Delta}
\end{equation}
Then:
\begin{enumerate}
 \item if \(\Delta>1\), \(F(x)>0\) for every \(x>0\), so
 \(\mathfrak D_E\) has no positive zero apart from the massless factor at
 \(x=0\);
 \item if \(\Delta=1\), \(F\) has one double zero at \(x=x_*\);
 \item if \(0<\Delta<1\), \(F\) has exactly two simple positive zeros
 \(x_-<x_*<x_+\).
\end{enumerate}
\end{theorem}

\begin{proof}
By \cref{lem:v2v28-minimum}, the sign of the unique minimum is the sign of
\[
 M^2-c\mu^2e^{-1-a/c}
 =
 c\mu^2e^{-1-a/c}(\Delta-1).
\]
Combine this with the positive endpoint limits and strict convexity.
\end{proof}

\subsection{Closed form and residues}
\label{sec:v2v28-Lambert}

Let \(W_0\) and \(W_{-1}\) be the two real Lambert branches on
\([-e^{-1},0)\).

\begin{theorem}[Lambert-\(W\) formula]
\label{thm:v2v28-Lambert}
When \(0<\Delta<1\), define
\[
 B=\frac{M^2}{c\mu^2},
 \qquad
 z=-Be^{a/c}\in(-e^{-1},0).
\]
The two positive roots are
\begin{equation}
 x_-=-\frac{M^2/c}{W_{-1}(z)},
 \qquad
 x_+=-\frac{M^2/c}{W_0(z)}.
\label{eq:v2v28-roots}
\end{equation}
At a simple root \(x_r\),
\begin{equation}
 \mathfrak D_E'(x_r)
 =
 x_r F'(x_r).
\label{eq:v2v28-residue}
\end{equation}
Consequently
\[
 \mathfrak D_E'(x_-)<0,
 \qquad
 \mathfrak D_E'(x_+)>0,
\]
so the two simple-pole residues of
\(\mathfrak D_E^{-1}\) have opposite signs.
\end{theorem}

\begin{proof}
Set \(y=x/\mu^2\) and \(u=a/c+\log y\).  The equation \(F(x)=0\) becomes
\[
 u e^u=-Be^{a/c}=z.
\]
Thus \(u=W_k(z)\), and
\[
 x=\mu^2e^{u-a/c}
 =-\frac{M^2/c}{W_k(z)}.
\]
The branch \(W_{-1}\) gives the smaller root and \(W_0\) the larger.
At a root, differentiating
\(\mathfrak D_E=xF\) gives
\eqref{eq:v2v28-residue}.  Since \(F\) decreases before its unique minimum
and increases after it, the two derivative signs follow.
\end{proof}

\begin{assessment}[Meaning of the opposite residues]
\label{ass:v2v28-residues}
If both Euclidean roots continued to isolated simple propagator poles, their
opposite residues would prevent both modes from having the same positive
quadratic sign.  This is the elementary higher-derivative ghost warning.
The statement is conditional on pole continuation: the complete retarded
polarization has a branch cut and cannot be replaced globally by the real
Euclidean logarithm.
\end{assessment}

\subsection{Subtraction-scale invariance}
\label{sec:v2v28-RG}

\begin{theorem}[The zero criterion is not removed by changing \(\mu\)]
\label{thm:v2v28-scale}
The function \eqref{eq:v2v28-denominator} is unchanged under
\begin{equation}
 \mu\mapsto\mu',
 \qquad
 a\mapsto a'
 =
 a+c\log(\mu'^2/\mu^2).
\label{eq:v2v28-running}
\end{equation}
Moreover
\begin{equation}
 \mu'^2e^{-a'/c}
 =
 \mu^2e^{-a/c},
 \qquad
 \Delta'=\Delta.
\label{eq:v2v28-invariant}
\end{equation}
Thus the trichotomy in \cref{thm:v2v28-zero-trichotomy} is invariant under a
mere change of subtraction scale.
\end{theorem}

\begin{proof}
Substitute \eqref{eq:v2v28-running}:
\[
 a'+c\log(x/\mu'^2)
 =
 a+c\log(x/\mu^2).
\]
The two identities in \eqref{eq:v2v28-invariant} follow directly.
\end{proof}

\begin{firewall}[Scale choice is not a physical coupling change]
\label{fw:v2v28-scale}
Holding \(a\) fixed while changing \(\mu\) changes the denominator and hence
the physical renormalization condition.  Running \(a\) by
\eqref{eq:v2v28-running} leaves the operator and its zeros unchanged.
Therefore an unwanted zero cannot be removed by relabelling the subtraction
scale.
\end{firewall}

\subsection{A coercive Euclidean regime}
\label{sec:v2v28-coercive}

\begin{corollary}[Quantitative positivity above threshold]
\label{cor:v2v28-coercive}
If \(\Delta>1\), then
\begin{equation}
 \mathfrak D_E(x)
 \ge
 \gamma x,
 \qquad
 \gamma
 =
 M^2-c\mu^2e^{-1-a/c}>0.
\label{eq:v2v28-coercive}
\end{equation}
Hence for every Euclidean source \(S\) for which the expressions are finite,
the solution
\(\widehat h(Q)=\widehat S(Q)/\mathfrak D_E(Q^2)\) obeys
\begin{equation}
 \|h\|_{\dot H^{r+2}}
 \le
 \gamma^{-1}\|S\|_{\dot H^r}.
\label{eq:v2v28-two-derivative}
\end{equation}
At sufficiently large \(Q\), one additionally has
\begin{equation}
 |\widehat h(Q)|
 \le
 \frac{C|\widehat S(Q)|}
 {Q^4\log(Q^2/\mu^2)}.
\label{eq:v2v28-four-log}
\end{equation}
\end{corollary}

\begin{proof}
The minimum formula gives \(F(x)\ge\gamma\), proving
\eqref{eq:v2v28-coercive}.  Multiply the Fourier relation by
\(|Q|^{r+2}\), use
\(\mathfrak D_E(Q^2)\ge\gamma Q^2\), and apply Plancherel to obtain
\eqref{eq:v2v28-two-derivative}.  The logarithmic term dominates the fixed
lower-order coefficients for large \(Q\), proving
\eqref{eq:v2v28-four-log}.
\end{proof}

\begin{assessment}[Euclidean inversion repairs the derivative count]
\label{ass:v2v28-derivatives}
In the zero-free regime, the complete principal operator has at least the
ordinary Einstein two-derivative inverse and gains nearly four derivatives
in the ultraviolet.  This is why the V27 polarization must be placed in the
left-hand operator rather than estimated as a source.  Euclidean coercivity,
however, is not a hyperbolic stability theorem.
\end{assessment}

\subsection{The retarded gate}
\label{sec:v2v28-retarded}

For timelike \(z>4m^2\), the causal boundary value has
\begin{equation}
 \operatorname{Im}\pi_2^{\mathrm{ret}}(z+i0)
 =
 -c_{\mathrm{conv}}\rho_2(z),
\label{eq:v2v28-imaginary}
\end{equation}
where the sign depends only on the retarded Fourier convention and
\(\rho_2(z)>0\) is fixed by
\eqref{eq:v2v27-spectral}.  The full denominator is analytic off its
two-particle cut, obeys the corresponding dispersion relation, and contains
the exact massive threshold rather than a pure logarithm.

\begin{firewall}[What V28 does not prove]
\label{fw:v2v28-retarded}
The inequality \(\Delta>1\) excludes positive zeros of the Euclidean
logarithmic model.  It does not exclude:
\begin{enumerate}
 \item complex zeros of the analytically continued denominator;
 \item zeros below the two-particle threshold;
 \item zeros introduced by the exact finite massive part;
 \item scalar or longitudinal gravitational poles absent from the TT model;
 \item nonlinear runaway solutions.
\end{enumerate}
A causal continuum theorem must analyze the complete subtracted spectral
function, specify its retarded contour, and prove bounds for that inverse.
\end{firewall}

\begin{programme}[Retarded inverse acquisition]
\label{prog:v2v28}
For each physical spin sector:
\begin{enumerate}
 \item retain the positive spectral density before tensor components are
 estimated separately;
 \item form the complete subtracted analytic denominator with the fixed
 renormalized gravitational coefficients;
 \item locate all first-sheet zeros and identify any second-sheet
 resonances;
 \item define the retarded inverse by a contour consistent with the
 spectral cut and initial data;
 \item prove a frequency-localized energy or spacetime estimate for that
 inverse;
 \item insert the nonlinear remainder only after this principal estimate is
 available.
\end{enumerate}
\end{programme}

\subsection{V28 conclusion and next step}
\label{sec:v2v28-conclusion}

For the Euclidean Einstein-plus-logarithm TT model, the dimensionless
quantity \(\Delta\) in \eqref{eq:v2v28-Delta} gives the complete positive
zero classification.  When two zeros occur, their exact Lambert-\(W\)
locations have opposite residue signs.  The classification is invariant
under consistent running of the finite curvature-squared coefficient.  In
the zero-free regime the Euclidean inverse recovers two derivatives
globally and nearly four at high frequency.

These results justify moving the polarization into the principal operator,
but they do not construct its causal inverse.  V29 will use the complete
positive spectral density to derive a subtracted Stieltjes representation
for the Euclidean form factor and test whether a sign condition stronger
than \(\Delta>1\) controls first-sheet zeros of the full massive
denominator.
\clearpage
\section*{Second Volume, Version V29: Complete Massive Stieltjes Denominator and Exact Euclidean Root Criterion}
\addcontentsline{toc}{section}{Second Volume, Version V29: Complete Massive Stieltjes Denominator and Exact Euclidean Root Criterion}

\subsection*{Purpose}

V28 classified the zeros of the ultraviolet logarithmic model.  This version
retains the complete positive massive spectral density instead.  After the
three local subtraction directions have been fixed, the nonlocal Euclidean
form factor is a Stieltjes-derived function \(H\).  Its first two
derivatives imply strict convexity of the reduced inverse propagator.  This
gives an exact root trichotomy and an implicit necessary-and-sufficient
zero-free condition for the full massive TT spectral denominator.  No
large-momentum replacement by a pure logarithm is used.

\subsection{Subtracted positive spectral function}
\label{sec:v2v29-Stieltjes}

Let \(s_0>0\) and let
\(\rho:[s_0,\infty)\to[0,\infty)\) be locally integrable.  Assume
\begin{equation}
 c_-s^2\le\rho(s)\le c_+s^2
 \quad (s\ge S_0)
\label{eq:v2v29-spectral-growth}
\end{equation}
for constants \(c_\pm>0\).  The scalar TT density in
\eqref{eq:v2v27-spectral} has these properties with \(s_0=4m^2\).

For \(x>0\), define
\begin{equation}
 H(x)
 =
 \int_{s_0}^{\infty}
 \frac{\rho(s)}{s^3}\frac{x}{s+x}\,\mathrm ds.
\label{eq:v2v29-H}
\end{equation}

\begin{lemma}[Well-definedness and derivatives]
\label{lem:v2v29-H}
The integral \eqref{eq:v2v29-H} is finite for every \(x>0\),
\(H(0+)=0\), and
\begin{align}
 H'(x)
 &=
 \int_{s_0}^{\infty}
 \frac{\rho(s)}{s^3}
 \frac{s}{(s+x)^2}\,\mathrm ds>0,
\label{eq:v2v29-H-prime}\\
 H''(x)
 &=
 -2\int_{s_0}^{\infty}
 \frac{\rho(s)}{s^3}
 \frac{s}{(s+x)^3}\,\mathrm ds<0.
\label{eq:v2v29-H-second}
\end{align}
Moreover
\begin{equation}
 c_1\log(1+x)-C_1
 \le H(x)\le
 C_2\log(1+x)+C_3.
\label{eq:v2v29-H-log}
\end{equation}
\end{lemma}

\begin{proof}
At large \(s\), the integrand in \eqref{eq:v2v29-H} is
\(O(x/s^2)\), so it is integrable.  The positive mass threshold controls
the lower endpoint.  Dominated convergence gives \(H(0+)=0\) and permits
the two displayed derivatives on compact \(x\)-intervals.

For \(S_0\le s\le x\), the factor \(x/(s+x)\) is bounded above and below by
positive constants, while \(\rho(s)/s^3\simeq1/s\).  This gives the
logarithmic lower bound for large \(x\).  Splitting the integral at
\(S_0\), \(x\), and infinity, and using \(x/(s+x)\le\min\{1,x/s\}\),
gives the upper bound.  Adjust the constants on bounded \(x\)-intervals.
\end{proof}

\begin{remark}[Why the mass threshold is stated]
\label{rem:v2v29-massless}
If a physical channel begins at \(s_0=0\), subtraction at zero momentum can
have an infrared logarithm.  Such a channel must be subtracted at a
nonzero reference point or treated with an infrared prescription.  V29's
zero criterion is proved for the massive spectral sector and is not applied
unchanged to photons or other massless thresholds.
\end{remark}

\subsection{The complete reduced denominator}
\label{sec:v2v29-denominator}

After fixing the local cosmological, Einstein, and curvature-squared
directions, write the Euclidean TT inverse as
\begin{equation}
 \mathfrak D_\rho(x)
 =
 xF_\rho(x),
 \qquad
 F_\rho(x)
 =
 M^2+x\big[a+H(x)\big],
\label{eq:v2v29-denominator}
\end{equation}
where \(M^2>0\) and \(a\in\mathbb R\) is the fixed finite local
fourth-order coefficient in this normalization.

\begin{lemma}[Strict convexity]
\label{lem:v2v29-convexity}
For every \(x>0\),
\begin{align}
 F_\rho'(x)
 &=
 a+H(x)+xH'(x),
\label{eq:v2v29-F-prime}\\
 F_\rho''(x)
 &=
 2\int_{s_0}^{\infty}
 \frac{\rho(s)}{s^3}
 \frac{s^2}{(s+x)^3}\,\mathrm ds
 >0.
\label{eq:v2v29-F-second}
\end{align}
Also
\[
 F_\rho(0+)=M^2,
 \qquad
 F_\rho(x)\longrightarrow+\infty
 \quad (x\to\infty).
\]
\end{lemma}

\begin{proof}
Differentiate \eqref{eq:v2v29-denominator}.  Using
\eqref{eq:v2v29-H-prime}--\eqref{eq:v2v29-H-second},
\begin{align*}
 2H'(x)+xH''(x)
 &=
 2\int
 \frac{\rho(s)}{s^3}
 \left[
 \frac{s}{(s+x)^2}
 -\frac{xs}{(s+x)^3}
 \right]\mathrm ds\\
 &=
 2\int
 \frac{\rho(s)}{s^3}
 \frac{s^2}{(s+x)^3}\,\mathrm ds.
\end{align*}
It is strictly positive because \(\rho\) is nonzero.  The endpoint
statements follow from \(H(0+)=0\) and
\eqref{eq:v2v29-H-log}.
\end{proof}

\begin{theorem}[Exact full-massive Euclidean zero criterion]
\label{thm:v2v29-zero}
The positive zeros of \(\mathfrak D_\rho\), apart from its Einstein factor
at \(x=0\), obey the following trichotomy.
\begin{enumerate}
 \item If \(a\ge0\), then \(F_\rho(x)>0\) for all \(x>0\).
 \item If \(a<0\), there is a unique \(x_*>0\) satisfying
 \begin{equation}
  a+H(x_*)+x_*H'(x_*)=0.
 \label{eq:v2v29-critical}
 \end{equation}
 At that point
 \begin{equation}
  F_\rho(x_*)=M^2-x_*^2H'(x_*).
 \label{eq:v2v29-minimum}
 \end{equation}
 Hence:
 \begin{align}
  M^2>x_*^2H'(x_*)&\quad\Longleftrightarrow\quad
  \text{no positive zero},\label{eq:v2v29-no-zero}\\
  M^2=x_*^2H'(x_*)&\quad\Longleftrightarrow\quad
  \text{one double positive zero},\label{eq:v2v29-double}\\
  M^2<x_*^2H'(x_*)&\quad\Longleftrightarrow\quad
  \text{exactly two simple positive zeros}.
 \label{eq:v2v29-two}
 \end{align}
\end{enumerate}
\end{theorem}

\begin{proof}
If \(a\ge0\), positivity of \(H\) gives
\(F_\rho(x)\ge M^2>0\).  Suppose \(a<0\).  Then
\(F_\rho'(0+)=a<0\), while
\eqref{eq:v2v29-H-log} implies \(F_\rho'(x)\to+\infty\).
Strict convexity gives one and only one critical point, which is the global
minimum.  At that point, \eqref{eq:v2v29-critical} gives
\(a+H(x_*)=-x_*H'(x_*)\); substitution yields
\eqref{eq:v2v29-minimum}.  The positive endpoint values, strict convexity,
and the sign of the unique minimum give the three alternatives.
\end{proof}

\begin{corollary}[Exact coercivity in the zero-free regime]
\label{cor:v2v29-coercive}
If the strict no-zero condition holds, put
\[
 \gamma=\inf_{x>0}F_\rho(x)>0.
\]
Then
\begin{equation}
 \mathfrak D_\rho(x)\ge\gamma x,
\label{eq:v2v29-coercive}
\end{equation}
and the Euclidean inverse maps
\(\dot H^r\) to \(\dot H^{r+2}\) with norm at most \(\gamma^{-1}\).
At high momentum it has the stronger pointwise decay
\[
 \mathfrak D_\rho(x)^{-1}
 =
 O\big(x^{-2}(\log x)^{-1}\big).
\]
\end{corollary}

\begin{proof}
The first claim is the definition of \(\gamma\) and
\eqref{eq:v2v29-denominator}.  Plancherel gives the Sobolev estimate exactly
as in \cref{cor:v2v28-coercive}.  Finally,
\eqref{eq:v2v29-H-log} gives
\(\mathfrak D_\rho(x)\simeq x^2\log x\) up to fixed positive constants at
large \(x\).
\end{proof}

\subsection{Positive addition of physical channels}
\label{sec:v2v29-addition}

\begin{theorem}[Adding positive spectral weight cannot create a Euclidean
zero]
\label{thm:v2v29-monotone}
Let \(\rho_1,\rho_2\ge0\) satisfy the hypotheses above, and keep \(M^2\) and
the physical local coefficient \(a\) fixed.  Then
\begin{equation}
 F_{\rho_1+\rho_2}(x)
 =
 F_{\rho_1}(x)+xH_{\rho_2}(x)
 \ge F_{\rho_1}(x).
\label{eq:v2v29-monotone}
\end{equation}
Thus if the \(\rho_1\) denominator is positive for every \(x>0\), the
denominator with the additional physical channel is also positive.
\end{theorem}

\begin{proof}
The map \(\rho\mapsto H_\rho\) is linear and has a nonnegative kernel in
\eqref{eq:v2v29-H}.  Substitute in
\eqref{eq:v2v29-denominator}.
\end{proof}

\begin{firewall}[The local coefficient must be held physically fixed]
\label{fw:v2v29-monotone}
Integrating out an additional field also changes the relation between a bare
coefficient and the renormalized \(a\).  The monotonicity theorem compares
theories at the same fixed renormalized gravitational coefficient.  It does
not say that an arbitrary bare-parameter prescription is monotone.
\end{firewall}

\begin{assessment}[Complete positive SM carrier]
\label{ass:v2v29-SM-carrier}
For a unitary matter theory, the physical stress spectral density is a sum
over positive intermediate-state norms.  At the level of the Euclidean TT
two-point function, all massive SM channels should therefore be assembled
in one nonnegative \(\rho_{\mathrm{SM}}\) before tensor estimates, in direct
analogy with the complete positive YM bank.  Gauge fixing and ghosts may
enter intermediate calculations, but the final physical spectral density
must be tested for positivity after the BRST cohomology projection.
\end{assessment}

\subsection{Relation to the V28 logarithmic criterion}
\label{sec:v2v29-relation}

\begin{proposition}[Ultraviolet reduction to the logarithmic model]
\label{prop:v2v29-log}
If
\[
 \rho(s)/s^2\longrightarrow c_{\mathrm{eff}}>0,
\]
then
\begin{equation}
 H(x)=c_{\mathrm{eff}}\log x+o(\log x).
\label{eq:v2v29-log-weak}
\end{equation}
If, in addition,
\begin{equation}
 \int_{S_0}^{\infty}
 \left|\frac{\rho(s)}{s^2}-c_{\mathrm{eff}}\right|
 \frac{\mathrm ds}{s}<\infty,
\label{eq:v2v29-log-rate}
\end{equation}
then there is a constant \(b_{\mathrm{eff}}\) such that
\begin{equation}
 H(x)
 =
 c_{\mathrm{eff}}\log(x/\mu^2)+b_{\mathrm{eff}}+o(1).
\label{eq:v2v29-log}
\end{equation}
The constant \(b_{\mathrm{eff}}\) is absorbed into \(a\).  Thus V28 is the
ultraviolet asymptotic of \eqref{eq:v2v29-denominator}, while
\cref{thm:v2v29-zero} retains the exact threshold region.
\end{proposition}

\begin{proof}
Write \(\rho(s)=s^2[c_{\mathrm{eff}}+r(s)]\), where \(r(s)\to0\),
and split the integral at a fixed large scale.  The constant part is compared
with
\[
 c_{\mathrm{eff}}\int
 \frac{x}{s(s+x)}\,\mathrm ds.
\]
Its primitive is
\(c_{\mathrm{eff}}[\log s-\log(s+x)]\), giving
\(c_{\mathrm{eff}}\log x\) plus a constant.  Splitting the remainder at
large fixed \(S\), using \(|r(s)|<\varepsilon\) for \(s>S\), and then
letting \(\varepsilon\downarrow0\) proves \eqref{eq:v2v29-log-weak}.
Under \eqref{eq:v2v29-log-rate}, dominated convergence applies to
\(r(s)x/[s(s+x)]\), producing a finite limiting constant and proving
\eqref{eq:v2v29-log}.
\end{proof}

\begin{firewall}[Asymptotic convergence needs a rate]
\label{fw:v2v29-asymptotic}
The final \(o(1)\) in \eqref{eq:v2v29-log} uses the explicit integrable
rate \eqref{eq:v2v29-log-rate}.  Without that rate, only
\eqref{eq:v2v29-log-weak} and the two-sided bound
\eqref{eq:v2v29-H-log} are proved; slower sublogarithmic terms may remain.  The exact root criterion
\cref{thm:v2v29-zero} does not require the sharper asymptotic.
\end{firewall}

\subsection{What positivity still does not control}
\label{sec:v2v29-retarded}

The analytic continuation of \eqref{eq:v2v29-H} is
\begin{equation}
 H(z)
 =
 \int_{s_0}^{\infty}
 \frac{\rho(s)}{s^3}\frac{z}{s+z}\,\mathrm ds,
 \qquad
 z\in\mathbb C\setminus(-\infty,-s_0].
\label{eq:v2v29-analytic}
\end{equation}
It maps the upper half-plane to itself, because each
\(z/(s+z)\) has positive imaginary part there.  The reduced denominator,
however, contains \(zH(z)\), which need not have a fixed imaginary sign.

\begin{firewall}[Stieltjes positivity is not yet first-sheet stability]
\label{fw:v2v29-first-sheet}
Positivity and strict convexity prove the complete statement on the positive
Euclidean axis.  They do not exclude complex zeros of
\[
 M^2+z[a+H(z)]
\]
in the cut plane.  Multiplication by \(z\) destroys the direct Herglotz sign
that \(H\) itself possesses.  A first-sheet zero theorem therefore requires
an additional argument; it cannot be inferred from
\(\rho\ge0\) alone.
\end{firewall}

\subsection{V29 conclusion and next step}
\label{sec:v2v29-conclusion}

The complete massive TT spectral density yields the finite positive
function \(H\) in \eqref{eq:v2v29-H}.  Its derivative identities make the
reduced Euclidean denominator strictly convex.  For \(a<0\), the single
critical point \(x_*\) gives the exact zero-free condition
\(M^2>x_*^2H'(x_*)\); equality gives a double zero and reversal gives two
simple zeros.  Adding positive physical spectral weight at fixed
renormalized \(a\) cannot create a Euclidean zero.

The remaining issue is genuinely complex analytic.  V30 is the next
mandatory companion audit and will then test a contour or argument-principle
criterion for zeros of the complete first-sheet function, retaining the
massless-channel infrared distinction and the physical BRST spectral
projection.
\clearpage
\section*{Second Volume, Version V30: Companion Audit and a Right-Half-Plane Zero Theorem}
\addcontentsline{toc}{section}{Second Volume, Version V30: Companion Audit and a Right-Half-Plane Zero Theorem}

\subsection*{Purpose and mandatory audit}

This is the required five-version companion audit following V25.  The
current sources inspected were:
\begin{center}
\begin{tabular}{p{0.47\textwidth}r p{0.32\textwidth}}
\toprule
file & bytes & SHA--256\\
\midrule
\texttt{Renewal\_NS\_Stability\_Research\_Notes\_V31.tex}
&887537&
\texttt{F1121B62238AD5491BB7CF03635EA9934942EDF573ED8FAAA9EE79E1D38A6696}\\
\texttt{Renewal\_Yang\_Mills\_Mass\_Gap\_Research\_Notes\_V22.tex}
&281756&
\texttt{C7E6FF1152510E6A61B8C131CFB9224B8F3130D265A6E7EA7C8378C5141B5C5A}\\
\bottomrule
\end{tabular}
\end{center}
The NS file and digest remain unchanged.  The YM folder also contains a
V21-named file with the same byte count and digest as V22.  The V22-named
source itself ends with the compact V21 record, so this audit attributes no
unwritten V22 theorem to the filename.

YM V20 closes complete nondual incidence-two \(SU(2)\) banks outside output
priority.  V21 then separates the unique protected all-singlet line before
applying a complete bank resolvent; the gapped complement is bounded without
output counting, while the protected line retains its external factor
\(z/(1-z)\).  The output-priority protected carrier remains open.

\begin{assessment}[V30 audit transfer]
\label{ass:v2v30-audit}
The spectral denominator has an analogous protected direction: its factor
\(z=0\) is the massless Einstein pole.  It must be factored out before the
matter polarization is tested for additional zeros.  The reduced function
\[
 F_\rho(z)=M^2+z[a+H(z)]
\]
retains the external Einstein coefficient \(M^2\), just as the YM protected
line retains its external payer.  Dividing away or estimating \(M^2\)
separately would destroy the sign argument below.
\end{assessment}

\subsection{Positive-real property of the massive spectral kernel}
\label{sec:v2v30-positive-real}

Retain the massive density and analytic function from V29:
\begin{equation}
 H(z)
 =
 \int_{s_0}^{\infty}
 \frac{\rho(s)}{s^3}\frac{z}{s+z}\,\mathrm ds,
 \qquad
 z\in\mathbb C\setminus(-\infty,-s_0].
\label{eq:v2v30-H}
\end{equation}

\begin{lemma}[The kernel is positive real on the right half-plane]
\label{lem:v2v30-positive-real}
If \(z=x+iy\) with \(x\ge0\) and \(z\ne0\), then
\begin{equation}
 \operatorname{Re}\frac{z}{s+z}
 =
 \frac{x(s+x)+y^2}{(s+x)^2+y^2}>0
\label{eq:v2v30-kernel-real}
\end{equation}
for every \(s>0\).  Consequently,
\begin{equation}
 \operatorname{Re}H(z)>0
\label{eq:v2v30-H-real}
\end{equation}
whenever \(\rho\) is nonzero.
\end{lemma}

\begin{proof}
Multiply \(z/(s+z)\) by
\((s+\overline z)/(s+\overline z)\) and take the real part.  The numerator
in \eqref{eq:v2v30-kernel-real} is positive for \(x\ge0\), \(z\ne0\).
Integrate it against the nonnegative, nonzero weight \(\rho(s)/s^3\).
\end{proof}

\begin{remark}[This is stronger than the Herglotz sign used in V29]
\label{rem:v2v30-sign}
For \(\operatorname{Im}z>0\), the same kernel also has positive imaginary
part.  V29 observed that multiplication by \(z\) loses a direct Herglotz
sign.  The right-half-plane argument instead divides the zero equation by
\(z\) and uses the positive real part of \(H\).
\end{remark}

\subsection{A first-sheet zero exclusion region}
\label{sec:v2v30-zero}

Define
\begin{equation}
 F_\rho(z)
 =
 M^2+z[a+H(z)],
 \qquad
 M^2>0.
\label{eq:v2v30-F}
\end{equation}

\begin{theorem}[No reduced zero in the closed right half-plane]
\label{thm:v2v30-no-zero}
If \(a\ge0\), then
\begin{equation}
 F_\rho(z)\ne0
 \qquad
 \text{for every }z\ne0
 \text{ with }\operatorname{Re}z\ge0.
\label{eq:v2v30-no-zero}
\end{equation}
Also \(F_\rho(0)=M^2\ne0\).  Hence the complete denominator
\[
 \mathfrak D_\rho(z)=zF_\rho(z)
\]
has exactly its protected Einstein zero at \(z=0\) in the closed right
half-plane and no additional zero there.
\end{theorem}

\begin{proof}
Suppose first that \(x=\operatorname{Re}z>0\) and \(F_\rho(z)=0\).
After division by \(z\),
\begin{equation}
 a+H(z)=-\frac{M^2}{z}.
\label{eq:v2v30-zero-equation}
\end{equation}
The real part of the left side is strictly positive by
\cref{lem:v2v30-positive-real} and \(a\ge0\).  The real part of the right
side is
\[
 -\frac{M^2x}{|z|^2}<0,
\]
a contradiction.

If \(x=0\) and \(z=iy\ne0\), the right side of
\eqref{eq:v2v30-zero-equation} is purely imaginary, whereas the left side
has strictly positive real part.  Thus no boundary zero exists.  The value
at the origin follows directly from \eqref{eq:v2v30-F}.
\end{proof}

\begin{corollary}[Sectorial lower bound]
\label{cor:v2v30-sector}
Fix \(0<\delta\le1\).  If
\[
 \operatorname{Re}z\ge\delta|z|>0,
\]
then
\begin{equation}
 |F_\rho(z)|\ge\delta M^2,
 \qquad
 |\mathfrak D_\rho(z)^{-1}|
 \le\frac{1}{\delta M^2|z|}.
\label{eq:v2v30-sector}
\end{equation}
\end{corollary}

\begin{proof}
Write
\[
 \frac{F_\rho(z)}z
 =
 \frac{M^2}{z}+a+H(z).
\]
Every term in its real part is nonnegative, and
\[
 \operatorname{Re}\frac{M^2}{z}
 =
 \frac{M^2\operatorname{Re}z}{|z|^2}
 \ge\frac{\delta M^2}{|z|}.
\]
Therefore
\[
 |F_\rho(z)|
 =
 |z|\left|\frac{F_\rho(z)}z\right|
 \ge\delta M^2.
\]
Multiply by \(|z|\) for the inverse bound.
\end{proof}

\begin{assessment}[A genuine complex zero theorem]
\label{ass:v2v30-complex}
V29 proved positivity only on the positive Euclidean axis.  V30 excludes
zeros on the entire closed right half-plane of the first-sheet cut domain
when \(a\ge0\).  The proof uses the complete massive spectral integral and
does not use the ultraviolet logarithmic approximation.
\end{assessment}

\subsection{Positive aggregation preserves the theorem}
\label{sec:v2v30-aggregation}

\begin{corollary}[Complete physical massive carrier]
\label{cor:v2v30-aggregation}
Let
\(\rho_{\mathrm{tot}}=\sum_A\rho_A\), with every
\(\rho_A\ge0\) satisfying the massive integrability assumptions.  At fixed
\(M^2>0\) and \(a\ge0\), the denominator built from
\(\rho_{\mathrm{tot}}\) obeys
\cref{thm:v2v30-no-zero,cor:v2v30-sector}.
No bound depends on the number of intermediate-state channels.
\end{corollary}

\begin{proof}
The sum may be inserted inside \eqref{eq:v2v30-H} by monotone convergence.
The real-part identity remains positive, and the proofs above are
unchanged.
\end{proof}

\begin{assessment}[Relation to the YM complete-bank rule]
\label{ass:v2v30-YM}
The proof sums positive physical spectral channels before estimating the
denominator.  It therefore has no channel-counting loss, paralleling the
complete-output YM resolvent bound.  BRST quartets and ghosts are not
separately positive channels; positivity is asserted only after projection
to the physical state space.
\end{assessment}

\subsection{The massless infrared distinction}
\label{sec:v2v30-massless}

For a massless channel with \(\rho(s)\simeq cs^2\) down to \(s=0\), the
zero-subtracted integral
\eqref{eq:v2v30-H} contains
\[
 \int_0\frac{x}{s(s+x)}\,\mathrm ds
\]
and is infrared divergent.  A subtraction at a nonzero spacelike reference
\(x_0\) produces a finite difference, but its kernel is no longer exactly
\eqref{eq:v2v30-kernel-real}.

\begin{firewall}[Massless channels are not inserted by continuity in the
threshold]
\label{fw:v2v30-massless}
The limit \(s_0\downarrow0\) of the V29--V30 zero-subtracted representation
is singular.  Photon, perturbative gluon, and other massless sectors require
a nonzero subtraction point together with their infrared state and
inclusive-observable prescription.  The massive right-half-plane theorem
does not automatically extend to them.
\end{firewall}

\subsection{Relation to retarded stability}
\label{sec:v2v30-retarded}

The analytic variable \(z\) is Euclidean squared momentum.  For a Laplace
frequency \(p\) and spatial frequency \(\xi\), the hyperbolic continuation
samples
\[
 z=p^2+|\xi|^2.
\]
Even when \(\operatorname{Re}p>0\), this \(z\) need not have nonnegative real
part because
\[
 \operatorname{Re}z
 =
 (\operatorname{Re}p)^2-(\operatorname{Im}p)^2+|\xi|^2.
\]

\begin{firewall}[Right-half \(z\) is not the full Laplace stability domain]
\label{fw:v2v30-Laplace}
\Cref{thm:v2v30-no-zero} excludes every zero whose squared-momentum image
has nonnegative real part.  It does not exclude a growing Laplace mode whose
image lies in the left half of the cut \(z\)-plane.  Therefore the theorem
is a strict advance in first-sheet localization, not yet a causal
well-posedness theorem.
\end{firewall}

\begin{corollary}[Localization of any remaining first-sheet zero]
\label{cor:v2v30-localization}
Under \(a\ge0\), every nonzero first-sheet zero of the complete massive
reduced denominator must satisfy
\begin{equation}
 \operatorname{Re}z<0
\label{eq:v2v30-left}
\end{equation}
and must lie off the excluded spectral cut.
\end{corollary}

\begin{proof}
Apply \cref{thm:v2v30-no-zero}.
\end{proof}

\subsection{V30 conclusion and next step}
\label{sec:v2v30-conclusion}

The mandatory audit records the unchanged NS V31 source and the YM
V21 protected-line separation contained in the latest V22-named file.
Following that rule, V30 keeps the Einstein zero factored from the complete
massive reduced denominator.  The subtracted spectral kernel has strictly
positive real part on the closed right half-plane.  For \(a\ge0\), this
excludes every additional zero there and gives the sectorial inverse bound
\eqref{eq:v2v30-sector}, uniformly after aggregation of all positive
massive physical channels.

Any remaining first-sheet zero lies in the left half-plane, where the
spectral cut and hyperbolic continuation matter.  V31 will formulate an
argument-principle count on that cut plane and derive a computable
sufficient condition, in terms of the boundary spectral phase, for the
absence of left-half-plane first-sheet zeros.
\clearpage
\section*{Second Volume, Version V31: Positive Spectral Density Can Still Produce Runaway Poles}
\addcontentsline{toc}{section}{Second Volume, Version V31: Positive Spectral Density Can Still Produce Runaway Poles}

\subsection*{Purpose}

V30 localized every possible additional zero to the left half of the
first-sheet \(z\)-plane and proposed an argument-principle exclusion.
Before attempting that proof, this version tests whether positivity of the
complete spectral density is sufficient.  It is not.  A one-atom positive
spectral measure gives an exactly solvable reduced denominator with a
complex conjugate pair of left-half-plane zeros.  Those zeros persist after
the atom is smoothed and a positive \(s^2\) ultraviolet tail is added.
Every nonreal \(z\)-zero has a Laplace square root with positive real part,
so the unreduced equation has a runaway homogeneous mode.  This no-go
forces the next step to use the actual SM spectral function and physical
couplings, or to adopt a controlled order-reduced theory.

\subsection{A positive one-atom spectral model}
\label{sec:v2v31-atom}

Extend V29's definition from densities to positive spectral measures and
choose
\begin{equation}
 \mathrm d\rho_0(s)
 =
 g s_0^3\delta_{s_0}(\mathrm ds),
 \qquad
 g>0,\quad s_0>0.
\label{eq:v2v31-atom}
\end{equation}
Then
\begin{equation}
 H_0(z)
 =
 g\frac{z}{s_0+z},
 \qquad
 z\in\mathbb C\setminus\{-s_0\}.
\label{eq:v2v31-H}
\end{equation}
The reduced denominator is
\begin{equation}
 F_0(z)
 =
 M^2+z\left[a+g\frac{z}{s_0+z}\right].
\label{eq:v2v31-F}
\end{equation}

\begin{lemma}[Exact polynomial numerator]
\label{lem:v2v31-polynomial}
Away from \(z=-s_0\),
\begin{equation}
 F_0(z)
 =
 \frac{P_0(z)}{s_0+z},
\label{eq:v2v31-rational}
\end{equation}
where
\begin{equation}
 P_0(z)
 =
 (a+g)z^2+(M^2+as_0)z+M^2s_0.
\label{eq:v2v31-polynomial}
\end{equation}
\end{lemma}

\begin{proof}
Multiply \eqref{eq:v2v31-F} by \(s_0+z\) and collect powers of \(z\).
\end{proof}

\begin{theorem}[Exact conjugate first-sheet zeros]
\label{thm:v2v31-roots}
Assume \(a+g>0\) and
\begin{equation}
 \mathscr D
 =
 (M^2+as_0)^2-4(a+g)M^2s_0<0.
\label{eq:v2v31-discriminant}
\end{equation}
Then \(F_0\) has two simple nonreal zeros
\begin{equation}
 z_\pm
 =
 \frac{
 -(M^2+as_0)\pm i\sqrt{-\mathscr D}
 }{2(a+g)}.
\label{eq:v2v31-roots}
\end{equation}
If \(a\ge0\), they satisfy
\begin{equation}
 \operatorname{Re}z_\pm<0.
\label{eq:v2v31-left}
\end{equation}
They lie off the spectral singularity and are zeros on the analytic first
sheet of the rational model.
\end{theorem}

\begin{proof}
Apply the quadratic formula to
\eqref{eq:v2v31-polynomial}.  A negative discriminant gives two simple
conjugate roots.  Under \(a\ge0\), the numerator of their real part is
strictly negative and the denominator is positive.  A nonreal root cannot
equal the real pole \(-s_0\).
\end{proof}

\begin{corollary}[The phenomenon occurs with \(a=0\)]
\label{cor:v2v31-a-zero}
At vanishing finite curvature-squared coefficient,
\begin{equation}
 F_0(z)=M^2+g\frac{z^2}{s_0+z}.
\end{equation}
It has a complex conjugate zero pair whenever
\begin{equation}
 4gs_0>M^2.
\label{eq:v2v31-simple-condition}
\end{equation}
\end{corollary}

\begin{proof}
With \(a=0\),
\(\mathscr D=M^4-4gM^2s_0\).
\end{proof}

\begin{assessment}[Compatibility with V30]
\label{ass:v2v31-V30}
The roots in \eqref{eq:v2v31-roots} have negative real part, exactly as
required by \cref{thm:v2v30-no-zero}.  Thus V30's right-half-plane theorem
cannot be extended to the full cut plane using spectral positivity alone.
\end{assessment}

\subsection{Opposite residues and the higher-derivative pair}
\label{sec:v2v31-residues}

\begin{proposition}[Residue identity]
\label{prop:v2v31-residue}
At either simple root,
\begin{equation}
 \operatorname*{Res}_{z=z_\pm}\frac1{F_0(z)}
 =
 \frac{s_0+z_\pm}{P_0'(z_\pm)},
 \qquad
 P_0'(z_\pm)
 =
 \pm i\sqrt{-\mathscr D}.
\label{eq:v2v31-residue}
\end{equation}
The two residues are conjugate.  When the discriminant is positive and both
roots are real, the signs of \(P_0'\) at the ordered roots are opposite.
\end{proposition}

\begin{proof}
Differentiate \eqref{eq:v2v31-polynomial}.  Substitution of
\eqref{eq:v2v31-roots} gives the displayed derivative.  The residue formula
follows from \eqref{eq:v2v31-rational}.
\end{proof}

\subsection{Persistence for smooth physical-type densities}
\label{sec:v2v31-persistence}

The atom is a diagnostic model rather than a claim about the exact SM
density.  Its zeros are stable under positive smoothing.

\begin{theorem}[Positive smooth perturbations retain the complex zeros]
\label{thm:v2v31-persistence}
Assume \eqref{eq:v2v31-discriminant}.  Let \(C_\pm\) be disjoint small
circles around \(z_\pm\), avoiding \(-s_0\).  There exists
\(\varepsilon_0>0\) such that if an analytic perturbation
\(\delta H\) satisfies
\begin{equation}
 \sup_{z\in C_+\cup C_-}|z\,\delta H(z)|
 <
 \varepsilon_0,
\label{eq:v2v31-small}
\end{equation}
then
\begin{equation}
 F_\varepsilon(z)
 =
 F_0(z)+z\,\delta H(z)
\label{eq:v2v31-perturbed}
\end{equation}
has one zero inside each \(C_\pm\).

Moreover, \(\delta H\) may be produced by replacing the atom with a
nonnegative smooth bump of the same mass near \(s_0\) and adding a
nonnegative smooth tail
\(\varepsilon\rho_{\mathrm{uv}}(s)\) satisfying
\(\rho_{\mathrm{uv}}(s)\simeq s^2\) at infinity, for sufficiently small
bump width and \(\varepsilon>0\).
\end{theorem}

\begin{proof}
Because the two roots are simple, shrink the circles so that
\[
 m_\pm=\inf_{z\in C_\pm}|F_0(z)|>0.
\]
Take \(\varepsilon_0<\min(m_+,m_-)\).  Rouche's theorem applied to
\(F_0\) and \(z\delta H\) gives the same number of zeros inside each
circle, namely one.

For a smooth bump converging weakly to
\(g s_0^3\delta_{s_0}\), its Stieltjes transform converges uniformly on
compact sets separated from \(-s_0\).  A positive ultraviolet tail with
the V29 integrability condition contributes
\[
 \delta H_{\mathrm{uv}}(z)
 =
 \varepsilon\int
 \frac{\rho_{\mathrm{uv}}(s)}{s^3}
 \frac{z}{s+z}\,\mathrm ds,
\]
which is uniformly \(O(\varepsilon)\) on the two fixed circles.  Choose the
bump width and then \(\varepsilon\) small enough for
\eqref{eq:v2v31-small}.
\end{proof}

\begin{corollary}[Positivity plus \(s^2\) growth is insufficient]
\label{cor:v2v31-no-universal}
There exist smooth strictly nonnegative massive spectral densities obeying
\eqref{eq:v2v29-spectral-growth} whose reduced denominator has nonreal
first-sheet zeros even when \(a=0\).
\end{corollary}

\begin{proof}
Choose the atomic parameters with the strict inequality
\eqref{eq:v2v31-simple-condition}, then apply
\cref{thm:v2v31-persistence}.
\end{proof}

\subsection{From a \(z\)-zero to a runaway mode}
\label{sec:v2v31-runaway}

For a spatial Fourier mode \(\xi\), the Laplace continuation uses
\begin{equation}
 z=p^2+|\xi|^2.
\label{eq:v2v31-Laplace}
\end{equation}

\begin{theorem}[Every nonreal first-sheet zero generates a growing square
root]
\label{thm:v2v31-runaway}
Let \(z_0\notin\mathbb R\) satisfy \(F(z_0)=0\).  At \(\xi=0\), the equation
\(p^2=z_0\) has two roots \(p_\pm=\pm\sqrt{z_0}\), exactly one of which has
positive real part.  The corresponding homogeneous mode
\begin{equation}
 h(t)=e^{p_+t}e_{\mathrm{TT}}
\label{eq:v2v31-runaway}
\end{equation}
grows exponentially.
\end{theorem}

\begin{proof}
A nonreal number has square roots with nonzero opposite real parts.  Choose
the one with positive real part.  Substitution of
\eqref{eq:v2v31-Laplace} into the Fourier--Laplace symbol makes the
denominator vanish, so the exponential is a homogeneous mode.
\end{proof}

\begin{firewall}[A contour prescription can change the problem]
\label{fw:v2v31-contour}
One may define a bounded solution by deforming a contour to exclude a
runaway pole or by imposing a future boundary condition.  Such a
prescription is not the standard retarded Cauchy problem: it can introduce
pre-response or remove freely specifiable initial data.  A continuum
theorem must state the prescription and prove its causal and stability
properties rather than silently omitting the pole.
\end{firewall}

\subsection{Consequences for the all-frequency branch}
\label{sec:v2v31-consequence}

\begin{theorem}[No positivity-only retarded stability theorem]
\label{thm:v2v31-no-positivity}
The assumptions
\[
 \rho\ge0,\qquad
 \rho(s)\simeq s^2,\qquad
 M^2>0,\qquad
 a\ge0
\]
do not imply absence of first-sheet zeros or stability of the unreduced
retarded semiclassical equation.
\end{theorem}

\begin{proof}
Use the smooth positive densities in
\cref{cor:v2v31-no-universal}.  Their zeros give growing Laplace roots by
\cref{thm:v2v31-runaway}.
\end{proof}

\begin{assessment}[What data are now indispensable]
\label{ass:v2v31-data}
The all-frequency semiclassical branch requires the actual renormalized SM
spectral density, the measured gravitational coefficients, and a global
zero analysis of their complete denominator.  General unitarity positivity
and ultraviolet power counting do not decide stability.  This is a sharper
gate than the proposed V30 boundary-phase criterion.
\end{assessment}

\begin{programme}[Revised acquisition order]
\label{prog:v2v31}
Proceed in the following order:
\begin{enumerate}
 \item compute or bound the complete free-SM spin-two and spin-zero spectral
 densities, including massive and massless thresholds;
 \item fix a nonzero subtraction prescription for the massless channels and
 match all local coefficients to physical renormalization data;
 \item test the resulting first-sheet denominator for zeros, analytically
 where possible and with certified numerics otherwise;
 \item if a first-sheet runaway remains, formulate a controlled
 low-energy order reduction and prove its error only below an explicit
 scale;
 \item return to nonlinear self-consistency only after one of these linear
 propagators is proved causal and stable.
\end{enumerate}
\end{programme}

\subsection{V31 conclusion and next step}
\label{sec:v2v31-conclusion}

A positive massive spectral density can generate nonreal first-sheet zeros.
The one-atom model gives their exact quadratic formula, and Rouche stability
shows that the phenomenon persists for smooth positive densities with the
same \(s^2\) ultraviolet growth as a stress spectral function.  Each
nonreal zero yields an exponentially growing Laplace mode.  Therefore
spectral positivity, right-half-plane exclusion, and Euclidean coercivity do
not close the retarded all-frequency problem.

V32 will construct the alternative low-energy branch.  It will define
order reduction as a finite Neumann expansion of the causal polarization on
a bandlimited source class and prove an explicit remainder bound in terms
of the dimensionless ratio
\(G_{\mathrm R}|K^2\log(-K^2/\mu^2)|\).
\clearpage
\section*{Second Volume, Version V32: Controlled Causal Order Reduction Below the Matter Threshold}
\addcontentsline{toc}{section}{Second Volume, Version V32: Controlled Causal Order Reduction Below the Matter Threshold}

\subsection*{Purpose}

V31 proves that a positive spectral density can still create first-sheet
runaway poles in the unreduced all-frequency equation.  This version
constructs the alternative low-energy branch.  The complete retarded matter
polarization is treated as a causal perturbation of the Einstein solution
operator.  If their dimensionless composition has norm below one, the exact
bandlimited solution is a convergent causal Neumann series.  Every finite
order is causal, and the discarded remainder has an explicit geometric
bound.  The result also gives regulator stability and a nonlinear
contraction criterion.

\subsection{Abstract causal operator theorem}
\label{sec:v2v32-abstract}

Let \(X_T\) be a Banach space of metric perturbations on a slab
\(I=[0,T]\), and let \(Y_T\) be a Banach source space.  Let
\[
 E:Y_T\longrightarrow X_T
\]
be the retarded solution operator for the gauge-fixed linearized Einstein
equation with zero initial data.  Let
\[
 \Pi:X_T\longrightarrow Y_T
\]
be the renormalized retarded linear matter-polarization operator, including
the fixed local finite terms.  Define
\begin{equation}
 K=E\Pi:X_T\longrightarrow X_T.
\label{eq:v2v32-K}
\end{equation}

\begin{theorem}[Causal Neumann construction]
\label{thm:v2v32-Neumann}
Assume
\begin{equation}
 \|K\|_{X_T\to X_T}\le\eta<1.
\label{eq:v2v32-small}
\end{equation}
Then for every \(S\in Y_T\), the equation
\begin{equation}
 h+Kh=ES
\label{eq:v2v32-equation}
\end{equation}
has the unique solution
\begin{equation}
 h
 =
 \sum_{j=0}^{\infty}(-K)^jES
\label{eq:v2v32-series}
\end{equation}
in \(X_T\), with
\begin{equation}
 \|h\|_{X_T}
 \le
 \frac{\|E\|}{1-\eta}\|S\|_{Y_T}.
\label{eq:v2v32-solution-bound}
\end{equation}
If \(E\) and \(\Pi\) are causal, then every partial sum and the limit are
causal.
\end{theorem}

\begin{proof}
The norm bound makes the geometric operator series
\(\sum_{j\ge0}(-K)^j\) converge absolutely to \((I+K)^{-1}\).
This proves existence, uniqueness, and
\eqref{eq:v2v32-solution-bound}.  A composition and finite sum of retarded
operators is retarded.  The future-supported subspace is closed in
\(X_T\), so the norm limit is retarded.
\end{proof}

\begin{definition}[Order-\(N\) causal reduction]
\label{def:v2v32-order}
Define
\begin{equation}
 h^{[N]}
 =
 \sum_{j=0}^{N}(-K)^jES.
\label{eq:v2v32-order}
\end{equation}
\end{definition}

\begin{theorem}[Explicit truncation error]
\label{thm:v2v32-error}
Under \eqref{eq:v2v32-small},
\begin{equation}
 \|h-h^{[N]}\|_{X_T}
 \le
 \frac{\eta^{N+1}}{1-\eta}
 \|E\|\,\|S\|_{Y_T}.
\label{eq:v2v32-error}
\end{equation}
The residual in \eqref{eq:v2v32-equation} is exactly
\begin{equation}
 (I+K)h^{[N]}-ES
 =
 (-1)^NK^{N+1}ES.
\label{eq:v2v32-residual}
\end{equation}
\end{theorem}

\begin{proof}
The error is the tail of the geometric series.  Multiplying the finite sum
by \(I+K\) telescopes and gives
\eqref{eq:v2v32-residual}.
\end{proof}

\begin{assessment}[Meaning of order reduction]
\label{ass:v2v32-meaning}
The approximation is defined by a finite composition of retarded operators,
not by discarding selected high derivatives inside the differential
equation.  Its causal prescription and error are therefore explicit.
Different truncation orders approximate one fixed full retarded equation as
long as \(\eta<1\).
\end{assessment}

\subsection{Massive spectral low-energy bound}
\label{sec:v2v32-massive}

Retain the massive TT spectral function of V29.  For
\(|z|\le\kappa^2<s_0/2\), the cut is avoided and
\begin{equation}
 H(z)
 =
 \int_{s_0}^{\infty}
 \frac{\rho(s)}{s^3}\frac{z}{s+z}\,\mathrm ds.
\label{eq:v2v32-H}
\end{equation}
Put
\begin{equation}
 C_\rho
 =
 2\int_{s_0}^{\infty}\frac{\rho(s)}{s^4}\,\mathrm ds<\infty.
\label{eq:v2v32-Crho}
\end{equation}

\begin{lemma}[Analytic threshold-disc estimate]
\label{lem:v2v32-disc}
For \(|z|\le\kappa^2<s_0/2\),
\begin{equation}
 |H(z)|\le C_\rho|z|.
\label{eq:v2v32-H-bound}
\end{equation}
Consequently the ratio of the polarization part to the Einstein part obeys
\begin{equation}
 |r(z)|
 =
 \left|
 \frac{z[a+H(z)]}{M^2}
 \right|
 \le
 \eta_\kappa
 :=
 \frac{\kappa^2}{M^2}
 \left(|a|+C_\rho\kappa^2\right).
\label{eq:v2v32-eta}
\end{equation}
\end{lemma}

\begin{proof}
Since \(|s+z|\ge s-|z|\ge s/2\),
\[
 \left|\frac{z}{s+z}\right|
 \le\frac{2|z|}{s}.
\]
Insert this in \eqref{eq:v2v32-H}; the result is
\eqref{eq:v2v32-H-bound}.  The ratio bound follows directly.
The integral in \eqref{eq:v2v32-Crho} converges because
\(\rho(s)=O(s^2)\) at infinity and \(s_0>0\).
\end{proof}

\begin{corollary}[Explicit perturbative energy scale]
\label{cor:v2v32-scale}
If
\begin{equation}
 \frac{\kappa^2}{M^2}
 \left(|a|+C_\rho\kappa^2\right)<1,
\label{eq:v2v32-scale}
\end{equation}
then the stationary TT Fourier multiplier is pointwise Neumann-convergent
throughout the disc \(|z|\le\kappa^2\), and its order-\(N\) relative
multiplier error is at most
\begin{equation}
 \frac{\eta_\kappa^{N+1}}{1-\eta_\kappa}.
\label{eq:v2v32-relative}
\end{equation}
\end{corollary}

\begin{proof}
Factor the denominator as
\[
 M^2z[1+r(z)]
\]
and apply the scalar geometric series uniformly on the disc.
\end{proof}

\begin{firewall}[A pointwise multiplier bound is not yet an energy bound]
\label{fw:v2v32-multiplier}
Equation \eqref{eq:v2v32-scale} verifies the algebraic small parameter on a
stationary band.  To invoke \cref{thm:v2v32-Neumann} on a finite curved
slab, one must prove the operator estimate
\(\|E\Pi\|_{X_T\to X_T}\le\eta\).  A supremum of a flat Fourier symbol does
not automatically control variable coefficients, boundaries, or the
Einstein characteristic cone.
\end{firewall}

\subsection{Regulator stability}
\label{sec:v2v32-regulator}

Let \(E_\Lambda,\Pi_\Lambda,S_\Lambda\) be cutoff operators and sources, and
put \(K_\Lambda=E_\Lambda\Pi_\Lambda\).

\begin{theorem}[Uniform reduced-continuum estimate]
\label{thm:v2v32-regulator}
Assume
\begin{align}
 \|K_\Lambda\|,\|K\|&\le\eta<1,
\label{eq:v2v32-uniform}\\
 \|E_\Lambda-E\|_{Y_T\to X_T}
 +\|K_\Lambda-K\|_{X_T\to X_T}
 +\|S_\Lambda-S\|_{Y_T}
 &\longrightarrow0.
\label{eq:v2v32-operator-convergence}
\end{align}
Let
\[
 h_\Lambda=(I+K_\Lambda)^{-1}E_\Lambda S_\Lambda,
 \qquad
 h=(I+K)^{-1}ES.
\]
Then
\begin{align}
 \|h_\Lambda-h\|_{X_T}
 \le&
 \frac1{1-\eta}
 \|E_\Lambda S_\Lambda-ES\|_{X_T}
 \nonumber\\
 &+
 \frac{\|E\|\,\|S\|_{Y_T}}{(1-\eta)^2}
 \|K_\Lambda-K\|_{X_T\to X_T}.
\label{eq:v2v32-regulator-rate}
\end{align}
In particular, the entire cutoff family converges.
\end{theorem}

\begin{proof}
The uniform Neumann bound gives
\[
 \|(I+K_\Lambda)^{-1}\|,
 \|(I+K)^{-1}\|
 \le(1-\eta)^{-1}.
\]
Use the resolvent identity
\[
 (I+K_\Lambda)^{-1}-(I+K)^{-1}
 =
 (I+K_\Lambda)^{-1}(K-K_\Lambda)(I+K)^{-1}
\]
and add and subtract
\((I+K_\Lambda)^{-1}ES\).
\end{proof}

\begin{corollary}[Truncation and cutoff limits commute uniformly]
\label{cor:v2v32-commute}
Under \cref{thm:v2v32-regulator}, let
\[
 h_\Lambda^{[N]}
 =
 \sum_{j=0}^N(-K_\Lambda)^jE_\Lambda S_\Lambda.
\]
The truncation error is bounded uniformly in \(\Lambda\) by
\[
 \frac{\eta^{N+1}}{1-\eta}
 \sup_\Lambda\|E_\Lambda\|\,\|S_\Lambda\|,
\]
and for each fixed \(N\),
\(h_\Lambda^{[N]}\to h^{[N]}\).
Thus either cutoff removal followed by \(N\to\infty\), or the reverse
order, gives \(h\).
\end{corollary}

\begin{proof}
Use \cref{thm:v2v32-error} uniformly.  Finite products converge in operator
norm for fixed \(N\).  A standard two-parameter epsilon argument gives the
commutation of limits.
\end{proof}

\subsection{A nonlinear reduced fixed point}
\label{sec:v2v32-nonlinear}

Let \(\mathcal N:\mathcal B_R\subset X_T\to Y_T\) collect the nonlinear
Einstein remainder and the nonlinear state-dependent mean-stress remainder
after the linear polarization \(\Pi\) has been extracted.  Consider
\begin{equation}
 h+Kh=E[S+\mathcal N(h)].
\label{eq:v2v32-nonlinear}
\end{equation}

\begin{theorem}[Low-energy nonlinear contraction]
\label{thm:v2v32-nonlinear}
Assume \(\|K\|\le\eta<1\), the map
\[
 \Psi(h)=(I+K)^{-1}E[S+\mathcal N(h)]
\]
preserves a closed ball \(\mathcal B_R\), and
\begin{equation}
 \|\mathcal N(h)-\mathcal N(\widetilde h)\|_{Y_T}
 \le L_N\|h-\widetilde h\|_{X_T},
 \qquad
 \frac{\|E\|L_N}{1-\eta}<1.
\label{eq:v2v32-nonlinear-small}
\end{equation}
Then \eqref{eq:v2v32-nonlinear} has a unique solution in
\(\mathcal B_R\).
\end{theorem}

\begin{proof}
The uniform inverse bound gives
\[
 \|\Psi(h)-\Psi(\widetilde h)\|_{X_T}
 \le
 \frac{\|E\|L_N}{1-\eta}
 \|h-\widetilde h\|_{X_T}.
\]
Apply Banach's fixed-point theorem.
\end{proof}

\begin{assessment}[The linear response is no longer charged as a source]
\label{ass:v2v32-linear-left}
Unlike the failed V26 map, \(\Pi h\) is inverted together with the Einstein
operator.  The nonlinear Lipschitz constant in
\eqref{eq:v2v32-nonlinear-small} applies only to the remainder after this
linear response is extracted.  This removes the \(N^3\log N\) derivative
loss found in V27 from the contraction ledger on the controlled band.
\end{assessment}

\subsection{Scope of the low-energy continuum}
\label{sec:v2v32-scope}

\begin{firewall}[Band limitation is part of the theorem]
\label{fw:v2v32-band}
The condition \eqref{eq:v2v32-scale} holds only below an explicit spectral
scale and away from unresolved massless infrared prescriptions.  It does
not define an all-frequency continuum and does not remove the first-sheet
poles constructed in V31 outside the perturbative band.  Calling the
order-reduced solution the full semiclassical solution without its
bandlimit and remainder estimate would change the claim.
\end{firewall}

\begin{programme}[Completion data for this branch]
\label{prog:v2v32-completion}
To turn \cref{thm:v2v32-regulator,thm:v2v32-nonlinear} into a curved
SM--Einstein theorem, prove:
\begin{enumerate}
 \item a covariant spectral or microlocal projector defining the low-energy
 class on the evolving geometry;
 \item a uniform causal bound for \(E_\Lambda\Pi_\Lambda\) below that scale;
 \item regulator convergence in the operator norms of
 \eqref{eq:v2v32-operator-convergence};
 \item conservation of the separated linear and nonlinear source pieces;
 \item the nonlinear remainder estimate
 \eqref{eq:v2v32-nonlinear-small};
 \item compatibility of adjacent slabs so the construction can be
 continued while the low-energy and hyperbolicity margins persist.
\end{enumerate}
\end{programme}

\subsection{V32 conclusion and next step}
\label{sec:v2v32-conclusion}

Below a massive spectral threshold, the dimensionless flat polarization
ratio obeys the explicit bound
\[
 \eta_\kappa
 =
 \frac{\kappa^2}{M^2}(|a|+C_\rho\kappa^2).
\]
Whenever the corresponding causal operator norm is below one, the exact
retarded solution is a convergent Neumann series about the Einstein
propagator.  Its order-\(N\) truncation error is geometric, cutoff removal
is stable and commutes with truncation, and a nonlinear remainder closes
under the explicit contraction condition
\eqref{eq:v2v32-nonlinear-small}.

This constructs a rigorous abstract low-energy continuum mechanism, not the
full all-frequency theory.  V33 will attack the first curved-space datum in
\cref{prog:v2v32-completion}: replace a sharp Fourier band by a covariant
spectral cutoff of a positive spatial Laplace-type operator and prove its
commutator and time-variation bounds on a uniformly hyperbolic slab.
\clearpage
\section*{Second Volume, Version V33: Covariant Moving Spectral Bands on a Hyperbolic Slab}
\addcontentsline{toc}{section}{Second Volume, Version V33: Covariant Moving Spectral Bands on a Hyperbolic Slab}

\subsection*{Purpose}

V32 used a flat band \(|z|\le\kappa^2\).  A self-consistent curved metric has
no preferred spatial Fourier transform, and its low-energy subspace changes
with time.  This version defines the band by smooth functional calculus of a
positive spatial Laplace-type operator.  It proves covariance, uniform
Sobolev bounds, time-variation estimates, and rapidly small leakage beyond a
separated spectral buffer.  The calculation also records a necessary cost:
the moving projector has an order-zero time derivative and is not small
merely because \(\kappa\) is large.

\subsection{Uniform geometric setup}
\label{sec:v2v33-setup}

Let \(\Sigma\) be a compact manifold without boundary and let
\(\gamma(t)\), \(t\in[0,T]\), be a family of Riemannian metrics with a common
ellipticity constant.  Conjugate all \(L^2(\Sigma,d\mu_{\gamma(t)})\) spaces
unitarily to one fixed density.  Let
\begin{equation}
 A(t)=I-\Delta_{\gamma(t)}+\mathcal R(t)
\label{eq:v2v33-A}
\end{equation}
be a positive self-adjoint Laplace-type operator on the relevant tensor
bundle, where the order-zero endomorphism \(\mathcal R\) is chosen so that
\begin{equation}
 A(t)\ge cI
\label{eq:v2v33-positive}
\end{equation}
uniformly.

Choose \(\chi\in C_c^\infty([0,\infty))\) with
\[
 0\le\chi\le1,\qquad
 \chi(\lambda)=1\ (0\le\lambda\le1),\qquad
 \chi(\lambda)=0\ (\lambda\ge2),
\]
and define
\begin{equation}
 P_\kappa(t)=\chi(A(t)/\kappa^2),
 \qquad
 \kappa\ge1.
\label{eq:v2v33-projector}
\end{equation}
This is a smooth cutoff rather than an idempotent projection.

\begin{assumption}[Coefficient control]
\label{ass:v2v33-coefficients}
For the Sobolev orders used below, assume the coefficients of \(A(t)\), their
spatial derivatives, and one time derivative are uniformly bounded in a
class sufficient for parameter-dependent pseudodifferential calculus.  Put
their controlling bound into \(V_T\).
\end{assumption}

\subsection{Covariance and Sobolev bounds}
\label{sec:v2v33-covariance}

\begin{theorem}[Spatial covariance]
\label{thm:v2v33-covariance}
If \(\varphi:\Sigma\to\Sigma\) is a spatial diffeomorphism and
\(U_\varphi\) is its natural unitary pullback on the bundle, then
\begin{equation}
 A_{\varphi^*\gamma}
 =
 U_\varphi^{-1}A_\gamma U_\varphi
 \quad\Longrightarrow\quad
 P_{\kappa,\varphi^*\gamma}
 =
 U_\varphi^{-1}P_{\kappa,\gamma}U_\varphi.
\label{eq:v2v33-covariance}
\end{equation}
\end{theorem}

\begin{proof}
The first identity is naturality of the Laplace-type operator.  The second
follows from uniqueness of the spectral functional calculus under unitary
conjugation.
\end{proof}

\begin{theorem}[Uniform Sobolev multiplier bounds]
\label{thm:v2v33-Sobolev}
For every fixed \(r\in\mathbb R\),
\begin{equation}
 \sup_{t,\kappa\ge1}
 \|P_\kappa(t)\|_{H^r\to H^r}
 \le C_{r,V_T}.
\label{eq:v2v33-Sobolev}
\end{equation}
For every \(m,N\ge0\),
\begin{equation}
 \|P_\kappa(t)u\|_{H^{r+m}}
 \le C_{r,m}\kappa^m\|u\|_{H^r},
\label{eq:v2v33-Bernstein}
\end{equation}
and
\begin{equation}
 \|(I-P_\kappa(t))u\|_{H^r}
 \le C_{r,N}\kappa^{-N}\|u\|_{H^{r+N}}
\label{eq:v2v33-tail}
\end{equation}
when the high-pass cutoff is interpreted with a nested smooth partition.
\end{theorem}

\begin{proof}
The parameter-dependent functional calculus makes
\(\chi(A/\kappa^2)\) a uniform order-zero pseudodifferential operator with
symbol supported where the spatial covariable has size \(O(\kappa)\).
The order-zero mapping theorem gives \eqref{eq:v2v33-Sobolev}.  Multiplying
its symbol by \(\langle\xi\rangle^m\) costs at most \(C\kappa^m\), proving
\eqref{eq:v2v33-Bernstein}.  On the complementary nested high symbol,
\(\langle\xi\rangle^{-N}\le C\kappa^{-N}\); the uniform symbolic remainder
has the same bound, proving \eqref{eq:v2v33-tail}.
\end{proof}

\subsection{Time variation}
\label{sec:v2v33-time}

For an almost-analytic extension \(\widetilde\chi\), the
Helffer--Sjostrand formula gives
\begin{equation}
 P_\kappa(t)
 =
 -\frac1\pi\int_{\mathbb C}
 \overline\partial\widetilde\chi(z)
 \left(z-\frac{A(t)}{\kappa^2}\right)^{-1}
 \,\mathrm d^2z.
\label{eq:v2v33-HS}
\end{equation}

\begin{theorem}[Derivative of the moving band]
\label{thm:v2v33-time}
The strong derivative exists on every Sobolev space in the controlled range
and equals
\begin{align}
 \dot P_\kappa(t)
 =
 -\frac1{\pi\kappa^2}\int_{\mathbb C}
 \overline\partial\widetilde\chi(z)
 \left(z-\frac{A}{\kappa^2}\right)^{-1}
 \dot A
 \left(z-\frac{A}{\kappa^2}\right)^{-1}
 \,\mathrm d^2z.
\label{eq:v2v33-Pdot}
\end{align}
It obeys
\begin{equation}
 \|\dot P_\kappa(t)u\|_{H^r}
 \le C_{r,V_T}\|u\|_{H^r},
\label{eq:v2v33-Pdot-zero}
\end{equation}
and, for \(0\le\theta\le2\),
\begin{equation}
 \|\dot P_\kappa(t)u\|_{H^r}
 \le C_{r,\theta,V_T}\kappa^{-\theta}
 \|u\|_{H^{r+\theta}}.
\label{eq:v2v33-Pdot-scale}
\end{equation}
\end{theorem}

\begin{proof}
Differentiate the resolvent identity in
\eqref{eq:v2v33-HS}; the derivative of
\((z-A/\kappa^2)^{-1}\) gives the two resolvents and
\(\dot A/\kappa^2\), proving \eqref{eq:v2v33-Pdot}.
Uniform parameter-dependent resolvent estimates and the fact that
\(\dot A\) has order two make the complete operator order zero on the
transition band \(|\xi|\simeq\kappa\), proving
\eqref{eq:v2v33-Pdot-zero}.  On that band,
\(\|u\|_{H^r}\le C\kappa^{-\theta}\|u\|_{H^{r+\theta}}\).
The off-band symbolic terms are rapidly smoothing; interpolation gives
\eqref{eq:v2v33-Pdot-scale} for \(0\le\theta\le2\).
\end{proof}

\begin{firewall}[Time variation is not cutoff-small at fixed regularity]
\label{fw:v2v33-time}
The factor \(\kappa^{-2}\) visible in
\eqref{eq:v2v33-Pdot} is cancelled by the order-two size of \(\dot A\) on
the transition shell.  Therefore
\eqref{eq:v2v33-Pdot-zero} is order one.  Claiming an unconditional
\(\kappa^{-1}\) or \(\kappa^{-2}\) gain in \(H^r\to H^r\) would lose the
motion of the principal metric symbol.
\end{firewall}

\subsection{Separated-buffer leakage}
\label{sec:v2v33-buffer}

Choose \(\chi_{\mathrm{out}}\in C_c^\infty([0,\infty))\) equal to one on a
neighbourhood of \(\operatorname{supp}\chi\), and supported in a strictly
larger compact interval.  Put
\[
 P_\kappa^{\mathrm{out}}(t)
 =
 \chi_{\mathrm{out}}(A(t)/\kappa^2),
 \qquad
 Q_\kappa^{\mathrm{far}}=I-P_\kappa^{\mathrm{out}}.
\]
At each fixed time,
\begin{equation}
 Q_\kappa^{\mathrm{far}}P_\kappa=0.
\label{eq:v2v33-nested}
\end{equation}

\begin{theorem}[Rapid leakage bound beyond the buffer]
\label{thm:v2v33-leakage}
For every \(r,N\), under sufficient coefficient regularity,
\begin{equation}
 \|Q_\kappa^{\mathrm{far}}(t)\dot P_\kappa(t)u\|_{H^r}
 \le
 C_{r,N,V_T}\kappa^{-N}\|u\|_{H^{r+N}}.
\label{eq:v2v33-leakage}
\end{equation}
Thus motion of the inner band can populate its transition buffer at order
one, but direct leakage to a spectrally separated far region is rapidly
small.
\end{theorem}

\begin{proof}
The principal parameter-dependent symbol of \(\dot P_\kappa\) is supported
where a derivative of the inner cutoff is active.  The symbol of
\(Q_\kappa^{\mathrm{far}}\) vanishes on a neighbourhood of that set.
The composition therefore has zero complete principal symbol at every
finite order.  Iterating the symbolic composition formula leaves a
parameter-smoothing remainder of size \(O(\kappa^{-N})\) from
\(H^{r+N}\) to \(H^r\).
\end{proof}

\subsection{Moving-band evolution identity}
\label{sec:v2v33-evolution}

Let \(u_\kappa=P_\kappa(t)u(t)\).

\begin{proposition}[Exact differentiated cutoff identity]
\label{prop:v2v33-evolution}
For every differentiable \(u\),
\begin{equation}
 \partial_tu_\kappa
 =
 P_\kappa\partial_tu+\dot P_\kappa u.
\label{eq:v2v33-evolution}
\end{equation}
Consequently,
\begin{align}
 \|\partial_tu_\kappa\|_{H^r}
 \le
 C\left(
 \|\partial_tu\|_{H^r}+\|u\|_{H^r}
 \right),
\label{eq:v2v33-evolution-bound}
\end{align}
and
\begin{equation}
 \|Q_\kappa^{\mathrm{far}}\partial_tu_\kappa\|_{H^r}
 \le
 C_{r,N}\kappa^{-N}\|u\|_{H^{r+N}}.
\label{eq:v2v33-far-evolution}
\end{equation}
\end{proposition}

\begin{proof}
Differentiate the product to get
\eqref{eq:v2v33-evolution}.  Apply
\cref{thm:v2v33-Sobolev,thm:v2v33-time} for the first bound.
For the second, use
\(Q_\kappa^{\mathrm{far}}P_\kappa=0\) and
\cref{thm:v2v33-leakage}.
\end{proof}

\begin{assessment}[The correct curved low-energy object]
\label{ass:v2v33-object}
A curved low-energy solution is not a field with a time-independent finite
Fourier support.  It is a section tracked by the moving covariant band,
together with an order-one transition-buffer term
\(\dot P_\kappa u\).  Only leakage beyond a separated outer band is rapidly
small.  The buffer must therefore be included in the V32 operator norm.
\end{assessment}

\subsection{Commutators with local coefficients}
\label{sec:v2v33-commutator}

\begin{theorem}[Multiplication commutator]
\label{thm:v2v33-commutator}
For a smooth bundle endomorphism \(b(t,x)\),
\begin{equation}
 \|[b,P_\kappa(t)]u\|_{H^r}
 \le
 C_{r,V_T}\kappa^{-1}
 \|b\|_{C^{r_*}}\|u\|_{H^r},
\label{eq:v2v33-commutator}
\end{equation}
for a finite \(r_*\) depending on \(r\) and the dimension.
For a first-order variable-coefficient differential operator \(B\),
the corresponding \(H^r\to H^r\) commutator is generally only \(O(1)\).
\end{theorem}

\begin{proof}
In semiclassical notation \(h=\kappa^{-1}\), the leading symbol of
\([b,\chi(h^2A)]\) is
\(ih\{b,\chi(a_2)\}\), so it has one factor \(h\); the uniform remainder
has the same bound by the parameter calculus.  If \(B\) has first-order
symbol of size \(h^{-1}\), that size cancels the commutator factor \(h\),
leaving an order-zero bound.
\end{proof}

\begin{firewall}[A covariant cutoff does not commute with evolution]
\label{fw:v2v33-commute}
Spatial covariance of \(P_\kappa(t)\) does not imply commutation with the
hyperbolic operator.  Lapse, shift, connection, and metric time variation
produce the order-one terms identified above.  The low-energy
Einstein--matter equation must include these commutators in its remainder
map; omitting them would invalidate the V32 contraction constant.
\end{firewall}

\subsection{Insertion into the V32 branch}
\label{sec:v2v33-insertion}

\begin{theorem}[Conditional curved-band reduction]
\label{thm:v2v33-insertion}
Let \(X_{T,\kappa}^{\mathrm{buf}}\) include the inner range of \(P_\kappa(t)\)
and the transition range of \(P_\kappa^{\mathrm{out}}(t)\).  Suppose the
retarded Einstein operator \(E_g\), the renormalized polarization
\(\Pi_g\), and all cutoff commutators define a causal operator
\[
 K_{g,\kappa}:X_{T,\kappa}^{\mathrm{buf}}
 \longrightarrow X_{T,\kappa}^{\mathrm{buf}}
\]
with
\begin{equation}
 \sup_{g\in\mathcal B_R}
 \|K_{g,\kappa}\|\le\eta<1.
\label{eq:v2v33-insertion}
\end{equation}
Then the V32 causal Neumann theorem, truncation estimate, and regulator
stability hold on this covariant moving band.
\end{theorem}

\begin{proof}
The covariance theorem makes the band definition coordinate independent.
The moving-band and commutator theorems identify every additional term
created by projection; by hypothesis all are included in
\(K_{g,\kappa}\).  Apply
\cref{thm:v2v32-Neumann,thm:v2v32-error,thm:v2v32-regulator}.
\end{proof}

\begin{firewall}[The decisive norm remains conditional]
\label{fw:v2v33-conditional}
V33 constructs the correct covariant band and proves its geometric
calculus.  It does not prove \eqref{eq:v2v33-insertion} for the SM
polarization on an evolving metric.  In particular, the order-one
\(\dot P_\kappa\) term must be made small by slab length, adiabatic metric
control, or a refined transport, not by a nonexistent inverse cutoff power.
\end{firewall}

\subsection{V33 conclusion and next step}
\label{sec:v2v33-conclusion}

The smooth spectral cutoff
\(P_\kappa(t)=\chi(A(t)/\kappa^2)\) is spatially covariant and uniformly
bounded on Sobolev spaces.  Its time derivative is an order-zero transition
operator.  Motion of the band can therefore populate a fixed buffer at
order one, while leakage beyond a separated outer band is rapidly small.
Multiplication commutators gain one cutoff power, but first-order evolution
commutators generally do not.

V34 will construct a parallel-transported spectral frame and compare it
with \(P_\kappa(t)\).  The aim is to move the order-one projector derivative
into an explicit skew connection, preserving the low-band norm while
isolating the genuine off-band transition that enters the contraction
constant.
\clearpage
\section*{Second Volume, Version V34: Kato Transport of a Gapped Covariant Band}
\addcontentsline{toc}{section}{Second Volume, Version V34: Kato Transport of a Gapped Covariant Band}

\subsection*{Purpose}

V33 proves that the derivative of a smooth moving spectral cutoff is order
zero.  This version asks whether that term can be absorbed into a
norm-preserving connection.  For a genuine orthogonal spectral projection
separated by a uniform gap, Kato's commutator connection gives an exact
unitary transport between the time-dependent low-energy spaces.  The
transport removes internal frame motion and leaves only the projected
physical evolution.  A Weyl-law audit then shows why a relative hard gap
cannot persist uniformly at arbitrarily high cutoffs.  The construction is
therefore valid for a fixed gapped physical band, while V33's smooth buffer
remains necessary for cutoff-uniform work.

\subsection{A gapped spectral projection}
\label{sec:v2v34-projection}

Retain the positive self-adjoint family \(A(t)\) of V33.  Fix
\(\kappa>0\) and suppose there is \(g_\kappa>0\) such that
\begin{equation}
 \operatorname{spec}A(t)\cap
 [\kappa^2-g_\kappa,\kappa^2+g_\kappa]
 =
 \varnothing
 \qquad (0\le t\le T).
\label{eq:v2v34-gap}
\end{equation}
Define the orthogonal Riesz projection
\begin{equation}
 \mathsf P(t)
 =
 \mathbf1_{[0,\kappa^2]}(A(t)).
\label{eq:v2v34-P}
\end{equation}

\begin{lemma}[Differentiated projection identities]
\label{lem:v2v34-identities}
If \(\mathsf P\) is strongly differentiable, then
\begin{align}
 \mathsf P\dot{\mathsf P}\mathsf P&=0,
\label{eq:v2v34-PPP}\\
 (I-\mathsf P)\dot{\mathsf P}(I-\mathsf P)&=0.
\label{eq:v2v34-QQQ}
\end{align}
Thus \(\dot{\mathsf P}\) is off diagonal relative to
\(\operatorname{Ran}\mathsf P\oplus\ker\mathsf P\).
\end{lemma}

\begin{proof}
Differentiate \(\mathsf P^2=\mathsf P\):
\[
 \dot{\mathsf P}\mathsf P+\mathsf P\dot{\mathsf P}
 =
 \dot{\mathsf P}.
\]
Multiply on both sides by \(\mathsf P\) to get
\eqref{eq:v2v34-PPP}; multiply by \(I-\mathsf P\) to get
\eqref{eq:v2v34-QQQ}.
\end{proof}

\subsection{The Kato connection}
\label{sec:v2v34-Kato}

Define
\begin{equation}
 \mathsf G(t)
 =
 [\dot{\mathsf P}(t),\mathsf P(t)].
\label{eq:v2v34-G}
\end{equation}

\begin{lemma}[Skewness and transport identity]
\label{lem:v2v34-G}
One has
\begin{equation}
 \mathsf G^*=-\mathsf G,
 \qquad
 [\mathsf G,\mathsf P]=\dot{\mathsf P}.
\label{eq:v2v34-G-identities}
\end{equation}
Also
\[
 \mathsf P\mathsf G\mathsf P=0,
 \qquad
 (I-\mathsf P)\mathsf G(I-\mathsf P)=0.
\]
\end{lemma}

\begin{proof}
Both \(\mathsf P\) and \(\dot{\mathsf P}\) are self-adjoint, so their
commutator is skew-adjoint.  Use
\cref{lem:v2v34-identities} and expand the double commutator:
\[
 [[\dot{\mathsf P},\mathsf P],\mathsf P]
 =
 \dot{\mathsf P}\mathsf P+\mathsf P\dot{\mathsf P}
 =
 \dot{\mathsf P}.
\]
The diagonal identities follow from
\eqref{eq:v2v34-PPP}--\eqref{eq:v2v34-QQQ}.
\end{proof}

\begin{theorem}[Exact unitary band transport]
\label{thm:v2v34-transport}
Let \(\mathsf U(t)\) solve
\begin{equation}
 \dot{\mathsf U}(t)=\mathsf G(t)\mathsf U(t),
 \qquad
 \mathsf U(0)=I.
\label{eq:v2v34-U}
\end{equation}
Then \(\mathsf U(t)\) is unitary and
\begin{equation}
 \mathsf P(t)\mathsf U(t)
 =
 \mathsf U(t)\mathsf P(0).
\label{eq:v2v34-intertwine}
\end{equation}
Hence \(\mathsf U(t)\) maps the initial low-energy space isometrically onto
the instantaneous low-energy space.
\end{theorem}

\begin{proof}
Skew-adjointness gives
\[
 \partial_t(\mathsf U^*\mathsf U)
 =
 \mathsf U^*(\mathsf G^*+\mathsf G)\mathsf U=0.
\]
Thus \(\mathsf U\) is unitary.  Differentiate
\(\mathsf U^{-1}\mathsf P\mathsf U\):
\[
 \partial_t(\mathsf U^{-1}\mathsf P\mathsf U)
 =
 \mathsf U^{-1}
 \big(\dot{\mathsf P}-[\mathsf G,\mathsf P]\big)
 \mathsf U=0
\]
by \eqref{eq:v2v34-G-identities}.  Its value is \(\mathsf P(0)\), which is
equivalent to \eqref{eq:v2v34-intertwine}.
\end{proof}

\begin{corollary}[Band-covariant time derivative]
\label{cor:v2v34-covariant}
Define
\begin{equation}
 \mathsf D_tu=\partial_tu-\mathsf G(t)u.
\label{eq:v2v34-Dt}
\end{equation}
If \(u(t)\in\operatorname{Ran}\mathsf P(t)\), then
\begin{equation}
 (I-\mathsf P(t))\mathsf D_tu(t)=0.
\label{eq:v2v34-Dt-band}
\end{equation}
For a purely transported vector \(u(t)=\mathsf U(t)u_0\),
\(\mathsf D_tu=0\).
\end{corollary}

\begin{proof}
Differentiate \(\mathsf Pu=u\) and project with \(I-\mathsf P\):
\[
 (I-\mathsf P)\partial_tu
 =
 (I-\mathsf P)\dot{\mathsf P}u.
\]
For \(u=\mathsf Pu\),
\[
 \mathsf Gu
 =
 (\dot{\mathsf P}\mathsf P-\mathsf P\dot{\mathsf P})u
 =
 (I-\mathsf P)\dot{\mathsf P}u.
\]
Subtract to prove \eqref{eq:v2v34-Dt-band}.  The transported case is
\eqref{eq:v2v34-U}.
\end{proof}

\begin{assessment}[What the connection removes]
\label{ass:v2v34-removes}
The order-zero motion of the hard low-energy subspace becomes the explicit
skew connection \(\mathsf G\).  It preserves norm and does not mix vectors
within either diagonal block.  In transported coordinates
\(v=\mathsf U^{-1}u\), the moving projection is the fixed projection
\(\mathsf P(0)\).  Physical evolution still acquires conjugated
commutators with \(\mathsf G\); they have not been estimated away.
\end{assessment}

\subsection{Quantitative gap dependence}
\label{sec:v2v34-gap-bound}

Assume the quadratic-form derivative is relatively bounded:
\begin{equation}
 \|A(t)^{-1/2}\dot A(t)A(t)^{-1/2}\|_{L^2\to L^2}
 \le V_T.
\label{eq:v2v34-relative}
\end{equation}

\begin{theorem}[Projection and connection cost]
\label{thm:v2v34-gap-bound}
Under \eqref{eq:v2v34-gap}--\eqref{eq:v2v34-relative},
\begin{equation}
 \|\dot{\mathsf P}(t)\|_{L^2\to L^2}
 \le
 C V_T\frac{\kappa^2+g_\kappa}{g_\kappa},
\qquad
 \|\mathsf G(t)\|
 \le2\|\dot{\mathsf P}(t)\|.
\label{eq:v2v34-gap-bound}
\end{equation}
In particular, a relative gap
\(g_\kappa\ge\delta\kappa^2\) gives
\(\|\mathsf G\|\le C_\delta V_T\).
\end{theorem}

\begin{proof}
In the spectral representation of \(A\), the off-diagonal matrix elements
of the differentiated spectral projection are divided differences:
\[
 (\dot{\mathsf P})_{\lambda\mu}
 =
 \frac{
 \mathbf1_{\{\lambda\le\kappa^2\}}
 -\mathbf1_{\{\mu\le\kappa^2\}}
 }{\lambda-\mu}
 (\dot A)_{\lambda\mu}.
\]
Only pairs on opposite sides of the gap occur, so
\(|\lambda-\mu|\ge2g_\kappa\).  The relative form bound gives
\[
 |(\dot A)_{\lambda\mu}|
 \le V_T\sqrt{\lambda\mu}
 \le C V_T(\kappa^2+g_\kappa)
\]
for the near-gap contribution; pairs farther away are bounded by the same
divided-difference argument after dyadic spectral decomposition.  This
proves the first estimate.  The commutator definition gives the second.
\end{proof}

\begin{firewall}[The gap appears in the denominator]
\label{fw:v2v34-gap}
Kato transport reorganizes projector motion but does not make it small.
When the spectral gap shrinks, the connection grows like
\(\kappa^2/g_\kappa\).  Any curved-band contraction estimate must retain
that factor.
\end{firewall}

\subsection{Weyl-density obstruction to a uniform high hard gap}
\label{sec:v2v34-Weyl}

On a three-dimensional compact slice, the eigenvalue counting function of a
Laplace-type operator obeys
\begin{equation}
 N_A(L)
 =
 C_\Sigma L^{3/2}+o(L^{3/2}).
\label{eq:v2v34-Weyl}
\end{equation}

\begin{theorem}[No fixed relative empty window at high energy]
\label{thm:v2v34-no-relative-gap}
For every fixed \(0<\delta<1\), the interval
\[
 [(1-\delta)\kappa^2,(1+\delta)\kappa^2]
\]
contains eigenvalues for all sufficiently large \(\kappa\).  Hence
\eqref{eq:v2v34-gap} cannot hold with
\(g_\kappa=\delta\kappa^2\) along an unbounded cutoff sequence.
Moreover, the mean consecutive eigenvalue spacing near
\(L=\kappa^2\) is of order \(\kappa^{-1}\), so a typical hard-gap Kato
factor \(\kappa^2/g_\kappa\) grows like \(\kappa^3\).
\end{theorem}

\begin{proof}
By \eqref{eq:v2v34-Weyl},
\begin{align*}
 &N_A((1+\delta)\kappa^2)
 -N_A((1-\delta)\kappa^2)\\
 &\qquad=
 C_\Sigma\big[(1+\delta)^{3/2}-(1-\delta)^{3/2}\big]\kappa^3
 +o(\kappa^3),
\end{align*}
which is positive and tends to infinity.  Thus the interval is not empty.
Differentiating the leading Weyl count heuristically, and equivalently
dividing a bounded \(L\)-window by its asymptotic eigenvalue count, gives
mean spacing
\[
 (N_A'(L))^{-1}\simeq C L^{-1/2}=C\kappa^{-1}.
\]
Substitution gives the stated typical Kato factor.  The mean-spacing
statement is an average and does not assert an upper bound on every
individual gap.
\end{proof}

\begin{assessment}[Fixed physical band versus cutoff removal]
\label{ass:v2v34-fixed}
For the V32 low-energy branch, \(\kappa\) is a fixed physical scale while
the ultraviolet regulator is removed.  A gapped Kato frame may therefore
be useful if the evolving spectrum avoids that fixed threshold.  It cannot
be used as a uniform family of hard projections with
\(\kappa\to\infty\).  The smooth inner/outer bands of V33 remain the robust
construction when eigenvalues cross or become dense.
\end{assessment}

\subsection{Insertion into the reduced evolution}
\label{sec:v2v34-insertion}

Let \(u=\mathsf Uv\).  For an evolution
\[
 \partial_tu=\mathcal L(t)u+f,
\]
one has the exact transported equation
\begin{equation}
 \partial_tv
 =
 \mathsf U^{-1}(\mathcal L-\mathsf G)\mathsf Uv
 +\mathsf U^{-1}f.
\label{eq:v2v34-transported-equation}
\end{equation}

\begin{theorem}[Conditional fixed-band reduction]
\label{thm:v2v34-reduction}
Suppose a fixed \(\kappa\) obeys \eqref{eq:v2v34-gap}, the transported
Einstein--polarization operator including
\(-\mathsf U^{-1}\mathsf G\mathsf U\) has norm at most \(\eta<1\) on the
fixed space \(\operatorname{Ran}\mathsf P(0)\), and the nonlinear remainder
has contraction constant below \(1-\eta\).  Then the V32 causal fixed-point,
truncation, and regulator convergence theorems hold in the Kato frame.
\end{theorem}

\begin{proof}
The unitary intertwiner identifies every instantaneous band with one fixed
Banach subspace and preserves its Hilbert norm.  Equation
\eqref{eq:v2v34-transported-equation} includes the complete connection cost.
Apply
\cref{thm:v2v32-Neumann,thm:v2v32-error,thm:v2v32-regulator,
thm:v2v32-nonlinear}
with the transported operator.
\end{proof}

\begin{firewall}[The fixed-band norm is still unproved for the SM]
\label{fw:v2v34-SM}
V34 removes an avoidable frame-motion ambiguity.  It does not prove the
small operator norm for the interacting SM polarization, and it requires a
gap that may fail during metric evolution.  If a crossing occurs, the
construction must return to V33's smooth buffered band rather than continue
a discontinuous hard projection.
\end{firewall}

\subsection{V34 conclusion and next step}
\label{sec:v2v34-conclusion}

A uniformly gapped covariant spectral projection has an exact unitary Kato
transport.  Its skew connection absorbs the order-zero motion of the
low-energy space and converts the projector to a fixed reference band.
The connection cost is proportional to
\(V_T\kappa^2/g_\kappa\).  Weyl density excludes a fixed relative hard gap
at unbounded cutoff, so the construction is appropriate only for a fixed
physical band with no crossings.

V35 is the next mandatory companion audit.  It will then combine the smooth
buffer and Kato alternatives into a precise continuation criterion: either
a fixed gap persists and the unitary frame is used, or a crossing is
handled inside the smooth transition buffer without changing the
low-energy error bound.
\clearpage
\section*{Second Volume, Version V35: Companion Audit and a Hybrid Gap--Crossing Continuation Theorem}
\addcontentsline{toc}{section}{Second Volume, Version V35: Companion Audit and a Hybrid Gap--Crossing Continuation Theorem}

\subsection*{Purpose and mandatory audit}

This is the required companion audit five versions after V30.  The current
sources inspected were:
\begin{center}
\begin{tabular}{p{0.47\textwidth}r p{0.32\textwidth}}
\toprule
file & bytes & SHA--256\\
\midrule
\texttt{Renewal\_NS\_Stability\_Research\_Notes\_V31.tex}
&887537&
\texttt{F1121B62238AD5491BB7CF03635EA9934942EDF573ED8FAAA9EE79E1D38A6696}\\
\texttt{Renewal\_Yang\_Mills\_Mass\_Gap\_Research\_Notes\_V24.tex}
&321073&
\texttt{48AD4A08C452CEA2732F6D756D9A06D6D21170FAC63B4857F7D56F22F205DCFE}\\
\bottomrule
\end{tabular}
\end{center}
The NS source remains unchanged.  YM V24 now closes the complete
incidence-two shared-output-resolvent payer for \(SU(2)\), including mixed
balanced and strongly unbalanced banks and the protected network inherited
from V22.  It uses a positive operator decomposition followed by one-sided
insertions; it does not multiply overlapping capacities.  The
junction-output payer, higher incidence, general groups, and the full
Yang--Mills continuum remain open.

\begin{assessment}[V35 audit transfer]
\label{ass:v2v35-audit}
The hard Kato band and smooth crossing band are not independent estimates to
be multiplied.  They are two descriptions inside one complete buffered
spectral carrier.  The correct analogue of the YM one-sided decomposition is
to keep the smooth outer band for the whole slab, use the Kato frame only on
the gapped core, and return to the smooth description on crossings.  This
prevents both a product of projection losses and a discontinuous change of
state space.
\end{assessment}

\subsection{One global buffered carrier}
\label{sec:v2v35-carrier}

Retain \(P_\kappa(t)\), \(P_\kappa^{\mathrm{out}}(t)\), and
\(Q_\kappa^{\mathrm{far}}(t)\) from V33.  The evolving unknown is always
represented in
\begin{equation}
 \mathcal X_\kappa(t)
 =
 \operatorname{Ran}P_\kappa^{\mathrm{out}}(t),
\label{eq:v2v35-carrier}
\end{equation}
with the transition region included.  Choose a gap threshold
\(g_0>0\), and define
\begin{align}
 \mathcal G
 &=
 \{t:\operatorname{dist}(\kappa^2,\operatorname{spec}A(t))\ge g_0\},
\label{eq:v2v35-gap-set}\\
 \mathcal C
 &=[0,T]\setminus\mathcal G.
\label{eq:v2v35-crossing-set}
\end{align}
On each component of \(\mathcal G\), let \(\mathsf P(t)\) be the hard
projection and \(\mathsf U(t)\) its Kato transport from V34.  On
\(\mathcal C\), retain only the smooth functional-calculus frame.

\begin{lemma}[No state-space jump at a switch]
\label{lem:v2v35-no-jump}
Choose the outer cutoff equal to one on every eigenvalue below
\(\kappa^2+g_0\).  Then, on \(\mathcal G\),
\begin{equation}
 P_\kappa^{\mathrm{out}}(t)\mathsf P(t)=\mathsf P(t).
\label{eq:v2v35-embedding}
\end{equation}
Thus entering or leaving a Kato interval changes coordinates on the hard
core but does not project or discard the global buffered state.
\end{lemma}

\begin{proof}
Both operators are functions of the same self-adjoint \(A(t)\).
The scalar outer cutoff equals one on the spectral support of
\(\mathsf P(t)\), so their functional-calculus product is \(\mathsf P(t)\).
\end{proof}

\subsection{Energy charged only on the crossing set}
\label{sec:v2v35-energy}

Use the transported Hilbert norm on the hard core of each gapped component
and an equivalent Sobolev graph norm on the global buffer.  Let
\(\mathcal E_\kappa(t)\) denote their uniformly equivalent combined energy.

\begin{theorem}[Hybrid moving-band energy inequality]
\label{thm:v2v35-energy}
Suppose the projected physical evolution, excluding frame motion, satisfies
\begin{equation}
 \frac{\mathrm d}{\mathrm dt}\mathcal E_\kappa(t)
 \le
 c_0(t)\mathcal E_\kappa(t)+\|f_\kappa(t)\|_{Y}
\label{eq:v2v35-physical-energy}
\end{equation}
in a fixed frame.  Under the V33--V34 coefficient hypotheses, the hybrid
frame satisfies
\begin{equation}
 \frac{\mathrm d}{\mathrm dt}\mathcal E_\kappa(t)
 \le
 \left[
 c_0(t)+C V(t)\mathbf1_{\mathcal C}(t)
 \right]\mathcal E_\kappa(t)
 +\|f_\kappa(t)\|_Y,
\label{eq:v2v35-hybrid-energy}
\end{equation}
where \(V\) controls \(\dot A\).  Consequently,
\begin{align}
 \mathcal E_\kappa(t)
 \le
 \exp\!\left(
 \int_0^t c_0(r)\,\mathrm dr+
 C\int_{\mathcal C\cap[0,t]}V(r)\,\mathrm dr
 \right)
 \left[
 \mathcal E_\kappa(0)+
 \int_0^t\|f_\kappa(r)\|_Y\,\mathrm dr
 \right].
\label{eq:v2v35-Gronwall}
\end{align}
The bound depends on the total crossing occupation, not the number of
switches.
\end{theorem}

\begin{proof}
On a gapped component, the Kato generator is skew in the transported
Hilbert structure by \cref{thm:v2v34-transport}; it contributes no real
energy growth.  The remaining physical terms give
\eqref{eq:v2v35-physical-energy}.  On \(\mathcal C\), V33 gives
\[
 \|\dot P_\kappa^{\mathrm{out}}u\|
 \le CV(t)\|u\|.
\]
This adds \(CV(t)\mathcal E_\kappa\) to the energy inequality.  The state
itself is unchanged at the boundaries by
\cref{lem:v2v35-no-jump}, so there is no multiplicative switch factor.
Gronwall's inequality proves \eqref{eq:v2v35-Gronwall}.
\end{proof}

\begin{corollary}[Crossing continuation criterion]
\label{cor:v2v35-continuation}
The moving-band linear evolution continues through arbitrarily many
crossings on \([0,T]\) provided
\begin{equation}
 \int_0^T c_0(t)\,\mathrm dt<\infty,
 \qquad
 \int_{\mathcal C}V(t)\,\mathrm dt<\infty,
\label{eq:v2v35-crossing-criterion}
\end{equation}
and the underlying metric remains uniformly hyperbolic with the V33
coefficient bounds.
\end{corollary}

\begin{proof}
The right side of \eqref{eq:v2v35-Gronwall} remains finite and is independent
of the number of connected components of \(\mathcal C\).
\end{proof}

\begin{assessment}[Why the crossing occupation is the right scalar]
\label{ass:v2v35-occupation}
A hard-gap estimate charges \(V\kappa^2/g_\kappa\), which diverges at a
crossing.  The smooth frame charges only \(V\), but it is needed only on
\(\mathcal C\).  The hybrid ledger therefore replaces the singular inverse
gap by the finite occupation
\(\int_{\mathcal C}V\).  This resembles the NS companion's use of
restricted occupation rather than a globally repeated worst-case debit,
without importing any NS regularity claim.
\end{assessment}

\subsection{Far-band leakage through crossings}
\label{sec:v2v35-leakage}

\begin{theorem}[Integrated far leakage]
\label{thm:v2v35-leakage}
For every \(r,N\), the part of frame motion that enters the separated far
band satisfies
\begin{equation}
 \int_0^T
 \|Q_\kappa^{\mathrm{far}}
 \partial_t(P_\kappa u)\|_{H^r}\,\mathrm dt
 \le
 C_{r,N}\kappa^{-N}
 \int_{\mathcal C}
 V(t)\|u(t)\|_{H^{r+N}}\,\mathrm dt
 +\mathcal R_{\mathcal G,N},
\label{eq:v2v35-leakage}
\end{equation}
where \(\mathcal R_{\mathcal G,N}\) contains only physical evolution
commutators on the gapped set and vanishes for pure Kato transport.
\end{theorem}

\begin{proof}
On \(\mathcal C\), apply
\cref{thm:v2v33-leakage,prop:v2v33-evolution} and integrate.  On
\(\mathcal G\), pure Kato transport remains exactly in the hard band by
\eqref{eq:v2v34-intertwine}; only the separately projected physical
evolution can reach the far band, which is the stated remainder.
\end{proof}

\begin{firewall}[Higher regularity is visible]
\label{fw:v2v35-leakage}
The rapid factor \(\kappa^{-N}\) in
\eqref{eq:v2v35-leakage} is paid by \(H^{r+N}\) regularity.  It cannot be
summed with \(N\) chosen arbitrarily unless the solution has the
corresponding uniform higher norms.  At the base Sobolev level, only the
order-zero buffer estimate is available.
\end{firewall}

\subsection{Hybrid retarded polarization operator}
\label{sec:v2v35-polarization}

Let \(E_\kappa^{\mathrm{hyb}}\) be the retarded Einstein evolution operator
on the global buffer, with the Kato connection used on \(\mathcal G\) and
the smooth-frame term used on \(\mathcal C\).  Let
\(\Pi_\kappa^{\mathrm{ret}}\) be the low-energy renormalized matter
polarization and define
\begin{equation}
 K_\kappa^{\mathrm{hyb}}
 =
 E_\kappa^{\mathrm{hyb}}\Pi_\kappa^{\mathrm{ret}}.
\label{eq:v2v35-K}
\end{equation}

\begin{theorem}[Hybrid causal Neumann criterion]
\label{thm:v2v35-Neumann}
Assume
\begin{align}
 \|E_\kappa^{\mathrm{hyb}}\|_{Y_T\to X_T}
 &\le
 C_E
 \exp\!\left(
 C\int_{\mathcal C}V
 \right),
\label{eq:v2v35-E}\\
 \|K_\kappa^{\mathrm{hyb}}\|_{X_T\to X_T}
 &\le\eta<1.
\label{eq:v2v35-eta}
\end{align}
Then the complete linear low-energy equation has a unique causal solution
on the moving band, its order-\(N\) reduction error is bounded by
\begin{equation}
 \frac{\eta^{N+1}}{1-\eta}
 C_Ee^{C\int_{\mathcal C}V}\|S\|_{Y_T},
\label{eq:v2v35-error}
\end{equation}
and no factor depending on the number of gap crossings occurs.
\end{theorem}

\begin{proof}
The energy theorem gives \eqref{eq:v2v35-E}.  Apply
\cref{thm:v2v32-Neumann,thm:v2v32-error} to
\(E_\kappa^{\mathrm{hyb}}\) and
\(K_\kappa^{\mathrm{hyb}}\).
\end{proof}

\begin{theorem}[Hybrid nonlinear and regulator continuation]
\label{thm:v2v35-nonlinear}
Suppose, uniformly in the ultraviolet regulator,
\begin{enumerate}
 \item the hyperbolicity, coefficient, and crossing-occupation bounds hold;
 \item \(E_{\kappa,\Lambda}^{\mathrm{hyb}}\to
 E_\kappa^{\mathrm{hyb}}\) and
 \(K_{\kappa,\Lambda}^{\mathrm{hyb}}\to
 K_\kappa^{\mathrm{hyb}}\) in the V32 operator norms, the regulated
 sources converge in \(Y_T\), and the nonlinear remainders
 \(\mathcal N_\Lambda\) converge uniformly to \(\mathcal N\) on the
 common ball;
 \item \(\sup_\Lambda\|K_{\kappa,\Lambda}^{\mathrm{hyb}}\|\le\eta<1\);
 \item the nonlinear remainder is Lipschitz with constant \(L_N\) satisfying
 \begin{equation}
  \frac{
  C_Ee^{C\int_{\mathcal C}V}L_N
  }{1-\eta}<1.
 \label{eq:v2v35-nonlinear}
 \end{equation}
\end{enumerate}
Then the regulated nonlinear moving-band solutions exist uniquely on the
common slab and converge in \(X_T\) to the unique reduced continuum
solution.
\end{theorem}

\begin{proof}
The first item gives a uniform hybrid Einstein bound.  The second and third
give linear cutoff convergence by
\cref{thm:v2v32-regulator}.  The fourth makes the nonlinear map a uniform
contraction as in \cref{thm:v2v32-nonlinear}.  For two regulators, subtract
their fixed-point equations; the contraction term is absorbed on the left,
while the operator, source, and nonlinear-remainder differences tend to zero uniformly on the common ball.  Hence the fixed
points form a Cauchy family and converge to the unique limiting fixed point.
\end{proof}

\subsection{What remains to verify}
\label{sec:v2v35-gate}

\begin{firewall}[The hybrid theorem does not prove its SM hypotheses]
\label{fw:v2v35-hypotheses}
V35 proves that crossings do not create a combinatorial multiplicity when
one global buffer is retained.  It does not prove:
\begin{enumerate}
 \item a regulator-uniform bound on the crossing occupation;
 \item the causal polarization norm \eqref{eq:v2v35-eta};
 \item convergence of the interacting SM polarization in operator norm;
 \item the nonlinear estimate \eqref{eq:v2v35-nonlinear};
 \item continuation beyond loss of hyperbolicity or the low-energy margin.
\end{enumerate}
These are analytic inputs, not consequences of switching frames.
\end{firewall}

\begin{programme}[Next acquisition]
\label{prog:v2v35}
The next calculation should estimate the crossing occupation from eigenvalue
motion.  A usable result must bound
\[
 |\mathcal C|
 =
 \left|
 \left\{
 t:\operatorname{dist}(\kappa^2,\operatorname{spec}A(t))<g_0
 \right\}
 \right|
\]
or its \(V(t)\)-weighted version by spectral-flow data without counting the
same eigenvalue repeatedly.  Tangential crossings and eigenvalue clusters
must be treated explicitly.
\end{programme}

\subsection{V35 conclusion and next step}
\label{sec:v2v35-conclusion}

The mandatory audit records the unchanged NS V31 note and the YM V24
complete incidence-two shared-payer theorem.  Following its one-carrier
assembly rule, V35 keeps one smooth spectral buffer throughout the slab and
uses Kato transport only as a coordinate frame on gapped components.
Projector motion is charged by
\(\int_{\mathcal C}V\), not by the number of crossings or a singular inverse
gap.  Under a uniform polarization contraction, this yields causal
order-reduction, nonlinear existence, and regulator convergence across the
crossings.

The construction remains conditional on a crossing-occupation estimate and
the interacting SM operator norms.  V36 will derive a coarea bound for
transversal eigenvalue crossings and isolate the exact additional datum
needed for tangential or clustered spectral flow.
\clearpage
\section*{Second Volume, Version V36: Coarea Control of Spectral Crossing Occupation}
\addcontentsline{toc}{section}{Second Volume, Version V36: Coarea Control of Spectral Crossing Occupation}

\subsection*{Purpose}

V35 reduced the frame-motion cost to the weighted occupation
\(\int_{\mathcal C}V\), where the crossing strip is the set of times at
which the moving elliptic spectrum lies within \(g_0\) of \(\kappa^2\).
This note proves an exact one-dimensional coarea identity for that
occupation.  It gives a useful bound under quantitative transversality,
shows why endpoint spectral flow and upper speed estimates are
insufficient, and states finite-type substitutes for tangential contacts
and eigenvalue clusters.  The result is conditional: the required
transversality or finite-type datum is not inferred from the Standard
Model--Einstein equations here.

\subsection{Active branches and counted occupation}
\label{sec:v2v36-branches}

Set \(a=\kappa^2\) and choose an auxiliary spectral window \(G>g>0\).
Work on a compact time interval \(I=[0,T]\).  Assume that the part of the
spectrum of \(A(t)\) meeting \((a-G,a+G)\) is represented by finitely many
Lipschitz eigenvalue
branches
\[
 \lambda_1,\ldots,\lambda_J:I\longrightarrow\mathbb R.
\]
This is a local spectral regularity hypothesis.  It holds, for example,
for a norm-resolvent \(C^1\) compact-resolvent family after an isolated
finite spectral window has been selected.  Define
\begin{align}
 I_j(g)&=\{t\in I:|\lambda_j(t)-a|<g\},
 \label{eq:v2v36-Ij}\\
 \mathcal C_g&=\bigcup_{j=1}^J I_j(g),
 \label{eq:v2v36-Cg}\\
 \Omega_V(g)&=\sum_{j=1}^J\int_{I_j(g)}V(t)\,\mathrm dt,
 \label{eq:v2v36-counted}
\end{align}
where \(V\ge0\) is measurable and integrable.  Simultaneous branches are
counted separately in \(\Omega_V\), so
\begin{equation}
 \int_{\mathcal C_g}V(t)\,\mathrm dt\le \Omega_V(g).
\label{eq:v2v36-union}
\end{equation}

For almost every regular level \(y\), put
\begin{equation}
 \Theta_V(y)
 =
 \sum_{j=1}^J
 \sum_{\substack{t\in I_j(g):\lambda_j(t)=y\\
                         \dot\lambda_j(t)\ne0}}
 \frac{V(t)}{|\dot\lambda_j(t)|},
 \qquad |y-a|<g,
\label{eq:v2v36-density}
\end{equation}
and define the critical occupation
\begin{equation}
 \Omega_V^{\mathrm{crit}}(g)
 =
 \sum_{j=1}^J
 \int_{I_j(g)\cap\{\dot\lambda_j=0\}}V(t)\,\mathrm dt.
\label{eq:v2v36-critical}
\end{equation}
Derivatives and critical sets are understood up to null sets.

\begin{theorem}[Exact coarea decomposition]
\label{thm:v2v36-coarea}
Under the preceding hypotheses,
\begin{equation}
 \Omega_V(g)
 =
 \Omega_V^{\mathrm{crit}}(g)
 +\int_{a-g}^{a+g}\Theta_V(y)\,\mathrm dy.
\label{eq:v2v36-coarea}
\end{equation}
Consequently, if
\begin{equation}
 \Omega_V^{\mathrm{crit}}(g)=0,
 \qquad
 \mathop{\rm ess\,sup}_{|y-a|<g}\Theta_V(y)\le \Theta_* ,
\label{eq:v2v36-transverse-data}
\end{equation}
then
\begin{equation}
 \int_{\mathcal C_g}V(t)\,\mathrm dt
 \le 2g\Theta_*.
\label{eq:v2v36-main-bound}
\end{equation}
This bound uses the level-set density itself and does not introduce a
factor equal to the number of components of \(\mathcal C_g\).
\end{theorem}

\begin{proof}
Apply the one-dimensional area formula to each Lipschitz map
\(t\mapsto\lambda_j(t)\), on the measurable set
\(I_j(g)\cap\{\dot\lambda_j\ne0\}\), with integrand
\(V/|\dot\lambda_j|\).  It gives
\[
 \int_{I_j(g)\cap\{\dot\lambda_j\ne0\}}V(t)\,\mathrm dt
 =
 \int_{a-g}^{a+g}
 \sum_{\substack{t\in I_j(g):\lambda_j(t)=y\\
                         \dot\lambda_j(t)\ne0}}
 \frac{V(t)}{|\dot\lambda_j(t)|}\,\mathrm dy.
\]
Add the integral over the critical set and sum over \(j\).  This proves
\eqref{eq:v2v36-coarea}.  Combine it with
\eqref{eq:v2v36-union} and integrate the essential-supremum bound over an
interval of length \(2g\) to obtain \eqref{eq:v2v36-main-bound}.
\end{proof}

\begin{remark}[Relabelling invariance]
\label{rem:v2v36-relabelling}
The separate branch terms depend on a labelling through eigenvalue
collisions, but their sum in \eqref{eq:v2v36-coarea} does not: it is the
area formula on the union of the eigenvalue graphs, counted with spectral
multiplicity.  Thus a permutation of absolutely continuous branch labels
does not change \(\Omega_V\), \(\Omega_V^{\mathrm{crit}}\), or
\(\Theta_V\) almost everywhere.
\end{remark}

\subsection{Concrete transverse criteria}
\label{sec:v2v36-transverse}

\begin{corollary}[Multiplicity and slope criterion]
\label{cor:v2v36-slope}
Suppose that, for almost every \(y\in(a-g,a+g)\), branch \(j\) has at most
\(m_j\) regular preimages in \(I_j(g)\), and that
\begin{equation}
 V(t)\le \gamma_j|\dot\lambda_j(t)|
 \quad\hbox{for almost every }t\in I_j(g).
\label{eq:v2v36-weighted-slope}
\end{equation}
Then
\begin{equation}
 \int_{\mathcal C_g}V(t)\,\mathrm dt
 \le 2g\sum_{j=1}^Jm_j\gamma_j.
\label{eq:v2v36-multiplicity-bound}
\end{equation}
In particular, if every active branch is monotone in the strip, one may
take \(m_j=1\).  For unweighted occupation, the hypotheses
\(|\dot\lambda_j|\ge\nu_j>0\) give
\begin{equation}
 |\mathcal C_g|
 \le 2g\sum_{j=1}^J\frac{m_j}{\nu_j}.
\label{eq:v2v36-time-bound}
\end{equation}
\end{corollary}

\begin{proof}
Condition \eqref{eq:v2v36-weighted-slope} makes the critical occupation
zero and bounds each regular summand in \eqref{eq:v2v36-density} by
\(\gamma_j\).  Hence
\(\Theta_V\le\sum_jm_j\gamma_j\).  Apply
\cref{thm:v2v36-coarea}.  The unweighted statement is the case
\(V=1\) and \(\gamma_j=\nu_j^{-1}\).
\end{proof}

\begin{corollary}[Scale-adapted transversality]
\label{cor:v2v36-scaled}
If
\begin{equation}
 |\dot\lambda_j(t)|
 \ge \vartheta_j\kappa^2 V(t)
 \quad\hbox{on }I_j(g),
 \qquad \vartheta_j>0,
\label{eq:v2v36-scaled}
\end{equation}
and the regular level multiplicities are bounded by \(m_j\), then
\begin{equation}
 \int_{\mathcal C_g}V(t)\,\mathrm dt
 \le
 \frac{2g}{\kappa^2}
 \sum_{j=1}^J\frac{m_j}{\vartheta_j}.
\label{eq:v2v36-scaled-bound}
\end{equation}
Thus a strip of relative width \(g/\kappa^2\) has a correspondingly small
frame-motion charge, provided the relative transverse speed is uniform.
\end{corollary}

\begin{proof}
Use \(\gamma_j=(\vartheta_j\kappa^2)^{-1}\) in
\cref{cor:v2v36-slope}.
\end{proof}

\subsection{What the evolution of the operator supplies}
\label{sec:v2v36-Hadamard}

\begin{proposition}[Feynman--Hellmann datum]
\label{prop:v2v36-FH}
Let \(A(t)\) be differentiable as a self-adjoint compact-resolvent family
with common form domain.  At a time where \(\lambda_j(t)\) is simple, with
normalized eigenvector \(\phi_j(t)\),
\begin{equation}
 \dot\lambda_j(t)
 =
 \langle\phi_j(t),\dot A(t)\phi_j(t)\rangle.
\label{eq:v2v36-FH}
\end{equation}
At an eigenvalue of multiplicity \(q\), the first derivatives of the
branches through the cluster are the eigenvalues of the Hermitian form
\begin{equation}
 P_*\dot A(t)P_*:\operatorname{Ran}P_*\longrightarrow
 \operatorname{Ran}P_*,
\label{eq:v2v36-cluster-derivative}
\end{equation}
whenever those branch derivatives exist.
\end{proposition}

\begin{proof}
Differentiate \(A\phi_j=\lambda_j\phi_j\), pair with \(\phi_j\), and use
self-adjointness to cancel the two terms containing \(\dot\phi_j\).  At a
multiple eigenvalue, restrict the differentiated eigenvalue equation to
the eigenspace.  Diagonalizing the resulting finite-dimensional Hermitian
form gives the possible first derivatives.
\end{proof}

The V33 coefficient estimate bounds the absolute value of
\eqref{eq:v2v36-FH} from above by a multiple of \(\kappa^2V\) on the
active window.  The lower bound \eqref{eq:v2v36-scaled} is different: it
requires a non-cancellation condition for the matrix element of
\(\dot A\).  An upper operator-speed estimate cannot imply it.

\begin{proposition}[Endpoint flow and upper speed do not control occupation]
\label{prop:v2v36-obstruction}
There is no bound for \(\int_{\mathcal C_g}V\) in terms only of the
endpoint spectral flow through \(a\), an upper bound for \(\|\dot A\|\),
and \(g\), uniformly in \(T\).
\end{proposition}

\begin{proof}
On \(\mathbb C^2\), for arbitrary \(T>0\), take
\[
 A_T(t)=
 \begin{pmatrix}
 a+g/2&0\\0&a+3g+t
 \end{pmatrix},
 \qquad 0\le t\le T,
 \qquad V(t)=1.
\]
The first eigenvalue stays inside the strip for the entire interval, so
\(\mathcal C_g=[0,T]\) and \(\int_{\mathcal C_g}V=T\).  No eigenvalue
crosses \(a\), while \(\|\dot A_T\|=1\) independently of \(T\).
Therefore the proposed data stay fixed as the occupation becomes
arbitrarily large.
\end{proof}

\begin{firewall}[Transversality is an additional dynamical input]
\label{fw:v2v36-transversality}
The Feynman--Hellmann formula identifies the quantity that must be
controlled, but does not give its sign or a positive lower bound.  Gauge
symmetry, geometric symmetry, or cancellation in
\(\langle\phi,\dot A\phi\rangle\) may pin a mode in the crossing strip.
Such a mode contributes to \(\Omega_V^{\mathrm{crit}}\) and is invisible
to the regular coarea density.
\end{firewall}

\subsection{Finite-type tangencies and clusters}
\label{sec:v2v36-finite-type}

\begin{proposition}[Finite-type branch occupation]
\label{prop:v2v36-finite-type}
Suppose, for branch \(j\), that there are finitely many centres
\(z_{j\ell}\), neighbourhoods \(U_{j\ell}\), constants
\(c_{j\ell}>0\), and integers \(m_{j\ell}\ge1\) such that
\begin{align}
 I_j(g)&\subset\bigcup_\ell U_{j\ell},
 \label{eq:v2v36-cover}\\
 |\lambda_j(t)-a|
 &\ge c_{j\ell}|t-z_{j\ell}|^{m_{j\ell}}
 \quad(t\in U_{j\ell}).
\label{eq:v2v36-finite-type}
\end{align}
Then
\begin{equation}
 |I_j(g)|
 \le
 2\sum_\ell
 \left(\frac{g}{c_{j\ell}}\right)^{1/m_{j\ell}}.
\label{eq:v2v36-finite-type-bound}
\end{equation}
If \(V\in L^\infty(I)\), the same right side multiplied by
\(\|V\|_{L^\infty}\) bounds \(\int_{I_j(g)}V\).
\end{proposition}

\begin{proof}
For \(t\in I_j(g)\cap U_{j\ell}\),
\eqref{eq:v2v36-finite-type} implies
\(|t-z_{j\ell}|<(g/c_{j\ell})^{1/m_{j\ell}}\).  Cover
\(I_j(g)\) by these intervals and sum their lengths.  The weighted result
follows from the \(L^\infty\) bound.
\end{proof}

Linear transversality is the case \(m=1\).  A quadratic tangency costs
\(O(g^{1/2})\), and a contact of order \(m\) costs \(O(g^{1/m})\).  Thus
tangencies need not be excluded, but their order and leading coefficient
must be controlled uniformly.

\begin{proposition}[Basis-independent cluster criterion]
\label{prop:v2v36-cluster}
Let an isolated rank-\(q\) spectral cluster have eigenvalues
\(\mu_1(t),\ldots,\mu_q(t)\) satisfying
\(|\mu_i(t)-a|\le R\).  On its spectral bundle define the invariant
cluster determinant
\begin{equation}
 D(t)=\det\bigl(A(t)|_{\operatorname{Ran}P_{\rm cl}(t)}-aI\bigr)
 =\prod_{i=1}^q(\mu_i(t)-a).
\label{eq:v2v36-determinant}
\end{equation}
Suppose the cluster crossing set is covered by finitely many
neighbourhoods \(U_\ell\) on which
\begin{equation}
 |D(t)|\ge c_\ell|t-z_\ell|^{m_\ell}.
\label{eq:v2v36-det-type}
\end{equation}
Then
\begin{equation}
 \left|
 \left\{t:\min_i|\mu_i(t)-a|<g\right\}
 \right|
 \le
 2\sum_\ell
 \left(\frac{gR^{q-1}}{c_\ell}\right)^{1/m_\ell}.
\label{eq:v2v36-cluster-bound}
\end{equation}
The corresponding weighted bound follows after multiplication by
\(\|V\|_{L^\infty}\).
\end{proposition}

\begin{proof}
If one cluster eigenvalue is within \(g\) of \(a\), then
\[
 |D(t)|=\prod_{i=1}^q|\mu_i(t)-a|<gR^{q-1}.
\]
Together with \eqref{eq:v2v36-det-type}, this places \(t\) within
\((gR^{q-1}/c_\ell)^{1/m_\ell}\) of the appropriate centre.  Sum the
lengths of the covering intervals.  The determinant is unchanged by a
unitary change of frame on the cluster, so the criterion does not require
a branch choice at a degeneracy.
\end{proof}

\subsection{Insertion into the hybrid Einstein estimate}
\label{sec:v2v36-insertion}

\begin{corollary}[Quantitative crossing continuation]
\label{cor:v2v36-continuation}
Under the hypotheses of \cref{thm:v2v35-energy}, suppose either the
coarea data \eqref{eq:v2v36-transverse-data} hold uniformly, or the
finite-type bounds of \cref{prop:v2v36-finite-type,prop:v2v36-cluster}
give a uniform number \(B_\kappa\) such that
\begin{equation}
 \int_{\mathcal C_g}V(t)\,\mathrm dt\le B_\kappa.
\label{eq:v2v36-Bk}
\end{equation}
Then the frame factor in the V35 energy, retarded evolution, reduction
error, and nonlinear contraction bounds is at most
\begin{equation}
 \exp(CB_\kappa).
\label{eq:v2v36-frame-factor}
\end{equation}
If the spectral data and \(B_\kappa\) are regulator-uniform, crossing
occupation causes no loss in the regulator limit.
\end{corollary}

\begin{proof}
Take \(g=g_0\) from \eqref{eq:v2v35-gap-set} and choose the auxiliary
window \(G>g_0\).  Then \(\mathcal C_g=\mathcal C\).
Insert \eqref{eq:v2v36-Bk} into
\eqref{eq:v2v35-Gronwall}, \eqref{eq:v2v35-E},
\eqref{eq:v2v35-error}, and \eqref{eq:v2v35-nonlinear}.  Each occurrence
of the crossing occupation is then bounded by the same scalar
\(B_\kappa\), independently of the number of switching intervals.
\end{proof}

\begin{firewall}[What V36 does and does not close]
\label{fw:v2v36-status}
V36 closes the crossing-occupation step only after a transversal density or
finite-type cluster bound has been supplied.  For the full
SM--Einstein continuum one must still prove such a bound from the evolving
geometric operator, uniformly in the regulator, and must also establish
the interacting causal polarization and nonlinear estimates listed in
\cref{fw:v2v35-hypotheses}.  The finite-dimensional obstruction in
\cref{prop:v2v36-obstruction} rules out obtaining the missing lower bound
from operator regularity alone.
\end{firewall}

\subsection{V36 conclusion and next step}
\label{sec:v2v36-conclusion}

The crossing charge has an exact regular part and a singular part.  The
regular part is the coarea integral of \(V/|\dot\lambda|\) over spectral
levels; the singular part is occupation carried by stationary branches.
Transverse branches give an \(O(g)\) bound, finite-order tangencies give
\(O(g^{1/m})\), and an isolated cluster can be handled without choosing
eigenvectors by a finite-type bound on its determinant.

V37 will go upstream to the geometric variation of the actual Laplace-type
operators on an Einstein foliation.  It will derive the Hadamard form of
\(\langle\phi,\dot A\phi\rangle\), separate lapse, shift, and metric
deformation terms, and determine which sign or finite-type condition could
yield the V36 spectral datum without assuming genericity by declaration.
\clearpage
\section*{Second Volume, Version V37: Geometric Hadamard Variation and a Concrete Transversality Criterion}
\addcontentsline{toc}{section}{Second Volume, Version V37: Geometric Hadamard Variation and a Concrete Transversality Criterion}

\subsection*{Purpose and relation to earlier notes}

V36 identified the regular crossing density
\(V/|\dot\lambda|\) and showed that no upper bound for \(\dot A\) can
replace a lower bound on the motion of active eigenvalues.  This note
computes that motion for the Laplace-type operator used in V33.  Here
``Hadamard variation'' means the shape derivative of an eigenvalue.  It is
distinct from the Hadamard subtraction of quantum two-point functions in
V26 and from the boundary microlocal analysis in the first volume.

The calculation is first made for a scalar Schrödinger operator on the
closed spatial slice assumed in V33.  It is then extended to a
gauge-covariant Laplace-type operator.  The ADM shift cancels exactly by
naturality.  A near-homothetic metric deformation gives a sufficient
transversality condition, while the remaining Standard Model terms enter
an explicit matrix-element error budget.

\subsection{Variation of the scalar quadratic form}
\label{sec:v2v37-scalar}

Let \((\Sigma,h(t))\) be a closed \(d\)-dimensional Riemannian manifold,
with \(h\) of class \(C^1\) in time, and set
\begin{equation}
 S_{ij}=\frac12\dot h_{ij},
 \qquad \operatorname{tr}S=h^{ij}S_{ij}.
\label{eq:v2v37-S}
\end{equation}
Consider
\begin{equation}
 A(t)=-\Delta_{h(t)}+U(t),
\label{eq:v2v37-scalar-A}
\end{equation}
where \(U\) is real.  The identity term and the scalar part of
\(\mathcal R\) in \eqref{eq:v2v33-A} may be included in \(U\).

\begin{theorem}[Scalar Hadamard formula]
\label{thm:v2v37-scalar-Hadamard}
Suppose \(A(t)\phi(t)=\lambda(t)\phi(t)\), the eigenvalue branch is
differentiable at the time in question, and
\(\int_\Sigma|\phi|^2\,\mathrm d\mu_h=1\).  Then
\begin{align}
 \dot\lambda
 =\int_\Sigma\Bigl\{
 &-2S^{ij}\operatorname{Re}
   (\nabla_i\overline\phi\,\nabla_j\phi)
 +(\operatorname{tr}S)
   \bigl(|\nabla\phi|^2+(U-\lambda)|\phi|^2\bigr)
 +\dot U|\phi|^2
 \Bigr\}\,\mathrm d\mu_h.
\label{eq:v2v37-scalar-Hadamard}
\end{align}
The same formula gives every differentiable branch through a multiple
eigenvalue after \(\phi\) is chosen to diagonalize the compressed
variation on that eigenspace.
\end{theorem}

\begin{proof}
Let
\begin{align*}
 q_t(u)&=\int_\Sigma
 \bigl(|\nabla u|_h^2+U|u|^2\bigr)\,\mathrm d\mu_h,\\
 n_t(u)&=\int_\Sigma|u|^2\,\mathrm d\mu_h.
\end{align*}
For an eigenfunction, stationarity of the Rayleigh quotient gives
\(\dot\lambda=\dot q_t(\phi)-\lambda\dot n_t(\phi)\); the terms containing
\(\dot\phi\) cancel by the eigenvalue equation.  Since
\begin{equation}
 \dot h^{ij}=-2S^{ij},
 \qquad
 \partial_t\mathrm d\mu_h=(\operatorname{tr}S)\,\mathrm d\mu_h,
\label{eq:v2v37-volume-variation}
\end{equation}
differentiating the two forms gives exactly
\eqref{eq:v2v37-scalar-Hadamard}.  On a multiple eigenspace this quadratic
form is the compressed variation from
\cref{prop:v2v36-FH}; diagonalizing it gives the branch derivatives.
\end{proof}

Define the spectral stress of a normalized eigenfunction by
\begin{equation}
 \mathbb T_{ij}[\phi]
 =-2\operatorname{Re}(\nabla_i\overline\phi\,\nabla_j\phi)
 +h_{ij}\bigl(|\nabla\phi|^2+(U-\lambda)|\phi|^2\bigr).
\label{eq:v2v37-spectral-stress}
\end{equation}
Then the formula takes the compact form
\begin{equation}
 \dot\lambda
 =\int_\Sigma
 \left(S^{ij}\mathbb T_{ij}[\phi]+\dot U|\phi|^2\right)
 \,\mathrm d\mu_h.
\label{eq:v2v37-stress-form}
\end{equation}

\subsection{ADM lapse and exact cancellation of the shift}
\label{sec:v2v37-ADM}

Fix the sign convention
\begin{equation}
 \dot h_{ij}=-2N K_{ij}+(\mathcal L_\beta h)_{ij}.
\label{eq:v2v37-ADM}
\end{equation}
For a scalar potential define its material derivative
\begin{equation}
 \mathring U=(\partial_t-\mathcal L_\beta)U.
\label{eq:v2v37-material-U}
\end{equation}

\begin{theorem}[Shift-free physical variation]
\label{thm:v2v37-shift-free}
For the natural scalar operator \eqref{eq:v2v37-scalar-A},
\begin{align}
 \dot\lambda
 =\int_\Sigma\Bigl\{
 &2NK^{ij}\operatorname{Re}
   (\nabla_i\overline\phi\,\nabla_j\phi)
 -N(\operatorname{tr}K)
   \bigl(|\nabla\phi|^2+(U-\lambda)|\phi|^2\bigr)
 +\mathring U|\phi|^2
 \Bigr\}\,\mathrm d\mu_h.
\label{eq:v2v37-shift-free}
\end{align}
There is no term involving \(\beta\).  In particular, a pure shift changes
the coordinate representation of the eigenspace but not its eigenvalues.
\end{theorem}

\begin{proof}
Split \(S=-NK+\tfrac12\mathcal L_\beta h\) and
\(\dot U=\mathring U+\mathcal L_\beta U\) in
\eqref{eq:v2v37-scalar-Hadamard}.  The terms containing \(\beta\) are the
first variation along the pullback by the flow of \(\beta\).  Naturality
gives
\[
 A_{\Phi_s^*h,\Phi_s^*U}
 =\Phi_s^*A_{h,U}(\Phi_s^{-1})^*,
\]
so this variation is a commutator and has zero diagonal matrix element on
an eigenfunction.  The remaining terms are obtained by substituting
\(S=-NK\), which proves \eqref{eq:v2v37-shift-free}.
\end{proof}

\begin{remark}[Where lapse gradients can enter]
\label{rem:v2v37-lapse-gradient}
The intrinsic operator \(-\Delta_h+U\) depends on the lapse only through
the physical velocity \(-2NK\) and through the material evolution of
\(U\).  Spatial derivatives of \(N\) appear if one instead builds lapse
weights into the elliptic operator, or when spacetime wave operators are
split into space and time.  They are absent from
\eqref{eq:v2v37-shift-free} because the V33 cutoff operator is intrinsic to
each slice.
\end{remark}

\subsection{A gauge-covariant Laplace-type formula}
\label{sec:v2v37-bundle}

Let \(E\to\Sigma\) be a Hermitian bundle, identified in time by a unitary
trivialization, and consider
\begin{equation}
 A=\nabla^*\nabla+\mathcal E,
\label{eq:v2v37-bundle-A}
\end{equation}
where \(\nabla\) is a unitary connection and \(\mathcal E\) a self-adjoint
endomorphism.  Write
\begin{equation}
 B_i=\partial_t\nabla_i,
 \qquad F=\partial_t\mathcal E
\label{eq:v2v37-BF}
\end{equation}
in that trivialization.

\begin{theorem}[Bundle Hadamard formula]
\label{thm:v2v37-bundle-Hadamard}
If \(A\psi=\lambda\psi\) and \(\|\psi\|_{L^2(h)}=1\), then at every
differentiability time of the branch,
\begin{align}
 \dot\lambda
 =\int_\Sigma\Bigl\{
 &-2S^{ij}\operatorname{Re}
   \langle\nabla_i\psi,\nabla_j\psi\rangle
 +(\operatorname{tr}S)
   \bigl(|\nabla\psi|^2+
   \langle(\mathcal E-\lambda)\psi,\psi\rangle\bigr)
 \notag\\
 &+2\operatorname{Re}
   h^{ij}\langle B_i\psi,\nabla_j\psi\rangle
 +\langle F\psi,\psi\rangle
 \Bigr\}\,\mathrm d\mu_h.
\label{eq:v2v37-bundle-Hadamard}
\end{align}
\end{theorem}

\begin{proof}
Differentiate
\[
 q_t(\psi)=\int_\Sigma
 \left(h^{ij}\langle\nabla_i\psi,\nabla_j\psi\rangle
 +\langle\mathcal E\psi,\psi\rangle\right)\,\mathrm d\mu_h.
\]
The inverse-metric and density terms are the first line of
\eqref{eq:v2v37-bundle-Hadamard}.  Differentiating the two connection
factors gives twice the real part of the \(B_i\) term, and differentiating
the endomorphism gives the \(F\) term.  As in the scalar proof, the
eigenvalue equation cancels all terms containing \(\dot\psi\), including
the derivative of its normalization.
\end{proof}

For a covariantly constructed bundle operator, replace ordinary time
derivatives in \eqref{eq:v2v37-BF} by the material derivatives obtained
after subtracting the natural lifted Lie derivative along \(\beta\).
Exactly as in \cref{thm:v2v37-shift-free}, the shift part is a unitary
conjugation and drops out.  The remaining terms are the physical metric
deformation \(-NK\), the material gauge-connection velocity, and the
material endomorphism velocity.  Higgs, curvature, mass, and gauge-fixing
contributions to the V33 operator occur in the latter two rows.

\begin{corollary}[Explicit lower-bound ledger]
\label{cor:v2v37-ledger}
For an active normalized bundle eigenvector, write
\begin{equation}
 \dot\lambda=G_{\rm met}+G_{\rm conn}+G_{\rm end},
\label{eq:v2v37-ledger}
\end{equation}
where the three quantities are respectively the physical metric, material
connection, and material endomorphism integrals in
\eqref{eq:v2v37-bundle-Hadamard}.  If one term, say \(G_{\rm met}\), obeys
\begin{equation}
 |G_{\rm met}|
 \ge \vartheta\kappa^2V+|G_{\rm conn}|+|G_{\rm end}|,
 \qquad \vartheta>0,
\label{eq:v2v37-dominance}
\end{equation}
through the crossing strip, then the V36 scale-adapted transversality
condition \eqref{eq:v2v36-scaled} holds.
\end{corollary}

\begin{proof}
The reverse triangle inequality in \eqref{eq:v2v37-ledger} gives
\(|\dot\lambda|\ge\vartheta\kappa^2V\).
\end{proof}

Condition \eqref{eq:v2v37-dominance} is a checkable matrix-element
criterion.  It also displays a possible failure: connection and
endomorphism motion may cancel the geometric contribution even when all
three terms separately have the expected size.

\subsection{A concrete near-homothetic criterion}
\label{sec:v2v37-homothetic}

The scalar Laplacian already supplies a nontrivial positive case.  Put
\(U=0\), decompose
\begin{equation}
 K=Hh+\widehat K,
 \qquad H=\frac1d\operatorname{tr}K,
 \qquad \operatorname{tr}\widehat K=0,
\label{eq:v2v37-K-decomposition}
\end{equation}
and set
\begin{equation}
 q=NH,
 \qquad W=N\widehat K.
\label{eq:v2v37-qW}
\end{equation}

\begin{theorem}[Near-homothetic spectral transversality]
\label{thm:v2v37-homothetic}
Let \(d\ge2\), let \(-\Delta_h\phi=\lambda\phi\) with
\(\lambda>0\), and suppose there is a scalar \(q_0(t)\), constant on
\(\Sigma\) at each time, such that
\begin{equation}
 \|q-q_0\|_{L^\infty(\Sigma)}\le\varepsilon,
 \qquad
 \|W\|_{L^\infty(\Sigma;\mathrm{op})}\le\delta.
\label{eq:v2v37-near-homothetic}
\end{equation}
Then
\begin{equation}
 |\dot\lambda-2q_0\lambda|
 \le\bigl((2d-2)\varepsilon+2\delta\bigr)\lambda.
\label{eq:v2v37-homothetic-error}
\end{equation}
In dimension three, if
\begin{equation}
 b(t):=2|q_0(t)|-4\varepsilon(t)-2\delta(t)>0,
\label{eq:v2v37-b}
\end{equation}
then every differentiable positive eigenvalue branch satisfies
\begin{equation}
 |\dot\lambda|\ge b(t)\lambda,
 \qquad
 \operatorname{sign}\dot\lambda=operatorname{sign}q_0.
\label{eq:v2v37-homothetic-lower}
\end{equation}
\end{theorem}

\begin{proof}
Insert \(NK=qh+W\) and \(U=0\) into
\eqref{eq:v2v37-shift-free}.  Since
\(\int|\nabla\phi|^2=\lambda\) and \(\int|\phi|^2=1\),
\begin{align*}
 \dot\lambda
 ={}&(2-d)\int_\Sigma q|\nabla\phi|^2\,\mathrm d\mu_h
 +d\lambda\int_\Sigma q|\phi|^2\,\mathrm d\mu_h\\
 &+2\int_\Sigma W^{ij}
 \operatorname{Re}(\nabla_i\overline\phi\,\nabla_j\phi)
 \,\mathrm d\mu_h.
\end{align*}
Replacing \(q\) by the spatial constant \(q_0\) gives
\(2q_0\lambda\).  The two errors involving \(q-q_0\) are bounded by
\(((d-2)+d)\varepsilon\lambda\), and the trace-free term by
\(2\delta\lambda\).  This proves
\eqref{eq:v2v37-homothetic-error}.  For \(d=3\), the strict inequality in
\eqref{eq:v2v37-b} and the reverse triangle inequality give both claims in
\eqref{eq:v2v37-homothetic-lower}.
\end{proof}

\begin{corollary}[Occupation bound on a near-homothetic slab]
\label{cor:v2v37-occupation}
Let \(a=\kappa^2\), \(0<g<a\), and suppose the hypotheses of
\cref{thm:v2v37-homothetic} hold in dimension three for all active scalar
branches.  If
\begin{equation}
 b(t)\ge\theta V(t),
 \qquad \theta>0,
\label{eq:v2v37-bV}
\end{equation}
then, for \(J\) active branches,
\begin{equation}
 \int_{\mathcal C_g}V(t)\,\mathrm dt
 \le \frac{2gJ}{\theta(a-g)}.
\label{eq:v2v37-occupation}
\end{equation}
If \(q_0\) is continuous, \eqref{eq:v2v37-b} prevents it from changing
sign, so every active branch is monotone and no additional level
multiplicity factor is needed.
\end{corollary}

\begin{proof}
On the strip, \(\lambda\ge a-g\), so
\eqref{eq:v2v37-homothetic-lower,eq:v2v37-bV} give
\[
 |\dot\lambda|\ge\theta(a-g)V.
\]
Continuity and \eqref{eq:v2v37-b} make the sign common and fixed; hence
each branch has at most one preimage of a regular level.  Apply
\cref{cor:v2v36-slope} with
\(m_j=1\) and \(\gamma_j=[\theta(a-g)]^{-1}\).
\end{proof}

\begin{assessment}[Meaning of the positive case]
\label{ass:v2v37-positive-case}
Uniform contraction or expansion of a slice moves every nonzero scalar
Laplacian eigenvalue coherently.  Small spatial variation of \(NH\) and a
small shear preserve that direction and produce the V36 occupation bound.
This is a genuine populated class of geometries; it is not a claim that a
general Einstein--Standard Model evolution remains in that class.
\end{assessment}

\subsection{Limits of the criterion}
\label{sec:v2v37-limits}

\begin{firewall}[Einstein evolution does not force spectral transversality]
\label{fw:v2v37-Einstein}
The constraint and evolution equations do not impose
\eqref{eq:v2v37-b} on an arbitrary foliation.  A static slice has
\(K=0\), a maximal slice may have \(H=0\), and anisotropic shear can cancel
the trace contribution on a particular eigenfunction.  In the bundle
sectors, gauge-connection and endomorphism variations add further
possible cancellations.  Therefore V37 proves a sufficient geometric
criterion and an exact diagnostic formula, not a regulator-uniform SM
transversality theorem.
\end{firewall}

\begin{firewall}[Boundary and zero-mode qualifications]
\label{fw:v2v37-boundary}
The derivation uses the closed-slice setup of V33.  If \(\Sigma\) has a
boundary, varying domains and boundary conditions add Hadamard boundary
terms and must be combined with the domain analysis of the first volume.
Zero modes do not satisfy the positive-eigenvalue estimate in
\cref{thm:v2v37-homothetic}; gauge and constraint zero modes must be
removed or treated in their reduced complex before applying a spectral
crossing estimate at positive \(\kappa^2\).
\end{firewall}

\subsection{V37 conclusion and next step}
\label{sec:v2v37-conclusion}

The eigenvalue velocity is the pairing of the physical ADM deformation
with a spectral stress, plus material connection and endomorphism terms.
The shift cancels exactly.  Near spatially uniform nonzero expansion or
contraction, the metric term dominates and yields the V36 coarea datum
with an explicit error margin.  Static, maximal, anisotropic, and
matter-cancelling configurations remain outside that theorem.

V38 will address this failure upstream by replacing pointwise
transversality with an integrated alternative.  The target is a finite-type
criterion obtained from one or more time derivatives of the compressed
cluster operator, together with a quantitative sublevel estimate that
remains valid when \(\dot\lambda\) vanishes at isolated times.
\clearpage
\section*{Second Volume, Version V38: Higher-Order Spectral Contacts and Uniformly Integrable Crossing Weights}
\addcontentsline{toc}{section}{Second Volume, Version V38: Higher-Order Spectral Contacts and Uniformly Integrable Crossing Weights}

\subsection*{Purpose}

V37 gave a first-derivative criterion for transverse spectral motion.  It
necessarily fails at an isolated tangency, even when the eigenvalue leaves
the crossing strip rapidly enough for the V35 energy bound.  This note
replaces first-order transversality by a finite-order derivative condition.
The scalar estimate is an elementary one-dimensional sublevel theorem.  An
isolated spectral cluster is first transported to a fixed finite-dimensional
space and then treated through its determinant, avoiding denominators from
gaps inside the cluster.

There is a second issue.  A small unweighted crossing set need not have
small \(V\)-weighted occupation for a regulator-dependent family.  The
exact additional hypothesis is uniform integrability of the coefficient
weights.  This note records that condition and inserts the resulting bound
into the V35 hybrid continuation theorem.

\subsection{A higher-derivative sublevel lemma}
\label{sec:v2v38-sublevel}

\begin{lemma}[One-dimensional finite-type sublevel bound]
\label{lem:v2v38-sublevel}
Let \(f\in C^m(J)\) on an interval \(J\), where \(m\ge1\), and suppose
\begin{equation}
 |f^{(m)}(t)|\ge c>0
 \qquad(t\in J).
\label{eq:v2v38-mth-lower}
\end{equation}
Then, for every \(\epsilon>0\),
\begin{equation}
 \left|\left\{t\in J:|f(t)|<\epsilon\right\}\right|
 \le 2m\left(\frac{\epsilon}{c}\right)^{1/m}.
\label{eq:v2v38-sublevel}
\end{equation}
The constant is independent of the length of \(J\).
\end{lemma}

\begin{proof}
Let \(E=\{t\in J:|f(t)|<\epsilon\}\) and \(L=|E|\).  It is enough to
prove the result for compact subsets of \(E\) and pass to their increasing
union.  An elementary quantile selection gives points
\(x_0<\cdots<x_m\) in such a compact subset with
\begin{equation}
 x_j-x_i\ge\frac{j-i}{m}L
 \qquad(0\le i<j\le m),
\label{eq:v2v38-separated-points}
\end{equation}
up to replacing \(L\) by an arbitrarily smaller number.  Indeed, choose
successive points after each additional portion of set measure \(L/m\);
an interval containing that portion must have at least the same length.

The divided-difference formula and \(|f(x_i)|<\epsilon\) give
\begin{align}
 |[x_0,\ldots,x_m]f|
 &\le
 \epsilon\sum_{i=0}^m
 \frac1{\prod_{j\ne i}|x_i-x_j|}\\
 &\le
 \frac{\epsilon m^m}{L^m}
 \sum_{i=0}^m\frac1{i!(m-i)!}
 =\frac{\epsilon(2m)^m}{m!L^m}.
\label{eq:v2v38-divided}
\end{align}
By the generalized mean-value theorem, the left side equals
\(|f^{(m)}(\xi)|/m!\) for some \(\xi\) between the selected points.
Use \eqref{eq:v2v38-mth-lower} and rearrange to obtain
\eqref{eq:v2v38-sublevel}.
\end{proof}

\begin{corollary}[Finite-order eigenvalue contact]
\label{cor:v2v38-branch}
Let \(\lambda_j\in C^{m_j}(J_j)\) and suppose
\begin{equation}
 |\lambda_j^{(m_j)}(t)|\ge c_j>0
 \qquad(t\in J_j).
\label{eq:v2v38-branch-derivative}
\end{equation}
Then
\begin{equation}
 \left|\left\{t\in J_j:
 |\lambda_j(t)-\kappa^2|<g\right\}\right|
 \le 2m_j\left(\frac g{c_j}\right)^{1/m_j}.
\label{eq:v2v38-branch-bound}
\end{equation}
If the full crossing set is covered by finitely many such branch intervals,
the sum of their right sides bounds its measure.
\end{corollary}

\begin{proof}
Apply \cref{lem:v2v38-sublevel} to
\(f=\lambda_j-\kappa^2\), then use the union bound.
\end{proof}

This statement permits \(\dot\lambda_j=0\).  A quadratic contact is handled
with \(m_j=2\), and a contact of order \(m_j\) has occupation
\(O(g^{1/m_j})\).  What must remain uniform is the order, the derivative
floor, and the number of intervals in the finite cover.

\subsection{Why repeated simple-eigenvalue differentiation is unstable}
\label{sec:v2v38-simple}

\begin{proposition}[Second derivative of a simple eigenvalue]
\label{prop:v2v38-second}
Let \(A(t)\) be a \(C^2\) self-adjoint family on a fixed Hilbert space,
and let \(\lambda_j\) be simple at \(t\), with orthonormal eigenbasis
\((\phi_k)\).  Then
\begin{equation}
 \ddot\lambda_j
 =\langle\phi_j,\ddot A\phi_j\rangle
 +2\sum_{k\ne j}
 \frac{|\langle\phi_k,\dot A\phi_j\rangle|^2}
 {\lambda_j-\lambda_k}.
\label{eq:v2v38-second}
\end{equation}
\end{proposition}

\begin{proof}
Choose the phase so that
\(\langle\phi_j,\dot\phi_j\rangle=0\).  Projecting the differentiated
eigenvalue equation onto \(\phi_k\), \(k\ne j\), gives
\[
 \langle\phi_k,\dot\phi_j\rangle
 =\frac{\langle\phi_k,\dot A\phi_j\rangle}
 {\lambda_j-\lambda_k}.
\]
Differentiate the Feynman--Hellmann identity
\(\dot\lambda_j=\langle\phi_j,\dot A\phi_j\rangle\), insert this expansion,
and take the real part.  The two eigenvector-derivative terms coincide and
give \eqref{eq:v2v38-second}.
\end{proof}

Formula \eqref{eq:v2v38-second} is useful only while the selected
eigenvalue remains separated from every other eigenvalue.  It develops
inverse internal gaps precisely where a cluster forms.  The next
construction removes that artificial singularity.

\subsection{Kato compression of an isolated cluster}
\label{sec:v2v38-cluster}

Let \(A(t)\) be a \(C^M\) self-adjoint compact-resolvent family.  Suppose a
rank-\(q\) spectral cluster is separated from the rest of the spectrum by
an external gap \(\delta_{\rm ext}>0\) on an interval \(J\).  Let
\(P_{\rm cl}(t)\) be its Riesz projection and let \(U(t)\) be the Kato
transport from a fixed reference fibre \(E_0=\operatorname{Ran}P_{\rm cl}(t_0)\).
Define the finite Hermitian matrix
\begin{equation}
 B(t)=U(t)^*A(t)U(t)|_{E_0}
\label{eq:v2v38-B}
\end{equation}
and, for \(a=\kappa^2\),
\begin{equation}
 D_a(t)=\det(B(t)-aI).
\label{eq:v2v38-D}
\end{equation}

\begin{lemma}[Regularity and frame invariance]
\label{lem:v2v38-cluster-regularity}
Under the preceding hypotheses, \(B\) and \(D_a\) are \(C^M\) on \(J\).
The scalar \(D_a\) is independent of the unitary frame chosen on the
cluster.  Its zeros are exactly the times at which the cluster spectrum
contains \(a\), counted with algebraic multiplicity.
\end{lemma}

\begin{proof}
The external Riesz contour and the resolvent derivative identity show that
\(P_{\rm cl}\) is \(C^M\).  Its Kato ordinary differential equation then
gives a \(C^M\) unitary transport, so \(B\) and its determinant are
\(C^M\).  A different unitary frame conjugates \(B\), leaving its
determinant and eigenvalues unchanged.  Finally,
\(D_a=\prod_{i=1}^q(\mu_i-a)\) for the cluster eigenvalues \(\mu_i\).
\end{proof}

\begin{theorem}[Higher-derivative cluster occupation]
\label{thm:v2v38-cluster}
Assume \(|\mu_i(t)-a|\le R\) on \(J\), and for some
\(1\le m\le M\),
\begin{equation}
 |D_a^{(m)}(t)|\ge c>0
 \qquad(t\in J).
\label{eq:v2v38-det-derivative}
\end{equation}
Then
\begin{equation}
 \left|\left\{t\in J:
 \min_{1\le i\le q}|\mu_i(t)-a|<g\right\}\right|
 \le
 2m\left(\frac{gR^{q-1}}c\right)^{1/m}.
\label{eq:v2v38-cluster-bound}
\end{equation}
No gap between eigenvalues inside the cluster occurs in this estimate.
\end{theorem}

\begin{proof}
If one cluster eigenvalue lies within \(g\) of \(a\), then
\[
 |D_a(t)|=\prod_{i=1}^q|\mu_i(t)-a|<gR^{q-1}.
\]
Apply \cref{lem:v2v38-sublevel} to \(D_a\) with
\(\epsilon=gR^{q-1}\).  Only the external gap was used to construct the
smooth compression; the proof never divides by an internal gap.
\end{proof}

\begin{proposition}[Populated matrix finite type]
\label{prop:v2v38-matrix-type}
Suppose at \(t_0\)
\begin{equation}
 B(t)=aI+(t-t_0)^r C+O(|t-t_0|^{r+1}),
 \qquad \det C\ne0,
\label{eq:v2v38-matrix-jet}
\end{equation}
in operator norm on the \(q\)-dimensional cluster.  Then
\begin{equation}
 D_a(t)
 =(t-t_0)^{rq}\det C+O(|t-t_0|^{rq+1}),
\label{eq:v2v38-det-jet}
\end{equation}
so
\begin{equation}
 D_a^{(rq)}(t_0)=(rq)!\det C\ne0.
\label{eq:v2v38-det-nonzero}
\end{equation}
After shrinking the interval around \(t_0\),
\cref{thm:v2v38-cluster} applies with \(m=rq\).
\end{proposition}

\begin{proof}
Factor \((t-t_0)^r\) from each of the \(q\) columns of
\(B(t)-aI\).  The remaining determinant is
\(\det(C+O(|t-t_0|))=\det C+O(|t-t_0|)\), which proves
\eqref{eq:v2v38-det-jet}.  Differentiation at \(t_0\) gives
\eqref{eq:v2v38-det-nonzero}; continuity supplies a positive lower bound
for \(|D_a^{(rq)}|\) on a smaller interval.
\end{proof}

Thus a \(q\)-fold cluster crossing with an invertible first compressed
velocity has determinant type \(q\), while an order-\(r\) matrix tangency
has determinant type \(rq\).  This loss is real: the determinant is the
product of all \(q\) spectral distances.

\subsection{From small crossing time to weighted occupation}
\label{sec:v2v38-uniform-integrability}

For a family of nonnegative weights \((V_\Lambda)\) on \([0,T]\), define
its uniform-integrability modulus
\begin{equation}
 \omega_V(\delta)
 =\sup_\Lambda
 \sup_{\substack{E\subset[0,T]\ \mathrm{measurable}\\|E|\le\delta}}
 \int_E V_\Lambda(t)\,\mathrm dt.
\label{eq:v2v38-UI-modulus}
\end{equation}
The family is uniformly integrable exactly when
\(\omega_V(\delta)\to0\) as \(\delta\downarrow0\), together with a finite
uniform \(L^1\) bound.

\begin{theorem}[Uniform weighted crossing estimate]
\label{thm:v2v38-weighted}
Suppose the active scalar branches and isolated clusters admit a finite
cover by intervals satisfying
\cref{cor:v2v38-branch,thm:v2v38-cluster}, with constants uniform in the
regulator.  Let \(L_\kappa(g)\) be the sum of the corresponding right-hand
sides.  Then
\begin{equation}
 |\mathcal C_{g,\Lambda}|\le L_\kappa(g)
\label{eq:v2v38-L}
\end{equation}
and, for uniformly integrable \((V_\Lambda)\),
\begin{equation}
 \sup_\Lambda
 \int_{\mathcal C_{g,\Lambda}}V_\Lambda(t)\,\mathrm dt
 \le\omega_V(L_\kappa(g)).
\label{eq:v2v38-weighted}
\end{equation}
In particular, the weighted occupation tends to zero with \(g\) whenever
\(L_\kappa(g)\to0\).
\end{theorem}

\begin{proof}
The scalar and cluster sublevel estimates, followed by the union bound,
give \eqref{eq:v2v38-L}.  The definition
\eqref{eq:v2v38-UI-modulus} applied to
\(E=\mathcal C_{g,\Lambda}\) gives
\eqref{eq:v2v38-weighted}.  Uniform integrability and
\(L_\kappa(g)\to0\) give the last claim.
\end{proof}

\begin{corollary}[An \(L^p\) sufficient condition]
\label{cor:v2v38-Lp}
If, for some \(p>1\),
\(\sup_\Lambda\|V_\Lambda\|_{L^p}\le M_p\), then
\begin{equation}
 \sup_\Lambda
 \int_{\mathcal C_{g,\Lambda}}V_\Lambda
 \le M_p L_\kappa(g)^{1-1/p}.
\label{eq:v2v38-Lp}
\end{equation}
For \(p=\infty\), the exponent is one.
\end{corollary}

\begin{proof}
Apply Holder's inequality on the crossing set and then use
\eqref{eq:v2v38-L}.
\end{proof}

\begin{proposition}[A uniform \(L^1\) bound is insufficient]
\label{prop:v2v38-L1-failure}
There are sets \(C_n\) with \(|C_n|\to0\) and nonnegative weights \(V_n\)
with \(\|V_n\|_{L^1}=1\) such that
\begin{equation}
 \int_{C_n}V_n=1
\label{eq:v2v38-L1-failure}
\end{equation}
for every \(n\).
\end{proposition}

\begin{proof}
On \([0,1]\), take \(C_n=[0,n^{-1}]\) and
\(V_n=n\mathbf1_{C_n}\).  Then \(|C_n|=n^{-1}\) and both stated
integrals equal one.
\end{proof}

\begin{firewall}[Time concentration is a separate regulator defect]
\label{fw:v2v38-UI}
Finite-type spectral motion controls the measure of the crossing set.  It
does not prevent the coefficient velocity from concentrating on that set.
For regulator-uniform V35 bounds, uniform integrability of
\((V_\Lambda)\), or a stronger \(L^p\) estimate, is an independent compactness
input.  A uniform \(L^1\) bound alone does not supply it.
\end{firewall}

\subsection{Insertion into the SM--Einstein continuation ledger}
\label{sec:v2v38-insertion}

\begin{corollary}[Finite-type hybrid continuation]
\label{cor:v2v38-continuation}
Assume the V35 hyperbolicity and physical energy hypotheses.  At
\(g=g_0\), suppose the finite-order branch and cluster data are uniform in
the regulator, and suppose \((V_\Lambda)\) is uniformly integrable.  Then
the crossing factor in the regulated hybrid evolution is bounded by
\begin{equation}
 \exp\!\left(C\omega_V(L_\kappa(g_0))\right).
\label{eq:v2v38-hybrid-factor}
\end{equation}
Consequently, if the polarization and nonlinear hypotheses of
\cref{thm:v2v35-nonlinear} also hold, the regulated solutions converge on
the common slab through all finite-type crossings covered above.
\end{corollary}

\begin{proof}
Use \eqref{eq:v2v38-weighted} as the number \(B_\kappa\) in
\cref{cor:v2v36-continuation}, then invoke
\cref{thm:v2v35-nonlinear}.
\end{proof}

\begin{firewall}[Remaining analytic inputs]
\label{fw:v2v38-status}
V38 does not prove that the compressed SM cluster determinants have
uniform finite type.  It shows exactly what is sufficient and removes
spurious inverse gaps within a cluster.  External cluster separation,
uniform bounds on the relevant time jets of \(A\), uniform integrability
of the geometric coefficient velocity, the causal SM polarization norm,
and the nonlinear contraction estimate remain to be obtained from the
coupled dynamics.
\end{firewall}

\subsection{V38 conclusion and next step}
\label{sec:v2v38-conclusion}

An isolated zero of \(\dot\lambda\) is no longer an obstruction.  A
uniform nonzero derivative of some finite order gives an explicit sublevel
bound.  For clusters, the determinant of the Kato-compressed operator
provides a smooth, basis-independent scalar and avoids internal spectral
denominators.  Uniform integrability is the exact condition that converts
the resulting small crossing time into a regulator-uniform weighted
occupation.

V39 will examine the remaining external-gap assumption.  It will build a
nested spectral-window decomposition and determine whether overlapping
clusters can be merged with a bounded-rank cost, or whether Weyl density
forces the programme back to the smooth V33 functional calculus at high
cutoff.
\clearpage
\section*{Second Volume, Version V39: Local Cluster Covers and Regulator-Stable Finite Type}
\addcontentsline{toc}{section}{Second Volume, Version V39: Local Cluster Covers and Regulator-Stable Finite Type}

\subsection*{Purpose}

V38 treated a spectral cluster only when it was separated from the rest of
the spectrum on a whole interval.  A global external gap is too strong for
a moving high-frequency spectrum.  This note proves that the argument only
needs local gaps near exact resonance times.  Compact resolvent isolates a
finite cluster at each such time, compactness gives a finite time cover,
and a sufficiently thin crossing strip lies inside that cover.

The second result transfers this local construction to ultraviolet
regulators under uniform differentiated norm-resolvent convergence.  The
rank of every fixed-energy cluster then stabilizes, and a nonzero finite
jet of the continuum cluster determinant remains nonzero for all large
regulators.  This supplies a rigorous route to the uniform finite-type
input of V38 at fixed physical \(\kappa\).  It does not produce estimates
uniform as \(\kappa\to\infty\).

\subsection{Continuity of the distance to a fixed energy}
\label{sec:v2v39-distance}

Let \(A(t)\), \(t\in I=[0,T]\), be self-adjoint, lower semibounded, and
compact-resolvent on a fixed Hilbert space.  Fix \(c>0\) so that
\(A(t)+c\ge I\), and assume
\begin{equation}
 t\longmapsto (A(t)+c)^{-1}
\quad\hbox{is norm continuous}.
\label{eq:v2v39-norm-resolvent}
\end{equation}
Let \(a=\kappa^2\) and define
\begin{equation}
 d_a(t)=\operatorname{dist}(a,\operatorname{spec}A(t)),
 \qquad
 Z_a=\{t\in I:d_a(t)=0\}.
\label{eq:v2v39-Za}
\end{equation}

\begin{lemma}[Continuity of the local spectral distance]
\label{lem:v2v39-distance}
Under \eqref{eq:v2v39-norm-resolvent}, \(d_a\) is continuous.  Hence
\(Z_a\) is compact.
\end{lemma}

\begin{proof}
The spectral mapping theorem sends an eigenvalue \(\lambda\) of \(A(t)\)
to \((\lambda+c)^{-1}\) for the compact self-adjoint resolvent.  Norm
continuity of compact self-adjoint operators gives continuity, with
multiplicity, of every finite group of nonzero resolvent eigenvalues.
Only eigenvalues of \(A(t)\) in a fixed bounded neighbourhood of \(a\) can
realize \(d_a(t)\); this group corresponds to resolvent eigenvalues bounded
away from zero.  Its distance from \((a+c)^{-1}\), and therefore
\(d_a(t)\), varies continuously.  The zero set of a continuous function on
compact \(I\) is compact.
\end{proof}

\subsection{A finite local cluster cover}
\label{sec:v2v39-cover}

\begin{theorem}[Local isolation near all exact resonances]
\label{thm:v2v39-cover}
Assume \eqref{eq:v2v39-norm-resolvent}.  For every \(t_*\in Z_a\), there
are a time neighbourhood \(U_*\), a bounded spectral interval
\(J_*=(a-r_*,a+r_*)\), and a contour \(\Gamma_*\) such that:
\begin{enumerate}
 \item \(a\) belongs to the cluster inside \(J_*\) at \(t_*\);
 \item \(\Gamma_*\subset\rho(A(t))\) for every \(t\in U_*\);
 \item the Riesz projection
 \begin{equation}
  P_*(t)=\frac1{2\pi i}\int_{\Gamma_*}(z-A(t))^{-1}\,\mathrm dz
 \label{eq:v2v39-Riesz}
 \end{equation}
 has a finite constant rank on \(U_*\).
\end{enumerate}
Moreover, there are finitely many such neighbourhoods
\(U_1,\ldots,U_L\) and a number \(g_*>0\) such that
\begin{equation}
 \mathcal C_g
 =\{t:d_a(t)<g\}
 \subset\bigcup_{\ell=1}^L U_\ell
 \qquad(0<g<g_*).
\label{eq:v2v39-cover}
\end{equation}
Thus a single global external gap is unnecessary.
\end{theorem}

\begin{proof}
At \(t_*\), compact resolvent makes \(a\) an isolated eigenvalue of finite
multiplicity.  Choose \(r_*>0\) so that the endpoints of \(J_*\) and a
small enclosing contour \(\Gamma_*\) avoid the spectrum of \(A(t_*)\).
Norm-resolvent continuity and the resolvent identity keep
\(\Gamma_*\) in the resolvent set after the time neighbourhood is shrunk.
The Riesz projections then vary in operator norm.  Projections at distance
less than one have the same rank, so the rank is finite and constant on
\(U_*\).

The sets \(U_*\) cover compact \(Z_a\); retain a finite subcover
\(U_1,\ldots,U_L\).  Its union \(U\) is an open neighbourhood of \(Z_a\).
On the compact complement \(I\setminus U\), the continuous positive
function \(d_a\) has a positive minimum, say \(2g_*\).  If \(g<g_*\), no
point of the crossing strip lies in that complement, proving
\eqref{eq:v2v39-cover}.
\end{proof}

\begin{remark}[Overlapping clusters are merged locally]
\label{rem:v2v39-merge}
If two eigenvalue groups collide near the same resonance time, choose
\(J_*\) and \(\Gamma_*\) around their union.  Internal crossings do not
affect the rank or smoothness of the Riesz projection.  Only separation of
the merged finite cluster from the remaining spectrum is used.  Different
members of the finite time cover may use different ranks and contours.
\end{remark}

\begin{corollary}[Local determinants suffice]
\label{cor:v2v39-local-det}
Suppose \(A\) is \(C^M\) in norm-resolvent sense.  On every \(U_\ell\),
Kato-transport \(P_\ell(t)\) to a fixed fibre and form the determinant
\(D_{a,\ell}\) of \eqref{eq:v2v38-D}.  If each crossing portion admits a
finite refinement on which
\begin{equation}
 |D_{a,\ell}^{(m_{\ell r})}(t)|\ge c_{\ell r}>0,
 \qquad 1\le m_{\ell r}\le M,
\label{eq:v2v39-local-type}
\end{equation}
then the sum of the V38 cluster bounds over this finite refinement controls
\(|\mathcal C_g|\) for every sufficiently small \(g\).
\end{corollary}

\begin{proof}
Apply \cref{thm:v2v38-cluster} on every refined interval and use the finite
union bound.  By \cref{thm:v2v39-cover}, no part of \(\mathcal C_g\) lies
outside the selected neighbourhoods when \(g<g_*\).
\end{proof}

\subsection{Transfer to ultraviolet regulators}
\label{sec:v2v39-regulator}

Let \(A_\Lambda(t)\) be self-adjoint compact-resolvent regulators on the
same Hilbert space.  A discretization on varying spaces requires the same
hypotheses after insertion of specified quasi-unitary identification maps;
that extension is not implicit here.  Assume
\begin{equation}
 \sup_{t\in I}
 \left\|(A_\Lambda(t)+c)^{-1}-(A(t)+c)^{-1}\right\|
 \longrightarrow0.
\label{eq:v2v39-uniform-NR}
\end{equation}

\begin{theorem}[Uniform transfer of local clusters]
\label{thm:v2v39-transfer}
For each of the finitely many continuum contours \(\Gamma_\ell\) in
\cref{thm:v2v39-cover}, there is \(\Lambda_0\) such that, for
\(\Lambda\ge\Lambda_0\),
\begin{enumerate}
 \item \(\Gamma_\ell\subset\rho(A_\Lambda(t))\) uniformly for
       \(t\in U_\ell\);
 \item the regulated Riesz projections \(P_{\ell,\Lambda}(t)\) satisfy
 \begin{equation}
  \sup_{t\in U_\ell}
  \|P_{\ell,\Lambda}(t)-P_\ell(t)\|\longrightarrow0;
 \label{eq:v2v39-P-convergence}
 \end{equation}
 \item \(\operatorname{rank}P_{\ell,\Lambda}
       =\operatorname{rank}P_\ell\).
 \item for some \(g_{**}>0\), every sufficiently large regulator obeys
 \begin{equation}
  \mathcal C_{g,\Lambda}\subset\bigcup_{\ell=1}^L U_\ell
  \qquad(0<g<g_{**}).
 \label{eq:v2v39-reg-cover}
 \end{equation}
\end{enumerate}
In particular, no regulator-dependent cluster multiplicity appears at a
fixed physical energy.
\end{theorem}

\begin{proof}
Map the contour by \(z\mapsto(z+c)^{-1}\).  Its compact image has positive
distance from the spectrum of the continuum compact resolvents, uniformly
on the closure of a slightly smaller \(U_\ell\).  Uniform norm convergence
in \eqref{eq:v2v39-uniform-NR} preserves this resolvent separation for all
large \(\Lambda\).  The contour resolvent identity then gives uniform
convergence of the Riesz integrals, proving
\eqref{eq:v2v39-P-convergence}.  Once the norm difference is less than one,
the two projections have isomorphic ranges and equal finite rank.
On the compact complement of the time cover, the continuum spectrum has a
positive distance from \(a\).  Uniform norm-resolvent convergence preserves
half of this local spectral exclusion for all large \(\Lambda\), which
proves \eqref{eq:v2v39-reg-cover}.
\end{proof}

To transfer finite type, assume differentiated convergence.  In a Kato
frame on each local cluster, let \(B_{\ell,\Lambda}\) and
\(B_\ell\) be the regulated and continuum matrices.

\begin{theorem}[Persistence of determinant finite type]
\label{thm:v2v39-jet-transfer}
Suppose, after compatible Kato transport,
\begin{equation}
 B_{\ell,\Lambda}\longrightarrow B_\ell
 \quad\hbox{in }C^M(U_\ell;\operatorname{Herm}(q_\ell)).
\label{eq:v2v39-CM}
\end{equation}
If on a refined interval \(U_{\ell r}\)
\begin{equation}
 |D_{a,\ell}^{(m)}(t)|\ge c>0,
 \qquad m\le M,
\label{eq:v2v39-limit-type}
\end{equation}
then, for all sufficiently large \(\Lambda\),
\begin{equation}
 |D_{a,\ell,\Lambda}^{(m)}(t)|\ge c/2
 \qquad(t\in U_{\ell r}).
\label{eq:v2v39-reg-type}
\end{equation}
Consequently, the V38 sublevel constants are uniform in the regulator.
\end{theorem}

\begin{proof}
The determinant and each of its derivatives through order \(M\) are
polynomials in the matrix entries and their derivatives.  Therefore
\eqref{eq:v2v39-CM} implies
\(D_{a,\ell,\Lambda}\to D_{a,\ell}\) in \(C^M\).  For large
\(\Lambda\), the uniform difference of the two \(m\)-th derivatives is
less than \(c/2\).  The reverse triangle inequality gives
\eqref{eq:v2v39-reg-type}, and
\cref{thm:v2v38-cluster} gives the final statement.
\end{proof}

\begin{proposition}[A sufficient operator convergence hypothesis]
\label{prop:v2v39-operator-jets}
Suppose the resolvents and their time derivatives through order \(M\)
converge uniformly on each contour \(\Gamma_\ell\), and the corresponding
compressed operators converge there through order \(M\).  Then Kato
transports may be chosen so that \eqref{eq:v2v39-CM} holds.
\end{proposition}

\begin{proof}
Differentiating the Riesz integral gives convergence of
\(P_{\ell,\Lambda}\) to \(P_\ell\) in \(C^M\).  The Kato generators
\([\dot P,P]\) therefore converge in \(C^{M-1}\).  Continuous dependence
for their finite-rank unitary ordinary differential equations gives
convergence of the transports.  Combining this with the assumed compressed
operator convergence yields \eqref{eq:v2v39-CM}.
\end{proof}

\subsection{Rank bounds at fixed physical energy}
\label{sec:v2v39-rank}

\begin{proposition}[Uniform fixed-energy rank bound]
\label{prop:v2v39-rank}
For a uniformly elliptic Laplace-type family on a compact
\(d\)-manifold, assume the uniform counting estimate
\begin{equation}
 N_{A(t)}(L)
 :=\operatorname{rank}\mathbf1_{(-\infty,L]}(A(t))
 \le C(1+L^{d/2}).
\label{eq:v2v39-Weyl-upper}
\end{equation}
Then every local cluster contained below \(a+R\) has
\begin{equation}
 q_\ell\le C(1+(a+R)^{d/2}).
\label{eq:v2v39-rank}
\end{equation}
If the regulators obey the same estimate, the bound is uniform in
\(\Lambda\).
\end{proposition}

\begin{proof}
The cluster range is a subspace of the spectral subspace below \(a+R\),
so its rank cannot exceed the counting function there.  Apply
\eqref{eq:v2v39-Weyl-upper} to the continuum and regulated operators.
\end{proof}

At fixed \(a=\kappa^2\), this is a finite number independent of the
ultraviolet cutoff.  In three dimensions the displayed coarse bound grows
like \(O(\kappa^3)\).  A local Weyl estimate can improve the count in a
thin shell, but neither estimate gives a fixed rank as
\(\kappa\to\infty\).  Moreover, the determinant type may grow with
\(q_\ell\), as \cref{prop:v2v38-matrix-type} already shows.

\begin{firewall}[Fixed physical cutoff and ultraviolet cutoff are distinct]
\label{fw:v2v39-two-cutoffs}
The regulator limit \(\Lambda\to\infty\) at fixed physical
\(\kappa\) can have stable local cluster ranks under
\cref{thm:v2v39-transfer}.  Sending \(\kappa\to\infty\) is a different
limit.  Weyl density then enlarges the active finite-dimensional space and
generally shrinks available external gaps.  V39 does not turn the hard-band
construction into a cutoff-uniform high-energy argument; the smooth V33
functional calculus remains necessary there.
\end{firewall}

\subsection{Gap-free assembly of the crossing estimate}
\label{sec:v2v39-assembly}

\begin{theorem}[Regulator-stable local finite-type continuation]
\label{thm:v2v39-continuation}
Assume:
\begin{enumerate}
 \item the continuum operator satisfies the local finite-type cover
       \eqref{eq:v2v39-local-type};
 \item the regulators satisfy the convergence hypotheses of
       \cref{thm:v2v39-transfer,thm:v2v39-jet-transfer};
 \item the coefficient velocities \((V_\Lambda)\) are uniformly
       integrable as in \eqref{eq:v2v38-UI-modulus}.
\end{enumerate}
Then, for sufficiently small fixed crossing width \(g\), there is a number
\(B_\kappa(g)<\infty\), independent of \(\Lambda\), such that
\begin{equation}
 \int_{\mathcal C_{g,\Lambda}}V_\Lambda(t)\,\mathrm dt
 \le B_\kappa(g),
 \qquad
 B_\kappa(g)\longrightarrow0\quad(g\downarrow0).
\label{eq:v2v39-B}
\end{equation}
Hence the crossing factor in the V35 energy theorem is bounded by
\(\exp(CB_\kappa(g))\) without a global external spectral gap.
\end{theorem}

\begin{proof}
The finite local cover comes from \cref{thm:v2v39-cover}.  Determinant
finite type persists uniformly by \cref{thm:v2v39-jet-transfer}; summing
the V38 sublevel bounds gives a regulator-independent
\(L_\kappa(g)\to0\).  Apply \cref{thm:v2v38-weighted} and set
\(B_\kappa(g)=\omega_V(L_\kappa(g))\).  The last statement follows from
\cref{thm:v2v35-energy}.
\end{proof}

\begin{firewall}[Spectral convergence remains an SM input]
\label{fw:v2v39-spectral-convergence}
Weak convergence of metrics or coefficients does not by itself give the
differentiated norm-resolvent convergence used above.  That convergence
must be proved from the regulator construction, with common domains or
explicit identification maps and without spectral pollution.  Likewise,
finite type of every continuum local determinant is a dynamical statement,
not a consequence of compact resolvent.
\end{firewall}

\subsection{V39 conclusion and next step}
\label{sec:v2v39-conclusion}

The external-gap hypothesis can be localized and no single gapped cluster
need persist through the entire slab.  Exact resonance times have a finite
cover by isolated merged clusters, a thin crossing strip lies inside that
cover, and differentiated norm-resolvent convergence transfers ranks and
finite determinant jets to the regulators.  This closes the logical gap
between the V38 finite-type theorem and a fixed-energy regulator limit,
subject to strong spectral convergence and uniform integrability.

V40 is the next mandatory companion audit.  After that audit it will turn
to the strongest remaining low-energy bottleneck: proving the causal
polarization contraction in \eqref{eq:v2v35-eta}, beginning with a
gauge-reduced retarded energy estimate and an explicit separation of local
renormalization terms from the genuinely nonlocal response.
\clearpage
\section*{Second Volume, Version V40: Companion Audit and a Causal Volterra Polarization Theorem}
\addcontentsline{toc}{section}{Second Volume, Version V40: Companion Audit and a Causal Volterra Polarization Theorem}

\subsection*{Purpose and mandatory companion audit}

This is the required audit five versions after V35.  The current companion
sources inspected were:
\begin{center}
\begin{tabular}{p{0.46\textwidth}r p{0.33\textwidth}}
\toprule
file & bytes & SHA--256\\
\midrule
\texttt{Renewal\_NS\_Stability\_Research\_Notes\_V31.tex}
&887537&
\texttt{F1121B62238AD5491BB7CF03635EA9934942EDF573ED8FAAA9EE79E1D38A6696}\\
\texttt{Renewal\_Yang\_Mills\_Mass\_Gap\_Research\_Notes\_V33.tex}
&408679&
\texttt{24EF2AF77D30F0BBA1796FAD54D9E62655C20A5F2F678ADB40E86052CD05163C}\\
\bottomrule
\end{tabular}
\end{center}
The NS note remains at V31.  It proves an exact dyadic first-moment
formation ladder and shows that a weaker Duhamel source norm restores the
missing shell weight; its open gate is a correlation on the actual
recent-entry set, and persistent intervals must not be charged again.

The current named YM note is V33, but it is byte-identical to V32, with the
same digest displayed above.  The new mathematical endpoint is therefore
V32: the complete pure incidence-three, unique-junction \(SU(2)\) sector is
closed.  Its proof first assembles the complete local payer, then inserts
balanced or unbalanced shared moments one-sidedly.  Mixed incidence two and
three, higher incidence, general groups, the full continuum, and the mass
gap remain open.  Both YM files were left untouched.

\begin{assessment}[Transfer from the V40 audit]
\label{ass:v2v40-audit}
Two rules transfer to the causal SM polarization problem.
\begin{enumerate}
 \item A retarded response must be estimated in the source norm matched to
       its Duhamel map, with pre-existing history retained as one known
       source.  Recharging the same history on every continuation slab is
       the analogue of the forbidden NS double payment.
 \item Local renormalization terms must be assembled into the complete
       effective Einstein operator before the nonlocal response is
       inserted.  Multiplying separate sector bounds would reproduce the
       overlapping-capacity error avoided in the YM proof.
\end{enumerate}
These rules point to a Volterra estimate for the genuinely nonlocal
retarded polarization, rather than a single undifferentiated norm for all
local and nonlocal terms.
\end{assessment}

\subsection{Gauge-reduced causal setup}
\label{sec:v2v40-gauge}

After the V35 hybrid spectral frame identifies the moving band with a
fixed Banach space \(H^s_{\rm red}\), let
\begin{equation}
 X_T=L^\infty([0,T];H^s_{\rm red}).
\label{eq:v2v40-X}
\end{equation}
The suffix ``red'' denotes a fixed gauge-reduced and constraint-compatible
space.  Let \(Q_X\) and \(Q_Y\) be bounded projections onto the physical
solution and source spaces.

\begin{proposition}[Constraint-compatible retarded reduction]
\label{prop:v2v40-gauge}
Let \(L_{\rm eff}\) be a gauge-fixed hyperbolic operator with retarded
solution map \(E_{\rm eff}:Y_T\to X_T\).  Suppose
\begin{equation}
 Q_XE_{\rm eff}=E_{\rm eff}Q_Y,
 \qquad
 Q_Y\Pi^{\rm ret}_{\rm nl}=\Pi^{\rm ret}_{\rm nl}Q_X.
\label{eq:v2v40-Ward}
\end{equation}
If the source is physical and the subsidiary constraint system has zero
initial data, every iterate of
\begin{equation}
 K=E_{\rm eff}\Pi^{\rm ret}_{\rm nl}
\label{eq:v2v40-K}
\end{equation}
lies in \(\operatorname{Ran}Q_X\).  Thus the causal resolvent may be
constructed entirely on the reduced space.
\end{proposition}

\begin{proof}
For \(h\in\operatorname{Ran}Q_X\), the second identity in
\eqref{eq:v2v40-Ward} places \(\Pi^{\rm ret}_{\rm nl}h\) in
\(\operatorname{Ran}Q_Y\), and the first places its retarded solution in
\(\operatorname{Ran}Q_X\).  Induction gives the claim for all iterates.
The homogeneous subsidiary equation preserves the zero constraints.
\end{proof}

The identities \eqref{eq:v2v40-Ward} are the operator form of the Ward and
constraint-propagation requirements.  They must hold after regularization
and renormalization; gauge fixing alone does not prove them.

\subsection{Separate local renormalization before estimating memory}
\label{sec:v2v40-local}

Write the renormalized linear response schematically as
\begin{equation}
 \Pi^{\rm ret}
 =\Pi_{\rm abs}+R_{\rm loc}+\Pi^{\rm ret}_{\rm nl}.
\label{eq:v2v40-split}
\end{equation}
Here \(\Pi_{\rm abs}\) contains the local terms included in the chosen
renormalized gravitational couplings and in the hyperbolic principal
operator.  Higher-derivative local terms are first causally order-reduced
on the low band; the remaining pointwise-in-time operator is denoted
\(R_{\rm loc}(t)\).  Define
\begin{equation}
 L_{\rm eff}=L_0-\Pi_{\rm abs}.
\label{eq:v2v40-Leff}
\end{equation}

\begin{lemma}[One-sided local absorption]
\label{lem:v2v40-local}
Suppose the retarded equation after applying \(E_{\rm eff}\) is
\begin{equation}
 h=h_0+K_{\rm loc}h+K_{\rm nl}h,
 \qquad
 (K_{\rm loc}h)(t)=K_{\rm loc}(t)h(t),
\label{eq:v2v40-local-equation}
\end{equation}
and
\begin{equation}
 \mathop{\rm ess\,sup}_{0\le t\le T}
 \|K_{\rm loc}(t)\|\le\eta_{\rm loc}<1.
\label{eq:v2v40-local-small}
\end{equation}
Then
\begin{equation}
 J(t)=(I-K_{\rm loc}(t))^{-1},
 \qquad
 \|J(t)\|\le(1-\eta_{\rm loc})^{-1},
\label{eq:v2v40-J}
\end{equation}
and the equation is equivalently
\begin{equation}
 h=Jh_0+JK_{\rm nl}h.
\label{eq:v2v40-one-sided}
\end{equation}
If \(K_{\rm nl}\) has a causal kernel bounded by a scalar kernel \(b\), the
new kernel is bounded by \(b/(1-\eta_{\rm loc})\).
\end{lemma}

\begin{proof}
The pointwise Neumann series gives \eqref{eq:v2v40-J}.  Move the local term
to the left of \eqref{eq:v2v40-local-equation} and multiply once by
\(J(t)\).  The kernel statement follows from the operator-norm bound for
\(J(t)\).
\end{proof}

This is the appropriate place for the YM audit lesson: the complete local
block is inverted once, and its norm is inserted on one side of the
nonlocal kernel.  No product of overlapping sector capacities is taken.

\begin{firewall}[Local fourth-order terms cannot be silently absorbed]
\label{fw:v2v40-local}
Curvature-squared counterterms vary to fourth-order differential
operators.  They may be included in \(L_{\rm eff}\) only with a proved
well-posed causal prescription, or reduced on the low band with a
controlled error before entering \(R_{\rm loc}\).  Treating them as an
ordinary second-order change of the Einstein operator would hide the
higher-derivative modes that V32 was designed to control.
\end{firewall}

\subsection{The fractional Volterra estimate}
\label{sec:v2v40-Volterra}

The decisive observation is that a genuinely retarded response is
triangular in time.  Its Neumann series can converge even when its global
operator norm exceeds one.

\begin{theorem}[Fractional Volterra resolvent]
\label{thm:v2v40-Volterra}
Let \(K:X_T\to X_T\) be linear and causal.  Suppose, for some
\(0<\rho\le1\) and \(C_K\ge0\),
\begin{equation}
 \|(Kh)(t)\|_{H^s_{\rm red}}
 \le C_K\int_0^t(t-\tau)^{\rho-1}
 \|h(\tau)\|_{H^s_{\rm red}}\,\mathrm d\tau.
\label{eq:v2v40-kernel}
\end{equation}
Then, for every \(n\ge0\),
\begin{equation}
 \|K^n\|_{X_T\to X_T}
 \le
 \frac{[C_K\Gamma(\rho)T^\rho]^n}{\Gamma(n\rho+1)}.
\label{eq:v2v40-iterate}
\end{equation}
Consequently the series
\begin{equation}
 \mathcal R_K=(I-K)^{-1}=\sum_{n=0}^\infty K^n
\label{eq:v2v40-resolvent}
\end{equation}
converges in operator norm for every finite \(T\), and
\begin{equation}
 \|\mathcal R_K\|_{X_T\to X_T}
 \le E_\rho(C_K\Gamma(\rho)T^\rho),
 \qquad
 E_\rho(z)=\sum_{n=0}^\infty\frac{z^n}{\Gamma(n\rho+1)}.
\label{eq:v2v40-Mittag-Leffler}
\end{equation}
No assumption \(\|K\|<1\) is required.
\end{theorem}

\begin{proof}
Let \(k_\rho(t)=t^{\rho-1}\mathbf1_{t>0}\).  The beta integral gives
\begin{equation}
 k_\rho^{*n}(t)
 =\frac{\Gamma(\rho)^n}{\Gamma(n\rho)}
 t^{n\rho-1}
 \qquad(n\ge1).
\label{eq:v2v40-convolution}
\end{equation}
Iterating \eqref{eq:v2v40-kernel}, then integrating
\eqref{eq:v2v40-convolution} from zero to \(T\), gives
\eqref{eq:v2v40-iterate}.  The scalar series in
\eqref{eq:v2v40-Mittag-Leffler} converges for every finite argument, so the
operator series converges absolutely.  Multiplying its partial sums by
\(I-K\) leaves \(I-K^{N+1}\); the bound
\eqref{eq:v2v40-iterate} makes the last term tend to zero.  Hence the sum
is the two-sided inverse of \(I-K\).
\end{proof}

For \(\rho=1\), this is the familiar factorial estimate
\(\|K^n\|\le(C_KT)^n/n!\).  The range \(0<\rho<1\) permits an integrable
diagonal singularity \((t-\tau)^{-(1-\rho)}\).

\begin{corollary}[Non-geometric order-reduction remainder]
\label{cor:v2v40-tail}
For the solution \(h=\mathcal R_Kh_0\), the order-\(N\) causal truncation
\(h^{[N]}=\sum_{n=0}^NK^nh_0\) satisfies
\begin{equation}
 \|h-h^{[N]}\|_{X_T}
 \le
 \left[
 \sum_{n=N+1}^\infty
 \frac{[C_K\Gamma(\rho)T^\rho]^n}{\Gamma(n\rho+1)}
 \right]\|h_0\|_{X_T}.
\label{eq:v2v40-tail}
\end{equation}
The bracket tends to zero for every finite \(T\).
\end{corollary}

\begin{proof}
Subtract the partial series from \eqref{eq:v2v40-resolvent}, apply
\eqref{eq:v2v40-iterate}, and sum.
\end{proof}

This improves the geometric V32 error bound when the renormalized
polarization really has the Volterra kernel property.  The V32 criterion
\(\|K\|<1\) remains useful as a simple quantitative estimate, but it is not
necessary for linear existence or convergence of causal order reduction.

\subsection{Continuation without recharging past history}
\label{sec:v2v40-history}

\begin{proposition}[One-use Duhamel history]
\label{prop:v2v40-history}
Let \(h=h_0+Kh\) with a kernel satisfying
\eqref{eq:v2v40-kernel}.  At a continuation time \(t_0\), define
\begin{equation}
 h_{\rm hist}(t)
 =h_0(t)+\int_0^{t_0}k(t,\tau)h(\tau)\,\mathrm d\tau,
 \qquad t\ge t_0.
\label{eq:v2v40-history}
\end{equation}
Then the unknown future obeys
\begin{equation}
 h(t)=h_{\rm hist}(t)
 +\int_{t_0}^tk(t,\tau)h(\tau)\,\mathrm d\tau.
\label{eq:v2v40-restart}
\end{equation}
The past interval enters once, as a known source.  On a slab of length
\(\Delta\), the future Volterra operator has norm at most
\begin{equation}
 \frac{C_K\Delta^\rho}{\rho}.
\label{eq:v2v40-slab-small}
\end{equation}
\end{proposition}

\begin{proof}
Split the causal integral at \(t_0\).  The first part is determined by the
already constructed solution and gives \eqref{eq:v2v40-history}; the second
contains the new unknown.  Integrating
\((t-\tau)^{\rho-1}\) over an interval of length at most \(\Delta\) gives
\eqref{eq:v2v40-slab-small}.
\end{proof}

The proposition is the precise SM analogue of the NS audit rule: past
ancestry is retained, but it is not charged again as part of the new
unknown's contraction constant.

\subsection{Regulator stability and the nonlinear map}
\label{sec:v2v40-regulator}

\begin{theorem}[Regulator convergence of Volterra resolvents]
\label{thm:v2v40-regulator}
Let \(K_\Lambda\) and \(K\) satisfy
\eqref{eq:v2v40-kernel} with the same \(C_K,\rho,T\), and suppose
\begin{equation}
 \|K_\Lambda-K\|_{X_T\to X_T}\longrightarrow0,
 \qquad
 h_{0,\Lambda}\longrightarrow h_0\quad\hbox{in }X_T.
\label{eq:v2v40-K-convergence}
\end{equation}
Then
\begin{equation}
 (I-K_\Lambda)^{-1}h_{0,\Lambda}
 \longrightarrow
 (I-K)^{-1}h_0
 \quad\hbox{in }X_T.
\label{eq:v2v40-solution-convergence}
\end{equation}
It is sufficient for the first convergence in
\eqref{eq:v2v40-K-convergence} that kernel representatives obey
\begin{equation}
 \sup_{0\le t\le T}
 \int_0^t\|k_\Lambda(t,\tau)-k(t,\tau)\|\,\mathrm d\tau
 \longrightarrow0.
\label{eq:v2v40-kernel-convergence}
\end{equation}
\end{theorem}

\begin{proof}
The Volterra theorem bounds all regulated and limiting resolvents by the
same Mittag--Leffler number \(M_T\).  The resolvent identity gives
\[
 (I-K_\Lambda)^{-1}-(I-K)^{-1}
 =(I-K_\Lambda)^{-1}(K_\Lambda-K)(I-K)^{-1}.
\]
Thus the operator difference is at most
\(M_T^2\|K_\Lambda-K\|\), which tends to zero; combine this with source
convergence.  The kernel condition bounds the operator norm directly on
\(L^\infty_tH^s\).
\end{proof}

\begin{theorem}[Nonlinear Volterra fixed point]
\label{thm:v2v40-nonlinear}
Let \(\mathcal R_K=(I-K)^{-1}\) be as above, and suppose
\(\mathcal N:X_T\to X_T\) maps a closed ball of radius \(R\) into
\(X_T\), with
\begin{equation}
 \|\mathcal N(h)-\mathcal N(\widetilde h)\|_{X_T}
 \le L_R\|h-\widetilde h\|_{X_T}.
\label{eq:v2v40-nonlinear-Lip}
\end{equation}
If
\begin{align}
 M_TL_R&<1,
 \label{eq:v2v40-nonlinear-small}\\
 M_T\bigl(\|h_0\|_{X_T}+\|\mathcal N(0)\|_{X_T}\bigr)
 &\le R(1-M_TL_R),
\label{eq:v2v40-ball}
\end{align}
where \(M_T=E_\rho(C_K\Gamma(\rho)T^\rho)\), then
\begin{equation}
 h=\mathcal R_K\bigl(h_0+\mathcal N(h)\bigr)
\label{eq:v2v40-nonlinear-equation}
\end{equation}
has a unique solution in the ball.
\end{theorem}

\begin{proof}
The right side of \eqref{eq:v2v40-nonlinear-equation} has Lipschitz
constant at most \(M_TL_R<1\).  For \(\|h\|\le R\), its norm is at most
\[
 M_T\bigl(\|h_0\|+\|\mathcal N(0)\|+L_RR\bigr)\le R
\]
by \eqref{eq:v2v40-ball}.  Banach's fixed-point theorem applies.
\end{proof}

The nonlinear smallness can be imposed slab by slab using
\cref{prop:v2v40-history}, provided the history source and the solution
ball remain controlled.  This is a continuation criterion, not a global
a priori bound.

\subsection{Insertion into the hybrid moving-band estimate}
\label{sec:v2v40-hybrid}

\begin{theorem}[Conditional causal SM polarization criterion]
\label{thm:v2v40-SM}
Assume:
\begin{enumerate}
 \item the V39 local finite-type and uniform-integrability hypotheses give
       a crossing bound \(B_\kappa\);
 \item the complete local renormalized response has been treated as in
       \cref{lem:v2v40-local}, preserving hyperbolicity and the Ward
       identities;
 \item on the gauge-reduced moving band, the remaining retarded
       polarization satisfies \eqref{eq:v2v40-kernel} with constants
       uniform in the ultraviolet regulator;
 \item the kernel and source convergence in
       \eqref{eq:v2v40-K-convergence} hold.
\end{enumerate}
Then the regulated complete linear low-energy solutions exist on every
finite controlled slab, their causal order-reduction series converge with
the tail \eqref{eq:v2v40-tail}, and the solutions converge in \(X_T\) as
the regulator is removed.  The constant \(C_K\) may contain the hybrid
Einstein factor \(C_E\exp(CB_\kappa)\), but no separate crossing-count
factor occurs.
\end{theorem}

\begin{proof}
The V39 bound controls the hybrid retarded Einstein evolution through V35.
One-sided local absorption leaves the genuinely retarded kernel with only
the factor in \eqref{eq:v2v40-J}.  Apply
\cref{thm:v2v40-Volterra,cor:v2v40-tail} uniformly, then use
\cref{thm:v2v40-regulator}.
\end{proof}

\begin{firewall}[The weakly singular kernel is the new analytic gate]
\label{fw:v2v40-kernel}
V40 proves that global norm smallness is unnecessary once the reduced
polarization has an integrable causal kernel.  It does not prove that the
interacting renormalized SM stress response satisfies
\eqref{eq:v2v40-kernel}.  Retarded quantum kernels are distributions with
light-cone and diagonal singularities.  Local Hadamard counterterms must be
removed, Ward identities preserved, and the remainder shown to map the
chosen Sobolev band with some \(\rho>0\), uniformly in the regulator.  A
critical nonintegrable diagonal singularity corresponds to \(\rho=0\) and
is outside the theorem.
\end{firewall}

\subsection{V40 conclusion and next step}
\label{sec:v2v40-conclusion}

The mandatory audit records the unchanged NS V31 note and the YM advance
to the pure incidence-three closure contained identically in V32 and V33.
Their transferable lesson leads to a sharper SM result: after local
renormalization is assembled once and gauge constraints are reduced, an
integrable retarded polarization kernel has a Mittag--Leffler Volterra
resolvent on every finite slab.  Linear causal existence and convergence of
the reduction series therefore do not require the global contraction
\(\|K\|<1\).  Past history becomes one known Duhamel source at continuation
and is not charged repeatedly.

The main remaining gate is now concrete: prove the weakly singular
Sobolev kernel estimate for the renormalized interacting SM stress
response.  V41 will go upstream to the spectral representation of a causal
two-point response.  It will separate contact terms, derive a sufficient
time-kernel bound from a weighted spectral measure, and identify the exact
high-frequency moment whose finiteness yields \(\rho>0\).
\clearpage
\section*{Second Volume, Version V41: Subtracted Spectral Densities and Integrable Retarded Memory}
\addcontentsline{toc}{section}{Second Volume, Version V41: Subtracted Spectral Densities and Integrable Retarded Memory}

\subsection*{Purpose}

V40 reduced linear causal existence to one analytic statement: after local
renormalization and gauge reduction, the nonlocal polarization must have an
integrable time-diagonal kernel.  This note gives a sufficient spectral
criterion for that statement.  Polynomial frequency terms are identified
as contact terms and assigned to the local V40 block.  The remaining
operator-valued spectral density is estimated by oscillatory integration by
parts.

The decisive exponent is the high-frequency order of the density after all
dispersion subtractions and low-band order reduction.  Positive decay of
that density gives a weakly singular but integrable retarded kernel.  The
result is abstract and sector-independent.  Verifying the required symbol
bounds for the interacting, gauge-reduced Standard Model remains open.

\subsection{Operator-valued retarded sine transforms}
\label{sec:v2v41-sine}

Let \(H^s_{\rm red}\) be the reduced Sobolev space of V40, and let
\(G:(0,\infty)\to\mathcal L(H^s_{\rm red})\) be strongly \(C^1\) for
\(\omega\ge1\).  Define, initially as an improper Bochner integral,
\begin{equation}
 k_G(t)=\mathbf1_{t>0}\int_0^\infty
 \sin(\omega t)G(\omega)\,\mathrm d\omega.
\label{eq:v2v41-kG}
\end{equation}

\begin{theorem}[Spectral symbol implies integrable memory]
\label{thm:v2v41-symbol}
Assume
\begin{align}
 M_{\rm IR}
 &:=\int_0^1\|G(\omega)\|\,\mathrm d\omega<\infty,
\label{eq:v2v41-IR}\\
 \|G(\omega)\|&\le M_q\omega^{-q},
 \qquad
 \|G'(\omega)\|\le M_q\omega^{-q-1}
 \quad(\omega\ge1)
\label{eq:v2v41-symbol-bound}
\end{align}
for some \(q>0\).  Then the integral
\eqref{eq:v2v41-kG} exists for every \(t>0\), and on \(0<t\le T\):
\begin{enumerate}
 \item if \(0<q<1\),
 \begin{equation}
  \|k_G(t)\|
  \le C_q(M_{\rm IR}+M_q)(1+T^{1-q})t^{q-1};
 \label{eq:v2v41-fractional}
 \end{equation}
 \item if \(q=1\),
 \begin{equation}
  \|k_G(t)\|
  \le C(M_{\rm IR}+M_1)(1+|\log t|),
 \label{eq:v2v41-log}
 \end{equation}
 and, for every \(0<\rho<1\), the right side is bounded by
 \(C_{\rho,T}(M_{\rm IR}+M_1)t^{\rho-1}\);
 \item if \(q>1\),
 \begin{equation}
  \|k_G(t)\|\le M_{\rm IR}+\frac{M_q}{q-1}.
 \label{eq:v2v41-bounded}
 \end{equation}
\end{enumerate}
Thus \(k_G\) satisfies the V40 kernel hypothesis with
\(\rho=q\) for \(0<q<1\), with any \(\rho<1\) for \(q=1\), and with
\(\rho=1\) for \(q>1\).
\end{theorem}

\begin{proof}
The low-frequency integral is at most \(M_{\rm IR}\).  For
\(0<t\le1\), split the high-frequency integral at \(R=t^{-1}\).  On
\([1,R]\), direct integration of \eqref{eq:v2v41-symbol-bound} gives
\begin{equation}
 \int_1^R\|G(\omega)\|\,\mathrm d\omega
 \le
 \begin{cases}
  M_q(R^{1-q}-1)/(1-q),&0<q<1,\\
  M_1\log R,&q=1,\\
  M_q/(q-1),&q>1.
 \end{cases}
\label{eq:v2v41-low-piece}
\end{equation}
For the tail, integration by parts on \([R,S]\) gives
\begin{align*}
 \int_R^S\sin(\omega t)G(\omega)\,\mathrm d\omega
 ={}&\left[-\frac{\cos(\omega t)}tG(\omega)\right]_R^S
 +\frac1t\int_R^S\cos(\omega t)G'(\omega)\,\mathrm d\omega.
\end{align*}
The boundary term at \(S\) tends to zero, and the remaining terms are at
most
\begin{equation}
 \frac{M_qR^{-q}}t+
 \frac{M_q}{t}\int_R^\infty\omega^{-q-1}\,\mathrm d\omega
 \le C_qM_qt^{q-1}.
\label{eq:v2v41-tail-piece}
\end{equation}
Combine \eqref{eq:v2v41-low-piece,eq:v2v41-tail-piece}.  For
\(1<t\le T\), split at frequency one and use the same integration by parts
from one to infinity, absorbing the resulting constant into the displayed
\(T\)-dependent bound.  Finally,
\(1+|\log t|\le C_{\rho,T}t^{\rho-1}\) on \((0,T]\), proving the
borderline claim.
\end{proof}

\begin{remark}[The derivative bound supplies the cancellation]
\label{rem:v2v41-derivative}
The pointwise decay of \(G\) alone does not control the non-absolutely
convergent tail when \(q\le1\).  The derivative bound in
\eqref{eq:v2v41-symbol-bound} is the regularity that makes Dirichlet
cancellation quantitative in operator norm.
\end{remark}

\subsection{Contact terms and dispersion subtractions}
\label{sec:v2v41-subtractions}

Let \(p\) be the Laplace frequency with \(\operatorname{Re}p>0\), and let
\(\sigma(\omega)\) be an operator-valued spectral density.  The elementary
identity
\begin{equation}
 \frac1{p^2+\omega^2}
 =\sum_{r=0}^{N-1}\frac{(-p^2)^r}{\omega^{2r+2}}
 +\frac{(-p^2)^N}{\omega^{2N}(p^2+\omega^2)}
\label{eq:v2v41-subtraction-identity}
\end{equation}
separates \(N\) local Taylor coefficients from a nonlocal remainder.

\begin{proposition}[Local/nonlocal spectral split]
\label{prop:v2v41-split}
Suppose a reduced polarization admits the subtracted representation
\begin{equation}
 \widehat\Pi^{\rm ret}(p)
 =P_{N-1}(p^2)
 +(-p^2)^N
 \int_0^\infty
 \frac{\sigma(\omega)}{\omega^{2N}(p^2+\omega^2)}
 \,\mathrm d\omega,
\label{eq:v2v41-dispersion}
\end{equation}
in the relevant Sobolev operator topology.  Then:
\begin{enumerate}
 \item the polynomial \(P_{N-1}(p^2)\) is supported at \(t=0\) after
       inverse Laplace transform and belongs to the V40 local block;
 \item before the exterior factor \((-p^2)^N\) is applied, the retarded
       nonlocal kernel is
 \begin{equation}
  g_N(t)=\mathbf1_{t>0}\int_0^\infty
  \sin(\omega t)G_N(\omega)\,\mathrm d\omega,
  \qquad
  G_N(\omega)=\frac{\sigma(\omega)}{\omega^{2N+1}};
 \label{eq:v2v41-GN}
 \end{equation}
 \item derivatives of the Heaviside factor generated by
       \((-p^2)^N\) are additional contact terms.  After those terms are
       absorbed locally and the remaining time derivatives are causally
       order-reduced on the low band, the nonlocal memory has the same time
       singularity as \(g_N\), multiplied by the norm of the reduced
       low-band differential operator.
\end{enumerate}
\end{proposition}

\begin{proof}
The inverse Laplace transform of a polynomial in \(p\) is a finite linear
combination of derivatives of \(\delta(t)\), proving the first item.  The
identity
\[
 \mathcal L^{-1}\left[(p^2+\omega^2)^{-1}\right](t)
 =\mathbf1_{t>0}\frac{\sin(\omega t)}\omega
\]
gives the second.  Multiplication by powers of \(p\) differentiates in
time and produces the usual initial contact terms.  Once these local terms
are separated, causal low-band order reduction replaces the exterior
derivatives by the specified bounded low-band operator; it does not alter
the scalar time kernel.
\end{proof}

\begin{firewall}[Order reduction is an assumption in the last step]
\label{fw:v2v41-order-reduction}
The algebraic split \eqref{eq:v2v41-subtraction-identity} is exact.  The
replacement of the exterior \(2N\) time derivatives by a bounded low-band
operator requires the V32 order-reduction hypotheses and a bound for its
remainder.  Without that step, differentiating \(g_N\) can restore the
subtracted diagonal singularities.
\end{firewall}

\subsection{The exact high-frequency exponent}
\label{sec:v2v41-exponent}

\begin{theorem}[Subtraction count versus spectral growth]
\label{thm:v2v41-exponent}
Assume, for \(\omega\ge1\),
\begin{equation}
 \|\sigma(\omega)\|\le C_\sigma\omega^m,
 \qquad
 \|\sigma'(\omega)\|\le C_\sigma\omega^{m-1},
\label{eq:v2v41-sigma-order}
\end{equation}
and assume the infrared part of \(G_N\) in
\eqref{eq:v2v41-GN} is integrable.  Put
\begin{equation}
 q_N=2N+1-m.
\label{eq:v2v41-qN}
\end{equation}
If \(q_N>0\), then \(G_N\) satisfies
\eqref{eq:v2v41-symbol-bound} with \(q=q_N\), and the nonlocal kernel has
the V40 exponent supplied by \cref{thm:v2v41-symbol}.  Explicitly:
\begin{enumerate}
 \item \(0<q_N<1\) gives \(\rho=q_N\);
 \item \(q_N=1\) gives any \(0<\rho<1\), with a logarithmic true bound;
 \item \(q_N>1\) gives the bounded case \(\rho=1\).
\end{enumerate}
Thus the strict threshold for this operator-symbol argument is
\begin{equation}
 2N+1>m.
\label{eq:v2v41-threshold}
\end{equation}
\end{theorem}

\begin{proof}
By \eqref{eq:v2v41-GN,eq:v2v41-sigma-order},
\[
 \|G_N(\omega)\|
 \le C_\sigma\omega^{m-2N-1}=C_\sigma\omega^{-q_N}.
\]
Differentiating the quotient gives two terms, both bounded by a constant
times \(\omega^{-q_N-1}\).  Apply
\cref{thm:v2v41-symbol}.
\end{proof}

The exponent \(m\) is measured after inserting the chosen spatial
Sobolev weights and physical gauge projections.  It is not merely the
engineering dimension of an unprojected correlation function.

\begin{corollary}[Causal polarization kernel bound]
\label{cor:v2v41-V40}
Suppose the local/nonlocal split of \cref{prop:v2v41-split} is compatible
with the V40 Ward projectors, the low-band order-reduction operator has norm
at most \(C_{\kappa,N}\), and \eqref{eq:v2v41-threshold} holds.  Then the
reduced nonlocal polarization kernel \(p_{\kappa,N}^{\rm ret}\) satisfies
\begin{equation}
 \|p^{\rm ret}_{\kappa,N}(t,\tau)\|_{H^s_{\rm red}\to H^s_{\rm red}}
 \le C_{\kappa,N,T}(t-\tau)^{\rho-1},
 \qquad 0\le\tau<t\le T,
\label{eq:v2v41-V40-kernel}
\end{equation}
with \(\rho\) as in \cref{thm:v2v41-exponent}, provided the slowly varying
background dependence is uniform in \(\tau\).  If the effective Einstein
propagator obeys
\begin{equation}
 \|e_{\rm eff}(t,s)\|\le C_{E,T},
 \qquad 0\le s\le t\le T,
\label{eq:v2v41-E-propagator}
\end{equation}
then the kernel of \(K=E_{\rm eff}\Pi^{\rm ret}_{\rm nl}\) satisfies
\begin{align}
 k_K(t,\tau)
 &=\int_\tau^t e_{\rm eff}(t,s)
 p^{\rm ret}_{\kappa,N}(s,\tau)\,\mathrm ds,
\label{eq:v2v41-composed-kernel}\\
 \|k_K(t,\tau)\|
 &\le\frac{C_{E,T}C_{\kappa,N,T}}\rho
 (t-\tau)^\rho.
\label{eq:v2v41-composed-bound}
\end{align}
In particular, \cref{thm:v2v40-Volterra} applies to \(K\), with its
bounded-kernel exponent \(\rho_K=1\).
\end{corollary}

\begin{proof}
Apply the low-band operator to the bound from
\cref{thm:v2v41-symbol}; its norm changes only the constant.  Uniformity in
the base time \(\tau\) gives \eqref{eq:v2v41-V40-kernel}.  Compose the two
causal kernels and integrate
\((s-\tau)^{\rho-1}\) from \(\tau\) to \(t\).  This gives
\eqref{eq:v2v41-composed-kernel,eq:v2v41-composed-bound}.  On a finite slab,
\((t-\tau)^\rho\le T^\rho\), so the composed kernel is bounded.
\end{proof}

\subsection{Regulator-uniform spectral convergence}
\subsection{Regulator-uniform spectral convergence}
\label{sec:v2v41-regulator}

For regulated effective densities \(G_{N,\Lambda}\), define
\begin{align}
 \varepsilon_\Lambda
 :={}&
 \int_0^1\|G_{N,\Lambda}(\omega)-G_N(\omega)\|\,\mathrm d\omega
 \notag\\
 &+\sup_{\omega\ge1}\omega^q
 \|G_{N,\Lambda}(\omega)-G_N(\omega)\|
 +\sup_{\omega\ge1}\omega^{q+1}
 \|G'_{N,\Lambda}(\omega)-G'_N(\omega)\|.
\label{eq:v2v41-epsilon}
\end{align}

\begin{theorem}[Spectral-symbol convergence implies kernel convergence]
\label{thm:v2v41-regulator}
Let \(q>0\), and suppose the regulated densities obey uniform versions of
\eqref{eq:v2v41-IR,eq:v2v41-symbol-bound}.  If
\(\varepsilon_\Lambda\to0\), then their retarded kernels converge in the
Volterra operator norm on every finite slab:
\begin{equation}
 \sup_{0\le t\le T}
 \int_0^t
 \|k_{G_{N,\Lambda}}(t-\tau)-k_{G_N}(t-\tau)\|
 \,\mathrm d\tau
 \longrightarrow0.
\label{eq:v2v41-kernel-convergence}
\end{equation}
\end{theorem}

\begin{proof}
Apply the proof of \cref{thm:v2v41-symbol} to the difference of the two
densities.  It gives
\[
 \|k_{G_{N,\Lambda}}(r)-k_{G_N}(r)\|
 \le C_{q,T}\varepsilon_\Lambda r^{\rho-1}
 \]
for the exponent specified there.  Integrating from zero to \(T\) gives a
constant times \(\varepsilon_\Lambda T^\rho/\rho\), which tends to zero.
\end{proof}

If the regulated effective Einstein propagators also converge in the
bounded causal-kernel norm of \eqref{eq:v2v41-E-propagator}, composition
preserves this convergence.  Together with
\cref{thm:v2v40-regulator}, this gives regulator convergence of the
complete causal linear solution.

\subsection{Sharp obstruction at nondecaying effective density}
\label{sec:v2v41-obstruction}

\begin{proposition}[No uniform integrable majorant at order zero]
\label{prop:v2v41-order-zero}
Let
\begin{equation}
 G_\Lambda(\omega)=\mathbf1_{[0,\Lambda]}(\omega)I.
\label{eq:v2v41-flat-density}
\end{equation}
Then
\begin{equation}
 k_\Lambda(t)=\frac{1-\cos(\Lambda t)}tI,
\label{eq:v2v41-flat-kernel}
\end{equation}
and
\begin{equation}
 \int_0^T\|k_\Lambda(t)\|\,\mathrm dt
 \longrightarrow\infty
 \qquad(\Lambda\to\infty).
\label{eq:v2v41-flat-divergence}
\end{equation}
Thus an effective spectral density with no positive decay need not satisfy
any regulator-uniform V40 kernel bound.
\end{proposition}

\begin{proof}
Integrating \(\sin(\omega t)\) in \(\omega\) gives
\eqref{eq:v2v41-flat-kernel}.  With \(x=\Lambda t\),
\[
 \int_0^T\frac{|1-\cos(\Lambda t)|}{t}\,\mathrm dt
 =\int_0^{\Lambda T}\frac{1-\cos x}{x}\,\mathrm dx.
\]
The integral over each interval
\([2\pi n+\pi/2,2\pi n+3\pi/2]\) is bounded below by a constant multiple
of \(n^{-1}\).  The harmonic sum diverges, proving
\eqref{eq:v2v41-flat-divergence}.
\end{proof}

\begin{firewall}[Renormalizability does not yet prove the symbol estimate]
\label{fw:v2v41-SM}
Power-counting renormalizability identifies a finite family of local
counterterms in perturbation theory.  It does not by itself prove the
operator-valued bounds \eqref{eq:v2v41-sigma-order}, infrared integrability,
Ward-compatible spectral projection, or convergence
\eqref{eq:v2v41-epsilon} for the interacting SM stress response on a
time-dependent curved background.  Those are the precise analytic inputs
still missing.  Positivity of a scalar spectral measure is also
insufficient because the relevant gauge-reduced response is matrix-valued
and contains derivative insertions.
\end{firewall}

\subsection{V41 conclusion and next step}
\label{sec:v2v41-conclusion}

After \(N\) dispersion subtractions, a spectral density of Sobolev operator
order \(m\) produces an effective sine density of order
\(-q_N=-(2N+1-m)\).  The strict inequality \(2N+1>m\) yields an integrable
retarded kernel; the borderline is logarithmic, and a nondecaying effective
density admits a regulator sequence with divergent time-kernel norm.
Contact terms belong to the local V40 block, while the remaining memory is
handled by the Volterra resolvent.

V42 will specialize the spectral symbol ledger to free massive scalar,
Dirac, and gauge sectors on an ultrastatic reference slab.  The purpose is
to calculate the actual high-frequency orders after stress-tensor
insertions and gauge reduction, determine the minimal subtraction count in
each sector, and separate what survives perturbatively from what still
requires an interacting curved-background estimate.
\clearpage
\section*{Second Volume, Version V42: Free Standard Model Stress Spectra and the Three-Subtraction Threshold}
\addcontentsline{toc}{section}{Second Volume, Version V42: Free Standard Model Stress Spectra and the Three-Subtraction Threshold}

\subsection*{Purpose and scope}

V41 expressed the retarded-kernel gate in terms of the high-frequency
growth of a gauge-reduced stress spectral measure.  This note calculates
that growth for free scalar, Dirac, and gauge fields on the flat
ultrastatic slab
\begin{equation}
 (\mathbb R\times\mathbb T^3,-\mathrm dt^2+\mathrm dx^2).
\label{eq:v2v42-slab}
\end{equation}
The external metric perturbation is restricted to a fixed spatial Fourier
band \(|\ell|\le\kappa\).  Momentum conservation then reduces the
two-particle sum to one free lattice momentum.  In every physical sector a
stress insertion has matrix element of order one power of the particle
energy.  The resulting dispersion measure has cumulative growth
\(O(R^6)\).

It follows that three dispersion subtractions make the nonlocal spectral
kernel absolutely summable.  These are precisely the subtraction orders
whose local polynomial contains the cosmological, Einstein, and
curvature-squared response rows.  The calculation proves the V40 kernel
bound for the free flat reference theory after low-band order reduction.
It does not establish the corresponding interacting curved-background
estimate.

\subsection{Physical one-particle sectors}
\label{sec:v2v42-sectors}

Let \(a\) label one of the following free fields:
\begin{enumerate}
 \item a real scalar of mass \(m_a\ge0\), with any fixed curvature
       improvement coefficient;
 \item a Dirac field of mass \(m_a\ge0\);
 \item a free gauge field after restriction to the two transverse
       polarizations, with a finite internal multiplicity.
\end{enumerate}
For momentum \(p\in(2\pi\mathbb Z)^3\), write
\begin{equation}
 E_{a,p}=(|p|^2+m_a^2)^{1/2}.
\label{eq:v2v42-energy}
\end{equation}
Massless scalar constant modes and flat gauge-connection modes are removed
from the positive-energy Fock sector and treated as finite-dimensional
infrared variables.  This does not change any ultraviolet estimate below.

For a fixed external spatial momentum \(\ell\), the Fourier coefficient of
the stress tensor creates a particle pair with momenta \(p\) and
\(\ell-p\).  Put
\begin{equation}
 \Omega_{a,p,\ell}=E_{a,p}+E_{a,\ell-p}.
\label{eq:v2v42-pair-energy}
\end{equation}
Let
\(M^{(a)}_{\mu\nu}(p,\ell-p)\) denote the vacuum-to-pair matrix element of
the chosen physical stress component, including polarization indices.

\begin{lemma}[Uniform stress matrix-element order]
\label{lem:v2v42-matrix-element}
For every fixed \(\kappa\), there is \(C_{a,\kappa}\) such that, for
\(|\ell|\le\kappa\),
\begin{equation}
 \|M^{(a)}(p,\ell-p)\|
 \le C_{a,\kappa}(1+\Omega_{a,p,\ell}).
\label{eq:v2v42-matrix-bound}
\end{equation}
This holds for the scalar, Dirac, and physical gauge sectors listed above.
\end{lemma}

\begin{proof}
For a scalar, every quadratic stress term contains either two derivatives
or a mass squared.  Two derivatives contribute \(O(\Omega^2)\), while the
two normalized scalar mode functions contribute
\(O((E_{a,p}E_{a,\ell-p})^{-1/2})=O(\Omega^{-1})\) at high energy.  The
matrix element is therefore \(O(\Omega)\); improvement terms have the same
two-derivative order.

For a Dirac field, normalized spinor wave functions have order zero after
the conventional \((2E)^{-1/2}\) factor is combined with the spinor norm.
The stress bilinear contains one derivative, hence is \(O(\Omega)\).
For a physical gauge field, each normalized vector potential contributes
\(O(E^{-1/2})\), each field strength contributes one derivative, and the
quadratic tensor \(F_{\mu\alpha}F_\nu{}^\alpha\) is again
\(O(\Omega)\).  The finite spin, polarization, and internal multiplicities
only change the constant.  Bounded external \(\ell\) and the finitely many
low modes are absorbed into \(C_{a,\kappa}\).
\end{proof}

\subsection{The two-particle dispersion measure}
\label{sec:v2v42-measure}

Let \(\mathcal V_\ell\) be the finite-dimensional space of metric tensor
polarizations at external momentum \(\ell\).  Define the positive
operator-valued dispersion measure
\begin{equation}
 \mathrm d\Sigma_{a,\ell}(\Omega)
 =\sum_{p,\alpha}
 2\Omega_{a,p,\ell}
 M^{(a)}_{p,\ell,\alpha}{}^*
 M^{(a)}_{p,\ell,\alpha}
 \delta_{\Omega_{a,p,\ell}}(\mathrm d\Omega),
\label{eq:v2v42-measure}
\end{equation}
where \(\alpha\) ranges over the finite physical pair polarizations.  The
factor \(2\Omega\) is the change from the invariant mass variable to the
positive time frequency in the dispersion denominator.

\begin{lemma}[One free lattice momentum]
\label{lem:v2v42-count}
Uniformly for \(|\ell|\le\kappa\),
\begin{equation}
 \#\{p:\Omega_{a,p,\ell}\le R\}
 \le C_\kappa(1+R^3).
\label{eq:v2v42-pair-count}
\end{equation}
\end{lemma}

\begin{proof}
The triangle inequality gives
\(\Omega_{a,p,\ell}\ge |p|+|\ell-p|\ge2|p|-|\ell|\).
Thus \(\Omega\le R\) implies \(|p|\le(R+\kappa)/2\).  The number of
lattice points in a three-dimensional ball of that radius is
\(O_\kappa(1+R^3)\).
\end{proof}

\begin{theorem}[Sixth-order cumulative stress measure]
\label{thm:v2v42-growth}
For each free sector and fixed \(\kappa\),
\begin{equation}
 \left\|\Sigma_{a,\ell}([0,R])\right\|
 \le C_{a,\kappa}(1+R^6),
 \qquad |\ell|\le\kappa.
\label{eq:v2v42-growth}
\end{equation}
The finite sum over all free Standard Model species satisfies the same
bound with a larger constant.
\end{theorem}

\begin{proof}
For every atom with \(\Omega\le R\),
\cref{lem:v2v42-matrix-element} gives
\[
 \|2\Omega M^*M\|
 \le C_{a,\kappa}(1+R^3).
\]
There are at most \(C_\kappa(1+R^3)\) atoms after the finite polarization
multiplicity is included, by \cref{lem:v2v42-count}.  Positivity permits
the operator norm to be bounded by the sum of atom norms, yielding
\eqref{eq:v2v42-growth}.  The Standard Model has finitely many free field
components, so their measures add without changing the exponent.
\end{proof}

\begin{remark}[Gauge reduction precedes positivity]
\label{rem:v2v42-gauge}
The measure \eqref{eq:v2v42-measure} is formed after passing to transverse
physical gauge polarizations.  In a covariant gauge, ghost and longitudinal
terms enforce Ward identities but do not define separate positive measures.
They must be combined before taking the physical spectral norm, in accord
with the V40 audit rule.
\end{remark}

\subsection{Three subtractions give an absolutely summable kernel}
\label{sec:v2v42-three}

Choose an ultraviolet partition at a fixed positive frequency \(\mu\).
The finitely many modes below \(2\mu\) are left unsubtracted.  On the high
part, apply \eqref{eq:v2v41-subtraction-identity}.  After \(N\)
subtractions the nonlocal sine measure has the moment
\begin{equation}
 \mathfrak M_{a,N}(\mu)
 =\int_{[\mu,\infty)}
 \Omega^{-(2N+1)}\,mathrm d
 \|\Sigma_{a,\ell}\|(\Omega).
\label{eq:v2v42-moment}
\end{equation}

\begin{theorem}[Three-subtraction threshold]
\label{thm:v2v42-three}
For the free stress measures in \cref{thm:v2v42-growth},
\begin{equation}
 \sup_{|\ell|\le\kappa}\mathfrak M_{a,3}(\mu)<\infty.
\label{eq:v2v42-three-moment}
\end{equation}
Consequently the high-frequency nonlocal kernel
\begin{equation}
 g_{a,3}^{\rm hi}(t,\ell)
 =\mathbf1_{t>0}
 \int_{[\mu,\infty)}
 \frac{\sin(\Omega t)}{\Omega^7}
 \,\mathrm d\Sigma_{a,\ell}(\Omega)
\label{eq:v2v42-high-kernel}
\end{equation}
converges absolutely and obeys
\begin{equation}
 \sup_{t\ge0,\,|\ell|\le\kappa}
 \|g_{a,3}^{\rm hi}(t,\ell)\|
 \le C_{a,\kappa,\mu}.
\label{eq:v2v42-high-bound}
\end{equation}
The unsubtracted low-frequency part is a finite sum and is also bounded.
\end{theorem}

\begin{proof}
Let \(F(R)=\|\Sigma_{a,\ell}([\mu,R])\|\).  Positivity and
\eqref{eq:v2v42-growth} give \(F(R)\le C(1+R^6)\), uniformly in the
external band.  Stieltjes integration by parts yields
\begin{align*}
 \int_{[\mu,R]}\Omega^{-7}\,\mathrm dF(\Omega)
 =R^{-7}F(R)-\mu^{-7}F(\mu^-)
 +7\int_\mu^R\Omega^{-8}F(\Omega)\,\mathrm d\Omega.
\end{align*}
The boundary term at infinity vanishes, and the last integrand is
\(O(\Omega^{-2})\).  This proves
\eqref{eq:v2v42-three-moment}.  Since \(|\sin|\le1\), the same moment
dominates \eqref{eq:v2v42-high-kernel} uniformly in time.  Only finitely
many lattice pairs lie below the fixed cutoff, so the low part is bounded.
\end{proof}

\begin{corollary}[Free-sector V40 kernel]
\label{cor:v2v42-V40}
Suppose the exterior sixth-order time operator produced by the three
subtractions is causally order-reduced on the fixed band, with norm
\(C_{\kappa,3}\), and all contact terms are included in the V40 local
block.  Then every free scalar, Dirac, and physical gauge sector has a
bounded nonlocal polarization kernel.  After composition with the bounded
retarded Einstein propagator, the Volterra operator satisfies
\cref{thm:v2v40-Volterra} with \(\rho=1\).
\end{corollary}

\begin{proof}
The low and high parts are bounded by
\cref{thm:v2v42-three}.  Low-band order reduction multiplies this bound by
\(C_{\kappa,3}\).  Apply the causal composition estimate
\eqref{eq:v2v41-composed-bound} in its bounded-kernel case.
\end{proof}

\subsection{Why the ultraviolet partition is needed for massless fields}
\label{sec:v2v42-infrared}

Applying the Taylor identity \eqref{eq:v2v41-subtraction-identity} down to
\(\Omega=0\) introduces negative spectral moments.  This is harmless for a
massive sector above threshold but can be false for a massless sector.

\begin{proposition}[High-only subtraction preserves the infrared response]
\label{prop:v2v42-high-only}
Let \(\chi_{\rm hi}(\Omega)\) vanish on \([0,\mu]\) and equal one on
\([2\mu,\infty)\).  Split the exact dispersion integral into low and high
parts, and apply \eqref{eq:v2v41-subtraction-identity} only after
multiplication by \(\chi_{\rm hi}\).  Then:
\begin{enumerate}
 \item the low-frequency spectral response is unchanged;
 \item the transition interval contributes a smooth bounded kernel;
 \item the high-frequency remainder has the moment
       \eqref{eq:v2v42-three-moment};
 \item the resulting polynomial coefficients are local renormalization
       terms and may be fixed by the same renormalization conditions as in
       V40.
\end{enumerate}
\end{proposition}

\begin{proof}
The partition of unity is exact.  The subtraction identity is applied
pointwise only on the high support, where every negative power of
\(\Omega\) is bounded at its lower endpoint.  The compact transition
support gives a smooth finite spectral sum.  The high tail is exactly the
one estimated in \cref{thm:v2v42-three}.  The finite Taylor powers in the
Laplace frequency invert to contact terms.
\end{proof}

\begin{assessment}[The local polynomial matches gravitational counterterms]
\label{ass:v2v42-counterterms}
For \(N=3\), the local polynomial in \(p^2\) has degrees zero, one, and
two.  In a covariant tensor reconstruction these are the response orders
associated with the cosmological term, the Einstein--Hilbert term, and the
curvature-squared terms.  This correspondence explains why the first
absolutely summable nonlocal remainder occurs after the standard local
four-derivative gravitational ambiguity has been separated.  The statement
classifies derivative order; it does not select numerical renormalized
couplings.
\end{assessment}

\subsection{Sharpness in a nonzero shear channel}
\label{sec:v2v42-sharpness}

\begin{proposition}[Two subtractions are not absolutely summable]
\label{prop:v2v42-sharp}
For each nontrivial free sector, there is a physical stress polarization
and a conic set of lattice momenta for which
\begin{equation}
 \|M^{(a)}(p,\ell-p)v\|^2\ge c_a|p|^2
\label{eq:v2v42-lower-matrix}
\end{equation}
at large \(|p|\).  Consequently the cumulative scalar measure in that
polarization has a lower bound \(cR^6\), and
\begin{equation}
 \int_{[\mu,\infty)}\Omega^{-5}
 \,\mathrm d\langle v,\Sigma_{a,\ell}(\Omega)v\rangle
 =\infty.
\label{eq:v2v42-two-diverges}
\end{equation}
Thus the two-subtraction remainder is not absolutely summable in general.
\end{proposition}

\begin{proof}
Choose a transverse-traceless spatial shear polarization.  For the scalar,
the leading pair matrix element is a nonzero quadratic form in the
components of \(p\), divided by the mode-normalization energy, and hence is
bounded below by \(c|p|\) away from its algebraic zero cone.  The analogous
leading Dirac and gauge shear amplitudes are nonzero homogeneous functions
of degree one after physical polarization is imposed.  A fixed conic
subset of lattice momenta avoids their zero sets and contains
\(cR^3+o(R^3)\) points below radius \(R\).  Each dispersion atom includes
the additional factor \(2\Omega\), so its scalar weight is at least
\(cR^3\) on a dyadic shell.  Summing shells gives cumulative growth at
least \(cR^6\).

For the two-subtraction moment, each dyadic shell then contributes at least
a constant multiple of \(R^{-5}R^6\), after taking the difference of
successive cumulative bounds; equivalently, Stieltjes comparison gives a
divergent integral with density order \(R^5\).  Hence
\eqref{eq:v2v42-two-diverges}.
\end{proof}

\begin{firewall}[The sharpness statement is channelwise]
\label{fw:v2v42-sharpness}
Special tensor components can have cancellations; for example, an
integrated Hamiltonian component may be diagonal in particle number.  The
claim is that the complete physical stress response contains at least one
nonzero shear channel with sixth-order cumulative growth.  Therefore a
uniform tensor estimate cannot rely on cancellations confined to a
different component.
\end{firewall}

\subsection{Regulator removal for the free reference theory}
\label{sec:v2v42-regulator}

Let \(\Sigma_{a,\ell,\Lambda}\) retain only pair energies below a regulator
\(\Lambda\).  Absolute summability gives a direct convergence result.

\begin{theorem}[Uniform convergence of free regulated memory]
\label{thm:v2v42-regulator}
For \(N=3\),
\begin{equation}
 \sup_{t\ge0,\,|\ell|\le\kappa}
 \|g_{a,3,\Lambda}^{\rm hi}(t,\ell)
       -g_{a,3}^{\rm hi}(t,\ell)\|
 \longrightarrow0.
\label{eq:v2v42-regulator}
\end{equation}
The finite sum over free Standard Model species converges with the same
property.
\end{theorem}

\begin{proof}
The difference is the tail of \eqref{eq:v2v42-high-kernel} above
\(\Lambda\).  Its norm is at most the tail of the finite positive moment
\eqref{eq:v2v42-three-moment}, which tends to zero independently of
\(t\) and uniformly in the finite external band.  Sum the finitely many
species.
\end{proof}

\begin{firewall}[Free flat closure is not interacting curved closure]
\label{fw:v2v42-status}
The proof uses exact free-particle pairs, flat spatial momentum
conservation, a time-translation-invariant vacuum, and a physical
one-particle gauge reduction.  On a curved evolving slice, mode mixing
replaces the single lattice sum; in an interacting theory, multiparticle
continua, anomalous logarithms, composite-operator mixing, and BRST Ward
identities must be controlled.  The three-subtraction result identifies
the expected ultraviolet threshold and proves it in the reference model;
it does not provide the uniform interacting estimates required by
\cref{fw:v2v41-SM}.
\end{firewall}

\subsection{V42 conclusion and next step}
\label{sec:v2v42-conclusion}

On the free flat ultrastatic reference slab, every scalar, Dirac, and
physical gauge stress sector has dispersion-measure growth
\(O(R^6)\) on a fixed external spatial band.  Three subtractions leave the
finite moment \(\int\Omega^{-7}\,\mathrm d\|\Sigma\|\), hence a bounded
retarded memory kernel and uniform regulator removal.  Two subtractions
fail in a nonzero shear channel.  A high-frequency partition keeps
massless infrared response intact while assigning only contact terms to
the local gravitational block.

V43 will test how much of this survives perturbation theory.  It will
formulate a Ward-compatible weighted spectral bound allowing powers of
\(\log\Omega\), prove that three subtractions still suffice for every fixed
logarithmic order, and isolate the uniform-in-order or nonperturbative
estimate needed for the complete interacting Standard Model.
\clearpage
\section*{Second Volume, Version V43: Logarithmic Spectral Growth, Ward Assembly, and the Nonperturbative Dini Gate}
\addcontentsline{toc}{section}{Second Volume, Version V43: Logarithmic Spectral Growth, Ward Assembly, and the Nonperturbative Dini Gate}

\subsection*{Purpose}

V42 proved that a free stress dispersion measure with cumulative growth
\(R^6\) has a finite three-subtraction moment.  Renormalized perturbation
theory generally adds logarithms at high frequency.  This note proves that
every fixed logarithmic power remains harmless for the same subtraction
count.  It also states the exact cumulative Dini condition needed for a
regulator-uniform or resummed interacting response.

The Ward identity must be imposed before the spectral norm is estimated.
Gauge, longitudinal, and ghost contributions are not assigned independent
positive spectral budgets.  Their complete renormalized sum descends to
the physical quotient, where the Wightman spectral measure is positive.
This is the response-theoretic version of the one-sided complete-payer rule
used in V40.

\subsection{The exact cumulative moment criterion}
\label{sec:v2v43-Dini}

Let \(\Sigma\) be a positive operator-valued measure on a finite-dimensional
external polarization space, and set
\begin{equation}
 F(R)=\|\Sigma([\mu,R])\|,
 \qquad R\ge\mu>0.
\label{eq:v2v43-F}
\end{equation}

\begin{theorem}[Cumulative Dini criterion for three subtractions]
\label{thm:v2v43-Dini}
Suppose
\begin{align}
 R^{-7}F(R)&\longrightarrow0,
\label{eq:v2v43-boundary}\\
 \mathfrak D_3(\Sigma)
 &:=\int_\mu^\infty R^{-8}F(R)\,\mathrm dR<\infty.
\label{eq:v2v43-Dini}
\end{align}
Then
\begin{equation}
 \int_{[\mu,\infty)}\Omega^{-7}
 \,\mathrm d\|\Sigma\|(\Omega)<\infty,
\label{eq:v2v43-moment}
\end{equation}
and the three-subtraction sine kernel is uniformly bounded in time.
Conversely, for a scalar positive measure, finiteness of
\eqref{eq:v2v43-moment} implies \eqref{eq:v2v43-Dini} and the boundary
condition after modification on a bounded interval.
\end{theorem}

\begin{proof}
For a positive operator-valued measure on a finite-dimensional space, its
total variation is controlled by a dimensional constant times its trace
measure.  It is therefore enough to prove the scalar statement for each
matrix entry, or for the trace.  Stieltjes integration by parts gives
\begin{equation}
 \int_{[\mu,R]}\Omega^{-7}\,\mathrm dF(\Omega)
 =R^{-7}F(R)-\mu^{-7}F(\mu^-)
 +7\int_\mu^R\Omega^{-8}F(\Omega)\,\mathrm d\Omega.
\label{eq:v2v43-parts}
\end{equation}
The hypotheses make the right side converge.  Absolute convergence and
\(|\sin|\le1\) bound the time kernel.

Conversely, if the weighted scalar measure is finite, Tonelli's theorem
gives
\begin{align*}
 \int_\mu^\infty R^{-8}F(R)\,\mathrm dR
 &=\int_{[\mu,\infty)}
   \left(\int_\Omega^\infty R^{-8}\,\mathrm dR\right)
   \,\mathrm dF(\Omega)\\
 &=\frac17\int_{[\mu,\infty)}\Omega^{-7}\,\mathrm dF(\Omega).
\end{align*}
The weighted tail also forces \(R^{-7}F(R)\to0\) along all large \(R\)
after the irrelevant bounded initial mass is removed.
\end{proof}

\begin{remark}[Why positivity is used]
\label{rem:v2v43-positivity}
For a signed or non-normal operator-valued distribution, the cumulative
operator norm need not control total variation.  The theorem is applied
after physical Ward reduction, where the spectral measure of the
gauge-invariant stress observable is positive.  Before that reduction one
must assume a total-variation version of \eqref{eq:v2v43-Dini} directly.
\end{remark}

\subsection{Fixed logarithmic orders are harmless}
\label{sec:v2v43-logs}

\begin{corollary}[Any fixed logarithmic power]
\label{cor:v2v43-logs}
If, for fixed \(L\ge0\),
\begin{equation}
 F(R)\le C R^6(1+\log R)^L
 \qquad(R\ge2),
\label{eq:v2v43-log-growth}
\end{equation}
then \eqref{eq:v2v43-moment} holds.  More precisely,
\begin{equation}
 \int_2^\infty R^{-8}F(R)\,\mathrm dR
 \le C\int_2^\infty R^{-2}(1+\log R)^L\,\mathrm dR<\infty.
\label{eq:v2v43-log-integral}
\end{equation}
\end{corollary}

\begin{proof}
The boundary quantity is bounded by
\(CR^{-1}(1+\log R)^L\), which tends to zero.  The integral in
\eqref{eq:v2v43-log-integral} is finite after the substitution
\(u=\log R\), because it becomes an exponentially decaying integral times
a polynomial.  Apply \cref{thm:v2v43-Dini}.
\end{proof}

\begin{corollary}[A small anomalous power is also allowed]
\label{cor:v2v43-power}
If
\begin{equation}
 F(R)\le C R^{6+\delta}
 \qquad\hbox{for some }0\le\delta<1,
\label{eq:v2v43-subcritical-power}
\end{equation}
then three subtractions still give a bounded memory kernel.  The endpoint
\(F(R)\simeq R^7\) fails the Dini condition.
\end{corollary}

\begin{proof}
The Dini integrand is \(O(R^{-2+\delta})\), which is integrable exactly for
\(\delta<1\).  At \(\delta=1\), it has the nonintegrable order \(R^{-1}\).
\end{proof}

\subsection{Ward-compatible spectral assembly}
\label{sec:v2v43-Ward}

Let \(\mathcal H_{\rm gf}\) be a gauge-fixed state space and let
\(Q_{\rm phys}\) denote the quotient map onto physical cohomology.  Write
the complete renormalized stress response as
\begin{equation}
 \Pi_{\rm tot}
 =\Pi_{\rm gauge}+\Pi_{\rm matter}+\Pi_{\rm ghost}
 +\Pi_{\rm ct}.
\label{eq:v2v43-total}
\end{equation}

\begin{proposition}[Ward reduction before total variation]
\label{prop:v2v43-Ward}
Suppose the renormalized Ward identity gives
\begin{equation}
 Q_{\rm phys}\Pi_{\rm tot}
 =Q_{\rm phys}\Pi_{\rm tot}Q_{\rm phys},
\label{eq:v2v43-Ward}
\end{equation}
and the induced physical stress observable is self-adjoint in a positive
Hilbert space.  Then its nonlocal Wightman spectral measure
\(\Sigma_{\rm phys}\) is positive.  If its cumulative function satisfies
\eqref{eq:v2v43-Dini}, the corresponding retarded response obeys the V42
three-subtraction kernel bound.  No separate positivity or Dini estimate
is required for the ghost and longitudinal summands in
\eqref{eq:v2v43-total}.
\end{proposition}

\begin{proof}
Equation \eqref{eq:v2v43-Ward} makes the complete response descend to the
physical quotient.  For a physical polarization \(v\), the Wightman
spectral weight is a sum of squares
\[
 \mathrm d\langle v,\Sigma_{\rm phys}v\rangle
 =\sum_n|\langle n,T(v)\Omega\rangle|^2
 \delta_{E_n-E_0},
\]
and is nonnegative.  Apply \cref{thm:v2v43-Dini} to this assembled measure
and then the local/nonlocal split of V41.  Ghost and longitudinal terms are
parts of the Ward representative and need not be positive separately.
\end{proof}

\begin{firewall}[An anomalous Ward term is local only if proved local]
\label{fw:v2v43-anomaly}
If regularization violates the Ward identity, the defect may be moved into
the contact block only after it is shown to be a local covariant
polynomial and cancelled by an allowed counterterm.  A nonlocal Ward defect
cannot be discarded or estimated as a separate positive sector.  The
absence of a diffeomorphism anomaly for the intended complete SM matter
content must be implemented by the regulator, not merely assumed from a
formal diagram cancellation.
\end{firewall}

\subsection{Fixed perturbative order}
\label{sec:v2v43-perturbative}

Let \(g\) collectively denote the renormalized SM couplings, and write a
formal expansion of the physical spectral measure
\begin{equation}
 \Sigma_{\rm phys}(g)
 \sim\sum_{r=0}^\infty g^r\Sigma_r.
\label{eq:v2v43-series}
\end{equation}

\begin{theorem}[Finite-order logarithmic stability]
\label{thm:v2v43-finite-order}
Fix \(R_*\in\mathbb N\).  Suppose the Ward-completed coefficient at each
order \(0\le r\le R_*\) has total variation satisfying
\begin{equation}
 F_r(R)le C_rR^6(1+\log R)^{L_r}
\label{eq:v2v43-order-bound}
\end{equation}
for finite \(C_r,L_r\).  Then the truncated interacting response
\begin{equation}
 \Sigma^{[R_*]}(g)=\sum_{r=0}^{R_*}g^r\Sigma_r
\label{eq:v2v43-truncated}
\end{equation}
has a finite three-subtraction total-variation moment for every fixed
\(g\).  Its retarded nonlocal kernel therefore satisfies the V40 Volterra
hypothesis after the local contact terms and low-band exterior derivatives
are treated as in V41.
\end{theorem}

\begin{proof}
Let \(L_*\) be the maximum of the finitely many \(L_r\).  The total
variation of the finite sum is bounded by
\[
 \sum_{r=0}^{R_*}|g|^rC_rR^6(1+\log R)^{L_*}.
\]
Apply \cref{cor:v2v43-logs} and sum the finitely many moment bounds.  V41
then converts the spectral moment to the retarded kernel estimate.
\end{proof}

The theorem is conditional on the standard fixed-order renormalized
amplitude estimate \eqref{eq:v2v43-order-bound}; it does not derive that
estimate on an evolving curved background.  Its content is that no
additional subtraction is forced by any finite logarithmic degree.

\subsection{The uniform-in-order and nonperturbative gate}
\label{sec:v2v43-nonperturbative}

\begin{theorem}[A sufficient resummed majorant]
\label{thm:v2v43-resummed}
Suppose the complete physical spectral measure exists and, uniformly in
the ultraviolet regulator,
\begin{align}
 F_\Lambda(R)&\le R^6\mathcal L(R),
\label{eq:v2v43-majorant}\\
 \int_\mu^\infty R^{-2}\mathcal L(R)\,\mathrm dR&<\infty,
\label{eq:v2v43-majorant-Dini}\\
 R^{-1}\mathcal L(R)&\longrightarrow0.
\label{eq:v2v43-majorant-boundary}
\end{align}
Then the three-subtraction moments are uniformly bounded.  If in addition
their weighted tails are tight,
\begin{equation}
 \lim_{R\to\infty}\sup_\Lambda
 \int_{[R,\infty)}\Omega^{-7}
 \,\mathrm d\|\Sigma_\Lambda\|(\Omega)=0,
\label{eq:v2v43-tight}
\end{equation}
and the measures converge on bounded frequency intervals, the nonlocal
kernels converge uniformly in time on finite slabs.
\end{theorem}

\begin{proof}
Insert \eqref{eq:v2v43-majorant} into
\eqref{eq:v2v43-Dini}; the remaining integral is precisely
\eqref{eq:v2v43-majorant-Dini}, while
\eqref{eq:v2v43-majorant-boundary} supplies the boundary condition.
For convergence, split the weighted spectral integral at \(R\).  On the
bounded interval use convergence of the measures; on the tail use
\eqref{eq:v2v43-tight}.  Since \(|\sin|\le1\), the estimate is uniform in
time.
\end{proof}

\begin{proposition}[Finite-order control does not imply a resummed bound]
\label{prop:v2v43-no-resummation}
Knowledge that every coefficient in \eqref{eq:v2v43-series} separately
satisfies \eqref{eq:v2v43-order-bound} does not imply
\eqref{eq:v2v43-majorant-Dini} for the sum.
\end{proposition}

\begin{proof}
The constants \(C_r\) and logarithmic degrees \(L_r\) may grow without any
summable majorant.  For example, bounds of the form
\(C_r=r!\) are compatible with every fixed-order statement but make
\(\sum_r|g|^rC_r\) divergent for all \(g\ne0\).  Even a formally resummed
growth \(F(R)\simeq R^7/\log R\) has finite values at every finite
frequency but violates the Dini condition through
\(\int^infty(R\log R)^{-1}\mathrm dR\).  Thus a uniform summability or
nonperturbative spectral estimate is additional information.
\end{proof}

\subsection{Insertion into the causal continuum branch}
\label{sec:v2v43-insertion}

\begin{corollary}[What fixed-order perturbation theory would yield]
\label{cor:v2v43-continuum}
At any fixed perturbative order satisfying
\eqref{eq:v2v43-order-bound}, assume also the V39 crossing estimate, the
V40 hyperbolic and local-block hypotheses, and regulator convergence of the
weighted spectral measures.  Then the gauge-reduced causal low-energy
Einstein--SM equation at that perturbative order has the V40 Volterra
resolvent, convergent order reduction, and regulator limit on every finite
controlled slab.
\end{corollary}

\begin{proof}
Use \cref{thm:v2v43-finite-order} for the polarization kernel, V39 for the
moving-band Einstein propagator, and
\cref{thm:v2v40-Volterra,thm:v2v40-regulator} for causal existence and
convergence.
\end{proof}

\begin{firewall}[This is not a nonperturbative SM construction]
\label{fw:v2v43-status}
V43 proves stability of the three-subtraction threshold under every fixed
logarithmic perturbative order.  It neither constructs the interacting SM
Hilbert space and stress spectral measure nor proves the uniform Dini
majorant \eqref{eq:v2v43-majorant-Dini}.  It also does not transfer the flat
spectral estimate to a general evolving metric.  Those two steps are
independent: curved microlocal control concerns local ultraviolet
propagation, while the Dini majorant concerns summation over all
interaction orders and spectral states.
\end{firewall}

\subsection{V43 conclusion and next step}
\label{sec:v2v43-conclusion}

Three subtractions tolerate \(R^6\) growth times any fixed logarithmic
power, and even a power loss smaller than one.  The exact positive-measure
criterion is the cumulative Dini integral
\(\int R^{-8}F(R)\,\mathrm dR\).  Ward completion must precede positivity
and total variation.  Fixed-order perturbation theory can satisfy the
criterion without providing any uniform-in-order or nonperturbative
majorant.

V44 will address the geometric half of the remaining problem.  It will
replace flat momentum conservation by a local pseudodifferential stress
kernel on a bounded-geometry slab, separate the universal diagonal symbol
from smooth state dependence, and identify the uniform symbol seminorms
needed to carry the V43 Dini estimate through a slowly evolving metric.
\clearpage
\section*{Second Volume, Version V44: Curved Hadamard Scaling and the Order-Minus-Two Memory Symbol}
\addcontentsline{toc}{section}{Second Volume, Version V44: Curved Hadamard Scaling and the Order-Minus-Two Memory Symbol}

\subsection*{Purpose and terminology}

V42--V43 used time-translation-invariant spectral measures.  A general
Einstein foliation has no conserved frequency and no global spectral
measure.  The appropriate replacement is the local Fourier symbol in the
relative time variable.  This note derives its ultraviolet order for free
scalar, Dirac, and gauge sectors from the curved-spacetime Hadamard
parametrix.

Here ``Hadamard'' again refers to the microlocal short-distance form of a
quantum two-point function, as in V26.  It is unrelated to the eigenvalue
shape derivative called a Hadamard formula in V37.  The main result is that
the Ward-completed free stress response has diagonal scaling degree at most
eight in four dimensions.  Its local ambiguity has differential order at
most four.  After three causal subtractions, the high relative-time symbol
has order minus two, which is integrable in frequency and gives a bounded
memory kernel.

\subsection{Uniform local geometric setup}
\label{sec:v2v44-setup}

Let \((M,g)\) be a four-dimensional globally hyperbolic slab foliated by
compact Cauchy surfaces.  Fix a finite atlas of causally convex normal
neighbourhoods and a subordinate partition of unity.  Assume bounded
geometry through the finite jet order used below: the metric and inverse,
lapse, shift, curvature, bundle connections, and their required
derivatives are uniformly bounded, with a common hyperbolicity and
injectivity-radius constant.

For a free field sector \(a\), write its two-point function as
\begin{equation}
 W_a(x,y)=H_a(x,y)+S_a(x,y),
\label{eq:v2v44-HS}
\end{equation}
where \(H_a\) is a local Hadamard parametrix and \(S_a\) is smooth.  For a
family of states or regulators, assume the required \(C^r\) seminorms of
\(S_a\) are uniform on the slab.  This uniformity is stronger than the
statement that every member of the family is individually Hadamard.
Assume the state is quasifree, or more generally that the connected
four-point remainder beyond its Wick contractions is smooth with the same
uniform seminorm control imposed below.

\subsection{Scaling degree of the free stress response}
\label{sec:v2v44-scaling}

Recall that the scaling degree of a distribution at the diagonal is the
infimum of \(\delta\) for which its rescaling by relative coordinates with
weight \(\lambda^\delta\) tends to zero as \(\lambda\downarrow0\).  In four
dimensions the scalar and gauge-potential Hadamard two-point functions have
scaling degree two, while the Dirac two-point function has scaling degree
three.

\begin{lemma}[Free-sector scaling table]
\label{lem:v2v44-table}
For the connected two-point function of two stress tensors, away from the
diagonal, the leading Wick contractions have the following maximal scaling
degrees:
\begin{center}
\begin{tabular}{c c c c}
\toprule
sector & two propagators & stress derivatives & total\\
\midrule
scalar & \(2+2\) & \(4\) & \(8\)\\
Dirac & \(3+3\) & \(2\) & \(8\)\\
gauge potential and ghosts & \(2+2\) & \(4\) & \(8\)\\
\bottomrule
\end{tabular}
\end{center}
Curvature improvement, mass, and lower-order connection terms do not raise
these totals.
\end{lemma}

\begin{proof}
For a Gaussian field, the connected stress--stress function is a finite
sum of products of two differentiated two-point functions.  Multiplication
adds scaling degrees when the Hadamard wavefront condition permits the
product, and each relative derivative raises scaling degree by at most one.

A scalar stress tensor is quadratic with at most two derivatives at each
insertion, so the two scalar propagators contribute four and the four
derivatives contribute at most four.  A Dirac stress tensor is bilinear
with one derivative at each insertion; two Dirac propagators contribute
six and the two derivatives contribute two.  A gauge stress tensor is
quadratic in field strengths, giving four derivatives of two
gauge-potential propagators.  Gauge-fixing and ghost contractions have no
higher differential order.  Mass, curvature, and improvement terms replace
some derivatives by smooth coefficients and therefore cannot increase the
scaling degree.
\end{proof}

\begin{theorem}[Diagonal extension and local ambiguity]
\label{thm:v2v44-extension}
Each Ward-completed free stress response on \(M\setminus\operatorname{Diag}\)
extends across the diagonal with scaling degree at most eight.  The
difference between two extensions with the same covariance and field
content is a finite sum
\begin{equation}
 \sum_{|\alpha|\le4}C_\alpha(x)\nabla^\alpha\delta(x,y),
\label{eq:v2v44-local-ambiguity}
\end{equation}
subject to the tensor symmetries and Ward identities.  Hence all extension
ambiguities belong to the local differential block of V40 and have order
at most four.
\end{theorem}

\begin{proof}
The extension theorem for distributions of finite scaling degree says that
a distribution of scaling degree \(s\) in codimension four has extension
ambiguity supported on the diagonal with derivative order at most
\(\lfloor s-4\rfloor\).  By \cref{lem:v2v44-table}, \(s\le8\), giving
order at most four.  Local covariance and the Ward identities restrict the
coefficient tensors \(C_\alpha\), but cannot increase their differential
order.  These diagonal terms invert to contact operators in the retarded
equation, so they are assigned to the V40 local block.
\end{proof}

\subsection{Universal singular part and state-dependent remainder}
\label{sec:v2v44-state}

\begin{proposition}[State dependence is lower order]
\label{prop:v2v44-state}
Insert \eqref{eq:v2v44-HS} into the free Wick contractions.  The term
quadratic in \(H_a\) has the maximal scaling degree in
\cref{lem:v2v44-table}.  Every term containing one factor \(S_a\) has
scaling degree at most six, and the term quadratic in \(S_a\) is smooth.
The corresponding bounds are uniform for a family with uniform required
\(C^r\) seminorms of \(S_a\).
\end{proposition}

\begin{proof}
The smooth factor \(S_a\) has scaling degree zero.  For the scalar and gauge
sectors, a cross term has one singular propagator of scaling degree two and
at most four derivatives, hence scaling degree at most six.  For the Dirac
sector it has one singular propagator of degree three and at most two
derivatives, hence degree at most five.  The weaker common bound six is
therefore valid.  A product of smooth kernels remains smooth after finitely
many derivatives.  Uniform seminorms of \(S_a\) make all constants uniform.
\end{proof}

Thus the order-four ultraviolet symbol is fixed by the local parametrix.
Changing the Hadamard state changes lower-order nonlocal terms and smooth
terms, but not the leading subtraction count.

\subsection{Relative-time symbol after a spatial low-band cutoff}
\label{sec:v2v44-symbol}

Use a foliation chart with central and relative times
\begin{equation}
 u=\frac{t+t'}2,
 \qquad r=t-t'.
\label{eq:v2v44-times}
\end{equation}
Insert the V33 smooth spatial cutoffs on both external tensor variables,
with fixed physical \(\kappa\).  Let
\(a_{\kappa}(u,x,y,\omega)\) denote the local Fourier amplitude in \(r\)
of the Ward-reduced retarded stress response after the diagonal contact
terms \eqref{eq:v2v44-local-ambiguity} are separated.

The explicit Hadamard expansion is log-polyhomogeneous.  The uniform form
needed below is the following finite family of seminorm estimates:
\begin{equation}
 \|\partial_\omega^j a_\kappa(u,x,y,\omega)\|
 \le C_{j,\kappa}
 \langle\omega\rangle^{4-j}
 (1+\log\langle\omega\rangle)^L,
 \qquad 0\le j\le j_*,
\label{eq:v2v44-order-four}
\end{equation}
for a finite \(j_*\) large enough for the operator bounds.

\begin{theorem}[Free Hadamard symbol order]
\label{thm:v2v44-order-four}
For the free scalar, Dirac, and Ward-completed gauge sectors, the localized
Hadamard-parametrix contribution satisfies
\eqref{eq:v2v44-order-four} in every normal chart, with finite \(L\).  The
state-dependent cross terms have order at most two, and the smooth terms
have order minus infinity.  On a bounded-geometry slab the constants are
uniform over the finite atlas, provided the state-remainder seminorms in
\cref{prop:v2v44-state} are uniform.
\end{theorem}

\begin{proof}
In normal coordinates, each Hadamard coefficient is smooth with seminorms
controlled by finitely many metric, curvature, mass, and connection jets.
The singular scalar model distributions are boundary values of powers and
logarithms of the squared geodesic distance.  Products and the stress
derivatives in \cref{lem:v2v44-table} give almost-homogeneous distributions
of scaling degree at most eight.  Their localized Fourier transforms are
log-polyhomogeneous symbols of order \(8-4=4\); differentiating in the dual
variable lowers symbol order by one.  This gives
\eqref{eq:v2v44-order-four}.  The degree-six cross terms from
\cref{prop:v2v44-state} give order at most \(6-4=2\), and a localized
smooth kernel has rapidly decreasing Fourier transform.  A finite
partition of unity and the bounded-geometry seminorms make the constants
uniform.
\end{proof}

\begin{firewall}[The symbol estimate needs uniform state seminorms]
\label{fw:v2v44-state}
The Hadamard wavefront condition is qualitative.  A sequence of Hadamard
states can have unbounded smooth remainders \(S_a\).  Uniformity in a
regulator or nonlinear iteration therefore requires explicit bounds on the
finite \(C^r\) seminorms used in \cref{prop:v2v44-state}; the microlocal
spectrum condition alone does not provide them.
\end{firewall}

\subsection{Three subtractions produce an integrable symbol}
\label{sec:v2v44-minus-two}

Choose a smooth high relative-frequency cutoff that vanishes for
\(|\omega|\le\mu\).  Apply three causal subtractions to the high-frequency
amplitude and put every polynomial through degree four into the local V40
block.  Before the exterior sixth-order operator is restored, define
\begin{equation}
 b_{\kappa,3}(u,x,y,\omega)
 =\chi_{\rm hi}(\omega)(i\omega)^{-6}
 a_\kappa^{\rm rem}(u,x,y,\omega).
\label{eq:v2v44-b}
\end{equation}

\begin{theorem}[Order-minus-two memory symbol]
\label{thm:v2v44-minus-two}
Under \eqref{eq:v2v44-order-four}, for every required finite collection of
derivatives,
\begin{equation}
 \|\partial_\omega^j b_{\kappa,3}(u,x,y,\omega)\|
 \le C'_{j,\kappa}
 \langle\omega\rangle^{-2-j}
 (1+\log\langle\omega\rangle)^L.
\label{eq:v2v44-minus-two}
\end{equation}
In particular,
\begin{equation}
 \sup_{u,x,y}
 \int_{\mathbb R}\|b_{\kappa,3}(u,x,y,\omega)\|\,\mathrm d\omega
 <\infty.
\label{eq:v2v44-L1-symbol}
\end{equation}
The inverse relative-time kernel
\begin{equation}
 g_{\kappa,3}(u,x,y,r)
 =\frac1{2\pi}\int_{\mathbb R}
 e^{ir\omega}b_{\kappa,3}(u,x,y,\omega)\,\mathrm d\omega
\label{eq:v2v44-g}
\end{equation}
is therefore bounded uniformly in \(r\).
\end{theorem}

\begin{proof}
On the high-frequency support, division by \((i\omega)^6\) lowers the
order-four symbol to order minus two.  Leibniz' rule shows that every
\(\omega\)-derivative has the bound
\eqref{eq:v2v44-minus-two}; derivatives of the cutoff are compactly
supported.  The function
\(\langle\omega\rangle^{-2}(1+\log\langle\omega\rangle)^L\) is integrable
on \(\mathbb R\), proving \eqref{eq:v2v44-L1-symbol}.  The Bochner Fourier
integral then gives
\(\|g_{\kappa,3}\|\le(2\pi)^{-1}\|b_{\kappa,3}\|_{L^1_\omega}\).
\end{proof}

\begin{corollary}[Curved free-sector Volterra bound]
\label{cor:v2v44-Volterra}
Assume the spatially localized amplitudes in
\cref{thm:v2v44-minus-two} define bounded operators on
\(H^s_{\rm red}\), uniformly in the central time, and that the exterior
sixth-order operator is causally order-reduced with a bounded remainder.
Then the nonlocal free stress response has a bounded causal kernel on the
slab.  Composing it with the V35 hybrid Einstein propagator gives the V40
Volterra estimate with \(\rho=1\).
\end{corollary}

\begin{proof}
The operator-valued version of \eqref{eq:v2v44-L1-symbol} and Fourier
inversion give a uniform bound for the relative-time kernel.  Multiplication
by the bounded order-reduction operator changes only the constant.  The
retarded support restricts to \(t>t'\), and the composition calculation in
\eqref{eq:v2v41-composed-kernel,eq:v2v41-composed-bound} applies.
\end{proof}

\subsection{Regulator convergence at the symbol level}
\label{sec:v2v44-regulator}

\begin{theorem}[Dominated symbol convergence]
\label{thm:v2v44-regulator}
Let \(b_{\kappa,3,\Lambda}\) be regulated subtracted amplitudes.  Suppose
they satisfy the uniform order-minus-two majorant in
\eqref{eq:v2v44-minus-two} and converge almost everywhere in \(\omega\),
uniformly in the bounded geometric variables, to
\(b_{\kappa,3}\).  Then
\begin{equation}
 \sup_{u,x,y,r}
 \|g_{\kappa,3,\Lambda}(u,x,y,r)
       -g_{\kappa,3}(u,x,y,r)\|
 \longrightarrow0.
\label{eq:v2v44-kernel-convergence}
\end{equation}
If the corresponding spatial operator kernels converge in the same
dominated topology and the effective Einstein propagators converge, the
complete Volterra operators converge in the norm required by
\eqref{eq:v2v40-K-convergence}.
\end{theorem}

\begin{proof}
The order-minus-two logarithmic majorant is integrable in \(\omega\).
Dominated convergence makes the \(L^1_\omega\) norm of the amplitude
difference tend to zero uniformly in the remaining variables.  Fourier
inversion gives \eqref{eq:v2v44-kernel-convergence}.  The asserted spatial
operator convergence and bounded causal composition then give the V40
operator-norm convergence.
\end{proof}

\subsection{What changes in perturbation theory}
\label{sec:v2v44-interacting}

At a fixed perturbative order, renormalized time-ordered products are again
extended from configuration space away from partial diagonals.  Power
counting and the Ward identities predict the same order-four complete
stress response, with a larger finite logarithmic degree.  V43 shows that
such logarithms do not change the three-subtraction threshold.

\begin{theorem}[Conditional fixed-order curved transfer]
\label{thm:v2v44-fixed-order}
Fix a perturbative order.  Suppose the Ward-completed renormalized stress
response on the bounded-geometry slab satisfies the uniform
log-polyhomogeneous estimate \eqref{eq:v2v44-order-four}, and suppose its
state-dependent smooth remainders have the uniform seminorms specified
above.  Then three subtractions give the order-minus-two estimate
\eqref{eq:v2v44-minus-two}, a bounded retarded memory kernel, and the V40
causal Volterra resolvent on the fixed moving spatial band.
\end{theorem}

\begin{proof}
The proof of \cref{thm:v2v44-minus-two} uses only the stated symbol
seminorms, not Gaussianity.  Apply that theorem, then
\cref{cor:v2v44-Volterra,thm:v2v40-Volterra}.
\end{proof}

\begin{firewall}[The interacting symbol estimate is not proved here]
\label{fw:v2v44-interacting}
V44 proves the curved estimate for free Hadamard sectors and a conditional
fixed-order transfer.  It does not construct interacting time-ordered
products with regulator-uniform Ward identities and symbol seminorms, nor
does it control their sum over perturbative orders.  Overlapping partial
diagonals require a complete causal-renormalization argument; local scaling
degree alone does not provide the uniform constants or the nonperturbative
Dini majorant of V43.
\end{firewall}

\subsection{V44 conclusion and next step}
\label{sec:v2v44-conclusion}

The free curved stress response has universal diagonal scaling degree at
most eight in every scalar, Dirac, and Ward-completed gauge sector.  Its
local ambiguity has order at most four.  On a fixed spatial band, the
nonlocal relative-time amplitude is log-polyhomogeneous of order four;
three subtractions leave an order-minus-two frequency symbol and hence a
bounded causal memory kernel.  Smooth state dependence is lower order but
must be uniformly bounded across regulators.

V45 is the next mandatory companion audit.  It will then formulate the
overlapping-diagonal forest estimate needed to make the V44 fixed-order
symbol bound uniform over renormalized SM graphs, while keeping the V43
nonperturbative summability gate explicit.
\clearpage
\section*{Second Volume, Version V45: Companion Audit and a Fixed-Order Renormalization-Forest Symbol Theorem}
\addcontentsline{toc}{section}{Second Volume, Version V45: Companion Audit and a Fixed-Order Renormalization-Forest Symbol Theorem}

\subsection*{Purpose and mandatory companion audit}

This is the required audit five versions after V40.  The current companion
sources inspected were:
\begin{center}
\begin{tabular}{p{0.46\textwidth}r p{0.33\textwidth}}
\toprule
file & bytes & SHA--256\\
\midrule
\texttt{Renewal\_NS\_Stability\_Research\_Notes\_V31.tex}
&887537&
\texttt{F1121B62238AD5491BB7CF03635EA9934942EDF573ED8FAAA9EE79E1D38A6696}\\
\texttt{Renewal\_Yang\_Mills\_Mass\_Gap\_Research\_Notes\_V36.tex}
&484962&
\texttt{090039B90E3D803D1E0B352D7402D47FDE773CDC75092F58D960BF00A2416570}\\
\bottomrule
\end{tabular}
\end{center}
The NS note is unchanged.  The YM programme has advanced to V36.  V35
constructs a one-use first-entry automaton for the \([3,3]\) sector.  V36
then proves an exact first-four-block telescope, but a scalar Rademacher
regression shows that the restricted factor \(R=M-K_4\) need not be a
contraction and that a product of such factors can grow exponentially.
Accordingly, V36 withdraws the global reading of its earlier V34
embedded-\([4]\) statement.  The valid route returns upstream to the
complete connected quartic operator before the norm and reduces its
normalized global estimate to finite graph-cut bounds.  Global quartic and
all-order MTA, continuum Yang--Mills, and a mass gap remain open.  The
parallel notes were left untouched.

\begin{assessment}[V45 audit transfer]
\label{ass:v2v45-audit}
The analogue for SM renormalization is exact.  A subtraction operator
\((I-T_\gamma)\) or a restricted graph family need not be a contraction
when normed separately.  One must first form the complete forest-renormalized
amplitude, including all compatible counterterm insertions and the Ward
completion, and only then take the physical Sobolev symbol norm.  A raw sum
over overlapping divergent-subgraph witnesses risks the same multiplicity
as the invalid YM witness-pair sum.
\end{assessment}

\subsection{Compatible forests and exact renormalized assembly}
\label{sec:v2v45-forest}

Fix a connected stress-response graph \(\Gamma\) at a finite perturbative
order.  Let \(\mathfrak D(\Gamma)\) be its proper superficially divergent
subgraphs.  Two members are compatible if they are disjoint or one contains
the other.  A forest is a pairwise compatible subset
\(F\subset\mathfrak D(\Gamma)\).  Let \(T_\gamma\) denote the local
covariant Taylor or extension projector through the superficial degree of
\(\gamma\).

\begin{definition}[Forest-renormalized amplitude]
\label{def:v2v45-R}
The assembled amplitude is
\begin{equation}
 R_\Gamma I_\Gamma
 =(I-T_\Gamma)\sum_{F\in\mathfrak F(\Gamma)}
 \left(\prod_{\gamma\in F}^{\longrightarrow}(-T_\gamma)\right)
 I_\Gamma,
\label{eq:v2v45-forest}
\end{equation}
where nested Taylor maps are ordered from inner to outer subgraphs.  The
empty forest contributes \(I_\Gamma\).  If the overall graph is
superficially convergent, set \(T_\Gamma=0\).
\end{definition}

\begin{lemma}[No incompatible double subtraction]
\label{lem:v2v45-compatible}
Every term in \eqref{eq:v2v45-forest} corresponds to one compatible nesting
or disjoint collection of subtraction regions.  Two overlapping subgraphs
that are not nested never occur in the same product.  The forest formula is
therefore an algebraic assembly before norms, not a product of independent
absolute subtraction costs.
\end{lemma}

\begin{proof}
The first two claims are the definition of a forest.  Expanding the
Bogoliubov recursion by selecting, at each stage, a disjoint family of
maximal proper divergent subgraphs produces each compatible nested family
once, ordered by inclusion.  Conversely, removing the maximal members of a
forest gives exactly one term in that recursion.  This bijection proves
\eqref{eq:v2v45-forest} without assigning independent norms to the factors.
\end{proof}

\begin{remark}[Finite order versus all orders]
\label{rem:v2v45-finite}
At a fixed perturbative order, a graph has finitely many divergent
subgraphs and hence finitely many forests.  The number of graphs and forests
may grow rapidly with perturbative order.  Nothing in the forest identity
provides a uniform all-order summability bound.
\end{remark}

\subsection{Taylor remainders and power-counting improvement}
\label{sec:v2v45-Taylor}

For a one-variable model, the exact Taylor remainder is
\begin{equation}
 (I-T^d)f(p)
 =\frac{p^{d+1}}{d!}
 \int_0^1(1-s)^d f^{(d+1)}(sp)\,\mathrm ds.
\label{eq:v2v45-Taylor}
\end{equation}
The covariant multivariable version has the same integral form with
multi-indices.  It is this identity, together with subgraph power counting,
that improves ultraviolet degree.  The separate operator \(I-T^d\) has no
reason to have norm at most one.

\begin{lemma}[Subtraction-factor regression]
\label{lem:v2v45-regression}
There is no abstract contraction estimate for \(I-T\) based only on
idempotence and a finite bound for \(\|T\|\).  On Euclidean \(\mathbb R^2\),
the operator
\begin{equation}
 T=\begin{pmatrix}1&\sqrt8\\0&0\end{pmatrix}
\label{eq:v2v45-T}
\end{equation}
satisfies \(T^2=T\) and
\begin{equation}
 \|T\|=\|I-T\|=3.
\label{eq:v2v45-regression}
\end{equation}
On the Hilbert tensor product of \(n\) independent copies,
\(\|(I-T)^{\otimes n}\|=3^n\).
\end{lemma}

\begin{proof}
Direct multiplication gives \(T^2=T\).  The matrix \(T\) has the one
nonzero singular value \((1+8)^{1/2}=3\).  The complement
\(I-T=\left(\begin{smallmatrix}0&-\sqrt8\\0&1\end{smallmatrix}\right)\)
also has the one nonzero singular value three.  Operator norms multiply on
Hilbert tensor products.  This genuine nonorthogonal projection example is
the linear-operator form of the YM V36 regression and shows why forest terms
must not be bounded by multiplying norms of restricted factors.
\end{proof}

\subsection{A conditional forest-to-symbol theorem}
\subsection{A conditional forest-to-symbol theorem}
\label{sec:v2v45-symbol}

Work in one of the bounded-geometry charts of V44, after the fixed spatial
low-band cutoffs have been inserted.  Let \(p\) denote the relative
four-momentum and let \(\omega(\gamma)\) be the superficial degree of a
subgraph.  The following hypotheses are the local form of Weinberg power
counting with parameters:
\begin{enumerate}
 \item every propagator and vertex amplitude has uniform differentiated
       symbol bounds in the background variables;
 \item after all proper divergent subgraphs of a graph are renormalized,
       every scaling region associated with a superficially convergent
       subgraph has a strictly negative degree;
 \item the covariant Taylor maps \(T_\gamma\) are bounded on the finite jet
       seminorms and produce only local vertices of the allowed
       renormalized action;
 \item all estimates are uniform over the compact background parameter
       set and the fixed external spatial band.
\end{enumerate}

\begin{theorem}[Forest power counting with external parameters]
\label{thm:v2v45-forest-symbol}
Under the preceding hypotheses, the assembled renormalized amplitude of a
fixed graph obeys
\begin{equation}
 \|\partial_p^\alpha R_\Gamma I_\Gamma(p)\|
 \le C_{\Gamma,\alpha}
 \langle p\rangle^{\omega(\Gamma)-|\alpha|}
 (1+\log\langle p\rangle)^{L_\Gamma}
\label{eq:v2v45-graph-symbol}
\end{equation}
for a finite logarithmic degree \(L_\Gamma\).  The constants are uniform in
the stated background parameters.  Local counterterm terms are included in
the assembled expression and are not normed as separate subtraction
factors.
\end{theorem}

\begin{proof}
Induct on the number of proper divergent subgraphs.  For a primitive graph,
the asserted degree is the parameter-uniform power-counting hypothesis;
radial integration at equality can add a finite logarithm.  For a graph
with subdivergences, use \cref{lem:v2v45-compatible} to group each forest by
its maximal disjoint subgraphs.  By induction, replacing such a subgraph by
its assembled local Taylor jet produces an allowed vertex with its stated
degree and uniform seminorms.

In every scaling region where a proper subgraph momentum becomes large,
the integral Taylor formula \eqref{eq:v2v45-Taylor} supplies the exact
remainder power required to make the already-renormalized subgraph degree
negative.  The second hypothesis then makes that region integrable.  The
region where all loop momenta scale together has the overall degree
\(\omega(\Gamma)\).  Differentiating in external momentum lowers that
degree by \(|\alpha|\); borderline radial integrations contribute only a
finite logarithmic power.  There are finitely many scaling regions and
forests at fixed graph order, so their assembled bounds give
\eqref{eq:v2v45-graph-symbol}.  Uniform parameter seminorms permit all
derivatives and finite sums with the same background-independent constant.
\end{proof}

\begin{firewall}[The theorem assumes the parameter-uniform Weinberg bounds]
\label{fw:v2v45-Weinberg}
V45 records the exact induction needed to pass from subgraph estimates to
the complete graph symbol.  It does not prove the stated Weinberg bounds
for every SM graph on a time-dependent curved background.  In particular,
massless infrared regions, gauge degeneracy, and partial-diagonal
extensions require separate estimates.  The theorem prevents a false
factorwise norm argument; it does not replace those analytic inputs.
\end{firewall}

\subsection{Stress response degree and Ward completion}
\label{sec:v2v45-stress}

For two insertions of a dimension-four stress tensor in four spacetime
dimensions, the overall position-space scaling degree is eight and the
Fourier degree is four, as proved for free sectors in V44.  Renormalizable
interaction vertices do not change the overall stress-response degree when
the subgraph hypotheses of \cref{thm:v2v45-forest-symbol} hold.

Let \(\mathfrak G_r\) be the finite set of graphs, including ghosts and
counterterm graphs, at perturbative order \(r\).  Define the complete Ward
sum
\begin{equation}
 \mathcal A_r^{\rm phys}
 =Q_{\rm phys}
 \left(\sum_{\Gamma\in\mathfrak G_r}
 R_\Gamma I_\Gamma\right)Q_{\rm phys}.
\label{eq:v2v45-Ward-sum}
\end{equation}

\begin{theorem}[Fixed-order Ward-completed order-four symbol]
\label{thm:v2v45-Ward-symbol}
Fix \(r\).  Suppose every graph in \(\mathfrak G_r\) satisfies the
hypotheses of \cref{thm:v2v45-forest-symbol}, and suppose the local
counterterms have been chosen so that the complete sum satisfies the
renormalized Ward identity.  Then, on the physical quotient,
\begin{equation}
 \|\partial_p^\alpha\mathcal A_r^{\rm phys}(p)\|
 \le C_{r,\alpha}
 \langle p\rangle^{4-|\alpha|}
 (1+\log\langle p\rangle)^{L_r}.
\label{eq:v2v45-order-four}
\end{equation}
Consequently three high-frequency subtractions give the V44
order-minus-two memory symbol at perturbative order \(r\).
\end{theorem}

\begin{proof}
Apply \cref{thm:v2v45-forest-symbol} to every graph.  The overall degree is
at most four.  Sum the finitely many amplitudes and local counterterms
inside \eqref{eq:v2v45-Ward-sum}, then apply the bounded physical
projection.  Taking the maximum logarithmic degree and summing the finite
constants gives \eqref{eq:v2v45-order-four}.  Division of the high-frequency
remainder by \((i\omega)^6\) lowers order four to order minus two, as in
\cref{thm:v2v44-minus-two}.
\end{proof}

\begin{proposition}[Why the Ward sum precedes graph norms]
\label{prop:v2v45-cancellation}
Termwise symbol norms can lose a Ward cancellation.  Indeed, for any
operator-valued high-order symbol \(L(p)\) and lower-order symbol \(S(p)\),
the two representatives
\begin{equation}
 A_1(p)=L(p),
 \qquad
 A_2(p)=-L(p)+S(p)
\label{eq:v2v45-cancellation}
\end{equation}
have a sum of the order of \(S\), while
\(\|A_1\|+\|A_2\|\) retains the order of \(L\).
\end{proposition}

\begin{proof}
The assembled sum is exactly \(S\).  The triangle estimate taken before
addition has no access to the cancellation of \(L\).
\end{proof}

In a gauge theory, the longitudinal and ghost representatives can behave
like \(A_1,A_2\).  The physical estimate is therefore taken only after the
renormalized Ward sum is formed.  Positivity, when used for the spectral
Dini theorem of V43, also applies only at that stage.

\subsection{Finite-order causal consequence}
\label{sec:v2v45-causal}

\begin{corollary}[Forest criterion for the causal Volterra branch]
\label{cor:v2v45-causal}
At a fixed perturbative order, assume:
\begin{enumerate}
 \item the parameter-uniform forest hypotheses of
       \cref{thm:v2v45-forest-symbol};
 \item the Ward completion of \cref{thm:v2v45-Ward-symbol};
 \item uniform state-remainder seminorms as in
       \cref{fw:v2v44-state};
 \item the V39 moving-band crossing bound and convergence of the effective
       Einstein propagator.
\end{enumerate}
Then three subtractions give a bounded gauge-reduced nonlocal memory kernel,
and the complete fixed-order low-energy Einstein--SM equation has the V40
causal Volterra resolvent and regulator limit on every controlled finite
slab.
\end{corollary}

\begin{proof}
The forest and Ward theorems give the uniform order-four amplitude.  V44
turns it into an order-minus-two integrable symbol and a bounded memory
kernel.  V39 controls the moving band, and
\cref{thm:v2v40-Volterra,thm:v2v40-regulator} give the causal solution and
regulator convergence.
\end{proof}

\subsection{Why this does not close the all-order programme}
\label{sec:v2v45-all-order}

\begin{theorem}[Exact all-order majorant needed]
\label{thm:v2v45-majorant}
Let \(C_r\) denote a valid constant in
\eqref{eq:v2v45-order-four}, enlarged to include the finite graph and
forest sums at order \(r\).  A sufficient condition for an absolutely
summed order-four response at coupling \(g\) is
\begin{equation}
 \sum_{r=0}^\infty |g|^rC_r<\infty
\label{eq:v2v45-majorant}
\end{equation}
with a common logarithmic or Dini majorant of the type in
\cref{thm:v2v43-resummed}.  Fixed-order forest bounds alone do not imply
\eqref{eq:v2v45-majorant}.
\end{theorem}

\begin{proof}
Under \eqref{eq:v2v45-majorant}, the Weierstrass test sums the physical
symbol bounds and their required derivatives uniformly.  The common Dini
majorant permits the V43 weighted spectral or V44 frequency integration.
Conversely, the forest theorem places no restriction on the growth of
\(C_r\); factorial graph proliferation is compatible with every one of its
fixed-order conclusions.  Therefore no all-order sum follows without an
additional summability construction.
\end{proof}

\begin{firewall}[The full SM--Einstein continuum remains open]
\label{fw:v2v45-status}
V45 supplies a correct fixed-order assembly theorem and rules out
factorwise subtraction norms.  It does not prove uniform Weinberg estimates
for the full curved SM, BRST restoration at every order, convergence or
Borel summability of the graph series, the nonperturbative Dini majorant,
or the nonlinear a priori estimates needed to continue the Einstein
geometry for arbitrary time.  These remain distinct gates.
\end{firewall}

\subsection{V45 conclusion and next step}
\label{sec:v2v45-conclusion}

The companion audit records the YM V36 correction: restricted products can
grow even when the complete moment is controlled, so the complete connected
operator must be assembled before norms.  Applied to SM renormalization,
this yields the compatible-forest rule and a conditional fixed-order
theorem: parameter-uniform Weinberg bounds plus Ward completion give the
order-four curved stress symbol, and three subtractions give the bounded
V40 memory kernel.

The remaining obstruction has moved to uniformity in perturbative order.
V46 will examine a Borel-majorant route: factorially weighted graph bounds,
sectorial Borel transforms, and the conditions under which a Borel sum
preserves Ward identities, the V43 Dini moment, and causal retarded support.
\clearpage
\section*{Second Volume, Version V46: A Banach-Valued Borel Route and Its Nonperturbative Obstruction}
\addcontentsline{toc}{section}{Second Volume, Version V46: A Banach-Valued Borel Route and Its Nonperturbative Obstruction}

\subsection*{Purpose}

V45 reduced the all-order perturbative problem to growth of the complete
Ward-assembled forest constants.  Factorial growth is expected and rules
out ordinary absolute summation.  This note formulates a Borel alternative
in a response Banach space that already contains the three pieces required
by the causal continuum argument: local contact coefficients, the
three-subtracted ultraviolet amplitude, and the retarded Volterra kernel.

The main theorem shows that a sectorially continuable Banach-valued Borel
transform preserves Ward identities, retarded support, and the V43--V44
Dini/symbol norm.  Uniform Borel convergence also passes through the
ultraviolet regulator limit.  A sharp scalar example then shows why a
factorial graph bound alone is insufficient: its Borel transform can have a
pole on the positive integration ray, producing an exponentially small
ambiguity invisible to every perturbative coefficient.

\subsection{The causal response Banach space}
\label{sec:v2v46-space}

Fix a physical spatial band \(\kappa\) and a finite slab \([0,T]\).  Let
\(\mathcal C_4\) be the finite-dimensional space of local physical contact
operators of differential order at most four.  Let \(\mathcal A_1\) be the
space of countably additive operator-valued measures
\(b(u,\mathrm d\omega)\) of bounded variation, uniformly in the background
parameter \(u\), with norm
\begin{equation}
 \|b\|_{\mathcal A_1}
 =\sup_{u\in[0,T]}
 \int_{\mathbb R}\mathrm d
 |b(u,\cdot)|_{H^s_{\rm red}\to H^s_{\rm red}}(\omega).
\label{eq:v2v46-A1}
\end{equation}
The V44 amplitudes enter as absolutely continuous measures
\(b(u,\omega)\,\mathrm d\omega\), while the V42 compact-space spectra may
be atomic.  The finitely many background and spatial operator seminorms
needed for composition are included in the variation norm.  Let
\(\mathcal V_T\) be the Banach
space of causal kernels with
\begin{equation}
 \|k\|_{\mathcal V_T}
 =\sup_{0\le t\le T}
 \int_0^t\|k(t,\tau)\|_{H^s_{\rm red}\to H^s_{\rm red}}
 \,\mathrm d\tau.
\label{eq:v2v46-V}
\end{equation}
Define
\begin{equation}
 \mathfrak B_{\kappa,T}
 =\mathcal C_4\oplus\mathcal A_1\oplus\mathcal V_T
\label{eq:v2v46-B}
\end{equation}
with the sum norm.

\begin{lemma}[Completeness and continuous reconstruction]
\label{lem:v2v46-complete}
The space \(\mathfrak B_{\kappa,T}\) is Banach.  Fourier--Stieltjes reconstruction of the
\(\mathcal A_1\) component, causal restriction, bounded low-band order
reduction, and composition with a uniformly bounded retarded Einstein
propagator define continuous linear maps into \(\mathcal V_T\).
\end{lemma}

\begin{proof}
The contact space is finite dimensional.  The bounded-variation measure space with values in
bounded operators is Banach, as is the essential-supremum space of its
background-parameter families.  The causal kernel norm is the operator
kernel form of an \(L^\infty_t\) Volterra norm and is complete after
identification almost everywhere.  Finite direct sums of Banach spaces are
Banach.

Fourier--Stieltjes reconstruction has pointwise operator norm at most
\((2\pi)^{-1}\|b\|_{\mathcal A_1}\).  Causal restriction does not increase
that bound.  Bounded order reduction changes it by its operator norm, and
the causal composition estimate of
\eqref{eq:v2v41-composed-kernel,eq:v2v41-composed-bound} is linear and
bounded on a finite slab.
\end{proof}

\begin{lemma}[Dini measures embed continuously]
\label{lem:v2v46-Dini}
Let \(\Sigma\) be a signed finite-dimensional operator measure on
\([\mu,\infty)\) with
\begin{equation}
 \|\Sigma\|_{D,3}
 :=\int_{[\mu,\infty)}\Omega^{-7}
 \,\mathrm d|\Sigma|(\Omega)<\infty.
\label{eq:v2v46-Dini}
\end{equation}
Then its three-subtracted sine kernel belongs to the reconstructed
\(\mathcal A_1\)-to-\(\mathcal V_T\) response class, with norm bounded by a
constant times \(\|\Sigma\|_{D,3}\).
\end{lemma}

\begin{proof}
The total-variation integral dominates
\(\int\Omega^{-7}\sin(\Omega r)\,\mathrm d\Sigma\) uniformly in \(r\).
The V43 Dini identity shows that this is the same weighted moment controlled
by the cumulative condition.  Bounded order reduction and retarded
composition then use \cref{lem:v2v46-complete}.
\end{proof}

Let \(\mathcal W:\mathfrak B_{\kappa,T}\to\mathfrak D\) be the bounded
linear Ward-defect map into a defect Banach space.  Let
\(\mathfrak B_{\rm ret}\) be the closed subspace whose reconstructed
distribution has retarded support.  Define the physical causal subspace
\begin{equation}
 \mathfrak B_{\rm phys,ret}
 =\ker\mathcal W\cap\mathfrak B_{\rm ret}.
\label{eq:v2v46-physical}
\end{equation}

\begin{lemma}[Closedness of the causal Ward conditions]
\label{lem:v2v46-closed}
The space \(\mathfrak B_{\rm phys,ret}\) is closed in
\(\mathfrak B_{\kappa,T}\).
\end{lemma}

\begin{proof}
The kernel of a bounded linear map is closed.  Convergence in the response
norm gives distributional convergence of reconstructed kernels.  Support
in the closed retarded cone is preserved under distributional limits, so
\(\mathfrak B_{\rm ret}\) is closed.  Their intersection is closed.
\end{proof}

\subsection{Banach-valued Borel summation}
\label{sec:v2v46-Borel}

Use one complex bookkeeping coupling \(g\) along a fixed ray in the
finite-dimensional SM coupling space.  A multidirectional construction
requires a corresponding multivariate Borel domain and is not hidden in
this notation.  Let
\begin{equation}
 \mathcal R(g)\sim\sum_{n=0}^\infty g^n\mathcal R_n,
 \qquad
 \mathcal R_n\in\mathfrak B_{\rm phys,ret}.
\label{eq:v2v46-formal}
\end{equation}

\begin{theorem}[Sectorial Banach-valued Borel sum]
\label{thm:v2v46-Borel}
Assume the Gevrey-one estimate
\begin{equation}
 \|\mathcal R_n\|_{\mathfrak B}
 \le C_0A^n n!
\label{eq:v2v46-Gevrey}
\end{equation}
and define the local Borel transform
\begin{equation}
 \widehat{\mathcal R}(\zeta)
 =\sum_{n=0}^\infty\frac{\mathcal R_n}{n!}\zeta^n.
\label{eq:v2v46-Borel-transform}
\end{equation}
Suppose it admits analytic continuation to a sector containing the positive
real axis and, on that axis,
\begin{equation}
 \|\widehat{\mathcal R}(\zeta)\|_{\mathfrak B}
 \le C_Be^{b\zeta},
 \qquad \zeta\ge0.
\label{eq:v2v46-exponential}
\end{equation}
Then, for real \(0<g<b^{-1}\) when \(b>0\), the Bochner integral
\begin{equation}
 \mathcal R_B(g)
 =\frac1g\int_0^\infty
 e^{-\zeta/g}\widehat{\mathcal R}(\zeta)\,\mathrm d\zeta
\label{eq:v2v46-Laplace}
\end{equation}
exists in \(\mathfrak B_{\kappa,T}\) and obeys
\begin{equation}
 \|\mathcal R_B(g)\|_{\mathfrak B}
 \le\frac{C_B}{1-bg}.
\label{eq:v2v46-Borel-bound}
\end{equation}
It belongs to \(\mathfrak B_{\rm phys,ret}\) and has
\eqref{eq:v2v46-formal} as its Gevrey-one asymptotic expansion in every
closed subsector where the same analytic continuation and exponential bound
hold.
\end{theorem}

\begin{proof}
Estimate the integrand by
\(C_Bg^{-1}e^{-(g^{-1}-b)\zeta}\).  Its scalar integral is
\(C_B/(1-bg)\), proving Bochner integrability and
\eqref{eq:v2v46-Borel-bound}.

In the disk \(|\zeta|<A^{-1}\), every coefficient lies in the closed
subspace \(\mathfrak B_{\rm phys,ret}\), so the convergent Borel series does
also.  The Ward defect of the analytic continuation vanishes by the
identity theorem, and the continued function remains in the closed
retarded subspace by continuation through overlapping Banach-valued power
series.  The Bochner integral of a closed-subspace-valued function remains
in that subspace.

For the asymptotic statement, subtract the first \(N\) terms of
\eqref{eq:v2v46-Borel-transform} inside \eqref{eq:v2v46-Laplace}.  The
monomial identity
\begin{equation}
 \frac1g\int_0^\infty e^{-\zeta/g}\zeta^n\,\mathrm d\zeta
 =n!g^n
\label{eq:v2v46-monomial}
\end{equation}
recovers the partial perturbation series.  Cauchy estimates near the origin
and the sectorial exponential bound on the remaining contour give
\(C'A'^NN!|g|^N\) for the remainder, which is the Gevrey-one asymptotic
condition.
\end{proof}

\begin{corollary}[Ward identity, causality, and subtraction commute]
\label{cor:v2v46-commute}
Under \cref{thm:v2v46-Borel},
\begin{equation}
 \mathcal W\mathcal R_B(g)=0,
 \qquad
 \operatorname{supp}\mathcal R_B(g)\subseteq\{t\ge\tau\}.
\label{eq:v2v46-Ward-causal}
\end{equation}
Every continuous linear contact-separation, three-subtraction, Fourier
reconstruction, or retarded-composition map commutes with the Borel--Laplace
integral.
\end{corollary}

\begin{proof}
The first two claims are part of the closed-subspace conclusion of
\cref{thm:v2v46-Borel}.  A continuous linear map passes through a Bochner
integral, proving the final claim.
\end{proof}

\subsection{Preservation of the Dini and Volterra estimates}
\label{sec:v2v46-Dini-preservation}

\begin{theorem}[Borel-summed causal memory bound]
\label{thm:v2v46-memory}
Assume the coefficients in \eqref{eq:v2v46-formal} are represented in the
three-subtracted Dini or V44 symbol component of
\(\mathfrak B_{\kappa,T}\), and the hypotheses of
\cref{thm:v2v46-Borel} hold in that norm.  Then the Borel-summed response has
a finite reconstructed Volterra norm.  If its composed kernel is denoted
\(K_B(g)\), then
\begin{equation}
 \|K_B(g)\|_{\mathcal V_T}
 \le C_{\kappa,T}\frac{C_B}{1-bg}.
\label{eq:v2v46-memory}
\end{equation}
The causal equation \(h=h_0+K_B(g)h\) has the V40 Volterra resolvent on
every finite slab, without requiring the right side of
\eqref{eq:v2v46-memory} to be less than one.
\end{theorem}

\begin{proof}
The reconstructed-kernel map is continuous by
\cref{lem:v2v46-complete,lem:v2v46-Dini}.  Apply it to
\eqref{eq:v2v46-Borel-bound}.  Retarded support follows from
\cref{cor:v2v46-commute}.  The V40 fractional Volterra theorem then gives
the resolvent; bounded memory corresponds to its exponent \(\rho=1\).
\end{proof}

\subsection{Uniform removal of the ultraviolet regulator}
\label{sec:v2v46-regulator}

Let \(\widehat{\mathcal R}_\Lambda\) be regulated Borel transforms.

\begin{theorem}[Regulator limit through the Borel integral]
\label{thm:v2v46-regulator}
Suppose all \(\widehat{\mathcal R}_\Lambda\) are analytic in one common
sector, obey the uniform bound
\begin{equation}
 \|\widehat{\mathcal R}_\Lambda(\zeta)\|_{\mathfrak B}
 \le C_Be^{b\zeta},
 \qquad\zeta\ge0,
\label{eq:v2v46-uniform-exponential}
\end{equation}
and converge in \(\mathfrak B\) at every \(\zeta\ge0\) to
\(\widehat{\mathcal R}(\zeta)\).  Then, for every fixed
\(0<g<b^{-1}\),
\begin{equation}
 \|\mathcal R_{B,\Lambda}(g)-\mathcal R_B(g)\|_{\mathfrak B}
 \longrightarrow0.
\label{eq:v2v46-regulator}
\end{equation}
Consequently the reconstructed Volterra kernels and their linear causal
solutions converge in the V40 norms, provided the effective Einstein
propagators and sources converge as in V40.
\end{theorem}

\begin{proof}
The difference of the Borel integrands tends to zero pointwise and is
bounded by \(2C_Bg^{-1}e^{-(g^{-1}-b)\zeta}\), which is integrable.
Banach-valued dominated convergence gives
\eqref{eq:v2v46-regulator}.  Continuous reconstruction and
\cref{thm:v2v40-regulator} give the causal solution convergence.
\end{proof}

\subsection{Positivity is not a linear identity}
\label{sec:v2v46-positivity}

\begin{proposition}[A sufficient positivity condition]
\label{prop:v2v46-positive}
Suppose the spectral-measure component of
\(\widehat{\mathcal R}(\zeta)\) is a positive operator-valued measure for
every \(\zeta\ge0\).  Then its Borel--Laplace integral is a positive
operator-valued measure and the V43 positive Dini theorem applies.
Coefficientwise Ward identities and retarded support alone do not imply
this positivity hypothesis.
\end{proposition}

\begin{proof}
For every physical vector \(v\) and measurable frequency set \(E\), the
scalar integrand
\(\langle v,\widehat\Sigma(\zeta;E)v\rangle\) is nonnegative.  Integration
against the positive weight \(g^{-1}e^{-\zeta/g}\) preserves
nonnegativity.  Ward identities and support constraints are linear
equalities and place no sign restriction on this quadratic form.
\end{proof}

Thus Borel summation automatically preserves the linear structural
conditions needed for gauge reduction and causality, but physical
positivity or unitarity needs a compatible nonperturbative construction.

\subsection{A positive-axis obstruction}
\label{sec:v2v46-obstruction}

\begin{proposition}[Factorial control does not imply Borel summability]
\label{prop:v2v46-pole}
Let \(x\ne0\) belong to any Banach space and set
\begin{equation}
 \mathcal R_n=A^n n!\,x,
 \qquad A>0.
\label{eq:v2v46-pole-coefficients}
\end{equation}
Then \eqref{eq:v2v46-Gevrey} holds, but
\begin{equation}
 \widehat{\mathcal R}(\zeta)=\frac{x}{1-A\zeta}
\label{eq:v2v46-pole}
\end{equation}
has a pole on the positive Borel ray.  The ordinary integral
\eqref{eq:v2v46-Laplace} is undefined.  Upper and lower lateral contours
differ by a nonzero multiple of
\begin{equation}
 g^{-1}e^{-1/(Ag)}x.
\label{eq:v2v46-ambiguity}
\end{equation}
\end{proposition}

\begin{proof}
The Borel series is geometric and sums to
\eqref{eq:v2v46-pole} inside \(|\zeta|<A^{-1}\).  Its pole lies on the
integration path.  Deforming the path above or below the pole changes the
integral by \(2\pi i\) times the residue of
\(g^{-1}e^{-\zeta/g}x/(1-A\zeta)\), which is a nonzero constant multiple of
\eqref{eq:v2v46-ambiguity}.  Every Taylor coefficient at \(g=0\) misses
this exponentially small difference.
\end{proof}

\begin{firewall}[A lateral prescription is new physical data]
\label{fw:v2v46-lateral}
When the physical Borel ray contains singularities, choosing an upper,
lower, principal-value, or transseries completion is not determined by the
perturbative coefficients.  The choice must be fixed by additional
nonperturbative data and then checked for Ward identities, causality,
reality, positivity, and the Dini response bound.  Calling a lateral sum
``the continuum'' without those checks would leave the principal
construction ambiguous by terms such as \eqref{eq:v2v46-ambiguity}.
\end{firewall}

\subsection{Conditional insertion into the SM--Einstein branch}
\label{sec:v2v46-insertion}

\begin{theorem}[Borel route to the linear causal continuum]
\label{thm:v2v46-continuum}
Assume:
\begin{enumerate}
 \item the V39 moving-band and regulator-stable crossing hypotheses;
 \item the V40 local hyperbolic absorption and constraint propagation;
 \item the V45 forest coefficients satisfy the Gevrey bound
       \eqref{eq:v2v46-Gevrey} in \(\mathfrak B_{\kappa,T}\);
 \item their Borel transforms satisfy the common sectorial continuation,
       exponential bound, and regulator convergence of
       \cref{thm:v2v46-Borel,thm:v2v46-regulator};
 \item any positive-axis ambiguity has been fixed by a physical completion
       satisfying the same Ward, causal, Dini, reality, and positivity
       conditions.
\end{enumerate}
Then the Borel-completed gauge-reduced linear low-energy Einstein--SM
response has a regulator-independent causal Volterra solution on every
finite controlled slab.
\end{theorem}

\begin{proof}
The Borel theorem gives a physical retarded response in the complete
response norm.  The memory theorem reconstructs its Volterra kernel, and
the regulator theorem passes to the ultraviolet limit.  V39 controls the
moving band and V40 supplies the causal resolvent and constraint
propagation.  The fifth hypothesis removes the ambiguity exhibited by
\cref{prop:v2v46-pole} and supplies the non-linear sign properties not
implied by the Banach-valued sum.
\end{proof}

\begin{firewall}[The hypotheses are not established for the SM]
\label{fw:v2v46-status}
V46 proves a functional-analytic Borel implication.  It does not prove
Gevrey-one forest bounds in the V43 Dini/V44 symbol norm, analytic
continuation of the SM Borel transform, absence or cancellation of
positive-axis renormalons, spectral positivity of a chosen completion, or
uniform Lipschitz dependence on the evolving metric.  The theorem therefore
does not constitute a nonperturbative SM construction.
\end{firewall}

\subsection{V46 conclusion and next step}
\label{sec:v2v46-conclusion}

A sectorial Banach-valued Borel sum would preserve the three-subtracted
response norm, Ward identities, causal support, and regulator convergence.
These are closed linear or norm conditions.  Positivity requires a
positive measure-valued Borel continuation, and a singularity on the
physical Borel ray creates exponentially small completion data not fixed
by perturbation theory.  Factorial graph bounds alone therefore do not
close the continuum programme.

V47 will analyze the ambiguity itself.  It will formulate the space of
allowed exponentially small retarded additions, prove which local or
spectral conditions preserve the Ward and Dini ledgers, and determine
whether a normalization condition can select a unique completion without
silently importing the desired continuum measure.
\clearpage
\section*{Second Volume, Version V47: Exponentially Flat Response Ambiguities and a Uniqueness Criterion}
\addcontentsline{toc}{section}{Second Volume, Version V47: Exponentially Flat Response Ambiguities and a Uniqueness Criterion}

\subsection*{Purpose}

V46 showed that a singularity on the positive Borel ray produces an
exponentially small ambiguity.  Such a term has zero formal power series at
the origin, so no finite or infinite list of perturbative coefficients can
detect it.  This note determines which of the continuum ledgers constrain
that ambiguity.

Ward identities, retarded support, and the three-subtracted Dini norm are
closed linear or norm conditions.  They are preserved by any exponentially
small addition already lying in the physical causal response space.
Spectral positivity narrows the allowed cone but need not make it a single
point.  Finite local normalization conditions cannot fix an arbitrary
infinite-dimensional nonlocal ambiguity.  They do give uniqueness if a
separately justified finite-dimensional transseries space is supplied and
the resulting normalization matrix is invertible.

\subsection{The flat ambiguity space}
\label{sec:v2v47-flat}

Use the physical causal Banach space
\(\mathfrak B_{\rm phys,ret}\) of V46.  For a sector
\begin{equation}
 S_\theta=\{g\ne0:|\arg g|<\theta,\ |g|<g_0\},
\label{eq:v2v47-sector}
\end{equation}
define \(\mathfrak F_c\) to be the space of analytic maps
\(\Delta:S_\theta\to\mathfrak B_{\rm phys,ret}\) such that
\begin{equation}
 \|\Delta(g)\|_{\mathfrak B}
 \le C_\Delta e^{-c/|g|}
\label{eq:v2v47-flat}
\end{equation}
on every closed subsector, for some \(c>0\).

\begin{lemma}[Exponential flatness]
\label{lem:v2v47-flat}
Every \(\Delta\in\mathfrak F_c\) has the zero Gevrey asymptotic power
series:
\begin{equation}
 \|\Delta(g)\|_{\mathfrak B}=O(|g|^N)
 \qquad\hbox{for every }N\ge0.
\label{eq:v2v47-zero-series}
\end{equation}
\end{lemma}

\begin{proof}
For \(x>0\), the function \(x^{-N}e^{-c/x}\) tends to zero as
\(x\downarrow0\) and is bounded on \((0,g_0]\).  Multiply
\eqref{eq:v2v47-flat} by \(|g|^{-N}\).
\end{proof}

\begin{proposition}[Structural ledgers do not remove flat additions]
\label{prop:v2v47-ledgers}
Let \(\mathcal R(g)\in\mathfrak B_{\rm phys,ret}\) be a completed response
and let \(X\in\mathfrak B_{\rm phys,ret}\).  Then
\begin{equation}
 \mathcal R_a(g)=\mathcal R(g)+a e^{-c/g}X,
 \qquad a\in\mathbb C,quad g>0,
\label{eq:v2v47-family}
\end{equation}
has the same perturbative expansion, Ward identity, retarded support, and
finite Dini/Volterra norm as \(\mathcal R\).  Its norm differs by at most
\(|a|e^{-c/g}\|X\|_{\mathfrak B}\).
\end{proposition}

\begin{proof}
The perturbative statement is \cref{lem:v2v47-flat}.  The physical causal
space is linear and closed by \cref{lem:v2v46-closed}, so
\(\mathcal R_a\) remains in it.  The response-norm estimate is the triangle
inequality.
\end{proof}

\subsection{Positivity does not generally select one completion}
\label{sec:v2v47-positive}

\begin{proposition}[A positive spectral ray of ambiguities]
\label{prop:v2v47-positive-ray}
Let \(P\ge0\) be a nonzero physical polarization operator and choose
\(\Omega_0>0\).  Define the positive spectral response datum
\begin{equation}
 \mathrm d\Sigma_X(\Omega)=P\,\delta_{\Omega_0}(\mathrm d\Omega).
\label{eq:v2v47-atom}
\end{equation}
It has finite three-subtracted Dini norm and reconstructs a retarded kernel.
If \(\Sigma\ge0\), then
\begin{equation}
 \Sigma_a(g)=\Sigma(g)+a e^{-c/g}\Sigma_X
\label{eq:v2v47-positive-family}
\end{equation}
is positive for every \(a\ge0\), has the same perturbative series, and
satisfies the same Ward, causal, and Dini conditions.
\end{proposition}

\begin{proof}
The weighted total variation of the atom is
\(\Omega_0^{-7}\|P\|<\infty\).  Its sine reconstruction is supported at
positive relative time after retarded restriction.  The sum of positive
operator-valued measures with nonnegative coefficients is positive.  Apply
\cref{prop:v2v47-ledgers} for the remaining assertions.
\end{proof}

Thus positivity is essential for physical interpretation but is not, by
itself, a normalization condition for the nonperturbative spectral weight.

\subsection{Finite local conditions cannot fix an arbitrary nonlocal space}
\label{sec:v2v47-no-finite}

Let \(L_1,\ldots,L_m\) be continuous linear normalization functionals on
\(\mathfrak B_{\rm phys,ret}\).  They may encode prescribed local
couplings, finitely many Euclidean values, or finitely many spectral
moments.

\begin{theorem}[Finite-codimension nonuniqueness]
\label{thm:v2v47-nonunique}
If an allowed ambiguity subspace \(\mathfrak X\subset
\mathfrak B_{\rm phys,ret}\) is infinite dimensional, then
\begin{equation}
 \mathfrak X\cap\bigcap_{j=1}^m\ker L_j
\label{eq:v2v47-kernel-intersection}
\end{equation}
is nonzero and, in fact, has finite codimension at most \(m\) in
\(\mathfrak X\).  Hence finitely many linear normalization conditions do
not select a unique completion among all exponentially flat additions
\(e^{-c/g}X\), \(X\in\mathfrak X\).
\end{theorem}

\begin{proof}
The map
\[
 L:\mathfrak X\longrightarrow\mathbb C^m,
 \qquad L(X)=(L_1X,\ldots,L_mX),
\]
has rank at most \(m\).  The quotient
\(\mathfrak X/\ker L\) is isomorphic to its finite-dimensional range.
Therefore \(\ker L\) has codimension at most \(m\) and cannot be zero when
\(\mathfrak X\) is infinite dimensional.
\end{proof}

In particular, fixing the cosmological, Newton, and curvature-squared
contact coefficients does not determine an arbitrary nonlocal spectral
atom or continuum.  Local renormalization conditions and nonperturbative
state selection are different tasks.

\subsection{A finite-dimensional uniqueness theorem}
\label{sec:v2v47-finite}

Suppose independent nonperturbative analysis identifies a finite family
\(X_1,\ldots,X_m\in\mathfrak B_{\rm phys,ret}\) and positive actions
\(c_1,\ldots,c_m\).  Consider
\begin{equation}
 \Delta_{\boldsymbol\sigma}(g)
 =\sum_{j=1}^m\sigma_j e^{-c_j/g}X_j.
\label{eq:v2v47-transseries}
\end{equation}

\begin{theorem}[Normalization-matrix criterion]
\label{thm:v2v47-matrix}
Let \(L_1,\ldots,L_m\) be continuous linear functionals and define, at one
fixed physical \(g_*>0\),
\begin{equation}
 \mathsf M_{ij}=e^{-c_j/g_*}L_i(X_j).
\label{eq:v2v47-M}
\end{equation}
For prescribed normalization values \(d_i\), the equations
\begin{equation}
 L_i(\mathcal R(g_*)+\Delta_{\boldsymbol\sigma}(g_*))=d_i
 \qquad(1\le i\le m)
\label{eq:v2v47-normalization}
\end{equation}
select a unique coefficient vector \(\boldsymbol\sigma\) if and only if
\(\mathsf M\) is invertible.  In that case
\begin{equation}
 \boldsymbol\sigma
 =\mathsf M^{-1}
 \bigl(\boldsymbol d-L(\mathcal R(g_*))\bigr).
\label{eq:v2v47-sigma}
\end{equation}
\end{theorem}

\begin{proof}
Substitution of \eqref{eq:v2v47-transseries} into
\eqref{eq:v2v47-normalization} gives the finite linear system
\(\mathsf M\boldsymbol\sigma=\boldsymbol d-L(\mathcal R(g_*))\).
It has one solution for every right side exactly when \(\mathsf M\) is
invertible, and then the solution is \eqref{eq:v2v47-sigma}.
\end{proof}

\begin{firewall}[The finite transseries basis must be derived]
\label{fw:v2v47-basis}
The theorem does not justify truncating the ambiguity space to
\(\operatorname{span}\{X_1,\ldots,X_m\}\).  The actions, response vectors,
and admissible Stokes parameters must come from a constructive definition,
resurgent equation, saddle analysis, or other nonperturbative input.  An
invertible normalization matrix proves uniqueness only inside the supplied
finite-dimensional model.
\end{firewall}

\subsection{Reality and median lateral sums}
\label{sec:v2v47-reality}

\begin{proposition}[What median summation preserves]
\label{prop:v2v47-median}
Suppose real perturbative coefficients have lateral Borel sums
\(\mathcal R_+(g)\) and \(\mathcal R_-(g)\) satisfying
\begin{equation}
 \mathcal R_-(g)=\overline{\mathcal R_+(g)}
\label{eq:v2v47-conjugate}
\end{equation}
and both belong to \(\mathfrak B_{\rm phys,ret}\).  Then
\begin{equation}
 \mathcal R_{\rm med}(g)
 =\frac12(\mathcal R_+(g)+\mathcal R_-(g))
\label{eq:v2v47-median}
\end{equation}
is real, Ward compatible, retarded, and Dini bounded.  It is spectrally
positive if both lateral spectral measures are positive; conjugacy alone
does not imply positivity.
\end{proposition}

\begin{proof}
Equation \eqref{eq:v2v47-conjugate} makes the average real.  The physical
causal response space is linear and closed, and its norm obeys the convex
bound.  The cone of positive measures is convex, proving the conditional
positivity statement.  Complex conjugacy imposes reality, not a sign on
quadratic spectral forms.
\end{proof}

\subsection{Propagation of ambiguity through the causal solution}
\label{sec:v2v47-propagation}

\begin{theorem}[Linear Volterra sensitivity]
\label{thm:v2v47-sensitivity}
Let
\begin{equation}
 h_j=(I-K_j)^{-1}h_{0,j},
 \qquad j=1,2,
\label{eq:v2v47-solutions}
\end{equation}
and suppose the V40 Volterra resolvents satisfy
\(\|(I-K_j)^{-1}\|\le M_j\).  Then
\begin{equation}
 \|h_1-h_2\|_{X_T}
 \le M_1\left(
 \|h_{0,1}-h_{0,2}\|_{X_T}
 +\|K_1-K_2\|\,\|h_2\|_{X_T}
 \right).
\label{eq:v2v47-sensitivity}
\end{equation}
In particular, an \(O(e^{-c/g})\) response ambiguity produces an
\(O(e^{-c/g})\) solution ambiguity on a slab with uniform
\(M_j\) and bounded \(h_2\).
\end{theorem}

\begin{proof}
Subtract
\((I-K_1)h_1=h_{0,1}\) and
\((I-K_2)h_2=h_{0,2}\) to obtain
\[
 (I-K_1)(h_1-h_2)
 =h_{0,1}-h_{0,2}+(K_1-K_2)h_2.
\]
Apply \((I-K_1)^{-1}\) and take norms.
\end{proof}

\begin{corollary}[Nonlinear sensitivity inside a contraction ball]
\label{cor:v2v47-nonlinear}
Suppose the two nonlinear maps have a common contraction constant
\(q<1\) after their Volterra resolvents are applied.  If their data and
response maps differ by at most \(\varepsilon\) uniformly on the common
ball, their fixed points differ by at most
\begin{equation}
 \frac{C\varepsilon}{1-q}.
\label{eq:v2v47-nonlinear}
\end{equation}
Thus an exponentially flat completion ambiguity persists at the same
exponential order in the nonlinear solution while the contraction margin
is uniform.
\end{corollary}

\begin{proof}
Subtract the two fixed-point equations.  Bound the difference of the maps
at one fixed point by \(C\varepsilon\), bound the common Lipschitz part by
\(q\) times the fixed-point difference, and absorb it on the left.
\end{proof}

\begin{firewall}[Causal evolution does not select the Borel completion]
\label{fw:v2v47-evolution}
The Volterra theorem propagates a chosen response uniquely, but it does not
choose among responses that differ by exponentially flat physical causal
terms.  The sensitivity estimate shows that the ambiguity reaches the
metric and matter solution rather than disappearing under retarded
evolution.
\end{firewall}

\subsection{V47 conclusion and next step}
\label{sec:v2v47-conclusion}

The perturbative series, Ward identity, causal support, finite Dini norm,
and even spectral positivity can leave a nontrivial ray of exponentially
small response additions.  Finitely many local conditions cannot fix an
infinite-dimensional nonlocal ambiguity.  A finite normalization matrix
does give uniqueness after a finite transseries basis has been derived by
independent nonperturbative analysis.  Any unresolved ambiguity propagates
through the linear and nonlinear Volterra equations at the same
exponential order.

V48 will test a stronger selection mechanism: Euclidean reflection
positivity and reconstruction.  It will state the precise conditions under
which a regulator-independent Euclidean measure determines a unique
positive spectral response and hence removes the lateral ambiguity, while
separating this implication from the still-open construction of that
measure for the chiral Standard Model coupled to dynamical geometry.
\clearpage
\section*{Second Volume, Version V48: Reflection Positivity, Spectral Reconstruction, and Euclidean Selection}
\addcontentsline{toc}{section}{Second Volume, Version V48: Reflection Positivity, Spectral Reconstruction, and Euclidean Selection}

\subsection*{Purpose and scope}

V47 proved that perturbative data, Ward identities, causality, the Dini
bound, and even positivity can leave exponentially small response
ambiguities.  This note asks whether a complete Euclidean construction can
select one response.  On a reflection-invariant stationary background, the
answer is affirmative under the Osterwalder--Schrader hypotheses: the full
noncoincident Euclidean stress correlator reconstructs a positive
Hamiltonian and a unique positive spectral measure.  Its analytic
continuation then fixes the nonlocal retarded response.

The theorem is an implication, not a construction of the Euclidean
Standard Model coupled to gravity.  It is stated after spatial smearing on
the fixed V33 band, so the remaining Euclidean time separation is a genuine
operator-valued function for \(s>0\).  Coincidence-supported contact terms
are excluded from that function and remain fixed by the local
renormalization conditions of V40.

\subsection{Reflection-positive Euclidean stress kernel}
\label{sec:v2v48-OS}

Let \(V_\kappa\) be the finite physical polarization space obtained after
spatial low-band projection and Ward reduction on an ultrastatic Euclidean
cylinder.  Let
\begin{equation}
 C_E(s):V_\kappa\longrightarrow V_\kappa,
 \qquad s>0,
\label{eq:v2v48-CE}
\end{equation}
be the noncoincident connected Euclidean stress two-point kernel.  Assume
Euclidean time-translation invariance, Hermiticity, regularity for
\(s>0\), and reflection positivity:
\begin{equation}
 \sum_{i,j=1}^n
 \langle v_i,C_E(s_i+s_j)v_j\rangle\ge0
\label{eq:v2v48-RP}
\end{equation}
for every \(n\), \(s_i>0\), and \(v_i\in V_\kappa\).  Include the standard
OS time-translation axiom so that positive Euclidean translations descend
to a strongly continuous contraction semigroup.

\begin{theorem}[OS two-point reconstruction]
\label{thm:v2v48-reconstruction}
Under the preceding hypotheses there are a Hilbert space
\(\mathcal H_{\rm OS}\), a nonnegative self-adjoint operator \(H\), and,
for every \(\varepsilon>0\), a bounded map
\(J_\varepsilon:V_\kappa\to\mathcal H_{\rm OS}\) such that
\begin{equation}
 C_E(s+2\varepsilon)
 =J_\varepsilon^*e^{-sH}J_\varepsilon,
 \qquad s>0.
\label{eq:v2v48-semigroup}
\end{equation}
There is a unique positive locally finite operator-valued measure
\(\Sigma\) on \([0,\infty)\), with finite exponential moments, satisfying
\begin{equation}
 C_E(s)=\int_{[0,\infty)}e^{-s\Omega}
 \,\mathrm d\Sigma(\Omega),
 \qquad s>0.
\label{eq:v2v48-Laplace}
\end{equation}
For every \(\varepsilon>0\), its regularization obeys
\begin{equation}
 e^{-2\varepsilon\Omega}\,\mathrm d\Sigma(\Omega)
 =J_\varepsilon^*\mathbf1_{\mathrm d\Omega}(H)J_\varepsilon.
\label{eq:v2v48-Sigma}
\end{equation}
The unregularized stress insertion need not define a bounded vector at
Euclidean time zero.
\end{theorem}

\begin{proof}
On the vector space generated by symbols \([s,v]\), \(s>0\), define
\[
 \langle[s,v],[t,w]\rangle_{\rm OS}
 =\langle v,C_E(s+t)w\rangle.
\]
Reflection positivity makes this form nonnegative.  Quotient by its null
space and complete.  Positive time translation
\([s,v]\mapsto[s+a,v]\) descends, by the OS translation axiom, to a
strongly continuous symmetric contraction semigroup.  Hence it has the
form \(e^{-aH}\) for a nonnegative self-adjoint generator \(H\).  For
\(\varepsilon>0\), set \(J_\varepsilon v=[\varepsilon,v]\).  Its norm is
controlled by \(C_E(2\varepsilon)\), so it is bounded on the
finite-dimensional space \(V_\kappa\), and the OS inner product gives
\eqref{eq:v2v48-semigroup}.

The spectral theorem gives the finite regularized measure
\(\mathrm d\Sigma_\varepsilon
=J_\varepsilon^*\mathbf1_{\mathrm d\Omega}(H)J_\varepsilon\).
The semigroup relation implies
\(\mathrm d\Sigma_{\varepsilon_2}
=e^{-2(\varepsilon_2-\varepsilon_1)\Omega}
\mathrm d\Sigma_{\varepsilon_1}\) for \(\varepsilon_2>\varepsilon_1\).
Therefore
\(\mathrm d\Sigma=e^{2\varepsilon\Omega}\mathrm d\Sigma_\varepsilon\)
is independent of \(\varepsilon\) on bounded spectral intervals, is
locally finite, and has every positive exponential damping as a finite
measure.  This proves \eqref{eq:v2v48-Laplace,eq:v2v48-Sigma}.

For uniqueness, suppose two positive locally finite operator measures with
finite exponential moments have the same Laplace transform.  Pair their
difference with arbitrary vectors and multiply by
\(e^{-s_0\Omega}\), \(s_0>0\), to obtain a finite signed scalar measure.
Its integrals against every exponential \(e^{-s\Omega}\), \(s\ge0\),
vanish.  After the change of variable \(x=e^{-\Omega}\), polynomials in
\(x\) are dense in the continuous functions on compact subintervals, and
tight truncation handles infinity.  The signed measure is zero.
Polarization in the external vectors proves equality of the operator
measures.
\end{proof}

\subsection{Unique retarded continuation}
\subsection{Unique retarded continuation}
\label{sec:v2v48-retarded}

For spatially smeared physical stresses, the reconstructed Wightman
function is
\begin{equation}
 W(t)=\int_{[0,\infty)}e^{-it\Omega}
 \,\mathrm d\Sigma(\Omega)
\label{eq:v2v48-Wightman}
\end{equation}
in the distributional sense.  With the convention
\(R(t)=i\mathbf1_{t>0}[T(t),T(0)]\), its nonlocal part is
\begin{equation}
 R(t)=2\mathbf1_{t>0}
 \int_{[0,\infty)}\sin(\Omega t)
 \,\mathrm d\Sigma(\Omega),
\label{eq:v2v48-retarded}
\end{equation}
up to the fixed tensor-index convention.

\begin{corollary}[Euclidean uniqueness removes a lateral ambiguity]
\label{cor:v2v48-unique-retarded}
Let two completed responses have the same noncoincident Euclidean kernel
\(C_E(s)\) for every \(s>0\), satisfy the OS hypotheses, and have identical
local contact coefficients.  Then their positive spectral measures and
their complete retarded responses coincide.
\end{corollary}

\begin{proof}
The uniqueness part of \cref{thm:v2v48-reconstruction} makes the two
spectral measures equal.  Equation \eqref{eq:v2v48-retarded} then makes the
nonlocal retarded parts equal, while equality of the contact coefficients
handles the diagonal-supported part.
\end{proof}

\begin{proposition}[An asymptotic Euclidean series is insufficient]
\label{prop:v2v48-asymptotic}
Equality of the perturbative or short-distance asymptotic expansions of
two Euclidean kernels does not imply the conclusion of
\cref{cor:v2v48-unique-retarded}.  Indeed, for the positive atom
\(\Sigma_X=P\delta_{\Omega_0}\) of V47,
\begin{equation}
 C_{E,a}(s;g)
 =C_E(s;g)+a e^{-c/g}e^{-s\Omega_0}P,
 \qquad a\ge0,
\label{eq:v2v48-flat-Euclidean}
\end{equation}
is reflection positive whenever \(C_E\) is, and has the same power series
in \(g\) as \(C_E\).
\end{proposition}

\begin{proof}
The added kernel is the Laplace transform of the positive measure
\(a e^{-c/g}P\delta_{\Omega_0}\), so reflection positivity is preserved.
Exponential flatness is \cref{lem:v2v47-flat}.
\end{proof}

Thus Euclidean selection works only when a complete Euclidean correlator or
measure is constructed, rather than when its formal expansion is specified.

\subsection{An exact Euclidean form of the Dini moment}
\label{sec:v2v48-Dini}

Fix \(\mu>0\), let
\(\Sigma^{\rm hi}=\mathbf1_{[\mu,\infty)}\Sigma\), and define
\begin{equation}
 C_E^{\rm hi}(s)
 =\int_{[\mu,\infty)}e^{-s\Omega}
 \,\mathrm d\Sigma(\Omega).
\label{eq:v2v48-high-CE}
\end{equation}

\begin{theorem}[Euclidean moment identity]
\label{thm:v2v48-Dini}
In the weak operator sense,
\begin{equation}
 \int_{[\mu,\infty)}\Omega^{-7}
 \,\mathrm d\Sigma(\Omega)
 =\frac1{6!}\int_0^\infty s^6C_E^{\rm hi}(s)\,\mathrm ds.
\label{eq:v2v48-Dini}
\end{equation}
Both sides are positive.  In particular, if the external polarization
space is finite dimensional and
\begin{align}
 \operatorname{tr}C_E^{\rm hi}(s)
 &\le C s^{-6}(1+|\log s|)^L,
 &&0<s\le1,
\label{eq:v2v48-short}\\
 \operatorname{tr}C_E^{\rm hi}(s)
 &\le C e^{-\mu(s-1)},
 &&s\ge1,
\label{eq:v2v48-long}
\end{align}
then the V43 three-subtracted Dini moment is finite.
\end{theorem}

\begin{proof}
For \(\Omega>0\),
\begin{equation}
 \Omega^{-7}=\frac1{6!}\int_0^\infty s^6e^{-s\Omega}\,\mathrm ds.
\label{eq:v2v48-Gamma}
\end{equation}
Pair with a vector and use Tonelli's theorem for the positive scalar
measure to interchange the integrals.  Polarization gives
\eqref{eq:v2v48-Dini}.  Under \eqref{eq:v2v48-short}, multiplication by
\(s^6\) leaves an integrable logarithmic power at zero.  The high spectral
support gives exponential decay; the stated bound makes the large-time
integral finite.  In finite dimension, the trace controls the positive
operator measure and hence its total-variation Dini norm.
\end{proof}

\begin{remark}[The free scaling reappears in Euclidean time]
\label{rem:v2v48-scaling}
The V42 cumulative growth \(F(R)=O(R^6\log^L R)\) corresponds to the
short-time estimate \(C_E^{\rm hi}(s)=O(s^{-6}\log^L(1/s))\).  The factor
\(s^6\) in \eqref{eq:v2v48-Dini} is exactly the gain from three
subtractions.  No additional ultraviolet power is lost in the Euclidean
reconstruction.
\end{remark}

\subsection{Regulator-stable OS reconstruction}
\label{sec:v2v48-regulator}

Let \(C_{E,\Lambda}\) be reflection-positive regulated kernels with
positive reconstructed measures \(\Sigma_\Lambda\).  Put
\begin{equation}
 \nu_\Lambda(\mathrm d\Omega)
 =\mathbf1_{[\mu,\infty)}(\Omega)
 \Omega^{-7}\Sigma_\Lambda(\mathrm d\Omega).
\label{eq:v2v48-nu}
\end{equation}

\begin{lemma}[Laplace transform of the weighted measure]
\label{lem:v2v48-weighted-Laplace}
For \(s>0\),
\begin{equation}
 \int e^{-s\Omega}\,\mathrm d\nu_\Lambda(\Omega)
 =\frac1{6!}\int_s^\infty
 (u-s)^6C_{E,\Lambda}^{\rm hi}(u)\,\mathrm du.
\label{eq:v2v48-weighted-Laplace}
\end{equation}
\end{lemma}

\begin{proof}
For fixed \(\Omega\), substitute \(v=u-s\) and use
\eqref{eq:v2v48-Gamma}:
\[
 \frac1{6!}\int_s^\infty(u-s)^6e^{-u\Omega}\,\mathrm du
 =e^{-s\Omega}\Omega^{-7}.
\]
Integrate against the positive measure and apply Tonelli.
\end{proof}

\begin{theorem}[OS limit and convergence of subtracted memory]
\label{thm:v2v48-regulator}
Assume:
\begin{enumerate}
 \item \(C_{E,\Lambda}(s)\to C_E(s)\) and
       \(C_{E,\Lambda}^{\rm hi}(s)\to C_E^{\rm hi}(s)\) for every
       \(s>0\), in the physical finite-band operator topology;
 \item the bounds \eqref{eq:v2v48-short,eq:v2v48-long} hold uniformly;
 \item the weighted measures are uniformly tight:
 \begin{equation}
  \lim_{R\to\infty}\sup_\Lambda
  \|\nu_\Lambda([R,\infty))\|=0;
 \label{eq:v2v48-tight}
 \end{equation}
 \item the local contact coefficients converge.
\end{enumerate}
Then \(C_E\) is reflection positive, its OS measure \(\Sigma\) is unique,
and \(\nu_\Lambda\) converges weakly to
\(\nu=\mathbf1_{[\mu,\infty)}\Omega^{-7}\Sigma\).  The three-subtracted
retarded kernels converge uniformly for \(t\) in compact intervals.
Together with convergence of the effective Einstein propagator, the V40
linear causal solutions converge.
\end{theorem}

\begin{proof}
Every finite reflection-positive quadratic form
\eqref{eq:v2v48-RP} remains nonnegative under the stated pointwise limit,
so \(C_E\) is reflection positive.  The reconstruction theorem gives its
unique measure.

Uniform Euclidean bounds dominate the right side of
\eqref{eq:v2v48-weighted-Laplace}; dominated convergence therefore gives
convergence of the Laplace transforms of \(\nu_\Lambda\).  Uniform
tightness and uniqueness of the Laplace transform identify every weak
subsequence limit with \(\nu\), hence the full family converges weakly.

On a fixed compact frequency interval, weak convergence of finite
operator measures makes their sine transforms converge uniformly for
\(t\) in a compact interval, by a finite-net argument.  The uniformly
small tails in \eqref{eq:v2v48-tight} complete the estimate.  These sine
transforms are exactly the three-subtracted retarded kernels.  Apply
\cref{thm:v2v40-regulator} after composing with the convergent Einstein
propagators.
\end{proof}

\begin{firewall}[Uniform bounded Dini mass is weaker than tightness]
\label{fw:v2v48-tight}
A uniform bound on \(\nu_\Lambda([\mu,\infty))\) does not prevent that mass
from escaping to infinite frequency.  The tail condition
\eqref{eq:v2v48-tight} is needed for uniform convergence of the retarded
kernels; it is the Euclidean counterpart of the regulator-tightness
condition in V43.
\end{firewall}

\subsection{Limits of Euclidean selection for the full programme}
\label{sec:v2v48-limits}

\begin{firewall}[Stationarity and reflection are genuine restrictions]
\label{fw:v2v48-stationary}
The semigroup and one-variable spectral measure use Euclidean time
translation and a reflection hypersurface.  A general evolving Einstein
background has neither.  Local Hadamard continuation from V44 can control
ultraviolet singularities there, but it does not supply one global OS
Hamiltonian or select a nonperturbative state throughout the evolution.
\end{firewall}

\begin{firewall}[The required Euclidean SM--gravity measure is unconstructed]
\label{fw:v2v48-construction}
For the chiral Standard Model, a regulator must implement anomaly
cancellation, the phase of the fermion determinant, gauge reduction, and
reflection positivity on the physical observable algebra.  Coupling to a
dynamical Euclidean metric adds diffeomorphism gauge fixing and the
unbounded conformal direction of the Einstein action.  Formal Wick
rotation or perturbative anomaly cancellation does not prove existence,
normalization, tightness, or reflection positivity of the required
regulator-independent measure.
\end{firewall}

\begin{theorem}[Exact Euclidean selection criterion]
\label{thm:v2v48-selection}
A Euclidean route removes the V47 nonperturbative response ambiguity if all
of the following are supplied:
\begin{enumerate}
 \item a regulator-independent reflection-positive physical Euclidean
       stress kernel for every \(s>0\);
 \item convergent local contact coefficients fixed by renormalization
       conditions;
 \item the short-time and tightness estimates of
       \cref{thm:v2v48-Dini,thm:v2v48-regulator};
 \item an analytic-continuation or reconstruction theorem identifying the
       resulting Lorentzian observable with the one entering the coupled
       Einstein equation.
\end{enumerate}
Under these hypotheses the positive spectral measure, nonlocal retarded
response, and V40 linear causal solution are unique on the stationary
reference slab.
\end{theorem}

\begin{proof}
OS reconstruction and Laplace uniqueness give the unique positive measure.
The Dini identity and tightness give the regulated kernel limit.  Contact
conditions fix the local ambiguity.  The fourth hypothesis identifies this
unique response with the Lorentzian source, and V40 gives the unique causal
solution.
\end{proof}

\subsection{V48 conclusion and next step}
\label{sec:v2v48-conclusion}

The complete noncoincident reflection-positive Euclidean stress correlator
determines a unique positive spectral measure and retarded response.  Its
three-subtracted Dini moment is exactly the sixth Euclidean-time moment
\((6!)^{-1}\int s^6C_E^{\rm hi}(s)\,\mathrm ds\).  Uniform OS convergence,
short-time control, and weighted spectral tightness pass that response
through the regulator limit.  Equality of perturbative expansions is not
enough, because a positive exponentially flat spectral atom preserves all
such expansions.

The unresolved step is constructive: no reflection-positive,
regulator-independent Euclidean measure for the chiral SM coupled to a
dynamical Einstein metric has been produced.  V49 will go upstream to the
measure-existence question, separating the matter theory on a fixed
reflection-positive background from dynamical gravity and identifying the
minimal tightness, gauge-cohomology, and conformal-mode estimates needed
before OS reconstruction can be invoked.
\clearpage
\section*{Second Volume, Version V49: Fixed-Background OS Compactness and the Euclidean Conformal Obstruction}
\addcontentsline{toc}{section}{Second Volume, Version V49: Fixed-Background OS Compactness and the Euclidean Conformal Obstruction}

\subsection*{Purpose}

V48 showed that a complete reflection-positive Euclidean stress correlator
would select a unique positive retarded response.  This note separates two
measure-existence problems that must not be conflated.  On a fixed compact
reflection-symmetric background, tight positive regulator measures have
reflection-positive subsequential limits.  Composite stress correlators
pass to that limit only with an additional uniform-integrability estimate.
For a dynamical Euclidean metric, the Einstein action has an explicit
negative conformal kinetic direction and does not by itself define a
normalizable positive Gibbs weight.  Thus fixed-background matter
compactness does not solve the gravity measure problem.

\subsection{Tightness of fixed-background matter regulators}
\label{sec:v2v49-tightness}

Let \(X=H^{-s}(M;E)\) be a separable Hilbert space of distributional field
configurations on a compact Euclidean manifold, and let
\(X_0=H^{-s_0}(M;E)\) with \(s>s_0\), so the embedding
\begin{equation}
 X_0\hookrightarrow X
\label{eq:v2v49-compact-embedding}
\end{equation}
is compact.  Let \(\mu_\Lambda\) be probability measures describing a
positive regulator on the physical observable configuration space.

\begin{theorem}[Sobolev-moment tightness]
\label{thm:v2v49-tightness}
If, for some \(p>0\),
\begin{equation}
 \sup_\Lambda\int_X
 \|\Phi\|_{X_0}^p\,\mathrm d\mu_\Lambda(\Phi)<\infty,
\label{eq:v2v49-moment}
\end{equation}
where the norm is set to infinity off \(X_0\), then
\((\mu_\Lambda)\) is tight on \(X\).  Hence every regulator sequence has a
weakly convergent subsequence.
\end{theorem}

\begin{proof}
Let the uniform moment in \eqref{eq:v2v49-moment} be \(C\).  Markov's
inequality gives
\[
 \mu_\Lambda(\|\Phi\|_{X_0}>R)\le CR^{-p}.
\]
The closure in \(X\) of the radius-\(R\) ball of \(X_0\) is compact by
\eqref{eq:v2v49-compact-embedding}.  Given \(\varepsilon>0\), choose
\(R\) with \(CR^{-p}<\varepsilon\).  This one compact set has measure at
least \(1-\varepsilon\) for every \(\Lambda\), proving tightness.
Prokhorov's theorem on the Polish space \(X\) gives subsequential weak
compactness.
\end{proof}

\subsection{Passage of reflection positivity}
\label{sec:v2v49-RP}

Let \(\theta\) be the Euclidean time reflection and let
\(\mathcal A_+\) be the algebra of bounded continuous cylinder observables
supported in positive Euclidean time.

\begin{theorem}[Reflection positivity is closed under weak convergence]
\label{thm:v2v49-RP}
Suppose every \(\mu_\Lambda\) satisfies
\begin{equation}
 \int_X\overline{F(\theta\Phi)}F(\Phi)
 \,\mathrm d\mu_\Lambda(\Phi)\ge0
 \qquad(F\in\mathcal A_+),
\label{eq:v2v49-RP}
\end{equation}
and \(\mu_\Lambda\Rightarrow\mu\) weakly on \(X\).  If reflection acts
continuously on \(X\), then \(\mu\) satisfies the same inequality.
Euclidean symmetries and linear Ward identities tested by bounded
continuous cylinder observables pass to the limit in the same way.
\end{theorem}

\begin{proof}
For \(F\in\mathcal A_+\), the function
\(\overline{F\circ\theta},F\) is bounded and continuous.  Its integrals
converge under weak convergence.  Taking the limit in
\eqref{eq:v2v49-RP} preserves nonnegativity.  The same argument applies to
equalities formed from bounded continuous test observables.
\end{proof}

For gauge theories, this theorem is applied to the gauge-invariant physical
observable algebra.  A gauge-fixed ghost functional is not a positive
probability measure, and reflection positivity of its separate factors is
neither required nor generally meaningful.

\subsection{Composite stress insertions require uniform integrability}
\label{sec:v2v49-stress}

Weak convergence of measures controls bounded continuous observables, but
the renormalized stress tensor is unbounded and depends on the regulator.

\begin{theorem}[Passage of stress two-point functions]
\label{thm:v2v49-stress}
Let \(T_\Lambda(f)\) be regulated, spatially and temporally smeared physical
stress observables, and suppose for each finite list of test functions:
\begin{enumerate}
 \item \((\Phi,T_\Lambda(f_1),\ldots,T_\Lambda(f_m))\) converges in law to
       \((\Phi,T(f_1),\ldots,T(f_m))\);
 \item for some \(\delta>0\),
 \begin{equation}
  \sup_\Lambda
  \int|T_\Lambda(f)|^{2+\delta}\,\mathrm d\mu_\Lambda<\infty.
 \label{eq:v2v49-stress-moment}
 \end{equation}
\end{enumerate}
Then
\begin{equation}
 \int T_\Lambda(f)T_\Lambda(h)\,\mathrm d\mu_\Lambda
 \longrightarrow
 \int T(f)T(h)\,\mathrm d\mu
\label{eq:v2v49-stress-limit}
\end{equation}
for all such \(f,h\).  Reflection positivity of the stress kernel passes to
the limit.
\end{theorem}

\begin{proof}
Holder's inequality and \eqref{eq:v2v49-stress-moment} give a uniform
\(L^{1+\delta/2}\) bound for
\(T_\Lambda(f)T_\Lambda(h)\).  The products are therefore uniformly
integrable.  Joint convergence in law and truncation at bounded product
size give convergence of expectations, proving
\eqref{eq:v2v49-stress-limit}.  The limiting reflection-positive quadratic
forms follow by applying this convergence to finite linear combinations of
positive-time smearings.
\end{proof}

\begin{proposition}[Weak field convergence alone is insufficient]
\label{prop:v2v49-weak-insufficient}
There are random variables \(X_n\to0\) in probability whose second moments
do not tend to zero.
\end{proposition}

\begin{proof}
Let \(X_n=n^{1/2}\) with probability \(n^{-1}\) and zero otherwise.  Then
\(X_n\to0\) in probability but \(\mathbb E X_n^2=1\).  Thus a concentrated
composite insertion can survive weak convergence unless uniform
integrability is proved.
\end{proof}

The stress moment in \eqref{eq:v2v49-stress-moment} is separate from the
field-space tightness moment \eqref{eq:v2v49-moment}.  Renormalization can
create composite concentration even when the underlying field measures are
tight.

\subsection{Full-sequence convergence and uniqueness}
\label{sec:v2v49-uniqueness}

\begin{theorem}[Subsequence compactness plus uniqueness]
\label{thm:v2v49-uniqueness}
Assume \((\mu_\Lambda)\) is tight and every subsequential limit has the same
Schwinger functional on a convergence-determining cylinder algebra.  Then
the full regulator family converges weakly to that unique measure.  If the
stress hypotheses of \cref{thm:v2v49-stress} are uniform, the full stress
kernel converges as well.
\end{theorem}

\begin{proof}
If the full sequence did not converge, some subsequence would remain
outside a weak neighbourhood of the proposed limit.  Tightness gives a
further convergent subsequence.  Its limit has the same determining
Schwinger functional and therefore is the proposed measure, a
contradiction.  Uniform integrability makes the stress argument independent
of the selected subsequence.
\end{proof}

\begin{firewall}[Compactness does not prove uniqueness]
\label{fw:v2v49-uniqueness}
The moment estimate \eqref{eq:v2v49-moment} produces limit points, not a
unique phase.  Uniqueness requires an independent argument, such as a
contraction, correlation inequality, phase specification, or complete set
of Schwinger equations with a uniqueness theorem.  Choosing one convergent
subsequence would reintroduce the nonperturbative ambiguity isolated in
V47.
\end{firewall}

\subsection{The Euclidean Einstein conformal direction}
\label{sec:v2v49-conformal}

Let \((M,\bar g)\) be a closed four-dimensional Riemannian manifold and put
\begin{equation}
 g=e^{2\varphi}\bar g.
\label{eq:v2v49-conformal}
\end{equation}
Use the Euclidean Einstein--Hilbert convention
\begin{equation}
 S_E(g)=-\frac1{16\pi G}
 \int_M(R_g-2\Lambda)\,\mathrm d\mu_g.
\label{eq:v2v49-SE}
\end{equation}

\begin{lemma}[Exact conformal formula]
\label{lem:v2v49-conformal}
On a closed four-manifold,
\begin{equation}
 S_E(e^{2\varphi}\bar g)
 =-\frac1{16\pi G}
 \int_M e^{2\varphi}
 \left(\bar R+6|\bar\nabla\varphi|^2\right)
 \,\mathrm d\mu_{\bar g}
 +\frac{\Lambda}{8\pi G}
 \int_M e^{4\varphi}\,\mathrm d\mu_{\bar g}.
\label{eq:v2v49-conformal-action}
\end{equation}
\end{lemma}

\begin{proof}
The conformal identities in dimension four are
\[
 R_{e^{2\varphi}\bar g}
 =e^{-2\varphi}
 \left(\bar R-6\bar\Delta\varphi
 -6|\bar\nabla\varphi|^2\right),
 \qquad
 \mathrm d\mu_g=e^{4\varphi}\mathrm d\mu_{\bar g}.
\]
Integration by parts gives
\(\int e^{2\varphi}\bar\Delta\varphi
=-2\int e^{2\varphi}|\bar\nabla\varphi|^2\).  Substitute into
\eqref{eq:v2v49-SE}.
\end{proof}

\begin{theorem}[Unbounded conformal action]
\label{thm:v2v49-unbounded}
The functional \(S_E\) is unbounded below along smooth conformal metrics,
even when \(\varphi\) remains uniformly bounded in \(L^\infty\).
\end{theorem}

\begin{proof}
Choose a coordinate ball and a nonzero compactly supported smooth
oscillatory sequence
\begin{equation}
 \varphi_n(x)=\varepsilon\chi(x)\sin(nx^1),
 \qquad 0<\varepsilon<1,
\label{eq:v2v49-sequence}
\end{equation}
with fixed cutoff \(\chi\).  Then \(e^{2\varphi_n}\) and
\(e^{4\varphi_n}\) are bounded above and below independently of \(n\),
while
\(\int|\bar\nabla\varphi_n|^2\ge c\varepsilon^2n^2-O(n)\).
The curvature-without-gradient and cosmological integrals in
\eqref{eq:v2v49-conformal-action} remain bounded.  Its negative gradient
term therefore drives \(S_E(e^{2\varphi_n}\bar g)\to-\infty\).
\end{proof}

\begin{corollary}[No coercivity from the bare Einstein weight]
\label{cor:v2v49-no-Gibbs}
Along the real conformal contour, the factor \(e^{-S_E(g)}\) fails to
decay along all bounded-amplitude smooth conformal directions, and the
bare action supplies no coercive sublevel-set estimate there.  Therefore
the tightness and normalization hypotheses required for a probability
measure cannot be deduced from the bare Einstein action alone.  Any
measure construction must supply additional input---for example a
contour, stabilization, constraint, or independent nonperturbative
definition---and separately prove reflection positivity and its relation
to Lorentzian gravity.
\end{corollary}

\begin{proof}
Along \eqref{eq:v2v49-sequence}, the formal Boltzmann weight grows like
\(e^{c n^2}\) while the amplitudes remain in a bounded conformal range.
Hence the bare action has no lower coercive bound on this family and its
exponential cannot furnish a decaying domination estimate along it.
\end{proof}

\begin{firewall}[Matter determinants do not automatically stabilize gravity]
\label{fw:v2v49-matter-stabilization}
Integrating out matter adds a nonlocal effective action and local
curvature counterterms.  It may change the high-frequency conformal
quadratic form, but stabilization requires a proved positive lower bound
for the complete Ward-renormalized action with the chosen contour.  Power
counting or the sign of one anomaly coefficient does not establish such a
bound, and higher-derivative stabilization must also be reconciled with the
causal order-reduction ledger of V40--V44.
\end{firewall}

\subsection{Exact separation of the two Euclidean routes}
\label{sec:v2v49-separation}

\begin{theorem}[Conditional fixed-background matter route]
\label{thm:v2v49-matter-route}
On a fixed reflection-positive background, suppose:
\begin{enumerate}
 \item a positive physical regulator measure satisfies the tightness bound
       \eqref{eq:v2v49-moment};
 \item reflection positivity and physical Ward identities hold at every
       regulator;
 \item the composite stress bounds and convergence in
       \cref{thm:v2v49-stress} hold;
 \item the short-time Dini and weighted-tail estimates of V48 hold;
 \item the limiting Schwinger functional is unique.
\end{enumerate}
Then the full regulator family determines a unique reflection-positive
stress kernel, a unique positive spectral response, and the V48--V40
linear causal matter backreaction on that fixed background.
\end{theorem}

\begin{proof}
Tightness and uniqueness give full measure convergence by
\cref{thm:v2v49-tightness,thm:v2v49-uniqueness}.  Stress uniform
integrability gives the limiting stress kernel.  Reflection positivity
passes by \cref{thm:v2v49-RP}, and V48 reconstructs the spectral measure
and retarded response.  Its Dini and tightness bounds feed the V40 Volterra
solution.
\end{proof}

\begin{firewall}[This theorem freezes the metric]
\label{fw:v2v49-fixed}
The measure theorem treats matter on a prescribed Euclidean geometry.  It
does not integrate over metrics, prove the semiclassical Einstein fixed
point, or supply uniform estimates as the background itself changes.  The
conformal theorem shows why simply adjoining the bare Einstein Gibbs factor
does not close that gap.
\end{firewall}

\subsection{V49 conclusion and next step}
\label{sec:v2v49-conclusion}

Fixed-background Euclidean matter has a clear compactness route: a stronger
Sobolev moment gives tightness, reflection positivity is closed under weak
limits, composite stress convergence requires its own uniform-integrability
bound, and uniqueness upgrades subsequences to the full regulator limit.
Dynamical Euclidean gravity fails at the first coercivity step because its
real conformal kinetic direction drives the Einstein action to minus
infinity.

V50 is the next mandatory companion audit.  After that audit the programme
will compare two upstream alternatives: a proved stabilized or complex
Euclidean gravity contour with OS reconstruction, and a purely Lorentzian
state-selection compactness theorem that avoids a positive Euclidean
metric measure.  Neither route will be treated as available without its
own tightness, gauge, and causal estimates.
\clearpage
\section*{Second Volume, Version V50: Companion Audit and the Local Stabilizer--OS Obstruction}
\addcontentsline{toc}{section}{Second Volume, Version V50: Companion Audit and the Local Stabilizer--OS Obstruction}

\subsection*{Purpose}

V49 isolated the negative conformal kinetic direction of the Euclidean
Einstein action.  The most immediate proposed repair is to add a positive
local fourth-order term.  This note tests that repair before importing it
into the continuum programme.  The conclusion is a precise Gaussian
obstruction: a nonzero scalar covariance with a positive
Kallen--Lehmann measure cannot decay faster than one inverse power of
Euclidean momentum squared, whereas a strictly fourth-order local
stabilizer produces two inverse powers.  Thus positivity of the Euclidean
Gaussian quadratic form and Osterwalder--Schrader positivity are different
requirements.

This is also the mandatory fifth-version companion audit.  The audit changes
the acquisition order.  Stronger ultraviolet decay is not accepted as a
substitute for a positive spectral measure, and all gauge and partition
sectors must be assembled before positivity or norm estimates are taken.

\subsection{V50 companion audit}
\label{sec:v2v50-audit}

The Navier--Stokes file inspected was
\[
 \texttt{Renewal\_NS\_Stability\_Research\_Notes\_V31.tex},
\]
of length \(887537\) bytes and SHA--256 digest
\[
 \texttt{F1121B62238AD5491BB7CF03635EA9934942EDF573ED8FAAA9EE79E1D38A6696}.
\]
Its V31 formation theorem proves that heat evolution converts the
scale-matched source \(A^{-1/4}\mathcal B_{\rm NS}(u)\) into the critical
dyadic first moment.  Replacing it by the weaker, energy-paid
\(A^{-3/4}\mathcal B_{\rm NS}(u)\) merely transfers the missing two powers
of frequency into the shell weights.  The available global upper bound then
returns to the continuation-class quantity
\(\int\|\Lambda u\|_2^4\,\mathrm dt\).  The relevant lesson here is logical:
a weak tightness or smoothing estimate does not create the weighted moment
needed at the endpoint; that moment needs its own uniformly integrable
stock.

The Yang--Mills files named V42 and V43 were byte-identical at the audit
time.  Each had length \(558513\) bytes and SHA--256 digest
\[
 \texttt{A376E6C33D593AEB98389113527109492A7B9AD0EAB360B77D5A3483B26671D6}.
\]
Their last mathematical chapter is V42.  It expresses the connected fourth
cumulant through the single fifteen-element partition lattice \(\Pi_4\),
keeps unmarked coordinate factors inside the exact global state, and takes
one fixed state norm only after assembly.  This removes a false
support-dependent tensor expansion.  The relevant lesson here is again
logical: positivity of separate gauge, ghost, or partition pieces is not
the criterion.  Positivity must be tested on the complete physical
observable carrier after all signed pieces have been combined.

\begin{assessment}[Acquisition rule imported by the audit]
\label{ass:v2v50-audit-rule}
For the SM--Einstein programme the two lessons give the following rule.
A proposed regulator must provide, on the complete physical observable
algebra,
\begin{enumerate}
 \item one globally assembled reflection-positive quadratic form;
 \item the exact weighted stress moments used by the V48 Dini identity;
 \item uniform integrability for every composite insertion whose limit is
       taken.
\end{enumerate}
Improved decay of a signed or gauge-dependent component proves none of
these items by itself.
\end{assessment}

\subsection{The conformal Hessian on a flat reference metric}
\label{sec:v2v50-hessian}

Let \((M,\bar g)\) be a flat compact four-manifold without boundary and put
\(A=-\Delta_{\bar g}\ge0\).  Use the convention from V49,
\[
 S_E(g)
 =
 -\frac1{16\pi G}\int_M(R_g-2\Lambda)\,\mathrm d\mu_g .
\]
For a smooth real mean-zero function \(\psi\), set
\(g_\epsilon=e^{2\epsilon\psi}\bar g\).

\begin{lemma}[Exact quadratic conformal Hessian]
\label{lem:v2v50-conformal-hessian}
As \(\epsilon\to0\),
\begin{align}
 S_E(g_\epsilon)
 &=
 S_E(\bar g)
 +\epsilon^2\left[
  -\frac{3}{8\pi G}\langle\psi,A\psi\rangle
  +\frac{\Lambda}{\pi G}\|\psi\|_2^2
 \right]
 +O_\psi(\epsilon^3).
\label{eq:v2v50-conformal-expansion}
\end{align}
If one adds the quadratic local stabilizer
\[
 S_{\rm stab}(\varphi)=\frac a2\|A\varphi\|_2^2,
 \qquad a>0,
\]
then the combined Hessian on mean-zero conformal modes is
\begin{equation}
 P_a(A)
 =
 aA^2-\frac{3}{4\pi G}A+\frac{2\Lambda}{\pi G}.
\label{eq:v2v50-hessian-polynomial}
\end{equation}
\end{lemma}

\begin{proof}
Equation~\eqref{eq:v2v49-conformal-action} with
\(\bar R=0\) gives
\[
 S_E(g_\epsilon)
 =
 -\frac{3}{8\pi G}
 \int_M e^{2\epsilon\psi}\epsilon^2|\nabla\psi|^2
 \,\mathrm d\mu_{\bar g}
 +\frac{\Lambda}{8\pi G}
 \int_M e^{4\epsilon\psi}\,\mathrm d\mu_{\bar g}.
\]
Taylor expansion and \(\int_M\psi\,\mathrm d\mu_{\bar g}=0\) yield
\[
 \int_M e^{4\epsilon\psi}\,\mathrm d\mu_{\bar g}
 =
 \operatorname{Vol}_{\bar g}(M)
 +8\epsilon^2\|\psi\|_2^2+O_\psi(\epsilon^3).
\]
The gradient term equals
\(\epsilon^2\|\nabla\psi\|_2^2+O_\psi(\epsilon^3)\).
Since \(\|\nabla\psi\|_2^2=\langle\psi,A\psi\rangle\), this proves
\eqref{eq:v2v50-conformal-expansion}.  Writing the quadratic part as
\(\frac12\langle\varphi,P_a(A)\varphi\rangle\) gives
\eqref{eq:v2v50-hessian-polynomial}.
\end{proof}

\begin{corollary}[What the local term repairs]
\label{cor:v2v50-elliptic-repair}
The polynomial \(P_a(z)\) tends to \(+\infty\) as \(z\to+\infty\).
Consequently the added term reverses the high-frequency sign of the
conformal Hessian.  After any remaining finitely many nonpositive modes are
removed or shifted so that \(P_a(A)>0\), it defines an ordinary positive
Gaussian quadratic form with covariance \(P_a(A)^{-1}\).
\end{corollary}

\begin{proof}
The leading coefficient is \(a>0\).  On a compact manifold only finitely
many eigenvalues of \(A\) lie below any fixed threshold.  Strict positivity
on the retained spectrum is exactly the condition needed for the Gaussian
quadratic form.
\end{proof}

\subsection{Positive spectral measures cannot produce fourth-order decay}
\label{sec:v2v50-spectral-obstruction}

\begin{lemma}[Positive Stieltjes tails have an inverse-first-order lower bound]
\label{lem:v2v50-Stieltjes-lower}
Let \(\rho\) be a nonzero positive locally finite measure on
\([0,\infty)\), and suppose
\begin{equation}
 F(z)=\int_{[0,\infty)}\frac{\mathrm d\rho(m^2)}{z+m^2}<\infty
 \qquad(z>0).
\label{eq:v2v50-Stieltjes}
\end{equation}
Then there are \(R<\infty\) and \(c>0\) such that
\begin{equation}
 F(z)\ge\frac{c}{z+R}
 \qquad(z>0).
\label{eq:v2v50-Stieltjes-lower}
\end{equation}
In particular \(F(z)\) cannot be \(O(z^{-1-\delta})\) for any
\(\delta>0\).
\end{lemma}

\begin{proof}
Because \(\rho\ne0\) and is locally finite, some finite interval
\([0,R]\) has mass \(c:=\rho([0,R])>0\).  Positivity of the integrand gives
\[
 F(z)
 \ge
 \int_{[0,R]}\frac{\mathrm d\rho(m^2)}{z+m^2}
 \ge\frac{c}{z+R}.
\]
Multiplying by \(z^{1+\delta}\) proves the last assertion.
\end{proof}

\begin{definition}[Spectral-positive scalar covariance]
\label{def:v2v50-spectral-positive}
A translation- and Euclidean-invariant scalar covariance is called
spectral-positive here if, away from contact terms, its Fourier multiplier
has the representation
\[
 \widehat C(p)=F(p^2),
 \qquad
 F(z)=\int_{[0,\infty)}\frac{\mathrm d\rho(m^2)}{z+m^2}
\]
for a positive locally finite measure \(\rho\).  This is the
Kallen--Lehmann form supplied by OS reconstruction for a nonzero scalar
observable under the usual regularity assumptions.
\end{definition}

\begin{theorem}[Local higher-derivative scalar no-go]
\label{thm:v2v50-polynomial-no-go}
Let \(P\) be a real polynomial of degree \(d\ge2\), positive on
\((0,\infty)\), with positive leading coefficient.  Suppose that, with an
explicit infrared prescription if needed, the multiplier
\begin{equation}
 \widehat C(p)=\frac1{P(p^2)}
\label{eq:v2v50-polynomial-covariance}
\end{equation}
defines a nonzero positive-definite tempered Gaussian covariance.  Then it
is not spectral-positive.  Hence it cannot be the two-point function of a
nonzero OS-positive scalar observable satisfying the reconstruction
hypotheses in Definition~\ref{def:v2v50-spectral-positive}.
\end{theorem}

\begin{proof}
By hypothesis the infrared definition is a positive Euclidean covariance.
If it also had
the representation \eqref{eq:v2v50-Stieltjes}, then
Lemma~\ref{lem:v2v50-Stieltjes-lower} would give
\[
 \frac1{P(z)}\ge\frac c{z+R}.
\]
But \(P(z)=a_dz^d(1+o(1))\), and therefore
\(1/P(z)=O(z^{-d})=o(z^{-1})\), a contradiction.
\end{proof}

\begin{corollary}[The two-pole sign is explicit when the roots are real]
\label{cor:v2v50-two-pole}
If
\[
 P(z)=a(z+m_1^2)(z+m_2^2),
 \qquad
 a>0,\quad 0\le m_1<m_2,
\]
then
\begin{equation}
 \frac1{P(z)}
 =
 \frac1{a(m_2^2-m_1^2)}
 \left(
  \frac1{z+m_1^2}-\frac1{z+m_2^2}
 \right).
\label{eq:v2v50-negative-residue}
\end{equation}
The second spectral residue is negative.  A repeated root gives a double
pole, which likewise cannot be the Stieltjes transform of a positive
measure.
\end{corollary}

\begin{proof}
The displayed identity is partial fractions.  Positive atomic spectral
measures have nonnegative residues at their simple negative-axis poles.
For a positive measure, the transform in
\eqref{eq:v2v50-Stieltjes} has only first-order atomic poles; a double pole
cannot occur.
\end{proof}

\begin{theorem}[Operator-valued form of the obstruction]
\label{thm:v2v50-operator-no-go}
Let \(\Sigma\) be a nonzero positive operator-valued locally finite measure
on \([0,\infty)\), acting on a finite-dimensional physical multiplet
\(\mathcal H\), and set
\[
 C(z)=\int_{[0,\infty)}\frac{\mathrm d\Sigma(m^2)}{z+m^2}.
\]
Then \(\|C(z)\|\) is not \(O(z^{-1-\delta})\) for any \(\delta>0\).
Consequently a finite globally assembled physical covariance satisfying
\(\|C(z)\|=O(z^{-2})\) cannot have a positive
operator-valued Kallen--Lehmann measure unless it vanishes.
\end{theorem}

\begin{proof}
Choose \(v\in\mathcal H\) such that the scalar positive measure
\(\rho_v(B)=\langle v,\Sigma(B)v\rangle\) is nonzero.  Applying
Lemma~\ref{lem:v2v50-Stieltjes-lower} to
\(\langle v,C(z)v\rangle\) gives
\[
 \|C(z)\|\,\|v\|^2
 \ge \langle v,C(z)v\rangle
 \ge\frac c{z+R}.
\]
This contradicts faster decay.  The last assertion is the case
\(\delta=1\).
\end{proof}

\begin{assessment}[Why global assembly does not remove the obstruction]
\label{ass:v2v50-global-assembly}
Signed gauge-fixed components can cancel their \(z^{-1}\) terms and leave a
\(z^{-2}\) effective propagator.  The YM audit forbids reading that
componentwise cancellation as positivity.  If the assembled field is a
nonzero physical observable in a positive reconstructed Hilbert space,
Theorem~\ref{thm:v2v50-operator-no-go} applies to its positive
operator-valued spectral measure.  Thus the cancellation either occurs
outside the physical positive algebra, introduces a null observable, or
signals failure of OS positivity.
\end{assessment}

\subsection{Application to the real conformal stabilizer}
\label{sec:v2v50-conformal-application}

\begin{theorem}[The fourth-order conformal repair misses the OS gate]
\label{thm:v2v50-conformal-no-go}
Consider the translation-invariant tangent Gaussian model of the real
conformal mode with inverse covariance
\[
 P_a(p^2)
 =
 a(p^2)^2-\frac{3}{4\pi G}p^2+\frac{2\Lambda}{\pi G},
 \qquad a>0.
\]
Exactly one of the following holds.
\begin{enumerate}
 \item \(P_a\) is nonpositive at some retained momentum, so it does not
       define a strictly positive real Gaussian quadratic form there.
 \item \(P_a\) is positive on all retained momenta, so it defines a positive
       Euclidean Gaussian covariance, but that covariance is not
       spectral-positive and cannot furnish an OS-positive nonzero scalar
       conformal observable under the standard reconstruction hypotheses.
\end{enumerate}
Thus adding this local fourth-order term does not simultaneously establish
real-contour Gaussian coercivity and the OS positivity required by the V48
spectral-response route.
\end{theorem}

\begin{proof}
The alternatives exhaust positivity of the polynomial on the retained
spectrum.  In the second alternative,
\[
 P_a(z)^{-1}=a^{-1}z^{-2}(1+O(z^{-1}))
\]
as \(z\to\infty\).  Apply
Theorem~\ref{thm:v2v50-polynomial-no-go}.
\end{proof}

\begin{firewall}[Scope of the no-go theorem]
\label{fw:v2v50-scope}
Theorem~\ref{thm:v2v50-conformal-no-go} concerns a real scalar conformal
coordinate treated as a nonzero physical Gaussian observable with a local
polynomial inverse covariance and unit numerator.  It does not rule out:
\begin{enumerate}
 \item removal of the conformal coordinate by a proved constraint or
       physical-state quotient;
 \item a complex contour whose final gauge-invariant Schwinger functions
       are independently shown to be OS positive;
 \item an enlarged auxiliary-field system whose complete physical
       covariance has a positive spectral measure and \(z^{-1}\) leading
       behaviour;
 \item a nonlocal definition with a positive Stieltjes covariance;
 \item a Lorentzian construction that never asserts a positive probability
       measure on real Euclidean metrics.
\end{enumerate}
Each alternative needs its own construction.  Naming a contour or a
constraint is not a proof of positivity.
\end{firewall}

\begin{firewall}[Ultraviolet decay is not positive compactness stock]
\label{fw:v2v50-decay-not-stock}
The \(p^{-4}\) covariance improves some negative-Sobolev moments, but any
spectral decomposition producing that improvement must leave the class of
positive measures.  The NS audit
shows the analogous danger in a different setting: a weaker norm may be
finite while the weighted endpoint moment remains unpaid.  Therefore the
fourth-order decay cannot be inserted into the V49 tightness theorem as if
it also supplied the positive stress spectral moment used in V48.
\end{firewall}

\subsection{V50 conclusion and revised acquisition order}
\label{sec:v2v50-conclusion}

\begin{theorem}[What V50 proves]
\label{thm:v2v50-conclusion}
Relative to V49, this note proves:
\begin{enumerate}
 \item the exact conformal Hessian
       \eqref{eq:v2v50-hessian-polynomial};
 \item a positive-measure lower bound showing that every nonzero
       Kallen--Lehmann covariance has at least inverse-first-order
       high-momentum size;
 \item scalar and finite-multiplet no-go theorems for strictly local
       higher-derivative covariances with inverse-second-order decay;
 \item the resulting failure of a naive fourth-order real conformal
       stabilizer to pass both the coercivity and OS gates.
\end{enumerate}
\end{theorem}

\begin{proof}
The four assertions are respectively
Lemma~\ref{lem:v2v50-conformal-hessian},
Lemma~\ref{lem:v2v50-Stieltjes-lower},
Theorems~\ref{thm:v2v50-polynomial-no-go} and
\ref{thm:v2v50-operator-no-go}, and
Theorem~\ref{thm:v2v50-conformal-no-go}.
\end{proof}

The programme will not use local fourth-order damping as an unproved bridge
from Euclidean coercivity to Lorentzian positivity.  The next acquisition
is upstream: formulate compactness directly for positive Lorentzian states
on a local physical observable algebra, prove that positivity, Ward
identities, microlocal regularity, and stress uniform integrability pass to
weak-star limits, and only then couple the selected stress response to the
V40 causal Einstein solver.  A Euclidean route remains admissible only
after an explicit contour or constraint construction passes the complete
physical OS form.
\clearpage
\section*{Second Volume, Version V51: Lorentzian Physical-State Compactness and Hadamard Closure}
\addcontentsline{toc}{section}{Second Volume, Version V51: Lorentzian Physical-State Compactness and Hadamard Closure}

\subsection*{Purpose}

V50 showed that a local fourth-order real conformal Gaussian cannot be used
as an unproved bridge between Euclidean coercivity and OS positivity.  This
note develops the alternative announced there: compactness directly on a
Lorentzian physical observable algebra.  The result is not a construction
of the interacting Standard Model state.  It is a precise compactness
compiler which identifies the uniform estimates that would make such a
construction converge.

Three topologies must be kept separate.  Pointwise convergence on bounded
observables preserves positivity but does not control unbounded stress
moments.  Distributional convergence of \(n\)-point functions does not by
itself preserve a wavefront-set bound.  Finally, the existence of a
Hadamard limit does not by itself justify restriction of a composite kernel
to the diagonal.  The theorems below give the additional bound required at
each passage.

\subsection{A common physical topological algebra}
\label{sec:v2v51-algebra}

Fix a globally hyperbolic spacetime \((M,g)\).  Assume that all regulated
physical observables have been embedded into one Hausdorff locally convex
unital \(*\)-algebra \(\mathfrak A_{\rm phys}\).  Its topology is required
to contain the smeared fields, their finite products needed for the
renormalized stress tensor, and the Ward-completed physical combinations.
This common-algebra hypothesis is part of the theorem, not a consequence
of gauge fixing.

A Lorentzian state on \(\mathfrak A_{\rm phys}\) is a continuous complex
linear functional \(\omega\) satisfying
\[
 \omega(\mathbf1)=1,
 \qquad
 \omega(B^*B)\ge0
 \quad(B\in\mathfrak A_{\rm phys}).
\]
Let \(\sigma(\mathfrak A_{\rm phys}',\mathfrak A_{\rm phys})\) denote
pointwise convergence on the algebra.

\begin{theorem}[Equicontinuous physical-state compactness]
\label{thm:v2v51-state-compactness}
Let \((\omega_\lambda)\) be states on
\(\mathfrak A_{\rm phys}\).  Suppose there is a continuous seminorm \(p\)
such that
\begin{equation}
 |\omega_\lambda(B)|\le p(B)
 \qquad
 (\lambda,\ B\in\mathfrak A_{\rm phys}).
\label{eq:v2v51-equicontinuity}
\end{equation}
Then the family has a cluster point in
\(\sigma(\mathfrak A_{\rm phys}',\mathfrak A_{\rm phys})\).  Every cluster
point \(\omega\) is a continuous positive normalized state and obeys
\begin{equation}
 |\omega(B)|\le p(B).
\label{eq:v2v51-limit-bound}
\end{equation}
If the seminormed quotient
\(\mathfrak A_{\rm phys}/\ker p\), completed in \(p\), is separable, then
the relevant compact polar is metrizable and every sequence
\((\omega_n)\) has a pointwise convergent subsequence.
\end{theorem}

\begin{proof}
Let
\[
 U=\{B:p(B)\le1\}.
\]
Equation~\eqref{eq:v2v51-equicontinuity} places every \(\omega_\lambda\)
in the polar
\[
 U^\circ
 =
 \{\ell\in\mathfrak A_{\rm phys}':|\ell(B)|\le1
   \text{ for }B\in U\}.
\]
The Alaoglu--Bourbaki theorem makes \(U^\circ\) compact in the weak dual
topology.  Hence the family has a cluster point \(\omega\in U^\circ\), and
membership in the polar gives \eqref{eq:v2v51-limit-bound}.

For each fixed \(B\), evaluation
\(\ell\mapsto\ell(B)\) is weak-dual continuous.  Therefore
\[
 \omega(\mathbf1)=\lim_\lambda\omega_\lambda(\mathbf1)=1,
 \qquad
 \omega(B^*B)=\lim_\lambda\omega_\lambda(B^*B)\ge0.
\]
Thus normalization and positivity are closed.

For the last assertion, choose a countable \(p\)-dense subset of the
completed quotient.  On \(U^\circ\), values on that subset determine the
functional because
\[
 |\ell(B-C)|\le p(B-C).
\]
The weak topology on \(U^\circ\) is therefore generated by countably many
evaluations, so the compact set is metrizable and hence sequentially
compact.
\end{proof}

\begin{corollary}[Closed Ward and causal identities]
\label{cor:v2v51-closed-identities}
Let \(\omega_n\to\omega\) pointwise under
Theorem~\ref{thm:v2v51-state-compactness}.  If
\(D_\alpha\in\mathfrak A_{\rm phys}\) and
\[
 \omega_n(D_\alpha)\longrightarrow0
\]
for every fixed Ward-defect element \(D_\alpha\), then
\(\omega(D_\alpha)=0\).  The same statement holds for fixed commutator,
field-equation, BRST-cohomology, and local-covariance relations represented
by elements of the common algebra.
\end{corollary}

\begin{proof}
Pointwise convergence gives
\(\omega(D_\alpha)=\lim_n\omega_n(D_\alpha)=0\).
\end{proof}

\begin{firewall}[A common gauge-fixed representative is insufficient]
\label{fw:v2v51-common-algebra}
Theorem~\ref{thm:v2v51-state-compactness} applies after the physical
quotient and its regulator embeddings have been constructed.  A family of
BRST gauge-fixed algebras with cutoff-dependent cohomology is not yet a
common \(\mathfrak A_{\rm phys}\).  One must prove compatible comparison
maps, independence of representatives, and convergence of the Ward defects
before invoking the theorem.
\end{firewall}

\subsection{Why bounded-observable compactness does not control stress}
\label{sec:v2v51-unbounded}

\begin{proposition}[Bounded weak convergence can lose every quadratic moment]
\label{prop:v2v51-bounded-counterexample}
Let
\[
 \mu_n=\left(1-\frac1n\right)\delta_0+\frac1n\delta_n
\]
on \(\mathbb R\).  Then
\[
 \int f\,\mathrm d\mu_n\longrightarrow f(0)
\]
for every bounded continuous \(f\), while
\[
 \int x\,\mathrm d\mu_n=1,
 \qquad
 \int x^2\,\mathrm d\mu_n=n.
\]
Thus weak compactness of states on a \(C^*\)-algebra of bounded
observables does not imply convergence, or even boundedness, of a stress
second moment.
\end{proposition}

\begin{proof}
The bounded-observable claim follows from
\[
 \left|\int f\,\mathrm d\mu_n-f(0)\right|
 =
 \frac1n|f(n)-f(0)|
 \le\frac{2\|f\|_\infty}{n}.
\]
The moment identities are direct.
\end{proof}

\begin{assessment}[The seminorm must contain the composite insertions]
\label{ass:v2v51-stress-seminorm}
For the continuum programme, the seminorm \(p\) in
\eqref{eq:v2v51-equicontinuity} must control the renormalized elements
\[
 T(f),\qquad T(f)T(h),
\]
and the Ward-completed products entering the response.  Otherwise
Theorem~\ref{thm:v2v51-state-compactness} produces only a state on bounded
observables and says nothing about the stress source required by the
Einstein equation.  This is the Lorentzian form of the V49
uniform-integrability gate.
\end{assessment}

\subsection{Uniform microlocal estimates and wavefront closure}
\label{sec:v2v51-microlocal}

Let \(X\) be a smooth manifold and let
\(\Gamma\subset T^*X\setminus0\) be a closed cone.  For a distribution
\(u_n\), a chart cutoff \(\chi\), and a closed coordinate cone \(V\), write
\(\widehat{\chi u_n}\) for the localized Fourier transform.

\begin{definition}[Uniform exclusion of a microlocal cone]
\label{def:v2v51-uniform-microlocal}
The sequence \((u_n)\) is uniformly regular outside \(\Gamma\) if, whenever
\[
 (\operatorname{supp}\chi\times V)\cap\Gamma=\varnothing,
\]
then for every \(N\ge0\) there is \(C_{\chi,V,N}<\infty\), independent of
\(n\), such that
\begin{equation}
 \sup_{\xi\in V}
 (1+|\xi|)^N
 |\widehat{\chi u_n}(\xi)|
 \le C_{\chi,V,N}.
\label{eq:v2v51-uniform-wavefront}
\end{equation}
\end{definition}

\begin{theorem}[Uniform wavefront bounds are closed]
\label{thm:v2v51-wavefront-closure}
Suppose \(u_n\to u\) in \(\mathcal D'(X)\) and
\((u_n)\) satisfies
\eqref{eq:v2v51-uniform-wavefront}.  Then
\begin{equation}
 \operatorname{WF}(u)\subseteq\Gamma.
\label{eq:v2v51-wavefront-limit}
\end{equation}
\end{theorem}

\begin{proof}
Fix \(\chi,V,N\) as in the definition.  For every fixed \(\xi\), the
function
\[
 x\longmapsto\chi(x)e^{-ix\cdot\xi}
\]
is a fixed test function.  Distributional convergence therefore gives
\[
 \widehat{\chi u_n}(\xi)\longrightarrow
 \widehat{\chi u}(\xi).
\]
Taking the pointwise limit in
\eqref{eq:v2v51-uniform-wavefront} yields
\[
 (1+|\xi|)^N|\widehat{\chi u}(\xi)|
 \le C_{\chi,V,N}
 \quad(\xi\in V).
\]
This holds for every \(N\), which is exactly microlocal regularity outside
\(\Gamma\).
\end{proof}

\begin{corollary}[Hadamard closure with the fixed commutator]
\label{cor:v2v51-Hadamard-closure}
Let \(W_n\) be two-point distributions for one fixed normally hyperbolic
field, suppose
\[
 W_n-W_n^{\mathrm T}=iE,
 \qquad
 \operatorname{WF}(W_n)\subseteq\mathcal C^+,
\]
where \(E\) is the causal propagator and \(\mathcal C^+\) is the future
Hadamard cone.  Assume \(W_n\to W\) distributionally and the uniform
microlocal estimates of
Definition~\ref{def:v2v51-uniform-microlocal} hold with
\(\Gamma=\mathcal C^+\).  Then
\[
 W-W^{\mathrm T}=iE,
 \qquad
 \operatorname{WF}(W)=\mathcal C^+.
\]
Thus the limit is Hadamard.
\end{corollary}

\begin{proof}
The antisymmetric identity passes to the distributional limit.
Theorem~\ref{thm:v2v51-wavefront-closure} gives
\(\operatorname{WF}(W)\subseteq\mathcal C^+\), and transposition gives
\(\operatorname{WF}(W^{\mathrm T})\subseteq\mathcal C^-\).

The causal propagator has wavefront set
\(\mathcal C^+\cup\mathcal C^-\).  If a point of \(\mathcal C^+\) were
absent from \(\operatorname{WF}(W)\), it would also be absent from
\(\operatorname{WF}(W^{\mathrm T})\), because the future and past cones are
disjoint.  The identity
\(iE=W-W^{\mathrm T}\) and the wavefront-set sum rule would then make \(E\)
regular there, a contradiction.  Hence
\(\mathcal C^+\subseteq\operatorname{WF}(W)\), proving equality.
\end{proof}

\begin{firewall}[Qualitative Hadamard membership is not uniformity]
\label{fw:v2v51-qualitative}
Smooth approximate identities have empty wavefront set term by term and
can converge distributionally to a delta distribution.  Therefore
\(\operatorname{WF}(W_n)\subseteq\mathcal C^+\) for each \(n\) is not enough
for Corollary~\ref{cor:v2v51-Hadamard-closure}.  The uniform constants in
\eqref{eq:v2v51-uniform-wavefront} are an independent regulator estimate.
\end{firewall}

\subsection{Diagonal continuity for the renormalized stress}
\label{sec:v2v51-diagonal}

Let \(H_g\) be one fixed local Hadamard parametrix and write
\[
 S_n=W_n-H_g.
\]
For Hadamard \(W_n\), the remainder \(S_n\) is smooth near the diagonal.
Let \(\mathcal D_{\mu\nu}\) be the point-split stress operator, of
differential order at most two, including the fixed local curvature
counterterms.

\begin{lemma}[Distributional convergence plus one spare derivative]
\label{lem:v2v51-Cr-closure}
Let \(K\) be compact and suppose \(S_n\to S\) in
\(\mathcal D'\) on a neighborhood of \(K\).  If, for an integer \(r\ge0\),
\begin{equation}
 \sup_n\|S_n\|_{C^{r+1}(K)}<\infty,
\label{eq:v2v51-Cr-bound}
\end{equation}
then
\begin{equation}
 S_n\longrightarrow S
 \quad\hbox{in }C^r(K),
\label{eq:v2v51-Cr-convergence}
\end{equation}
and \(S\) is \(C^r\) on \(K\).
\end{lemma}

\begin{proof}
The bound \eqref{eq:v2v51-Cr-bound} makes all derivatives through order
\(r\) uniformly bounded and equicontinuous.  Arzela--Ascoli gives a
\(C^r(K)\)-convergent subsequence of every subsequence.  Its \(C^r\) limit
has the same distributional limit \(S\).  Thus every convergent subfamily
has the unique limit \(S\).  If the full sequence failed to converge in
\(C^r\), a subsequence staying a fixed distance from \(S\) would have a
further \(C^r\)-convergent subsequence, contradicting uniqueness.
\end{proof}

\begin{theorem}[Continuity of point-split stress expectations]
\label{thm:v2v51-stress-limit}
Let \(K_0\subset M\) be compact and apply
Lemma~\ref{lem:v2v51-Cr-closure} on a compact neighborhood of the diagonal
over \(K_0\), with \(r\ge2\).  Then
\begin{equation}
 \left.
 \mathcal D_{\mu\nu}S_n(x,y)
 \right|_{y=x}
 \longrightarrow
 \left.
 \mathcal D_{\mu\nu}S(x,y)
 \right|_{y=x}
\label{eq:v2v51-stress-diagonal}
\end{equation}
uniformly on \(K_0\).  Hence the renormalized stress one-point functions,
with the same allowed local counterterms, converge uniformly on \(K_0\).

For the connected stress two-point function the same conclusion holds if
the subtracted connected four-point remainders converge distributionally
and obey the corresponding uniform \(C^{r+1}\) estimate near the double
diagonal, with \(r\) at least the total point-split differential order.
\end{theorem}

\begin{proof}
Differentiation through order \(r\) and restriction of a \(C^r\) function
to the diagonal are continuous operations.  Apply
\eqref{eq:v2v51-Cr-convergence}.  The local counterterms are fixed smooth
tensors, so they pass unchanged.  The four-point assertion is the same
argument on \(M^4\) after the universal Hadamard subtractions have been
made.
\end{proof}

\begin{firewall}[The microlocal cone does not pay the diagonal seminorm]
\label{fw:v2v51-diagonal-bound}
Theorem~\ref{thm:v2v51-wavefront-closure} controls directions of
singularity.  It does not imply the numerical
\(C^{r+1}\) bound in \eqref{eq:v2v51-Cr-bound}.  V44 identified these smooth
state-remainder seminorms; V51 proves that they are sufficient for passage
to the stress limit but does not derive them for the interacting chiral
SM.
\end{firewall}

\subsection{Causal response is closed under distributional limits}
\label{sec:v2v51-causal}

Let \(\mathcal J^+\subset M\times M\) be the closed retarded relation and
let \(R_n\in\mathcal D'(M\times M)\) be the Ward-reduced retarded stress
responses.

\begin{proposition}[Retarded support and Ward closure]
\label{prop:v2v51-retarded-closure}
Suppose
\[
 R_n\longrightarrow R
 \quad\hbox{in }\mathcal D'(M\times M),
 \qquad
 \operatorname{supp}R_n\subseteq\mathcal J^+,
\]
and let \(\mathcal W\) be any fixed differential Ward operator.  If
\(\mathcal W R_n\to0\) distributionally, then
\begin{equation}
 \operatorname{supp}R\subseteq\mathcal J^+,
 \qquad
 \mathcal W R=0.
\label{eq:v2v51-retarded-Ward}
\end{equation}
\end{proposition}

\begin{proof}
If a test function is supported outside \(\mathcal J^+\), every \(R_n\)
annihilates it, and so does the limit.  This proves the support statement.
Differential operators are continuous on distributions, hence
\[
 \mathcal W R=\lim_n\mathcal W R_n=0.
\]
\end{proof}

\begin{assessment}[What remains beyond causal closure]
\label{ass:v2v51-Dini}
Proposition~\ref{prop:v2v51-retarded-closure} preserves support and the Ward
identity, but not the V48 weighted spectral moment.  To obtain the bounded
three-subtraction memory kernel, the positive physical spectral measures
must additionally satisfy the V48 weighted tightness and Dini hypotheses.
The bounded-state topology cannot replace that estimate, by
Proposition~\ref{prop:v2v51-bounded-counterexample}.
\end{assessment}

\subsection{A fixed-background Lorentzian compactness compiler}
\label{sec:v2v51-compiler}

\begin{theorem}[Conditional Lorentzian state-to-response limit]
\label{thm:v2v51-compiler}
On one fixed globally hyperbolic bounded-geometry slab, suppose a regulator
sequence supplies:
\begin{enumerate}
 \item compatible embeddings into a common
       \(\mathfrak A_{\rm phys}\) and the equicontinuity bound
       \eqref{eq:v2v51-equicontinuity};
 \item vanishing physical Ward, field-equation, and commutator defects;
 \item distributional convergence along subsequences of the two- and
       four-point kernels, together with the uniform microlocal estimates
       \eqref{eq:v2v51-uniform-wavefront};
 \item the uniform smooth-remainder bounds
       \eqref{eq:v2v51-Cr-bound} at the orders required by the stress
       one- and two-point functions;
 \item the V48 positive physical spectral weighted-tightness and Dini
       bounds;
 \item uniqueness of every subsequential state on
       \(\mathfrak A_{\rm phys}\).
\end{enumerate}
Then the full regulator sequence converges pointwise to one positive
physical Hadamard state.  Its renormalized stress one- and two-point
functions are the regulator limits, its response is Ward-compatible and
retarded, and its three-subtraction memory kernel satisfies the V48 bound.
Consequently this selected quantum source is admissible for the fixed
background V40 causal linear-response solver.
\end{theorem}

\begin{proof}
Theorem~\ref{thm:v2v51-state-compactness} gives state cluster points.
Corollary~\ref{cor:v2v51-closed-identities} preserves the algebraic
identities.  Theorem~\ref{thm:v2v51-wavefront-closure} and
Corollary~\ref{cor:v2v51-Hadamard-closure} give the Hadamard property.
Theorem~\ref{thm:v2v51-stress-limit} identifies the composite stress
limits, and Proposition~\ref{prop:v2v51-retarded-closure} preserves causal
support and Ward compatibility.  The V48 hypotheses give the spectral
Dini identity and bounded subtracted memory.  If the full sequence did not
converge to the unique state, some subnet outside a fixed weak
neighborhood would have a cluster point different from it, contradicting
uniqueness.  The final assertion is exactly the input condition of the V40
fixed-background Volterra construction.
\end{proof}

\begin{firewall}[The metric is still fixed]
\label{fw:v2v51-fixed-metric}
Theorem~\ref{thm:v2v51-compiler} does not yet select states uniformly as
\(g\) varies, identify algebras on different metrics, or solve the
semiclassical Einstein fixed-point problem.  Pointwise compactness for each
background can choose incompatible subsequences.  The next step requires a
background-parametric equicontinuity theorem before the quantum source may
be inserted into a nonlinear metric iteration.
\end{firewall}

\subsection{V51 conclusion and next acquisition}
\label{sec:v2v51-conclusion}

V51 replaces the unavailable positive Euclidean metric measure by a
rigorous fixed-background Lorentzian compactness route.  Positivity and
linear physical identities are closed under an equicontinuous weak-dual
limit.  Hadamard regularity survives only under uniform microlocal Fourier
bounds.  Composite stress tensors survive only under separate uniform
smooth-remainder bounds, and causal support survives distributional
limits.  These facts assemble into
Theorem~\ref{thm:v2v51-compiler}, while leaving the interacting estimates
explicitly open.

V52 will address the next upstream obstruction: construct a single
subsequence for a compact family of backgrounds by combining uniform
equicontinuity in the physical algebra with an Arzela--Ascoli modulus in
the metric parameter.  Only such a continuous state-selection map can
enter the nonlinear Einstein fixed-point problem without a hidden
background-dependent choice.
\clearpage
\section*{Second Volume, Version V52: Background-Parametric State Compactness and the Einstein Fixed-Point Compiler}
\addcontentsline{toc}{section}{Second Volume, Version V52: Background-Parametric State Compactness and the Einstein Fixed-Point Compiler}

\subsection*{Purpose}

V51 compactified regulator states on one fixed Lorentzian background.  A
nonlinear Einstein iteration requires more: the regulator limit must be
chosen simultaneously and continuously for every metric in the iteration
set.  Choosing a different subsequence at each metric does not define a
source map and cannot enter a fixed-point theorem.

This note proves a background-parametric compactness theorem.  It uses a
common physical-algebra trivialization, one equicontinuity seminorm, and a
uniform modulus in the metric parameter.  A parallel vector-valued
Arzela--Ascoli argument gives a continuous renormalized stress map.  These
two limits are then combined with a compact invariant Einstein solution
map.  The resulting fixed-point theorem is conditional on explicit
uniform estimates; it does not assert that the interacting SM regulator
already supplies them.

\subsection{A compact class of backgrounds and a common algebra chart}
\label{sec:v2v52-backgrounds}

Let \(\mathcal K\) be a compact metric space of Lorentzian metrics on one
fixed slab \(M_T\).  Assume every \(g\in\mathcal K\) is globally hyperbolic,
the time orientation and Cauchy surfaces are common, and the metrics stay
inside one bounded-geometry and strict-hyperbolicity class.

For each \(g\), let \(\mathfrak A_{\rm phys}(g)\) be the physical observable
algebra.  A \emph{common algebra chart} is a family of continuous unital
\(*\)-isomorphisms
\begin{equation}
 \tau_g:\mathfrak A_{\rm phys}(g)\longrightarrow\mathfrak A_*,
\label{eq:v2v52-algebra-chart}
\end{equation}
to one Hausdorff locally convex unital \(*\)-algebra
\(\mathfrak A_*\).  All states below are pulled back to
\(\mathfrak A_*\).  Continuity in \(g\) is measured only after this
specified identification.

\begin{firewall}[There is no canonical comparison hidden in the notation]
\label{fw:v2v52-no-canonical-chart}
Observable algebras on different metrics are not canonically identical.
Equation~\eqref{eq:v2v52-algebra-chart} must come from a proved local
covariance, time-slice, relative-Cauchy, or gauge-fixed comparison
construction.  Changing \(\tau_g\) changes the meaning of the metric
modulus below.  The theorem does not manufacture this chart.
\end{firewall}

Let \(p\) be a continuous seminorm on \(\mathfrak A_*\), and suppose the
completion of \(\mathfrak A_*/\ker p\) is separable.  Let
\[
 \mathscr S_p
 =
 \{\omega\in\mathfrak A_*':
   \omega\text{ is a state and }|\omega(B)|\le p(B)\}.
\]
By Theorem~\ref{thm:v2v51-state-compactness}, \(\mathscr S_p\) is compact
and metrizable in its weak-dual topology.

Choose a countable \(p\)-dense set \((B_j)_{j\ge1}\) in the closed
\(p\)-unit ball and define
\begin{equation}
 d_*(\omega,\eta)
 =
 \sum_{j=1}^\infty2^{-j}
 \min\{1,|\omega(B_j)-\eta(B_j)|\}.
\label{eq:v2v52-state-metric}
\end{equation}
This metric induces the weak-dual topology on \(\mathscr S_p\).

\subsection{One subsequence for every background}
\label{sec:v2v52-parametric-state}

Let
\[
 \omega_n:\mathcal K\longrightarrow\mathscr S_p
\]
be the \(n\)-th regulated state family.  Assume every \(\omega_n\) is
continuous and there is a modulus
\(\eta:[0,\infty)\to[0,\infty)\), independent of \(n\), with
\(\eta(r)\to0\) as \(r\downarrow0\), such that
\begin{equation}
 |\omega_n(g)(B)-\omega_n(h)(B)|
 \le p(B)\eta(d_{\mathcal K}(g,h))
\label{eq:v2v52-background-modulus}
\end{equation}
for all \(g,h\in\mathcal K\), \(B\in\mathfrak A_*\), and \(n\).

\begin{lemma}[The algebra modulus controls the state-space metric]
\label{lem:v2v52-metric-modulus}
For all \(g,h\in\mathcal K\),
\begin{equation}
 d_*(\omega_n(g),\omega_n(h))
 \le\min\{1,\eta(d_{\mathcal K}(g,h))\}.
\label{eq:v2v52-star-modulus}
\end{equation}
\end{lemma}

\begin{proof}
Since \(p(B_j)\le1\), every summand in
\eqref{eq:v2v52-state-metric} is at most
\[
 2^{-j}\min\{1,\eta(d_{\mathcal K}(g,h))\}.
\]
Sum the geometric series.
\end{proof}

\begin{theorem}[Background-parametric state compactness]
\label{thm:v2v52-parametric-state}
Under \eqref{eq:v2v52-background-modulus}, a subsequence
\((\omega_{n_k})\) and a continuous map
\[
 \omega_\infty:\mathcal K\longrightarrow\mathscr S_p
\]
exist such that
\begin{equation}
 \sup_{g\in\mathcal K}
 d_*(\omega_{n_k}(g),\omega_\infty(g))
 \longrightarrow0.
\label{eq:v2v52-uniform-state}
\end{equation}
For every fixed \(B\in\mathfrak A_*\),
\begin{equation}
 \sup_{g\in\mathcal K}
 |\omega_{n_k}(g)(B)-\omega_\infty(g)(B)|
 \longrightarrow0.
\label{eq:v2v52-uniform-evaluation}
\end{equation}
\end{theorem}

\begin{proof}
The target \((\mathscr S_p,d_*)\) is compact.  By
Lemma~\ref{lem:v2v52-metric-modulus}, the maps \(\omega_n\) are
equicontinuous.  The metric-space Arzela--Ascoli theorem therefore gives a
uniformly convergent subsequence with continuous limit
\(\omega_\infty\), proving \eqref{eq:v2v52-uniform-state}.

Fix \(B\) and \(\varepsilon>0\).  Choose a finite linear combination
\(C\) of the \(B_j\) such that \(p(B-C)<\varepsilon\).  The uniform
\(d_*\)-convergence gives uniform convergence of each of the finitely many
evaluations forming \(C\), while the \(p\)-bound gives
\[
 |(\omega_{n_k}(g)-\omega_\infty(g))(B-C)|
 \le2p(B-C)<2\varepsilon.
\]
Let \(k\to\infty\), then \(\varepsilon\downarrow0\).
\end{proof}

\begin{corollary}[Uniform closure of physical identities]
\label{cor:v2v52-uniform-Ward}
Let \(D_\alpha(g)\in\mathfrak A_*\) be a background-dependent defect family
with compact range in the \(p\)-topology.  If
\begin{equation}
 \sup_{g\in\mathcal K}
 |\omega_n(g)(D_\alpha(g))|
 \longrightarrow0,
\label{eq:v2v52-uniform-defect}
\end{equation}
then
\[
 \omega_\infty(g)(D_\alpha(g))=0
 \qquad(g\in\mathcal K).
\]
\end{corollary}

\begin{proof}
Given \(\varepsilon>0\), compactness of the range provides finitely many
\(p\)-balls of radius \(\varepsilon\) centered at
\(D_\alpha(g_1),\dots,D_\alpha(g_N)\).  The \(p\)-bound controls the error
from replacing \(D_\alpha(g)\) by a center, uniformly in \(g\).
Equation~\eqref{eq:v2v52-uniform-evaluation} controls the finitely many
centers, and \eqref{eq:v2v52-uniform-defect} controls the regulated defect.
Let first \(n\to\infty\), then \(\varepsilon\downarrow0\).
\end{proof}

\begin{corollary}[Pointwise uniqueness upgrades the whole family]
\label{cor:v2v52-full-family}
Suppose every uniformly convergent subsequence in
Theorem~\ref{thm:v2v52-parametric-state} has, at each \(g\), the same
prescribed state \(\omega(g)\).  Then
\[
 \sup_{g\in\mathcal K}d_*(\omega_n(g),\omega(g))\longrightarrow0,
\]
and \(g\mapsto\omega(g)\) is continuous.
\end{corollary}

\begin{proof}
If uniform convergence failed, some subsequence would remain at distance
at least \(\varepsilon>0\) from \(\omega\).  Theorem
\ref{thm:v2v52-parametric-state} gives a uniformly convergent further
subsequence.  Its pointwise limit must equal \(\omega\), a contradiction.
Continuity follows as a uniform limit of continuous maps.
\end{proof}

\begin{proposition}[Pointwise compactness alone is inadequate]
\label{prop:v2v52-xn}
The maps \(f_n(x)=x^n\) from \([0,1]\) to \([0,1]\) take values in a compact
space and converge at every \(x\), but their limit is discontinuous at
\(x=1\), and no subsequence converges uniformly.
\end{proposition}

\begin{proof}
The pointwise limit is zero on \([0,1)\) and one at \(1\).  A uniform limit
of continuous maps is continuous, so no subsequence can converge uniformly
to this pointwise limit.  Every subsequence has the same pointwise limit.
\end{proof}

\subsection{Uniform compactness of the renormalized stress family}
\label{sec:v2v52-stress-family}

Let \(K_0\Subset M_T\) and let \(\mathcal U\) be a fixed compact coordinate
neighborhood of the diagonal over \(K_0\).  Transport tensor indices using
the common background chart and write
\[
 S_n(g)=W_n(g)-H_g
\]
for the smooth two-point remainder.  Fix \(r\ge2\).

\begin{theorem}[Parametric diagonal compactness]
\label{thm:v2v52-parametric-stress}
Suppose
\begin{align}
 \sup_{n,g\in\mathcal K}
 \|S_n(g)\|_{C^{r+1}(\mathcal U)}
 &<\infty,
\label{eq:v2v52-uniform-Cr}\\
 \|S_n(g)-S_n(h)\|_{C^r(\mathcal U)}
 &\le\zeta(d_{\mathcal K}(g,h)),
\label{eq:v2v52-Cr-modulus}
\end{align}
where \(\zeta(a)\to0\) as \(a\downarrow0\).  Then a subsequence converges
in
\[
 C(\mathcal K;C^r(\mathcal U)).
\]
If the distributional state limit identifies every subsequential
remainder uniquely, the full sequence converges in this space.

Assume also that the point-split operators
\(\mathcal D_g\) and local counterterms depend continuously on \(g\) in
their required \(C^r\) coefficient norms.  Then the renormalized stresses
\begin{equation}
 \Theta_n(g)
 =
 \left.\mathcal D_gS_n(g)\right|_{\rm diag}
 +C_{\rm loc}(g)
\label{eq:v2v52-stress-map}
\end{equation}
converge uniformly in
\(C(\mathcal K;C^{r-2}(K_0))\) to a continuous map
\(\Theta_\infty\).
\end{theorem}

\begin{proof}
The bounded inclusion
\[
 C^{r+1}(\mathcal U)\hookrightarrow C^r(\mathcal U)
\]
is compact.  Thus the values of all \(S_n\) lie in one relatively compact
subset of \(C^r(\mathcal U)\).  Equation
\eqref{eq:v2v52-Cr-modulus} is equicontinuity in the metric parameter.
The vector-valued Arzela--Ascoli theorem gives compactness in
\(C(\mathcal K;C^r(\mathcal U))\).  Uniqueness upgrades subsequential
convergence by the argument of Corollary~\ref{cor:v2v52-full-family}.

A differential operator of order two maps \(C^r\) continuously to
\(C^{r-2}\), restriction to the diagonal is continuous, and the
coefficients vary continuously by hypothesis.  Applying these operations
to the uniform limit proves convergence and continuity of
\eqref{eq:v2v52-stress-map}.
\end{proof}

\begin{corollary}[The same statement for stress response]
\label{cor:v2v52-response-family}
Theorem~\ref{thm:v2v52-parametric-stress} applies to the subtracted
connected four-point remainder on \(M_T^4\), with \(r\) at least the total
point-split differential order.  Together with uniform versions of the V51
microlocal estimates and the V48 weighted spectral tightness over
\(g\in\mathcal K\), it produces a continuous family of Ward-reduced
retarded response kernels and bounded three-subtraction memory operators.
\end{corollary}

\begin{proof}
Use the same compact embedding and vector-valued Arzela--Ascoli theorem on
the double-diagonal neighborhood.  V51 preserves the wavefront cone,
retarded support, and Ward identity for each \(g\).  Uniform V48 weighted
tightness controls the memory norm uniformly in \(g\).
\end{proof}

\begin{firewall}[The metric modulus is a new analytic estimate]
\label{fw:v2v52-modulus-open}
Neither the fixed-background Hadamard condition nor pointwise regulator
convergence implies
\eqref{eq:v2v52-background-modulus} or
\eqref{eq:v2v52-Cr-modulus}.  These estimates require quantitative control
of relative Cauchy evolution, equivalently of integrated stress response
to a metric perturbation.  Assuming them without proof would insert the
desired continuity of the quantum source at the exact place where it is
needed.
\end{firewall}

\subsection{A continuum fixed-point passage}
\label{sec:v2v52-fixed-point}

Let \(X\) be a Banach space of gauge-fixed metrics on \(M_T\), and now
assume that the background class \(\mathcal K\subset X\) is nonempty,
compact, and convex.  Let \(Y\) be the source space containing the
renormalized stress and the V48 memory term.  Let
\[
 \mathcal E_n:\mathcal K\times Y\longrightarrow\mathcal K
\]
be the cutoff causal Einstein solution map with the prescribed initial
data, harmonic constraint correction, local running couplings, and
regulated nonlocal response.

\begin{theorem}[Uniform fixed-point limit]
\label{thm:v2v52-fixed-point}
Suppose:
\begin{enumerate}
 \item every \(\mathcal E_n\) is continuous;
 \item the state and stress hypotheses of
       Theorems~\ref{thm:v2v52-parametric-state} and
       \ref{thm:v2v52-parametric-stress} hold, giving continuous
       \(\Theta_n:\mathcal K\to Y\);
 \item the maps
       \[
        \Phi_n(g):=\mathcal E_n(g,\Theta_n(g))
       \]
       converge uniformly in \(X\) to a continuous
       \(\Phi:\mathcal K\to\mathcal K\).
\end{enumerate}
Then every \(\Phi_n\) has a fixed point \(g_n\in\mathcal K\).  Some
subsequence converges in \(X\) to \(g_\infty\in\mathcal K\), and
\begin{equation}
 g_\infty=\Phi(g_\infty).
\label{eq:v2v52-limit-fixed-point}
\end{equation}
If the limiting fixed point is unique, the full sequence \(g_n\) converges
to it.
\end{theorem}

\begin{proof}
Each \(\Phi_n\) is a continuous self-map of the nonempty compact convex
subset \(\mathcal K\) of a Banach space.  Schauder's theorem gives
\(g_n=\Phi_n(g_n)\).  Compactness of \(\mathcal K\) gives a subsequence,
still denoted \(g_n\), with \(g_n\to g_\infty\).  Uniform convergence and
continuity of \(\Phi\) imply
\begin{align*}
 \|g_\infty-\Phi(g_\infty)\|_X
 &\le
 \|g_\infty-g_n\|_X
 +\|\Phi_n(g_n)-\Phi(g_n)\|_X\\
 &\quad+
 \|\Phi(g_n)-\Phi(g_\infty)\|_X
 \longrightarrow0.
\end{align*}
This proves \eqref{eq:v2v52-limit-fixed-point}.  If the fixed point is
unique, every convergent subsequence has the same limit; compactness then
forces convergence of the full sequence.
\end{proof}

\begin{corollary}[One route from the V15 slab to a compact iteration set]
\label{cor:v2v52-V15-route}
Suppose the cutoff solutions have the uniform V15 order-\(s\) hyperbolic
bound and a uniform time-regularity bound strong enough for compact
embedding into
\[
 X=C([0,T_*];H^{s-1}).
\]
If their closed convex hull stays strictly inside the common
hyperbolicity region and the causal solution map preserves that hull, its
closure is an admissible compact convex set \(\mathcal K\) for
Theorem~\ref{thm:v2v52-fixed-point}.
\end{corollary}

\begin{proof}
The assumed space-time compact embedding makes the family relatively
compact in \(X\).  In a Banach space, the closed convex hull of a compact
set is compact.  The strict interior and invariance assumptions place the
solution map on that compact convex set.
\end{proof}

\begin{firewall}[The fixed-point theorem is a compiler, not the missing estimate]
\label{fw:v2v52-fixed-point-open}
Theorem~\ref{thm:v2v52-fixed-point} does not prove that
\(\Phi_n(\mathcal K)\subseteq\mathcal K\), that the V15 high norm remains
uniform under quantum backreaction, or that
\(\Phi_n\to\Phi\) uniformly.  Those are the nonlinear continuum estimates.
The theorem proves that once they and the state-family moduli are supplied,
no further background-dependent subsequence obstruction remains.
\end{firewall}

\subsection{V52 conclusion and next acquisition}
\label{sec:v2v52-conclusion}

V52 proves that one regulator subsequence can serve an entire compact metric
family if the pulled-back physical states obey one algebra seminorm and one
background modulus.  Uniform diagonal bounds produce a continuous
renormalized stress map, and uniform convergence of the resulting causal
Einstein maps carries cutoff fixed points to a continuum fixed point.
Every remaining hypothesis is stated in a topology strong enough for the
operation it controls.

The sharp new bottleneck is
\eqref{eq:v2v52-background-modulus}.  V53 will derive a quantitative
candidate for that modulus from relative Cauchy evolution: an integrated
Ward-completed retarded stress response gives Lipschitz dependence of
physical observables on the background whenever its operator norm is
uniformly integrable.  This will connect the V48 spectral moment and V40
Volterra kernel directly to the parametric compactness gate rather than
assuming continuity twice.
\clearpage
\section*{Second Volume, Version V53: Relative Cauchy Response Generates the Background Modulus}
\addcontentsline{toc}{section}{Second Volume, Version V53: Relative Cauchy Response Generates the Background Modulus}

\subsection*{Purpose}

V52 required a regulator-uniform modulus for the state and stress as the
background metric varies.  V27 already derived the local-contact plus
retarded-commutator formula for the metric derivative of the mean stress.
This note converts that differential response into the precise
equicontinuity used by V52.

The weak state topology and the strong stress topology require different
bounds.  For weak state compactness it is enough to control the response
of a countable dense family of observables, with a different finite
constant for each observable.  For the Einstein source one needs a single
operator norm in the chosen source space.  Keeping these two statements
separate avoids imposing an unnecessarily strong norm on the entire
observable algebra while still paying the strong estimate required by the
metric equation.

\subsection{Quasiconvex background charts}
\label{sec:v2v53-paths}

Let \(\mathcal K\) be the background class of V52, contained in a Banach
space \(X\) of metric perturbations.  Assume that \(\mathcal K\) is
\(Q\)-quasiconvex: every \(g,h\in\mathcal K\) can be joined by an absolutely
continuous curve
\[
 \gamma:[0,1]\longrightarrow\mathcal K,
 \qquad
 \gamma(0)=g,\quad\gamma(1)=h,
\]
whose \(X\)-length obeys
\begin{equation}
 \int_0^1\|\dot\gamma(s)\|_X\,\mathrm ds
 \le Q\|g-h\|_X.
\label{eq:v2v53-quasiconvex}
\end{equation}
A convex set has \(Q=1\) by the straight segment.

Retain the common algebra chart \(\tau_g\), the seminorm \(p\), the dense
\(p\)-unit-ball sequence \((B_j)\), and the state metric \(d_*\) from V52.
For the cutoff state \(\omega_n(g)\), define the complete transported
metric-response functional
\[
 \mathfrak R_{n,g}(B;h).
\]
It includes the relative-Cauchy commutator with the Ward-completed stress,
the derivative of the algebra chart, and any explicit background
dependence of the representative \(B\).

\subsection{Countable response bounds imply weak-state equicontinuity}
\label{sec:v2v53-state-modulus}

\begin{theorem}[Response-to-state-modulus theorem]
\label{thm:v2v53-state-modulus}
Suppose that for every \(n,j\) and every admissible path \(\gamma\), the
function
\[
 F_{n,j}(s):=\omega_n(\gamma(s))(B_j)
\]
is absolutely continuous and satisfies, for almost every \(s\),
\begin{align}
 \frac{\mathrm d}{\mathrm ds}F_{n,j}(s)
 &=
 \mathfrak R_{n,\gamma(s)}(B_j;\dot\gamma(s)),
\label{eq:v2v53-rce-identity}\\
 |\mathfrak R_{n,g}(B_j;k)|
 &\le L_j\|k\|_X,
\label{eq:v2v53-test-response}
\end{align}
where \(0\le L_j<\infty\) is independent of \(n\) and \(g\).  Define
\begin{equation}
 \eta_*(r)
 :=
 \sum_{j=1}^\infty2^{-j}\min\{1,L_jr\}.
\label{eq:v2v53-eta}
\end{equation}
Then \(\eta_*(r)\to0\) as \(r\downarrow0\), and
\begin{equation}
 d_*(\omega_n(g),\omega_n(h))
 \le
 \eta_*(Q\|g-h\|_X)
\label{eq:v2v53-state-Lipschitz}
\end{equation}
for every \(n\) and \(g,h\in\mathcal K\).  Consequently the
background-parametric compactness conclusion of
Theorem~\ref{thm:v2v52-parametric-state} holds without assuming the
stronger all-observable estimate
\eqref{eq:v2v52-background-modulus}.
\end{theorem}

\begin{proof}
Integrate \eqref{eq:v2v53-rce-identity} along a path satisfying
\eqref{eq:v2v53-quasiconvex}.  Equation
\eqref{eq:v2v53-test-response} gives
\[
 |F_{n,j}(1)-F_{n,j}(0)|
 \le
 L_j\int_0^1\|\dot\gamma(s)\|_X\,\mathrm ds
 \le QL_j\|g-h\|_X.
\]
Insert this estimate into the definition
\eqref{eq:v2v52-state-metric} and sum to obtain
\eqref{eq:v2v53-state-Lipschitz}.

For each fixed \(j\),
\(\min\{1,L_jr\}\to0\) as \(r\downarrow0\), and every summand is bounded by
\(2^{-j}\).  Dominated convergence on the counting measure proves
\(\eta_*(r)\to0\).  Thus the maps \(g\mapsto\omega_n(g)\) are
equicontinuous into the compact metric state space
\((\mathscr S_p,d_*)\), which is the exact input used by the
Arzela--Ascoli proof of Theorem~\ref{thm:v2v52-parametric-state}.
\end{proof}

\begin{corollary}[A single derivation norm gives a Lipschitz modulus]
\label{cor:v2v53-uniform-derivation}
If the complete response is represented by a derivation
\(\delta_{n,g;k}\) and
\begin{equation}
 p(\delta_{n,g;k}B)
 \le L\|k\|_Xp(B)
\label{eq:v2v53-derivation-bound}
\end{equation}
uniformly in \(n,g\), then
\[
 d_*(\omega_n(g),\omega_n(h))
 \le\min\{1,LQ\|g-h\|_X\}.
\]
\end{corollary}

\begin{proof}
The state bound gives
\[
 |\mathfrak R_{n,g}(B;k)|
 =
 |\omega_n(g)(\delta_{n,g;k}B)|
 \le p(\delta_{n,g;k}B).
\]
Use \(p(B_j)\le1\) in
Theorem~\ref{thm:v2v53-state-modulus}, with \(L_j=L\).
\end{proof}

\begin{lemma}[An integrated response stock is sufficient]
\label{lem:v2v53-integrated}
The pointwise bound \eqref{eq:v2v53-test-response} may be replaced by the
following path-segment estimate.  Suppose that for every admissible
subpath of \(X\)-length at most \(r\),
\begin{equation}
 \left|
 \int_a^b
 \mathfrak R_{n,\gamma(s)}(B_j;\dot\gamma(s))\,\mathrm ds
 \right|
 \le L_j\alpha(r),
\label{eq:v2v53-integrated-response}
\end{equation}
where \(\alpha(r)\to0\) uniformly in \(n,\gamma,a,b\).  Then
\[
 d_*(\omega_n(g),\omega_n(h))
 \le
 \sum_{j=1}^\infty2^{-j}
 \min\{1,L_j\alpha(Q\|g-h\|_X)\}.
\]
\end{lemma}

\begin{proof}
Apply \eqref{eq:v2v53-integrated-response} to the full connecting path and
repeat the last summation and dominated-convergence argument in
Theorem~\ref{thm:v2v53-state-modulus}.
\end{proof}

\begin{assessment}[Why countably many constants suffice]
\label{ass:v2v53-countable}
No uniform bound on all \(L_j\) is needed for weak-state equicontinuity.
The summable weights in \(d_*\) and dominated convergence absorb their
growth.  This does not weaken the V51 equicontinuity bound \(p\), which
still controls compactness and extension from the dense test algebra.
It only separates compactness in the weak state topology from the stronger
operator bound required for the stress source.
\end{assessment}

\subsection{The relative-Cauchy formula and gauge directions}
\label{sec:v2v53-rce}

At finite cutoff the transported observable derivative has the schematic
identity
\begin{equation}
 \mathfrak R_{n,g}(B;h)
 =
 -\frac i2
 \int_{M_T}
 \omega_n(g)\!\left(
 [T_n^{\mu\nu}(x),B]
 \right)
 h_{\mu\nu}(x)\,\mathrm d\operatorname{vol}_g(x)
 +
 \omega_n(g)(\nabla_h^{\rm chart}B),
\label{eq:v2v53-rce}
\end{equation}
with retarded support determined by the relative-Cauchy transport.  The
last term is the derivative of the chosen common algebra chart and of any
explicit metric dependence.  Formula \eqref{eq:v2v53-rce} is not used
without proving it for the regulator; it specifies the object whose norm
is required in \eqref{eq:v2v53-test-response}.

\begin{proposition}[Ward conservation removes compactly supported vertical directions]
\label{prop:v2v53-vertical}
Let \(h=\mathcal L_\xi g\), where \(\xi\) is compactly supported in the
interior of the slab.  If the renormalized physical stress satisfies
\[
 \nabla_\mu T_n^{\mu\nu}=0
\]
as an operator identity on the physical algebra and the chart connection
is diffeomorphism covariant, then
\begin{equation}
 \mathfrak R_{n,g}(B;\mathcal L_\xi g)=0
\label{eq:v2v53-vertical-zero}
\end{equation}
for every diffeomorphism-invariant physical observable \(B\).
Thus the modulus descends to the local metric quotient.
\end{proposition}

\begin{proof}
Using symmetry of the stress tensor and compact support,
\begin{align*}
 \frac12\int T_n^{\mu\nu}
 (\mathcal L_\xi g)_{\mu\nu}\,\mathrm d\operatorname{vol}_g
 &=
 \int T_n^{\mu\nu}\nabla_\mu\xi_\nu\,
 \mathrm d\operatorname{vol}_g\\
 &=
 -\int(\nabla_\mu T_n^{\mu\nu})\xi_\nu\,
 \mathrm d\operatorname{vol}_g
 =0.
\end{align*}
The commutator term in \eqref{eq:v2v53-rce} therefore vanishes.  Covariance
of the chart makes its vertical derivative equal to the infinitesimal
diffeomorphism action on \(B\), which is zero for a physical invariant.
\end{proof}

\begin{firewall}[Expectation-level conservation is not enough for every observable]
\label{fw:v2v53-operator-Ward}
The identity
\(\nabla_\mu\omega_n(T_n^{\mu\nu})=0\) controls the mean stress but does not
make the commutator in \eqref{eq:v2v53-rce} vanish in a gauge direction.
Proposition~\ref{prop:v2v53-vertical} needs the Ward identity inside
time-ordered and commutator insertions, including contact terms.  This is
the complete physical Ward condition emphasized in V43 and the V50 audit.
\end{firewall}

\subsection{The strong stress modulus from the V27 response}
\label{sec:v2v53-stress-modulus}

Let \(X_s\) be a metric space strong enough for four derivatives and let
\(Y_{s-4}\) be the chosen stress-source space.  Write the cutoff derivative
from V27 as
\begin{equation}
 D\Theta_n[g](h)
 =
 L_{n,g}h+\mathcal R_{n,g}^{\rm ret}h,
\label{eq:v2v53-stress-derivative}
\end{equation}
where \(L_{n,g}\) contains the Ward-completed local contact and running
coupling terms, and \(\mathcal R_{n,g}^{\rm ret}\) is the renormalized
nonlocal response.

\begin{theorem}[Response norm implies the V52 stress modulus]
\label{thm:v2v53-stress-modulus}
Assume \(\mathcal K\) is \(Q\)-quasiconvex in \(X_s\), every
\(\Theta_n:\mathcal K\to Y_{s-4}\) is differentiable along admissible
paths, and
\begin{align}
 \sup_{n,g\in\mathcal K}
 \|L_{n,g}\|_{\mathcal L(X_s,Y_{s-4})}
 &\le C_{\rm loc},
\label{eq:v2v53-local-bound}\\
 \sup_{n,g\in\mathcal K}
 \|\mathcal R_{n,g}^{\rm ret}\|_{\mathcal L(X_s,Y_{s-4})}
 &\le C_{\rm mem}.
\label{eq:v2v53-memory-bound}
\end{align}
Then
\begin{equation}
 \|\Theta_n(g)-\Theta_n(h)\|_{Y_{s-4}}
 \le
 Q(C_{\rm loc}+C_{\rm mem})
 \|g-h\|_{X_s}.
\label{eq:v2v53-stress-Lipschitz}
\end{equation}
In particular \eqref{eq:v2v52-Cr-modulus} holds in any local
\(C^r\) topology continuously dominated by \(Y_{s-4}\).
\end{theorem}

\begin{proof}
Along a connecting path \(\gamma\), the Banach-space fundamental theorem
of calculus and \eqref{eq:v2v53-stress-derivative} give
\[
 \Theta_n(h)-\Theta_n(g)
 =
 \int_0^1
 \bigl(L_{n,\gamma(s)}
       +\mathcal R_{n,\gamma(s)}^{\rm ret}\bigr)
 \dot\gamma(s)\,\mathrm ds.
\]
Take the \(Y_{s-4}\)-norm, apply
\eqref{eq:v2v53-local-bound}--\eqref{eq:v2v53-memory-bound}, and then
\eqref{eq:v2v53-quasiconvex}.
\end{proof}

\begin{corollary}[How the V48--V40 estimates enter]
\label{cor:v2v53-V48-V40}
Suppose uniformly over \(n\) and \(g\in\mathcal K\):
\begin{enumerate}
 \item the complete physical spectral measure satisfies the V48 Dini and
       weighted-tightness bounds;
 \item the three-subtracted kernel acts with norm at most \(M\) on the
       V40 Volterra domain after its exterior derivatives are restored;
 \item the finite local running couplings and bounded-geometry coefficient
       jets give a contact-operator bound \(C\).
\end{enumerate}
Then \eqref{eq:v2v53-local-bound} and
\eqref{eq:v2v53-memory-bound} hold with
\(C_{\rm loc}=C\) and \(C_{\rm mem}=M\), and therefore
\[
 \operatorname{Lip}(\Theta_n)
 \le Q(C+M).
\]
\end{corollary}

\begin{proof}
The local block is a finite differential operator of net order at most
four, so the stated coefficient bound gives its
\(\mathcal L(X_s,Y_{s-4})\)-norm.  V48 makes the three-subtracted memory
kernel integrable, and the stated V40 domain estimate includes the exterior
derivatives, giving norm \(M\).  Apply
Theorem~\ref{thm:v2v53-stress-modulus}.
\end{proof}

\begin{firewall}[Dini integrability alone is not the full source norm]
\label{fw:v2v53-Dini-domain}
The V48 scalar Dini integral bounds the subtracted memory kernel before it
is composed with all exterior derivatives and tensor projectors.  The
second hypothesis of Corollary~\ref{cor:v2v53-V48-V40} is therefore
essential.  Dropping it would confuse an \(L^1\) kernel bound with the
complete operator norm from \(X_s\) to \(Y_{s-4}\).
\end{firewall}

\subsection{Uniform convergence of the Einstein maps}
\label{sec:v2v53-Einstein-map}

V52 assumed uniform convergence of the composite cutoff maps
\(\Phi_n(g)=\mathcal E_n(g,\Theta_n(g))\).  The next lemma separates this
into a classical solver error and the quantum-source error.

\begin{lemma}[Source convergence composes uniformly with the causal solver]
\label{lem:v2v53-composition}
Let \(\Theta_n,\Theta:\mathcal K\to Y\) satisfy
\[
 \sup_{g\in\mathcal K}
 \|\Theta_n(g)-\Theta(g)\|_Y\longrightarrow0.
\]
Suppose \(\mathcal E_n,\mathcal E:\mathcal K\times Y\to X\) obey, on a
common set containing all source ranges,
\begin{align}
 \sup_{g,y}\|\mathcal E_n(g,y)-\mathcal E(g,y)\|_X
 &\le\varepsilon_n,\qquad \varepsilon_n\to0,
\label{eq:v2v53-solver-consistency}\\
 \|\mathcal E(g,y)-\mathcal E(g,z)\|_X
 &\le C_E\|y-z\|_Y.
\label{eq:v2v53-solver-source-Lipschitz}
\end{align}
Then
\begin{equation}
 \sup_{g\in\mathcal K}
 \|\mathcal E_n(g,\Theta_n(g))
   -\mathcal E(g,\Theta(g))\|_X
 \le
 \varepsilon_n+
 C_E\sup_g\|\Theta_n(g)-\Theta(g)\|_Y
 \longrightarrow0.
\label{eq:v2v53-composite-convergence}
\end{equation}
\end{lemma}

\begin{proof}
Add and subtract \(\mathcal E(g,\Theta_n(g))\), then use
\eqref{eq:v2v53-solver-consistency} and
\eqref{eq:v2v53-solver-source-Lipschitz}.
\end{proof}

\begin{corollary}[Response bounds discharge the continuity part of the V52 compiler]
\label{cor:v2v53-fixed-point-interface}
Assume the state convergence of V52, the stress convergence of
Theorem~\ref{thm:v2v52-parametric-stress}, the uniform response bounds of
Theorem~\ref{thm:v2v53-stress-modulus}, and the classical solver estimates
\eqref{eq:v2v53-solver-consistency}--%
\eqref{eq:v2v53-solver-source-Lipschitz}.  If the composite maps preserve
the compact convex set \(\mathcal K\), then the cutoff fixed points have a
subsequence converging to a continuum fixed point as in
Theorem~\ref{thm:v2v52-fixed-point}.
\end{corollary}

\begin{proof}
The response theorem supplies equicontinuity, the V52 stress theorem
supplies uniform source convergence, and
Lemma~\ref{lem:v2v53-composition} supplies the uniform composite-map
convergence required in V52.  Apply
Theorem~\ref{thm:v2v52-fixed-point}.
\end{proof}

\subsection{Sharp obstruction and V53 conclusion}
\label{sec:v2v53-conclusion}

The example \(f_n(x)=x^n\) from V52 has
\[
 \sup_{x\in[0,1]}|f_n'(x)|=n.
\]
Thus compact pointwise target values do not prevent the response norm from
diverging, and the loss of equicontinuity is exactly visible in the
derivative.  This shows that a uniform relative-Cauchy response estimate is
not technical decoration; it excludes the known background-selection
failure.

V53 proves that countable physical response bounds generate the weak state
modulus needed for a single background-uniform subsequence.  A complete
stress-response operator bound generates the strong source modulus, and
uniform source convergence composes with the causal Einstein solver to
supply the V52 fixed-point convergence hypothesis.  Ward conservation also
removes compactly supported diffeomorphism directions, provided it holds
inside the complete physical insertions.

The remaining hard estimate is now localized: construct the common chiral
SM physical algebra and prove the regulator-uniform bounds
\eqref{eq:v2v53-test-response} and
\eqref{eq:v2v53-memory-bound} on one invariant hyperbolic metric class.
V54 will go upstream to the algebra chart itself.  It will formulate a
cocycle criterion for cutoff comparison maps and prove that summable
relative-Cauchy shell derivations produce a regulator-independent physical
algebra trivialization, while recording the precise anomaly obstruction.
\clearpage
\section*{Second Volume, Version V54: Summable Shell Transports and the Physical Algebra Chart}
\addcontentsline{toc}{section}{Second Volume, Version V54: Summable Shell Transports and the Physical Algebra Chart}

\subsection*{Purpose}

V52 assumed a common physical algebra chart across backgrounds, and V53
included the derivative of that chart in the relative-Cauchy response.
This note gives a sufficient construction.  After the finite-dimensional
local running part has been extracted, suppose each ultraviolet shell
produces a bounded Ward-compatible derivation on one reference cylinder
algebra.  Absolute summability of those derivations makes the ordered
product of shell transports converge in operator norm, uniformly over the
metric class.  The limit is an invertible \(*\)-automorphism and supplies
the required chart.

The statement also distinguishes equality on the gauge-fixed algebra from
equality on the physical cohomology.  Two regulator charts can differ as
operators yet induce the same physical chart when their difference has a
summable chain homotopy.  Conversely, a shell transport which is not a
chain map does not act on BRST cohomology.  This is where a genuine anomaly
obstructs the construction.

\subsection{The renormalized reference algebra}
\label{sec:v2v54-reference}

Let \(\mathfrak A_*\) be a unital Banach \(*\)-algebra with submultiplicative
norm and isometric involution.  It is the completion of a common algebraic
cylinder algebra after local counterterms and running couplings have been
placed in fixed renormalization coordinates.  Let \(\mathcal K\) be the
compact background class of V52.

For each shell \(n\ge1\) and \(g\in\mathcal K\), let
\[
 \delta_{n,g}:\mathfrak A_*\longrightarrow\mathfrak A_*
\]
be a bounded \(*\)-derivation:
\begin{align}
 \delta_{n,g}(AB)
 &=
 \delta_{n,g}(A)B+A\delta_{n,g}(B),
\label{eq:v2v54-Leibniz}\\
 \delta_{n,g}(A^*)
 &=
 \delta_{n,g}(A)^*.
\label{eq:v2v54-star}
\end{align}
Define the shell transport
\begin{equation}
 \alpha_{n,g}:=\exp(\delta_{n,g})
 =
 \sum_{k=0}^\infty\frac{\delta_{n,g}^k}{k!}.
\label{eq:v2v54-shell-exponential}
\end{equation}

\begin{lemma}[A bounded derivation exponentiates to an automorphism]
\label{lem:v2v54-exponential}
For every \(n,g\), \(\alpha_{n,g}\) is a bounded unital
\(*\)-automorphism with inverse \(\exp(-\delta_{n,g})\), and
\begin{equation}
 \|\alpha_{n,g}^{\pm1}\|
 \le e^{\|\delta_{n,g}\|}.
\label{eq:v2v54-exponential-bound}
\end{equation}
\end{lemma}

\begin{proof}
For \(t\in\mathbb R\), put \(\alpha(t)=e^{t\delta_{n,g}}\).  The two
functions
\[
 t\longmapsto\alpha(t)(AB)
 \quad\hbox{and}\quad
 t\longmapsto\alpha(t)(A)\alpha(t)(B)
\]
solve the same Banach-space initial-value problem
\[
 F'(t)=\delta_{n,g}F(t),\qquad F(0)=AB,
\]
where the Leibniz rule is used for the second function.  Uniqueness gives
multiplicativity.  The same argument using
\eqref{eq:v2v54-star} gives preservation of the involution, and
\(\delta_{n,g}(\mathbf1)=0\) gives unitality.  The exponential series gives
the norm bound and
\(e^{\delta_{n,g}}e^{-\delta_{n,g}}=I\).
\end{proof}

\subsection{Uniform convergence of the shell product}
\label{sec:v2v54-product}

Order the shells from infrared to ultraviolet and set
\begin{equation}
 U_{N,g}
 :=
 \alpha_{N,g}\alpha_{N-1,g}\cdots\alpha_{1,g}.
\label{eq:v2v54-partial-product}
\end{equation}
No commutativity between different shell derivations is assumed.

\begin{theorem}[Uniform infinite-product algebra chart]
\label{thm:v2v54-product}
Suppose
\begin{equation}
 S:=
 \sum_{n=1}^\infty
 \sup_{g\in\mathcal K}\|\delta_{n,g}\|
 <\infty.
\label{eq:v2v54-summability}
\end{equation}
Then \(U_{N,g}\) converges in operator norm, uniformly in \(g\), to a
bounded unital \(*\)-automorphism \(U_g\).  The inverse products converge
uniformly to \(U_g^{-1}\), and
\begin{equation}
 \sup_{g\in\mathcal K}
 \bigl(\|U_g\|+\|U_g^{-1}\|\bigr)
 \le2e^S.
\label{eq:v2v54-chart-bound}
\end{equation}
More precisely, for \(N>M\),
\begin{equation}
 \sup_g\|U_{N,g}-U_{M,g}\|
 \le
 e^S
 \left[
  \exp\!\left(
   \sum_{n=M+1}^N\sup_g\|\delta_{n,g}\|
  \right)-1
 \right].
\label{eq:v2v54-product-tail}
\end{equation}
\end{theorem}

\begin{proof}
Write
\[
 U_{N,g}
 =
 \bigl(\alpha_{N,g}\cdots\alpha_{M+1,g}\bigr)U_{M,g}.
\]
For bounded operators \(D_1,\dots,D_k\), the product estimate
\begin{equation}
 \left\|
 \prod_{j=k}^1e^{D_j}-I
 \right\|
 \le
 \exp\!\left(\sum_{j=1}^k\|D_j\|\right)-1
\label{eq:v2v54-product-identity-bound}
\end{equation}
follows by expanding the product of exponential series and deleting the
constant term.  Also \(\|U_{M,g}\|\le e^S\) by
Lemma~\ref{lem:v2v54-exponential}.  These facts give
\eqref{eq:v2v54-product-tail}.  The tail in
\eqref{eq:v2v54-summability} tends to zero, so \(U_{N,g}\) is uniformly
Cauchy.

Each \(U_{N,g}\) is multiplicative, unital, and involution preserving.
Those identities pass to the operator-norm limit because multiplication
in \(\mathfrak A_*\) is continuous.  Apply the same argument to
\[
 V_{N,g}
 =
 \alpha_{1,g}^{-1}\cdots\alpha_{N,g}^{-1}.
\]
The limits \(U_g,V_g\) satisfy \(U_gV_g=V_gU_g=I\), by passage to the
finite-product identities.  Thus \(V_g=U_g^{-1}\), and the norm bounds
follow from \eqref{eq:v2v54-exponential-bound}.
\end{proof}

\begin{corollary}[Construction of the V52 chart]
\label{cor:v2v54-chart}
Let the completed algebra at background \(g\) be the transport of the
reference completion by \(U_g\).  Then
\begin{equation}
 \tau_g:=U_g^{-1}:
 \mathfrak A_{\rm phys}(g)\longrightarrow\mathfrak A_*
\label{eq:v2v54-chart}
\end{equation}
is a uniformly bounded common algebra chart.  Finite-cutoff observables
transported through \(U_{N,g}\) converge in norm to their continuum
representatives.
\end{corollary}

\begin{proof}
Theorem~\ref{thm:v2v54-product} supplies an invertible
\(*\)-automorphism.  Transporting the multiplication, involution, and norm
through it gives the background completion and the chart.  For
\(A\in\mathfrak A_*\),
\[
 \|U_{N,g}A-U_gA\|
 \le\|U_{N,g}-U_g\|\,\|A\|\longrightarrow0
\]
uniformly in \(g\).
\end{proof}

\subsection{Continuity and differentiation in the background}
\label{sec:v2v54-background}

\begin{lemma}[Difference of exponentials]
\label{lem:v2v54-exp-difference}
For bounded operators \(A,B\),
\begin{equation}
 \|e^A-e^B\|
 \le e^{\|A\|+\|B\|}\|A-B\|.
\label{eq:v2v54-exp-difference}
\end{equation}
\end{lemma}

\begin{proof}
Duhamel's identity gives
\[
 e^A-e^B
 =
 \int_0^1e^{(1-s)A}(A-B)e^{sB}\,\mathrm ds.
\]
Take norms and use
\(\|e^{tA}\|\le e^{t\|A\|}\).
\end{proof}

\begin{theorem}[A summable shell modulus gives a continuous chart]
\label{thm:v2v54-chart-modulus}
In addition to \eqref{eq:v2v54-summability}, suppose
\begin{equation}
 \|\delta_{n,g}-\delta_{n,h}\|
 \le\ell_n\|g-h\|_X,
 \qquad
 \sum_{n=1}^\infty\ell_n<\infty.
\label{eq:v2v54-shell-modulus}
\end{equation}
Then
\begin{equation}
 \|U_g-U_h\|
 \le
 e^S
 \left(\sum_{n=1}^\infty\ell_n\right)
 \|g-h\|_X.
\label{eq:v2v54-chart-Lipschitz}
\end{equation}
The inverse chart has the analogous bound with another constant depending
only on \(S\).
\end{theorem}

\begin{proof}
The finite product difference telescopes exactly:
\begin{align*}
 U_{N,g}-U_{N,h}
 =
 \sum_{k=1}^N&
 \alpha_{N,g}\cdots\alpha_{k+1,g}
 (\alpha_{k,g}-\alpha_{k,h})\\
 &\cdot
 \alpha_{k-1,h}\cdots\alpha_{1,h}.
\end{align*}
Put \(s_k=\sup_g\|\delta_{k,g}\|\).  The Duhamel formula in the proof of
Lemma~\ref{lem:v2v54-exp-difference} gives the sharper shellwise estimate
\[
 \|\alpha_{k,g}-\alpha_{k,h}\|
 \le e^{s_k}\ell_k\|g-h\|_X.
\]
The two products outside this difference have combined norm at most
\(e^{S-s_k}\).  Hence the \(k\)-th telescoping term is bounded by
\(e^S\ell_k\|g-h\|_X\).  Sum in \(k\), then pass uniformly to the limit.
Apply the same proof to the reverse products of
\(e^{-\delta_{n,g}}\) for the inverse.
\end{proof}

\begin{remark}[Only finiteness of the constant is used later]
\label{rem:v2v54-constant}
The constant in \eqref{eq:v2v54-chart-Lipschitz} is not used sharply.  The
essential facts are independence from the cutoff and summability of
\((\ell_n)\).  This bound controls the chart-connection term in
\eqref{eq:v2v53-rce}.
\end{remark}

\subsection{BRST descent}
\label{sec:v2v54-BRST}

Let \((\mathfrak A_*,s)\) now be a differential graded Banach
\(*\)-algebra.  Assume \(s\) is a bounded degree-one derivation,
\(s^2=0\), and \(\operatorname{im}s\) is closed in \(\ker s\) at degree
zero.  Then
\[
 H^0_s(\mathfrak A_*)
 =
 \ker(s:\mathfrak A_*^0\to\mathfrak A_*^1)/
 \operatorname{im}(s:\mathfrak A_*^{-1}\to\mathfrak A_*^0)
\]
is a Hausdorff Banach physical space.

\begin{theorem}[The shell product descends to physical cohomology]
\label{thm:v2v54-BRST-descent}
Under Theorem~\ref{thm:v2v54-product}, suppose
\begin{equation}
 [s,\delta_{n,g}]=0
\label{eq:v2v54-chain-derivation}
\end{equation}
for every \(n,g\).  Then every \(\alpha_{n,g}\), \(U_{N,g}\), and \(U_g\)
commutes with \(s\).  Hence \(U_g\) induces a bounded invertible map
\begin{equation}
 [U_g]:H^0_s(\mathfrak A_*)\longrightarrow
 H^0_s(\mathfrak A_*).
\label{eq:v2v54-physical-chart}
\end{equation}
\end{theorem}

\begin{proof}
Equation~\eqref{eq:v2v54-chain-derivation} implies
\(s\delta_{n,g}^k=\delta_{n,g}^ks\) for every \(k\), hence
\(s\alpha_{n,g}=\alpha_{n,g}s\).  Products commute with \(s\), and the
identity passes to the operator-norm limit because \(s\) is bounded.
A chain automorphism maps cycles to cycles and boundaries to boundaries;
its inverse does the same.  It therefore induces the asserted cohomology
automorphism.
\end{proof}

\begin{proposition}[The anomaly obstruction is prior to compactness]
\label{prop:v2v54-anomaly}
Let \(\alpha\) be a shell transport.  If there is a cycle \(A\) with
\begin{equation}
 sA=0,
 \qquad
 s\alpha(A)\ne0,
\label{eq:v2v54-anomaly-witness}
\end{equation}
then the assignment \([A]\mapsto[\alpha(A)]\) is not defined on BRST
cohomology.  Thus \(\alpha\) itself defines no cohomology map; norm
convergence of a product containing \(\alpha\) does not by itself prove
that the product descends.
\end{proposition}

\begin{proof}
A cohomology class must be represented by a cycle.  The second condition
in \eqref{eq:v2v54-anomaly-witness} says that \(\alpha(A)\) is not a cycle,
so it represents no class.  A later compensating correction would have to
be proved algebraically; convergence alone supplies no such correction.
\end{proof}

\begin{firewall}[Local anomaly cancellation must precede the product]
\label{fw:v2v54-anomaly}
The first-volume BV analysis identified local ghost-number-one anomaly
classes and the need for uniformly bounded primitives.  V54 assumes that
the allowed local counterterms have already corrected each shell transport
to satisfy \eqref{eq:v2v54-chain-derivation}.  A formal species
cancellation without a regulator-compatible primitive does not supply a
chain map and cannot be inserted into
Theorem~\ref{thm:v2v54-BRST-descent}.
\end{firewall}

\subsection{Regulator independence on physical cohomology}
\label{sec:v2v54-independence}

Let
\[
 A_{N,g}=a_{N,g}\cdots a_{1,g},
 \qquad
 B_{N,g}=b_{N,g}\cdots b_{1,g}
\]
be two summably generated shell products.  Assume all \(a_{n,g}\) and
\(b_{n,g}\) are chain automorphisms commuting with \(s\).

\begin{theorem}[Summable shell homotopies give one physical chart]
\label{thm:v2v54-homotopy}
Suppose there are bounded degree-minus-one maps \(k_{n,g}\) such that
\begin{equation}
 b_{n,g}-a_{n,g}=sk_{n,g}+k_{n,g}s
\label{eq:v2v54-shell-homotopy}
\end{equation}
and
\begin{equation}
 \sum_{n=1}^\infty\sup_{g\in\mathcal K}\|k_{n,g}\|<\infty.
\label{eq:v2v54-homotopy-sum}
\end{equation}
Assume the partial products and their inverses are uniformly bounded and
converge in operator norm to \(A_g\) and \(B_g\).  Then there is a bounded
degree-minus-one map \(K_g\) such that
\begin{equation}
 B_g-A_g=sK_g+K_gs.
\label{eq:v2v54-limit-homotopy}
\end{equation}
Consequently
\begin{equation}
 [B_g]=[A_g]
 \quad\hbox{on }H^0_s(\mathfrak A_*).
\label{eq:v2v54-same-physical}
\end{equation}
\end{theorem}

\begin{proof}
The noncommutative product difference has the exact telescope
\begin{equation}
 B_{N,g}-A_{N,g}
 =
 \sum_{n=1}^N
 b_{N,g}\cdots b_{n+1,g}
 (b_{n,g}-a_{n,g})
 a_{n-1,g}\cdots a_{1,g}.
\label{eq:v2v54-product-telescope}
\end{equation}
All factors outside \(k_{n,g}\) commute with \(s\).  Therefore
\eqref{eq:v2v54-shell-homotopy} turns
\eqref{eq:v2v54-product-telescope} into
\[
 B_{N,g}-A_{N,g}=sK_{N,g}+K_{N,g}s,
\]
where
\[
 K_{N,g}
 =
 \sum_{n=1}^N
 b_{N,g}\cdots b_{n+1,g}
 k_{n,g}
 a_{n-1,g}\cdots a_{1,g}.
\]
For each fixed \(n\), the left tail converges as \(N\to\infty\), because
\[
 b_{N,g}\cdots b_{n+1,g}=B_{N,g}B_{n,g}^{-1}.
\]
Uniform product bounds and \eqref{eq:v2v54-homotopy-sum} give a summable
majorant independent of \(N,g\).  Dominated convergence for the operator
series yields \(K_{N,g}\to K_g\).  Pass to the limit and use boundedness of
\(s\) to obtain \eqref{eq:v2v54-limit-homotopy}.  Chain-homotopic maps
induce the same map on cohomology.
\end{proof}

\begin{assessment}[The exact meaning of regulator independence]
\label{ass:v2v54-independence}
Two regulators need not give identical gauge-fixed representatives.
Theorem~\ref{thm:v2v54-homotopy} proves the physically relevant statement:
after fixed renormalization conditions, a summable chain homotopy makes
their continuum comparison maps identical on BRST cohomology.  A finite
nontrivial cohomology automorphism left by the comparison must instead be
fixed by a physical normalization condition; summable closeness alone
does not force it to be the identity.
\end{assessment}

\subsection{Why raw marginal shells cannot be inserted}
\label{sec:v2v54-marginal}

\begin{proposition}[Vanishing shell size is not enough]
\label{prop:v2v54-harmonic}
Let \(H=H^*\in M_2(\mathbb C)\) have two distinct eigenvalues and define
\[
 \delta_n(A)=\frac{i}{n}[H,A].
\]
Then \(\|\delta_n\|\to0\), but the ordered products
\[
 U_N
 =
 \prod_{n=1}^N e^{\delta_n}
 =
 \operatorname{Ad}\!\left(
  e^{\,iH\sum_{n=1}^N1/n}
 \right)
\]
need not converge.  Thus shellwise decay without a summable tail does not
construct a chart.
\end{proposition}

\begin{proof}
The derivations commute, giving the displayed product.  On an off-diagonal
matrix unit between eigenvalues with gap \(\lambda\ne0\), \(U_N\) acts by
the phase
\[
 \exp\!\left(i\lambda\sum_{n=1}^N\frac1n\right).
\]
Write \(h_N=\sum_{n=1}^Nn^{-1}=\log N+\gamma+o(1)\).  For any angle
\(\theta\), choose integers
\[
 N_k=\left\lfloor
 \exp\!\left(\frac{2\pi k+\theta}{|\lambda|}-\gamma\right)
 \right\rfloor .
\]
Then \(e^{i\lambda h_{N_k}}\) tends to the phase with angle
\(\operatorname{sgn}(\lambda)\theta\).  Two distinct choices of
\(\theta\) give distinct subsequential limits.  Therefore \(U_N\) does
not converge on that matrix unit.
\end{proof}

\begin{firewall}[Local running terms must be extracted first]
\label{fw:v2v54-local-running}
Marginal local shells commonly have size \(1/n\) or a constant size in
logarithmic scale.  Proposition~\ref{prop:v2v54-harmonic} shows why they
cannot be hidden in \(\delta_{n,g}\).  They must first be accumulated in
the finite-dimensional local coupling transport fixed by renormalization
conditions.  The summability hypothesis
\eqref{eq:v2v54-summability} applies only to the Ward-compatible
nonlocal or irrelevant remainder after that extraction.
\end{firewall}

\subsection{V54 conclusion and next acquisition}
\label{sec:v2v54-conclusion}

V54 constructs the common algebra chart required in V52 from a uniformly
summable product of renormalized shell derivations.  A summable metric
modulus makes the chart Lipschitz and controls the chart term in V53.
Commutation with the BRST differential makes the limit descend to the
physical algebra, while a summable chain homotopy proves regulator
independence on physical cohomology.  The harmonic-series example shows
that merely vanishing shell changes do not suffice.

The theorem leaves two analytic tasks explicit: extract the full
dimension-four local running block with fixed physical normalization, and
prove absolute summability of the complete Ward-compatible remainder in a
Banach algebra norm containing the stress insertions.  V55 is the next
mandatory companion audit.  It will compare the newest NS and YM advances
with this shell-product criterion and choose whether the next estimate
should target absolute summability, a weaker conditional convergence
mechanism, or the local anomaly primitive.
\clearpage
\section*{Second Volume, Version V55: Companion Audit and Ordered Abel-Tail Convergence}
\addcontentsline{toc}{section}{Second Volume, Version V55: Companion Audit and Ordered Abel-Tail Convergence}

\subsection*{Purpose and mandatory audit}

V54 gave an absolutely summable shell-derivation criterion for the common
physical algebra chart.  Absolute summability is robust but may be too
strong at a marginal or oscillatory ultraviolet boundary.  This mandatory
companion audit indicates the correct weakening: preserve the complete
ordered carrier, telescope its signed first-order tail, and take a norm
only on the resulting tail--increment product.

The Navier--Stokes file inspected was
\[
 \texttt{Renewal\_NS\_Stability\_Research\_Notes\_V31.tex},
\]
with \(887537\) bytes and SHA--256 digest
\[
 \texttt{F1121B62238AD5491BB7CF03635EA9934942EDF573ED8FAAA9EE79E1D38A6696}.
\]
Its last result remains the V31 formation analysis: a weak heat source
does not generate the missing critical first moment for free.  The sharper
remaining object is a stopping-time-correlated occupation, not the product
of separate upper bounds.  This supports looking for an ordered
tail--increment stock rather than demanding the sum of all shell norms.

The Yang--Mills files named V48 and V49 were byte-identical at the audit
time.  Each had \(643955\) bytes and SHA--256 digest
\[
 \texttt{91720B37083B85FDB60603F7F078D2F7C1E9D229250F673A920DFC839A05E162}.
\]
Their last mathematical chapter is V48.  V47 constructs one exact
connected-graph identity on the fifteen-state partition carrier.  V48
then proves that an incidence-three invariant supplies only one capacity
power, compresses every consecutive repeated-coordinate run, and reduces
all remaining protected overlaps to one incidence-four crossed repeat.
The transferable rule is that two cancellations using the same local
factor cannot be multiplied.  One must first assemble a single global
operator identity and only then estimate it.

\begin{assessment}[V55 audit decision]
\label{ass:v2v55-audit}
For the SM algebra chart, the shell increments are noncommuting operators
on one complete physical carrier.  Sectorwise conditional convergence is
therefore insufficient, just as separate YM cut capacities are
insufficient.  The next legitimate weakening of V54 is an exact Abel
identity for the complete ordered product, controlled by one tail in the
same operator order.  Absolute summability remains a sufficient special
case.
\end{assessment}

\subsection{An exact Abel identity for ordered products}
\label{sec:v2v55-Abel}

Let \(\mathcal B\) be a unital Banach algebra and let
\(D_n\in\mathcal B\).  Set
\begin{equation}
 P_0=I,
 \qquad
 P_n=(I+D_n)(I+D_{n-1})\cdots(I+D_1).
\label{eq:v2v55-product}
\end{equation}
Suppose the series \(\sum_{n\ge1}D_n\) converges in the Banach-algebra
norm, and define its ordered additive tail
\begin{equation}
 R_n:=\sum_{k=n+1}^\infty D_k.
\label{eq:v2v55-tail}
\end{equation}

\begin{lemma}[Exact ordered Abel identity]
\label{lem:v2v55-Abel}
For integers \(n>m\ge0\),
\begin{equation}
\boxed{
 P_n-P_m
 =
 R_mP_m-R_nP_{n-1}
 +\sum_{k=m+1}^{n-1}R_kD_kP_{k-1}.}
\label{eq:v2v55-Abel}
\end{equation}
No commutativity is assumed.
\end{lemma}

\begin{proof}
The product recurrence gives
\[
 P_k-P_{k-1}=D_kP_{k-1},
\]
hence
\[
 P_n-P_m=\sum_{k=m+1}^nD_kP_{k-1}.
\]
Since \(D_k=R_{k-1}-R_k\),
\begin{align*}
 \sum_{k=m+1}^nD_kP_{k-1}
 &=
 \sum_{k=m+1}^n
 (R_{k-1}-R_k)P_{k-1}\\
 &=
 R_mP_m-R_nP_{n-1}
 +\sum_{k=m+1}^{n-1}R_k(P_k-P_{k-1}).
\end{align*}
Insert \(P_k-P_{k-1}=D_kP_{k-1}\).
\end{proof}

\begin{theorem}[Conditional ordered-product convergence]
\label{thm:v2v55-product}
Assume
\begin{align}
 \sum_{n=1}^\infty D_n
 &\quad\hbox{converges in }\mathcal B,
\label{eq:v2v55-additive-convergence}\\
 \sum_{n=1}^\infty\|R_nD_n\|
 &<\infty.
\label{eq:v2v55-tail-increment}
\end{align}
Then \(P_n\) converges in the Banach-algebra norm.
\end{theorem}

\begin{proof}
First prove uniform boundedness.  Choose \(m_0\) so large that
\[
 \varepsilon:=
 \sup_{k\ge m_0}\|R_k\|
 <\frac14,
 \qquad
 A:=
 \sum_{k=m_0+1}^\infty\|R_kD_k\|
 <\frac14.
\]
For \(n>m_0\), Lemma~\ref{lem:v2v55-Abel} gives
\[
 \|P_n\|
 \le
 (1+\varepsilon)\|P_{m_0}\|
 +(\varepsilon+A)
 \max_{m_0\le j\le n}\|P_j\|.
\]
Taking the maximum over \(n\) and using
\(\varepsilon+A<1\) yields
\begin{equation}
 \sup_{n\ge0}\|P_n\|=:M<\infty.
\label{eq:v2v55-product-bound}
\end{equation}

Apply the Abel identity again.  For \(n>m\),
\[
 \|P_n-P_m\|
 \le
 M\left(
  \|R_m\|+\|R_n\|
  +\sum_{k=m+1}^{n-1}\|R_kD_k\|
 \right).
\]
The right side tends to zero uniformly in \(n>m\), by
\eqref{eq:v2v55-additive-convergence} and
\eqref{eq:v2v55-tail-increment}.  Thus \((P_n)\) is Cauchy.
\end{proof}

\begin{assessment}[The new stock is ordered and correlated]
\label{ass:v2v55-stock}
The quantity in \eqref{eq:v2v55-tail-increment} is not
\[
 \sum_n\|R_n\|\,\|D_n\|.
\]
It is the norm of the complete ordered product \(R_nD_n\).  Cancellation
inside that product is therefore retained until the correct stage.  This
is the algebra-chart analogue of the NS restricted formation occupation
and the YM single global graph capacity.
\end{assessment}

\subsection{Uniform families and the V52 chart}
\label{sec:v2v55-uniform}

Let \(D_{n,g}\in\mathcal B\) depend on \(g\in\mathcal K\), define
\(R_{n,g}\) and \(P_{n,g}\) as above, and assume every
\(g\mapsto D_{n,g}\) is continuous.

\begin{theorem}[Uniform Abel-tail criterion]
\label{thm:v2v55-uniform}
Suppose the additive series converges uniformly in \(g\), and
\begin{align}
 \sup_{g\in\mathcal K}\sup_{k\ge m}\|R_{k,g}\|
 &\longrightarrow0,
\label{eq:v2v55-uniform-tail}\\
 \sum_{n=1}^\infty
 \sup_{g\in\mathcal K}\|R_{n,g}D_{n,g}\|
 &<\infty.
\label{eq:v2v55-uniform-tail-increment}
\end{align}
Then \(P_{n,g}\) converges in norm, uniformly in \(g\), to a continuous
map \(g\mapsto P_g\).
\end{theorem}

\begin{proof}
The proof of Theorem~\ref{thm:v2v55-product} is uniform after replacing
each tail and tail--increment norm by the corresponding supremum over
\(g\).  It yields one bound \(M\) and a uniform Cauchy estimate.
Every finite product \(P_{n,g}\) is continuous in \(g\), so its uniform
limit is continuous.
\end{proof}

\begin{corollary}[V54 absolute summability is a special case]
\label{cor:v2v55-absolute}
If
\[
 \sum_n\sup_g\|D_{n,g}\|<\infty,
\]
then the hypotheses of Theorem~\ref{thm:v2v55-uniform} hold.
\end{corollary}

\begin{proof}
The additive series converges absolutely and uniformly, while
\[
 \sum_n\sup_g\|R_{n,g}D_{n,g}\|
 \le
 \sum_n
 \left(\sum_{k>n}\sup_g\|D_{k,g}\|\right)
 \sup_g\|D_{n,g}\|
 <\infty.
\]
The last double sum is bounded by one half of the square of the
\(\ell^1\) norm, up to the diagonal convention.
\end{proof}

\subsection{Invertibility requires the inverse order}
\label{sec:v2v55-inverse}

Convergence of invertible factors does not alone imply that their product
limit is invertible.  Let
\[
 E_n:=(I+D_n)^{-1}-I,
 \qquad
 S_n:=\sum_{k=n+1}^\infty E_k.
\]
The inverse partial product is
\begin{equation}
 P_n^{-1}
 =
 (I+E_1)(I+E_2)\cdots(I+E_n),
\label{eq:v2v55-inverse-product}
\end{equation}
whose order is reversed.

\begin{theorem}[Two-sided Abel criterion for an algebra chart]
\label{thm:v2v55-invertible}
Assume the hypotheses of Theorem~\ref{thm:v2v55-product}, every
\(I+D_n\) is invertible, and
\begin{align}
 \sum_{n=1}^\infty E_n
 &\quad\hbox{converges in }\mathcal B,
\label{eq:v2v55-inverse-additive}\\
 \sum_{n=1}^\infty\|E_nS_n\|
 &<\infty.
\label{eq:v2v55-inverse-tail}
\end{align}
Then \(P_n\to P\), \(P_n^{-1}\to Q\), and
\[
 PQ=QP=I.
\]
Hence \(P\) is invertible.  The same conclusion holds uniformly over
\(\mathcal K\) under the uniform versions of all four tail conditions.
\end{theorem}

\begin{proof}
Apply Theorem~\ref{thm:v2v55-product} to the sequence \(E_n\) in the
opposite Banach algebra, whose multiplication is
\(A\circ B=BA\).  Its left ordered product is precisely
\eqref{eq:v2v55-inverse-product}, and its tail--increment product in the
opposite algebra is \(E_nS_n\).  Thus \(P_n^{-1}\to Q\).
Passing to the limit in
\(P_nP_n^{-1}=P_n^{-1}P_n=I\) gives the result.
\end{proof}

\begin{firewall}[One-sided convergence is not an automorphism theorem]
\label{fw:v2v55-one-sided}
Theorem~\ref{thm:v2v55-product} gives a bounded endomorphism limit when the
factors are algebra endomorphisms, but it does not prove that the limit is
onto.  The inverse-tail conditions in
Theorem~\ref{thm:v2v55-invertible}, or a separate uniform inverse bound and
surjectivity argument, are required before the limit may be used as the
V52 algebra chart.
\end{firewall}

\subsection{Application to shell automorphisms}
\label{sec:v2v55-automorphisms}

Return to the V54 shell automorphisms
\[
 \alpha_{n,g}=e^{\delta_{n,g}},
 \qquad
 D_{n,g}:=\alpha_{n,g}-I.
\]

\begin{corollary}[Conditional physical chart]
\label{cor:v2v55-physical-chart}
Suppose the uniform forward and inverse Abel criteria hold for
\[
 D_{n,g}=\alpha_{n,g}-I,
 \qquad
 E_{n,g}=\alpha_{n,g}^{-1}-I.
\]
Then the ordered products converge uniformly to invertible unital
\(*\)-automorphisms \(U_g\), continuously in \(g\).  If every
\(\alpha_{n,g}\) commutes with the BRST differential, \(U_g\) descends to
the physical cohomology.  If two such regulator products obey the
summable chain-homotopy hypothesis of V54, they induce the same physical
chart.
\end{corollary}

\begin{proof}
Theorems~\ref{thm:v2v55-uniform} and
\ref{thm:v2v55-invertible} give uniform convergence and invertibility.
Multiplicativity, the unit, and the involution pass through the norm limit.
BRST descent and regulator independence are
Theorems~\ref{thm:v2v54-BRST-descent} and
\ref{thm:v2v54-homotopy}.
\end{proof}

\begin{firewall}[The Abel tail must be formed after Ward assembly]
\label{fw:v2v55-Ward-assembly}
The increments \(D_{n,g}\) in
\eqref{eq:v2v55-tail-increment} are the complete physical shell
automorphisms after gauge, ghost, matter, contact, and anomaly-counterterm
pieces have been combined.  Applying the criterion to separate signed
sectors and multiplying the resulting estimates would repeat the
multi-cut error excluded by YM V48.  The BRST chain-map property must also
hold before the tail is formed.
\end{firewall}

\subsection{A genuinely conditional example}
\label{sec:v2v55-example}

\begin{proposition}[The Abel criterion is strictly weaker than absolute summability]
\label{prop:v2v55-conditional-example}
Let \(H=H^*\in M_2(\mathbb C)\), let
\(2/3<a\le1\), and set
\[
 \delta_n(A)
 =
 i(-1)^nn^{-a}[H,A],
 \qquad
 \alpha_n=e^{\delta_n}.
\]
Then
\[
 \sum_n\|\delta_n\|=\infty
\]
when \(H\) is not scalar, but the products
\(\alpha_N\cdots\alpha_1\) converge to an inner
\(*\)-automorphism.  Their increments satisfy the forward and inverse
Abel criteria.
\end{proposition}

\begin{proof}
All derivations commute, and the alternating scalar series converges.
Therefore
\[
 \alpha_N\cdots\alpha_1
 =
 \exp\!\left(
  i\sum_{n=1}^N(-1)^nn^{-a}\operatorname{ad}_H
 \right)
\]
converges in operator norm.  Non-scalarity of \(H\) gives
\(\|\delta_n\|=c_Hn^{-a}\), whose sum diverges for \(a\le1\).

For the explicit Abel hypotheses, Taylor expansion gives
\[
 D_n=\delta_n+Q_n,
 \qquad
 \|Q_n\|\le C_Hn^{-2a}.
\]
The \(Q_n\) series is absolutely summable because \(2a>1\).  The
alternating-series estimate gives
\[
 \left\|\sum_{k>n}\delta_k\right\|\le C_Hn^{-a}.
\]
The Taylor remainder tail is bounded by
\[
 \sum_{k>n}\|Q_k\|
 \le C_H n^{1-2a}.
\]
Consequently
\[
 \|R_n\|
 \le C_H\bigl(n^{-a}+n^{1-2a}\bigr),
 \qquad
 \|R_nD_n\|
 \le C_H\bigl(n^{-2a}+n^{1-3a}\bigr).
\]
The last sequence is summable exactly under the stated restriction
\(a>2/3\).  Replacing \(\delta_n\) by \(-\delta_n\) gives the same
estimates for the inverse increments.
\end{proof}

\begin{assessment}[What the example proves and does not prove]
\label{ass:v2v55-example}
The example proves that V54 absolute summability is not logically
necessary.  Its cancellation is unusually transparent because the
derivations commute.  The interacting SM shell maps need not commute, so
their convergence must be established through the complete ordered
quantity \(R_nD_n\), not by citing the scalar alternating-series test.
\end{assessment}

\subsection{V55 conclusion and next acquisition}
\label{sec:v2v55-conclusion}

V55 proves an exact noncommutative Abel identity and a conditional
ordered-product convergence theorem.  The criterion allows a convergent
signed additive shell series whose norm series diverges, provided the
single correlated tail--increment stock is summable.  A separate
opposite-algebra criterion makes the limit invertible.  Uniform versions
produce a continuous background chart, and the V54 BRST and chain-homotopy
theorems then give a regulator-independent physical chart.

The next analytic target is more precise than absolute summability.
V56 will express the complete renormalized SM shell increment as an
adjacent cutoff difference plus a local running block and a remainder.  It
will derive a resolvent-based bound for the ordered Abel stock
\(\sum_n\|R_nD_n\|\), keeping the Ward-completed species sum intact.  If
the marginal local block does not cancel, it will be returned to the
finite-dimensional coupling transport rather than charged to the
nonlocal product.
\clearpage
\section*{Second Volume, Version V56: Cumulative Generators and a Two-Sided Abel-Stock Compiler}
\addcontentsline{toc}{section}{Second Volume, Version V56: Cumulative Generators and a Two-Sided Abel-Stock Compiler}

\subsection*{Purpose}

V55 replaced absolute shell summability by forward and inverse ordered
Abel-tail criteria.  This note derives those criteria from a more natural
Wilsonian object: a convergent cumulative renormalized generator.  The
linear increments telescope, while exponentiating them creates a
quadratic remainder.  Square summability pays that remainder.  The only
noncommutative first-order quantities left are the two ordered products of
the cumulative tail with the current increment.

This is a compiler theorem.  It identifies a sufficient stock for the
complete Ward-assembled SM shell transport, but does not claim that the
interacting shell maps already obey it.  The V9 resolvent theorem verifies
a stronger special case for one commuting transition-supported regression;
its unresolved off-diagonal and running-insertion terms remain unresolved.

\subsection{Cumulative derivations and shell increments}
\label{sec:v2v56-cumulative}

Let \(\mathfrak A_*\) be the V54 Banach \(*\)-algebra and let
\[
 \mathcal B=\mathcal L(\mathfrak A_*)
\]
with composition as multiplication.  Let \(K_n\in\mathcal B\) be bounded
\(*\)-derivations and suppose
\begin{equation}
 K_n\longrightarrow K_\infty
 \quad\hbox{in }\mathcal B.
\label{eq:v2v56-generator-limit}
\end{equation}
Put
\begin{align}
 q_n&:=K_n-K_{n-1},
\label{eq:v2v56-q}\\
 h_n&:=K_\infty-K_n
     =\sum_{k=n+1}^\infty q_k,
\label{eq:v2v56-h}\\
 d_n&:=\|q_n\|,
\qquad
 \eta_n:=\sum_{k=n+1}^\infty d_k^2.
\label{eq:v2v56-eta}
\end{align}
The shell automorphism and its forward increment are
\begin{equation}
 \alpha_n=e^{q_n},
 \qquad
 D_n=\alpha_n-I.
\label{eq:v2v56-D}
\end{equation}
Because \(q_n\) is a bounded \(*\)-derivation,
Lemma~\ref{lem:v2v54-exponential} makes \(\alpha_n\) a
\(*\)-automorphism.

\begin{lemma}[Quadratic exponential remainder]
\label{lem:v2v56-exp-remainder}
If \(\sup_nd_n\le b<\infty\), then
\begin{align}
 D_n&=q_n+N_n,
\label{eq:v2v56-forward-remainder}\\
 e^{-q_n}-I&=-q_n+\widetilde N_n,
\label{eq:v2v56-inverse-remainder}\\
 \|N_n\|+\|\widetilde N_n\|
 &\le C_b d_n^2,
\label{eq:v2v56-remainder-bound}
\end{align}
where \(C_b=e^b\) is sufficient.
\end{lemma}

\begin{proof}
The exponential series gives
\[
 e^{q_n}-I-q_n
 =
 \sum_{r=2}^\infty\frac{q_n^r}{r!},
\]
whose norm is at most
\[
 d_n^2\sum_{r=2}^\infty\frac{b^{r-2}}{r!}
 \le e^b d_n^2.
\]
Replace \(q_n\) by \(-q_n\) for the inverse.
\end{proof}

\subsection{The forward ordered stock}
\label{sec:v2v56-forward}

Assume
\begin{equation}
 \sum_{n=1}^\infty d_n^2<\infty.
\label{eq:v2v56-square}
\end{equation}
Then \(\sum N_n\) and \(\sum\widetilde N_n\) converge absolutely.  The
additive series of forward increments satisfies
\[
 \sum_{n=1}^ND_n
 =
 K_N-K_0+\sum_{n=1}^NN_n,
\]
so it converges.  Its tail is
\begin{equation}
 R_n:=\sum_{k=n+1}^\infty D_k
 =
 h_n+\mathcal N_n,
\qquad
 \mathcal N_n:=\sum_{k=n+1}^\infty N_k,
\label{eq:v2v56-forward-tail}
\end{equation}
with
\begin{equation}
 \|\mathcal N_n\|\le C_b\eta_n.
\label{eq:v2v56-Ntail}
\end{equation}

\begin{lemma}[Forward Abel-stock reduction]
\label{lem:v2v56-forward-stock}
Under \eqref{eq:v2v56-generator-limit} and
\eqref{eq:v2v56-square},
\begin{align}
 \|R_nD_n\|
 \le{}&
 \|h_nq_n\|
 +C_b\|h_n\|d_n^2
 +C_b\eta_nd_n
 +C_b^2\eta_nd_n^2.
\label{eq:v2v56-forward-stock}
\end{align}
\end{lemma}

\begin{proof}
Use
\[
 R_nD_n
 =
 (h_n+\mathcal N_n)(q_n+N_n)
\]
and expand the four terms in their displayed operator order.  Apply
\eqref{eq:v2v56-remainder-bound} and
\eqref{eq:v2v56-Ntail}.
\end{proof}

\subsection{The inverse ordered stock}
\label{sec:v2v56-inverse}

Set
\[
 E_n:=\alpha_n^{-1}-I=e^{-q_n}-I,
 \qquad
 S_n:=\sum_{k=n+1}^\infty E_k.
\]
By the preceding argument,
\begin{equation}
 S_n=-h_n+\widetilde{\mathcal N}_n,
\qquad
 \|\widetilde{\mathcal N}_n\|\le C_b\eta_n.
\label{eq:v2v56-inverse-tail}
\end{equation}

\begin{lemma}[Inverse Abel-stock reduction]
\label{lem:v2v56-inverse-stock}
Under the same hypotheses,
\begin{align}
 \|E_nS_n\|
 \le{}&
 \|q_nh_n\|
 +C_b d_n\eta_n
 +C_b\|h_n\|d_n^2
 +C_b^2d_n^2\eta_n.
\label{eq:v2v56-inverse-stock}
\end{align}
\end{lemma}

\begin{proof}
Expand, in the inverse order required by V55,
\[
 E_nS_n
 =
 (-q_n+\widetilde N_n)
 (-h_n+\widetilde{\mathcal N}_n).
\]
Bound the four terms by
\eqref{eq:v2v56-remainder-bound} and
\eqref{eq:v2v56-inverse-tail}.
\end{proof}

\subsection{The two-sided compiler}
\label{sec:v2v56-compiler}

\begin{theorem}[Cumulative-generator criterion for an invertible chart]
\label{thm:v2v56-compiler}
Suppose \eqref{eq:v2v56-generator-limit} and
\eqref{eq:v2v56-square} hold, and assume
\begin{align}
 \sum_{n=1}^\infty
 \bigl(\|h_nq_n\|+\|q_nh_n\|\bigr)
 &<\infty,
\label{eq:v2v56-ordered-correlations}\\
 \sum_{n=1}^\infty\eta_nd_n
 &<\infty.
\label{eq:v2v56-quadratic-tail}
\end{align}
Then the ordered shell products
\begin{equation}
 U_N=e^{q_N}e^{q_{N-1}}\cdots e^{q_1}
\label{eq:v2v56-shell-product}
\end{equation}
converge in operator norm to an invertible unital
\(*\)-automorphism \(U\), and \(U_N^{-1}\) converges to \(U^{-1}\).
\end{theorem}

\begin{proof}
Convergence of \(K_n\) makes \((h_n)\) bounded.  Hence
\[
 \sum_n\|h_n\|d_n^2<\infty
\]
by \eqref{eq:v2v56-square}.  Also \((\eta_n)\) is bounded, so
\[
 \sum_n\eta_nd_n^2<\infty.
\]
Lemma~\ref{lem:v2v56-forward-stock},
\eqref{eq:v2v56-ordered-correlations}, and
\eqref{eq:v2v56-quadratic-tail} give
\[
 \sum_n\|R_nD_n\|<\infty.
\]
The additive forward series converges by
\eqref{eq:v2v56-forward-tail}.  Therefore
Theorem~\ref{thm:v2v55-product} gives \(U_N\to U\).

For the inverse increments, the additive series converges by
\eqref{eq:v2v56-inverse-remainder}, and
Lemma~\ref{lem:v2v56-inverse-stock} gives
\[
 \sum_n\|E_nS_n\|<\infty.
\]
The opposite-algebra theorem
\ref{thm:v2v55-invertible} gives convergence of \(U_N^{-1}\) and
invertibility of \(U\).  Multiplicativity, the unit, and the involution
pass through the norm limits as in V54.
\end{proof}

\begin{assessment}[Only two first-order operator products remain]
\label{ass:v2v56-two-products}
All exponential corrections are paid by the scalar stocks
\eqref{eq:v2v56-square} and
\eqref{eq:v2v56-quadratic-tail}.  Noncommutativity survives only in
\[
 h_nq_n
 \quad\hbox{and}\quad
 q_nh_n.
\]
They cannot be replaced by one another and cannot be bounded by multiplying
independent sector estimates.  Both must be formed on the complete
Ward-assembled physical carrier.
\end{assessment}

\begin{corollary}[Uniform background family]
\label{cor:v2v56-uniform}
Let \(K_{n,g}\to K_{\infty,g}\) uniformly over
\(g\in\mathcal K\), with every finite-stage map continuous in \(g\).
If the sums in
\eqref{eq:v2v56-square},
\eqref{eq:v2v56-ordered-correlations}, and
\eqref{eq:v2v56-quadratic-tail} remain finite after placing
\(\sup_{g\in\mathcal K}\) inside every summand, then
\(U_{N,g}\) and \(U_{N,g}^{-1}\) converge uniformly to continuous inverse
charts \(U_g\) and \(U_g^{-1}\).
\end{corollary}

\begin{proof}
The lemmas give the uniform forward and inverse Abel criteria of V55.
Apply Theorems~\ref{thm:v2v55-uniform} and
\ref{thm:v2v55-invertible}.  Uniform limits of continuous maps are
continuous.
\end{proof}

\begin{corollary}[BRST descent remains shellwise]
\label{cor:v2v56-BRST}
If
\[
 [s,K_n]=[s,K_\infty]=0
\]
for all \(n\), then every \(q_n\) commutes with \(s\), and the limiting
chart in Theorem~\ref{thm:v2v56-compiler} descends to
\(H_s^0(\mathfrak A_*)\).
\end{corollary}

\begin{proof}
Subtract the two adjacent commutator identities to obtain
\([s,q_n]=0\).  Apply Theorem~\ref{thm:v2v54-BRST-descent}.
\end{proof}

\subsection{A checkable power-law corollary}
\label{sec:v2v56-power}

\begin{corollary}[Logarithmic-shell power threshold]
\label{cor:v2v56-power}
Assume \(K_n\to K_\infty\) and suppose for some \(a>2/3\),
\begin{equation}
 d_n\le Cn^{-a},
 \qquad
 \|h_nq_n\|+\|q_nh_n\|\le Cn^{-2a}.
\label{eq:v2v56-power-assumptions}
\end{equation}
Then all hypotheses of
Theorem~\ref{thm:v2v56-compiler} hold.
\end{corollary}

\begin{proof}
Since \(2a>1\), \(\sum d_n^2<\infty\), and
\[
 \eta_n\le C\sum_{k>n}k^{-2a}
 \le C'n^{1-2a}.
\]
Therefore
\[
 \sum_n\eta_nd_n
 \le C''\sum_n n^{1-3a}<\infty
\]
exactly when \(a>2/3\).  The ordered correlations are summable because
\(2a>1\).
\end{proof}

\begin{firewall}[Norm products are a fallback, not the definition]
\label{fw:v2v56-norm-fallback}
The second inequality in
\eqref{eq:v2v56-power-assumptions} is a bound on already assembled ordered
products.  The stronger estimates
\(\|h_n\|\|q_n\|\le Cn^{-2a}\) would suffice, but are not required by the
compiler.  Using separate gauge, ghost, fermion, Higgs, and contact norms
would discard the Ward cancellation that the V55 audit requires us to
retain.
\end{firewall}

\subsection{Interface with the V9 resolvent regression}
\label{sec:v2v56-V9}

Use one fixed BRST-compatible local projector
\(P_{\rm loc}\) and let \(\mathcal J_n\) denote a cumulative
Ward-completed shell generator in a chosen Banach operator topology.  Set
\begin{equation}
 K_n=(I-P_{\rm loc})\mathcal J_n,
 \qquad
 q_n=K_n-K_{n-1}.
\label{eq:v2v56-projected-generator}
\end{equation}
Because the same projector is used at every shell,
\[
 q_n
 =
 (I-P_{\rm loc})(\mathcal J_n-\mathcal J_{n-1})
\]
without an extra projector-variation term.

\begin{proposition}[What the commuting V9 estimate already supplies]
\label{prop:v2v56-V9}
In the finite-dimensional commuting transition-supported regression of
V9, suppose the projected adjacent generator is the row
\(Q_rJ_n\) from \eqref{eq:v2v9-J-remainder}.  Put
\begin{align}
 \varepsilon_1
 &=r+1-(d+\mu-2),
\label{eq:v2v56-epsilon1}\\
 \varepsilon_2
 &=r+1-(d+\rho+\mu-3).
\label{eq:v2v56-epsilon2}
\end{align}
If \(\varepsilon_1,\varepsilon_2>0\) and
\(\Lambda_n\simeq2^n\), then
\begin{equation}
 \|q_n\|
 \le C\bigl(\Lambda_n^{-\varepsilon_1}
             +\Lambda_n^{-\varepsilon_2}\bigr),
\label{eq:v2v56-V9-bound}
\end{equation}
so \(\sum_n\|q_n\|<\infty\).  This regression therefore satisfies the
stronger V54 criterion and hence the V56 compiler.
\end{proposition}

\begin{proof}
Equation~\eqref{eq:v2v9-J-remainder} gives
\eqref{eq:v2v56-V9-bound} on every fixed external compact set.  Positive
powers of a dyadic inverse scale form geometric series.  Thus the
increments are absolutely summable, \(K_n\) converges, and all V56 stocks
are finite.
\end{proof}

\begin{firewall}[This does not close the complete SM shell]
\label{fw:v2v56-V9-scope}
Proposition~\ref{prop:v2v56-V9} uses the V9 commuting spectral regression
and treats only its first adjacent covariance row.  The complete
interacting update still requires:
\begin{enumerate}
 \item off-diagonal localization of
       \(\Delta F_nC_nF_n\) and \(F_{n-1}C_n\Delta F_n\);
 \item a quasilocal bound for
       \(F_{n-1}C_n\Delta V_nC_{n-1}F_{n-1}\);
 \item the running insertion difference left open in V9--V10;
 \item assembly of all species, ghosts, contact terms, and the anomaly
       primitive before the ordered products are estimated.
\end{enumerate}
None of these follows from
\eqref{eq:v2v56-V9-bound}.
\end{firewall}

\subsection{The local running block}
\label{sec:v2v56-local}

The dispersion identity of V41 and the shell Taylor projector separate
the cumulative generator as
\begin{equation}
 \mathcal J_n
 =
 P_{\rm loc}\mathcal J_n
 +(I-P_{\rm loc})\mathcal J_n.
\label{eq:v2v56-local-split}
\end{equation}
The first term contains the cosmological, Einstein, gauge, Yukawa, Higgs,
and allowed curvature counterterm coordinates.  Its marginal increments
need not tend to zero in an absolutely summable norm.

\begin{firewall}[A marginal beta function is not an Abel remainder]
\label{fw:v2v56-beta}
If \(P_{\rm loc}\mathcal J_n-P_{\rm loc}\mathcal J_{n-1}\) has an
unsummed \(n^{-1}\) component, it belongs to the finite-dimensional running
coupling equation.  Moving it into \(q_n\) can violate
\eqref{eq:v2v56-generator-limit} or the ordered stock conditions.  It may
be removed from the physical remainder only by fixed renormalization
conditions and a BRST-compatible local projector, not by declaring it
small shellwise.
\end{firewall}

\subsection{V56 conclusion and next acquisition}
\label{sec:v2v56-conclusion}

V56 reduces the forward and inverse chart convergence to a convergent
cumulative generator, square-summable increments, one scalar quadratic
tail, and the two ordered correlations \(h_nq_n\) and \(q_nh_n\).
This strictly extends the V54 absolute criterion while preserving the
complete operator order demanded by V55.  The existing V9 theorem verifies
a stronger special case but not the complete interacting SM shell.

The next upstream task is now forced by
Firewall~\ref{fw:v2v56-V9-scope}.  V57 will attack the first noncommuting
resolvent term by a block-resolvent localization theorem.  It will identify
a quantitative off-diagonal mixing norm under which a transition-shell
inverse retains its ultraviolet factor, and it will give a finite-matrix
counterexample showing why global coercivity alone cannot provide that
factor.
\clearpage
\section*{Second Volume, Version V57: Schur-Complement Localization of the Noncommuting Shell Resolvent}
\addcontentsline{toc}{section}{Second Volume, Version V57: Schur-Complement Localization of the Noncommuting Shell Resolvent}

\subsection*{Purpose}

V9 left the first two noncommuting adjacent-resolvent terms open because a
global inverse bound can carry a transition-shell vector into low modes and
destroy its inverse-frequency gain.  V56 identified this as the first
analytic input needed for the complete Abel stock.  This note gives the
exact block criterion.

The result is a Feshbach--Schur estimate.  The inverse restricted to the
transition shell retains a factor \(\Lambda^{-1}\) when the shell block is
coercive, the complementary block is invertible, and the round trip from
the shell through the complement has norm strictly below one.  No
commutation with the reference spectral operator is required.  A
two-dimensional example proves that global coercivity alone cannot replace
the round-trip condition.

\subsection{Block notation}
\label{sec:v2v57-blocks}

Let \(\mathcal H\) be a Hilbert space, let \(P=P^*=P^2\), and put
\(Q=I-P\).  For a bounded self-adjoint invertible operator \(A\), write
\begin{align}
 A_{PP}&=PAP|_{P\mathcal H},
&
 A_{PQ}&=PAQ|_{Q\mathcal H},
\label{eq:v2v57-blocks1}\\
 A_{QP}&=QAP|_{P\mathcal H},
&
 A_{QQ}&=QAQ|_{Q\mathcal H}.
\label{eq:v2v57-blocks2}
\end{align}
The finite-cutoff setting of V9 is covered directly.  The same formulas
hold for closed operators on a common form domain once the displayed block
maps and inverses are bounded in the stated graph norms.

\begin{lemma}[Exact shell-column formula]
\label{lem:v2v57-column}
Assume \(A_{QQ}\) is invertible and define the Schur complement
\begin{equation}
 \mathscr S_P
 =
 A_{PP}-A_{PQ}A_{QQ}^{-1}A_{QP}.
\label{eq:v2v57-Schur}
\end{equation}
If \(\mathscr S_P\) is invertible, then
\begin{align}
 PA^{-1}P&=\mathscr S_P^{-1},
\label{eq:v2v57-PPP}\\
 QA^{-1}P&=-A_{QQ}^{-1}A_{QP}\mathscr S_P^{-1}.
\label{eq:v2v57-QPP}
\end{align}
Consequently, if
\begin{equation}
 M_P:=\|A_{QQ}^{-1}A_{QP}\|,
\label{eq:v2v57-M}
\end{equation}
then
\begin{equation}
 \|A^{-1}P\|
 \le\sqrt{1+M_P^2}\,\|\mathscr S_P^{-1}\|.
\label{eq:v2v57-column-bound}
\end{equation}
Since \(A^{-1}\) is self-adjoint,
\begin{equation}
 \|PA^{-1}\|=\|A^{-1}P\|.
\label{eq:v2v57-row-column}
\end{equation}
\end{lemma}

\begin{proof}
Given \(y\in P\mathcal H\), solve
\[
 A\binom{x_P}{x_Q}=\binom y0.
\]
The \(Q\)-row gives
\[
 x_Q=-A_{QQ}^{-1}A_{QP}x_P.
\]
Substitution in the \(P\)-row gives
\(\mathscr S_Px_P=y\).  This proves
\eqref{eq:v2v57-PPP}--\eqref{eq:v2v57-QPP}.  The two components are
orthogonal, so
\[
 \|A^{-1}Py\|^2
 \le
 (1+M_P^2)\|\mathscr S_P^{-1}\|^2\|y\|^2.
\]
Finally,
\((A^{-1}P)^*=PA^{-1}\).
\end{proof}

\subsection{A quantitative feedback condition}
\label{sec:v2v57-feedback}

Assume \(A_{PP}\) is invertible and define the dimensionless return operator
\begin{equation}
 \mathscr B_P
 :=
 A_{PP}^{-1}A_{PQ}A_{QQ}^{-1}A_{QP}.
\label{eq:v2v57-return}
\end{equation}
Then
\[
 \mathscr S_P=A_{PP}(I-\mathscr B_P).
\]

\begin{theorem}[Transition-shell inverse localization]
\label{thm:v2v57-localization}
Let \(\Lambda\ge1\).  Suppose
\begin{align}
 \|A_{PP}^{-1}\|
 &\le\frac{c_P}{\Lambda},
\label{eq:v2v57-shell-coercivity}\\
 \|\mathscr B_P\|
 &\le\theta<1,
\label{eq:v2v57-feedback-bound}\\
 \|A_{QQ}^{-1}A_{QP}\|
 &\le M.
\label{eq:v2v57-one-way-mixing}
\end{align}
Then
\begin{equation}
\boxed{
 \|A^{-1}P\|+\|PA^{-1}\|
 \le
 \frac{2c_P\sqrt{1+M^2}}{(1-\theta)\Lambda}.}
\label{eq:v2v57-localization}
\end{equation}
\end{theorem}

\begin{proof}
The Neumann series and
\eqref{eq:v2v57-feedback-bound} give
\[
 \|(I-\mathscr B_P)^{-1}\|
 \le\frac1{1-\theta}.
\]
Since
\[
 \mathscr S_P^{-1}
 =
 (I-\mathscr B_P)^{-1}A_{PP}^{-1},
\]
equation \eqref{eq:v2v57-shell-coercivity} yields
\[
 \|\mathscr S_P^{-1}\|
 \le\frac{c_P}{(1-\theta)\Lambda}.
\]
Apply Lemma~\ref{lem:v2v57-column} and
\eqref{eq:v2v57-row-column}.
\end{proof}

\begin{assessment}[The two mixing quantities have different roles]
\label{ass:v2v57-two-mixing}
The return norm \(\theta\) protects invertibility of the Schur complement.
The one-way norm \(M\) controls how much complementary mass appears in the
inverse column after that complement is inverted.  A bound on either one
alone does not imply \eqref{eq:v2v57-localization}.  Both are measured
after the complete interaction and gauge blocks have been assembled.
\end{assessment}

\subsection{Commutator form for a reference spectral shell}
\label{sec:v2v57-commutator}

Let \(D_0=D_0^*\) commute with \(P\), and write
\[
 A=D_0+V.
\]
Then
\begin{equation}
 A_{QP}=QVP=Q[V,P]P,
 \qquad
 A_{PQ}=-P[V,P]Q.
\label{eq:v2v57-offdiagonal}
\end{equation}
Thus off-diagonal mixing is a cutoff commutator, not the full size of
\(V\).

\begin{corollary}[A lower-order commutator criterion]
\label{cor:v2v57-lower-order}
Assume
\begin{align}
 \|A_{PP}^{-1}\|&\le c_P\Lambda^{-1},
\label{eq:v2v57-cor-P}\\
 \|A_{QQ}^{-1}\|&\le c_Q,
\label{eq:v2v57-cor-Q}\\
 \|Q[V,P]P\|&\le v\Lambda^\sigma
\label{eq:v2v57-cor-mix}
\end{align}
and the same bound for its adjoint, with \(0\le\sigma<1/2\).  Then
\begin{align}
 M&\le c_Qv\Lambda^\sigma,
\label{eq:v2v57-cor-M}\\
 \theta&\le c_Pc_Qv^2\Lambda^{2\sigma-1}.
\label{eq:v2v57-cor-theta}
\end{align}
For all sufficiently large \(\Lambda\),
\begin{equation}
 \|A^{-1}P\|+\|PA^{-1}\|
 \le C\Lambda^{\sigma-1}.
\label{eq:v2v57-cor-localization}
\end{equation}
In the order-zero mixing case \(\sigma=0\), the full
\(\Lambda^{-1}\) shell gain is retained.
\end{corollary}

\begin{proof}
Equations \eqref{eq:v2v57-offdiagonal} and
\eqref{eq:v2v57-cor-Q} give \eqref{eq:v2v57-cor-M}.  Submultiplicativity
in \eqref{eq:v2v57-return} gives
\[
 \|\mathscr B_P\|
 \le
 c_P\Lambda^{-1}\,
 v\Lambda^\sigma\,
 c_Q\,
 v\Lambda^\sigma,
\]
which is \eqref{eq:v2v57-cor-theta}.  Since \(2\sigma-1<0\), this is below
one half for sufficiently large \(\Lambda\).  Apply
Theorem~\ref{thm:v2v57-localization} and use
\(\sqrt{1+c_Q^2v^2\Lambda^{2\sigma}}\le C\Lambda^\sigma\).
\end{proof}

\begin{remark}[What smooth spectral cutoffs are expected to buy]
\label{rem:v2v57-smooth}
For a first-order interaction, its full norm on the shell may be
\(O(\Lambda)\), while a commutator with a smooth cutoff can have lower
effective order because differentiating the cutoff symbol contributes an
inverse power of \(\Lambda\).  Corollary~\ref{cor:v2v57-lower-order}
requires this reduction as an operator estimate; symbolic order counting
alone is not a proof of \eqref{eq:v2v57-cor-mix}.
\end{remark}

\subsection{Noncommuting extension of the V9 covariance estimate}
\label{sec:v2v57-V9}

Let \(F_n\), \(\Delta F_n\), \(C_n=(D_0+V_n)^{-1}\), and the transition
projection \(P_n\) be as in V9, but do not assume
\([D_0,V_n]=0\).  To avoid collision with the V9 notation for the
regulated covariance, write
\[
 \mathscr A_n=F_nC_nF_n.
\]
Assume
\begin{align}
 \Delta F_n&=P_n\Delta F_n=\Delta F_nP_n,
\label{eq:v2v57-transition-F}\\
 \operatorname{rank}P_n&\le W\Lambda_n^d,
\label{eq:v2v57-rank}\\
 \|C_nP_n\|+\|P_nC_n\|
 &\le c_C\Lambda_n^{-1},
\label{eq:v2v57-C-localized}
\end{align}
where \eqref{eq:v2v57-C-localized} may be obtained from
Theorem~\ref{thm:v2v57-localization}.

\begin{theorem}[The two cutoff terms need no commutation]
\label{thm:v2v57-two-cutoff}
Under
\eqref{eq:v2v57-transition-F}--%
\eqref{eq:v2v57-C-localized},
\begin{align}
 \|\Delta F_nC_nF_n\|_{\rm HS}
 &\le c_C\sqrt W\,\Lambda_n^{d/2-1},
\label{eq:v2v57-left-cutoff}\\
 \|F_{n-1}C_n\Delta F_n\|_{\rm HS}
 &\le c_C\sqrt W\,\Lambda_n^{d/2-1}.
\label{eq:v2v57-right-cutoff}
\end{align}
\end{theorem}

\begin{proof}
The transition multiplier has
\[
 \|\Delta F_n\|_{\rm HS}
 \le\sqrt{\operatorname{rank}P_n}
 \le\sqrt W\,\Lambda_n^{d/2},
\]
and \(\|F_j\|\le1\).  Insert \(P_n\) next to the inverse:
\begin{align*}
 \Delta F_nC_nF_n
 &=\Delta F_nP_nC_nF_n,\\
 F_{n-1}C_n\Delta F_n
 &=F_{n-1}C_nP_n\Delta F_n.
\end{align*}
Schatten ideal inequalities and
\eqref{eq:v2v57-C-localized} give the two estimates.
\end{proof}

\begin{corollary}[A transition-supported running step]
\label{cor:v2v57-running}
Assume additionally
\begin{equation}
 \Delta V_n=P_n\Delta V_nP_n,
 \qquad
 \|\Delta V_n\|\le v\Lambda_n^\rho,
\label{eq:v2v57-running-support}
\end{equation}
and that \eqref{eq:v2v57-C-localized} holds for both \(C_n\) and
\(C_{n-1}\).  Then
\begin{equation}
 \|F_{n-1}C_n\Delta V_nC_{n-1}F_{n-1}\|_{\rm HS}
 \le
 c_C^2v\sqrt W\,\Lambda_n^{d/2+\rho-2}.
\label{eq:v2v57-running-HS}
\end{equation}
\end{corollary}

\begin{proof}
Insert \(P_n\) on both sides of \(\Delta V_n\).  The resulting operator
has rank at most \(\operatorname{rank}P_n\) and operator norm at most
\[
 (c_C\Lambda_n^{-1})
 (v\Lambda_n^\rho)
 (c_C\Lambda_n^{-1}).
\]
The Hilbert--Schmidt norm is bounded by the operator norm times
\(\sqrt{\operatorname{rank}P_n}\).
\end{proof}

\begin{corollary}[Noncommuting projected-row threshold]
\label{cor:v2v57-projected-row}
Under Theorem~\ref{thm:v2v57-two-cutoff},
Corollary~\ref{cor:v2v57-running}, and the gauge-commutator and insertion
bounds of V9, the Taylor-projected adjacent row obeys the same powers and
the same summability threshold
\eqref{eq:v2v9-J-threshold} as the commuting V9 regression.
\end{corollary}

\begin{proof}
The exact three-term identity
\eqref{eq:v2v9-general-adjacent} does not use commutation.  Theorem
\ref{thm:v2v57-two-cutoff} supplies the V9 power
\(\Lambda_n^{d/2-1}\) for each cutoff term, and
Corollary~\ref{cor:v2v57-running} supplies
\(\Lambda_n^{d/2+\rho-2}\) for the running term.  Repeat the Schatten
H\"older and Taylor-remainder proof of
Theorem~\ref{thm:v2v9-projected-J}.
\end{proof}

\subsection{Global coercivity does not localize the inverse}
\label{sec:v2v57-counterexample}

\begin{proposition}[A rotated two-level counterexample]
\label{prop:v2v57-counterexample}
Let
\[
 U=\frac1{\sqrt2}
 \begin{pmatrix}1&-1\\1&1\end{pmatrix},
 \qquad
 A_\Lambda
 =
 U\begin{pmatrix}1&0\\0&\Lambda\end{pmatrix}U^*,
 \qquad
 P=\begin{pmatrix}0&0\\0&1\end{pmatrix}.
\]
Then
\begin{equation}
 \|A_\Lambda^{-1}\|=1,
\label{eq:v2v57-global-inverse}
\end{equation}
but
\begin{equation}
 \|A_\Lambda^{-1}P\|
 =
 \sqrt{\frac{1+\Lambda^{-2}}2}
 \longrightarrow\frac1{\sqrt2}.
\label{eq:v2v57-bad-column}
\end{equation}
Thus a cutoff-independent global inverse bound does not imply
\(O(\Lambda^{-1})\) localization on the reference transition shell.
\end{proposition}

\begin{proof}
Unitary conjugation gives
\[
 A_\Lambda^{-1}
 =
 U\begin{pmatrix}1&0\\0&\Lambda^{-1}\end{pmatrix}U^*.
\]
Its norm is one.  The norm of \(A_\Lambda^{-1}P\) is the Euclidean norm of
the second column of \(A_\Lambda^{-1}\), namely
\[
 \frac12
 \sqrt{(1-\Lambda^{-1})^2+(1+\Lambda^{-1})^2}
 =
 \sqrt{\frac{1+\Lambda^{-2}}2}.
\]
\end{proof}

\begin{assessment}[How the feedback condition fails]
\label{ass:v2v57-counterexample}
For the reference block \(P\), both diagonal compressions of
\(A_\Lambda\) are of order \(\Lambda\), but the off-diagonal block is also
of order \(\Lambda\).  The return norm
\(\|\mathscr B_P\|\) tends to one, so the Schur complement falls to order
one.  The counterexample therefore violates exactly
\eqref{eq:v2v57-feedback-bound}; it is not excluded by global coercivity.
\end{assessment}

\begin{firewall}[Approximate shell support creates a separate error]
\label{fw:v2v57-approximate}
Corollary~\ref{cor:v2v57-running} assumes the exact support
\(P_n\Delta V_nP_n\).  If
\[
 \Delta V_n=P_n\Delta V_nP_n+\mathcal E_n,
\]
the remainder contributes
\[
 F_{n-1}C_n\mathcal E_nC_{n-1}F_{n-1}.
\]
Neither transition rank nor
\eqref{eq:v2v57-C-localized} controls this term because
\(\mathcal E_n\) may couple to low modes on both sides.  A quasilocal
off-shell norm for \(\mathcal E_n\) remains necessary.
\end{firewall}

\subsection{V57 conclusion and next acquisition}
\label{sec:v2v57-conclusion}

V57 closes the two noncommuting cutoff terms of the V9 adjacent covariance
regression under an explicit Schur-feedback condition, and it closes the
running covariance term when the running step is transition-supported.
The result identifies the exact off-diagonal mixing norm needed to retain
the ultraviolet inverse factor.  The rotated two-level model proves that a
global inverse estimate alone cannot supply it.

This does not yet close the complete SM shell increment.  The next
unresolved term is the running composite insertion identified in V9--V10,
together with the off-shell remainder in
Firewall~\ref{fw:v2v57-approximate}.  V58 will formulate a connected
Duhamel estimate for the insertion difference in the common physical
Banach algebra, preserving the species and Ward assembly before the
two-sided Abel stock is taken.
\clearpage
\section*{Second Volume, Version V58: Connected Duhamel Control of the Running Insertion}
\addcontentsline{toc}{section}{Second Volume, Version V58: Connected Duhamel Control of the Running Insertion}

\subsection*{Purpose}

V57 controls the noncommuting covariance change when the transition-shell
resolvent satisfies an explicit Schur-feedback estimate.  The other
adjacent term identified in V9 is the change of the composite insertion
itself.  V10 gave its exact normalized Duhamel formula: two direct
covariances and one third connected cumulant.  The missing step is a norm
that respects the fact that all three variables use the same resampled
shell.

This note supplies that norm.  A maximum-index filtration estimate bounds
the third cumulant by four common-bridge stocks.  Integrating those stocks
along the normalized interpolation controls the complete running
insertion.  The local Taylor projector is applied only after the connected
species sum has been formed.  The theorem is rigorous for a positive
finite-cutoff interpolation and has an abstract cumulant-norm version for
a signed BV or Berezin functional.  It does not assert the required
interacting SM cluster bounds.

\subsection{Filtration seminorms for one connected bridge}
\label{sec:v2v58-filtration}

Fix one of the finite filtrations used in V10.  For
\(X\in L^3(\Omega,\mu)\), write
\[
 d_mX=\mathbb E[X\mid\mathcal F_m]
      -\mathbb E[X\mid\mathcal F_{m-1}],
 \qquad
 \widehat X_m=\mathbb E[X\mid\mathcal F_m]-\mathbb EX.
\]
Define
\begin{align}
 \mathsf S_2(X)
 &:=
 \left(\sum_m\|d_mX\|_3^2\right)^{1/2},
\label{eq:v2v58-S2}\\
 \mathsf S_3(X)
 &:=
 \left(\sum_m\|d_mX\|_3^3\right)^{1/3},
\label{eq:v2v58-S3}\\
 \mathsf M_3(X)
 &:=
 \sup_m\|\widehat X_m\|_3.
\label{eq:v2v58-M3}
\end{align}
All three quantities are invariant under adding a constant to \(X\).

\begin{theorem}[Common-bridge third-cumulant bound]
\label{thm:v2v58-common-bridge}
For \(F,G,H\in L^3\),
\begin{align}
 |\kappa_3(F,G,H)|
 \le{}&
 \mathsf S_2(F)\mathsf S_2(G)\mathsf M_3(H)
 \notag\\
 &+
 \mathsf S_2(F)\mathsf S_2(H)\mathsf M_3(G)
 \notag\\
 &+
 \mathsf S_2(G)\mathsf S_2(H)\mathsf M_3(F)
 \notag\\
 &+
 \mathsf S_3(F)\mathsf S_3(G)\mathsf S_3(H).
\label{eq:v2v58-common-bridge}
\end{align}
\end{theorem}

\begin{proof}
Use the exact maximum-index decomposition
\eqref{eq:v2v10-third-martingale}.  For its first term,
Cauchy--Schwarz in the common index \(m\) gives
\begin{align*}
 \sum_m
 \|d_mF\|_3\|d_mG\|_3\|\widehat H_{m-1}\|_3
 &\le
 \mathsf M_3(H)
 \left(\sum_m\|d_mF\|_3^2\right)^{1/2}
 \left(\sum_m\|d_mG\|_3^2\right)^{1/2}.
\end{align*}
The next two terms are identical after permutation.  For the last term,
H\"older's inequality on the counting measure gives
\[
 \sum_m
 \|d_mF\|_3\|d_mG\|_3\|d_mH\|_3
 \le
 \mathsf S_3(F)\mathsf S_3(G)\mathsf S_3(H).
\]
\end{proof}

\begin{assessment}[No independent bridge count]
\label{ass:v2v58-one-bridge}
Each product in \eqref{eq:v2v58-common-bridge} is formed only after the
same maximum filtration index has been selected.  Replacing it by three
independent sums over \(m\) would introduce configurations that are absent
from the connected cumulant and can lose the shell power.  This is the
insertion analogue of the one-capacity rule in the V55 YM audit.
\end{assessment}

\subsection{Uniform connected insertion kernels}
\label{sec:v2v58-kernel}

For the \(n\)-th interpolation, let
\(\langle\cdot\rangle_{n,t}\), \(0\le t\le1\), be the normalized state in
V10, with
\[
 S_{n,t}=S_{n-1}+t\Delta S_n.
\]
Let \(\mathcal U_F,\mathcal U_G\) be fixed balanced sets of external test
data.  The complete Ward-assembled observables are denoted
\[
 F_{n,t}(f),\qquad G_{n,t}(g),
 \qquad f\in\mathcal U_F,\quad g\in\mathcal U_G.
\]
When a positive physical representation has already been constructed, they
are the complete observables induced by the matter, gauge, ghost, contact,
boundary, and antifield row after physical reduction.  Before that reduction
the signed-functional theorem below, rather than the probability estimate,
is the applicable statement.

Define the connected insertion seminorm
\begin{equation}
 \|B_{n,t}\|_{\mathcal Z}
 :=
 \sup_{\substack{f\in\mathcal U_F\\g\in\mathcal U_G}}
 \left|
 \operatorname{Cov}_{n,t}
 \bigl(F_{n,t}(f),G_{n,t}(g)\bigr)
 \right|.
\label{eq:v2v58-Znorm}
\end{equation}
For a random variable \(X\), write
\(X^\circ=X-\langle X\rangle_{n,t}\).

For fixed \(n,t,f,g\), put
\begin{align}
 \mathfrak A_{n,t}(f,g)
 :={}&
 \|\dot F_{n,t}(f)^\circ\|_2
 \|G_{n,t}(g)^\circ\|_2
 \notag\\
 &+
 \|F_{n,t}(f)^\circ\|_2
 \|\dot G_{n,t}(g)^\circ\|_2
 \notag\\
 &+\beta\,
 \mathfrak M_{n,t}
 \bigl(F_{n,t}(f),G_{n,t}(g),\Delta S_n\bigr),
\label{eq:v2v58-A}
\end{align}
where
\begin{align}
 \mathfrak M(F,G,H)
 :={}&
 \mathsf S_2(F)\mathsf S_2(G)\mathsf M_3(H)
 \notag\\
 &+
 \mathsf S_2(F)\mathsf S_2(H)\mathsf M_3(G)
 \notag\\
 &+
 \mathsf S_2(G)\mathsf S_2(H)\mathsf M_3(F)
 \notag\\
 &+
 \mathsf S_3(F)\mathsf S_3(G)\mathsf S_3(H).
\label{eq:v2v58-M}
\end{align}
The seminorms in \eqref{eq:v2v58-M} are computed in the interpolated
probability state at \((n,t)\).

\begin{theorem}[Integrated running-insertion estimate]
\label{thm:v2v58-insertion}
Assume the V10 differentiation formula is justified for every external
test pair and
\begin{equation}
 c_n:=
 \int_0^1
 \sup_{\substack{f\in\mathcal U_F\\g\in\mathcal U_G}}
 \mathfrak A_{n,t}(f,g)\,\mathrm dt
 <\infty.
\label{eq:v2v58-cn}
\end{equation}
Then
\begin{equation}
 \|B_{n,1}-B_{n,0}\|_{\mathcal Z}
 \le c_n.
\label{eq:v2v58-insertion-bound}
\end{equation}
If \(\sum_nc_n<\infty\), the cumulative connected insertions are Cauchy in
\(\mathcal Z\).
\end{theorem}

\begin{proof}
The exact adjacent formula
\eqref{eq:v2v10-adjacent-insertion} gives, for each \(f,g\),
two covariance terms and the third cumulant with \(\Delta S_n\).
Cauchy--Schwarz in the probability state bounds the direct terms by the
first two lines of \eqref{eq:v2v58-A}.  Theorem
\ref{thm:v2v58-common-bridge} bounds the cumulant by
\(\mathfrak M\).  Integrate in \(t\), take the supremum over the external
tests, and use \eqref{eq:v2v58-cn}.  Summability gives the Cauchy criterion.
\end{proof}

\begin{corollary}[Dyadic cluster gain]
\label{cor:v2v58-dyadic}
If, for some \(\varepsilon>0\),
\begin{equation}
 \sup_{0\le t\le1}
 \sup_{f,g}
 \mathfrak A_{n,t}(f,g)
 \le C\Lambda_n^{-\varepsilon},
 \qquad
 \Lambda_n\simeq2^n,
\label{eq:v2v58-dyadic}
\end{equation}
then the running insertion difference is absolutely summable in
\(\mathcal Z\).
\end{corollary}

\begin{proof}
Equation~\eqref{eq:v2v58-cn} gives
\(c_n\le C\Lambda_n^{-\varepsilon}\), and the dyadic series is geometric.
\end{proof}

\subsection{Projection after complete connected assembly}
\label{sec:v2v58-projection}

Let \(P_{\rm loc}:\mathcal Z\to\mathcal Z_{\rm loc}\) be the single fixed
BRST-compatible local Taylor projector used in V54--V56.

\begin{proposition}[The nonlocal insertion increment]
\label{prop:v2v58-projected}
Under Theorem~\ref{thm:v2v58-insertion},
\begin{align}
 (I-P_{\rm loc})(B_{n,1}-B_{n,0})
 =
 \int_0^1(I-P_{\rm loc})\Bigl[
 &\operatorname{Cov}_{n,t}(\dot F_{n,t},G_{n,t})
 \notag\\
 &+\operatorname{Cov}_{n,t}(F_{n,t},\dot G_{n,t})
 \notag\\
 &-\beta\kappa_{3,n,t}
   (F_{n,t},G_{n,t},\Delta S_n)
 \Bigr]\,\mathrm dt.
\label{eq:v2v58-projected}
\end{align}
If the norm of the complete projected integrand obeys
\[
 \sup_{f,g}
 \left\|(I-P_{\rm loc})[\cdots]\right\|_{\mathcal Z}
 \le a_n(t)
\]
and \(\sum_n\int_0^1a_n(t)\,\mathrm dt<\infty\), then the projected
insertion increments are absolutely summable.
\end{proposition}

\begin{proof}
Apply the one fixed bounded linear projector to the exact V10 identity.
Linearity permits it to pass through the integral.  The asserted estimate
then follows from the Bochner triangle inequality after the complete
projected bracket has been formed.
\end{proof}

\begin{firewall}[The displayed scalar majorant is a fallback]
\label{fw:v2v58-project-after}
The estimate \eqref{eq:v2v58-A} norms the three connected terms
separately and is therefore a sufficient fallback.  The sharper
proposition allows cancellation among their complete Ward-assembled
sum before the \(\mathcal Z\)-norm.  Projecting or norming each species,
orientation, or contact term first can destroy the local Slavnov
cancellation and is not authorized by
Proposition~\ref{prop:v2v58-projected}.
\end{firewall}

\subsection{Signed and Berezin interpolations}
\label{sec:v2v58-signed}

The \(L^p\) proof above uses positivity of the interpolated measure.  A
finite-cutoff chiral BV integral contains Berezin and gauge-fixed signed
components.  For that case define the connected bilinear and trilinear
operator norms directly:
\begin{align}
 \|\mathcal C_{2,n,t}\|
 &:=
 \sup_{\substack{\|X\|_{\mathcal X_1}\le1\\\|Y\|_{\mathcal X_2}\le1}}
 |\operatorname{Cov}_{n,t}(X,Y)|,
\label{eq:v2v58-C2}\\
 \|\mathcal C_{3,n,t}\|
 &:=
 \sup_{\substack{\|X\|_{\mathcal X_1}\le1,\ \|Y\|_{\mathcal X_2}\le1\\
                    \|H\|_{\mathcal X_3}\le1}}
 |\kappa_{3,n,t}(X,Y,H)|.
\label{eq:v2v58-C3}
\end{align}

\begin{theorem}[Abstract connected-functional version]
\label{thm:v2v58-abstract}
Suppose the normalized Duhamel identity
\eqref{eq:v2v10-adjacent-insertion} holds algebraically for the complete
finite-cutoff functional, and the norms
\eqref{eq:v2v58-C2}--\eqref{eq:v2v58-C3} are finite.  Then
\begin{align}
 \|B_{n,1}-B_{n,0}\|_{\mathcal Z}
 \le\int_0^1\Bigl[
 &\|\mathcal C_{2,n,t}\|
  \bigl(
   \|\dot F_{n,t}\|_{\mathcal X_1}\|G_{n,t}\|_{\mathcal X_2}
   +
   \|F_{n,t}\|_{\mathcal X_1}\|\dot G_{n,t}\|_{\mathcal X_2}
  \bigr)
 \notag\\
 &+
 \beta\|\mathcal C_{3,n,t}\|
 \|F_{n,t}\|_{\mathcal X_1}
 \|G_{n,t}\|_{\mathcal X_2}
 \|\Delta S_n\|_{\mathcal X_3}
 \Bigr]\,\mathrm dt,
\label{eq:v2v58-abstract-bound}
\end{align}
provided the right side is uniform over the external test unit balls.
\end{theorem}

\begin{proof}
Apply the defining bilinear and trilinear norms to the three terms in the
exact Duhamel identity and integrate.
\end{proof}

\begin{assessment}[What replaces positivity in the BV branch]
\label{ass:v2v58-BV}
The abstract theorem does not infer
\(\|\mathcal C_3\|\) from a signed total-variation measure.  The needed
bound must come from exact Grassmann traces, bosonic cluster estimates,
and their Ward-completed mixed cumulants.  V10--V11 prove the purely
Gaussian quadratic regressions.  Higher interacting cumulants remain a
separate analytic input.
\end{assessment}

\subsection{Interface with the V56 Abel compiler}
\label{sec:v2v58-Abel}

Let the complete nonlocal shell derivation split algebraically as
\begin{equation}
 q_n=q_n^{\rm cov}+q_n^{\rm ins}+q_n^{\rm rep},
\label{eq:v2v58-q-split}
\end{equation}
where the terms denote the adjacent covariance row, the running insertion,
and the Ward/anomaly repair.  This identity is taken before any norm.

\begin{theorem}[A sufficient absolutely summable insertion branch]
\label{thm:v2v58-Abel-interface}
Suppose:
\begin{enumerate}
 \item V57 and the fixed Taylor subtraction give
       \(\sum_n\|q_n^{\rm cov}\|<\infty\);
 \item Proposition~\ref{prop:v2v58-projected} or
       Theorem~\ref{thm:v2v58-abstract} gives
       \(\sum_n\|q_n^{\rm ins}\|<\infty\);
 \item the corrected local cohomology construction gives
       \(\sum_n\|q_n^{\rm rep}\|<\infty\).
\end{enumerate}
Then \(\sum_n\|q_n\|<\infty\), and the V54 product theorem constructs the
physical algebra chart.  More generally, if any branch is only
conditionally convergent, the conclusion follows only after the complete
\(q_n\) and its cumulative tail satisfy the two ordered V56 stocks.
\end{theorem}

\begin{proof}
The triangle inequality after the algebraic assembly gives
\[
 \sum_n\|q_n\|
 \le
 \sum_n\left(
 \|q_n^{\rm cov}\|+
 \|q_n^{\rm ins}\|+
 \|q_n^{\rm rep}\|
 \right)<\infty.
\]
Apply Theorem~\ref{thm:v2v54-product}.  The conditional statement is
exactly Theorem~\ref{thm:v2v56-compiler}; separate conditional convergence
does not control its ordered cross terms.
\end{proof}

\begin{firewall}[No interacting cluster theorem has been smuggled in]
\label{fw:v2v58-open}
Theorem~\ref{thm:v2v58-insertion} converts explicit cluster seminorms into
a shell bound.  It does not prove
\eqref{eq:v2v58-dyadic} for the interacting Higgs--Yukawa--gauge measure,
nor does the abstract version prove
\(\|\mathcal C_{3,n,t}\|<\infty\) uniformly in the cutoff.  Those are the
current running-insertion estimates.  The Gaussian powers in V10--V11 fix
the expected local subtractions but do not settle this interacting gate.
\end{firewall}

\subsection{V58 conclusion and next acquisition}
\label{sec:v2v58-conclusion}

V58 supplies the missing normed form of the V10 running-insertion identity.
The positive branch is controlled by one maximum-index martingale stock;
the signed BV branch is controlled by an explicit connected-functional
norm.  A fixed local projector acts only after the complete connected
Duhamel bracket is assembled.  Absolute summability of this insertion
branch combines with the V57 covariance branch in V54, while conditional
convergence must use the full two-sided V56 stock.

The next unresolved covariance contribution is the off-shell running
remainder in Firewall~\ref{fw:v2v57-approximate}.  V59 will introduce a
weighted block norm for quasilocal shell leakage and prove when
\(C_n\mathcal E_nC_{n-1}\) remains summable despite low-mode couplings.
It will include a rank-one leakage example showing the exact decay
threshold.
\clearpage
\section*{Second Volume, Version V59: Weighted Quasilocal Control of Off-Shell Running Leakage}
\addcontentsline{toc}{section}{Second Volume, Version V59: Weighted Quasilocal Control of Off-Shell Running Leakage}

\subsection*{Purpose}

V57 treated a running step supported exactly in the transition projection
\(P_n\).  If
\[
 \Delta V_n=P_n\Delta V_nP_n+\mathcal E_n,
\]
the remainder \(\mathcal E_n\) may visit low modes, where a global inverse
has no ultraviolet gain.  This note gives the correct weighted norm for
that leakage.

The central identity inserts one spectral weight on each side of the
complete Ward-assembled leakage operator.  All block cancellations remain
inside that weighted operator until its norm is taken.  A two-level model
then shows the exact distinction among low--low, low--shell, and
shell--shell leakage: they receive respectively zero, one, and two inverse
shell factors.  Thus quasilocality must specify which weighted block norm
decays; unweighted smallness is not enough.

\subsection{A spectral inverse-weight frame}
\label{sec:v2v59-weight}

At shell \(n\), let
\[
 I=\sum_{a\in\mathcal I_n}P_{n,a}
\]
be a finite orthogonal block decomposition in the regulated Hilbert space.
Choose positive inverse scales \(r_{n,a}>0\) and define
\begin{equation}
 \mathsf R_n
 =
 \sum_{a\in\mathcal I_n}r_{n,a}P_{n,a}.
\label{eq:v2v59-R}
\end{equation}
For a low massive block one may take \(r_{n,a}\asymp1\); for the active
transition shell, V57 gives \(r_{n,a}\asymp\Lambda_n^{-1}\).  At finite
cutoff \(\mathsf R_n\) is invertible.

Let \(F_n\) be the low-pass contraction and let \(C_n\) be the complete
regulated inverse.  The weighted inverse-frame hypothesis is
\begin{align}
 \|F_{n-1}C_n\mathsf R_n^{-1}\|
 +
 \|\mathsf R_n^{-1}C_nF_{n-1}\|
 &\le c_C,
\label{eq:v2v59-frame-n}\\
 \|F_{n-1}C_{n-1}\mathsf R_n^{-1}\|
 +
 \|\mathsf R_n^{-1}C_{n-1}F_{n-1}\|
 &\le c_C.
\label{eq:v2v59-frame-prev}
\end{align}
These are global assembled estimates.  They may be proved blockwise using
V57, but they are not consequences of a global inverse norm.

\begin{theorem}[Weighted leakage sandwich]
\label{thm:v2v59-sandwich}
For any complete leakage operator \(\mathcal E_n\),
\begin{equation}
\boxed{
 \|F_{n-1}C_n\mathcal E_nC_{n-1}F_{n-1}\|_{\rm HS}
 \le
 c_C^2
 \|\mathsf R_n\mathcal E_n\mathsf R_n\|_{\rm HS}.}
\label{eq:v2v59-sandwich}
\end{equation}
The same statement holds in operator norm or in any two-sided operator
ideal norm.
\end{theorem}

\begin{proof}
Insert
\(\mathsf R_n^{-1}\mathsf R_n=I\) on both sides of
\(\mathcal E_n\):
\begin{align*}
 F_{n-1}C_n\mathcal E_nC_{n-1}F_{n-1}
 =
 &(F_{n-1}C_n\mathsf R_n^{-1})
 (\mathsf R_n\mathcal E_n\mathsf R_n)\\
 &\cdot
 (\mathsf R_n^{-1}C_{n-1}F_{n-1}).
\end{align*}
Apply the two-sided ideal inequality and
\eqref{eq:v2v59-frame-n}--\eqref{eq:v2v59-frame-prev}.
\end{proof}

\begin{assessment}[Where cancellation is retained]
\label{ass:v2v59-global}
The norm on the right of \eqref{eq:v2v59-sandwich} is taken after
\(\mathcal E_n\) has been assembled from every matter, gauge, ghost,
contact, and anomaly-repair contribution.  It is not the sum of their
weighted norms.  This preserves the V55 audit rule and allows Ward
cancellation inside each spectral block and between blocks before the
single ideal norm.
\end{assessment}

\subsection{A blockwise sufficient estimate}
\label{sec:v2v59-block}

Write
\[
 \mathcal E_{n,ab}
 =
 P_{n,a}\mathcal E_nP_{n,b}.
\]
Then
\begin{equation}
 \mathsf R_n\mathcal E_n\mathsf R_n
 =
 \sum_{a,b}r_{n,a}r_{n,b}\mathcal E_{n,ab}.
\label{eq:v2v59-block-expansion}
\end{equation}

\begin{corollary}[Absolute block fallback]
\label{cor:v2v59-block}
If the block series is absolutely convergent in the chosen ideal, then
\begin{equation}
 \|\mathsf R_n\mathcal E_n\mathsf R_n\|_{\rm HS}
 \le
 \sum_{a,b}
 r_{n,a}r_{n,b}
 \|\mathcal E_{n,ab}\|_{\rm HS}.
\label{eq:v2v59-block-bound}
\end{equation}
Consequently the right side controls the leakage sandwich.
\end{corollary}

\begin{proof}
Use \eqref{eq:v2v59-block-expansion} and the triangle inequality, then
Theorem~\ref{thm:v2v59-sandwich}.
\end{proof}

\begin{firewall}[The block fallback can be strictly weaker]
\label{fw:v2v59-block-fallback}
Equation~\eqref{eq:v2v59-block-bound} discards cancellation among blocks
and species.  It is only a sufficient estimate.  The primary analytic
target is the single weighted norm in
\eqref{eq:v2v59-sandwich}; the block sum should be used only when its loss
is harmless.
\end{firewall}

\subsection{The exact low--shell model}
\label{sec:v2v59-two-level}

Let
\[
 \mathcal H=\mathbb Ce_L\oplus\mathbb Ce_S,
 \qquad
 C_\Lambda=
 \begin{pmatrix}1&0\\0&\Lambda^{-1}\end{pmatrix},
 \qquad
 \mathsf R_\Lambda=
 \begin{pmatrix}1&0\\0&\Lambda^{-1}\end{pmatrix}.
\]
Then \(C_\Lambda\mathsf R_\Lambda^{-1}=I\), so the weighted frame constant
is one.

\begin{proposition}[Exact rank-one leakage weights]
\label{prop:v2v59-rank-one}
For scalar \(\varepsilon_\Lambda\), define
\begin{align*}
 \mathcal E_{LL}
 &=\varepsilon_\Lambda|e_L\rangle\langle e_L|,\\
 \mathcal E_{LS}
 &=\varepsilon_\Lambda|e_L\rangle\langle e_S|,\\
 \mathcal E_{SS}
 &=\varepsilon_\Lambda|e_S\rangle\langle e_S|.
\end{align*}
Then
\begin{align}
 \|C_\Lambda\mathcal E_{LL}C_\Lambda\|_{\rm HS}
 &=|\varepsilon_\Lambda|,
\label{eq:v2v59-LL}\\
 \|C_\Lambda\mathcal E_{LS}C_\Lambda\|_{\rm HS}
 &=|\varepsilon_\Lambda|\Lambda^{-1},
\label{eq:v2v59-LS}\\
 \|C_\Lambda\mathcal E_{SS}C_\Lambda\|_{\rm HS}
 &=|\varepsilon_\Lambda|\Lambda^{-2}.
\label{eq:v2v59-SS}
\end{align}
If \(\Lambda_n\simeq2^n\) and
\(|\varepsilon_{\Lambda_n}|\asymp\Lambda_n^\rho\), the three dyadic
leakage series are absolutely summable respectively when
\begin{equation}
 \rho<0,\qquad \rho<1,\qquad \rho<2.
\label{eq:v2v59-thresholds}
\end{equation}
These thresholds are exact.
\end{proposition}

\begin{proof}
Left and right multiplication by \(C_\Lambda\) contributes its diagonal
weight at the range and source of the rank-one operator.  This gives
\eqref{eq:v2v59-LL}--\eqref{eq:v2v59-SS}.  A dyadic power series
\(\sum_n\Lambda_n^\alpha\) converges exactly when \(\alpha<0\), yielding
\eqref{eq:v2v59-thresholds}.
\end{proof}

\begin{firewall}[Small support is not a frequency weight]
\label{fw:v2v59-low-low}
A rank-one low--low leakage has the smallest possible rank, but receives no
inverse shell factor.  If its amplitude stays of order one, its dyadic sum
diverges.  Thus finite rank, small spatial support, or a global inverse
bound cannot replace decay of the low--low weighted block.
\end{firewall}

\subsection{Insertion into the projected adjacent row}
\label{sec:v2v59-row}

Let \(\mathcal G\) denote the V9 smooth gauge generator, so that
\begin{equation}
 \|[F_n,\mathcal G]\|_{\rm HS}
 \le c_{\mathcal G}\Lambda_n^{d/2-1}
 \|\nabla\mathcal G\|_\infty.
\label{eq:v2v59-gauge}
\end{equation}
Let \(B_n(k)\) have the V9 insertion order \(\mu\), and subtract its external
Taylor jet through degree \(r\).

\begin{theorem}[Weighted leakage summability threshold]
\label{thm:v2v59-row}
Suppose
\begin{equation}
 \|\mathsf R_n\mathcal E_n\mathsf R_n\|_{\rm HS}
 \le C_E\Lambda_n^\chi.
\label{eq:v2v59-weighted-order}
\end{equation}
Then the leakage contribution
\[
 J_n^{\rm leak}(k)
 =
 \operatorname{Tr}\!\left(
  F_{n-1}C_n\mathcal E_nC_{n-1}F_{n-1}
  [F_n,\mathcal G]B_n(k)
 \right)
\]
obeys, after subtraction through degree \(r\),
\begin{equation}
 |Q_rJ_n^{\rm leak}(k)|
 \le
 C|k|^{r+1}\|\nabla\mathcal G\|_\infty
 \Lambda_n^{\chi+d/2-1+\mu-(r+1)}.
\label{eq:v2v59-row-bound}
\end{equation}
For dyadic \(\Lambda_n\), this contribution is absolutely summable if
\begin{equation}
 r+1>\chi+\frac d2+\mu-1.
\label{eq:v2v59-row-threshold}
\end{equation}
\end{theorem}

\begin{proof}
Theorem~\ref{thm:v2v59-sandwich} and
\eqref{eq:v2v59-weighted-order} give
\[
 \|F_{n-1}C_n\mathcal E_nC_{n-1}F_{n-1}\|_{\rm HS}
 \le c_C^2C_E\Lambda_n^\chi.
\]
Differentiate only the external insertion, use Schatten H\"older with
the Hilbert--Schmidt commutator
\eqref{eq:v2v59-gauge}, and apply the V9 Taylor remainder.  The exponents
add to the power in \eqref{eq:v2v59-row-bound}.  A dyadic power is
summable exactly when that exponent is negative.
\end{proof}

\begin{corollary}[Recovery of the transition-supported power]
\label{cor:v2v59-transition}
If \(\mathcal E_n\) is supported on a transition block of rank
\(O(\Lambda_n^d)\), has operator order \(\rho\), and the transition inverse
weight is \(r_n\asymp\Lambda_n^{-1}\), then
\[
 \chi=\frac d2+\rho-2.
\]
Condition~\eqref{eq:v2v59-row-threshold} becomes
\[
 r+1>d+\rho+\mu-3,
\]
which is exactly the running-row threshold of V9 and V57.
\end{corollary}

\begin{proof}
The weighted transition block has operator norm
\(O(\Lambda_n^{\rho-2})\) and rank \(O(\Lambda_n^d)\), hence
Hilbert--Schmidt norm
\(O(\Lambda_n^{d/2+\rho-2})\).  Substitute this \(\chi\).
\end{proof}

\subsection{A three-zone sufficient criterion}
\label{sec:v2v59-zones}

Let
\[
 I=P_L+P_S+P_H
\]
be low, transition-shell, and high zones, with inverse weights
\[
 r_L\asymp1,
 \qquad
 r_S\asymp\Lambda_n^{-1},
 \qquad
 0<r_H\le C\Lambda_n^{-1}.
\]

\begin{corollary}[Three-zone leakage ledger]
\label{cor:v2v59-zones}
A sufficient bound for \eqref{eq:v2v59-weighted-order} is
\begin{align}
 &\|P_L\mathcal E_nP_L\|_{\rm HS}
 \notag\\
 &+\Lambda_n^{-1}
 \bigl(
  \|P_L\mathcal E_nP_S\|_{\rm HS}
 +\|P_S\mathcal E_nP_L\|_{\rm HS}
 \bigr)
 \notag\\
 &+\Lambda_n^{-2}
 \|P_S\mathcal E_nP_S\|_{\rm HS}
 \notag\\
 &+
 \sum_{\substack{a,b\in\{L,S,H\}\\a=H\ {\rm or}\ b=H}}
 r_ar_b\|P_a\mathcal E_nP_b\|_{\rm HS}
 \le C\Lambda_n^\chi.
\label{eq:v2v59-zone-ledger}
\end{align}
\end{corollary}

\begin{proof}
This is Corollary~\ref{cor:v2v59-block} with the displayed weights.  Terms
involving \(H\) are left with their actual weights because high-zone
coercivity may be stronger than the common upper bound.
\end{proof}

\begin{assessment}[The low--low block is the sharp new gate]
\label{ass:v2v59-low-gate}
The transition and high blocks inherit inverse-frequency factors from V57.
Only the low--low component in
\eqref{eq:v2v59-zone-ledger} has no automatic shell decay.  A complete SM
proof must show that this block is local and removed by
\(P_{\rm loc}\), cancels by the Ward identity, or decays through a genuine
quasilocal estimate.  Calling it an off-shell remainder does not choose
among these alternatives.
\end{assessment}

\subsection{Interface with the V56 chart stock}
\label{sec:v2v59-Abel}

Let \(q_n^{\rm leak}\) be the nonlocal shell derivation produced by the
projected row \(Q_rJ_n^{\rm leak}\).

\begin{corollary}[Weighted decay feeds the algebra chart]
\label{cor:v2v59-Abel}
If the exponent in \eqref{eq:v2v59-row-bound} is
\(-\varepsilon<0\), then
\[
 \|q_n^{\rm leak}\|
 \le C\Lambda_n^{-\varepsilon}
\]
in the chosen physical Banach-algebra norm, and
\(\sum_n\|q_n^{\rm leak}\|<\infty\).  It may therefore be inserted in the
absolute V54 branch of the complete chart.  At the endpoint exponent zero,
no such conclusion holds; a signed cumulative limit and both ordered V56
correlations must be proved instead.
\end{corollary}

\begin{proof}
The dyadic inverse power is summable.  The endpoint statement follows
because a uniform \(O(1)\) bound supplies neither additive convergence nor
the Abel-tail correlations.
\end{proof}

\begin{firewall}[The weight frame is an assumption to be proved]
\label{fw:v2v59-frame-open}
The estimates
\eqref{eq:v2v59-frame-n}--\eqref{eq:v2v59-frame-prev} require a compatible
V57 Schur bound on every block retained by \(F_{n-1}\), not just on the
active shell.  The weighted sandwich theorem is exact once that frame
exists, but it does not construct the frame for the interacting chiral SM
operator.
\end{firewall}

\subsection{V59 conclusion and next acquisition}
\label{sec:v2v59-conclusion}

V59 replaces the vague requirement of quasilocal running leakage by the
single assembled norm
\(\|\mathsf R_n\mathcal E_n\mathsf R_n\|_{\rm HS}\).  Its weights record
exactly how many shell inverse factors survive each block transition.  The
result yields a sharp Taylor-subtraction threshold and feeds every
strictly decaying leakage row into the V54 chart theorem.  The low--low
block is isolated as the only zone without automatic ultraviolet decay.

V60 is the next mandatory companion audit.  After inspecting the newest NS
and YM notes, it will decide whether to attack the low--low block through a
locality theorem, a Ward cancellation, or a correlated endpoint Abel
estimate.
\clearpage
\section*{Second Volume, Version V60: Companion Audit and Local Extraction of the Low--Low Leakage}
\addcontentsline{toc}{section}{Second Volume, Version V60: Companion Audit and Local Extraction of the Low--Low Leakage}

\subsection*{Purpose and mandatory audit}

V59 isolated the low--low block of the running leakage as the only spectral
zone with no automatic inverse shell factor.  This mandatory audit selects
the locality route for that block.  The conclusion is finite-order: if the
complete low--low shell row has the standard derivative scaling of a
quasilocal amplitude of superficial degree \(\rho\), subtraction through
Taylor degree \(r\) produces the exact gain
\(\Lambda_n^{\rho-r-1}\).  The threshold \(r+1>\rho\) is sharp.

The Navier--Stokes file inspected was
\[
 \texttt{Renewal\_NS\_Stability\_Research\_Notes\_V31.tex},
\]
with \(887537\) bytes and SHA--256 digest
\[
 \texttt{F1121B62238AD5491BB7CF03635EA9934942EDF573ED8FAAA9EE79E1D38A6696}.
\]
Its current endpoint lesson is unchanged: a weaker smoothing norm does not
create the critical first moment, and the missing weighted occupation must
be estimated explicitly.  Accordingly, V60 does not infer low--low decay
from finite rank or from the word quasilocal; it assumes and uses a
specific external-derivative estimate.

The Yang--Mills file inspected was
\[
 \texttt{Renewal\_Yang\_Mills\_Mass\_Gap\_Research\_Notes\_V52.tex},
\]
with \(698394\) bytes and SHA--256 digest
\[
 \texttt{1714D80B4B1BBF099E60FFD3F62BBADD46D713C738A058BBB2C2FD2A8EA4D905}.
\]
YM V52 returns to the exact one-link Wilson density, proves every fixed
Casimir moment by self-adjoint integration by parts, closes the remaining
private terminal degrees through six, and thereby closes global quartic
MTA.  It explicitly leaves order-uniform MTA open.  The transferable rule
is to prove the finite degree actually used.  For the four-dimensional
SM low--low row, the corresponding target is the finite local jet through
degree four, not an all-order derivative atlas.

\begin{assessment}[V60 audit decision]
\label{ass:v2v60-audit}
The low--low block will be split after complete Ward and species assembly
into:
\begin{enumerate}
 \item its finite allowed local Taylor jet, accumulated in the running
       coupling block; and
 \item a Taylor remainder with a proved negative shell degree.
\end{enumerate}
Ward cancellation is used to classify the jet, not assumed to annihilate
the entire low--low amplitude.  Conditional Abel cancellation is reserved
for the exact endpoint where the Taylor remainder does not decay.
\end{assessment}

\subsection{Banach-valued Taylor extraction}
\label{sec:v2v60-Taylor}

Let \(\mathcal Z\) be the physical Banach space of complete shell kernels,
after all tensor, species, ghost, boundary, and contact contributions have
been assembled.  Let \(k\) lie in a ball
\(B_\kappa\subset\mathbb R^m\), with fixed physical
\(\kappa\), and let
\[
 \mathscr L_n:B_\kappa\longrightarrow\mathcal Z
\]
be the complete low--low leakage row at shell \(\Lambda_n\).  Define its
Fr\'echet Taylor polynomial
\begin{equation}
 \mathsf T_r\mathscr L_n(k)
 =
 \sum_{j=0}^r
 \frac1{j!}D^j\mathscr L_n(0)[k^{\otimes j}].
\label{eq:v2v60-Taylor}
\end{equation}

\begin{lemma}[Exact integral remainder]
\label{lem:v2v60-remainder}
If \(\mathscr L_n\in C^{r+1}(B_\kappa;\mathcal Z)\), then
\begin{align}
 \mathscr L_n(k)-\mathsf T_r\mathscr L_n(k)
 =
 \frac1{r!}\int_0^1
 (1-t)^r
 D^{r+1}\mathscr L_n(tk)
 [k^{\otimes(r+1)}]\,\mathrm dt.
\label{eq:v2v60-integral-remainder}
\end{align}
Consequently,
\begin{equation}
 \|\mathscr L_n(k)-\mathsf T_r\mathscr L_n(k)\|_{\mathcal Z}
 \le
 \frac{|k|^{r+1}}{(r+1)!}
 \sup_{|q|\le\kappa}
 \|D^{r+1}\mathscr L_n(q)\|_{\mathcal L^{r+1}}.
\label{eq:v2v60-remainder-bound}
\end{equation}
\end{lemma}

\begin{proof}
Apply the one-variable Banach-space Taylor formula to
\(f(t)=\mathscr L_n(tk)\) between \(t=0\) and \(t=1\).
Its \(j\)-th derivative is
\[
 f^{(j)}(t)=D^j\mathscr L_n(tk)[k^{\otimes j}].
\]
The integral of \((1-t)^r/r!\) is \(1/(r+1)!\), giving the norm bound.
\end{proof}

\subsection{The quasilocal derivative hypothesis}
\label{sec:v2v60-quasilocal}

The quantitative quasilocality statement needed here is: for some
superficial shell degree \(\rho\), logarithmic order \(L\ge0\), and fixed
integer \(r\ge0\),
\begin{equation}
 \sup_{|q|\le\kappa}
 \|D^{r+1}\mathscr L_n(q)\|_{\mathcal L^{r+1}}
 \le
 C\Lambda_n^{\rho-r-1}
 (1+\log\Lambda_n)^L.
\label{eq:v2v60-derivative}
\end{equation}
All norms are taken after complete Ward assembly.

Let \(P_{\rm loc}\) be the fixed bounded BRST-compatible projector onto the
allowed local jets.  We extend it coefficientwise to
\(\mathcal Z\)-valued polynomials in the external momentum.  Assume
\begin{equation}
 P_{\rm loc}\mathsf T_r\mathscr L_n
 =
 \mathsf T_r\mathscr L_n.
\label{eq:v2v60-jet-in-local}
\end{equation}

\begin{theorem}[Low--low locality theorem]
\label{thm:v2v60-locality}
Under \eqref{eq:v2v60-derivative} and
\eqref{eq:v2v60-jet-in-local},
\begin{equation}
 \sup_{|k|\le\kappa}
 \|(I-P_{\rm loc})\mathscr L_n(k)\|_{\mathcal Z}
 \le
 C_r\kappa^{r+1}
 \Lambda_n^{\rho-r-1}
 (1+\log\Lambda_n)^L.
\label{eq:v2v60-locality}
\end{equation}
For dyadic \(\Lambda_n\simeq2^n\), the projected low--low shell row is
absolutely summable whenever
\begin{equation}
 r+1>\rho.
\label{eq:v2v60-threshold}
\end{equation}
\end{theorem}

\begin{proof}
Equation~\eqref{eq:v2v60-jet-in-local} gives
\[
 (I-P_{\rm loc})\mathscr L_n
 =
 (I-P_{\rm loc})
 (\mathscr L_n-\mathsf T_r\mathscr L_n).
\]
Apply Lemma~\ref{lem:v2v60-remainder} and
\eqref{eq:v2v60-derivative}.  This proves
\eqref{eq:v2v60-locality}, with the operator norm of
\(I-P_{\rm loc}\) absorbed into \(C_r\).  If
\(\varepsilon=r+1-\rho>0\), then
\[
 \Lambda_n^{-\varepsilon}(1+\log\Lambda_n)^L
 \asymp
 2^{-\varepsilon n}(1+n)^L,
\]
which is summable.
\end{proof}

\begin{corollary}[The four-dimensional finite-order target]
\label{cor:v2v60-four-dimensional}
If the complete low--low row has superficial degree
\(\rho\le4\), and
\eqref{eq:v2v60-derivative} holds through derivative order five, then
subtraction through \(r=4\) gives
\begin{equation}
 \|(I-P_{\rm loc})\mathscr L_n\|_{\mathcal Z}
 \le
 C_\kappa\Lambda_n^{-1}
 (1+\log\Lambda_n)^L,
\label{eq:v2v60-four-dimensional}
\end{equation}
which is dyadically summable.
\end{corollary}

\begin{proof}
Set \(r=4\) in
Theorem~\ref{thm:v2v60-locality}.  The worst case is \(\rho=4\).
\end{proof}

\begin{firewall}[Power counting is not the derivative estimate]
\label{fw:v2v60-power-count}
Renormalizability predicts \(\rho\le4\) for the complete local
four-dimensional row, but that statement alone does not prove
\eqref{eq:v2v60-derivative} uniformly in the regulator and background.
The fifth external derivative must be estimated in the physical Banach
norm after the interacting interpolation and Ward completion.  This is the
analytic acquisition required by
Corollary~\ref{cor:v2v60-four-dimensional}.
\end{firewall}

\subsection{Sharpness at the Taylor endpoint}
\label{sec:v2v60-sharp}

\begin{proposition}[The threshold cannot be weakened by Taylor's theorem]
\label{prop:v2v60-sharp}
Fix a unit vector \(e\in\mathcal Z\), one scalar external variable \(k\),
and put
\begin{equation}
 \mathscr L_n(k)
 =
 \Lambda_n^\rho
 \left(\frac{k}{\Lambda_n}\right)^{r+1}e.
\label{eq:v2v60-sharp-family}
\end{equation}
Then \(\mathsf T_r\mathscr L_n=0\),
\eqref{eq:v2v60-derivative} holds with \(L=0\), and
\begin{equation}
 \|\mathscr L_n(k)\|_{\mathcal Z}
 =
 |k|^{r+1}\Lambda_n^{\rho-r-1}.
\label{eq:v2v60-sharp-size}
\end{equation}
At the endpoint \(r+1=\rho\), the row is independent of \(n\) for fixed
nonzero \(k\), so its dyadic series diverges.
\end{proposition}

\begin{proof}
The function is a monomial of degree \(r+1\), so every derivative through
degree \(r\) at zero vanishes.  Its \((r+1)\)-st derivative is
\((r+1)!\Lambda_n^{\rho-r-1}e\), which gives the claimed derivative
bound and exact remainder.  The endpoint conclusion is immediate.
\end{proof}

\begin{assessment}[What endpoint cancellation would have to prove]
\label{ass:v2v60-endpoint}
When \(r+1=\rho\), Taylor extraction alone gives an \(O(1)\) shell row.
The V55--V56 conditional route would then require convergence of the
complete signed additive series and both ordered tail correlations.
Neither follows from \eqref{eq:v2v60-derivative}.  Thus an endpoint use of
the Abel compiler is a separate theorem, not an alternative reading of
Taylor's bound.
\end{assessment}

\subsection{Ward compatibility of the local jet}
\label{sec:v2v60-Ward}

Let \(s\) be the bounded BRST differential on the common algebra, extended
continuously to \(\mathcal Z\).  Suppose it commutes with external
Fr\'echet differentiation.

\begin{lemma}[Taylor extraction preserves BRST closure]
\label{lem:v2v60-BRST}
If
\[
 s\mathscr L_n(k)=0
 \qquad(k\in B_\kappa),
\]
then
\begin{equation}
 s\,\mathsf T_r\mathscr L_n=0,
 \qquad
 s(\mathscr L_n-\mathsf T_r\mathscr L_n)=0.
\label{eq:v2v60-BRST}
\end{equation}
\end{lemma}

\begin{proof}
For every \(j\),
\[
 sD^j\mathscr L_n(0)
 =
 D^j(s\mathscr L_n)(0)=0.
\]
Apply \(s\) term by term in
\eqref{eq:v2v60-Taylor}; subtracting gives the second identity.
\end{proof}

\begin{firewall}[A closed jet need not be an allowed invariant jet]
\label{fw:v2v60-cohomology}
Lemma~\ref{lem:v2v60-BRST} proves cocycle closure, not
\eqref{eq:v2v60-jet-in-local}.  The latter requires the local
Slavnov-cohomology classification, anomaly cancellation with a
regulator-compatible primitive, and the chosen physical normalization
conditions.  A componentwise momentum Taylor projector need not preserve
the tensor Ward identity; the full invariant projector must be applied
after assembly.
\end{firewall}

\subsection{Interface with the V59 weighted leakage}
\label{sec:v2v60-V59}

Let \(P_L\) be the fixed physical low projection in V59.  Form the
complete low--low leakage row before projection:
\begin{equation}
 \mathscr L_n(k)
 =
 \operatorname{Tr}\!\left[
 F_{n-1}C_n
 P_L\mathcal E_n(k)P_L
 C_{n-1}F_{n-1}
 \,\mathcal V_n(k)
 \right]_{\rm conn,Ward},
\label{eq:v2v60-low-row}
\end{equation}
where \(\mathcal V_n\) denotes the remaining gauge commutator and composite
insertions and the subscript requires the connected physical assembly.

\begin{theorem}[Strict local remainder closes the low--low branch]
\label{thm:v2v60-low-close}
Suppose the row \eqref{eq:v2v60-low-row} satisfies
\eqref{eq:v2v60-derivative} with \(r+1>\rho\), and its Taylor jet satisfies
\eqref{eq:v2v60-jet-in-local}.  Then its nonlocal shell derivations obey
\[
 \sum_n
 \|(I-P_{\rm loc})\mathscr L_n\|_{\mathcal Z}<\infty.
\]
Consequently the low--low leakage contributes to the absolute V54 branch
of the common algebra chart.
\end{theorem}

\begin{proof}
Apply Theorem~\ref{thm:v2v60-locality} and sum the dyadic bound.  The
resulting shell derivation series is absolutely convergent, so
Theorem~\ref{thm:v2v54-product} applies to this branch.
\end{proof}

\begin{firewall}[The trace notation does not separate species]
\label{fw:v2v60-complete-row}
The bracket in \eqref{eq:v2v60-low-row} is one complete connected
Ward row.  The derivative estimate may be proved by expanding it, but
cancelling orientations, chiral species, ghosts, and contact terms must be
reassembled before the \(\mathcal Z\)-norm.  The theorem does not authorize
separate positive bounds followed by a product of cancellations.
\end{firewall}

\subsection{V60 conclusion and next acquisition}
\label{sec:v2v60-conclusion}

V60 proves that the low--low leakage is harmless after a finite local
extraction if its complete fifth external derivative has the expected
four-dimensional shell gain.  The threshold is sharp, logarithmic powers
do not affect strict dyadic summability, and BRST closure passes to the
Taylor jet.  The local cohomology classification and the uniform derivative
bound remain explicit inputs.

The next step is to prove the derivative estimate in a controlled model
rather than assume it.  V61 will differentiate a low-external-momentum
resolvent word and show that each derivative inserts one additional
transition-shell inverse.  It will identify the commutator and domain
hypotheses needed to obtain
\(\|D^5\mathscr L_n\|\lesssim\Lambda_n^{-1}\) for a degree-four row, and
will state precisely which interacting vertices remain outside that
regression.
\clearpage
\section*{Second Volume, Version V61: Differentiated Resolvent Words and the Fifth-Derivative Regression}
\addcontentsline{toc}{section}{Second Volume, Version V61: Differentiated Resolvent Words and the Fifth-Derivative Regression}

\subsection*{Purpose}

V60 reduced the nonlocal low--low leakage to one finite analytic target:
for an assembled row of superficial degree at most four, prove that five
external-momentum derivatives cost five powers of the transition scale.
V9 supplied only a first adjacent resolvent identity, while V57 localized
an undifferentiated noncommuting resolvent column.  Neither result by
itself controls repeated external differentiation.

This note proves the exact higher derivative formula for a compressed
resolvent, establishes the expected degree reduction for a finite
transition-supported resolvent word, and exhibits the precise obstruction
when a derivative leaves the transition shell.  It also gives an equivalent
Schur-complement criterion that will be used to attack that obstruction
upstream.

\subsection{Ordered partitions and compressed inverse differentiation}
\label{sec:v2v61-inverse}

For \(m\ge1\), let \(\operatorname{OP}(m)\) denote the set of ordered
partitions
\[
 \pi=(B_1,\ldots,B_j)
\]
of \(\{1,\ldots,m\}\) into nonempty disjoint blocks.  The order of the
blocks matters.  Within a block, arguments are inserted in increasing
index order; symmetry of a Fr\'echet derivative makes this convention
immaterial.

Let \(\mathcal H\) be a Hilbert space, let \(P=P^*=P^2\), let
\(B_\kappa\subset\mathbb R^q\) be a ball, and let
\[
 A:B_\kappa\longrightarrow\mathcal L(\mathcal H)
\]
be \(C^M\), with every \(A(k)\) invertible.  Write
\[
 C(k)=A(k)^{-1},
 \qquad
 G(k)=PC(k)P|_{P\mathcal H}.
\]
For directions \(h_i\in\mathbb R^q\) and a block
\(B=\{i_1,\ldots,i_a\}\), put
\[
 D_BA(k)=D^aA(k)[h_{i_1},\ldots,h_{i_a}].
\]

\begin{lemma}[Exact compressed inverse derivative under shell incidence]
\label{lem:v2v61-inverse-formula}
Suppose that, for \(1\le a\le M\),
\begin{equation}
 D^aA(k)[h_1,\ldots,h_a]
 =
 P D^aA(k)[h_1,\ldots,h_a]P.
\label{eq:v2v61-shell-incidence}
\end{equation}
Then, for \(1\le m\le M\),
\begin{equation}
 D^mG(k)[h_1,\ldots,h_m]
 =
 \sum_{\substack{\pi=(B_1,\ldots,B_j)\\
                  \in\operatorname{OP}(m)}}
 (-1)^j
 G(k)D_{B_1}A(k)G(k)\cdots
 D_{B_j}A(k)G(k).
\label{eq:v2v61-inverse-formula}
\end{equation}
\end{lemma}

\begin{proof}
Differentiate \(A(k)C(k)=I\) once to obtain
\[
 DC(k)[h]=-C(k)DA(k)[h]C(k).
\]
Repeated differentiation gives the same formula with \(C\) in place of
\(G\) and a sum over ordered partitions.  This can be proved by induction:
when the new direction differentiates a resolvent factor, it creates a new
singleton block in one of the ordered gaps and changes the sign; when it
differentiates \(D_BA\), it adjoins the new index to that block without
changing the number of blocks.  Every ordered partition of
\(\{1,\ldots,m+1\}\) is obtained exactly once by one of these two
operations.

Multiply the full formula by \(P\) on both sides.  Condition
\eqref{eq:v2v61-shell-incidence} inserts \(P\) on both sides of every
factor \(D_BA\).  Hence each intervening \(PCP\) is \(G\), which proves
\eqref{eq:v2v61-inverse-formula}.
\end{proof}

\begin{theorem}[One inverse scale per external derivative]
\label{thm:v2v61-inverse-degree}
Let \(\Lambda\ge1\).  Under the hypotheses of
Lemma~\ref{lem:v2v61-inverse-formula}, suppose
\begin{align}
 \sup_{|k|\le\kappa}\|G(k)\|
 &\le g\Lambda^{-1},
\label{eq:v2v61-G0}\\
 \sup_{|k|\le\kappa}\|D^aA(k)\|_{\mathcal L^a}
 &\le a_a\Lambda^{1-a},
 \qquad 1\le a\le M.
\label{eq:v2v61-A-degree}
\end{align}
Then, for \(0\le m\le M\),
\begin{equation}
 \sup_{|k|\le\kappa}
 \|D^mG(k)\|_{\mathcal L^m}
 \le g_m\Lambda^{-m-1},
\label{eq:v2v61-G-degree}
\end{equation}
where \(g_m\) depends only on \(m,g,a_1,\ldots,a_m\), and not on
\(\Lambda\).
\end{theorem}

\begin{proof}
The case \(m=0\) is \eqref{eq:v2v61-G0}.  In a term of
\eqref{eq:v2v61-inverse-formula} with \(j\) blocks, there are \(j+1\)
factors \(G\).  If the block sizes are
\(b_1,\ldots,b_j\), then \(\sum_\ell b_\ell=m\), and the scale exponent
is
\[
 -(j+1)+\sum_{\ell=1}^j(1-b_\ell)
 =-m-1.
\]
The number of ordered partitions is finite for fixed \(m\).  Summing their
coefficient bounds gives \(g_m\).
\end{proof}

\begin{remark}[Affine first-order symbols]
\label{rem:v2v61-affine}
If \(A(k)\) is affine in \(k\), only singleton blocks survive in
\eqref{eq:v2v61-inverse-formula}.  There are \(m!\) ordered singleton
partitions, and every derivative inserts one additional compressed
resolvent.  The factorial changes the fixed-order constant but not the
power \(\Lambda^{-m-1}\).
\end{remark}

\subsection{Finite resolvent words}
\label{sec:v2v61-words}

The transition projection below is a loop-momentum projection \(P_n\).
It is distinct from the fixed physical low projection \(P_L\) that
classifies the leakage insertion in V59--V60.  A low--low insertion with
respect to \(P_L\) may still occur inside a loop word supported at
transition scale \(\Lambda_n\).

For each \(n\), let \(G_n=P_nA_n^{-1}P_n\) satisfy
\eqref{eq:v2v61-G-degree}.  Let
\[
 B_{\ell,n}:B_\kappa\longrightarrow
 \mathcal L(P_n\mathcal H),
 \qquad 1\le\ell\le s,
\]
be \(C^M\) vertex families with degrees \(\beta_\ell\):
\begin{equation}
 \sup_{|k|\le\kappa}
 \|D^aB_{\ell,n}(k)\|_{\mathcal L^a}
 \le b_{\ell,a}\Lambda_n^{\beta_\ell-a},
 \qquad 0\le a\le M.
\label{eq:v2v61-B-degree}
\end{equation}
Let \(\mathcal Z\) be a Banach space.  To include an ordinary trace,
finite tensor contraction, and a finite already-assembled species sum,
assume a fixed \(2s\)-linear contraction \(\tau_n\) obeys
\begin{equation}
 \|\tau_n(X_1,\ldots,X_{2s})\|_{\mathcal Z}
 \le w\Lambda_n^d\prod_{i=1}^{2s}\|X_i\|.
\label{eq:v2v61-trace-bound}
\end{equation}
Define the complete word
\begin{equation}
 \mathscr W_n(k)
 =
 \tau_n\bigl(
 B_{1,n}(k),G_n(k),\ldots,
 B_{s,n}(k),G_n(k)
 \bigr).
\label{eq:v2v61-word}
\end{equation}

\begin{theorem}[External degree of a transition-supported word]
\label{thm:v2v61-word-degree}
Under
\eqref{eq:v2v61-G-degree}--\eqref{eq:v2v61-trace-bound}, put
\begin{equation}
 \rho=d+\sum_{\ell=1}^s\beta_\ell-s.
\label{eq:v2v61-rho}
\end{equation}
For every \(0\le m\le M\),
\begin{equation}
 \sup_{|k|\le\kappa}
 \|D^m\mathscr W_n(k)\|_{\mathcal L^m(\mathbb R^q;\mathcal Z)}
 \le C_m\Lambda_n^{\rho-m}.
\label{eq:v2v61-word-degree}
\end{equation}
The constant is uniform in \(n\).
\end{theorem}

\begin{proof}
Apply the Fr\'echet product rule to
\eqref{eq:v2v61-word}.  A typical term assigns derivative orders
\[
 u_1,v_1,\ldots,u_s,v_s\ge0,
 \qquad
 \sum_{\ell=1}^s(u_\ell+v_\ell)=m,
\]
to the vertex and resolvent factors.  Its scale exponent, using
\eqref{eq:v2v61-B-degree},
\eqref{eq:v2v61-G-degree}, and
\eqref{eq:v2v61-trace-bound}, is
\begin{align*}
 d+\sum_{\ell=1}^s
 \bigl((\beta_\ell-u_\ell)+(-1-v_\ell)\bigr)
 &=
 d+\sum_{\ell=1}^s\beta_\ell-s-m\\
 &=\rho-m.
\end{align*}
The multinomial coefficients and the number of derivative assignments are
finite for fixed \(m\) and \(s\).  Their sum is absorbed into \(C_m\).
\end{proof}

\begin{corollary}[The controlled fifth-derivative acquisition]
\label{cor:v2v61-fifth}
If a complete finite sum of words
\eqref{eq:v2v61-word} has \(\rho\le4\), satisfies the hypotheses with
\(M\ge5\), and is assembled before the
\(\mathcal Z\)-norm, then
\begin{equation}
 \sup_{|k|\le\kappa}
 \|D^5\mathscr L_n(k)\|_{\mathcal L^5(\mathbb R^q;\mathcal Z)}
 \le C\Lambda_n^{-1}.
\label{eq:v2v61-fifth}
\end{equation}
Consequently its degree-four Taylor remainder satisfies the summable
low--low estimate of
Corollary~\ref{cor:v2v60-four-dimensional}.
\end{corollary}

\begin{proof}
Apply Theorem~\ref{thm:v2v61-word-degree} to each word and use
\(\rho-5\le-1\).  The finite complete sum preserves the estimate.  The last
claim is exactly the V60 integral-remainder theorem.
\end{proof}

\begin{firewall}[The word theorem is a controlled regression]
\label{fw:v2v61-regression}
Theorem~\ref{thm:v2v61-word-degree} proves the desired derivative count
when differentiated kinetic and vertex factors obey the stated
transition-supported symbol bounds.  It does not prove those bounds for
the full curved, running, gauge-fixed Standard Model operator.  In
particular, a derivative of the full operator may leave \(P_n\mathcal H\),
and a differentiated vertex may contain a composite contact distribution
whose extension must first be included in the complete Ward row.
\end{firewall}

\subsection{Why undifferentiated localization is insufficient}
\label{sec:v2v61-counterexample}

\begin{counterexample}[An order-zero derivative can make a low excursion]
\label{sep:v2v61-low-excursion}
Let \(\mathcal H=\mathbb C^2\),
\(P=\operatorname{diag}(1,0)\), and, for \(|k|\le\kappa\), set
\begin{equation}
 A_\Lambda(k)
 =
 \begin{pmatrix}
  \Lambda & k\\
  k&1
 \end{pmatrix},
 \qquad
 \Lambda>\kappa^2.
\label{eq:v2v61-counter-A}
\end{equation}
Then \(A_\Lambda(k)\) is positive and invertible,
\(\|A_\Lambda'(k)\|=1\), and \(A_\Lambda''(k)=0\).  Nevertheless,
\begin{equation}
 G_\Lambda(k)
 =
 PA_\Lambda(k)^{-1}P
 =
 \frac1{\Lambda-k^2}P,
\label{eq:v2v61-counter-G}
\end{equation}
so
\begin{equation}
 \|G_\Lambda(0)\|=\Lambda^{-1},
 \qquad
 \|G_\Lambda''(0)\|=2\Lambda^{-2}.
\label{eq:v2v61-counter-loss}
\end{equation}
The expected two-derivative bound \(O(\Lambda^{-3})\) fails by one
inverse power.
\end{counterexample}

\begin{proof}
The determinant is \(\Lambda-k^2>0\), and the displayed inverse follows
from the \(2\)-by-\(2\) inverse formula.  Differentiating
\((\Lambda-k^2)^{-1}\) twice at zero gives \(2\Lambda^{-2}\).
The first derivative of \(A_\Lambda\) is the off-diagonal flip of norm one.
Thus the failure is not caused by a large differentiated operator.
It is caused by the itinerary
\[
 P A_\Lambda(0)^{-1}P\,
 P A_\Lambda'(0)Q\,
 Q A_\Lambda(0)^{-1}Q\,
 Q A_\Lambda'(0)P\,
 P A_\Lambda(0)^{-1}P,
\]
whose complementary inverse is \(O(1)\), rather than
\(O(\Lambda^{-1})\).
\end{proof}

\begin{assessment}[New content relative to V57]
\label{ass:v2v61-relative-V57}
V57 controls \(A^{-1}P\) at zero derivative through a Schur feedback norm.
Counterexample~\ref{sep:v2v61-low-excursion} satisfies the desired
undifferentiated compressed estimate, but its second derivative loses a
power.  Therefore V57 cannot simply be differentiated without estimates
on the differentiated feedback paths.
\end{assessment}

\subsection{The effective shell operator}
\label{sec:v2v61-effective}

Let \(Q=I-P\).  Suppose \(QA(k)Q\) is invertible and define
\begin{equation}
 \mathscr S(k)
 =
 PA(k)P
 -
 PA(k)Q\,[QA(k)Q]^{-1}\,QA(k)P.
\label{eq:v2v61-Schur}
\end{equation}
Block Gaussian elimination gives the same algebraic Schur identity as
Lemma~\ref{lem:v2v57-column}; self-adjointness is not needed for this
identity once the displayed block inverses exist.
\begin{equation}
 G(k)=PA(k)^{-1}P=\mathscr S(k)^{-1}.
\label{eq:v2v61-G-Schur}
\end{equation}

\begin{theorem}[Effective-shell derivative criterion]
\label{thm:v2v61-effective}
Suppose \(\mathscr S:B_\kappa\to\mathcal L(P\mathcal H)\) is \(C^M\) and
\begin{align}
 \sup_{|k|\le\kappa}\|\mathscr S(k)^{-1}\|
 &\le g\Lambda^{-1},
\label{eq:v2v61-S-inverse}\\
 \sup_{|k|\le\kappa}\|D^a\mathscr S(k)\|_{\mathcal L^a}
 &\le s_a\Lambda^{1-a},
 \qquad 1\le a\le M.
\label{eq:v2v61-S-degree}
\end{align}
Then \(G=\mathscr S^{-1}\) satisfies
\eqref{eq:v2v61-G-degree}, without the shell-incidence assumption
\eqref{eq:v2v61-shell-incidence}.
\end{theorem}

\begin{proof}
Apply the full-space inverse derivative formula from the proof of
Lemma~\ref{lem:v2v61-inverse-formula} to \(\mathscr S\) on
\(P\mathcal H\).  The exponent calculation in
Theorem~\ref{thm:v2v61-inverse-degree} is unchanged.
\end{proof}

\begin{remark}[Diagnosis of the counterexample]
\label{rem:v2v61-counter-Schur}
For \eqref{eq:v2v61-counter-A},
\[
 \mathscr S_\Lambda(k)=\Lambda-k^2.
\]
Thus \(D^2\mathscr S_\Lambda=-2\), whereas
\eqref{eq:v2v61-S-degree} requires \(O(\Lambda^{-1})\) at derivative order
two.  The Schur criterion detects exactly the missing inverse scale in the
low excursion.
\end{remark}

\subsection{V61 conclusion and next acquisition}
\label{sec:v2v61-conclusion}

V61 proves the fifth-derivative estimate for every complete finite
transition-supported resolvent-word regression of superficial degree at
most four.  It also proves that undifferentiated shell localization does
not suffice: external derivatives can travel through a low complementary
inverse and lose the scale needed by V60.

The remaining acquisition is now sharper than a generic quasilocality
claim.  One must prove
\[
 \|D^a\mathscr S_n\|\lesssim\Lambda_n^{1-a},
 \qquad 1\le a\le5,
\]
for the complete effective transition-shell operator, or prove an
equivalent weighted estimate for every differentiated excursion in its
Schur self-energy.  V62 will differentiate the Schur self-energy, classify
its ordered shell--complement itineraries, and identify a weighted
off-diagonal hypothesis that supplies the missing inverse power without
assuming exact shell incidence.
\clearpage
\section*{Second Volume, Version V62: Weighted Schur Excursions and the Differentiated Feedback Budget}
\addcontentsline{toc}{section}{Second Volume, Version V62: Weighted Schur Excursions and the Differentiated Feedback Budget}

\subsection*{Purpose}

V61 showed that an undifferentiated shell inverse can retain its
\(\Lambda_n^{-1}\) scale while its second external derivative loses a
power through a low complementary excursion.  This note differentiates
that excursion exactly.  The result is a scale budget for the Schur
self-energy
\[
 \Sigma_n=X_nH_n^{-1}Y_n.
\]
If the two crossings have differentiated degrees \(\alpha\) and \(\beta\),
and the complementary inverse has differentiated degree \(\gamma\), the
effective shell operator has the required first-order symbol bounds
precisely under the sufficient budget
\[
 \alpha+\beta+\gamma\le1.
\]
For a self-adjoint low excursion, \(\gamma=0\), so each crossing may have
degree at most \(1/2\).  For a complementary shell inverse,
\(\gamma=-1\), each crossing may have degree as large as one.

\subsection{Differentiated inverse on the complementary block}
\label{sec:v2v62-complement}

Let \(B_\kappa\subset\mathbb R^q\), let
\[
 H_n:B_\kappa\longrightarrow\mathcal L(\mathcal K_n)
\]
be \(C^M\), and suppose every \(H_n(k)\) is invertible.  Put
\(R_n(k)=H_n(k)^{-1}\).

\begin{lemma}[Scale-dual inverse differentiation]
\label{lem:v2v62-complement-inverse}
Fix \(\gamma\in\mathbb R\).  Suppose, uniformly for
\(|k|\le\kappa\),
\begin{align}
 \|R_n(k)\|
 &\le r_0\Lambda_n^\gamma,
\label{eq:v2v62-R0}\\
 \|D^aH_n(k)\|_{\mathcal L^a}
 &\le h_a\Lambda_n^{-\gamma-a},
 \qquad 1\le a\le M.
\label{eq:v2v62-H-degree}
\end{align}
Then, for \(0\le m\le M\),
\begin{equation}
 \|D^mR_n(k)\|_{\mathcal L^m}
 \le r_m\Lambda_n^{\gamma-m},
\label{eq:v2v62-R-degree}
\end{equation}
with constants independent of \(n\).
\end{lemma}

\begin{proof}
The case \(m=0\) is
\eqref{eq:v2v62-R0}.  Apply the ordered-partition inverse formula proved
in Lemma~\ref{lem:v2v61-inverse-formula}, now with the identity projection.
A term indexed by \(j\) blocks of total size \(m\) contains \(j+1\)
factors \(R_n\) and \(j\) differentiated factors \(H_n\).  Its exponent is
\[
 (j+1)\gamma+
 \sum_{\ell=1}^j(-\gamma-|B_\ell|)
 =\gamma-m.
\]
There are finitely many ordered partitions at fixed \(m\), so their
constant bounds may be summed.
\end{proof}

\begin{assessment}[Meaning of the complement exponent]
\label{ass:v2v62-gamma}
The exponent \(\gamma\) measures the complementary inverse, not the
complementary kinetic operator.  A low block with a uniformly bounded
inverse has \(\gamma=0\).  A transition or high block with one inverse
frequency has \(\gamma=-1\).  Equation
\eqref{eq:v2v62-H-degree} is an actual differentiated estimate: a global
bound on \(H_n\) does not imply it.
\end{assessment}

\subsection{Exact differentiation of one Schur excursion}
\label{sec:v2v62-excursion}

Let
\begin{align*}
 X_n(k)&:\mathcal K_n\longrightarrow\mathcal P_n,\\
 Y_n(k)&:\mathcal P_n\longrightarrow\mathcal K_n
\end{align*}
be \(C^M\), and define
\begin{equation}
 \Sigma_n(k)=X_n(k)R_n(k)Y_n(k).
\label{eq:v2v62-Sigma}
\end{equation}
For a map
\(\varepsilon:\{1,\ldots,m\}\to\{X,R,Y\}\), let
\(I_\varepsilon,J_\varepsilon,K_\varepsilon\) be the inverse images of
\(X,R,Y\), respectively.  A derivative indexed by an empty set means the
undifferentiated factor.

\begin{lemma}[Exact three-factor assignment formula]
\label{lem:v2v62-assignment}
For directions \(u_1,\ldots,u_m\in\mathbb R^q\),
\begin{equation}
 D^m\Sigma_n(k)[u_1,\ldots,u_m]
 =
 \sum_{\varepsilon:\{1,\ldots,m\}\to\{X,R,Y\}}
 D_{I_\varepsilon}X_n(k)\,
 D_{J_\varepsilon}R_n(k)\,
 D_{K_\varepsilon}Y_n(k).
\label{eq:v2v62-assignment}
\end{equation}
In particular, the sum contains exactly \(3^m\) terms before identical
symmetric terms are collected.
\end{lemma}

\begin{proof}
Apply the multilinear Leibniz rule.  Each labelled direction acts on
exactly one of the three ordered factors.  Thus an assignment
\(\varepsilon\) records one and only one term.  Conversely every Leibniz
term determines such an assignment.
\end{proof}

\begin{theorem}[Weighted excursion budget]
\label{thm:v2v62-budget}
Suppose
\eqref{eq:v2v62-R-degree} holds and, for \(0\le a\le M\),
\begin{align}
 \|D^aX_n(k)\|_{\mathcal L^a}
 &\le x_a\Lambda_n^{\alpha-a},
\label{eq:v2v62-X-degree}\\
 \|D^aY_n(k)\|_{\mathcal L^a}
 &\le y_a\Lambda_n^{\beta-a}.
\label{eq:v2v62-Y-degree}
\end{align}
Then
\begin{equation}
 \|D^m\Sigma_n(k)\|_{\mathcal L^m}
 \le C_m\Lambda_n^{\alpha+\beta+\gamma-m},
 \qquad 0\le m\le M.
\label{eq:v2v62-Sigma-degree}
\end{equation}
Consequently,
\begin{equation}
 \alpha+\beta+\gamma\le1
\label{eq:v2v62-budget}
\end{equation}
implies the first-order effective-symbol estimate
\begin{equation}
 \|D^m\Sigma_n(k)\|_{\mathcal L^m}
 \le C_m\Lambda_n^{1-m}.
\label{eq:v2v62-Sigma-first-order}
\end{equation}
\end{theorem}

\begin{proof}
For a term in
\eqref{eq:v2v62-assignment}, put
\[
 i=|I_\varepsilon|,
 \qquad j=|J_\varepsilon|,
 \qquad \ell=|K_\varepsilon|.
\]
Then \(i+j+\ell=m\).  Its exponent is
\[
 (\alpha-i)+(\gamma-j)+(\beta-\ell)
 =\alpha+\beta+\gamma-m.
\]
Sum the \(3^m\) uniform bounds.  If
\eqref{eq:v2v62-budget} holds and \(\Lambda_n\ge1\), the last exponent is
at most \(1-m\).
\end{proof}

\begin{proposition}[Sharpness of the exponent budget]
\label{prop:v2v62-sharp}
Let \(f\in C^M([-\kappa,\kappa])\), and suppose
\((f^2)^{(m_0)}(0)\ne0\) for some \(m_0\le M\).  In one-dimensional
blocks define
\[
 X_n(k)=\Lambda_n^\alpha f(k/\Lambda_n),
 \quad
 R_n(k)=\Lambda_n^\gamma,
 \quad
 Y_n(k)=\Lambda_n^\beta f(k/\Lambda_n).
\]
Then
\begin{equation}
 \Sigma_n^{(m_0)}(0)
 =
 (f^2)^{(m_0)}(0)
 \Lambda_n^{\alpha+\beta+\gamma-m_0}.
\label{eq:v2v62-sharp}
\end{equation}
Thus the factor hypotheses alone cannot yield
\(O(\Lambda_n^{1-m_0})\) when
\(\alpha+\beta+\gamma>1\).  Equality in
\eqref{eq:v2v62-budget} gives exactly the endpoint first-order degree.
\end{proposition}

\begin{proof}
The self-energy is
\[
 \Sigma_n(k)
 =
 \Lambda_n^{\alpha+\beta+\gamma}
 f(k/\Lambda_n)^2.
\]
Differentiate \(m_0\) times and set \(k=0\).
\end{proof}

\subsection{Closure of the effective shell derivative}
\label{sec:v2v62-effective}

Write the block operator as
\[
 A_n(k)
 =
 \begin{pmatrix}
  K_n(k)&X_n(k)\\
  Y_n(k)&H_n(k)
 \end{pmatrix}
\]
on \(\mathcal P_n\oplus\mathcal K_n\), and let
\begin{equation}
 \mathscr S_n(k)=K_n(k)-\Sigma_n(k).
\label{eq:v2v62-effective}
\end{equation}

\begin{theorem}[Weighted feedback closes the differentiated resolvent]
\label{thm:v2v62-close}
Assume the hypotheses of
Theorem~\ref{thm:v2v62-budget} through order \(M\), assume
\(\alpha+\beta+\gamma\le1\), and suppose
\begin{align}
 \|D^aK_n(k)\|_{\mathcal L^a}
 &\le k_a\Lambda_n^{1-a},
 \qquad 0\le a\le M,
\label{eq:v2v62-K-degree}\\
 \|\mathscr S_n(k)^{-1}\|
 &\le g\Lambda_n^{-1}.
\label{eq:v2v62-S-inverse}
\end{align}
Then
\begin{align}
 \|D^a\mathscr S_n(k)\|_{\mathcal L^a}
 &\le s_a\Lambda_n^{1-a},
 \qquad 0\le a\le M,
\label{eq:v2v62-S-degree}\\
 \|D^a(P_nA_n(k)^{-1}P_n)\|_{\mathcal L^a}
 &\le g_a\Lambda_n^{-1-a},
 \qquad 0\le a\le M.
\label{eq:v2v62-G-close}
\end{align}
\end{theorem}

\begin{proof}
Equation
\eqref{eq:v2v62-S-degree} follows from
\eqref{eq:v2v62-effective},
\eqref{eq:v2v62-K-degree}, and
Theorem~\ref{thm:v2v62-budget}.  Block Gaussian elimination gives
\[
 P_nA_n(k)^{-1}P_n=\mathscr S_n(k)^{-1}.
\]
Apply the effective-shell derivative criterion
Theorem~\ref{thm:v2v61-effective} to
\eqref{eq:v2v62-S-degree} and
\eqref{eq:v2v62-S-inverse}.
\end{proof}

\begin{corollary}[Passage to the V60 low--low remainder]
\label{cor:v2v62-V60}
Let a complete degree-four shell row be a finite sum of the resolvent words
of V61.  Suppose every differentiated compressed resolvent in those words
is controlled by
Theorem~\ref{thm:v2v62-close} with \(M\ge5\), and every differentiated
vertex obeys
\eqref{eq:v2v61-B-degree}.  Then
\[
 \|D^5\mathscr L_n\|\le C\Lambda_n^{-1}.
\]
After the fixed BRST-compatible Taylor projection through degree four,
the nonlocal low--low row is absolutely summable.
\end{corollary}

\begin{proof}
Theorem~\ref{thm:v2v62-close} supplies the resolvent hypothesis of
Theorem~\ref{thm:v2v61-word-degree}.  Apply
Corollary~\ref{cor:v2v61-fifth} and then
Corollary~\ref{cor:v2v60-four-dimensional}.
\end{proof}

\subsection{Spectral-zone interpretation}
\label{sec:v2v62-zones}

Suppose \(A_n(k)\) is self-adjoint, so \(Y_n=X_n^*\), and the differentiated
bounds for the two crossings have the same exponent
\(\alpha=\beta\).  The budget becomes
\begin{equation}
 2\alpha+\gamma\le1.
\label{eq:v2v62-selfadjoint-budget}
\end{equation}
The principal cases are
\[
\begin{array}{c|c|c}
 \text{complementary inverse}&\gamma&
 \text{largest allowed crossing degree }\alpha\\ \hline
 \text{uniform low inverse}&0&1/2\\
 \text{transition or high inverse}&-1&1
\end{array}
\]
Thus a low excursion needs a genuine half-order improvement on each
crossing.  An excursion that already carries a complementary shell inverse
can tolerate first-order crossings.

More generally, let
\[
 Q_n=\bigoplus_{z=1}^NQ_{n,z}
\]
reduce \(H_n\), so that \(R_n=\bigoplus_zR_{n,z}\), and put
\(X_{n,z}=P_nA_nQ_{n,z}\),
\(Y_{n,z}=Q_{n,z}A_nP_n\).  Then
\begin{equation}
 \Sigma_n=\sum_{z=1}^N X_{n,z}R_{n,z}Y_{n,z}.
\label{eq:v2v62-zone-sum}
\end{equation}

\begin{corollary}[Finite spectral-zone budget]
\label{cor:v2v62-zones}
If each zone satisfies the differentiated hypotheses with exponents
\(\alpha_z,\beta_z,\gamma_z\) and
\[
 \alpha_z+\beta_z+\gamma_z\le1,
\]
then
\[
 \|D^m\Sigma_n\|\le C_m\Lambda_n^{1-m}
 \qquad(0\le m\le M).
\]
\end{corollary}

\begin{proof}
Apply Theorem~\ref{thm:v2v62-budget} to each term of
\eqref{eq:v2v62-zone-sum} and sum the fixed finite number of zones.
\end{proof}

\begin{firewall}[Reduction of the complementary block is substantive]
\label{fw:v2v62-reduction}
Corollary~\ref{cor:v2v62-zones} assumes that the chosen spectral zones
reduce \(H_n(k)\).  If \(H_n\) mixes those zones, its inverse contains
longer itineraries and
\eqref{eq:v2v62-zone-sum} is false.  One must then apply another Schur
decomposition or estimate the complete weighted complementary resolvent.
A formal insertion of zone projections does not make the inverse block
diagonal.
\end{firewall}

\begin{firewall}[Ward assembly precedes the excursion norm]
\label{fw:v2v62-Ward}
The blocks in this note are the complete gauge, ghost, matter, contact, and
counterterm blocks in one Ward row.  The exponent budget may be verified
by expanding that row, but cancellations cannot be replaced by products
of separate absolute species bounds.  Nor does the budget prove the local
cohomology statement
\eqref{eq:v2v60-jet-in-local}.
\end{firewall}

\subsection{V62 conclusion and next acquisition}
\label{sec:v2v62-conclusion}

V62 turns the differentiated Schur obstruction into an exact finite-order
budget.  It proves that low and high complementary sectors demand different
crossing estimates, and it closes the V60 fifth-derivative bound whenever
those estimates hold through order five.

The main unproved input is now the low-zone crossing estimate
\[
 \|D^a(P_nA_nQ_{n,L})\|
 \lesssim\Lambda_n^{1/2-a}
\]
or a stronger bound.  V63 will attack that estimate by a spectral-gap
commutator identity.  It will separate genuinely low modes from adjacent
transition modes and determine which commutators with the reference
kinetic operator yield the required half-order gain.
\clearpage
\section*{Second Volume, Version V63: Spectral-Gap Commutators for the Low-Zone Crossing}
\addcontentsline{toc}{section}{Second Volume, Version V63: Spectral-Gap Commutators for the Low-Zone Crossing}

\subsection*{Purpose}

V62 reduced the differentiated low excursion to a crossing estimate of
degree at most \(1/2\) on each side of a uniformly bounded low inverse.
This note proves a sharper estimate whenever the transition shell and the
low sector are separated in the spectrum of a fixed reference kinetic
operator.  The proof is an exact Sylvester inversion, not a symbolic
analogy.

The result also separates two notions that must not be merged.  A mode can
have frequency comparable to \(\Lambda_n\) and still be adjacent to the
chosen transition shell, in which case there is no gap of order
\(\Lambda_n\).  Only the genuinely low sector receives the commutator
gain proved here.

\subsection{Signed transition shells}
\label{sec:v2v63-shells}

Work first at finite regulator on a Hilbert space \(\mathcal H_n\).  Let
\(D_0=D_0^*\), let \(0\le c_L<c_-\le c_+\), and let
\(\Lambda_n\ge1\).  Define the mutually orthogonal spectral projections
\begin{align}
 P_{n,+}
 &=
 \mathbf1_{[c_-\Lambda_n,c_+\Lambda_n]}(D_0),
\label{eq:v2v63-Pplus}\\
 P_{n,-}
 &=
 \mathbf1_{[-c_+\Lambda_n,-c_-\Lambda_n]}(D_0),
\label{eq:v2v63-Pminus}\\
 Q_{n,L}
 &=
 \mathbf1_{[-c_L\Lambda_n,c_L\Lambda_n]}(D_0),
\label{eq:v2v63-QL}
\end{align}
and put \(P_n=P_{n,+}+P_{n,-}\).  Set
\begin{equation}
 \delta=c_--c_L>0.
\label{eq:v2v63-delta}
\end{equation}

\begin{lemma}[Positive-shell Sylvester inverse]
\label{lem:v2v63-positive}
Let \(T:\operatorname{Ran}Q_{n,L}\to\operatorname{Ran}P_{n,+}\) be
bounded.  The equation
\begin{equation}
 (P_{n,+}D_0P_{n,+})X
 -
 X(Q_{n,L}D_0Q_{n,L})
 =T
\label{eq:v2v63-Sylvester-plus}
\end{equation}
has the solution
\begin{equation}
 X
 =
 \int_0^\infty
 e^{-tP_{n,+}D_0P_{n,+}}\,
 T\,
 e^{tQ_{n,L}D_0Q_{n,L}}\,\mathrm dt,
\label{eq:v2v63-plus-integral}
\end{equation}
and
\begin{equation}
 \|X\|
 \le
 \frac{\|T\|}{\delta\Lambda_n}.
\label{eq:v2v63-plus-bound}
\end{equation}
\end{lemma}

\begin{proof}
The spectral bounds give
\begin{align*}
 \|e^{-tP_{n,+}D_0P_{n,+}}\|
 &\le e^{-c_-\Lambda_nt},\\
 \|e^{tQ_{n,L}D_0Q_{n,L}}\|
 &\le e^{c_L\Lambda_nt}.
\end{align*}
Thus the integral converges in operator norm and its norm is bounded by
\[
 \|T\|\int_0^\infty e^{-\delta\Lambda_nt}\,\mathrm dt
 =
 \frac{\|T\|}{\delta\Lambda_n}.
\]
Differentiating the integrand shows that the left side of
\eqref{eq:v2v63-Sylvester-plus} is the negative integral of its
\(t\)-derivative.  The boundary value at infinity is zero and the value at
zero is \(T\), which proves the equation.
\end{proof}

\begin{lemma}[Negative-shell Sylvester inverse]
\label{lem:v2v63-negative}
With \(T:\operatorname{Ran}Q_{n,L}\to\operatorname{Ran}P_{n,-}\), the
same Sylvester equation has the solution
\begin{equation}
 X
 =
 -\int_0^\infty
 e^{tP_{n,-}D_0P_{n,-}}\,
 T\,
 e^{-tQ_{n,L}D_0Q_{n,L}}\,\mathrm dt
\label{eq:v2v63-minus-integral}
\end{equation}
and satisfies
\begin{equation}
 \|X\|
 \le
 \frac{\|T\|}{\delta\Lambda_n}.
\label{eq:v2v63-minus-bound}
\end{equation}
\end{lemma}

\begin{proof}
Now
\[
 \|e^{tP_{n,-}D_0P_{n,-}}\|
 \le e^{-c_-\Lambda_nt},
 \qquad
 \|e^{-tQ_{n,L}D_0Q_{n,L}}\|
 \le e^{c_L\Lambda_nt}.
\]
The norm bound is the same.  The extra minus sign makes the integral of
the \(t\)-derivative equal the required Sylvester right side.
\end{proof}

\subsection{The exact commutator crossing estimate}
\label{sec:v2v63-crossing}

Let \(V_n(k)\) be a \(C^M\) bounded family at finite regulator.  The
projections above are independent of the external parameter \(k\).
Since \(P_{n,\pm}\) and \(Q_{n,L}\) commute with \(D_0\), one has
\begin{equation}
 \begin{split}
 &(P_{n,\pm}D_0P_{n,\pm})
   (P_{n,\pm}V_nQ_{n,L})\\
 &\qquad
 -(P_{n,\pm}V_nQ_{n,L})
   (Q_{n,L}D_0Q_{n,L})
 =
 P_{n,\pm}[D_0,V_n]Q_{n,L}.
 \end{split}
\label{eq:v2v63-crossing-Sylvester}
\end{equation}

\begin{theorem}[Gap-commutator crossing theorem]
\label{thm:v2v63-crossing}
For \(0\le a\le M\),
\begin{align}
 &\|D^a(P_nV_n(k)Q_{n,L})\|_{\mathcal L^a}
\notag\\
 &\qquad\le
 \frac{2}{\delta\Lambda_n}
 \max_{\sigma\in\{+,-\}}
 \|P_{n,\sigma}[D_0,D^aV_n(k)]Q_{n,L}\|_{\mathcal L^a}.
\label{eq:v2v63-crossing-bound}
\end{align}
The reverse crossing obeys
\begin{align}
 &\|D^a(Q_{n,L}V_n(k)P_n)\|_{\mathcal L^a}
\notag\\
 &\qquad\le
 \frac{2}{\delta\Lambda_n}
 \max_{\sigma\in\{+,-\}}
 \|Q_{n,L}[D_0,D^aV_n(k)]P_{n,\sigma}\|_{\mathcal L^a}.
\label{eq:v2v63-reverse-bound}
\end{align}
\end{theorem}

\begin{proof}
Because the projections and \(D_0\) do not depend on \(k\),
external differentiation commutes with every compression in
\eqref{eq:v2v63-crossing-Sylvester}.  Apply
Lemmas~\ref{lem:v2v63-positive} and
\ref{lem:v2v63-negative} to \(D^aV_n(k)\), in every collection of
directions.  Add the positive and negative shell bounds.  This gives
\eqref{eq:v2v63-crossing-bound}; the factor two is harmless and avoids
using a Hilbert-column square estimate.  Apply the same argument to the
adjoint orientation, or interchange the two spectral blocks, to obtain
\eqref{eq:v2v63-reverse-bound}.
\end{proof}

\begin{corollary}[The half-order target]
\label{cor:v2v63-half}
If
\begin{align}
 \|P_{n,\sigma}[D_0,D^aV_n(k)]Q_{n,L}\|_{\mathcal L^a}
 &\le c_a\Lambda_n^{3/2-a},
\label{eq:v2v63-half-comm}\\
 \|Q_{n,L}[D_0,D^aV_n(k)]P_{n,\sigma}\|_{\mathcal L^a}
 &\le c_a\Lambda_n^{3/2-a}
\label{eq:v2v63-half-comm-reverse}
\end{align}
for \(\sigma\in\{+,-\}\), then
\begin{align}
 \|D^a(P_nV_nQ_{n,L})\|_{\mathcal L^a}
 &\le C_a\Lambda_n^{1/2-a},
\label{eq:v2v63-half-X}\\
 \|D^a(Q_{n,L}V_nP_n)\|_{\mathcal L^a}
 &\le C_a\Lambda_n^{1/2-a}.
\label{eq:v2v63-half-Y}
\end{align}
Thus the two crossings satisfy the V62 low-zone budget
with \(\alpha=\beta=1/2\).
\end{corollary}

\begin{proof}
Insert
\eqref{eq:v2v63-half-comm} and
\eqref{eq:v2v63-half-comm-reverse} in
Theorem~\ref{thm:v2v63-crossing} and use the inverse gap
\((\delta\Lambda_n)^{-1}\).
\end{proof}

\begin{corollary}[An order-one commutator gives strict room]
\label{cor:v2v63-order-one}
If the two commutator blocks in
\eqref{eq:v2v63-crossing-bound}--\eqref{eq:v2v63-reverse-bound}
are bounded by
\[
 c_a\Lambda_n^{1-a},
\]
then both crossings are bounded by
\[
 C_a\Lambda_n^{-a}.
\]
They satisfy the V62 budget with \(\alpha=\beta=0\), leaving one full
power of strict room when the low inverse has \(\gamma=0\).
\end{corollary}

\begin{proof}
Apply Theorem~\ref{thm:v2v63-crossing}.  The shell gap removes one power
of \(\Lambda_n\).
\end{proof}

\begin{assessment}[What symbolic commutator lowering must mean]
\label{ass:v2v63-symbolic}
For first-order \(D_0\) and a first-order differential or
pseudodifferential perturbation, the principal order-two part of
\([D_0,V_n]\) is absent at the level of scalar commuting principal symbols.
Corollary~\ref{cor:v2v63-order-one} may then be expected.  The corollary,
however, requires the displayed compressed operator estimates through
external derivative order five.  Formal symbol order does not prove their
uniform constants, domains, or Ward-completed background dependence.
\end{assessment}

\subsection{Insertion into the low Schur self-energy}
\label{sec:v2v63-Schur}

Suppose \(A_n(k)=D_0+V_n(k)\).  Let
\[
 H_{n,L}(k)=Q_{n,L}A_n(k)Q_{n,L},
 \qquad
 R_{n,L}(k)=H_{n,L}(k)^{-1},
\]
and suppose
\begin{equation}
 \|D^aR_{n,L}(k)\|_{\mathcal L^a}
 \le r_a\Lambda_n^{-a},
 \qquad 0\le a\le5.
\label{eq:v2v63-low-inverse}
\end{equation}
This is the V62 complement estimate with \(\gamma=0\).  Define the low
self-energy
\begin{equation}
 \Sigma_{n,L}
 =
 P_nV_nQ_{n,L}\,
 R_{n,L}\,
 Q_{n,L}V_nP_n.
\label{eq:v2v63-low-Sigma}
\end{equation}

\begin{theorem}[Gap closure of the low excursion]
\label{thm:v2v63-low-close}
Assume
\eqref{eq:v2v63-low-inverse} and the commutator bounds
\eqref{eq:v2v63-half-comm}--\eqref{eq:v2v63-half-comm-reverse}
for \(0\le a\le5\).  Then
\begin{equation}
 \|D^m\Sigma_{n,L}(k)\|_{\mathcal L^m}
 \le C_m\Lambda_n^{1-m},
 \qquad 0\le m\le5.
\label{eq:v2v63-low-Sigma-degree}
\end{equation}
If the order-one commutator hypothesis of
Corollary~\ref{cor:v2v63-order-one} holds, the stronger bound
\begin{equation}
 \|D^m\Sigma_{n,L}(k)\|_{\mathcal L^m}
 \le C_m\Lambda_n^{-m}
\label{eq:v2v63-low-Sigma-strict}
\end{equation}
holds.
\end{theorem}

\begin{proof}
In the half-order case, use
Corollary~\ref{cor:v2v63-half} in the weighted excursion
Theorem~\ref{thm:v2v62-budget} with
\[
 \alpha=\beta=\frac12,
 \qquad
 \gamma=0.
\]
Their sum equals one.  In the order-one commutator case, use
\(\alpha=\beta=\gamma=0\); equation
\eqref{eq:v2v62-Sigma-degree} then gives
\eqref{eq:v2v63-low-Sigma-strict}.
\end{proof}

\begin{firewall}[The low inverse remains a separate acquisition]
\label{fw:v2v63-low-inverse}
The spectral-gap argument controls only the two crossings.  It does not
prove
\eqref{eq:v2v63-low-inverse}.  A uniformly bounded low inverse gives only
the case \(a=0\); its external derivatives can remain \(O(1)\), exactly as
the V61 counterexample warns.  The finite fifth-order problem must keep
track of any derivative deficit inside \(R_{n,L}\).
\end{firewall}

\subsection{Why adjacent modes receive no gap gain}
\label{sec:v2v63-adjacent}

\begin{counterexample}[Frequency size does not imply spectral separation]
\label{sep:v2v63-adjacent}
On \(\mathbb C^2\), let
\[
 D_{0,\Lambda}
 =
 \begin{pmatrix}
  \Lambda&0\\
  0&\Lambda-1
 \end{pmatrix},
 \qquad
 P=
 \begin{pmatrix}1&0\\0&0\end{pmatrix},
 \qquad
 Q=I-P,
\]
and let \(V\) be the off-diagonal flip.  Then
\begin{equation}
 \|PVQ\|=1,
 \qquad
 \|P[D_{0,\Lambda},V]Q\|=1.
\label{eq:v2v63-adjacent}
\end{equation}
Although both eigenvalues are of order \(\Lambda\), their gap is one, so
no inverse power of \(\Lambda\) follows.
\end{counterexample}

\begin{proof}
Both displayed compressed matrices have one nonzero entry of absolute
value one.  The exact Sylvester denominator is
\(\Lambda-(\Lambda-1)=1\).
\end{proof}

\begin{assessment}[Required spectral split]
\label{ass:v2v63-split}
The complementary space must be split at least into a genuinely low block,
an adjacent transition block, and a high block before the Schur estimate.
The low block is treated by the present gap theorem.  Adjacent transition
modes require the shell inverse budget \(\gamma=-1\) or a wider smooth
transition.  Combining them into one \(Q_n\) discards both mechanisms.
\end{assessment}

\begin{firewall}[Moving projections create derivative terms]
\label{fw:v2v63-moving-P}
The proof assumes that \(D_0\), \(P_{n,\pm}\), and \(Q_{n,L}\) are fixed
with respect to the external momentum.  If the projection itself depends
on the differentiated parameter, the product rule creates derivatives of
the projection.  Those terms require the Riesz-projector estimates of
V38--V39 and are not contained in
\eqref{eq:v2v63-crossing-bound}.
\end{firewall}

\subsection{V63 conclusion and next acquisition}
\label{sec:v2v63-conclusion}

V63 proves the required low-zone crossing estimate from a quantitative
spectral gap and differentiated commutator bounds.  It also proves that
adjacent modes cannot be treated by the same argument.  Together with V62,
this closes every differentiated low excursion for which the low inverse
obeys
\eqref{eq:v2v63-low-inverse}.

The remaining low-zone obstruction is finite and explicit: derivatives of
the low inverse may fail to gain one shell power each.  V64 will replace
the uniform exponent ansatz by a fifth-order derivative ledger.  It will
show exactly how surplus decay in the two gap crossings can pay a deficit
in \(D^jR_{n,L}\), and it will identify the weakest collection of
inequalities needed for the V60 fifth derivative.
\clearpage
\section*{Second Volume, Version V64: The Fifth-Order Excursion Ledger}
\addcontentsline{toc}{section}{Second Volume, Version V64: The Fifth-Order Excursion Ledger}

\subsection*{Purpose}

V62 used one affine degree law for every derivative of each factor in a
Schur excursion.  V63 showed why that can be too rigid: a spectral-gap
commutator may give surplus decay in the two crossings while derivatives
of the low inverse have a deficit.  Because the V60 target uses only
derivatives through order five, an all-order degree law is unnecessary.

This note gives the exact finite ledger.  It contains one inequality for
each distribution of at most five labelled derivatives among the left
crossing, low inverse, and right crossing.  A scalar jet construction proves
that none of these monomial entries can be omitted from a universal
factorwise criterion.  The ledger then reduces to a simple statement:
the sum of the two crossing surpluses must dominate every low-inverse
derivative deficit through order five.

\subsection{Raw derivative majorants}
\label{sec:v2v64-raw}

Let
\[
 \Sigma_n(k)=X_n(k)R_n(k)Y_n(k),
 \qquad
 k\in B_\kappa\subset\mathbb R^q,
\]
with all three factors \(C^5\).  Define the uniform derivative majorants
\begin{align}
 x_{i,n}
 &=
 \sup_{|k|\le\kappa}
 \|D^iX_n(k)\|_{\mathcal L^i},
\label{eq:v2v64-x}\\
 r_{j,n}
 &=
 \sup_{|k|\le\kappa}
 \|D^jR_n(k)\|_{\mathcal L^j},
\label{eq:v2v64-r}\\
 y_{\ell,n}
 &=
 \sup_{|k|\le\kappa}
 \|D^\ell Y_n(k)\|_{\mathcal L^\ell}.
\label{eq:v2v64-y}
\end{align}
For order zero the multilinear norm means the ordinary operator norm.

\begin{lemma}[Exact multinomial majorant]
\label{lem:v2v64-multinomial}
For \(0\le m\le5\),
\begin{equation}
 \sup_{|k|\le\kappa}
 \|D^m\Sigma_n(k)\|_{\mathcal L^m}
 \le
 \sum_{\substack{i,j,\ell\ge0\\i+j+\ell=m}}
 \frac{m!}{i!\,j!\,\ell!}\,
 x_{i,n}r_{j,n}y_{\ell,n}.
\label{eq:v2v64-multinomial}
\end{equation}
\end{lemma}

\begin{proof}
Use the labelled assignment formula
\eqref{eq:v2v62-assignment}.  Exactly
\(m!/(i!j!\ell!)\) assignments send \(i\) directions to \(X_n\), \(j\)
to \(R_n\), and \(\ell\) to \(Y_n\).  Bound each resulting ordered product
by submultiplicativity and take the supremum.
\end{proof}

\begin{theorem}[Finite fifth-order excursion ledger]
\label{thm:v2v64-ledger}
Suppose that for every triple \(i,j,\ell\ge0\) with
\(i+j+\ell\le5\), there is a constant \(c_{i,j,\ell}\), independent of
\(n\), such that
\begin{equation}
 x_{i,n}r_{j,n}y_{\ell,n}
 \le
 c_{i,j,\ell}\,
 \Lambda_n^{1-i-j-\ell}.
\label{eq:v2v64-ledger}
\end{equation}
Then
\begin{equation}
 \sup_{|k|\le\kappa}
 \|D^m\Sigma_n(k)\|_{\mathcal L^m}
 \le C_m\Lambda_n^{1-m},
 \qquad 0\le m\le5.
\label{eq:v2v64-ledger-conclusion}
\end{equation}
There are
\begin{equation}
 \sum_{m=0}^5\binom{m+2}{2}
 =
 \binom83
 =
 56
\label{eq:v2v64-count}
\end{equation}
ledger entries.
\end{theorem}

\begin{proof}
Insert
\eqref{eq:v2v64-ledger} into
\eqref{eq:v2v64-multinomial}.  Every term with total order \(m\) has the
same scale \(\Lambda_n^{1-m}\), and the finite multinomial sum gives
\(C_m\).  The number of nonnegative triples with total at most five is the
stars-and-bars count in
\eqref{eq:v2v64-count}.
\end{proof}

\subsection{Sharpness of the componentwise ledger}
\label{sec:v2v64-sharp}

\begin{proposition}[Every monomial entry is independently visible]
\label{prop:v2v64-sharp}
Fix a triple \((i,j,\ell)\) with \(i+j+\ell=m\le5\) and positive numbers
\(x,r,y\).  There are scalar polynomial germs \(X,R,Y\) at zero such that
\begin{align*}
 X^{(a)}(0)&=0\quad(a\ne i),&
 X^{(i)}(0)&=x,\\
 R^{(b)}(0)&=0\quad(b\ne j),&
 R^{(j)}(0)&=r,\\
 Y^{(c)}(0)&=0\quad(c\ne\ell),&
 Y^{(\ell)}(0)&=y,
\end{align*}
for all derivative orders at most \(m\), and
\begin{equation}
 (XRY)^{(m)}(0)
 =
 \frac{m!}{i!\,j!\,\ell!}\,xry.
\label{eq:v2v64-sharp}
\end{equation}
If a nonvanishing inverse factor is required, the construction may be
perturbed by a positive constant and restricted to a sufficiently small
neighbourhood, after which it is \(H^{-1}\) for the smooth scalar
\(H=R^{-1}\).
\end{proposition}

\begin{proof}
Before the optional perturbation, take
\[
 X(t)=\frac{x}{i!}t^i,
 \qquad
 R(t)=\frac{r}{j!}t^j,
 \qquad
 Y(t)=\frac{y}{\ell!}t^\ell.
\]
Their product is
\(xry\,t^m/(i!j!\ell!)\), proving
\eqref{eq:v2v64-sharp}.  To make \(R\) nonzero, add a positive constant.
If \(j>0\), choose \(X\) and \(Y\) as above; the added term has total
degree \(i+\ell<m\) and therefore vanishes after the \(m\)-th derivative.
If \(j=0\), choose the positive constant itself to be \(r\).
Continuity gives a neighbourhood on which \(R\) is invertible.
\end{proof}

\begin{assessment}[Meaning of sharpness]
\label{ass:v2v64-sharpness}
Proposition~\ref{prop:v2v64-sharp} does not say that every physical
Standard Model block realizes arbitrary jets.  It says that a proof using
only the separate derivative majorants
\eqref{eq:v2v64-x}--\eqref{eq:v2v64-y} cannot discard a failing ledger
entry: a scalar product can isolate it.  Ward identities or correlations
between factors may yield a stronger assembled estimate, but that is
additional structure.
\end{assessment}

\subsection{Crossing surplus and inverse deficit}
\label{sec:v2v64-surplus}

The half-order V62 baseline for a low excursion is
\[
 \|D^aX_n\|\sim\Lambda_n^{1/2-a},
 \qquad
 \|D^aY_n\|\sim\Lambda_n^{1/2-a},
 \qquad
 \|D^aR_n\|\sim\Lambda_n^{-a}.
\]
Allow fixed crossing surpluses
\(\sigma_X,\sigma_Y\ge0\), and let
\(\ell_j\ge0\) measure the loss in the \(j\)-th low-inverse derivative.

\begin{theorem}[Compressed fifth-order ledger]
\label{thm:v2v64-compressed}
Suppose, for \(0\le a,j\le5\),
\begin{align}
 x_{a,n}
 &\le x_a\Lambda_n^{1/2-a-\sigma_X},
\label{eq:v2v64-X-surplus}\\
 y_{a,n}
 &\le y_a\Lambda_n^{1/2-a-\sigma_Y},
\label{eq:v2v64-Y-surplus}\\
 r_{j,n}
 &\le r_j\Lambda_n^{-j+\ell_j}.
\label{eq:v2v64-R-deficit}
\end{align}
If
\begin{equation}
 \ell_j\le\sigma_X+\sigma_Y
 \qquad(0\le j\le5),
\label{eq:v2v64-surplus-condition}
\end{equation}
then the full ledger
\eqref{eq:v2v64-ledger} holds and hence
\eqref{eq:v2v64-ledger-conclusion} follows.
Within these separated power bounds, condition
\eqref{eq:v2v64-surplus-condition} is sharp for each \(j\).
\end{theorem}

\begin{proof}
For \(i+j+\ell=m\), the product exponent is
\begin{align*}
 &(1/2-i-\sigma_X)+(-j+\ell_j)
 +(1/2-\ell-\sigma_Y)\\
 &\qquad=
 1-m+\ell_j-\sigma_X-\sigma_Y.
\end{align*}
Condition
\eqref{eq:v2v64-surplus-condition} makes this no larger than \(1-m\).
Thus every ledger entry holds.  If the condition fails for one \(j\),
choose any \(i,\ell\) with \(i+j+\ell\le5\) and use the scalar jet
construction of Proposition~\ref{prop:v2v64-sharp} with scale-dependent
coefficients saturating
\eqref{eq:v2v64-X-surplus}--\eqref{eq:v2v64-R-deficit}.
The resulting derivative has exponent strictly larger than \(1-m\).
\end{proof}

\begin{corollary}[What the V63 strict commutator gain pays]
\label{cor:v2v64-V63}
Under the order-one commutator hypothesis of
Corollary~\ref{cor:v2v63-order-one}, the two crossings satisfy
\[
 x_{a,n},y_{a,n}\lesssim\Lambda_n^{-a}.
\]
Relative to the half-order baseline, this means
\[
 \sigma_X=\sigma_Y=\frac12.
\]
Consequently the low excursion closes through order five whenever
\begin{equation}
 \ell_j\le1
 \qquad(0\le j\le5).
\label{eq:v2v64-one-loss}
\end{equation}
\end{corollary}

\begin{proof}
The two surpluses sum to one.  Apply
Theorem~\ref{thm:v2v64-compressed}.
\end{proof}

\begin{counterexample}[A bounded differentiated low inverse can exceed the spare power]
\label{sep:v2v64-bounded-R}
Suppose only
\[
 r_{j,n}\le C_j
 \qquad(0\le j\le5).
\]
In the parametrization
\eqref{eq:v2v64-R-deficit}, this corresponds to
\(\ell_j=j\).  The V63 strict crossing gain pays the cases \(j=0,1\), but
does not pay \(j=2,3,4,5\).  In particular, these assumptions permit scalar data with
\[
 x_{0,n}r_{5,n}y_{0,n}
 \asymp 1,
\]
whereas the corresponding ledger entry requires the scale
\(\Lambda_n^{-4}\).
\end{counterexample}

\begin{proof}
With \(\sigma_X+\sigma_Y=1\),
condition
\eqref{eq:v2v64-surplus-condition} becomes \(j\le1\).  The last displayed
comparison is the triple \((0,5,0)\); choose the scalar jets in
Proposition~\ref{prop:v2v64-sharp} with all three indicated derivatives
equal to one.  The same proposition shows that the failure cannot be removed
from a factorwise argument.
\end{proof}

\subsection{Insertion into the effective resolvent}
\label{sec:v2v64-effective}

Let
\[
 \mathscr S_n=K_n-\Sigma_n,
 \qquad
 P_nA_n^{-1}P_n=\mathscr S_n^{-1}.
\]

\begin{theorem}[Ledger-to-fifth-resolvent compiler]
\label{thm:v2v64-compiler}
Assume the 56 ledger bounds
\eqref{eq:v2v64-ledger}, the diagonal bounds
\begin{equation}
 \|D^mK_n(k)\|_{\mathcal L^m}
 \le k_m\Lambda_n^{1-m},
 \qquad 0\le m\le5,
\label{eq:v2v64-K}
\end{equation}
and the effective inverse bound
\begin{equation}
 \|\mathscr S_n(k)^{-1}\|
 \le g\Lambda_n^{-1}.
\label{eq:v2v64-S0}
\end{equation}
Then
\begin{equation}
 \|D^m(P_nA_n(k)^{-1}P_n)\|_{\mathcal L^m}
 \le g_m\Lambda_n^{-1-m},
 \qquad 0\le m\le5.
\label{eq:v2v64-G}
\end{equation}
Every complete degree-four V61 resolvent word built from these compressed
resolvents has fifth derivative \(O(\Lambda_n^{-1})\), and its V60
nonlocal Taylor remainder is absolutely summable.
\end{theorem}

\begin{proof}
Theorem~\ref{thm:v2v64-ledger} and
\eqref{eq:v2v64-K} give
\[
 \|D^m\mathscr S_n\|\le C_m\Lambda_n^{1-m}
 \qquad(0\le m\le5).
\]
Together with
\eqref{eq:v2v64-S0}, the effective-shell derivative criterion
Theorem~\ref{thm:v2v61-effective} gives
\eqref{eq:v2v64-G}.  Apply the V61 word theorem and its fifth-derivative
corollary, followed by the V60 locality theorem.
\end{proof}

\begin{firewall}[A ledger is a proof obligation, not a cancellation]
\label{fw:v2v64-ledger}
Writing the 56 inequalities does not establish them.  Their value is that
a failed entry identifies the exact distribution of derivatives that must
be treated upstream.  If the entry \((0,j,0)\) fails, no redistribution of
the product rule can move those \(j\) derivatives onto the crossings.
Only an assembled identity, a stronger crossing surplus, or a direct
estimate of \(D^jR_n\) can close it.
\end{firewall}

\subsection{V64 conclusion and next acquisition}
\label{sec:v2v64-conclusion}

V64 replaces the unnecessarily uniform low-inverse hypothesis of V63 by
the exact finite criterion used by the fifth derivative.  It proves that
the strict V63 commutator estimate tolerates one missing inverse power in
each derivative of the low resolvent, but not an arbitrary bounded
derivative family.

V65 is a mandatory companion audit.  After that audit, the next analytic
choice will be made from the ledger rather than by analogy: either derive
low-resolvent derivative gains from a further spectral split, or bypass
the failing \((0,j,0)\) entries through an assembled Ward identity for the
complete Schur self-energy.
\clearpage
\section*{Second Volume, Version V65: Companion Audit and the Complete-Word Derivative Ledger}
\addcontentsline{toc}{section}{Second Volume, Version V65: Companion Audit and the Complete-Word Derivative Ledger}

\subsection*{Purpose and mandatory audit}

V64 produced the exact fifth-order ledger for one three-factor Schur
excursion.  This mandatory audit changes the next proof order.  The
Navier--Stokes programme again warns that a missing weighted moment cannot
be manufactured by a weaker smoothing estimate.  The Yang--Mills
programme now proves a one-payer compiler for quintic partition products,
but only after the connected object is assembled.  The SM analogue is to
place the derivative ledger on the complete Ward word and use its total
superficial-degree slack before demanding optimal bounds from every low
inverse factor.

The Navier--Stokes file inspected was
\[
 \texttt{Renewal\_NS\_Stability\_Research\_Notes\_V31.tex},
\]
with \(887537\) bytes and SHA--256 digest
\[
 \texttt{F1121B62238AD5491BB7CF03635EA9934942EDF573ED8FAAA9EE79E1D38A6696}.
\]
It is unchanged from the V60 audit.  Its sharp lesson remains that the
critical first radial moment has to be paid explicitly; heat smoothing
does not create it.  Applied here, bounded derivatives of a low resolvent
cannot be relabelled as inverse shell powers.

The active Yang--Mills file inspected was
\[
 \texttt{Renewal\_Yang\_Mills\_Mass\_Gap\_Research\_Notes\_V57.tex},
\]
with \(761879\) bytes and SHA--256 digest
\[
 \texttt{4935205BB613FA7F60B4F290E028F479DA6B545423FC0F5EBDE635B4C5900897}.
\]
The frozen first YM volume
\[
 \texttt{Renewal\_Yang\_Mills\_Mass\_Gap\_Research\_Notes\_V135\_f1.tex}
\]
was present with \(3083467\) bytes and digest
\[
 \texttt{82F6BBFFEECDAD3C5C63551660B19A0BDC7F3D6DFFE4C68B4E9BBA3A783A2B71};
\]
it was not treated as the current active note.

YM V57 combines the complete binary, ternary, and quartic bounds into a
quintic partition product with one and only one output resolvent.  Its
forest-plus-quotient-tree lemma preserves the exact normalized power.  It
does not yet prove the connected quotient-tree factor for every
mixed-support quintic fibre, so global quintic MTA and the continuum mass
gap remain open.

\begin{assessment}[V65 audit decision]
\label{ass:v2v65-audit}
The SM fifth derivative will be estimated on complete connected Ward
words.  Degree slack and exact cancellations will be used there before a
norm is taken.  The V64 three-factor ledger remains a valid sufficient
test, but a failure of one factorwise entry does not end the route if the
complete word has lower superficial degree or the offending allocation
cancels in an assembled Ward packet.
\end{assessment}

\begin{firewall}[Cross-program type boundary]
\label{fw:v2v65-types}
A Yang--Mills output payer is not a low fermion or gauge resolvent, and an
NS radial moment is not an external-momentum derivative.  The audit
transfers only proof structure: retain the missing weight honestly, form
the connected object, allocate the available degree once, and then take
the physical norm.
\end{firewall}

\subsection{A derivative ledger for a complete word}
\label{sec:v2v65-word}

Let \(N\ge1\), let
\[
 F_{\nu,n}:B_\kappa\subset\mathbb R^q
 \longrightarrow\mathcal L(\mathcal H_{\nu,n}),
 \qquad 1\le\nu\le N,
\]
be \(C^5\), and let the \(k\)-independent multilinear assembly map
\(\tau_n\) obey
\begin{equation}
 \|\tau_n(Z_1,\ldots,Z_N)\|_{\mathcal Z}
 \le w\Lambda_n^d\prod_{\nu=1}^N\|Z_\nu\|.
\label{eq:v2v65-assembly}
\end{equation}
This map may include a trace, tensor contractions, fixed contact
insertions, and a finite species sum that has already been assembled.
Define
\begin{equation}
 \mathscr W_n(k)
 =
 \tau_n(F_{1,n}(k),\ldots,F_{N,n}(k)).
\label{eq:v2v65-word}
\end{equation}

Assign a nominal degree \(a_\nu\) to each factor and a derivative
deviation \(e_{\nu,j}\in\mathbb R\) by assuming
\begin{equation}
 \sup_{|k|\le\kappa}
 \|D^jF_{\nu,n}(k)\|_{\mathcal L^j}
 \le c_{\nu,j}
 \Lambda_n^{a_\nu-j+e_{\nu,j}},
 \qquad 0\le j\le5.
\label{eq:v2v65-factor-degree}
\end{equation}
A positive \(e_{\nu,j}\) is a derivative deficit and a negative one is a
surplus.  Put
\begin{align}
 \rho
 &=
 d+\sum_{\nu=1}^Na_\nu,
\label{eq:v2v65-rho}\\
 E_m
 &=
 \max_{\substack{j_1,\ldots,j_N\ge0\\
                  j_1+\cdots+j_N=m}}
 \sum_{\nu=1}^Ne_{\nu,j_\nu}.
\label{eq:v2v65-Em}
\end{align}

\begin{theorem}[Complete-word derivative ledger]
\label{thm:v2v65-word-ledger}
Under
\eqref{eq:v2v65-assembly}--\eqref{eq:v2v65-factor-degree},
\begin{equation}
 \sup_{|k|\le\kappa}
 \|D^m\mathscr W_n(k)\|_{\mathcal L^m}
 \le C_m\Lambda_n^{\rho-m+E_m},
 \qquad 0\le m\le5.
\label{eq:v2v65-word-ledger}
\end{equation}
In particular, if
\begin{equation}
 \rho+E_5\le4,
\label{eq:v2v65-fifth-condition}
\end{equation}
then
\begin{equation}
 \|D^5\mathscr W_n\|_{\mathcal L^5}
 \le C\Lambda_n^{-1}.
\label{eq:v2v65-fifth}
\end{equation}
\end{theorem}

\begin{proof}
The labelled multilinear product rule assigns each of the five derivative
directions to one of the \(N\) factors.  After terms with the same
derivative counts are collected, a count vector
\((j_1,\ldots,j_N)\) with total \(m\) has multinomial coefficient
\(m!/\prod_\nu j_\nu!\).  Its scale exponent is
\[
 d+\sum_{\nu=1}^N
 (a_\nu-j_\nu+e_{\nu,j_\nu})
 =
 \rho-m+\sum_{\nu=1}^Ne_{\nu,j_\nu}.
\]
The largest deviation is \(E_m\); the number of count vectors is finite.
This proves
\eqref{eq:v2v65-word-ledger}.  Condition
\eqref{eq:v2v65-fifth-condition} makes the fifth-order exponent at most
\(-1\).
\end{proof}

\begin{corollary}[Superficial-degree slack pays derivative deficit]
\label{cor:v2v65-slack}
If \(\rho\le4-\sigma\) for some \(\sigma\ge0\), then the fifth derivative
has the required inverse scale whenever
\begin{equation}
 E_5\le\sigma.
\label{eq:v2v65-slack}
\end{equation}
Thus degree-three words tolerate one total power of derivative deficit,
degree-two words tolerate two, and only marginal degree-four words have no
automatic slack.
\end{corollary}

\begin{proof}
The hypotheses give
\(\rho+E_5\le4-\sigma+\sigma=4\).  Apply
Theorem~\ref{thm:v2v65-word-ledger}.
\end{proof}

\begin{corollary}[One-deficit allocation]
\label{cor:v2v65-one-deficit}
Suppose every nonzero fifth-derivative allocation in a word has at most one
factor with positive deviation, all other deviations sum to a nonpositive
number, and the positive deviation is at most \(4-\rho\).  Then
\eqref{eq:v2v65-fifth} holds.
\end{corollary}

\begin{proof}
Every allocation has total deviation at most \(4-\rho\), so
\(E_5\le4-\rho\).  Apply
\eqref{eq:v2v65-fifth-condition}.
\end{proof}

\subsection{Sharpness before assembled cancellation}
\label{sec:v2v65-sharp}

\begin{proposition}[A derivative allocation can be isolated]
\label{prop:v2v65-sharp}
Fix nonnegative \(j_1,\ldots,j_N\) with total \(m\le5\).  For arbitrary
positive scale-dependent coefficients \(z_{\nu,n}\), there are scalar
polynomial factors satisfying
\[
 D^aF_{\nu,n}(0)=0\quad(a\ne j_\nu),
 \qquad
 D^{j_\nu}F_{\nu,n}(0)=z_{\nu,n},
\]
through order \(m\), such that
\begin{equation}
 D^m\!\left(\prod_{\nu=1}^NF_{\nu,n}\right)(0)
 =
 \frac{m!}{\prod_\nu j_\nu!}
 \prod_{\nu=1}^Nz_{\nu,n}.
\label{eq:v2v65-sharp}
\end{equation}
Consequently the exponent \(E_m\) in
\eqref{eq:v2v65-word-ledger} is sharp for arguments based only on the
separate factor majorants.
\end{proposition}

\begin{proof}
Take
\[
 F_{\nu,n}(t)
 =
 \frac{z_{\nu,n}}{j_\nu!}t^{j_\nu}.
\]
Their product is a scalar multiple of \(t^m\), and differentiation gives
\eqref{eq:v2v65-sharp}.  Choosing the coefficients to saturate
\eqref{eq:v2v65-factor-degree} realizes the maximal exponent.
\end{proof}

\begin{firewall}[Factorwise sharpness is not Ward sharpness]
\label{fw:v2v65-sharpness}
Proposition~\ref{prop:v2v65-sharp} forbids deleting a bad allocation from a
pure product-majorant proof.  It does not forbid an exact cancellation
among complete species, ghost, contact, and counterterm words.  Such a
cancellation must be written before the norm; it cannot be inferred from
the separate bounds.
\end{firewall}

\subsection{Exact cancellation packets}
\label{sec:v2v65-packets}

Let the fifth derivative of a finite complete Ward row have an algebraic
decomposition
\begin{equation}
 D^5\mathscr L_n
 =
 \sum_{\omega\in\Omega_{\rm good}}G_{\omega,n}
 +
 \sum_{\pi\in\Pi}
 \left(\sum_{\omega\in\Omega_\pi}B_{\omega,n}\right),
\label{eq:v2v65-packets}
\end{equation}
where each inner sum is one tensor- and species-compatible packet.

\begin{lemma}[Cancellation precedes the ledger norm]
\label{lem:v2v65-packets}
Suppose
\begin{align}
 \|G_{\omega,n}\|_{\mathcal L^5}
 &\le C_\omega\Lambda_n^{-1},
 \qquad \omega\in\Omega_{\rm good},
\label{eq:v2v65-good}\\
 \sum_{\omega\in\Omega_\pi}B_{\omega,n}
 &=0,
 \qquad \pi\in\Pi.
\label{eq:v2v65-cancel}
\end{align}
Then
\begin{equation}
 \|D^5\mathscr L_n\|_{\mathcal L^5}
 \le C\Lambda_n^{-1}.
\label{eq:v2v65-packet-close}
\end{equation}
No bound on an individual \(B_{\omega,n}\) is required.
\end{lemma}

\begin{proof}
Use the exact identities
\eqref{eq:v2v65-cancel} in
\eqref{eq:v2v65-packets}.  Only the finite good sum remains.  Apply
\eqref{eq:v2v65-good} and the triangle inequality after the cancellation.
\end{proof}

\begin{counterexample}[Taking norms first destroys an exact packet]
\label{sep:v2v65-norm-first}
For any scalar \(M_n>0\), let
\[
 B_{1,n}=M_n,
 \qquad
 B_{2,n}=-M_n.
\]
Then \(B_{1,n}+B_{2,n}=0\), while
\[
 |B_{1,n}|+|B_{2,n}|=2M_n.
\]
Thus even an arbitrarily large failing factorwise ledger entry may be
absent from the assembled row, but this fact is invisible after separate
norms.
\end{counterexample}

\begin{proof}
Both identities are immediate.
\end{proof}

\subsection{Application to the low-inverse deficit}
\label{sec:v2v65-low}

For a word \(\omega\), let \(\rho_\omega\le4\) be its superficial degree
after forest subtraction and complete Ward assembly.  Fix a
fifth-derivative allocation
\(\mathbf j=(j_1,\ldots,j_N)\).  If \(\mathcal I_L\) indexes its
low-inverse factors, define the total V64 deficit
\begin{equation}
 L_{\omega,\mathbf j}
 =
 \sum_{\nu\in\mathcal I_L}\ell_{\nu,j_\nu}.
\label{eq:v2v65-total-low-deficit}
\end{equation}
Let \(s_{\omega,\mathbf j}\ge0\) be the sum of all crossing and other
factor surpluses in the same allocation.  Its total deviation is
\begin{equation}
 e_{\omega,\mathbf j}
 =
 L_{\omega,\mathbf j}-s_{\omega,\mathbf j}.
\label{eq:v2v65-low-deviation}
\end{equation}

\begin{theorem}[Marginal-word filter]
\label{thm:v2v65-marginal-filter}
The allocation satisfies the V60 fifth-derivative target whenever
\begin{equation}
 L_{\omega,\mathbf j}-s_{\omega,\mathbf j}
 \le
 4-\rho_\omega.
\label{eq:v2v65-marginal-condition}
\end{equation}
Therefore an unresolved low-inverse allocation requires separate analysis
only if it belongs to a surviving packet for which
\begin{equation}
 \rho_\omega
 >
 4-L_{\omega,\mathbf j}+s_{\omega,\mathbf j}.
\label{eq:v2v65-unresolved}
\end{equation}
In particular, the hardest packets are marginal
\(\rho_\omega=4\) packets with uncompensated total deficit and no exact
Ward cancellation.
\end{theorem}

\begin{proof}
Equation
\eqref{eq:v2v65-low-deviation} is the allocation's contribution to
\(E_5\).  Condition
\eqref{eq:v2v65-marginal-condition} is exactly
\(\rho_\omega+e_{\omega,\mathbf j}\le4\), so
Theorem~\ref{thm:v2v65-word-ledger} applies.  Negating that condition gives
\eqref{eq:v2v65-unresolved}.
\end{proof}

\begin{assessment}[Concrete change of direction]
\label{ass:v2v65-direction}
The next proof will not attempt a uniform
\(\|D^jR_{n,L}\|\lesssim\Lambda_n^{-j}\) theorem for every low block.
It will first classify the complete marginal degree-four Ward packets in
which the fifth derivative lands on a low inverse.  Degree-slack packets
close by Theorem~\ref{thm:v2v65-marginal-filter}; exact gauge packets close
by Lemma~\ref{lem:v2v65-packets}.  Only the surviving marginal packet must
supply a new analytic inverse power.
\end{assessment}

\begin{firewall}[The marginal classification is still open]
\label{fw:v2v65-open}
V65 proves the complete-word ledger and the reduction to marginal
uncancelled packets.  It does not enumerate those packets for the full
interacting Standard Model, prove the required renormalized Ward identity
at every regulator, or establish the local cohomology hypothesis of V60.
The SM--Einstein continuum therefore remains open.
\end{firewall}

\subsection{V65 conclusion and next acquisition}
\label{sec:v2v65-conclusion}

V65 incorporates the companion audit without importing mismatched
estimates.  It replaces the demand that every factor attain its nominal
derivative degree by the sharp condition
\(\rho+E_5\le4\) on each complete word.  This automatically removes all
lower-degree packets whose slack pays their derivative deficit and
preserves exact Ward cancellations before norms.

V66 will construct the marginal-packet classifier in a finite
BRST-resolvent word complex.  Its first task is algebraic: identify which
degree-four fifth-derivative allocations are BRST-exact or contact-local,
and isolate the first cohomologically nontrivial low-inverse packet.  Only
that packet will be assigned a new analytic estimate.
\clearpage
\section*{Second Volume, Version V66: Relative BRST Classification of Marginal Fifth-Derivative Packets}
\addcontentsline{toc}{section}{Second Volume, Version V66: Relative BRST Classification of Marginal Fifth-Derivative Packets}

\subsection*{Purpose}

V65 showed that lower-degree complete words close whenever their
superficial-degree slack pays their derivative deficit.  The remaining
objects are complete degree-four Ward packets whose fifth derivatives
contain uncompensated low-inverse losses.  This note constructs the finite
algebraic classifier for those packets.

The classifier is relative cohomology.  A BRST boundary vanishes in the
physical quotient, while a contact-local polynomial of external degree at
most four is annihilated by five derivatives.  Therefore the analytic
estimate is needed only on representatives of
\[
 H^0(\mathcal W/\mathcal L,d),
\]
where \(\mathcal W\) is the finite marginal resolvent-word complex and
\(\mathcal L\) is its contact-local subcomplex.  This is an application of
the V3 splitting framework to a new finite word space; it does not assert
that the resulting Standard Model relative cohomology vanishes.

\subsection{The marginal word complex}
\label{sec:v2v66-complex}

Fix one perturbative order, one shell \(n\), and one finite regulator.
There are finitely many graph types, species labels, contact completions,
and distributions of five labelled external derivatives among the factors
of a graph.  Let
\[
 \mathcal W_n^\bullet
 =
 \bigoplus_g\mathcal W_n^g
\]
be the finite-dimensional graded vector space spanned by those formal
degree-four packets, with grading equal to ghost number.

The complete regulated Ward rules define a degree-one differential
\begin{equation}
 d_n:\mathcal W_n^g\longrightarrow\mathcal W_n^{g+1},
 \qquad
 d_n^2=0.
\label{eq:v2v66-d}
\end{equation}
Here one formal packet includes every ghost, antifield, orientation,
species, and contact contribution required for the Ward rule to close.
Let
\(\mathcal L_n^\bullet\subset\mathcal W_n^\bullet\) be the subspace of
contact-local packets of external polynomial degree at most four.

\begin{definition}[Admissible marginal word complex]
\label{def:v2v66-admissible}
The pair \((\mathcal W_n,\mathcal L_n)\) is admissible if:
\begin{enumerate}
 \item \(d_n^2=0\);
 \item \(d_n\mathcal L_n\subset\mathcal L_n\);
 \item there is a bounded chain evaluation map
 \begin{equation}
  \mathcal E_n:
  \mathcal W_n^\bullet
  \longrightarrow
  C^5(B_\kappa;\mathcal Z^\bullet)
 \label{eq:v2v66-evaluation}
 \end{equation}
 satisfying
 \begin{equation}
  \mathcal E_nd_n=s\mathcal E_n;
 \label{eq:v2v66-chain-evaluation}
 \end{equation}
 \item for every \(\ell\in\mathcal L_n\),
 \(\mathcal E_n\ell\) is a
 \(\mathcal Z\)-valued polynomial in external momentum of total degree at
 most four.
\end{enumerate}
\end{definition}

The evaluation map performs the regulated operator products and complete
tensor contractions.  The word differential \(d_n\) is formal; the
physical BRST differential \(s\) acts on the evaluated Banach complex.

\begin{firewall}[Nilpotence is a regulated Ward theorem]
\label{fw:v2v66-nilpotence}
Definition~\ref{def:v2v66-admissible} does not obtain
\(d_n^2=0\) from notation.  At finite regulator it requires the quantum
master or Slavnov identity with its measure term and contact extensions.
If an anomaly remains, the proposed word differential is curved rather
than nilpotent and the relative cohomology below is not defined.
\end{firewall}

\subsection{The relative obstruction space}
\label{sec:v2v66-relative}

Because \(\mathcal L_n\) is a subcomplex, the quotient
\[
 \overline{\mathcal W}_n
 =
 \mathcal W_n/\mathcal L_n
\]
inherits a differential \(\bar d_n\).

\begin{definition}[Marginal fifth-derivative obstruction]
\label{def:v2v66-obstruction}
The obstruction space is
\begin{equation}
 \mathfrak O_n
 :=
 H^0(\overline{\mathcal W}_n,\bar d_n).
\label{eq:v2v66-obstruction}
\end{equation}
For a closed complete packet
\(w\in\ker(d_n:\mathcal W_n^0\to\mathcal W_n^1)\), write
\begin{equation}
 \mathfrak o_n(w)=[w+\mathcal L_n]\in\mathfrak O_n.
\label{eq:v2v66-class}
\end{equation}
\end{definition}

\begin{lemma}[Concrete zero-class criterion]
\label{lem:v2v66-zero}
For \(w\in\ker d_n\), the following are equivalent:
\begin{enumerate}[label=\textup{(\roman*)}]
 \item \(\mathfrak o_n(w)=0\);
 \item there exist \(u\in\mathcal W_n^{-1}\) and
       \(\ell\in\mathcal L_n^0\) such that
 \begin{equation}
  w=d_nu+\ell.
 \label{eq:v2v66-exact-local}
 \end{equation}
\end{enumerate}
In this representation \(d_n\ell=0\).
\end{lemma}

\begin{proof}
The class of \(w+\mathcal L_n\) is zero in quotient cohomology exactly
when it is a quotient boundary.  Thus, for some \(u\),
\[
 w+\mathcal L_n
 =
 d_nu+\mathcal L_n,
\]
which is equivalent to
\eqref{eq:v2v66-exact-local}.  Conversely that equation makes the quotient
class a boundary.  Finally,
\[
 d_n\ell=d_nw-d_n^2u=0.
\]
\end{proof}

Let
\[
 q_{\rm phys}:
 \ker(s:\mathcal Z^0\to\mathcal Z^1)
 \longrightarrow H_s^0(\mathcal Z)
\]
be the quotient map.  As in V54, assume the boundary space is closed so
that the physical quotient is Hausdorff and \(q_{\rm phys}\) is bounded.

\begin{theorem}[Exact and local packets require no fifth-derivative estimate]
\label{thm:v2v66-elimination}
Let \((\mathcal W_n,\mathcal L_n)\) be admissible and let
\(w\in\ker d_n\).  If
\(\mathfrak o_n(w)=0\), then
\begin{equation}
 D^5\!\left(
 q_{\rm phys}\mathcal E_nw
 \right)=0.
\label{eq:v2v66-elimination}
\end{equation}
\end{theorem}

\begin{proof}
By Lemma~\ref{lem:v2v66-zero}, write
\(w=d_nu+\ell\) with \(\ell\in\mathcal L_n^0\).  The chain property gives
\[
 \mathcal E_nw
 =
 s\mathcal E_nu+\mathcal E_n\ell.
\]
The physical quotient kills the first term.  The second is a polynomial
of external degree at most four by admissibility, so its fifth
Fr\'echet derivative is zero.  The bounded linear map \(q_{\rm phys}\)
commutes with differentiation, proving
\eqref{eq:v2v66-elimination}.
\end{proof}

\begin{corollary}[Only relative classes survive]
\label{cor:v2v66-only-classes}
The map
\begin{equation}
 w\longmapsto
 D^5(q_{\rm phys}\mathcal E_nw)
\label{eq:v2v66-fifth-map}
\end{equation}
on closed ghost-number-zero packets factors uniquely through a linear map
\begin{equation}
 \mathfrak D_n^{(5)}:
 \mathfrak O_n
 \longrightarrow
 \mathcal L^5(\mathbb R^q;H_s^0(\mathcal Z)).
\label{eq:v2v66-factor-map}
\end{equation}
\end{corollary}

\begin{proof}
Theorem~\ref{thm:v2v66-elimination} says that
\eqref{eq:v2v66-fifth-map} vanishes on every zero relative class.
The universal property of a quotient therefore gives the unique factored
map.
\end{proof}

\subsection{Finite matrix construction}
\label{sec:v2v66-matrix}

Choose bases of
\(\mathcal W_n^{-1},\mathcal W_n^0,\mathcal W_n^1\) and of
\(\mathcal L_n^0\).  Let \(D_{-1}\) and \(D_0\) be the matrices of the
two relevant differentials, and let \(J_L\) be the inclusion matrix for
\(\mathcal L_n^0\).

\begin{proposition}[Row-reduction classifier]
\label{prop:v2v66-row-reduction}
The obstruction dimension is
\begin{equation}
 \dim\mathfrak O_n
 =
 \dim\ker D_0
 -
 \dim\left(
 \operatorname{im}D_{-1}
 +
 (\operatorname{im}J_L\cap\ker D_0)
 \right).
\label{eq:v2v66-dimension}
\end{equation}
A basis of \(\mathfrak O_n\) is obtained by:
\begin{enumerate}
 \item computing a basis of \(\ker D_0\);
 \item computing bases of
       \(\operatorname{im}D_{-1}\) and
       \(\operatorname{im}J_L\cap\ker D_0\);
 \item extending a basis of their sum to a basis of \(\ker D_0\).
\end{enumerate}
The added vectors are representatives of the relative obstruction classes.
\end{proposition}

\begin{proof}
By Lemma~\ref{lem:v2v66-zero},
\[
 \mathfrak O_n
 \cong
 \frac{\ker D_0}{
 \operatorname{im}D_{-1}
 +
 (\mathcal L_n^0\cap\ker D_0)}.
\]
The inclusion matrix identifies
\(\mathcal L_n^0\) with \(\operatorname{im}J_L\).
The dimension formula and algorithm are elementary quotient-space linear
algebra.
\end{proof}

\begin{assessment}[Why individual derivative allocations are not classified]
\label{ass:v2v66-packets}
A single allocation need not lie in \(\ker D_0\): its BRST variation may
be cancelled by a ghost or contact allocation.  The kernel in
Proposition~\ref{prop:v2v66-row-reduction} must therefore be computed
after the complete Ward matrix is assembled.  Intersecting a list of
factorwise bad terms with the kernel before their contact partners are
included can create false obstructions.
\end{assessment}

\subsection{A four-generator regression}
\label{sec:v2v66-regression}

Consider the finite complex with
\[
 \mathcal W^{-1}=\Bbbk u,
 \qquad
 \mathcal W^0=\Bbbk b\oplus\Bbbk p\oplus\Bbbk\ell,
 \qquad
 \mathcal W^1=\Bbbk a,
\]
and differential
\begin{equation}
 du=b,
 \qquad
 db=0,
 \qquad
 dp=a,
 \qquad
 d\ell=-a,
 \qquad
 da=0.
\label{eq:v2v66-toy-d}
\end{equation}
Let the local subcomplex be
\[
 \mathcal L^0=\Bbbk\ell,
 \qquad
 \mathcal L^1=\Bbbk a.
\]

\begin{proposition}[Contact completion can expose a genuine relative class]
\label{prop:v2v66-toy}
The element \(b\) is exact, \(\ell\) is local, and
\(p+\ell\) is closed.  Moreover,
\begin{equation}
 H^0(\mathcal W/\mathcal L,d)
 =
 \Bbbk[p+\mathcal L].
\label{eq:v2v66-toy-H}
\end{equation}
Thus the incomplete word \(p\) is not a cycle, while its contact-completed
packet \(p+\ell\) is a cycle with a nonzero relative class.
\end{proposition}

\begin{proof}
The first two statements follow from
\(du=b\) and the definition of \(\mathcal L\).  Also
\[
 d(p+\ell)=a-a=0.
\]
In the quotient, \(a\) and \(\ell\) vanish.  Hence both \(b\) and \(p\)
are quotient cycles, but only \(b=du\) is a boundary.  This proves
\eqref{eq:v2v66-toy-H}.
\end{proof}

\begin{assessment}[The candidate surviving packet]
\label{ass:v2v66-survivor}
The regression identifies the form of the first honest analytic target:
a contact-completed, BRST-closed degree-four packet whose class in
\(\mathfrak O_n\) is nonzero.  In the SM word complex, the relevant
candidate is the transverse physical completion of the fifth derivative
of the low Schur self-energy.  V66 does not assert that this candidate is
nonzero; that is decided only by the matrix classifier with the actual
regulated Ward rules.
\end{assessment}

\subsection{Uniform finite reduction}
\label{sec:v2v66-uniform}

Even when every \(\mathfrak O_n\) is finite-dimensional, a cutoff-uniform
analytic reduction needs uniformly controlled representatives.

\begin{theorem}[Uniform obstruction-basis compiler]
\label{thm:v2v66-uniform}
Suppose there are fixed \(N<\infty\), closed representatives
\(o_{1,n},\ldots,o_{N,n}\in\mathcal W_n^0\), and a constant \(C_{\rm dec}\)
such that every closed marginal packet \(w_n\) of formal norm at most one
admits
\begin{equation}
 w_n
 =
 \sum_{a=1}^Nc_{a,n}o_{a,n}
 +d_nu_n+\ell_n,
 \qquad
 \ell_n\in\mathcal L_n^0,
\label{eq:v2v66-uniform-decomposition}
\end{equation}
with
\begin{equation}
 \sum_{a=1}^N|c_{a,n}|
 \le C_{\rm dec}.
\label{eq:v2v66-coefficients}
\end{equation}
If
\begin{equation}
 \left\|
 D^5(q_{\rm phys}\mathcal E_no_{a,n})
 \right\|
 \le C_a\Lambda_n^{-1}
\label{eq:v2v66-representative-bound}
\end{equation}
for every \(a\), then every such \(w_n\) satisfies
\begin{equation}
 \left\|
 D^5(q_{\rm phys}\mathcal E_nw_n)
 \right\|
 \le
 C_{\rm dec}\max_aC_a\,\Lambda_n^{-1}.
\label{eq:v2v66-uniform-conclusion}
\end{equation}
\end{theorem}

\begin{proof}
Apply \(D^5q_{\rm phys}\mathcal E_n\) to
\eqref{eq:v2v66-uniform-decomposition}.  The exact and local terms vanish
by Theorem~\ref{thm:v2v66-elimination}.  Bound the remaining finite sum by
\eqref{eq:v2v66-coefficients} and
\eqref{eq:v2v66-representative-bound}.
\end{proof}

\begin{firewall}[Finite dimension is not uniform conditioning]
\label{fw:v2v66-conditioning}
The dimension formula
\eqref{eq:v2v66-dimension} does not imply
\eqref{eq:v2v66-coefficients}.  The angle between boundaries, local
cycles, and chosen obstruction representatives may collapse with the
cutoff.  This is the normed splitting issue already isolated in V3.
A continuum argument needs a uniformly bounded relative projection or an
explicit uniformly conditioned basis.
\end{firewall}

\subsection{Interface with the V65 derivative ledger}
\label{sec:v2v66-V65}

Let \(\mathcal W_n\) be the smallest \(d_n\)-invariant, contact-complete
degree-four subcomplex generated by packets that fail the V65 marginal
condition \eqref{eq:v2v65-marginal-condition}.  Ward partners needed for
closure remain in the complex even when their own ledger bounds are good.
Packets satisfying that condition
have already acquired the required \(\Lambda_n^{-1}\) fifth derivative
and may be placed directly in the good part of
\eqref{eq:v2v65-packets}.

\begin{theorem}[Algebraic reduction of the V60 bottleneck]
\label{thm:v2v66-V60-reduction}
Assume:
\begin{enumerate}
 \item every lower-degree or surplus-paid word satisfies the V65
       complete-word ledger;
 \item the remaining degree-four word complex is admissible;
 \item the uniform obstruction-basis hypothesis
       \eqref{eq:v2v66-uniform-decomposition}--%
       \eqref{eq:v2v66-coefficients} holds; and
 \item every obstruction representative satisfies
       \eqref{eq:v2v66-representative-bound}.
\end{enumerate}
Then the complete physical low--low row obeys
\begin{equation}
 \|D^5(q_{\rm phys}\mathscr L_n)\|
 \le C\Lambda_n^{-1},
\label{eq:v2v66-V60-bound}
\end{equation}
and its degree-four nonlocal Taylor remainder is dyadically summable.
\end{theorem}

\begin{proof}
The first hypothesis controls all V65 good words.  Apply
Theorem~\ref{thm:v2v66-uniform} to the remaining closed marginal packet
and add the finite bounds.  This proves
\eqref{eq:v2v66-V60-bound}.  Corollary~\ref{cor:v2v60-four-dimensional} gives the summable Taylor remainder.
\end{proof}

\begin{firewall}[The Standard Model matrix has not been evaluated]
\label{fw:v2v66-matrix-open}
V66 constructs the exact finite classifier and proves what each zero class
buys analytically.  It does not supply the actual matrices
\(D_{-1},D_0,J_L\) for all regulated Standard Model vertices, establish
their uniform conditioning, or estimate a nonzero obstruction
representative.  These remain finite but substantive algebraic and
analytic tasks.
\end{firewall}

\subsection{V66 conclusion and next acquisition}
\label{sec:v2v66-conclusion}

V66 removes every marginal fifth-derivative packet that is a BRST boundary
modulo a degree-four contact-local term.  At each fixed perturbative order it proves that the physical analytic
problem reduces to a finite one: estimate only a uniformly conditioned basis of
the relative obstruction space
\(\mathfrak O_n\).

V67 will evaluate the classifier on the minimal two-crossing Schur sector.
It will write the BRST differential of
\(X_nR_{n,L}Y_n\), include the endpoint contact completion, and determine
whether the transverse fifth-derivative packet is exact, local, or a
genuine relative class under explicit finite Ward identities.
\clearpage
\section*{Second Volume, Version V67: BRST Variation of the Two-Crossing Schur Packet}
\addcontentsline{toc}{section}{Second Volume, Version V67: BRST Variation of the Two-Crossing Schur Packet}

\subsection*{Purpose}

V66 constructed the relative cohomology that classifies marginal
fifth-derivative packets.  This note evaluates its first algebraic sector:
one shell-to-low crossing, one low inverse, and one low-to-shell crossing.
The exact BRST variation contains three endpoint defects.  A contact
completion cancels those defects and leaves only an output commutator.
After an invariant trace or physical contraction, the completed packet is
closed; after five external derivatives, its degree-four contact completion
vanishes.

This calculation identifies the remaining dichotomy.  The transverse
two-crossing packet is removed if it has a ghost-number-minus-one primitive.
It is a genuine relative obstruction if a linear topology witness
annihilates all boundaries and local terms but not this packet.  The note
states both tests without assuming which one the full regulated Standard
Model satisfies.

\subsection{Block Ward identities}
\label{sec:v2v67-block-Ward}

Let
\(\mathcal P_n\oplus\mathcal K_{n,L}\) be the transition-shell and low
decomposition.  Let
\[
 X_n:\mathcal K_{n,L}\to\mathcal P_n,
 \qquad
 H_n:\mathcal K_{n,L}\to\mathcal K_{n,L},
 \qquad
 Y_n:\mathcal P_n\to\mathcal K_{n,L}
\]
be \(C^5\) external-momentum families, with \(H_n(k)\) invertible, and put
\begin{equation}
 R_n=H_n^{-1},
 \qquad
 \Sigma_n=X_nR_nY_n.
\label{eq:v2v67-Sigma}
\end{equation}

Let \(c_{P,n}\) and \(c_{L,n}\) be \(k\)-independent ghost-number-one
generators on the two blocks.  Suppose the regulated block Ward identities
have the form
\begin{align}
 sX_n
 &=
 c_{P,n}X_n-X_nc_{L,n}+\xi_n,
\label{eq:v2v67-sX}\\
 sY_n
 &=
 c_{L,n}Y_n-Y_nc_{P,n}+\eta_n,
\label{eq:v2v67-sY}\\
 sH_n
 &=
 c_{L,n}H_n-H_nc_{L,n}+\zeta_n.
\label{eq:v2v67-sH}
\end{align}
The defect blocks \(\xi_n,\eta_n,\zeta_n\) contain the regulated endpoint,
measure, and contact contributions not represented by homogeneous block
covariance.

\begin{lemma}[BRST variation of the low inverse]
\label{lem:v2v67-sR}
Under
\eqref{eq:v2v67-sH},
\begin{equation}
 sR_n
 =
 c_{L,n}R_n-R_nc_{L,n}-R_n\zeta_nR_n.
\label{eq:v2v67-sR}
\end{equation}
\end{lemma}

\begin{proof}
Apply \(s\) to \(H_nR_n=I\):
\[
 (sH_n)R_n+H_n(sR_n)=0.
\]
Multiplication by \(R_n\) on the left gives
\(sR_n=-R_n(sH_n)R_n\).  Insert
\eqref{eq:v2v67-sH} and use
\(R_nH_n=H_nR_n=I\).
\end{proof}

\begin{theorem}[Exact two-crossing Ward defect]
\label{thm:v2v67-defect}
Define
\begin{equation}
 \Delta_n
 =
 \xi_nR_nY_n
 +
 X_nR_n\eta_n
 -
 X_nR_n\zeta_nR_nY_n.
\label{eq:v2v67-Delta}
\end{equation}
Then
\begin{equation}
 s\Sigma_n
 =
 c_{P,n}\Sigma_n-\Sigma_nc_{P,n}
 +
 \Delta_n.
\label{eq:v2v67-sSigma}
\end{equation}
\end{theorem}

\begin{proof}
All three factors in \(\Sigma_n\) have ghost number zero, so the graded
Leibniz rule has no intermediate sign:
\[
 s\Sigma_n
 =
 (sX_n)R_nY_n+X_n(sR_n)Y_n+X_nR_n(sY_n).
\]
Insert
\eqref{eq:v2v67-sX},
\eqref{eq:v2v67-sY}, and
\eqref{eq:v2v67-sR}.  The two terms containing
\(X_nc_{L,n}R_nY_n\) cancel, as do the two containing
\(X_nR_nc_{L,n}Y_n\).  The outer terms give the
\(c_{P,n}\)-commutator and the remaining terms are exactly
\eqref{eq:v2v67-Delta}.
\end{proof}

\begin{assessment}[The inverse defect is forced]
\label{ass:v2v67-inverse-defect}
The sign and the two inverse factors in
\(-X_nR_n\zeta_nR_nY_n\) are forced by differentiating an inverse.  A Ward
completion that includes endpoint defects for \(X_n\) and \(Y_n\) but
omits the \(H_n\)-defect is not closed.
\end{assessment}

\subsection{Contact completion and invariant contraction}
\label{sec:v2v67-completion}

Let \(K_n(k):\mathcal P_n\to\mathcal P_n\) be a ghost-number-zero contact
completion satisfying
\begin{equation}
 sK_n
 =
 -\Delta_n+
 c_{P,n}K_n-K_nc_{P,n}.
\label{eq:v2v67-sK}
\end{equation}
Let
\[
 \tau_{P,n}:
 \mathcal L(\mathcal P_n)\longrightarrow\mathcal Z
\]
be a bounded degree-preserving contraction which is a chain map and
annihilates output commutators:
\begin{align}
 s\,\tau_{P,n}(T)
 &=
 \tau_{P,n}(sT),
\label{eq:v2v67-tau-chain}\\
 \tau_{P,n}(c_{P,n}T-Tc_{P,n})
 &=0.
\label{eq:v2v67-tau-invariant}
\end{align}
An ordinary graded trace is the basic finite-dimensional example.

\begin{theorem}[The contact-completed Schur packet is closed]
\label{thm:v2v67-completed-cycle}
Define
\begin{equation}
 \mathscr P_n(k)
 =
 \tau_{P,n}\bigl(\Sigma_n(k)+K_n(k)\bigr).
\label{eq:v2v67-packet}
\end{equation}
Then
\begin{equation}
 s\mathscr P_n(k)=0.
\label{eq:v2v67-packet-closed}
\end{equation}
\end{theorem}

\begin{proof}
Equations
\eqref{eq:v2v67-sSigma} and
\eqref{eq:v2v67-sK} give
\[
 s(\Sigma_n+K_n)
 =
 c_{P,n}(\Sigma_n+K_n)
 -
 (\Sigma_n+K_n)c_{P,n}.
\]
Apply
\eqref{eq:v2v67-tau-chain} and
\eqref{eq:v2v67-tau-invariant}.
\end{proof}

\begin{corollary}[Five derivatives remove a degree-four contact completion]
\label{cor:v2v67-fifth-cycle}
Assume \(K_n(k)\) is a polynomial of total external degree at most four,
and assume \(s\) commutes with external differentiation.  Then
\begin{align}
 D^5\mathscr P_n
 &=
 D^5\tau_{P,n}(\Sigma_n),
\label{eq:v2v67-fifth-equality}\\
 sD^5\tau_{P,n}(\Sigma_n)
 &=0.
\label{eq:v2v67-fifth-closed}
\end{align}
\end{corollary}

\begin{proof}
The fifth derivative of \(K_n\) is zero, proving
\eqref{eq:v2v67-fifth-equality}.  Differentiate
\eqref{eq:v2v67-packet-closed} five times and use the commutation
hypothesis.
\end{proof}

\begin{firewall}[Existence of the contact completion is the local anomaly test]
\label{fw:v2v67-contact-existence}
Equation
\eqref{eq:v2v67-sK} is solvable in the allowed local space exactly when
the corresponding local ghost-number-one defect has a permitted primitive.
This is the shellwise repair problem of V3.  Standard Model anomaly
cancellation constrains the defect but does not, by itself, prove a
regulator-uniform solution including metric, boundary, antifield, and
normalization terms.
\end{firewall}

\subsection{Relative class of the two-crossing packet}
\label{sec:v2v67-relative}

Place the complete formal packet corresponding to
\(\mathscr P_n\) in the admissible complex
\((\mathcal W_n,\mathcal L_n)\) of V66, and denote it by \(p_n\).
The contact word \(K_n\) belongs to \(\mathcal L_n^0\).  Hence
\begin{equation}
 \mathfrak o_n(p_n)
 =
 [p_n+\mathcal L_n]
 =
 [\tau_{P,n}(\Sigma_n)+\mathcal L_n]
 \in\mathfrak O_n.
\label{eq:v2v67-relative-class}
\end{equation}

\begin{theorem}[Primitive criterion for removal]
\label{thm:v2v67-primitive}
The relative class of the packet is zero, and hence the packet is
eliminated algebraically by the V66 theorem, if and only if there are a
formal ghost-number-minus-one word
\(u_n\in\mathcal W_n^{-1}\) and a local word
\(\ell_n\in\mathcal L_n^0\) such that
\begin{equation}
 p_n=d_nu_n+\ell_n.
\label{eq:v2v67-primitive}
\end{equation}
\end{theorem}

\begin{proof}
This is the concrete zero-class criterion
Lemma~\ref{lem:v2v66-zero} applied to the closed packet \(p_n\), followed
by the elimination
Theorem~\ref{thm:v2v66-elimination}.
\end{proof}

For nontriviality, it is useful to avoid computing the entire quotient.

\begin{theorem}[Linear witness for a genuine relative obstruction]
\label{thm:v2v67-witness}
Let
\(\lambda_n:\mathcal W_n^0\to\Bbbk\) be linear and suppose
\begin{align}
 \lambda_n(d_nu)&=0
 &&(u\in\mathcal W_n^{-1}),
\label{eq:v2v67-witness-boundary}\\
 \lambda_n(\ell)&=0
 &&(\ell\in\mathcal L_n^0).
\label{eq:v2v67-witness-local}
\end{align}
If \(p_n\in\ker d_n\) and
\begin{equation}
 \lambda_n(p_n)\ne0,
\label{eq:v2v67-witness-nonzero}
\end{equation}
then
\begin{equation}
 \mathfrak o_n(p_n)\ne0.
\label{eq:v2v67-class-nonzero}
\end{equation}
\end{theorem}

\begin{proof}
If the class were zero, Lemma~\ref{lem:v2v66-zero} would give
\(p_n=d_nu+\ell\).  Applying \(\lambda_n\) and using
\eqref{eq:v2v67-witness-boundary}--\eqref{eq:v2v67-witness-local} would
give \(\lambda_n(p_n)=0\), contradicting
\eqref{eq:v2v67-witness-nonzero}.
\end{proof}

\begin{corollary}[Topology-coefficient test]
\label{cor:v2v67-topology}
Suppose the coefficient of the two-crossing skeleton defines a linear
functional \(\lambda_n\) on formal ghost-number-zero packets, no allowed
local word contains that skeleton, and the coefficient vanishes on every
BRST boundary.  If the completed packet has a nonzero coefficient, its
relative class is nonzero.
\end{corollary}

\begin{proof}
The three hypotheses are respectively
\eqref{eq:v2v67-witness-local},
\eqref{eq:v2v67-witness-boundary}, and
\eqref{eq:v2v67-witness-nonzero}.  Apply
Theorem~\ref{thm:v2v67-witness}.
\end{proof}

\begin{firewall}[Topology alone need not descend through a boundary]
\label{fw:v2v67-topology}
A BRST differential can preserve, split, contract, or create formal
propagator incidences through antifield and contact identities.  Therefore
the coefficient of a drawn two-crossing graph is a valid witness only
after
\eqref{eq:v2v67-witness-boundary} has been checked on the actual
ghost-number-minus-one differential matrix.  Visual graph distinction is
not a cohomology proof.
\end{firewall}

\subsection{The minimal finite matrix test}
\label{sec:v2v67-matrix}

Choose a basis whose first ghost-number-zero vector is the completed packet
\(p_n\).  Let \(D_{-1,n}\) and \(D_{0,n}\) be the matrices from V66, and
let \(J_{L,n}\) include the local words.

\begin{proposition}[Exact decision alternatives]
\label{prop:v2v67-decision}
Assume \(D_{0,n}p_n=0\).  Exactly one of the following holds:
\begin{enumerate}
 \item the coordinate vector of \(p_n\) belongs to
 \begin{equation}
  \operatorname{im}D_{-1,n}
  +
  \bigl(\operatorname{im}J_{L,n}\cap\ker D_{0,n}\bigr),
 \label{eq:v2v67-zero-span}
 \end{equation}
 in which case its relative class is zero;
 \item it does not belong to that span, in which case row reduction
 supplies a functional \(\lambda_n\) satisfying
 \eqref{eq:v2v67-witness-boundary}--\eqref{eq:v2v67-witness-nonzero}.
\end{enumerate}
\end{proposition}

\begin{proof}
The subspace in
\eqref{eq:v2v67-zero-span} is exactly the zero-class subspace of
\(\ker D_{0,n}\) by
Proposition~\ref{prop:v2v66-row-reduction}.  If \(p_n\) is not in that subspace, then it is not in the larger
subspace \(\operatorname{im}D_{-1,n}+\operatorname{im}J_{L,n}\).
Indeed, an identity \(p_n=D_{-1,n}u+J_{L,n}\ell\), together with
\(D_{0,n}p_n=0\) and \(D_{0,n}D_{-1,n}=0\), would force
\(J_{L,n}\ell\in\ker D_{0,n}\), contradicting the first assertion.
Finite-dimensional linear algebra now separates \(p_n\) from this larger
proper subspace.  The separating functional vanishes on every boundary and
every local word, and may be normalized to equal one on \(p_n\).
\end{proof}

\begin{assessment}[What has and has not been decided]
\label{ass:v2v67-decision}
Equations
\eqref{eq:v2v67-sSigma}--\eqref{eq:v2v67-fifth-closed} prove that the
contact-completed fifth derivative is a legitimate closed packet whenever
the local repair exists.  Proposition~\ref{prop:v2v67-decision} gives an
exact finite decision procedure for whether it is exact modulo local
terms.  The numerical decision still requires the actual regulated
Standard Model Ward matrices.
\end{assessment}

\subsection{Analytic consequence of a nonzero class}
\label{sec:v2v67-analytic}

If the class is zero, Theorem~\ref{thm:v2v66-elimination} removes the
packet.  If it is nonzero, the contact completion no longer contributes
after five derivatives, and the analytic representative is
\[
 D^5\tau_{P,n}(X_nR_nY_n).
\]

\begin{theorem}[Ledger for the surviving two-crossing representative]
\label{thm:v2v67-ledger}
Let \(x_{i,n},r_{j,n},y_{\ell,n}\) be the V64 derivative majorants for
\(X_n,R_n,Y_n\).  If
\begin{equation}
 x_{i,n}r_{j,n}y_{\ell,n}
 \le C_{i,j,\ell}
 \Lambda_n^{-4}
 \qquad
 (i+j+\ell=5),
\label{eq:v2v67-fifth-ledger}
\end{equation}
then
\begin{equation}
 \left\|
 D^5\tau_{P,n}(X_nR_nY_n)
 \right\|
 \le C\|\tau_{P,n}\|\Lambda_n^{-4}.
\label{eq:v2v67-fifth-bound}
\end{equation}
If this self-energy is inserted as an order-one effective-shell factor in
a complete degree-four V61 word, the remaining nominal word factors supply
the three powers encoded in its superficial degree, and the full fifth
derivative has scale \(O(\Lambda_n^{-1})\) under the V65 complete-word
ledger.
\end{theorem}

\begin{proof}
The exact multinomial formula
\eqref{eq:v2v64-multinomial} has 21 triples at total derivative order five.
Apply
\eqref{eq:v2v67-fifth-ledger} to every term and then the bounded map
\(\tau_{P,n}\).  The last statement is
Theorem~\ref{thm:v2v65-word-ledger} with total degree four; it records that
the self-energy bound is one factor of the complete word rather than a
standalone degree-four amplitude.
\end{proof}

\begin{firewall}[The displayed self-energy scale is contextual]
\label{fw:v2v67-context}
The exponent \(-4\) in
\eqref{eq:v2v67-fifth-ledger} is the fifth derivative of an order-one
effective-shell factor.  The V60 target \(-1\) concerns a complete
degree-four trace word.  Confusing these two degrees either demands three
unneeded inverse powers from the full word or loses three powers inside
the effective resolvent.
\end{firewall}

\subsection{V67 conclusion and next acquisition}
\label{sec:v2v67-conclusion}

V67 computes the exact BRST defect of the minimal low Schur excursion,
including the inverse-defect term with two low resolvents.  It proves that
a degree-four contact completion makes the invariant packet closed and
then disappears from its fifth derivative.  It also supplies exact
primitive and witness tests for the packet's relative cohomology class.

The next step does not require guessing that class.  V68 will decompose the
completed packet into longitudinal and transverse output sectors under a
fixed physical Hodge projector.  It will prove which longitudinal
fifth-derivative terms are BRST exact under an explicit contracting
homotopy and isolate the transverse analytic representative that cannot be
removed by gauge descent.
\clearpage
\section*{Second Volume, Version V68: Hodge Contraction of the Longitudinal Fifth-Derivative Packet}
\addcontentsline{toc}{section}{Second Volume, Version V68: Hodge Contraction of the Longitudinal Fifth-Derivative Packet}

\subsection*{Purpose}

V67 proved that the contact-completed two-crossing Schur packet is BRST
closed and that its degree-four contact term disappears after five
external derivatives.  This note separates the resulting closed element
into its longitudinal exact part and its transverse physical part.

The separation requires an actual contracting homotopy.  At finite
regulator it is supplied by the Green operator of the BRST Laplacian when
that Laplacian has a gap away from its harmonic subspace.  A collapsing
gap gives an explicit obstruction to uniform control.  Under the gap
hypothesis, only the harmonic projection of the V67 fifth derivative
requires an analytic \(\Lambda_n^{-1}\) estimate in physical cohomology.

\subsection{Finite Hilbert BRST complex}
\label{sec:v2v68-Hilbert}

Let
\[
 \mathcal K^\bullet
 =
 \bigoplus_g\mathcal K^g
\]
be a finite-dimensional graded Hilbert space and let
\[
 s:\mathcal K^g\longrightarrow\mathcal K^{g+1},
 \qquad s^2=0.
\]
Write \(s^*\) for the Hilbert adjoint and define the BRST Laplacian
\begin{equation}
 \Delta_s=ss^*+s^*s.
\label{eq:v2v68-Laplacian}
\end{equation}
Let \(\Pi\) be the orthogonal projection onto
\(\ker\Delta_s\).

\begin{lemma}[Harmonic subspace]
\label{lem:v2v68-harmonic}
One has
\begin{equation}
 \ker\Delta_s
 =
 \ker s\cap\ker s^*.
\label{eq:v2v68-harmonic}
\end{equation}
Every BRST cohomology class has a unique representative in
\(\ker\Delta_s\).
\end{lemma}

\begin{proof}
For every \(z\),
\[
 \langle z,\Delta_sz\rangle
 =
 \|sz\|^2+\|s^*z\|^2.
\]
This proves
\eqref{eq:v2v68-harmonic}.  Finite-dimensional orthogonal decomposition
gives
\[
 \mathcal K^g
 =
 \operatorname{im}s
 \oplus\ker\Delta_s
 \oplus\operatorname{im}s^*.
\]
Indeed the first two summands lie in \(\ker s\), their orthogonal
complement is \(\operatorname{im}s^*\), and
\(\operatorname{im}s\) is orthogonal to \(\ker\Delta_s\).
A closed element has no \(\operatorname{im}s^*\) component, because if
\(z=su+h+s^*v\) and \(sz=0\), then
\[
 0=\langle v,s(s^*v)\rangle=\|s^*v\|^2.
\]
Thus \(z=su+h\).  The harmonic representative is unique because a harmonic
boundary \(su\) has
\(\|su\|^2=\langle s^*su,u\rangle=0\).
\end{proof}

\subsection{Green operator and contracting homotopy}
\label{sec:v2v68-Green}

Assume that the nonzero spectrum of \(\Delta_s\) is bounded below by
\begin{equation}
 \Delta_s\ge\gamma^2(I-\Pi)
\label{eq:v2v68-gap}
\end{equation}
for some \(\gamma>0\).  Define
\begin{equation}
 G_s
 =
 \Delta_s^{-1}(I-\Pi),
 \qquad
 h=s^*G_s.
\label{eq:v2v68-h}
\end{equation}

\begin{theorem}[Hodge contracting identity]
\label{thm:v2v68-contract}
Under
\eqref{eq:v2v68-gap},
\begin{align}
 sh+hs
 &=
 I-\Pi,
\label{eq:v2v68-contract}\\
 \|G_s\|
 &\le\gamma^{-2},
\label{eq:v2v68-G-bound}\\
 \|h\|
 &\le\gamma^{-1}.
\label{eq:v2v68-h-bound}
\end{align}
Consequently, for every \(z\in\ker s\),
\begin{equation}
 z=\Pi z+s(hz).
\label{eq:v2v68-closed-split}
\end{equation}
\end{theorem}

\begin{proof}
The Laplacian commutes with \(s\) and \(s^*\):
\[
 s\Delta_s=ss^*s=\Delta_ss,
 \qquad
 s^*\Delta_s=s^*ss^*=\Delta_ss^*,
\]
where \(s^2=(s^*)^2=0\).  Hence \(G_s\) commutes with both on the
orthogonal complement of the harmonic space.  Therefore
\[
 sh+hs
 =
 ss^*G_s+s^*G_ss
 =
 (ss^*+s^*s)G_s
 =
 I-\Pi.
\]
The spectral theorem and
\eqref{eq:v2v68-gap} give
\eqref{eq:v2v68-G-bound}.  Moreover,
\begin{align*}
 \|s^*G_sz\|^2
 &=
 \langle G_sz,ss^*G_sz\rangle\\
 &\le
 \langle G_sz,\Delta_sG_sz\rangle
 =
 \langle (I-\Pi)z,G_sz\rangle
 \le
 \gamma^{-2}\|z\|^2,
\end{align*}
which proves
\eqref{eq:v2v68-h-bound}.  If \(sz=0\), apply
\eqref{eq:v2v68-contract} to \(z\); the term \(hsz\) vanishes.
\end{proof}

\begin{assessment}[Algebraic exactness versus uniform descent]
\label{ass:v2v68-uniform}
For a fixed regulator, every nonharmonic closed element is exact by
\eqref{eq:v2v68-closed-split}.  A continuum construction needs the
constants to remain controlled.  The quantitative condition is a
cutoff-uniform lower bound for \(\gamma\), or another explicitly bounded
contracting homotopy.  Finite-dimensional Hodge theory alone supplies no
such bound.
\end{assessment}

\subsection{Sharp failure when the Hodge gap collapses}
\label{sec:v2v68-collapse}

\begin{counterexample}[Exact longitudinal vectors with divergent primitives]
\label{sep:v2v68-collapse}
Let
\[
 \mathcal K_n^{-1}=\Bbbk u,
 \qquad
 \mathcal K_n^0=\Bbbk v,
 \qquad
 s_nu=\varepsilon_nv,
 \qquad
 s_nv=0,
\]
where \(\varepsilon_n>0\) and \(\varepsilon_n\to0\), and give the displayed
bases unit norm.  Then \(v\) is exact for every \(n\), but its least-norm
primitive is \(\varepsilon_n^{-1}u\).  The nonzero eigenvalue of
\(\Delta_{s_n}\) is \(\varepsilon_n^2\), and
\begin{equation}
 \|h_n\|=\varepsilon_n^{-1}\longrightarrow\infty.
\label{eq:v2v68-collapse}
\end{equation}
\end{counterexample}

\begin{proof}
The adjoint satisfies \(s_n^*v=\varepsilon_nu\).  Hence
\(\Delta_{s_n}\) is multiplication by \(\varepsilon_n^2\) in both degrees.
The Green operator is multiplication by
\(\varepsilon_n^{-2}\), so
\(h_nv=s_n^*G_nv=\varepsilon_n^{-1}u\).
\end{proof}

\begin{firewall}[A zero physical class need not give a bounded gauge representative]
\label{fw:v2v68-zero-class}
The physical quotient sends every \(v=s_n(\varepsilon_n^{-1}u)\) in the
counterexample to zero, but constructing compatible gauge-fixed
representatives or summable chain homotopies still sees the divergent
primitive.  Quotient cancellation is sufficient for the physical
fifth-derivative norm; it is not sufficient for the V54 regulator-chart
homotopy.
\end{firewall}

\subsection{External derivatives and a fixed Hodge projector}
\label{sec:v2v68-external}

Let \(B_\kappa\subset\mathbb R^q\), and suppose
\[
 z_n:B_\kappa\longrightarrow\mathcal K_n^0
\]
is \(C^5\) and \(s_nz_n(k)=0\).  Assume \(s_n,\Pi_n,h_n\) do not depend
on the external momentum \(k\).

\begin{theorem}[Longitudinal fifth derivatives are exact]
\label{thm:v2v68-longitudinal}
For every \(k\),
\begin{align}
 D^5z_n(k)
 &=
 \Pi_nD^5z_n(k)
 +
 s_nh_nD^5z_n(k),
\label{eq:v2v68-fifth-split}\\
 (I-\Pi_n)D^5z_n(k)
 &=
 s_nh_nD^5z_n(k).
\label{eq:v2v68-longitudinal}
\end{align}
Thus the longitudinal fifth derivative is a BRST boundary, and the
physical class satisfies
\begin{equation}
 q_{\rm phys}D^5z_n(k)
 =
 q_{\rm phys}\Pi_nD^5z_n(k).
\label{eq:v2v68-physical-transverse}
\end{equation}
\end{theorem}

\begin{proof}
External differentiation commutes with the fixed differential, so
\(s_nD^5z_n=0\).  Apply
\eqref{eq:v2v68-closed-split} to \(D^5z_n\).  The quotient map kills the
exact second term.
\end{proof}

\begin{corollary}[Only the harmonic norm is needed in the physical quotient]
\label{cor:v2v68-harmonic-bound}
If
\begin{equation}
 \|\Pi_nD^5z_n(k)\|
 \le C\Lambda_n^{-1}
\label{eq:v2v68-harmonic-bound}
\end{equation}
uniformly on \(B_\kappa\), then
\begin{equation}
 \|q_{\rm phys}D^5z_n(k)\|
 \le
 \|q_{\rm phys}\|C\Lambda_n^{-1}.
\label{eq:v2v68-quotient-bound}
\end{equation}
No separate norm estimate on
\((I-\Pi_n)D^5z_n\) is needed for this quotient inequality.
\end{corollary}

\begin{proof}
Use
\eqref{eq:v2v68-physical-transverse} and boundedness of
\(q_{\rm phys}\).
\end{proof}

\begin{firewall}[A momentum-dependent transverse projector changes the formula]
\label{fw:v2v68-moving-Hodge}
If \(\Pi_n=\Pi_n(k)\), then
\[
 D^5(\Pi_nz_n)
 =
 \sum_{j=0}^5\binom5j
 (D^j\Pi_n)(D^{5-j}z_n),
\]
with the natural symmetrization in several variables.  The usual Fourier
projector containing \(k\otimes k/|k|^2\) is also singular at \(k=0\).
Theorem~\ref{thm:v2v68-longitudinal} therefore uses a fixed BRST Hodge
projector on the regulated complex, not an unexamined momentum-dependent
polarization formula.
\end{firewall}

\subsection{Application to the V67 Schur packet}
\label{sec:v2v68-V67}

Let
\[
 z_n(k)
 =
 \tau_{P,n}\bigl(\Sigma_n(k)+K_n(k)\bigr)
\]
be the completed cycle from
\eqref{eq:v2v67-packet}, and suppose \(K_n\) has external degree at most
four.  Corollary~\ref{cor:v2v67-fifth-cycle} gives
\begin{equation}
 D^5z_n
 =
 D^5\tau_{P,n}(X_nR_nY_n).
\label{eq:v2v68-V67-fifth}
\end{equation}

\begin{theorem}[Physical reduction of the two-crossing packet]
\label{thm:v2v68-V67-reduction}
Under the fixed Hodge hypotheses,
\begin{equation}
 q_{\rm phys}D^5z_n
 =
 q_{\rm phys}\Pi_n
 D^5\tau_{P,n}(X_nR_nY_n).
\label{eq:v2v68-V67-reduction}
\end{equation}
If
\begin{equation}
 \left\|
 \Pi_nD^5\tau_{P,n}(X_nR_nY_n)
 \right\|
 \le C\Lambda_n^{-1},
\label{eq:v2v68-transverse-target}
\end{equation}
then the completed two-crossing packet meets the V60 fifth-derivative
target in physical cohomology.
\end{theorem}

\begin{proof}
Apply
Theorem~\ref{thm:v2v68-longitudinal} to \(z_n\), then use
\eqref{eq:v2v68-V67-fifth}.  The bound follows from
Corollary~\ref{cor:v2v68-harmonic-bound}.
\end{proof}

\begin{assessment}[Relation to the relative obstruction class]
\label{ass:v2v68-relative}
A nonzero class in the formal relative space \(\mathfrak O_n\) need not
produce a nonzero physical fifth derivative: it may lie in the kernel of
the evaluated map
\(\mathfrak D_n^{(5)}\) from V66, or its evaluated derivative may be
BRST exact.  Before passing to the quotient, the Hodge representative
defines the further factored map
\begin{equation}
 \mathfrak H_n^{(5)}:\mathfrak O_n
 \longrightarrow\ker\Delta_{s_n}\cap\mathcal K_n^0,
 \qquad
 \mathfrak H_n^{(5)}([w])
 =
 \Pi_nD^5\mathcal E_nw.
\label{eq:v2v68-harmonic-map}
\end{equation}
This is well defined: a local change has zero fifth derivative, and an
exact change has zero harmonic projection.  The analytic representatives
are therefore a basis of
\(\operatorname{Ran}\mathfrak H_n^{(5)}\), which can be smaller than
\(\mathfrak O_n\).
\end{assessment}

\subsection{Uniform physical-image compiler}
\label{sec:v2v68-image}

\begin{theorem}[Estimate only a basis of the harmonic image]
\label{thm:v2v68-image}
Suppose the hypotheses of the V66 uniform obstruction-basis theorem hold,
and suppose its evaluated representatives take values in
\(\mathcal K_n^0\).  Assume additionally that their harmonic images admit
a fixed finite spanning family
\(v_{1,n},\ldots,v_{r,n}\subset\ker\Delta_{s_n}\) such that
\begin{align}
 \Pi_nD^5(\mathcal E_no_{a,n})
 &=
 \sum_{b=1}^r\mu_{ab,n}v_{b,n},
\label{eq:v2v68-image-decomposition}\\
 \sum_{b=1}^r|\mu_{ab,n}|
 &\le C_{\rm img},
\label{eq:v2v68-image-coefficients}\\
 \|v_{b,n}\|
 &\le C_b\Lambda_n^{-1}.
\label{eq:v2v68-image-bounds}
\end{align}
Then every normalized closed marginal packet has physical fifth derivative
bounded by \(C\Lambda_n^{-1}\).
\end{theorem}

\begin{proof}
Use the V66 uniformly bounded decomposition into obstruction
representatives, exact terms, and local terms.  The latter two vanish after
the physical fifth-derivative map.  Apply
\eqref{eq:v2v68-image-decomposition} to the representatives and bound the
two finite coefficient sums by the V66 decomposition constant and
\eqref{eq:v2v68-image-coefficients}.  Finally use
\eqref{eq:v2v68-image-bounds} and apply the bounded physical quotient map.
\end{proof}

\begin{firewall}[A Hodge projector does not prove the transverse estimate]
\label{fw:v2v68-transverse-open}
The contraction theorem removes the longitudinal boundary exactly, but it
does not imply
\eqref{eq:v2v68-transverse-target}.  The harmonic component is the
physical response and can carry the full marginal degree.  It must be
estimated by the resolvent-word, spectral-gap, and complete-word machinery
of V61--V65.
\end{firewall}

\subsection{V68 conclusion and next acquisition}
\label{sec:v2v68-conclusion}

V68 proves the longitudinal/transverse decomposition of the closed V67
fifth derivative from a quantitative BRST Hodge gap.  The longitudinal
part is exact; the contact-local part already vanished; and the only
physical analytic target is the harmonic projection of the two-crossing
resolvent derivative.  A collapsing BRST gap is recorded as a separate
uniformity obstruction.

V69 will attack
\eqref{eq:v2v68-transverse-target}.  It will expand the 21 fifth-derivative
Schur allocations after harmonic projection and use the V63 gap estimate
to close every allocation except those with at least two derivatives on
the low inverse.  Those residual allocations will be grouped before norms
to test whether the transverse Ward identity supplies an additional
resolvent difference.
\clearpage
\section*{Second Volume, Version V69: Two-Scale Closure of the Minimal Low Schur Packet}
\addcontentsline{toc}{section}{Second Volume, Version V69: Two-Scale Closure of the Minimal Low Schur Packet}

\subsection*{Purpose}

V68 isolated the harmonic fifth derivative of the two-crossing Schur
packet.  A one-scale estimate treats every low inverse as if it carried
only the shell scale \(\Lambda_n\), and then derivatives of that inverse
appear deficient.  The correct upstream decomposition retains the actual
low frequency \(\mu\).

For a dyadic low band of frequency \(\mu\), a spectral gap from the
transition shell gives one factor \(\Lambda_n^{-1}\) at each crossing.
The differentiated low inverse has its own elliptic scale
\(\mu^{-1-a}\).  Five derivatives of the three-factor excursion therefore
give
\[
 \Lambda_n^{-2}\mu^{-4}.
\]
Every term factors through the low band, whose rank is \(O(\mu^4)\) in
four dimensions.  The powers of \(\mu\) cancel.  Summing the logarithmically
many low bands gives
\(\Lambda_n^{-2}(1+\log\Lambda_n)\), which is stronger than the V60
target \(O(\Lambda_n^{-1})\).

\subsection{Two spectral scales}
\label{sec:v2v69-scales}

Let \(D_0=D_0^*\) at finite regulator.  Let \(P_n\) be the signed
transition-shell projection of V63 at scale \(\Lambda_n\).  For dyadic
\(\mu\) with
\begin{equation}
 1\le\mu,
 \qquad
 b\mu\le c_L\Lambda_n,
 \qquad
 c_L<c_-,
\label{eq:v2v69-range}
\end{equation}
let \(Q_\mu\) be a spectral projection of \(D_0\) supported in
\begin{equation}
 \{a\mu\le |D_0|\le b\mu\}
\label{eq:v2v69-low-band}
\end{equation}
for fixed \(0<a<b\), with the endpoint bands modified harmlessly to form
a finite partition.  Assume
\begin{equation}
 \operatorname{rank}Q_\mu
 \le W\mu^d.
\label{eq:v2v69-rank}
\end{equation}
Because of
\eqref{eq:v2v69-range}, the distance between the spectra on \(P_n\) and
\(Q_\mu\) is at least
\begin{equation}
 \delta\Lambda_n,
 \qquad
 \delta=c_--c_L>0.
\label{eq:v2v69-gap}
\end{equation}

Let
\[
 X_{n,\mu}=P_nV_nQ_\mu,
 \qquad
 Y_{n,\mu}=Q_\mu V_nP_n.
\]

\begin{lemma}[Band-sensitive gap crossing]
\label{lem:v2v69-crossing}
Suppose, for \(0\le a_0\le5\),
\begin{align}
 \|P_{n,\sigma}[D_0,D^{a_0}V_n(k)]Q_\mu\|_{\mathcal L^{a_0}}
 &\le c_{a_0}\mu^{1-a_0},
\label{eq:v2v69-comm-X}\\
 \|Q_\mu[D_0,D^{a_0}V_n(k)]P_{n,\sigma}\|_{\mathcal L^{a_0}}
 &\le c_{a_0}\mu^{1-a_0}
\label{eq:v2v69-comm-Y}
\end{align}
for \(\sigma\in\{+,-\}\).  Then
\begin{align}
 \|D^{a_0}X_{n,\mu}(k)\|_{\mathcal L^{a_0}}
 &\le C_{a_0}\Lambda_n^{-1}\mu^{1-a_0},
\label{eq:v2v69-X}\\
 \|D^{a_0}Y_{n,\mu}(k)\|_{\mathcal L^{a_0}}
 &\le C_{a_0}\Lambda_n^{-1}\mu^{1-a_0}.
\label{eq:v2v69-Y}
\end{align}
\end{lemma}

\begin{proof}
Repeat the positive- and negative-shell Sylvester integrals of
Lemmas~\ref{lem:v2v63-positive} and
\ref{lem:v2v63-negative}.  The spectral separation is now
\eqref{eq:v2v69-gap}, while the right side of the Sylvester equation is
bounded by
\eqref{eq:v2v69-comm-X} or
\eqref{eq:v2v69-comm-Y}.  Integration supplies
\((\delta\Lambda_n)^{-1}\).  The fixed positive and negative shell sum is
absorbed into \(C_{a_0}\).
\end{proof}

\begin{assessment}[The commutator is measured at the low frequency]
\label{ass:v2v69-low-frequency}
The right side of
\eqref{eq:v2v69-comm-X} is \(\mu^{1-a_0}\), not
\(\Lambda_n^{1-a_0}\).  This is justified only when the compressed
commutator acts as an order-one differentiated symbol on the input band
\(Q_\mu\).  A global shell norm would erase the gain central to V69.
\end{assessment}

\subsection{Differentiated inverse on one low band}
\label{sec:v2v69-low-inverse}

Let
\[
 H_{n,\mu}(k):
 Q_\mu\mathcal H\longrightarrow Q_\mu\mathcal H
\]
be \(C^5\) and invertible, and put
\(R_{n,\mu}=H_{n,\mu}^{-1}\).

\begin{lemma}[Low-band inverse calculus]
\label{lem:v2v69-low-inverse}
Suppose
\begin{align}
 \|R_{n,\mu}(k)\|
 &\le r_0\mu^{-1},
\label{eq:v2v69-R0}\\
 \|D^{a_0}H_{n,\mu}(k)\|_{\mathcal L^{a_0}}
 &\le h_{a_0}\mu^{1-a_0},
 \qquad 1\le a_0\le5.
\label{eq:v2v69-H}
\end{align}
Then
\begin{equation}
 \|D^{a_0}R_{n,\mu}(k)\|_{\mathcal L^{a_0}}
 \le r_{a_0}\mu^{-1-a_0},
 \qquad 0\le a_0\le5.
\label{eq:v2v69-R}
\end{equation}
\end{lemma}

\begin{proof}
Apply the ordered-partition inverse formula of V61 with scale \(\mu\).
A term with \(j\) differentiated \(H_{n,\mu}\)-blocks of total derivative
order \(a_0\) contains \(j+1\) inverse factors.  Its exponent is
\[
 -(j+1)+\sum_{\ell=1}^j(1-|B_\ell|)
 =
 -1-a_0.
\]
The number of partitions is finite through order five.
\end{proof}

\begin{firewall}[Compression and inversion do not commute]
\label{fw:v2v69-compressed-inverse}
The inverse in
\eqref{eq:v2v69-R0} is the inverse of the actual reduced low-band kinetic
block.  In general,
\[
 (Q_\mu A_nQ_\mu)^{-1}
 \ne
 Q_\mu A_n^{-1}Q_\mu.
\]
If \(A_n\) mixes different low bands, the correct object is a Schur
complement or a weighted full low resolvent.  The block-diagonal regression
below does not suppress that issue.
\end{firewall}

\subsection{Fifth derivative of one band excursion}
\label{sec:v2v69-band}

Define
\begin{equation}
 \Sigma_{n,\mu}(k)
 =
 X_{n,\mu}(k)R_{n,\mu}(k)Y_{n,\mu}(k).
\label{eq:v2v69-Sigma}
\end{equation}

\begin{theorem}[Two-scale fifth-derivative bound]
\label{thm:v2v69-band-bound}
Under
\eqref{eq:v2v69-X},
\eqref{eq:v2v69-Y}, and
\eqref{eq:v2v69-R},
\begin{equation}
 \|D^5\Sigma_{n,\mu}(k)\|_{\mathcal L^5}
 \le
 C\Lambda_n^{-2}\mu^{-4}.
\label{eq:v2v69-band-bound}
\end{equation}
More generally, for \(0\le m\le5\),
\begin{equation}
 \|D^m\Sigma_{n,\mu}(k)\|_{\mathcal L^m}
 \le
 C_m\Lambda_n^{-2}\mu^{1-m}.
\label{eq:v2v69-band-general}
\end{equation}
\end{theorem}

\begin{proof}
Use the exact assignment formula
\eqref{eq:v2v62-assignment}.  A term with derivative orders
\(i,j,\ell\) and \(i+j+\ell=m\) has norm at most
\begin{align*}
 &C\bigl(\Lambda_n^{-1}\mu^{1-i}\bigr)
 \mu^{-1-j}
 \bigl(\Lambda_n^{-1}\mu^{1-\ell}\bigr)\\
 &\qquad=
 C\Lambda_n^{-2}\mu^{1-m}.
\end{align*}
Sum the finitely many assignments.  Setting \(m=5\) gives
\eqref{eq:v2v69-band-bound}.
\end{proof}

\begin{assessment}[Resolution of the V64 apparent deficit]
\label{ass:v2v69-V64}
If every factor is measured only against \(\Lambda_n\), a bounded
\(D^jR_{n,L}\) creates the V64 deficit \(\ell_j=j\).  On a
\(\mu\)-band, Lemma~\ref{lem:v2v69-low-inverse} supplies
\(\mu^{-1-j}\), while Lemma~\ref{lem:v2v69-crossing} retains the two
separate inverse shell gaps.  The deficit was therefore partly an artifact
of collapsing two spectral scales into one norm.
\end{assessment}

\subsection{Low-rank trace payer}
\label{sec:v2v69-trace}

Every derivative term in
\(\Sigma_{n,\mu}\) factors through \(Q_\mu\mathcal H\).  Hence its rank is
at most \(\operatorname{rank}Q_\mu\).

\begin{theorem}[Four-dimensional bandwise trace cancellation]
\label{thm:v2v69-trace}
Let \(\operatorname{Tr}_{P_n}\) be the finite-regulator trace on the shell
block.  Under
\eqref{eq:v2v69-rank} and
Theorem~\ref{thm:v2v69-band-bound},
\begin{equation}
 \left|
 D^5\operatorname{Tr}_{P_n}\Sigma_{n,\mu}(k)
 \right|
 \le
 CW\Lambda_n^{-2}\mu^{d-4}.
\label{eq:v2v69-trace}
\end{equation}
For \(d=4\), every dyadic low band contributes at most
\(CW\Lambda_n^{-2}\).
\end{theorem}

\begin{proof}
For a finite-rank operator \(T\),
\[
 |\operatorname{Tr}T|
 \le\|T\|_1
 \le\operatorname{rank}(T)\|T\|.
\]
Apply this inequality separately to every term of the exact fifth
derivative.  Each term factors as an operator from \(P_n\) through
\(Q_\mu\) and back, so its rank is at most
\(\operatorname{rank}Q_\mu\).  Insert
\eqref{eq:v2v69-rank} and
\eqref{eq:v2v69-band-bound}.
\end{proof}

\begin{counterexample}[The shell rank loses the two-scale gain]
\label{sep:v2v69-shell-rank}
Replacing the factorization rank in the proof by the crude shell estimate
\(\operatorname{rank}P_n\lesssim\Lambda_n^4\) gives only
\[
 \Lambda_n^4
 \Lambda_n^{-2}\mu^{-4}
 =
 \Lambda_n^2\mu^{-4},
\]
which is not uniformly \(O(\Lambda_n^{-1})\) on low bands.  Thus the low
rank is a genuine payer, not a cosmetic improvement of the constant.
\end{counterexample}

\begin{proof}
The displayed multiplication is the crude trace-norm estimate.  For fixed
\(\mu\), it grows as \(\Lambda_n^2\).
\end{proof}

\subsection{Summation over the low bands}
\label{sec:v2v69-sum}

Assume for the moment that the low block reduces the dyadic projections,
so that
\begin{equation}
 \Sigma_{n,L}
 =
 \sum_{\substack{\mu\ {\rm dyadic}\\
                  1\le\mu,\ b\mu\le c_L\Lambda_n}}
 \Sigma_{n,\mu}.
\label{eq:v2v69-band-sum}
\end{equation}

\begin{theorem}[Minimal two-crossing packet closes in four dimensions]
\label{thm:v2v69-close}
Let \(d=4\), and assume all constants in
Lemmas~\ref{lem:v2v69-crossing} and
\ref{lem:v2v69-low-inverse} are uniform in \(\mu,n\).
Then
\begin{align}
 \left|
 D^5\operatorname{Tr}_{P_n}\Sigma_{n,L}(k)
 \right|
 &\le
 C\Lambda_n^{-2}(1+\log\Lambda_n)
\label{eq:v2v69-log}\\
 &\le
 C'\Lambda_n^{-1}.
\label{eq:v2v69-target}
\end{align}
For dyadic \(\Lambda_n\simeq2^n\), the sharper first line is absolutely
summable in \(n\).
\end{theorem}

\begin{proof}
There are at most \(C(1+\log\Lambda_n)\) dyadic bands in
\eqref{eq:v2v69-band-sum}.  Sum the \(d=4\) estimate from
Theorem~\ref{thm:v2v69-trace}.  Since
\((1+\log\Lambda)/\Lambda\) is bounded for \(\Lambda\ge1\), the first line
implies the second.  For dyadic scales,
\[
 \sum_n2^{-2n}(1+n)<\infty.
\]
\end{proof}

\begin{corollary}[Harmonic projection does not spoil the bound]
\label{cor:v2v69-harmonic}
Let \(\Pi_n\) be the V68 orthogonal harmonic projector acting on a finite
tensor coefficient space after the shell trace, and let the remaining
fixed contraction have norm at most \(C_\tau\).  Then the harmonic
minimal packet obeys the same bound
\[
 \left\|
 \Pi_nD^5\tau_n(\Sigma_{n,L})
 \right\|
 \le
 C_\tau C\Lambda_n^{-2}(1+\log\Lambda_n).
\]
It therefore satisfies
\eqref{eq:v2v68-transverse-target}.
\end{corollary}

\begin{proof}
Orthogonal projection has norm one.  Apply the bounded finite contraction
after
\eqref{eq:v2v69-log}.
\end{proof}

\begin{firewall}[Only the minimal block-diagonal packet closes]
\label{fw:v2v69-scope}
Theorem~\ref{thm:v2v69-close} treats a single low excursion after a
block-diagonal dyadic reduction and a finite shell trace.  It does not
control:
\begin{enumerate}
 \item mixing between distinct low bands in the full inverse;
 \item massless zero modes or a continuum infrared threshold below
       \(\mu=1\);
 \item additional unbounded composite vertices in the same Ward word;
 \item uniform summation over perturbative order; or
 \item the local relative-cohomology and Hodge-gap hypotheses of V66--V68.
\end{enumerate}
\end{firewall}

\subsection{V69 conclusion and next acquisition}
\label{sec:v2v69-conclusion}

V69 proves the harmonic fifth-derivative target for the minimal
two-crossing packet in a four-dimensional, dyadically reduced low block.
The proof retains both spectral scales: two shell-gap powers, five
low-frequency derivatives, and the low-band rank cancel to a summable
\(\Lambda_n^{-2}\log\Lambda_n\) bound.

V70 is a mandatory companion audit.  After inspecting the current NS and
YM notes, it will choose between the two remaining upstream obstructions:
off-diagonal mixing among low dyadic bands and the infrared block below the
first positive low frequency.  The choice will be made from any new
parallel result rather than from the one-scale ledger.
\clearpage
\section*{Second Volume, Version V70: Companion Audit and Assembled Mixed-Band Low Resolvent}
\addcontentsline{toc}{section}{Second Volume, Version V70: Companion Audit and Assembled Mixed-Band Low Resolvent}

\subsection*{Purpose and mandatory audit}

V69 closed the minimal two-crossing packet when the low kinetic operator
reduces a dyadic band decomposition.  The mandatory V70 audit shows why
the next step must retain the full low inverse before taking norms.

The Navier--Stokes file inspected was
\[
 \texttt{Renewal\_NS\_Stability\_Research\_Notes\_V31.tex},
\]
with \(887537\) bytes and SHA--256 digest
\[
 \texttt{F1121B62238AD5491BB7CF03635EA9934942EDF573ED8FAAA9EE79E1D38A6696}.
\]
It remains unchanged.  Its endpoint lesson continues to apply: a missing
weighted moment must be paid by an estimate in the norm where it is used.

The active Yang--Mills file inspected was
\[
 \texttt{Renewal\_Yang\_Mills\_Mass\_Gap\_Research\_Notes\_V59.tex},
\]
with \(796599\) bytes and SHA--256 digest
\[
 \texttt{64AB78AB9CEED172909A3BEBEEDFE6E4DD50E09529286982125B72A86728F1D0}.
\]
The frozen first YM volume
\[
 \texttt{Renewal\_Yang\_Mills\_Mass\_Gap\_Research\_Notes\_V135\_f1.tex}
\]
was present with \(3083467\) bytes and digest
\[
 \texttt{82F6BBFFEECDAD3C5C63551660B19A0BDC7F3D6DFFE4C68B4E9BBA3A783A2B71};
\]
it was left untouched.

YM V59 proves an exact endpoint-covariance leakage in individual graph
differences and an explicit Wilson counterexample to marking every edge of
such a graph.  It withdraws the earlier all-weak and global quartic MTA
claims, and therefore makes the paid quartic entry used in parts of YM V57
conditional.  The valid replacement is an assembled product--cumulant
identity followed by a finite threshold compiler.  Global quartic and
quintic MTA, continuum Yang--Mills, and a mass gap remain open.

\begin{assessment}[Correction to the V65 audit transfer]
\label{ass:v2v70-correction}
V65 used only the algebraic lesson that a payer must be allocated after
complete assembly; it did not import the YM quartic estimate.  That lesson
survives.  The stronger reading that every V57 lower-order payer is already
available does not survive YM V59, because its paid four-row entry is now
conditional.  The SM proof will likewise avoid replacing the full low
inverse by independently normed dyadic inverse blocks unless reduction is
proved.
\end{assessment}

\begin{assessment}[V70 audit decision]
\label{ass:v2v70-decision}
The low dyadic projections will be inserted around the full inverse,
not used to define separate inverses:
\[
 Q_\mu R_{n,L}Q_\nu,
 \qquad
 R_{n,L}=H_{n,L}^{-1}.
\]
All mixed-band pairs will be summed after one symmetric weighted estimate
on \(D^jR_{n,L}\).  This preserves band mixing exactly and turns V69's
block-diagonal theorem into a special case.
\end{assessment}

\begin{firewall}[Cross-program type boundary]
\label{fw:v2v70-types}
A Wilson graph orbit is not a spectral band.  The audit transfers the
logical correction only: an auxiliary decomposition may be normed
termwise only when it is an exact positive or reducing decomposition with
the required bound.  The mixed-band theorem below is proved directly from
operator blocks and does not import a Yang--Mills estimate.
\end{firewall}

\subsection{The assembled weighted low inverse}
\label{sec:v2v70-weight}

Let
\[
 Q_{n,L}=\sum_{\mu\in\mathfrak M_n}Q_\mu
\]
be a finite orthogonal dyadic partition of the positive-frequency low
space, with
\[
 \mathfrak M_n
 =
 \{\mu=2^r:1\le b\mu\le c_L\Lambda_n\}.
\]
Let
\begin{equation}
 M_n
 =
 \sum_{\mu\in\mathfrak M_n}\mu Q_\mu.
\label{eq:v2v70-M}
\end{equation}
The zero-mode block is excluded and remains an infrared problem.

Let
\[
 R_{n,L}(k)=H_{n,L}(k)^{-1}
\]
be the inverse of the complete low kinetic block, without assuming that
\(H_{n,L}\) reduces the \(Q_\mu\).

\begin{definition}[Symmetric differentiated low-resolvent bound]
\label{def:v2v70-weighted-R}
The full low inverse has differentiated order minus one through order five
if
\begin{equation}
 \left\|
 M_n^{(1+j)/2}
 D^jR_{n,L}(k)
 M_n^{(1+j)/2}
 \right\|_{\mathcal L^j}
 \le r_j,
 \qquad 0\le j\le5,
\label{eq:v2v70-weighted-R}
\end{equation}
uniformly in \(n\), \(k\), and derivative directions of unit norm.
\end{definition}

\begin{lemma}[Weighted bound implies every mixed block bound]
\label{lem:v2v70-block-R}
Under
\eqref{eq:v2v70-weighted-R},
\begin{equation}
 \|Q_\mu D^jR_{n,L}(k)Q_\nu\|_{\mathcal L^j}
 \le
 r_j(\mu\nu)^{-(1+j)/2}.
\label{eq:v2v70-block-R}
\end{equation}
\end{lemma}

\begin{proof}
On \(\operatorname{Ran}Q_\mu\),
\(M_n^{-(1+j)/2}\) is multiplication by
\(\mu^{-(1+j)/2}\).  Insert this inverse weight on both sides of the
operator in
\eqref{eq:v2v70-weighted-R} and compress by \(Q_\mu,Q_\nu\).
\end{proof}

\begin{assessment}[Why the weight is symmetric]
\label{ass:v2v70-symmetric}
A mixed inverse block has an incoming and an outgoing low frequency.  A
one-sided weight can charge only one of them and gives no symmetric
summation when \(\mu\) and \(\nu\) are exchanged.  The half-power on each
side is exactly what converts a derivative of order \(j\) into the kernel
factor
\((\mu\nu)^{-(1+j)/2}\).
\end{assessment}

\subsection{Exact mixed-band expansion}
\label{sec:v2v70-expansion}

Let
\[
 X_{n,\mu}=P_nV_nQ_\mu,
 \qquad
 Y_{n,\nu}=Q_\nu V_nP_n,
\]
and assume the V69 band-sensitive crossing estimates
\begin{align}
 \|D^iX_{n,\mu}(k)\|_{\mathcal L^i}
 &\le x_i\Lambda_n^{-1}\mu^{1-i},
\label{eq:v2v70-X}\\
 \|D^\ell Y_{n,\nu}(k)\|_{\mathcal L^\ell}
 &\le y_\ell\Lambda_n^{-1}\nu^{1-\ell}
\label{eq:v2v70-Y}
\end{align}
through order five.  Define the complete low self-energy
\begin{equation}
 \Sigma_{n,L}
 =
 P_nV_nQ_{n,L}\,
 R_{n,L}\,
 Q_{n,L}V_nP_n.
\label{eq:v2v70-Sigma}
\end{equation}
Orthogonality gives the exact identity
\begin{equation}
 \Sigma_{n,L}
 =
 \sum_{\mu,\nu\in\mathfrak M_n}
 X_{n,\mu}
 (Q_\mu R_{n,L}Q_\nu)
 Y_{n,\nu}.
\label{eq:v2v70-mixed-sum}
\end{equation}
No inverse of a compressed diagonal block appears.

\begin{lemma}[Mixed-band differentiated block]
\label{lem:v2v70-mixed-block}
For \(i+j+\ell=5\), a corresponding differentiated term in
\eqref{eq:v2v70-mixed-sum} has operator norm at most
\begin{equation}
 C_{i,j,\ell}\Lambda_n^{-2}
 \mu^{1-i-(1+j)/2}
 \nu^{1-\ell-(1+j)/2}.
\label{eq:v2v70-mixed-op}
\end{equation}
Its rank is at most
\begin{equation}
 W\min\{\mu,\nu\}^d
\label{eq:v2v70-mixed-rank}
\end{equation}
under the V69 rank bound.
\end{lemma}

\begin{proof}
Multiply
\eqref{eq:v2v70-X},
\eqref{eq:v2v70-block-R}, and
\eqref{eq:v2v70-Y}.  This gives
\eqref{eq:v2v70-mixed-op}.  The product maps \(P_n\mathcal H\) through
\(Q_\nu\mathcal H\), then \(Q_\mu\mathcal H\), and back.  Its rank is no
larger than either intermediate rank, proving
\eqref{eq:v2v70-mixed-rank}.
\end{proof}

\subsection{Four-dimensional double-band summation}
\label{sec:v2v70-sum}

\begin{theorem}[Assembled mixed-band fifth-derivative bound]
\label{thm:v2v70-mixed}
Let \(d=4\).  Under
\eqref{eq:v2v70-weighted-R},
\eqref{eq:v2v70-X}, and
\eqref{eq:v2v70-Y},
\begin{align}
 \left|
 D^5\operatorname{Tr}_{P_n}\Sigma_{n,L}(k)
 \right|
 &\le
 C\Lambda_n^{-2}
 \left(
  \Lambda_n^{1/2}+(1+\log\Lambda_n)^2
 \right)
\label{eq:v2v70-mixed-bound}\\
 &\le
 C'\Lambda_n^{-3/2}.
\label{eq:v2v70-mixed-simple}
\end{align}
In particular, the mixed-band low packet is absolutely summable for
dyadic \(\Lambda_n\simeq2^n\) and satisfies the V60
\(O(\Lambda_n^{-1})\) target.
\end{theorem}

\begin{proof}
Use the exact fifth-order product rule in
\eqref{eq:v2v70-mixed-sum}.  Fix derivative counts
\(i+j+\ell=5\).

First suppose \(\mu\le\nu\).  Trace norm is bounded by rank times operator
norm.  Lemma~\ref{lem:v2v70-mixed-block} with \(d=4\) gives
\begin{equation}
 C\Lambda_n^{-2}
 \left(\frac{\mu}{\nu}\right)^{p_{i,\ell}},
 \qquad
 p_{i,\ell}
 =
 2+\frac{\ell-i}{2}.
\label{eq:v2v70-ratio}
\end{equation}
Indeed, after substituting \(j=5-i-\ell\), the total homogeneity in
\(\mu,\nu\) is zero.  Since \(i,\ell\ge0\) and \(i+\ell\le5\),
\[
 -\frac12\le p_{i,\ell}\le\frac92.
\]
For dyadic bands, the double sum of
\((\mu/\nu)^p\) over \(\mu\le\nu\) is:
\begin{enumerate}
 \item \(O(1+\log\Lambda_n)\) when \(p>0\);
 \item \(O((1+\log\Lambda_n)^2)\) when \(p=0\);
 \item \(O(\Lambda_n^{1/2})\) when \(p=-1/2\).
\end{enumerate}
The last estimate is a pair of geometric sums:
\[
 \sum_{\nu\le C\Lambda_n}\nu^{1/2}
 \sum_{\mu\le\nu}\mu^{-1/2}
 \le C\Lambda_n^{1/2}.
\]
Here and below the sums are over dyadic values, so the inner sum is
uniformly bounded.

The region \(\nu\le\mu\) is identical with
\(i,\mu\) exchanged with \(\ell,\nu\).  There are only 21 derivative-count
triples and finitely many labelled assignments.  This proves
\eqref{eq:v2v70-mixed-bound}.  The elementary bound
\((1+\log\Lambda)^2\le C\Lambda^{1/2}\) for \(\Lambda\ge1\) gives
\eqref{eq:v2v70-mixed-simple}.  Finally,
\(\sum_n2^{-3n/2}<\infty\).
\end{proof}

\begin{assessment}[What mixing costs]
\label{ass:v2v70-mixing-cost}
The block-diagonal V69 result is
\(O(\Lambda_n^{-2}\log\Lambda_n)\).
Arbitrary mixed bands satisfying the symmetric kernel bound cost at most
one half power, giving \(O(\Lambda_n^{-3/2})\).  The loss comes only from
the two extreme allocations in which all five derivatives land on one
crossing and the low frequencies are strongly separated.
\end{assessment}

\begin{corollary}[Harmonic mixed-band packet]
\label{cor:v2v70-harmonic}
Let the post-trace tensor contraction be uniformly bounded and let
\(\Pi_n\) be the V68 orthogonal harmonic projector.  Then
\begin{equation}
 \left\|
 \Pi_nD^5\tau_n(\Sigma_{n,L})
 \right\|
 \le C\Lambda_n^{-3/2}.
\label{eq:v2v70-harmonic}
\end{equation}
Thus the full positive-frequency mixed-band minimal packet meets
\eqref{eq:v2v68-transverse-target}.
\end{corollary}

\begin{proof}
Apply the bounded contraction to
Theorem~\ref{thm:v2v70-mixed}; the orthogonal projector has norm one.
\end{proof}

\subsection{Necessity of a joint inverse estimate}
\label{sec:v2v70-necessity}

\begin{counterexample}[Diagonal inverse estimates do not control mixing]
\label{sep:v2v70-diagonal}
Let the low space have two one-dimensional bands
\(Q_\mu,Q_\nu\).  For any \(M>0\), the block operator
\[
 R=
 \begin{pmatrix}
  \mu^{-1}&M\\
  0&\nu^{-1}
 \end{pmatrix}
\]
has the expected diagonal inverse sizes, while
\[
 \|Q_\mu RQ_\nu\|=M.
\]
Thus separate bounds on
\(Q_\mu RQ_\mu\) and \(Q_\nu RQ_\nu\) give no estimate of the mixed
self-energy.
\end{counterexample}

\begin{proof}
The statements are the displayed matrix entries.  The matrix is invertible
for every finite \(M\), with inverse a valid low kinetic block.
\end{proof}

\begin{firewall}[The weighted inverse hypothesis is the new analytic gate]
\label{fw:v2v70-weighted-open}
Theorem~\ref{thm:v2v70-mixed} removes the block-diagonal assumption of V69
but assumes
\eqref{eq:v2v70-weighted-R}.  Neither diagonal band coercivity nor the
unweighted norm of \(R_{n,L}\) proves this symmetric estimate.  It must be
derived from the complete low kinetic operator, including its interband
couplings, domains, massless sectors, and background dependence.
\end{firewall}

\subsection{V70 conclusion and next acquisition}
\label{sec:v2v70-conclusion}

V70 records the YM V59 correction and applies its proof-order lesson to the
SM low resolvent.  The complete mixed-band inverse is retained, and one
symmetric weighted derivative bound closes the harmonic minimal packet
with the summable scale \(O(\Lambda_n^{-3/2})\).

V71 will attack
\eqref{eq:v2v70-weighted-R} upstream.  It will normalize the low kinetic
operator by \(M_n^{1/2}\), separate its diagonal elliptic part from
interband coupling, and prove a differentiated Neumann criterion in a
weighted operator algebra.  A finite matrix regression will show the
precise failure when the normalized coupling reaches unit norm.
\clearpage
\section*{Second Volume, Version V71: Sobolev-Scale Neumann Criterion for the Weighted Low Resolvent}
\addcontentsline{toc}{section}{Second Volume, Version V71: Sobolev-Scale Neumann Criterion for the Weighted Low Resolvent}

\subsection*{Purpose}

V70 reduced arbitrary low-band mixing to the symmetric estimates
\[
 \left\|
 M_n^{(1+j)/2}D^jR_{n,L}M_n^{(1+j)/2}
 \right\|\le r_j,
 \qquad 0\le j\le5.
\]
This note proves those estimates from two operator-scale hypotheses:
uniform ellipticity of a diagonal reference block and strict contraction
of the relative interband coupling on every Sobolev index used by five
derivatives.

The proof avoids manipulating symmetric weights inside a noncommutative
product.  Instead, it regards a derivative of order \(j\) of the inverse
as an operator of order \(-1-j\) on the full Hilbert scale.  The symmetric
V70 norm is then one particular source and target pair on that scale.

\subsection{The dyadic Hilbert scale}
\label{sec:v2v71-scale}

Let \(M_n\ge I\) be the positive dyadic weight from
\eqref{eq:v2v70-M} on the positive-frequency low space
\(\mathcal H_{n,L}\).  For \(s\in\mathbb R\), define
\begin{equation}
 \mathcal H_{n,L}^s
 =
 \mathcal H_{n,L},
 \qquad
 \|u\|_{n,s}
 =
 \|M_n^su\|_{\mathcal H_{n,L}}.
\label{eq:v2v71-scale}
\end{equation}
At finite regulator these are the same vector space with different norms.

An operator \(T\) has order at most \(r\) on a scale interval \(I\) if
\begin{equation}
 T:\mathcal H_{n,L}^s
 \longrightarrow
 \mathcal H_{n,L}^{s-r}
\label{eq:v2v71-order}
\end{equation}
is bounded uniformly whenever both indices lie in \(I\).

\begin{lemma}[Symmetric weights are endpoint scale norms]
\label{lem:v2v71-symmetric}
For \(a\ge0\) and any operator \(T\),
\begin{equation}
 \|T:\mathcal H_{n,L}^{-a}\to\mathcal H_{n,L}^{a}\|
 =
 \|M_n^aTM_n^a\|_{\mathcal H_{n,L}\to\mathcal H_{n,L}}.
\label{eq:v2v71-symmetric}
\end{equation}
\end{lemma}

\begin{proof}
Write \(v=M_n^{-a}u\), so \(u=M_n^av\).  Then
\[
 \frac{\|Tu\|_{n,a}}{\|u\|_{n,-a}}
 =
 \frac{\|M_n^aTM_n^av\|}{\|v\|}.
\]
Taking the supremum over nonzero \(v\) proves the identity.
\end{proof}

\subsection{Inverse differentiation on a scale}
\label{sec:v2v71-inverse}

Let
\[
 H_n:B_\kappa\subset\mathbb R^q
 \longrightarrow\mathcal L(\mathcal H_{n,L})
\]
be \(C^5\) at finite regulator, invertible for every \(k\), and put
\(R_n=H_n^{-1}\).

\begin{theorem}[Differentiated inverse gains one scale per derivative]
\label{thm:v2v71-scale-inverse}
Assume, uniformly in \(n,k\):
\begin{enumerate}
 \item for every \(s\in[-3,2]\),
 \begin{equation}
  R_n(k):
  \mathcal H_{n,L}^s\longrightarrow
  \mathcal H_{n,L}^{s+1}
  \quad\hbox{has norm at most }r_0;
 \label{eq:v2v71-R-order}
 \end{equation}
 \item for \(1\le a\le5\),
 \begin{equation}
  D^aH_n(k):
  \mathcal H_{n,L}^s\longrightarrow
  \mathcal H_{n,L}^{s+a-1}
 \label{eq:v2v71-H-order}
 \end{equation}
 has norm at most \(h_a\) whenever the source and target indices lie in
 \([-3,3]\).
\end{enumerate}
Then, for \(0\le m\le5\),
\begin{equation}
 D^mR_n(k):
 \mathcal H_{n,L}^s\longrightarrow
 \mathcal H_{n,L}^{s+m+1}
\label{eq:v2v71-DmR}
\end{equation}
has norm at most \(r_m\) whenever all intermediate indices lie in
\([-3,3]\).  In particular,
\begin{equation}
 \left\|
 M_n^{(m+1)/2}D^mR_n(k)M_n^{(m+1)/2}
 \right\|
 \le r_m.
\label{eq:v2v71-V70}
\end{equation}
\end{theorem}

\begin{proof}
For \(m=0\), this is
\eqref{eq:v2v71-R-order}.  For \(m\ge1\), use the ordered-partition
inverse formula from Lemma~\ref{lem:v2v61-inverse-formula} with the
identity projection.  A typical term is
\[
 R_n(D^{b_1}H_n)R_n\cdots
 (D^{b_j}H_n)R_n,
 \qquad
 b_1+\cdots+b_j=m.
\]
Read the composition from right to left.  Each \(R_n\) raises the Sobolev
index by one, and each \(D^{b_\ell}H_n\) changes it by
\(b_\ell-1\).  The total change is
\[
 (j+1)+\sum_{\ell=1}^j(b_\ell-1)
 =
 m+1.
\]
The stated interval condition permits each composition, and the product
of its operator norms is bounded uniformly.  Summing the finitely many
ordered partitions gives \(r_m\).

For the last assertion, choose
\(s=-(m+1)/2\).  The target in
\eqref{eq:v2v71-DmR} is \((m+1)/2\), and both endpoints lie in
\([-3,3]\).  Apply Lemma~\ref{lem:v2v71-symmetric}.
\end{proof}

\begin{assessment}[Why all intermediate scale indices are stated]
\label{ass:v2v71-intermediate}
A symmetric endpoint bound alone is not stable under the inverse
derivative formula: the factors are ordered and pass through intermediate
Sobolev exponents.  The theorem assumes exactly the mapping properties
needed to compose those factors.  At order five the largest absolute
endpoint is three.
\end{assessment}

\subsection{Relative Neumann ellipticity}
\label{sec:v2v71-Neumann}

Write
\begin{equation}
 H_n(k)=D_n+E_n(k),
\label{eq:v2v71-split}
\end{equation}
where \(D_n\) is the fixed diagonal elliptic part and \(E_n\) contains the
running diagonal correction and all interband coupling.

Assume:
\begin{enumerate}[label=\textup{(N\arabic*)}]
 \item For every \(s\in[-3,2]\),
 \begin{equation}
  D_n^{-1}:
  \mathcal H_{n,L}^s\longrightarrow
  \mathcal H_{n,L}^{s+1},
  \qquad
  \|D_n^{-1}\|\le d_0.
 \label{eq:v2v71-D-inverse}
 \end{equation}
 \item For every \(t\in[-3,3]\),
 \begin{equation}
  K_n(k):=D_n^{-1}E_n(k)
  \quad\hbox{satisfies}\quad
  \|K_n(k):\mathcal H_{n,L}^t\to
                   \mathcal H_{n,L}^t\|
  \le\theta<1.
 \label{eq:v2v71-relative}
 \end{equation}
 \item For \(1\le a\le5\),
 \begin{equation}
  D^aE_n(k):
  \mathcal H_{n,L}^s\longrightarrow
  \mathcal H_{n,L}^{s+a-1},
  \qquad
  \|D^aE_n(k)\|\le e_a,
 \label{eq:v2v71-E-derivatives}
 \end{equation}
 whenever both indices lie in \([-3,3]\).
\end{enumerate}

\begin{theorem}[Differentiated Neumann criterion]
\label{thm:v2v71-Neumann}
Under \textup{(N1)}--\textup{(N3)}, \(H_n(k)\) is invertible and
\begin{equation}
 R_n(k)
 =
 (I+K_n(k))^{-1}D_n^{-1}.
\label{eq:v2v71-R-Neumann}
\end{equation}
For \(s\in[-3,2]\),
\begin{equation}
 \|R_n(k):
 \mathcal H_{n,L}^s\to\mathcal H_{n,L}^{s+1}\|
 \le
 \frac{d_0}{1-\theta}.
\label{eq:v2v71-R-bound}
\end{equation}
Moreover, all V70 symmetric derivative estimates
\eqref{eq:v2v71-V70} hold through order five.
\end{theorem}

\begin{proof}
The factorization
\[
 H_n
 =
 D_n(I+D_n^{-1}E_n)
 =
 D_n(I+K_n)
\]
is exact.  On each \(\mathcal H_{n,L}^t\), the Neumann series gives
\[
 (I+K_n)^{-1}
 =
 \sum_{r=0}^\infty(-K_n)^r,
 \qquad
 \|(I+K_n)^{-1}\|
 \le(1-\theta)^{-1}.
\]
For input index \(s\), first apply \(D_n^{-1}\) to reach \(s+1\), then
apply the inverse Neumann factor on that space.  This proves
\eqref{eq:v2v71-R-Neumann}--\eqref{eq:v2v71-R-bound}.

Because \(D_n\) is independent of \(k\),
\(D^aH_n=D^aE_n\) for \(a\ge1\).  Thus
\eqref{eq:v2v71-E-derivatives} supplies the differentiated kinetic bounds
of
Theorem~\ref{thm:v2v71-scale-inverse}, while
\eqref{eq:v2v71-R-bound} supplies its inverse bound.  Apply that theorem.
\end{proof}

\begin{corollary}[Closure of the V70 analytic gate]
\label{cor:v2v71-V70}
Assume the V70 crossing and rank hypotheses in addition to
\textup{(N1)}--\textup{(N3)}.  Then the complete positive-frequency
mixed-band minimal packet obeys
\begin{equation}
 \left\|
 \Pi_nD^5\tau_n(\Sigma_{n,L})
 \right\|
 \le C\Lambda_n^{-3/2},
\label{eq:v2v71-packet}
\end{equation}
and is absolutely summable over dyadic shells.
\end{corollary}

\begin{proof}
Theorem~\ref{thm:v2v71-Neumann} proves
\eqref{eq:v2v70-weighted-R}.  Apply
Corollary~\ref{cor:v2v70-harmonic}.
\end{proof}

\begin{firewall}[This is a conditional closure, not ellipticity of the full SM operator]
\label{fw:v2v71-SM}
Corollary~\ref{cor:v2v71-V70} closes the V70 gate when the relative
coupling is a strict contraction on the entire scale interval and the
external derivatives have the stated order.  V71 does not prove those
facts for the full curved, gauge-fixed, running Standard Model operator.
Nor does it treat the excluded zero-mode block.
\end{firewall}

\subsection{Sharpness of the contraction denominator}
\label{sec:v2v71-sharp}

\begin{counterexample}[The Neumann denominator is sharp]
\label{sep:v2v71-theta}
On \(\mathbb C^2\) with the trivial weight \(M=I\), take
\[
 D=I,
 \qquad
 E_\theta=
 \begin{pmatrix}
  0&-\theta\\
  -\theta&0
 \end{pmatrix},
 \qquad
 0\le\theta<1.
\]
Then
\begin{equation}
 \|D^{-1}E_\theta\|=\theta,
 \qquad
 \|(D+E_\theta)^{-1}\|
 =
 \frac1{1-\theta}.
\label{eq:v2v71-sharp}
\end{equation}
At \(\theta=1\), \(D+E_\theta\) is singular.
\end{counterexample}

\begin{proof}
The eigenvalues of \(E_\theta\) are \(\pm\theta\), while those of
\(I+E_\theta\) are \(1\pm\theta\).  The smallest singular value is
\(1-\theta\).
\end{proof}

\begin{assessment}[Sufficiency versus necessity]
\label{ass:v2v71-necessity}
The condition \(\|K_n\|<1\) is a robust sufficient criterion, not a
necessary condition for invertibility.  The example proves that the
factor \((1-\theta)^{-1}\) and the failure at the unit threshold cannot be
improved without using additional sign, normality, or spectral
information.
\end{assessment}

\subsection{Why an unweighted contraction is insufficient}
\label{sec:v2v71-unweighted}

\begin{counterexample}[Small coupling at \(s=0\) can be large at \(s=-3\)]
\label{sep:v2v71-scale-loss}
Let
\[
 M_L=
 \begin{pmatrix}
  1&0\\0&L
 \end{pmatrix},
 \qquad
 K_L=
 \begin{pmatrix}
  0&L^{-1}\\0&0
 \end{pmatrix},
 \qquad
 L>1.
\]
Then
\begin{equation}
 \|K_L:\mathcal H^0\to\mathcal H^0\|
 =
 L^{-1},
\label{eq:v2v71-unweighted-small}
\end{equation}
but
\begin{equation}
 \|K_L:\mathcal H^{-3}\to\mathcal H^{-3}\|
 =
 \|M_L^{-3}K_LM_L^3\|
 =
 L^2.
\label{eq:v2v71-weighted-large}
\end{equation}
\end{counterexample}

\begin{proof}
The unweighted matrix has one entry of size \(L^{-1}\).  Conjugation by
\(M_L^{-3}\) on the left and \(M_L^3\) on the right multiplies its
high-to-low entry by \(L^3\), giving \(L^2\).
\end{proof}

\begin{firewall}[One Sobolev norm cannot certify five derivatives]
\label{fw:v2v71-one-norm}
The unweighted relative coupling can tend to zero while the negative-scale
coupling diverges.  Hence an \(L^2\) Neumann argument does not imply
\eqref{eq:v2v70-weighted-R}.  The scale interval in
\eqref{eq:v2v71-relative} is an analytic requirement, not bookkeeping.
\end{firewall}

\subsection{V71 conclusion and next acquisition}
\label{sec:v2v71-conclusion}

V71 proves the V70 symmetric differentiated inverse bound from a
first-order Hilbert-scale calculus and a strict relative Neumann condition
on indices \([-3,3]\).  This closes the complete positive-frequency
mixed-band minimal packet under explicit ellipticity, coupling, and
external-derivative hypotheses.

The next upstream task is to prove the relative contraction rather than
assume it.  V72 will express the dyadic matrix of
\(K_n=D_n^{-1}E_n\), derive a Schur-test criterion from off-diagonal
frequency decay, and determine the decay exponent needed uniformly on all
Sobolev indices \([-3,3]\).  It will retain diagonal running corrections
separately so that a large marginal coupling is not hidden inside an
off-diagonal estimate.
\clearpage
\section*{Second Volume, Version V72: Dyadic Schur Test for the Interband Coupling}
\addcontentsline{toc}{section}{Second Volume, Version V72: Dyadic Schur Test for the Interband Coupling}

\subsection*{Purpose}

V71 proves the weighted low-resolvent estimates once the relative
interband coupling is a strict contraction on every Hilbert-scale index
\(|s|\le3\).  This note gives a concrete dyadic criterion for that
contraction.

A diagonal running correction need not be small.  It is therefore absorbed
into an exactly inverted bandwise kinetic block.  Only the dressed
off-diagonal coupling is subjected to a Schur test.  If its blocks decay
like
\(2^{-\alpha|r-t|}\) between dyadic indices, then conjugation to
\(\mathcal H^s\) consumes \(2^{|s||r-t|}\).  Uniform control through
\(|s|=3\) requires \(\alpha>3\).  A triangular finite-matrix family proves
that the endpoint \(\alpha=3\) is insufficient for this general block
class.

\subsection{Block Schur bound}
\label{sec:v2v72-Schur}

Let
\[
 \mathcal H=\bigoplus_{r=0}^{N}\mathcal H_r
\]
be an orthogonal finite sum, and let \(T=(T_{rt})\) be a block operator,
where \(T_{rt}:\mathcal H_t\to\mathcal H_r\).  Put
\[
 a_{rt}=\|T_{rt}\|.
\]

\begin{lemma}[Block Schur test]
\label{lem:v2v72-Schur}
If
\begin{equation}
 A_{\rm row}
 =
 \sup_r\sum_ta_{rt}<\infty,
 \qquad
 A_{\rm col}
 =
 \sup_t\sum_ra_{rt}<\infty,
\label{eq:v2v72-row-col}
\end{equation}
then
\begin{equation}
 \|T\|_{\mathcal H\to\mathcal H}
 \le
 \sqrt{A_{\rm row}A_{\rm col}}.
\label{eq:v2v72-Schur}
\end{equation}
\end{lemma}

\begin{proof}
For \(x=(x_t)\), submultiplicativity and weighted Cauchy--Schwarz give
\begin{align*}
 \|(Tx)_r\|^2
 &\le
 \left(\sum_ta_{rt}\|x_t\|\right)^2\\
 &\le
 \left(\sum_ta_{rt}\right)
 \left(\sum_ta_{rt}\|x_t\|^2\right).
\end{align*}
Sum in \(r\), use the row bound, interchange the finite sums, and use the
column bound:
\[
 \|Tx\|^2
 \le
 A_{\rm row}
 \sum_t\left(\sum_ra_{rt}\right)\|x_t\|^2
 \le
 A_{\rm row}A_{\rm col}\|x\|^2.
\]
\end{proof}

\subsection{Weighted dyadic decay}
\label{sec:v2v72-decay}

Let \(Q_r\) be the band projection at frequency
\(\mu_r=2^r\), and let
\[
 M=\sum_{r=0}^N2^rQ_r.
\]
The norm of \(T\) on the Hilbert scale
\(\mathcal H^s\) is
\begin{equation}
 \|T:\mathcal H^s\to\mathcal H^s\|
 =
 \|M^sTM^{-s}\|.
\label{eq:v2v72-conjugation}
\end{equation}
Its \(rt\)-block is multiplied by \(2^{s(r-t)}\).

For \(\beta>0\), define
\begin{equation}
 S_\beta
 =
 2\sum_{m=1}^\infty2^{-\beta m}
 =
 \frac{2}{2^\beta-1}.
\label{eq:v2v72-Sbeta}
\end{equation}

\begin{theorem}[Dyadic off-diagonal contraction]
\label{thm:v2v72-decay}
Suppose \(T_{rr}=0\) and
\begin{equation}
 \|T_{rt}\|
 \le
 \varepsilon\,2^{-\alpha|r-t|}
 \qquad(r\ne t)
\label{eq:v2v72-decay}
\end{equation}
for some \(\alpha>3\).  Then, for every \(|s|\le3\),
\begin{equation}
 \|T:\mathcal H^s\to\mathcal H^s\|
 \le
 \varepsilon S_{\alpha-3}.
\label{eq:v2v72-scale-bound}
\end{equation}
In particular, if
\begin{equation}
 \varepsilon S_{\alpha-3}<1,
\label{eq:v2v72-contraction}
\end{equation}
then \(I+T\) is invertible on every such scale and
\begin{equation}
 \|(I+T)^{-1}:\mathcal H^s\to\mathcal H^s\|
 \le
 \frac1{1-\varepsilon S_{\alpha-3}}.
\label{eq:v2v72-inverse}
\end{equation}
\end{theorem}

\begin{proof}
By
\eqref{eq:v2v72-conjugation}, the conjugated block is bounded by
\begin{align*}
 2^{s(r-t)}
 \varepsilon2^{-\alpha|r-t|}
 &\le
 \varepsilon2^{-(\alpha-|s|)|r-t|}\\
 &\le
 \varepsilon2^{-(\alpha-3)|r-t|}.
\end{align*}
The sum over \(t\ne r\) and the sum over \(r\ne t\) are both bounded by
\(\varepsilon S_{\alpha-3}\), independently of \(N\).  Apply
Lemma~\ref{lem:v2v72-Schur}.  Under
\eqref{eq:v2v72-contraction}, the Neumann series gives
\eqref{eq:v2v72-inverse}.
\end{proof}

\subsection{Exact treatment of the diagonal running block}
\label{sec:v2v72-diagonal}

Let the complete positive-frequency low kinetic operator be
\begin{equation}
 H_n(k)=D_n(I+K_n(k)),
 \qquad
 K_n=D_n^{-1}(H_n-D_n),
\label{eq:v2v72-relative-K}
\end{equation}
where \(D_n\) is the diagonal elliptic reference from V71.  Decompose
\begin{align}
 K_n^{\rm diag}
 &=
 \sum_rQ_rK_nQ_r,
\label{eq:v2v72-Kdiag}\\
 K_n^{\rm off}
 &=
 K_n-K_n^{\rm diag},
\label{eq:v2v72-Koff}\\
 B_n
 &=
 I+K_n^{\rm diag}.
\label{eq:v2v72-B}
\end{align}
Assume every band block of \(B_n\) is invertible and
\begin{equation}
 \sup_{n,k,r}
 \|(Q_rB_n(k)Q_r)^{-1}\|
 \le b_0.
\label{eq:v2v72-diagonal-inverse}
\end{equation}
Because \(B_n\) is block diagonal, the same \(b_0\) bounds
\(B_n^{-1}\) on every \(\mathcal H^s\).

Define the dressed off-diagonal coupling
\begin{equation}
 L_n(k)=B_n(k)^{-1}K_n^{\rm off}(k).
\label{eq:v2v72-L}
\end{equation}
Then
\begin{equation}
 H_n
 =
 D_nB_n(I+L_n).
\label{eq:v2v72-factorization}
\end{equation}

\begin{theorem}[Dressed interband Neumann criterion]
\label{thm:v2v72-dressed}
Suppose
\begin{equation}
 \|Q_rL_n(k)Q_t\|
 \le
 \varepsilon2^{-\alpha|r-t|}
 \qquad(r\ne t)
\label{eq:v2v72-dressed-decay}
\end{equation}
uniformly in \(n,k\), where
\[
 \alpha>3,
 \qquad
 \vartheta:=\varepsilon S_{\alpha-3}<1.
\]
Suppose also
\begin{equation}
 D_n^{-1}:
 \mathcal H_{n,L}^s\longrightarrow
 \mathcal H_{n,L}^{s+1},
 \qquad
 \|D_n^{-1}\|\le d_0
\label{eq:v2v72-D-elliptic}
\end{equation}
for \(-3\le s\le2\).  Then \(H_n(k)\) is invertible and
\begin{equation}
 R_{n,L}
 =
 (I+L_n)^{-1}B_n^{-1}D_n^{-1}.
\label{eq:v2v72-R}
\end{equation}
Moreover,
\begin{equation}
 \|R_{n,L}:
 \mathcal H_{n,L}^s\to\mathcal H_{n,L}^{s+1}\|
 \le
 \frac{b_0d_0}{1-\vartheta}
\label{eq:v2v72-R-order}
\end{equation}
uniformly for \(-3\le s\le2\).
\end{theorem}

\begin{proof}
Equation
\eqref{eq:v2v72-factorization} follows from
\[
 B_n(I+B_n^{-1}K_n^{\rm off})
 =
 I+K_n^{\rm diag}+K_n^{\rm off}
 =
 I+K_n.
\]
Theorem~\ref{thm:v2v72-decay} applies to \(L_n\), so
\((I+L_n)^{-1}\) has scale norm at most
\((1-\vartheta)^{-1}\).  Apply the factors in
\eqref{eq:v2v72-R} from right to left: \(D_n^{-1}\) raises the scale
index by one, while the two remaining factors preserve it.  Their norms
give
\eqref{eq:v2v72-R-order}.
\end{proof}

\begin{assessment}[Large diagonal running is allowed]
\label{ass:v2v72-large-diagonal}
No smallness of \(K_n^{\rm diag}\) is assumed.  Its only requirement is
the bandwise coercivity
\eqref{eq:v2v72-diagonal-inverse}.  This prevents a marginal diagonal
coupling from being hidden inside an off-diagonal smallness condition.
The contraction is imposed only after the diagonal block has been dressed.
\end{assessment}

\subsection{Differentiated weighted inverse}
\label{sec:v2v72-differentiated}

\begin{theorem}[Schur decay closes the V70 weighted gate]
\label{thm:v2v72-V70}
Assume
Theorem~\ref{thm:v2v72-dressed} and suppose, for \(1\le a\le5\),
\begin{equation}
 D^aH_n(k):
 \mathcal H_{n,L}^s
 \longrightarrow
 \mathcal H_{n,L}^{s+a-1}
\label{eq:v2v72-H-derivatives}
\end{equation}
has norm at most \(h_a\) whenever both indices lie in \([-3,3]\).
Then
\begin{equation}
 \left\|
 M_n^{(a+1)/2}
 D^aR_{n,L}(k)
 M_n^{(a+1)/2}
 \right\|
 \le r_a,
 \qquad 0\le a\le5.
\label{eq:v2v72-weighted-R}
\end{equation}
Consequently, under the V70 crossing and rank hypotheses,
\begin{equation}
 \left\|
 \Pi_nD^5\tau_n(\Sigma_{n,L})
 \right\|
 \le C\Lambda_n^{-3/2}.
\label{eq:v2v72-packet}
\end{equation}
\end{theorem}

\begin{proof}
Theorem~\ref{thm:v2v72-dressed} supplies the order-minus-one inverse bound
\eqref{eq:v2v71-R-order}.  The present derivative hypothesis is
\eqref{eq:v2v71-H-order}.  Apply
Theorem~\ref{thm:v2v71-scale-inverse} to obtain
\eqref{eq:v2v72-weighted-R}.  This is the V70 symmetric hypothesis, so
Corollary~\ref{cor:v2v70-harmonic} gives
\eqref{eq:v2v72-packet}.
\end{proof}

\begin{firewall}[The decay hypothesis is not yet a pseudodifferential theorem]
\label{fw:v2v72-decay-open}
Equation
\eqref{eq:v2v72-dressed-decay} is a quantitative off-diagonal estimate
for the complete dressed coupling.  V72 does not derive it from locality,
smooth coefficients, or the curved Standard Model symbol.  It also leaves
the differentiated order bounds
\eqref{eq:v2v72-H-derivatives} as hypotheses.  Those are the next analytic
acquisitions.
\end{firewall}

\subsection{Sharpness of the exponent for the full scale interval}
\label{sec:v2v72-sharp}

\begin{counterexample}[The endpoint \(\alpha=3\) is not uniform]
\label{sep:v2v72-endpoint}
Let every band be one-dimensional, with indices
\(0,\ldots,N\), and define the strictly lower-triangular matrix
\begin{equation}
 T^{(N)}_{rt}
 =
 \begin{cases}
  \varepsilon2^{-3(r-t)},&r>t,\\
  0,&r\le t.
 \end{cases}
\label{eq:v2v72-endpoint-matrix}
\end{equation}
It satisfies the off-diagonal envelope with \(\alpha=3\).  On
\(\mathcal H^3\), however,
\begin{equation}
 (M^3T^{(N)}M^{-3})_{rt}
 =
 \varepsilon
 \qquad(r>t),
\label{eq:v2v72-endpoint-conjugated}
\end{equation}
and
\begin{equation}
 \|T^{(N)}:\mathcal H^3\to\mathcal H^3\|
 \ge c\varepsilon N.
\label{eq:v2v72-endpoint-growth}
\end{equation}
\end{counterexample}

\begin{proof}
The conjugation multiplies the \(rt\)-entry by
\(2^{3(r-t)}\), proving
\eqref{eq:v2v72-endpoint-conjugated}.  Apply the conjugated matrix to
the unit vector
\[
 x=\frac1{\sqrt{N+1}}(1,\ldots,1).
\]
Its \(r\)-th output is
\(\varepsilon r/\sqrt{N+1}\).  Hence
\[
 \|M^3T^{(N)}M^{-3}x\|^2
 =
 \frac{\varepsilon^2}{N+1}\sum_{r=0}^Nr^2
 \ge c^2\varepsilon^2N^2.
\]
\end{proof}

\begin{corollary}[Failure below the endpoint]
\label{cor:v2v72-below}
For \(\alpha<3\), even one block joining indices \(0\) and \(N\), with
norm \(\varepsilon2^{-\alpha N}\), has
\(\mathcal H^3\)-norm
\[
 \varepsilon2^{(3-\alpha)N},
\]
which diverges with \(N\).
\end{corollary}

\begin{proof}
Conjugation by \(M^3\) multiplies that block by \(2^{3N}\).
\end{proof}

\begin{assessment}[Sharp claim boundary]
\label{ass:v2v72-sharp}
The threshold \(\alpha>3\) is sharp for a theorem based only on the
two-sided block envelope
\eqref{eq:v2v72-decay} over all scale indices \(|s|\le3\).  Special
triangularity, cancellations, self-adjointness, or a smaller scale interval
may improve the threshold, but none is assumed here.
\end{assessment}

\subsection{V72 conclusion and next acquisition}
\label{sec:v2v72-conclusion}

V72 converts dyadic off-diagonal decay of the dressed relative coupling
into the full Sobolev-scale contraction required by V71.  It permits
arbitrarily large diagonal running corrections subject to bandwise
coercivity, and it proves the strict threshold \(\alpha>3\) for the
general block-envelope argument.

V73 will attack
\eqref{eq:v2v72-dressed-decay} by iterated spectral commutators with the
reference kinetic operator.  It will derive a ratio-decay bound between
well-separated dyadic bands and isolate the adjacent-band remainder, where
commutator division supplies no small parameter.
\clearpage
\section*{Second Volume, Version V73: Iterated Spectral Commutators and the Adjacent-Band Remainder}
\addcontentsline{toc}{section}{Second Volume, Version V73: Iterated Spectral Commutators and the Adjacent-Band Remainder}

\subsection*{Purpose}

V72 reduced the low-resolvent gate to decay of the dressed relative
coupling between dyadic bands.  This note derives that decay for
well-separated bands from iterated commutators with a positive reference
frequency operator.

The result is deliberately two-part.  Uniform \(N\)-fold commutator bounds
give order-\(N\) off-diagonal decay and make the separated high-frequency
tail arbitrarily small when \(N>3\).  They do not make adjacent-band
couplings small.  A nilpotent shift example has no far off-diagonal blocks
at all, yet its inverse grows with the number of bands.  Thus the next
gate is a finite-width adjacent-band coercivity estimate, not another
commutator iteration.

\subsection{Separated spectral bands}
\label{sec:v2v73-bands}

Let \(A=A^*\ge I\) on a finite-regulator Hilbert space.  For fixed
\(0<a<b\), let
\begin{equation}
 Q_r
 =
 \mathbf1_{[a2^r,b2^r]}(A),
 \qquad
 r=0,\ldots,N,
\label{eq:v2v73-Qr}
\end{equation}
with endpoint modifications and disjointization understood.  Choose an
integer \(g\ge1\) so large that
\begin{equation}
 b2^{-g}\le\frac a2.
\label{eq:v2v73-g}
\end{equation}
If \(r\ge t+g\), then
\begin{equation}
 \operatorname{dist}
 \bigl(
 \sigma(A|_{Q_r\mathcal H}),
 \sigma(A|_{Q_t\mathcal H})
 \bigr)
 \ge
 \frac a2\,2^r.
\label{eq:v2v73-gap}
\end{equation}

For a bounded operator \(T\), write
\[
 \operatorname{ad}_A(T)=[A,T],
 \qquad
 \operatorname{ad}_A^m(T)
 =
 [A,\operatorname{ad}_A^{m-1}(T)].
\]

\begin{lemma}[Repeated Sylvester identity]
\label{lem:v2v73-Sylvester}
Let
\[
 T_{rt}=Q_rTQ_t,
 \qquad
 A_r=Q_rAQ_r,
 \qquad
 A_t=Q_tAQ_t.
\]
Then
\begin{equation}
 (A_r\,\cdot-\cdot\,A_t)^m(T_{rt})
 =
 Q_r\operatorname{ad}_A^m(T)Q_t.
\label{eq:v2v73-repeated-Sylvester}
\end{equation}
If the two spectra are ordered and separated by \(\delta_{rt}>0\), then
\begin{equation}
 \|T_{rt}\|
 \le
 \delta_{rt}^{-m}
 \|Q_r\operatorname{ad}_A^m(T)Q_t\|.
\label{eq:v2v73-Sylvester-bound}
\end{equation}
\end{lemma}

\begin{proof}
The first identity follows by induction, using
\(AQ_r=Q_rA\) and \(AQ_t=Q_tA\).  For the second, suppose first that
\(\inf\sigma(A_r)>\sup\sigma(A_t)\).  The Sylvester inverse is
\[
 \mathcal S_{rt}^{-1}(B)
 =
 \int_0^\infty e^{-uA_r}Be^{uA_t}\,\mathrm du,
\]
after a common scalar shift if desired.  Its norm is at most
\(\delta_{rt}^{-1}\).  Apply this inverse \(m\) times to
\eqref{eq:v2v73-repeated-Sylvester}.  The reverse spectral order uses the
negative-shell integral of V63 and has the same norm.
\end{proof}

\begin{theorem}[Iterated commutators give separated-band decay]
\label{thm:v2v73-separated}
Suppose
\begin{equation}
 \|\operatorname{ad}_A^m(T)\|
 \le C_m
\label{eq:v2v73-commutator}
\end{equation}
for an integer \(m\ge1\).  If \(|r-t|\ge g\), then
\begin{align}
 \|Q_rTQ_t\|
 &\le
 C_m'
 \max\{2^r,2^t\}^{-m}
\label{eq:v2v73-absolute-decay}\\
 &\le
 C_m'2^{-m|r-t|}.
\label{eq:v2v73-ratio-decay}
\end{align}
\end{theorem}

\begin{proof}
Assume \(r\ge t+g\).  Equations
\eqref{eq:v2v73-gap} and
\eqref{eq:v2v73-Sylvester-bound} give
\[
 \|Q_rTQ_t\|
 \le
 (2/a)^m2^{-mr}C_m.
\]
This is
\eqref{eq:v2v73-absolute-decay}.  Since \(t\ge0\),
\(2^{-mr}\le2^{-m(r-t)}\), proving
\eqref{eq:v2v73-ratio-decay}.  Exchange \(r,t\) for the other ordering.
\end{proof}

\begin{assessment}[What commutator regularity means]
\label{ass:v2v73-regularity}
The hypothesis
\eqref{eq:v2v73-commutator} is natural for a sufficiently smooth
order-zero coefficient operator whose repeated commutators with the
first-order reference frequency remain order zero.  It is not implied by
boundedness of \(T\).  For a general order-zero pseudodifferential
operator, the uniform commutator bounds must be proved in the chosen
background-dependent calculus.
\end{assessment}

\subsection{Smallness of the separated high-frequency tail}
\label{sec:v2v73-tail}

Let
\[
 M=\sum_r2^rQ_r
\]
and let \(T^{\rm far}\) contain only blocks with
\(|r-t|\ge g\).  For an integer \(R\ge0\), let
\[
 Q_{\ge R}=\sum_{r\ge R}Q_r.
\]

\begin{theorem}[Far-tail contraction on the full V71 scale interval]
\label{thm:v2v73-tail}
Assume
\eqref{eq:v2v73-commutator} for some integer \(m>3\).  Then
\begin{equation}
 \sup_{|s|\le3}
 \left\|
 Q_{\ge R}T^{\rm far}Q_{\ge R}:
 \mathcal H^s\to\mathcal H^s
 \right\|
 \le
 C_{m,g}2^{-(m-3)R}.
\label{eq:v2v73-tail}
\end{equation}
Consequently the far high-frequency coupling is a strict contraction for
all sufficiently large fixed \(R\), uniformly in the ultraviolet cutoff
\(N\).
\end{theorem}

\begin{proof}
Consider \(r\ge t+g\), with \(r,t\ge R\).  By
\eqref{eq:v2v73-absolute-decay}, the \(rt\)-block after conjugation by
\(M^s\) is bounded by
\[
 C_m'2^{-mr}2^{s(r-t)}
 \le
 C_m'2^{-(m-3)r}2^{-3t}
\]
when \(s\ge0\); weakening signs gives the same final geometric majorant
for every \(|s|\le3\).  Its row and column sums over this region are
bounded by \(C_{m,g}2^{-(m-3)R}\).  The region
\(t\ge r+g\) is symmetric.  Apply the block Schur test
Lemma~\ref{lem:v2v72-Schur}.
\end{proof}

\begin{remark}[A direct safer majorant]
\label{rem:v2v73-majorant}
Without tracking the favorable factor \(2^{-3t}\), one may combine
\eqref{eq:v2v73-ratio-decay} with
\(2^{s(r-t)}\) and sum
\(2^{-(m-3)|r-t|}\).  This gives a cutoff-uniform bound but not the
high-tail smallness in
\eqref{eq:v2v73-tail}.  The absolute decay
\eqref{eq:v2v73-absolute-decay} supplies that extra conclusion.
\end{remark}

\subsection{The adjacent-band norm}
\label{sec:v2v73-adjacent}

Define the finite-width adjacent part
\begin{equation}
 T^{\rm adj}
 =
 \sum_{\substack{r,t\\0<|r-t|<g}}
 Q_rTQ_t.
\label{eq:v2v73-adjacent}
\end{equation}
Put
\begin{align}
 A_{\rm row}^{\rm adj}
 &=
 \sup_r
 \sum_{\substack{t\\0<|r-t|<g}}
 \|Q_rTQ_t\|,
\label{eq:v2v73-adj-row}\\
 A_{\rm col}^{\rm adj}
 &=
 \sup_t
 \sum_{\substack{r\\0<|r-t|<g}}
 \|Q_rTQ_t\|.
\label{eq:v2v73-adj-col}
\end{align}

\begin{lemma}[Adjacent-band scale bound]
\label{lem:v2v73-adjacent}
For every \(|s|\le3\),
\begin{equation}
 \|T^{\rm adj}:\mathcal H^s\to\mathcal H^s\|
 \le
 2^{3g}
 \sqrt{
 A_{\rm row}^{\rm adj}
 A_{\rm col}^{\rm adj}}.
\label{eq:v2v73-adjacent-bound}
\end{equation}
\end{lemma}

\begin{proof}
On the adjacent support,
\[
 2^{s(r-t)}
 \le2^{|s||r-t|}
 \le2^{3g}.
\]
Thus the conjugated row and column sums are bounded by
\(2^{3g}A_{\rm row}^{\rm adj}\) and
\(2^{3g}A_{\rm col}^{\rm adj}\).  Apply
Lemma~\ref{lem:v2v72-Schur}.
\end{proof}

\begin{theorem}[Separated commutators plus adjacent coercivity]
\label{thm:v2v73-combined}
Let \(T\) have zero diagonal.  Assume the \(m\)-fold commutator bound with
\(m>3\).  Absorb the fixed finite block
\(\bigoplus_{r<R}\mathcal H_r\), together with its internal couplings,
into the exactly inverted diagonal core.  On the remaining tail, suppose
\begin{equation}
 \vartheta_{\rm adj}
 :=
 2^{3g}
 \sqrt{
 A_{\rm row}^{\rm adj}
 A_{\rm col}^{\rm adj}}
 <1.
\label{eq:v2v73-adj-small}
\end{equation}
For \(R\) sufficiently large that
\begin{equation}
 \vartheta_{\rm adj}
 +
 C_{m,g}2^{-(m-3)R}
 <1,
\label{eq:v2v73-combined-small}
\end{equation}
the tail coupling is a strict contraction on every
\(\mathcal H^s\), \(|s|\le3\).
\end{theorem}

\begin{proof}
Split the tail coupling into its adjacent and far parts.  Apply
Lemma~\ref{lem:v2v73-adjacent} to the first and
Theorem~\ref{thm:v2v73-tail} to the second, then use the triangle
inequality.  Condition
\eqref{eq:v2v73-combined-small} makes the resulting norm strictly below
one.
\end{proof}

\begin{firewall}[Core-to-tail blocks must be included]
\label{fw:v2v73-core-tail}
Absorbing a finite infrared core does not permit deletion of its couplings
to the tail.  Enlarge the core by the fixed adjacency width and retain the
remaining core-to-tail blocks in the far operator, where
Theorem~\ref{thm:v2v73-tail} applies.  Uniform invertibility of the enlarged
core and its Schur complement is an additional finite-dimensional
coercivity requirement.
\end{firewall}

\subsection{Adjacent bands are an independent obstruction}
\label{sec:v2v73-adjacent-obstruction}

\begin{counterexample}[No far blocks, but no uniform inverse]
\label{sep:v2v73-shift}
Let every band be one-dimensional and let \(S_N\) be the upper shift on
\(\mathbb C^{N+1}\):
\[
 S_Ne_t=e_{t+1}\quad(0\le t<N),
 \qquad
 S_Ne_N=0.
\]
Set \(T_N=-S_N\).  Then \(T_N\) has only adjacent blocks, so every far
off-diagonal estimate holds with constant zero.  Nevertheless,
\begin{equation}
 (I+T_N)^{-1}
 =
 I+S_N+\cdots+S_N^N
\label{eq:v2v73-shift-inverse}
\end{equation}
and
\begin{equation}
 \|(I+T_N)^{-1}\|
 \ge c\sqrt N.
\label{eq:v2v73-shift-growth}
\end{equation}
\end{counterexample}

\begin{proof}
Nilpotence gives
\((I-S_N)^{-1}=\sum_{j=0}^NS_N^j\).
Apply this inverse to \(e_0\).  The result is
\(e_0+\cdots+e_N\), of norm \(\sqrt{N+1}\).
\end{proof}

\begin{assessment}[The next gate is not more smoothness]
\label{ass:v2v73-next-gate}
Iterating the commutator controls only separated bands.  The shift
counterexample is banded and has no separated tail to improve.  Closing
the remaining operator therefore requires a sign, coercivity, diagonal
dominance, or transfer estimate on the adjacent-band block itself.
\end{assessment}

\begin{firewall}[The commutator bound remains unproved for the SM coupling]
\label{fw:v2v73-SM}
V73 proves the consequence of
\eqref{eq:v2v73-commutator}; it does not establish that bound uniformly
for the dressed Standard Model coupling on curved backgrounds.  The
dressing inverse can itself affect commutator regularity and must be
included by the Leibniz rule before the theorem is applied.
\end{firewall}

\subsection{V73 conclusion and next acquisition}
\label{sec:v2v73-conclusion}

V73 derives separated dyadic decay from iterated spectral commutators and
proves that every far high-frequency block becomes small on the full V71
scale interval when more than three commutators are controlled.  It also
isolates the adjacent-band operator and proves by a nilpotent shift that
far decay alone cannot control its inverse.

V74 will attack the adjacent operator through a block energy inequality.
It will compare the off-diagonal quadratic form with the exactly dressed
diagonal form, prove a coercive inverse bound without requiring each
adjacent block to be small, and state the sectorial modification needed
for non-self-adjoint gauge-fixed operators.
\clearpage
\section*{Second Volume, Version V74: Adjacent-Band Energy Coercivity and Sectorial Inversion}
\addcontentsline{toc}{section}{Second Volume, Version V74: Adjacent-Band Energy Coercivity and Sectorial Inversion}

\subsection*{Purpose}

V73 proves that sufficiently many spectral commutators make the
well-separated high-band coupling small on every Hilbert-scale index
\(|s|\le3\).  The remaining adjacent-band operator need not have small
norm.  This note replaces norm smallness by a lower bound on its Hermitian
part.

At finite regulator, uniform accretivity of
\(I+T_s^{\rm adj}\), where
\(T_s^{\rm adj}=M^sT^{\rm adj}M^{-s}\), gives an inverse bound even when
the skew-adjoint part is arbitrarily large.  A pairwise block-energy ledger
provides a concrete sufficient condition.  The far commutator tail can
then be added as a small form perturbation.

\subsection{Finite-dimensional accretive inversion}
\label{sec:v2v74-accretive}

Let \(\mathcal H\) be a finite-dimensional complex Hilbert space.

\begin{lemma}[Accretive lower bound implies an inverse bound]
\label{lem:v2v74-accretive}
Let \(A\in\mathcal L(\mathcal H)\) and suppose
\begin{equation}
 \operatorname{Re}\langle u,Au\rangle
 \ge c\|u\|^2
 \qquad(u\in\mathcal H)
\label{eq:v2v74-accretive}
\end{equation}
for some \(c>0\).  Then \(A\) is invertible and
\begin{equation}
 \|A^{-1}\|\le c^{-1}.
\label{eq:v2v74-inverse}
\end{equation}
\end{lemma}

\begin{proof}
Cauchy--Schwarz and
\eqref{eq:v2v74-accretive} give
\[
 c\|u\|^2
 \le
 |\langle u,Au\rangle|
 \le
 \|u\|\|Au\|.
\]
Thus \(\|Au\|\ge c\|u\|\), so \(A\) is injective.  In finite dimension it
is therefore surjective, and the same inequality gives
\eqref{eq:v2v74-inverse}.
\end{proof}

\begin{theorem}[Scale-wise sectorial inverse]
\label{thm:v2v74-scale}
Let \(T\) act on the dyadic Hilbert scale of V71 and put
\[
 T_s=M^sTM^{-s}.
\]
Suppose, for every \(|s|\le3\),
\begin{equation}
 \frac{T_s+T_s^*}{2}
 \ge
 -\vartheta I
\label{eq:v2v74-sectorial}
\end{equation}
with one \(\vartheta<1\).  Then \(I+T\) is invertible on every
\(\mathcal H^s\), and
\begin{equation}
 \|(I+T)^{-1}:\mathcal H^s\to\mathcal H^s\|
 \le
 \frac1{1-\vartheta}.
\label{eq:v2v74-scale-inverse}
\end{equation}
\end{theorem}

\begin{proof}
Conjugation identifies \(I+T\) on \(\mathcal H^s\) with
\(I+T_s\) on the base Hilbert space.  Equation
\eqref{eq:v2v74-sectorial} gives
\[
 \operatorname{Re}\langle u,(I+T_s)u\rangle
 \ge
 (1-\vartheta)\|u\|^2.
\]
Apply Lemma~\ref{lem:v2v74-accretive}.
\end{proof}

\begin{assessment}[Only the negative Hermitian part is dangerous]
\label{ass:v2v74-Hermitian}
The skew-adjoint part of \(T_s\) does not enter
\(\operatorname{Re}\langle u,T_su\rangle\).  A large positive Hermitian
part is also harmless for the lower bound.  This is strictly more flexible
than requiring \(\|T_s\|<1\).
\end{assessment}

\subsection{A pairwise adjacent-band ledger}
\label{sec:v2v74-ledger}

Let
\[
 \mathcal H=\bigoplus_{r=0}^N\mathcal H_r
\]
and let \(T_s^{\rm adj}\) have blocks only when
\(|r-t|<g\).  Write its Hermitian part as
\begin{equation}
 S_s
 =
 \frac{T_s^{\rm adj}+(T_s^{\rm adj})^*}{2},
 \qquad
 S_{s,tr}=S_{s,rt}^*.
\label{eq:v2v74-S}
\end{equation}
Assume
\begin{equation}
 S_{s,rr}\ge-d_{s,r}I,
 \qquad
 d_{s,r}\ge0.
\label{eq:v2v74-diagonal}
\end{equation}
For each pair \(r<t\) with \(t-r<g\), choose
\(\eta_{s,rt}>0\), and define the vertex load
\begin{align}
 \mathfrak l_{s,r}
 =
 d_{s,r}
 &+
 \sum_{\substack{t>r\\t-r<g}}
 \eta_{s,rt}\|S_{s,rt}\|
\notag\\
 &+
 \sum_{\substack{t<r\\r-t<g}}
 \eta_{s,tr}^{-1}\|S_{s,tr}\|.
\label{eq:v2v74-load}
\end{align}

\begin{theorem}[Pairwise block-energy coercivity]
\label{thm:v2v74-ledger}
If
\begin{equation}
 \sup_{\substack{|s|\le3\\0\le r\le N}}
 \mathfrak l_{s,r}
 \le\vartheta<1,
\label{eq:v2v74-load-bound}
\end{equation}
then
\begin{equation}
 S_s\ge-\vartheta I
 \qquad(|s|\le3).
\label{eq:v2v74-S-lower}
\end{equation}
Consequently \(I+T^{\rm adj}\) obeys
\eqref{eq:v2v74-scale-inverse}.
\end{theorem}

\begin{proof}
Let \(u=(u_r)\).  The diagonal contribution satisfies
\[
 \langle u_r,S_{s,rr}u_r\rangle
 \ge-d_{s,r}\|u_r\|^2.
\]
For \(r<t\), Young's inequality gives
\begin{align*}
 2\operatorname{Re}
 \langle u_r,S_{s,rt}u_t\rangle
 &\ge
 -2\|S_{s,rt}\|\|u_r\|\|u_t\|\\
 &\ge
 -\|S_{s,rt}\|
 \left(
  \eta_{s,rt}\|u_r\|^2+
  \eta_{s,rt}^{-1}\|u_t\|^2
 \right).
\end{align*}
Sum over adjacent unordered pairs.  The coefficient charged to
\(\|u_r\|^2\) is exactly
\(\mathfrak l_{s,r}\).  Equation
\eqref{eq:v2v74-load-bound} therefore gives
\[
 \langle u,S_su\rangle
 \ge-\vartheta\sum_r\|u_r\|^2.
\]
Apply Theorem~\ref{thm:v2v74-scale}.
\end{proof}

\begin{remark}[Optimization of the edge weights]
\label{rem:v2v74-eta}
The parameters \(\eta_{s,rt}\) allocate the cost of an adjacent block
between its two endpoints.  Choosing all of them equal to one gives a
simple row-load test.  Unequal choices can exploit different diagonal
coercivities or band dimensions.  Their optimization is a finite convex
feasibility problem after the block norms are known.
\end{remark}

\subsection{Adding the far commutator tail}
\label{sec:v2v74-far}

Let
\[
 T=T^{\rm adj}+T^{\rm far}.
\]

\begin{theorem}[Energy-stable far-tail perturbation]
\label{thm:v2v74-far}
Suppose the adjacent Hermitian parts satisfy
\eqref{eq:v2v74-sectorial} with \(\vartheta<1\), and suppose
\begin{equation}
 \sup_{|s|\le3}
 \|M^sT^{\rm far}M^{-s}\|
 \le\varepsilon_{\rm far}
 <
 1-\vartheta.
\label{eq:v2v74-far-small}
\end{equation}
Then
\begin{equation}
 \frac{M^sTM^{-s}+(M^sTM^{-s})^*}{2}
 \ge
 -(\vartheta+\varepsilon_{\rm far})I
\label{eq:v2v74-full-lower}
\end{equation}
and
\begin{equation}
 \|(I+T)^{-1}:\mathcal H^s\to\mathcal H^s\|
 \le
 \frac1{1-\vartheta-\varepsilon_{\rm far}}
\label{eq:v2v74-full-inverse}
\end{equation}
for every \(|s|\le3\).
\end{theorem}

\begin{proof}
For the conjugated far operator \(F_s\),
\[
 \operatorname{Re}\langle u,F_su\rangle
 \ge-\|F_s\|\|u\|^2
 \ge-\varepsilon_{\rm far}\|u\|^2.
\]
Add this to the adjacent lower bound and apply
Theorem~\ref{thm:v2v74-scale}.
\end{proof}

\begin{corollary}[Commutator tail plus energy ledger]
\label{cor:v2v74-V73}
Assume the V73 \(m\)-fold commutator bound with \(m>3\).  If the adjacent
load bound
\eqref{eq:v2v74-load-bound} holds with \(\vartheta<1\), then, after a
fixed sufficiently large infrared core has been absorbed as in V73, the
remaining dressed coupling has a uniformly bounded inverse on all
\(\mathcal H^s\), \(|s|\le3\).
\end{corollary}

\begin{proof}
By Theorem~\ref{thm:v2v73-tail}, choose the fixed tail threshold so that
\[
 \varepsilon_{\rm far}
 \le C_{m,g}2^{-(m-3)R}
 <
 1-\vartheta.
\]
Apply Theorem~\ref{thm:v2v74-far}.
\end{proof}

\subsection{Self-adjoint and sectorial interpretations}
\label{sec:v2v74-interpretation}

If the dressed adjacent coupling is self-adjoint at \(s=0\), the energy
condition is the relative form bound
\begin{equation}
 \langle u,T^{\rm adj}u\rangle
 \ge-\vartheta\|u\|^2.
\label{eq:v2v74-selfadjoint}
\end{equation}
On nonzero \(s\), conjugation by \(M^s\) need not preserve
self-adjointness, which is why
\eqref{eq:v2v74-sectorial} is imposed on the conjugated operator itself.

For a non-self-adjoint gauge-fixed operator, write
\[
 T_s^{\rm adj}=S_s+A_s,
 \qquad
 S_s=S_s^*,
 \qquad
 A_s^*=-A_s.
\]
Only the lower spectrum of \(S_s\) enters the inverse estimate.  The
skew-adjoint part \(A_s\) may be large.

\begin{proposition}[Large skew coupling can be harmless]
\label{prop:v2v74-skew}
On \(\mathbb C^2\), let
\[
 T_L=
 \begin{pmatrix}
 0&L\\
 -L&0
 \end{pmatrix},
 \qquad
 L\ge0.
\]
Then
\begin{equation}
 \|T_L\|=L,
 \qquad
 \operatorname{Re}\langle u,T_Lu\rangle=0,
 \qquad
 \|(I+T_L)^{-1}\|
 =
 (1+L^2)^{-1/2}\le1.
\label{eq:v2v74-skew}
\end{equation}
\end{proposition}

\begin{proof}
The matrix is skew-adjoint and satisfies
\[
 (I+T_L)^*(I+T_L)=(1+L^2)I.
\]
The three claims follow.
\end{proof}

\begin{counterexample}[A negative symmetric adjacent mode reaches zero]
\label{sep:v2v74-negative}
Let
\[
 T_c=
 \begin{pmatrix}
 0&-c\\
 -c&0
 \end{pmatrix},
 \qquad
 0\le c<1.
\]
Then the Hermitian part has lowest eigenvalue \(-c\), and
\begin{equation}
 \|(I+T_c)^{-1}\|=(1-c)^{-1}.
\label{eq:v2v74-negative}
\end{equation}
At \(c=1\), the vector \((1,1)\) lies in the kernel.
\end{counterexample}

\begin{proof}
The eigenvalues of \(T_c\) are \(\pm c\), so those of \(I+T_c\) are
\(1\pm c\).
\end{proof}

\subsection{Insertion into the weighted resolvent theorem}
\label{sec:v2v74-V71}

Let the dressed factorization of V72 be
\[
 H_n=D_nB_n(I+T_n),
\]
where the exactly inverted diagonal/core factors have the order-minus-one
bound of
\eqref{eq:v2v72-R-order} before the last factor is applied.

\begin{theorem}[Energy criterion closes the base inverse mapping property]
\label{thm:v2v74-base}
Assume:
\begin{enumerate}
 \item the diagonal/core inverse maps
       \(\mathcal H^s\to\mathcal H^{s+1}\) uniformly;
 \item the adjacent part of \(T_n\) satisfies
       \eqref{eq:v2v74-load-bound};
 \item the far part satisfies
       \eqref{eq:v2v74-far-small}; and
 \item \(D^aH_n:\mathcal H^s\to\mathcal H^{s+a-1}\) is uniformly bounded
       through \(a=5\) on the V71 scale interval.
\end{enumerate}
Then the V70 symmetric differentiated inverse estimates hold, and the
harmonic minimal low packet is \(O(\Lambda_n^{-3/2})\).
\end{theorem}

\begin{proof}
Theorem~\ref{thm:v2v74-far} bounds the inverse of the final dressed factor
on every scale.  Compose it with the diagonal/core inverse to obtain
\eqref{eq:v2v71-R-order}.  Apply
Theorem~\ref{thm:v2v71-scale-inverse} for the differentiated inverse, then
Corollary~\ref{cor:v2v70-harmonic} for the packet.
\end{proof}

\begin{firewall}[The energy ledger is not yet verified for the SM blocks]
\label{fw:v2v74-SM}
V74 proves a criterion.  It does not show that the Hermitian parts of the
gauge, ghost, Higgs, fermion, and metric adjacent-band blocks satisfy
\eqref{eq:v2v74-load-bound} after complete Ward assembly.  Indefinite
Krein formulations must first be transferred to a positive Hilbert
majorant or treated by a separate definitizable-operator theorem.
\end{firewall}

\subsection{V74 conclusion and next acquisition}
\label{sec:v2v74-conclusion}

V74 replaces adjacent-band norm smallness by scale-wise energy coercivity.
It proves an explicit pairwise ledger, permits arbitrarily large
skew-adjoint mixing, and combines the resulting adjacent inverse with the
small far tail from V73.  Under the differentiated kinetic bounds, this
closes the V70 weighted-resolvent gate.

V75 is a mandatory companion audit.  After inspecting the current NS and
YM notes, it will choose the first physical adjacent sector in which to
verify the energy ledger.  The likely starting point is the Euclidean
fermion--Higgs block, where the positive squared Dirac operator provides a
Hilbert majorant and the chiral off-diagonal coupling can be tested
explicitly.
\clearpage
\section*{Second Volume, Version V75: Companion Audit and the Order-Zero Fermion--Higgs Tail}
\addcontentsline{toc}{section}{Second Volume, Version V75: Companion Audit and the Order-Zero Fermion--Higgs Tail}

\subsection*{Purpose and mandatory audit}

V74 gives an adjacent-band energy criterion for a general dressed
coupling.  This mandatory audit selects the first physical sector in which
the criterion can be verified more strongly: a fixed-metric Euclidean
Dirac block perturbed by order-zero gauge and Yukawa multiplication.
The reference Dirac inverse gains one Sobolev order, so the entire
high-frequency relative perturbation is \(O(R^{-1})\), including adjacent
bands.

The Navier--Stokes file inspected was
\[
 \texttt{Renewal\_NS\_Stability\_Research\_Notes\_V31.tex},
\]
with \(887537\) bytes and SHA--256 digest
\[
 \texttt{F1121B62238AD5491BB7CF03635EA9934942EDF573ED8FAAA9EE79E1D38A6696}.
\]
It remains unchanged.

The newest Yang--Mills file inspected was
\[
 \texttt{Renewal\_Yang\_Mills\_Mass\_Gap\_Research\_Notes\_V61.tex},
\]
with \(836836\) bytes and SHA--256 digest
\[
 \texttt{A43BD025326B26F32CF0204C5F10F478D658CA5F1813EB84C654B664253B5CBD}.
\]
A V60 predecessor was also present, with \(814488\) bytes and digest
\[
 \texttt{91ECEE09FEE6B521AC8301E258E4775239112FBC5810EC023E81524A3069B099}.
\]
The frozen YM first volume was unchanged.  All companion files were left
untouched.

YM V61 proves an exact linked-word identity, disintegrates every connected
word at its unique first globally connecting coordinate, and replaces the
restricted accepted suffix by the complete moment suffix, which is a
contraction.  This removes the earlier Green-transfer target rather than
estimating it.  The support-uniform first-connectivity source sum,
including the \([3,3]\) source, remains open; global quartic MTA has not
yet been restored.

\begin{assessment}[V75 audit transfer]
\label{ass:v2v75-audit}
The type-correct SM lesson is to avoid estimating repeated adjacent
transfers when an exact operator factorization exposes a complete inverse
or contractive suffix.  In the order-zero fermion--Higgs sector, the
factorization by the reference Dirac operator exposes one inverse
derivative before any dyadic path expansion.  No Yang--Mills estimate is
imported.
\end{assessment}

\begin{firewall}[Cross-program type boundary]
\label{fw:v2v75-types}
A Wilson complete-moment suffix is not a Dirac resolvent.  The common
content is the proof order: perform the exact global factorization first,
then use its contraction.  The fermionic contraction below is proved from
Sobolev order and does not follow from the YM result.
\end{firewall}

\subsection{High-frequency embedding}
\label{sec:v2v75-embedding}

Let \(M\ge I\) generate the Hilbert scale
\(\|u\|_s=\|M^su\|\), and let
\begin{equation}
 Q_{\ge R}=\mathbf1_{[R,\infty)}(M),
 \qquad R\ge1.
\label{eq:v2v75-QR}
\end{equation}

\begin{lemma}[One high-tail embedding power]
\label{lem:v2v75-embedding}
For every \(s\in\mathbb R\),
\begin{equation}
 \|Q_{\ge R}u\|_s
 \le
 R^{-1}\|Q_{\ge R}u\|_{s+1}.
\label{eq:v2v75-embedding}
\end{equation}
\end{lemma}

\begin{proof}
On the spectral support of \(Q_{\ge R}\), the multiplier \(M^{-1}\) has
norm at most \(R^{-1}\).  Since \(M\) commutes with its spectral
projection,
\[
 M^sQ_{\ge R}
 =
 M^{-1}Q_{\ge R}M^{s+1}.
\]
Take norms.
\end{proof}

\subsection{Reference Dirac factorization}
\label{sec:v2v75-Dirac}

Let \(\mathscr D_0\) be an invertible reference Euclidean Dirac-type
operator on the positive-frequency space, and take
\[
 M=|\mathscr D_0|.
\]
Assume its inverse has the uniform order-minus-one bound
\begin{equation}
 \|\mathscr D_0^{-1}:
 \mathcal H^s\to\mathcal H^{s+1}\|
 \le c_D
 \qquad(-3\le s\le3).
\label{eq:v2v75-D0-inverse}
\end{equation}
For the exact polar scale of \(\mathscr D_0\), one may take \(c_D=1\).

Let \(V(k)\) be a \(C^5\) order-zero perturbation satisfying
\begin{equation}
 \|V(k):\mathcal H^s\to\mathcal H^s\|
 \le v_s
 \qquad(-3\le s\le3).
\label{eq:v2v75-V0}
\end{equation}
Define the compressed high-tail operator
\begin{equation}
 \mathscr D_R(k)
 =
 Q_{\ge R}(\mathscr D_0+V(k))Q_{\ge R}
\label{eq:v2v75-DR}
\end{equation}
on \(Q_{\ge R}\mathcal H\).

\begin{theorem}[Order-zero perturbations are relatively small on the tail]
\label{thm:v2v75-tail}
Put
\begin{equation}
 K_R(k)
 =
 (Q_{\ge R}\mathscr D_0Q_{\ge R})^{-1}
 Q_{\ge R}V(k)Q_{\ge R}.
\label{eq:v2v75-KR}
\end{equation}
Then, for \(-3\le s\le3\),
\begin{equation}
 \|K_R(k):\mathcal H^s\to\mathcal H^s\|
 \le
 \frac{c_Dv_s}{R}.
\label{eq:v2v75-KR-bound}
\end{equation}
If
\begin{equation}
 \vartheta_R
 :=
 \frac{c_D\max_{|s|\le3}v_s}{R}
 <1,
\label{eq:v2v75-tail-small}
\end{equation}
then \(\mathscr D_R(k)\) is invertible and
\begin{equation}
 \|\mathscr D_R(k)^{-1}:
 \mathcal H^s\to\mathcal H^{s+1}\|
 \le
 \frac{c_D}{1-\vartheta_R}
\label{eq:v2v75-tail-inverse}
\end{equation}
for \(-3\le s\le2\).
\end{theorem}

\begin{proof}
The reference spectral projection commutes with
\(\mathscr D_0\) and its inverse.  For \(u\) in the high tail,
\eqref{eq:v2v75-D0-inverse} gives
\[
 \|\mathscr D_0^{-1}Q_{\ge R}VQ_{\ge R}u\|_{s+1}
 \le
 c_Dv_s\|u\|_s.
\]
The output is again in the high tail.  Apply
Lemma~\ref{lem:v2v75-embedding} to lower the output index from \(s+1\)
to \(s\).  This proves
\eqref{eq:v2v75-KR-bound}.

The exact factorization is
\[
 \mathscr D_R
 =
 (Q_{\ge R}\mathscr D_0Q_{\ge R})(I+K_R).
\]
Under
\eqref{eq:v2v75-tail-small}, the Neumann inverse of \(I+K_R\) is bounded
on every displayed scale by \((1-\vartheta_R)^{-1}\).  Compose it with
the reference inverse to obtain
\eqref{eq:v2v75-tail-inverse}.
\end{proof}

\begin{assessment}[Adjacent bands are included]
\label{ass:v2v75-adjacent}
The estimate is taken on the complete high-tail operator.  It neither
separates adjacent dyadic blocks nor iterates a transfer matrix.  Thus the
V73 adjacent-band shift obstruction cannot occur in this order-zero
sector once \(R\) satisfies
\eqref{eq:v2v75-tail-small}.
\end{assessment}

\subsection{Gauge and Yukawa multiplication}
\label{sec:v2v75-multiplication}

On a bounded-geometry Euclidean chart, let
\begin{equation}
 V(k)
 =
 \gamma^\mu A_\mu(k)
 +
 Y\Phi(k),
\label{eq:v2v75-physical-V}
\end{equation}
where \(A_\mu\) and \(\Phi\) are matrix-valued multiplication operators
and \(Y\) is a fixed finite internal Yukawa matrix.

\begin{lemma}[Uniform multiplication bound through three derivatives]
\label{lem:v2v75-multiplication}
Suppose the spectral scale of \(|\mathscr D_0|\) is uniformly equivalent
to the geometric Sobolev scale through orders \([-3,3]\), and suppose
\begin{equation}
 \sup_{|k|\le\kappa}
 \left(
 \|A_\mu(k)\|_{W^{3,\infty}}
 +
 \|\Phi(k)\|_{W^{3,\infty}}
 \right)
 \le C_{\rm coeff}.
\label{eq:v2v75-coeff}
\end{equation}
Then \(V(k)\) satisfies
\eqref{eq:v2v75-V0}, with constants uniform on the chart and compact
background family.
\end{lemma}

\begin{proof}
For nonnegative integer \(s\le3\), the Leibniz rule bounds every weak
derivative of \(fu\) through order \(s\) by a finite sum of
\(\|\partial^\alpha f\|_\infty\|\partial^\beta u\|_2\).  Thus
multiplication by a \(W^{3,\infty}\) matrix is bounded on geometric
\(H^s\).  Duality gives the result for negative integer \(s\), and complex
interpolation gives the intermediate real indices.  Uniform equivalence
of the geometric and spectral norms transfers the bounds.  The finite
gamma and Yukawa matrices change only the constant.
\end{proof}

\begin{corollary}[High-tail fermion--Higgs adjacent sector]
\label{cor:v2v75-fermion-tail}
Under
\eqref{eq:v2v75-coeff}, there is a fixed \(R_0\), independent of the
ultraviolet shell \(n\), such that the complete high-tail
fermion--gauge--Higgs block is invertible with
\eqref{eq:v2v75-tail-inverse} for every \(R\ge R_0\).
\end{corollary}

\begin{proof}
Lemma~\ref{lem:v2v75-multiplication} makes
\(\max_{|s|\le3}v_s\) finite and uniform.  Choose
\(R_0>c_D\max v_s\) and apply
Theorem~\ref{thm:v2v75-tail}.
\end{proof}

\subsection{External derivatives and the V70 weight}
\label{sec:v2v75-external}

Assume additionally, for \(1\le a\le5\),
\begin{equation}
 D^a(\mathscr D_0+V(k)):
 \mathcal H^s\longrightarrow
 \mathcal H^{s+a-1}
\label{eq:v2v75-external-order}
\end{equation}
is uniformly bounded whenever the indices lie in \([-3,3]\).
For an affine external-momentum Dirac insertion, derivatives above first
order vanish; composite vertices require their own estimate.

\begin{theorem}[Weighted differentiated inverse on the high fermion tail]
\label{thm:v2v75-weighted}
Under
Theorem~\ref{thm:v2v75-tail} and
\eqref{eq:v2v75-external-order},
\begin{equation}
 \left\|
 M^{(a+1)/2}
 D^a\mathscr D_R(k)^{-1}
 M^{(a+1)/2}
 \right\|
 \le C_a,
 \qquad 0\le a\le5.
\label{eq:v2v75-weighted}
\end{equation}
\end{theorem}

\begin{proof}
Equation
\eqref{eq:v2v75-tail-inverse} is the order-minus-one inverse hypothesis
of
Theorem~\ref{thm:v2v71-scale-inverse}.
Equation
\eqref{eq:v2v75-external-order} supplies its differentiated kinetic
hypothesis.  Apply that theorem.
\end{proof}

\begin{corollary}[Conditional fermion-sector contribution to V70]
\label{cor:v2v75-V70}
If the finite low core and its core-to-tail Schur feedback satisfy the V57
uniform coercivity conditions, and their external derivatives through order
five obey the V71 scale mappings, then the full positive-frequency
fermion--gauge--Higgs low inverse satisfies the V70 weighted hypothesis.
Under the V69 crossing and rank bounds, its harmonic minimal packet is
\(O(\Lambda_n^{-3/2})\).
\end{corollary}

\begin{proof}
The high tail is controlled by
Theorem~\ref{thm:v2v75-weighted}.  Apply the finite-core Schur formula and
the V57 feedback bound to join it to the exactly treated core.  Differentiate
the finite Schur formula and use the assumed V71 mappings for every core and
feedback derivative; the inverse-derivative theorem then preserves the scale
estimates through order five.  Apply
Corollary~\ref{cor:v2v70-harmonic}.
\end{proof}

\begin{firewall}[The finite core is not automatically uniform]
\label{fw:v2v75-core}
The core threshold \(R_0\) is fixed, but its inverse may still lose
uniformity under volume growth, background variation, chiral zero modes,
or a vanishing Higgs eigenvalue.  V75 assumes the V57 core Schur
conditions in
Corollary~\ref{cor:v2v75-V70}; it does not prove them.
\end{firewall}

\subsection{Why principal-order metric variation is different}
\label{sec:v2v75-order-one}

\begin{counterexample}[An order-one perturbation need not become small]
\label{sep:v2v75-order-one}
Let \(\mathscr D_0=M\) on a positive spectral space and let
\[
 V=cM
\]
for a constant \(c\ne0\).  Then on every high tail
\begin{equation}
 \mathscr D_0^{-1}V=cI.
\label{eq:v2v75-order-one}
\end{equation}
No choice of \(R\) makes the relative norm smaller than \(|c|\).
\end{counterexample}

\begin{proof}
The identity follows by cancellation of \(M^{-1}M\).
\end{proof}

\begin{assessment}[Physical scope]
\label{ass:v2v75-scope}
Gauge connections and Yukawa--Higgs terms are order-zero perturbations of
a fixed-metric Dirac operator and fall under the high-tail theorem when
their coefficient bounds hold.  A change of metric changes the principal
Dirac symbol and is order one.  Its adjacent-band control requires the
energy or sectorial method of V74, together with a quantitative cone of
admissible metrics.
\end{assessment}

\begin{firewall}[No full fermionic determinant theorem]
\label{fw:v2v75-determinant}
V75 proves a resolvent mapping criterion for the one-particle Euclidean
fermion block.  It does not construct the interacting fermion determinant,
sum fermion loops, prove the regulated anomaly primitive, or establish the
full SM continuum measure.
\end{firewall}

\subsection{V75 conclusion and next acquisition}
\label{sec:v2v75-conclusion}

V75 verifies the adjacent high-tail inverse mechanism for the
fixed-metric order-zero fermion--gauge--Higgs block.  The proof treats the
complete tail operator and gains smallness from one reference Dirac
inverse, not from a dyadic transfer expansion.  The remaining fermionic
issues are the finite core, external composite derivative bounds, and
determinant assembly.

V76 will address the principal-order metric perturbation.  It will compare
two uniformly equivalent Euclidean metrics at the symbol level, prove a
relative form cone for their squared Dirac operators, and determine when
the first-order Dirac inverse inherits uniform scale bounds through
\([-3,3]\).
\clearpage
\section*{Second Volume, Version V76: Metric Form Cone for the Euclidean Dirac Inverse}
\addcontentsline{toc}{section}{Second Volume, Version V76: Metric Form Cone for the Euclidean Dirac Inverse}

\subsection*{Purpose}

V75 closes the high-tail inverse for order-zero gauge and Yukawa
perturbations of a fixed-metric Dirac operator.  A metric change modifies
the principal symbol and does not become relatively small merely by moving
to high frequency.  This note treats that principal-order change through
a form cone.

After a fixed unitary identification of the spinor bundles, the difference
of two Dirac operators is a first-order coefficient perturbation.  A
small \(W^{1,\infty}\) metric cone gives a relative graph bound
\[
 \|(D_g-D_0)u\|
 \le a\|D_0u\|+b\|u\|.
\]
If the combined dimensionless constant
\(\sqrt{a^2+(b/m)^2}\) is below one, the massive squared Dirac forms are
uniformly equivalent.  Uniform higher-order elliptic estimates then lift
the inverse to the full V71 scale interval.  The massless kernel and
global spinor-identification problems remain separate.

\subsection{Fixed bundle and relative graph bound}
\label{sec:v2v76-fixed-bundle}

Let \(M\) be a compact Euclidean spin manifold, or one bounded-geometry
chart with fixed boundary conditions.  Fix a reference metric \(g_0\).
Assume a chosen unitary identification transports the spinor spaces for
metrics \(g\) in a neighbourhood of \(g_0\) to one Hilbert space
\(\mathcal H\), with common first-order form domain \(H^1\).  Denote the
transported self-adjoint Dirac operators by \(D_g\) and \(D_0\).

In a fixed local frame, write
\begin{equation}
 D_g-D_0
 =
 A_g^j\nabla_j^{0}+B_g,
\label{eq:v2v76-difference}
\end{equation}
where \(A_g^j\) is the principal coefficient difference and \(B_g\)
contains the transported spin-connection and density terms.

\begin{lemma}[Coefficient control gives a relative graph bound]
\label{lem:v2v76-graph}
Suppose
\begin{align}
 \max_j\|A_g^j\|_\infty
 &\le\varepsilon_1,
\label{eq:v2v76-A}\\
 \|B_g\|_\infty
 &\le\varepsilon_0,
\label{eq:v2v76-B}\\
 \|\nabla^0u\|
 &\le C_G(\|D_0u\|+\|u\|)
\label{eq:v2v76-Garding}
\end{align}
uniformly on the common domain.  Then
\begin{equation}
 \|(D_g-D_0)u\|
 \le
 a\|D_0u\|+b\|u\|,
\label{eq:v2v76-relative}
\end{equation}
where one may take
\begin{equation}
 a=C_1C_G\varepsilon_1,
 \qquad
 b=C_1C_G\varepsilon_1+\varepsilon_0,
\label{eq:v2v76-ab}
\end{equation}
and \(C_1\) depends only on the fixed frame dimension and bundle norm.
\end{lemma}

\begin{proof}
Equation
\eqref{eq:v2v76-difference} and Cauchy--Schwarz in the finite frame index
give
\[
 \|(D_g-D_0)u\|
 \le
 C_1\varepsilon_1\|\nabla^0u\|
 +
 \varepsilon_0\|u\|.
\]
Insert
\eqref{eq:v2v76-Garding} and collect the two terms.
\end{proof}

\begin{proposition}[A small metric cone supplies the coefficients]
\label{prop:v2v76-metric-cone}
Suppose the bundle identification, orthonormal frame, Clifford
multiplication, and spin connection depend \(C^1\)-smoothly on the metric
in a bounded-geometry \(W^{1,\infty}\) neighbourhood of \(g_0\).  Then
there is \(C_{\rm met}\) such that
\begin{equation}
 \varepsilon_1+\varepsilon_0
 \le
 C_{\rm met}\|g-g_0\|_{W^{1,\infty}}.
\label{eq:v2v76-metric-coeff}
\end{equation}
Consequently \(a,b\) in
\eqref{eq:v2v76-relative} tend to zero with the metric radius.
\end{proposition}

\begin{proof}
The principal coefficient is obtained from the inverse vielbein and
Clifford multiplication, so its difference is pointwise Lipschitz in
\(g-g_0\) on a uniformly positive metric cone.  The spin-connection
difference contains one derivative of the vielbein and algebraic inverse
metric factors; it is bounded by the \(W^{1,\infty}\) metric difference.
The unitary density correction has the same dependence.  Smoothness on
the closed bounded-geometry cone makes the local Lipschitz constants
uniform.  A finite chart partition and the fixed transition functions
give
\eqref{eq:v2v76-metric-coeff}.
\end{proof}

\begin{firewall}[Spinor identification is geometric data]
\label{fw:v2v76-identification}
A pointwise comparison of metric matrices does not canonically identify
their spinor Hilbert spaces or boundary domains.  V76 assumes a fixed
unitary identification with a common form domain.  Proving compatibility
of those identifications around nontrivial loops in metric space, with
boundaries and changing spin structures, is outside the present local
cone theorem.
\end{firewall}

\subsection{Massive squared-form comparison}
\label{sec:v2v76-form}

Fix \(m>0\) and define
\begin{align}
 q_0[u]
 &=
 \|D_0u\|^2+m^2\|u\|^2,
\label{eq:v2v76-q0}\\
 q_g[u]
 &=
 \|D_gu\|^2+m^2\|u\|^2.
\label{eq:v2v76-qg}
\end{align}

\begin{theorem}[Metric Dirac form cone]
\label{thm:v2v76-form-cone}
Assume
\eqref{eq:v2v76-relative} and put
\begin{equation}
 \kappa
 =
 \sqrt{a^2+(b/m)^2}.
\label{eq:v2v76-kappa}
\end{equation}
If \(\kappa<1\), then
\begin{equation}
 (1-\kappa)^2q_0[u]
 \le
 q_g[u]
 \le
 (1+\kappa)^2q_0[u]
\label{eq:v2v76-form-cone}
\end{equation}
for every \(u\) in the common domain.
\end{theorem}

\begin{proof}
Let
\[
 \mathcal A_0u=(D_0u,mu),
 \qquad
 \mathcal A_gu=(D_gu,mu)
\]
as maps into \(\mathcal H\oplus\mathcal H\).  Then
\[
 \|(\mathcal A_g-\mathcal A_0)u\|
 =
 \|(D_g-D_0)u\|.
\]
With \(X=\|D_0u\|\) and \(Y=m\|u\|\),
\eqref{eq:v2v76-relative} and Cauchy--Schwarz give
\[
 \|(D_g-D_0)u\|
 \le
 aX+\frac bmY
 \le
 \kappa(X^2+Y^2)^{1/2}
 =
 \kappa q_0[u]^{1/2}.
\]
The triangle and reverse-triangle inequalities yield
\[
 (1-\kappa)q_0[u]^{1/2}
 \le q_g[u]^{1/2}
 \le(1+\kappa)q_0[u]^{1/2}.
\]
Square the inequalities.
\end{proof}

\begin{corollary}[Uniform positive squared Dirac operator]
\label{cor:v2v76-square}
Let
\[
 L_g=D_g^2+m^2
\]
be defined by the closed form \(q_g\).  Under
Theorem~\ref{thm:v2v76-form-cone},
\begin{equation}
 L_g\ge(1-\kappa)^2L_0
\label{eq:v2v76-square-order}
\end{equation}
in form order, and
\begin{equation}
 \|L_0^{1/2}L_g^{-1/2}\|
 \le
 (1-\kappa)^{-1}.
\label{eq:v2v76-square-inverse}
\end{equation}
\end{corollary}

\begin{proof}
The first statement is the lower form inequality.  Given \(v\) in the
range of \(L_g^{1/2}\), put \(u=L_g^{-1/2}v\) in that inequality.  Then
extend from the form range by density:
\[
 \|L_0^{1/2}L_g^{-1/2}v\|
 \le
 (1-\kappa)^{-1}\|v\|.
\]
\end{proof}

\subsection{Massive first-order inverse}
\label{sec:v2v76-massive}

Define
\begin{equation}
 \mathscr D_g=D_g+im.
\label{eq:v2v76-massive-Dirac}
\end{equation}
Since \(D_g\) is self-adjoint,
\begin{equation}
 \|\mathscr D_gu\|^2
 =
 q_g[u].
\label{eq:v2v76-normal}
\end{equation}

\begin{theorem}[Base graph-norm inverse]
\label{thm:v2v76-base-inverse}
Under
Theorem~\ref{thm:v2v76-form-cone}, \(\mathscr D_g\) is invertible and
\begin{equation}
 \|\mathscr D_g^{-1}f\|_{q_0}
 \le
 \frac1{1-\kappa}\|f\|,
\label{eq:v2v76-base-inverse}
\end{equation}
where \(\|u\|_{q_0}=q_0[u]^{1/2}\).
\end{theorem}

\begin{proof}
The spectral theorem for self-adjoint \(D_g\) gives
\[
 \mathscr D_g^{-1}
 =
 (D_g-im)(D_g^2+m^2)^{-1},
\]
so the inverse exists and is bounded on \(\mathcal H\).  If
\(u=\mathscr D_g^{-1}f\), then
\eqref{eq:v2v76-normal} and the lower form inequality give
\[
 (1-\kappa)\|u\|_{q_0}
 \le
 q_g[u]^{1/2}
 =
 \|f\|.
\]
\end{proof}

\begin{assessment}[Why the mass appears in the cone radius]
\label{ass:v2v76-mass}
The lower-order coefficient \(b\) is measured against \(m\) in
\eqref{eq:v2v76-kappa}.  A small fermion mass therefore permits a smaller
metric and connection neighbourhood.  This dependence cannot be deleted
from the present graph-form proof.
\end{assessment}

\subsection{Lifting the inverse through the V71 scale interval}
\label{sec:v2v76-scale}

Let
\[
 M_0=(D_0^2+m^2)^{1/2}
\]
generate the reference Hilbert scale.  The \(q_0\)-norm is exactly the
\(\mathcal H^1\)-norm associated with \(M_0\).

\begin{theorem}[Uniform elliptic lift]
\label{thm:v2v76-elliptic-lift}
Assume the form cone with one \(\kappa<1\).  Suppose, for every integer
\(s=0,1,2,3\), the family satisfies the uniform graph estimate
\begin{equation}
 \|u\|_{s+1}
 \le
 C_s\bigl(
 \|\mathscr D_gu\|_s+\|u\|
 \bigr),
\label{eq:v2v76-positive-elliptic}
\end{equation}
and the adjoint family \(D_g-im\) satisfies the analogous estimates.
Then
\begin{equation}
 \mathscr D_g^{-1}:
 \mathcal H^s\longrightarrow\mathcal H^{s+1}
\label{eq:v2v76-scale-inverse}
\end{equation}
is uniformly bounded for every real \(-3\le s\le3\) for which the target
index is used.  The constants depend on
\((1-\kappa)^{-1}\) and the displayed elliptic constants.
\end{theorem}

\begin{proof}
At \(s=0\), Theorem~\ref{thm:v2v76-base-inverse} controls both
\(\|u\|\) and \(\|u\|_1\) by \(\|f\|\) for
\(\mathscr D_gu=f\).  For positive integer \(s\), apply
\eqref{eq:v2v76-positive-elliptic} to \(u=\mathscr D_g^{-1}f\):
\[
 \|u\|_{s+1}
 \le
 C_s(\|f\|_s+\|u\|)
 \le
 C_s'\|f\|_s.
\]
For \(r=0,1,2,3\), the same positive-index argument for
\((\mathscr D_g^*)^{-1}=(D_g-im)^{-1}\) gives
\((\mathscr D_g^*)^{-1}:\mathcal H^r\to\mathcal H^{r+1}\).
Taking Hilbert-space adjoints therefore gives
\(\mathscr D_g^{-1}:\mathcal H^{-r-1}\to\mathcal H^{-r}\).
Together with the nonnegative estimates this covers every required integer
source index from \(-3\) through \(3\).  Complex interpolation between
consecutive integer indices gives all intermediate real \(s\).
\end{proof}

\begin{firewall}[Form equivalence is not higher elliptic regularity]
\label{fw:v2v76-regularity}
Theorem~\ref{thm:v2v76-form-cone} proves only the base graph norm.
The estimates
\eqref{eq:v2v76-positive-elliptic} require uniform coefficient
regularity, boundary complementing conditions, and a common
pseudodifferential or elliptic calculus.  They are stated separately so
that first-order form coercivity is not misread as three additional
derivatives.
\end{firewall}

\subsection{Insertion into the weighted low-resolvent programme}
\label{sec:v2v76-V70}

Suppose external-momentum differentiation, distinct from variation of the
background metric, satisfies
\begin{equation}
 D_k^j\mathscr D_g(k):
 \mathcal H^s\longrightarrow\mathcal H^{s+j-1},
 \qquad
 1\le j\le5,
\label{eq:v2v76-external}
\end{equation}
uniformly on the indices used in V71.

\begin{theorem}[Metric-cone weighted inverse]
\label{thm:v2v76-weighted}
Under
Theorem~\ref{thm:v2v76-elliptic-lift} and
\eqref{eq:v2v76-external},
\begin{equation}
 \left\|
 M_0^{(j+1)/2}
 D_k^j\mathscr D_g(k)^{-1}
 M_0^{(j+1)/2}
 \right\|
 \le C_j,
 \qquad 0\le j\le5.
\label{eq:v2v76-weighted}
\end{equation}
Consequently the corresponding minimal harmonic low packet satisfies the
V70 estimate whenever the crossing and rank hypotheses hold.
\end{theorem}

\begin{proof}
The elliptic lift supplies the order-minus-one inverse on the full V71
scale interval.  Equation
\eqref{eq:v2v76-external} supplies the differentiated kinetic orders.
Apply
Theorem~\ref{thm:v2v71-scale-inverse}, then
Corollary~\ref{cor:v2v70-harmonic}.
\end{proof}

\begin{counterexample}[Principal-symbol cancellation destroys a uniform cone]
\label{sep:v2v76-principal}
Let \(D_0\) have positive eigenvalues \(\lambda_j\to\infty\), and on their
span take
\[
 D_g=D_0-D_0=0.
\]
Then
\[
 \frac{q_g[e_j]}{q_0[e_j]}
 =
 \frac{m^2}{\lambda_j^2+m^2}
 \longrightarrow0.
\]
Thus a relative principal coefficient of size one can destroy every
uniform comparison with the reference first-order graph norm, even though
the massive operator \(im\) remains pointwise invertible.
\end{counterexample}

\begin{proof}
The displayed ratio follows directly from
\eqref{eq:v2v76-q0}--\eqref{eq:v2v76-qg}.
\end{proof}

\begin{firewall}[Physical sectors not covered]
\label{fw:v2v76-scope}
V76 treats a local Euclidean metric cone, a common spinor domain, and a
massive Dirac shift.  It does not treat massless chiral zero modes,
Lorentzian hyperbolic propagation, determinant phases, gravitational
conformal instability, or global changes of spin structure.  It also
does not prove the uniform elliptic estimates for the full boundary
problem.
\end{firewall}

\subsection{V76 conclusion and next acquisition}
\label{sec:v2v76-conclusion}

V76 proves a quantitative form cone for principal-order metric
perturbations of the Euclidean Dirac operator and identifies the exact
dimensionless radius
\(\sqrt{a^2+(b/m)^2}<1\).  Under separately stated uniform elliptic
regularity, the first-order inverse has the scale bounds required by V71
and therefore the V70 mixed-band packet estimate.

V77 will remove the artificial massive shift on the complement of the
Dirac kernel.  It will formulate a reduced inverse with a moving Riesz
projection, prove the scale bound under a uniform nonzero spectral gap,
and show by a collapsing-eigenvalue example why the massless infrared
sector cannot be absorbed into the metric cone.
\clearpage
\section*{Second Volume, Version V77: The Reduced Massless Dirac Inverse and Its Moving Kernel}
\addcontentsline{toc}{section}{Second Volume, Version V77: The Reduced Massless Dirac Inverse and Its Moving Kernel}
\label{sec:v2v77}

\subsection{Purpose and relation to the earlier spectral analysis}
\label{sec:v2v77-purpose}

V76 obtained the inverse of \(D_g+im\) in a small metric form cone.  The
mass \(m\) simultaneously removed the kernel and supplied the scale against
which the zero-order perturbation was measured.  This version removes that
auxiliary shift without pretending that the infrared problem disappears.

V34--V39 developed moving Riesz projections for spectral clusters near a
positive shell threshold.  We use that construction here rather than repeat
it.  The new point is the operator at the spectral origin: the kernel is
retained as a finite-dimensional infrared sector, the Dirac operator is
inverted only on its orthogonal complement, and derivatives of both the
kernel projection and the reduced inverse are computed exactly.

\subsection{A uniform zero-mode gap on a fixed Hilbert scale}
\label{sec:v2v77-gap}

Let \(I\) be a compact interval of metric or background parameters.  Work
after the fixed-bundle identifications of V76 on a Hilbert scale
\(\{\mathcal H^s\}_{s\in\mathbb R}\).  Let
\[
 D(t):\mathcal H^1\subset\mathcal H\longrightarrow\mathcal H,
 \qquad t\in I,
\]
be self-adjoint with compact resolvent and common domain
\(\mathcal H^1\).  Assume that, for every integer
\(-4\le s\le4\),
\begin{equation}
 D(t):\mathcal H^{s+1}\longrightarrow\mathcal H^s
\label{eq:v2v77-scale-map}
\end{equation}
is a \(C^5\) family with uniformly bounded derivatives, and that the
uniform elliptic estimate
\begin{equation}
 \|u\|_{s+1}
 \le E_s\bigl(\|D(t)u\|_s+\|u\|_s\bigr)
\label{eq:v2v77-elliptic}
\end{equation}
holds on the same window.  Finally assume that there is a
\(\delta>0\) such that
\begin{equation}
 \operatorname{spec}D(t)\cap
 \bigl((-\delta,\delta)\setminus\{0\}\bigr)=\varnothing
 \qquad(t\in I).
\label{eq:v2v77-zero-gap}
\end{equation}
Thus zero modes are allowed, but every nonzero eigenvalue is uniformly
separated from zero.

Put \(\Gamma=\{z\in\mathbb C:|z|=\delta/2\}\), oriented positively, and
define
\begin{align}
 P(t)
 &=
 \frac1{2\pi i}\int_\Gamma(z-D(t))^{-1}\,\mathrm dz,
 \label{eq:v2v77-P}\\
 Q(t)&=I-P(t).
 \label{eq:v2v77-Q}
\end{align}

\begin{lemma}[Scale resolvent on the zero contour]
\label{lem:v2v77-contour}
For every integer \(-4\le s\le4\),
\begin{equation}
 \sup_{\substack{t\in I\\z\in\Gamma}}
 \|(z-D(t))^{-1}\|_{\mathcal H^s\to\mathcal H^{s+1}}
 \le C_{s,\delta}.
\label{eq:v2v77-contour}
\end{equation}
Consequently \(P(t)\) and \(Q(t)\) are uniformly bounded on each
\(\mathcal H^s\) in this interval.
\end{lemma}

\begin{proof}
Self-adjointness and
\eqref{eq:v2v77-zero-gap} give
\[
 \|(z-D(t))^{-1}\|_{\mathcal H\to\mathcal H}
 \le \operatorname{dist}(z,\operatorname{spec}D(t))^{-1}
 \le 2\delta^{-1}.
\]
For \(u=(z-D(t))^{-1}f\), the elliptic estimate gives
\[
 \|u\|_{s+1}
 \le E_s\bigl(\|f+zu\|_s+\|u\|_s\bigr)
 \le C_{s,\delta}\|f\|_s
\]
at nonnegative integer indices, by induction in \(s\).  Apply the same
argument to the adjoint resolvent and use Hilbert duality for the negative
indices.  The contour formula and the continuous embeddings
\(\mathcal H^{s+1}\hookrightarrow\mathcal H^s\) give the assertion for
\(P(t)\); subtracting from the identity gives it for \(Q(t)\).
\end{proof}

\begin{proposition}[The Riesz sector is exactly the kernel]
\label{prop:v2v77-kernel}
For every \(t\in I\),
\[
 \operatorname{Ran}P(t)=\ker D(t).
\]
Its dimension is finite and independent of \(t\).  Moreover
\begin{equation}
 D(t)P(t)=P(t)D(t)=0,
 \qquad
 D(t)Q(t)=Q(t)D(t).
\label{eq:v2v77-reduction}
\end{equation}
\end{proposition}

\begin{proof}
The contour encloses only the spectral point zero, so the spectral theorem
identifies \(P(t)\) with the orthogonal projection onto \(\ker D(t)\).
Compact resolvent makes this eigenspace finite-dimensional.
Lemma~\ref{lem:v2v77-contour} and the resolvent identity make \(P(t)\)
norm-continuous.  Two orthogonal projections at operator distance less than
one have equal rank; a finite cover of the connected interval \(I\) shows
that the rank is constant.  The reduction identities follow from spectral
calculus.
\end{proof}

\subsection{Construction and scale bound for the pseudoinverse}
\label{sec:v2v77-reduced}

Define the Borel function
\[
 \gamma(\lambda)=
 \begin{cases}
  0,&\lambda=0,\\
  \lambda^{-1},&\lambda\ne0,
 \end{cases}
\]
on \(\operatorname{spec}D(t)\), and set
\begin{equation}
 G(t)=\gamma(D(t)).
\label{eq:v2v77-G}
\end{equation}
This is the Moore--Penrose inverse of \(D(t)\).

\begin{theorem}[Uniform reduced massless inverse]
\label{thm:v2v77-reduced}
Under
\eqref{eq:v2v77-scale-map}--\eqref{eq:v2v77-zero-gap},
\begin{align}
 D(t)G(t)&=G(t)D(t)=Q(t),
 \label{eq:v2v77-pseudo-identities}\\
 P(t)G(t)&=G(t)P(t)=0,
 \label{eq:v2v77-pseudo-orthogonal}\\
 \|G(t)\|_{\mathcal H\to\mathcal H}
 &\le\delta^{-1}.
 \label{eq:v2v77-L2}
\end{align}
For every real \(-3\le s\le3\),
\begin{equation}
 G(t):\mathcal H^s\longrightarrow\mathcal H^{s+1}
\label{eq:v2v77-G-scale}
\end{equation}
is uniformly bounded.
\end{theorem}

\begin{proof}
The first three statements follow pointwise from the spectral theorem and
the fact that
\(|\lambda|\ge\delta\) on the nonzero spectrum.  Let
\(u=G(t)f\).  Then
\[
 D(t)u=Q(t)f,
 \qquad
 \|u\|\le\delta^{-1}\|f\|.
\]
For a nonnegative integer \(s\le3\), use
Lemma~\ref{lem:v2v77-contour} to bound \(Q(t)\) on \(\mathcal H^s\), and
apply \eqref{eq:v2v77-elliptic}:
\[
 \|u\|_{s+1}
 \le
 E_s\bigl(\|Q(t)f\|_s+\|u\|_s\bigr)
 \le C_{s,\delta}\|f\|_s.
\]
The induction begins with \eqref{eq:v2v77-L2}.  Since \(G(t)\) is
self-adjoint, the adjoints of the positive-index estimates give
\(G(t):\mathcal H^{-r-1}\to\mathcal H^{-r}\) for \(r=0,1,2,3\).
Interpolation supplies the remaining real indices.
\end{proof}

\begin{assessment}[What replaced the mass]
\label{ass:v2v77-replacement}
In V76 the inverse constant depended on \(m^{-1}\).  In the reduced
massless problem it depends on \(\delta^{-1}\).  The replacement is exact:
the kernel has been projected out, but the smallest nonzero Dirac
eigenvalue is still the infrared scale controlling the inverse.
\end{assessment}

\subsection{Differentiating the moving kernel}
\label{sec:v2v77-P-derivatives}

Write
\[
 R_z(t)=(z-D(t))^{-1},
 \qquad
 D^{(a)}(t)=\frac{\mathrm d^aD(t)}{\mathrm dt^a}.
\]
For an ordered composition
\(\boldsymbol a=(a_1,\ldots,a_\ell)\) of \(r\), put
\[
 |\boldsymbol a|=a_1+\cdots+a_\ell=r,
 \qquad
 \binom{r}{\boldsymbol a}
 =\frac{r!}{a_1!\cdots a_\ell!}.
\]

\begin{proposition}[Exact differentiated Riesz formula]
\label{prop:v2v77-P-derivatives}
For \(1\le r\le5\),
\begin{equation}
 P^{(r)}(t)
 =
 \frac1{2\pi i}
 \int_\Gamma
 \sum_{\ell=1}^r
 \ \sum_{\substack{a_1+\cdots+a_\ell=r\\a_j\ge1}}
 \binom{r}{a_1,\ldots,a_\ell}
 R_zD^{(a_1)}R_z\cdots
 D^{(a_\ell)}R_z\,\mathrm dz.
\label{eq:v2v77-P-derivatives}
\end{equation}
On every scale index needed below these derivatives are uniformly bounded,
with constants polynomial in the coefficient bounds, the elliptic
constants, and \(\delta^{-1}\).
\end{proposition}

\begin{proof}
Differentiating
\((z-D(t))R_z(t)=I\) once gives
\[
 R_z'(t)=R_z(t)D'(t)R_z(t).
\]
Induction with the Leibniz rule gives the ordered-composition formula for
\(R_z^{(r)}\), including the displayed multinomial coefficient.  The
contour is fixed by the uniform gap, so differentiation under the integral
is valid.  Apply
Lemma~\ref{lem:v2v77-contour} between the uniformly bounded maps
\(D^{(a)}:\mathcal H^{s+1}\to\mathcal H^s\), and integrate over the contour.
The one-index margin in
\eqref{eq:v2v77-scale-map} covers the endpoint compositions.
\end{proof}

At first order the contour formula can be rewritten without a contour.

\begin{lemma}[Off-diagonal motion of the kernel]
\label{lem:v2v77-P-prime}
The derivative of the kernel projection satisfies
\begin{equation}
 P'
 =
 -G D'P-PD'G.
\label{eq:v2v77-P-prime}
\end{equation}
In particular
\begin{equation}
 PP'P=QP'Q=0.
\label{eq:v2v77-P-offdiag}
\end{equation}
\end{lemma}

\begin{proof}
Differentiate \(DP=0\) and multiply on the left by \(G\):
\[
 QP'=-GD'P.
\]
Differentiating \(PD=0\) and multiplying on the right by \(G\) gives
\[
 P'Q=-PD'G.
\]
The derivative of \(P^2=P\) is off-diagonal:
\(P'=QP'P+PP'Q\).  Combining these identities proves the formula and
\eqref{eq:v2v77-P-offdiag}.
\end{proof}

\subsection{Exact derivative recursion for the Moore--Penrose inverse}
\label{sec:v2v77-G-derivatives}

\begin{theorem}[Differentiated reduced inverse]
\label{thm:v2v77-G-derivatives}
For every \(1\le r\le5\), the derivative \(G^{(r)}(t)\) is determined by
the lower derivatives through
\begin{align}
 G^{(r)}
 ={}&
 -G P^{(r)}
 -\sum_{a=1}^r\binom ra
   G D^{(a)}G^{(r-a)}
 \nonumber\\
 &-\sum_{a=1}^r\binom ra
   P^{(a)}G^{(r-a)}.
\label{eq:v2v77-G-recursion}
\end{align}
In particular,
\begin{align}
 G'
 &=-GD'G-GP'-P'G
 \label{eq:v2v77-G-prime-a}\\
 &=-GD'G+G^2D'P+PD'G^2.
 \label{eq:v2v77-G-prime-b}
\end{align}
For every \(r\le5\) and every real \(-3\le s\le3\) lying in the declared
scale window,
\begin{equation}
 \sup_{t\in I}
 \|G^{(r)}(t)\|_{\mathcal H^s\to\mathcal H^{s+1}}
 \le C_{r,s,\delta}.
\label{eq:v2v77-G-derivative-bound}
\end{equation}
\end{theorem}

\begin{proof}
Differentiate \(DG=Q=I-P\) \(r\) times:
\[
 D G^{(r)}
 =
 -P^{(r)}
 -\sum_{a=1}^r\binom ra D^{(a)}G^{(r-a)}.
\]
Multiplication by \(G\) on the left gives the \(QG^{(r)}\) component.
The \(PG^{(r)}\) component follows by differentiating \(PG=0\):
\[
 PG^{(r)}
 =
 -\sum_{a=1}^r\binom ra P^{(a)}G^{(r-a)}.
\]
Adding the two components proves
\eqref{eq:v2v77-G-recursion}.  At \(r=1\) this is
\eqref{eq:v2v77-G-prime-a}; substituting
\eqref{eq:v2v77-P-prime} gives
\eqref{eq:v2v77-G-prime-b}.

The bound follows by induction on \(r\).  Each \(G\) gains one scale
index, each \(D^{(a)}\) loses one, and each \(P^{(a)}\) is bounded on the
scale index on which it occurs.  Thus every term on the right of
\eqref{eq:v2v77-G-recursion} maps
\(\mathcal H^s\) to \(\mathcal H^{s+1}\).  Theorem
\ref{thm:v2v77-reduced} starts the induction and Proposition
\ref{prop:v2v77-P-derivatives} supplies the projector bounds.
Interpolation handles noninteger \(s\).
\end{proof}

\begin{firewall}[The fixed-kernel inverse formula is incomplete]
\label{fw:v2v77-fixed-P}
The formula \(G'=-GD'G\) is valid only when the null projection does not
move in the chosen trivialization.  In general the two terms containing
\(P'\) in \eqref{eq:v2v77-G-prime-a} are required.  Dropping them violates
the differentiated equations \(DG=GD=Q\).
\end{firewall}

\subsection{Sharp obstruction when the nonzero gap collapses}
\label{sec:v2v77-collapse}

\begin{counterexample}[A norm-small perturbation can destroy the reduced bound]
\label{sep:v2v77-collapse}
On \(\ell^2(\mathbb N_0)\), with standard basis
\(\{e_j\}_{j\ge0}\), let
\[
 D_\varepsilon e_0=0,
 \qquad
 D_\varepsilon e_1=\varepsilon e_1,
 \qquad
 D_\varepsilon e_j=j e_j\quad(j\ge2),
\]
for \(0<\varepsilon\le1\).  These self-adjoint operators have one common
domain, and changing \(\varepsilon\) is a bounded rank-one perturbation.
For \(P_\varepsilon=|e_0\rangle\langle e_0|\), however,
\[
 G_\varepsilon e_1=\varepsilon^{-1}e_1,
 \qquad
 \|G_\varepsilon\|=\varepsilon^{-1}.
\]
At \(\varepsilon=0\) the kernel dimension jumps.  Hence neither
coefficient closeness nor a common elliptic principal symbol produces a
uniform massless reduced inverse without
\eqref{eq:v2v77-zero-gap}.
\end{counterexample}

\begin{firewall}[Infinite-volume massless sectors]
\label{fw:v2v77-infinite-volume}
Compact resolvent is essential to the present finite-dimensional kernel
sector.  On an infinite-volume limit, continuous spectrum may reach zero;
then zero is not enclosed by a Riesz contour and the reduced inverse can
be unbounded even when the point kernel is trivial.  That problem requires
weighted limiting-absorption or inclusive infrared estimates, not the
finite-volume pseudoinverse theorem.
\end{firewall}

\begin{firewall}[Index protection does not prove the gap]
\label{fw:v2v77-index}
A chiral index can force a net number of zero modes, but it does not rule
out the creation of opposite-chirality pairs arbitrarily close to zero.
The uniform nonzero gap and the constant total kernel rank used here
therefore remain analytic hypotheses.
\end{firewall}

\subsection{V77 conclusion and next acquisition}
\label{sec:v2v77-conclusion}

V77 replaces the massive inverse by a uniform Moore--Penrose inverse on
the complement of the moving Dirac kernel.  It proves the full scale bound
and the exact derivative recursion through order five.  The constants
retain the unavoidable factor associated with the smallest nonzero Dirac
eigenvalue.  A collapsing eigenvalue proves that this dependence is sharp
at the level of uniform invertibility.

The complement \(Q(t)\mathcal H\) still moves with the background.
V78 will use Kato transport to identify it with a fixed reference
complement, compute the transported kinetic operator and connection
terms, and determine precisely when the V71 weighted derivative theorem
survives this infrared reduction.
\clearpage
\section*{Second Volume, Version V78: Kato Trivialization of the Reduced Dirac Complement}
\addcontentsline{toc}{section}{Second Volume, Version V78: Kato Trivialization of the Reduced Dirac Complement}
\label{sec:v2v78}

\subsection{The fixed-fibre problem}
\label{sec:v2v78-purpose}

V77 constructed \(G(k)\), the Moore--Penrose inverse of a Dirac family,
on the moving space \(Q(k)\mathcal H\).  V71, however, differentiates an
ordinary inverse on one fixed Hilbert scale.  Directly identifying these
two statements would omit derivatives of the moving projection.

This version makes the identification explicitly.  A local Kato unitary
carries \(Q(k)\mathcal H\) to a fixed reference complement.  The
transported Dirac block becomes an ordinary invertible family, and its
inverse is the transported Moore--Penrose inverse.  Under the V71
differential-order hypotheses, the Riesz projection gains enough scale
order that differentiating the Kato unitary does not spoil the weighted
inverse estimate.

\subsection{External derivative assumptions on an extended scale}
\label{sec:v2v78-assumptions}

Let \(B\subset\mathbb R^q\) be a compact star-shaped parameter chart with
centre \(k_0\).  Let
\(\{\mathcal H^s\}_{s\in\mathbb R}\) be a Hilbert scale.  We use an
extended interval \(J\) containing \([-8,8]\); this margin is only to
justify every intermediate composition through derivative order five.

Let \(D(k)\) be a \(C^5\) family of self-adjoint compact-resolvent
operators with common domain and suppose
\begin{equation}
 \operatorname{spec}D(k)\cap
 \bigl((-\delta,\delta)\setminus\{0\}\bigr)=\varnothing
 \qquad(k\in B)
\label{eq:v2v78-gap}
\end{equation}
for one \(\delta>0\).  On
\(\Gamma=\{z:|z|=\delta/2\}\), assume
\begin{equation}
 R_z(k)=(z-D(k))^{-1}:
 \mathcal H^s\longrightarrow\mathcal H^{s+1}
\label{eq:v2v78-resolvent-order}
\end{equation}
uniformly whenever \(s,s+1\in J\).  Finally, for every multi-index
\(\alpha\) with \(1\le|\alpha|\le5\), assume
\begin{equation}
 \partial_k^\alpha D(k):
 \mathcal H^s\longrightarrow
 \mathcal H^{s+|\alpha|-1}
\label{eq:v2v78-D-order}
\end{equation}
uniformly whenever the indicated indices lie in \(J\).

Define \(P(k)\), \(Q(k)\), and \(G(k)\) by
\eqref{eq:v2v77-P}, \eqref{eq:v2v77-Q}, and
\eqref{eq:v2v77-G}.  The gap makes the rank of \(P(k)\) constant on
\(B\).

\begin{lemma}[Every external derivative of the null projection gains scale]
\label{lem:v2v78-P-order}
For \(0\le|\alpha|\le5\),
\begin{equation}
 \partial_k^\alpha P(k):
 \mathcal H^s\longrightarrow
 \mathcal H^{s+|\alpha|+1}
\label{eq:v2v78-P-order}
\end{equation}
is uniformly bounded whenever the source, target, and intermediate
indices lie in \(J\).
\end{lemma}

\begin{proof}
For \(|\alpha|=0\), use the Riesz formula
\[
 P(k)=\frac1{2\pi i}\int_\Gamma R_z(k)\,\mathrm dz
\]
and \eqref{eq:v2v78-resolvent-order}.  For
\(r=|\alpha|\ge1\), polarize the ordered-composition formula
\eqref{eq:v2v77-P-derivatives}.  A typical term has \(\ell+1\)
resolvents and differentiated kinetic factors of orders
\(a_1,\ldots,a_\ell\), where
\(a_1+\cdots+a_\ell=r\).  Its total scale gain is
\[
 (\ell+1)+\sum_{j=1}^{\ell}(a_j-1)=r+1.
\]
All combinatorial sums are finite for \(r\le5\), and the contour length is
fixed.  The assumed uniform bounds prove the claim.
\end{proof}

\subsection{Radial Kato transport}
\label{sec:v2v78-Kato}

For \(k\in B\) put \(k_\rho=k_0+\rho(k-k_0)\), \(0\le\rho\le1\), and
define
\begin{equation}
 \mathcal K(\rho,k)
 =
 [\partial_\rho P(k_\rho),P(k_\rho)].
\label{eq:v2v78-K}
\end{equation}
Let \(U(\rho,k)\) solve
\begin{equation}
 \partial_\rho U(\rho,k)
 =
 \mathcal K(\rho,k)U(\rho,k),
 \qquad
 U(0,k)=I,
\label{eq:v2v78-U-ode}
\end{equation}
and write \(U(k)=U(1,k)\).

\begin{lemma}[Local fixed-scale Kato trivialization]
\label{lem:v2v78-U}
The operator \(U(k)\) is unitary on \(\mathcal H\) and
\begin{equation}
 P(k)U(k)=U(k)P_0,
 \qquad
 Q(k)U(k)=U(k)Q_0,
\label{eq:v2v78-U-intertwine}
\end{equation}
where \(P_0=P(k_0)\) and \(Q_0=I-P_0\).
It and its inverse are uniformly bounded on every \(\mathcal H^s\) used
below.  Moreover, for \(1\le|\alpha|\le5\),
\begin{equation}
 \partial_k^\alpha U(k),\
 \partial_k^\alpha U(k)^*:
 \mathcal H^s\longrightarrow
 \mathcal H^{s+|\alpha|+1}
\label{eq:v2v78-U-order}
\end{equation}
are uniformly bounded after restriction to indices in \(J\).
\end{lemma}

\begin{proof}
The generator is skew-adjoint because \(P\) and
\(\partial_\rho P\) are self-adjoint.  The proof of
Theorem~\ref{thm:v2v34-transport}, applied on each radial path, gives
unitarity and \eqref{eq:v2v78-U-intertwine}.

Lemma~\ref{lem:v2v78-P-order} and the Leibniz rule show that every
\(k\)-derivative of \(\mathcal K\) of order \(a\) gains at least
\(a+2\) scale indices.  One derivative may hit the affine factor
\(k-k_0\) in \(\partial_\rho P(k_\rho)\); this is the case giving the
lower gain \(a+2\).  Differentiate
\eqref{eq:v2v78-U-ode} in \(k\).  Variation of constants expresses each
nonzero derivative of \(U\) as an integral of a differentiated
\(\mathcal K\), lower derivatives of \(U\), and undifferentiated
propagators.  Induction gives the weaker but sufficient gain
\(|\alpha|+1\) in \eqref{eq:v2v78-U-order}.  The undifferentiated
generator is bounded on every fixed \(\mathcal H^s\), so Gronwall's
inequality gives the scale bounds for \(U\) and \(U^{-1}=U^*\).
\end{proof}

\begin{remark}[Path dependence]
\label{rem:v2v78-path}
The radial choice fixes one smooth local trivialization.  A different
path can change \(U(k)\) by a unitary acting inside the fixed kernel and
complement.  The estimates below are invariant under any such
scale-bounded change of frame; no flatness of the Kato connection is
assumed.
\end{remark}

\subsection{The transported complement operator}
\label{sec:v2v78-transported}

Set
\[
 \mathcal X^s=Q_0\mathcal H^s
\]
with its closed-subspace norm.  Since \(Q_0\) is scale-bounded, these
spaces form the fixed complement scale needed by V71.  Define
\begin{align}
 \widehat D(k)
 &=
 Q_0U(k)^*D(k)U(k)Q_0:
 \mathcal X^{s+1}\longrightarrow\mathcal X^s,
 \label{eq:v2v78-Dhat}\\
 \widehat G(k)
 &=
 Q_0U(k)^*G(k)U(k)Q_0.
 \label{eq:v2v78-Ghat}
\end{align}

\begin{proposition}[Exact inverse in the fixed complement]
\label{prop:v2v78-fixed-inverse}
For every \(k\in B\),
\begin{equation}
 \widehat D(k)^{-1}=\widehat G(k)
\label{eq:v2v78-fixed-inverse}
\end{equation}
on \(\mathcal X^0\), and
\begin{equation}
 \widehat G(k):
 \mathcal X^s\longrightarrow\mathcal X^{s+1}
\label{eq:v2v78-Ghat-order}
\end{equation}
is uniformly bounded for \(-3\le s\le2\).
\end{proposition}

\begin{proof}
The intertwining identities give
\(UQ_0=QU\).  Using \(DG=GD=Q\),
\[
 \widehat D\widehat G
 =
 Q_0U^*D(UQ_0U^*)GUQ_0
 =
 Q_0U^*DQGUQ_0
 =
 Q_0U^*DGUQ_0
 =
 Q_0.
\]
The reverse product is identical.  The scale estimate follows from the
V77 reduced inverse bound and the scale boundedness of \(U\), \(U^*\),
and \(Q_0\).
\end{proof}

Along a one-parameter Kato path, let
\(K=[\dot P,P]\) and \(\dot U=KU\).  Direct differentiation gives
\begin{equation}
 \frac{\mathrm d}{\mathrm dt}(U^*DU)
 =
 U^*\bigl(\dot D+[D,K]\bigr)U.
\label{eq:v2v78-covariant-D}
\end{equation}
The connection term vanishes after complement compression:
\begin{equation}
 Q[D,K]Q=0.
\label{eq:v2v78-connection-cancel}
\end{equation}
Indeed, \(K\) is off-diagonal for \(P\oplus Q\), while
\(DP=PD=0\) and \(D=QDQ\).  Thus the first transported derivative is
exactly the compressed physical derivative.  Higher derivatives contain
connection terms, but Lemma~\ref{lem:v2v78-U} makes them strictly more
regular in the scale count.

\begin{theorem}[V71 survives exact-kernel reduction]
\label{thm:v2v78-V71}
Under \eqref{eq:v2v78-gap}--\eqref{eq:v2v78-D-order}, for every
multi-index \(\alpha\) with \(1\le|\alpha|\le5\),
\begin{equation}
 \partial_k^\alpha\widehat D(k):
 \mathcal X^s\longrightarrow
 \mathcal X^{s+|\alpha|-1}
\label{eq:v2v78-Dhat-order}
\end{equation}
is uniformly bounded whenever all relevant indices lie in \([-3,3]\).
Consequently, for \(0\le m\le5\),
\begin{equation}
 D_k^m\widehat G(k):
 \mathcal X^s\longrightarrow
 \mathcal X^{s+m+1}
\label{eq:v2v78-Ghat-derivative}
\end{equation}
is uniformly bounded whenever all intermediate indices lie in
\([-3,3]\).

If \(M_Q\ge I\) generates the fixed complement scale
\(\{\mathcal X^s\}\), then
\begin{equation}
 \left\|
 M_Q^{(m+1)/2}
 D_k^m\widehat G(k)
 M_Q^{(m+1)/2}
 \right\|
 \le C_{m,\delta},
 \qquad 0\le m\le5.
\label{eq:v2v78-symmetric}
\end{equation}
\end{theorem}

\begin{proof}
Differentiate \eqref{eq:v2v78-Dhat}.  The term on which all derivatives
hit \(D\) has the mapping property
\eqref{eq:v2v78-D-order}.  In every other Leibniz term, at least one
derivative hits \(U\) or \(U^*\).  If \(b\) derivatives hit \(D\), that
factor gains \(b-1\) scale indices; the remaining differentiated unitary
factors gain at least \(|\alpha|-b+1\) in total by
\eqref{eq:v2v78-U-order}.  Such a term therefore gains at least
\(|\alpha|\), which is stronger than the required
\(|\alpha|-1\).  Undifferentiated unitary factors are scale-bounded.
This proves \eqref{eq:v2v78-Dhat-order}.

Proposition~\ref{prop:v2v78-fixed-inverse} supplies
\eqref{eq:v2v71-R-order} on the fixed scale, and
\eqref{eq:v2v78-Dhat-order} supplies
\eqref{eq:v2v71-H-order}.  Apply
Theorem~\ref{thm:v2v71-scale-inverse}.  Its symmetric-weight conclusion,
with \(M_Q\), is exactly \eqref{eq:v2v78-symmetric}.
\end{proof}

\begin{assessment}[Infrared and ultraviolet constants separate]
\label{ass:v2v78-constants}
The V71 scale gain is unchanged by the kernel reduction.  Its constants
now also contain the gap dependence from V77 and the Kato-frame bounds.
Thus the ultraviolet derivative order closes uniformly only while the
nonzero Dirac gap remains uniform.  Kato transport organizes the moving
subspace; it does not improve the infrared constant.
\end{assessment}

\subsection{A dispersive zero crossing is not a persistent kernel}
\label{sec:v2v78-crossing}

\begin{counterexample}[Rank change at a massless dispersion point]
\label{sep:v2v78-rank-change}
On \(\mathbb C^2\), let
\[
 D(k)=
 \begin{pmatrix}
  0&0\\
  0&k
 \end{pmatrix},
 \qquad |k|<1.
\]
For \(k\ne0\), the exact kernel has rank one; at \(k=0\), it has rank two.
There is no norm-continuous projection onto the exact kernel across
\(k=0\), and the nonzero gap equals \(|k|\).  Therefore the persistent
kernel hypotheses of V77--V78 cannot be used through this crossing.

A valid local decomposition must instead place the second eigenvector in
a finite-dimensional crossing cluster throughout a neighbourhood of
zero.  The complement is then uniformly invertible, while the scalar
matrix entry \(k\) remains in the low crossing problem and is not
inverted uniformly.
\end{counterexample}

\begin{firewall}[The exact-kernel theorem is not the massless dispersion theorem]
\label{fw:v2v78-dispersion}
Protected constraint or topological zero modes may form a constant-rank
kernel bundle and are covered by V78.  Physical massless dispersion modes
usually cross zero and change the exact-kernel rank.  They require a
fixed-rank spectral cluster containing the crossing, followed by a
finite-dimensional low-energy analysis.
\end{firewall}

\begin{firewall}[Uniform scale regularity remains an input]
\label{fw:v2v78-scale}
Norm continuity of \(P(k)\) on \(\mathcal H\) proves neither
\eqref{eq:v2v78-P-order} nor the scale bounds for \(U(k)\).
Those conclusions use the differentiated contour-resolvent bounds and
the extended scale window.  They must be verified for the regulated
boundary problem with constants uniform in the regulator.
\end{firewall}

\subsection{V78 conclusion and next acquisition}
\label{sec:v2v78-conclusion}

V78 converts the reduced Dirac operator on the moving kernel complement
into an ordinary invertible family on a fixed Hilbert scale.  It proves
that the differentiated Riesz projection and the Kato frame are
sufficiently regular to preserve the V71 five-derivative symmetric
weighted estimate.  The conclusion applies to a persistent
constant-rank exact kernel with a uniform nonzero gap.

A dispersive massless crossing falls outside that hypothesis.  V79 will
replace the exact kernel by a fixed-rank Riesz cluster containing every
eigenvalue that can cross zero.  It will invert only the uniformly gapped
external complement and retain the Kato-compressed finite matrix as an
explicit low-energy block.

\clearpage
\section*{Second Volume, Version V79: A Crossing Cluster and a Uniform External Dirac Inverse}
\addcontentsline{toc}{section}{Second Volume, Version V79: A Crossing Cluster and a Uniform External Dirac Inverse}
\label{sec:v2v79}

\subsection{Enlarge the infrared sector before inverting}
\label{sec:v2v79-purpose}

The rank-change example in V78 identifies the correct upstream move.  One
must not project only onto the exact kernel at each momentum.  Instead,
select one fixed-rank spectral cluster containing every eigenvalue that
can reach zero.  The cluster is transported as a whole, including all of
its internal crossings.  Only the external complement is inverted
uniformly.

V38 already proved smooth Kato compression and finite-type occupation
bounds for an isolated cluster.  V79 adds the operator statement required
here: an exact block decomposition of the massless Dirac family, a
uniform V71 inverse on the external block, and a proof that every
remaining zero singularity lies in the finite Hermitian cluster matrix.

\subsection{Cluster hypotheses}
\label{sec:v2v79-hypotheses}

Let \(B\subset\mathbb R^q\) and the Hilbert scale be as in V78.  Let
\(D(k)\) be a \(C^5\) self-adjoint compact-resolvent family.  Assume there
is a closed contour \(\Gamma_{\rm cl}\) contained in the resolvent set of
\(D(k)\) for every \(k\in B\), enclosing a finite spectral cluster and no
other spectrum.  Define
\begin{equation}
 P_{\rm cl}(k)
 =
 \frac1{2\pi i}
 \int_{\Gamma_{\rm cl}}(z-D(k))^{-1}\,\mathrm dz,
 \qquad
 Q_{\rm ext}(k)=I-P_{\rm cl}(k).
\label{eq:v2v79-projections}
\end{equation}
Assume the cluster contains every spectral point in
\((-\delta,\delta)\):
\begin{equation}
 \mathbf1_{(-\delta,\delta)}(D(k))
 \le
 P_{\rm cl}(k)
\label{eq:v2v79-containment}
\end{equation}
in the order of orthogonal projections.  Equivalently,
\begin{equation}
 \operatorname{spec}\bigl(
 D(k)|_{\operatorname{Ran}Q_{\rm ext}(k)}
 \bigr)
 \cap(-\delta,\delta)=\varnothing.
\label{eq:v2v79-external-gap}
\end{equation}
Internal eigenvalues of the cluster may cross zero or each other.

Assume on the extended scale interval \(J\) the contour resolvent has
order minus one and the differentiated kinetic family satisfies
\eqref{eq:v2v78-D-order}.  The proof of
Lemma~\ref{lem:v2v78-P-order} then applies verbatim to
\(P_{\rm cl}\):
\begin{equation}
 \partial_k^\alpha P_{\rm cl}(k):
 \mathcal H^s\longrightarrow
 \mathcal H^{s+|\alpha|+1},
 \qquad
 0\le|\alpha|\le5.
\label{eq:v2v79-P-order}
\end{equation}

Let \(U(k)\) be the radial Kato transport of \(P_{\rm cl}(k)\) from
\(k_0\), and put
\[
 P_0=P_{\rm cl}(k_0),
 \qquad
 Q_0=I-P_0,
 \qquad
 E_0=\operatorname{Ran}P_0,
 \qquad
 X_0=\operatorname{Ran}Q_0.
\]
Then
\begin{equation}
 P_{\rm cl}(k)U(k)=U(k)P_0,
 \qquad
 Q_{\rm ext}(k)U(k)=U(k)Q_0.
\label{eq:v2v79-intertwine}
\end{equation}

\subsection{Exact block diagonalization}
\label{sec:v2v79-block}

Define
\begin{align}
 B_{\rm cl}(k)
 &=
 P_0U(k)^*D(k)U(k)P_0|_{E_0},
 \label{eq:v2v79-B}\\
 H_{\rm ext}(k)
 &=
 Q_0U(k)^*D(k)U(k)Q_0|_{X_0}.
 \label{eq:v2v79-H}
\end{align}

\begin{theorem}[Exact crossing-cluster decomposition]
\label{thm:v2v79-block}
The space \(E_0\) has finite dimension independent of \(k\),
\(B_{\rm cl}(k)\) is a \(C^5\) Hermitian matrix family, and
\(H_{\rm ext}(k)\) is self-adjoint on the fixed complement.  Moreover,
\begin{equation}
 U(k)^*D(k)U(k)
 =
 B_{\rm cl}(k)\oplus H_{\rm ext}(k)
\label{eq:v2v79-direct-sum}
\end{equation}
on \(E_0\oplus X_0\), with no off-diagonal remainder.
\end{theorem}

\begin{proof}
A Riesz projection for a self-adjoint operator is orthogonal and commutes
with that operator.  Hence
\[
 P_{\rm cl}D Q_{\rm ext}
 =
 Q_{\rm ext}D P_{\rm cl}=0.
\]
Conjugate these identities by the intertwining unitary \(U(k)\) to obtain
\[
 P_0U^*DUQ_0=Q_0U^*DUP_0=0,
\]
which proves \eqref{eq:v2v79-direct-sum}.  Finite rank and constant rank
follow from compact resolvent, the fixed contour, and norm continuity of
the Riesz projection.  Hermiticity and self-adjointness are preserved by
unitary conjugation and reducing compression.  Differentiated contour
regularity and the Kato ordinary differential equation give \(C^5\)
dependence, as in Lemma~\ref{lem:v2v38-cluster-regularity}.
\end{proof}

\begin{assessment}[No Schur complement is needed]
\label{ass:v2v79-no-Schur}
Because \(P_{\rm cl}(k)\) is the exact spectral projection, the two
off-diagonal blocks vanish identically.  Introducing a Schur complement
here would add an artificial error term.  Schur estimates become relevant
only for approximate, regulator-mismatched, or spatially localized
projectors.
\end{assessment}

\subsection{The uniform external inverse}
\label{sec:v2v79-external}

Define by spectral calculus
\begin{equation}
 G_{\rm ext}(k)
 =
 Q_{\rm ext}(k)
 \bigl(D(k)|_{\operatorname{Ran}Q_{\rm ext}(k)}\bigr)^{-1}
 Q_{\rm ext}(k).
\label{eq:v2v79-Gext}
\end{equation}
This operator is zero on the cluster and is the ordinary inverse on its
orthogonal complement.  Put
\begin{equation}
 \widehat G_{\rm ext}(k)
 =
 Q_0U(k)^*G_{\rm ext}(k)U(k)Q_0.
\label{eq:v2v79-Ghat}
\end{equation}

\begin{theorem}[Uniform external inverse and V71 derivatives]
\label{thm:v2v79-external}
One has
\begin{align}
 H_{\rm ext}(k)^{-1}
 &=\widehat G_{\rm ext}(k),
 \label{eq:v2v79-H-inverse}\\
 \|\widehat G_{\rm ext}(k)\|_{\mathcal H\to\mathcal H}
 &\le\delta^{-1}.
 \label{eq:v2v79-H-L2}
\end{align}
Uniform elliptic estimates give
\begin{equation}
 \widehat G_{\rm ext}(k):
 Q_0\mathcal H^s\longrightarrow Q_0\mathcal H^{s+1}
\label{eq:v2v79-H-scale}
\end{equation}
for \(-3\le s\le2\).  Under the V78 extended-scale hypotheses,
\begin{equation}
 \partial_k^\alpha H_{\rm ext}(k):
 Q_0\mathcal H^s\longrightarrow
 Q_0\mathcal H^{s+|\alpha|-1},
 \qquad
 1\le|\alpha|\le5.
\label{eq:v2v79-H-derivative}
\end{equation}
Therefore
\begin{equation}
 D_k^m\widehat G_{\rm ext}(k):
 Q_0\mathcal H^s\longrightarrow
 Q_0\mathcal H^{s+m+1},
 \qquad
 0\le m\le5,
\label{eq:v2v79-G-derivative}
\end{equation}
whenever all intermediate indices lie in \([-3,3]\).

If \(M_{\rm ext}\) generates the fixed complement scale, then
\begin{equation}
 \left\|
 M_{\rm ext}^{(m+1)/2}
 D_k^m\widehat G_{\rm ext}(k)
 M_{\rm ext}^{(m+1)/2}
 \right\|
 \le C_{m,\delta},
 \qquad
 0\le m\le5.
\label{eq:v2v79-symmetric}
\end{equation}
All these estimates remain uniform when \(B_{\rm cl}(k)\) is singular.
\end{theorem}

\begin{proof}
The external gap
\eqref{eq:v2v79-external-gap} and the spectral theorem give
\(\|G_{\rm ext}\|\le\delta^{-1}\).  Conjugation by \(U(k)\) proves
\eqref{eq:v2v79-H-inverse} and
\eqref{eq:v2v79-H-L2}.  The elliptic proof of
Theorem~\ref{thm:v2v77-reduced} applies to the external reducing
subspace and proves \eqref{eq:v2v79-H-scale}.

Equation \eqref{eq:v2v79-P-order} gives the same scale estimates for the
cluster Kato unitary as
Lemma~\ref{lem:v2v78-U}.  Differentiating
\eqref{eq:v2v79-H} and repeating the scale count in
Theorem~\ref{thm:v2v78-V71} proves
\eqref{eq:v2v79-H-derivative}.  The V71 differentiated inverse theorem
then gives \eqref{eq:v2v79-G-derivative} and
\eqref{eq:v2v79-symmetric}.  None of these steps inverts
\(B_{\rm cl}(k)\), so its internal zero crossings do not enter the
constants.
\end{proof}

\subsection{All zero singularities are finite-dimensional}
\label{sec:v2v79-singularities}

Let
\begin{equation}
 d(k)=\det B_{\rm cl}(k).
\label{eq:v2v79-det}
\end{equation}

\begin{theorem}[Localization of the massless singularity]
\label{thm:v2v79-localization}
The following statements hold.
\begin{enumerate}
 \item \(0\in\operatorname{spec}D(k)\) if and only if \(d(k)=0\).
 \item If \(d(k)\ne0\), then
 \begin{equation}
  D(k)^{-1}
  =
  U(k)\bigl(
  B_{\rm cl}(k)^{-1}\oplus H_{\rm ext}(k)^{-1}
  \bigr)U(k)^*.
  \label{eq:v2v79-full-inverse}
 \end{equation}
 \item On \(\mathcal H\),
 \begin{equation}
  \|D(k)^{-1}\|
  =
  \max\left\{
  \|B_{\rm cl}(k)^{-1}\|,
  \|H_{\rm ext}(k)^{-1}\|
  \right\}.
  \label{eq:v2v79-inverse-max}
 \end{equation}
 Hence every loss of uniform invertibility at zero is carried by the
 finite matrix \(B_{\rm cl}(k)\).
\end{enumerate}
\end{theorem}

\begin{proof}
By \eqref{eq:v2v79-direct-sum}, the spectrum of \(D(k)\) is the union of
the spectra of its two blocks.  The external block has no spectrum in
\((-\delta,\delta)\), so it cannot contain zero.  The cluster block
contains zero exactly when its determinant vanishes.  If both blocks are
invertible, inversion of the orthogonal direct sum proves
\eqref{eq:v2v79-full-inverse}; the norm of an orthogonal direct sum is the
maximum of the block norms, proving
\eqref{eq:v2v79-inverse-max}.
\end{proof}

\begin{corollary}[Finite type controls the crossing set]
\label{cor:v2v79-finite-type}
Along a \(C^5\) one-parameter curve \(k=k(t)\), suppose the cluster
eigenvalues remain bounded by \(R\), and for some \(1\le r\le5\),
\begin{equation}
 \left|
 \frac{\mathrm d^r}{\mathrm dt^r}
 \det B_{\rm cl}(k(t))
 \right|
 \ge c>0.
\label{eq:v2v79-finite-type}
\end{equation}
Then the set on which a cluster eigenvalue has absolute value below
\(g\) has length at most
\begin{equation}
 2r\left(\frac{gR^{q-1}}c\right)^{1/r},
 \qquad q=\dim E_0.
\label{eq:v2v79-crossing-measure}
\end{equation}
The external inverse estimates of
Theorem~\ref{thm:v2v79-external} hold on the entire curve, including this
crossing set.
\end{corollary}

\begin{proof}
Apply Theorem~\ref{thm:v2v38-cluster} at threshold \(a=0\).  That theorem
uses no internal eigenvalue gap.  The external statement has already been
proved without inverting the cluster.
\end{proof}

\begin{firewall}[Finite type is not automatic]
\label{fw:v2v79-finite-type}
Smoothness of \(B_{\rm cl}\) does not imply
\eqref{eq:v2v79-finite-type}.  A cluster eigenvalue can vanish to infinite
order or remain identically zero.  The determinant-jet bound must come
from the reduced dynamics, a symmetry classification, or a separately
verified transversality hypothesis.
\end{firewall}

\subsection{Minimal separator}
\label{sec:v2v79-separator}

\begin{counterexample}[Omitting the crossing mode forces inverse blow-up]
\label{sep:v2v79-omit}
Let
\[
 D(k)=
 \begin{pmatrix}
  k&0\\0&1
 \end{pmatrix}.
\]
If the crossing mode is left in the sector to be inverted, the inverse
norm is at least \(|k|^{-1}\).  If
\(P_{\rm cl}=|e_1\rangle\langle e_1|\) in the displayed coordinate
ordering, then
\[
 B_{\rm cl}(k)=(k),
 \qquad
 H_{\rm ext}(k)=(1),
\]
and the external inverse has norm one through the crossing.  The
singularity has not been discarded; it has been isolated in the exact
one-dimensional low block.
\end{counterexample}

\begin{firewall}[Exact and approximate projectors]
\label{fw:v2v79-approximate}
The zero off-diagonal blocks in
\eqref{eq:v2v79-direct-sum} use the exact Riesz projection of the same
operator \(D(k)\).  A projector imported from a limiting operator or a
different regulator generally leaves off-diagonal errors.  Those errors
must satisfy a relative Schur or Neumann estimate before the conclusions
of V79 can be transferred.
\end{firewall}

\subsection{V79 conclusion and next acquisition}
\label{sec:v2v79-conclusion}

V79 gives the correct operator decomposition at a massless dispersion
point.  A fixed-rank Riesz cluster contains all modes that can reach zero.
Its external complement has a uniform scale inverse and the complete V71
five-derivative symmetric bounds.  Every remaining zero and inverse
blow-up is localized in the finite Hermitian matrix
\(B_{\rm cl}(k)\), where the V38 determinant finite-type theorem applies.

V80 is the next mandatory companion audit.  It will inspect the current
Yang--Mills and Navier--Stokes notes before deciding whether the next SM
step should estimate the finite crossing block, repair an
approximate-projector transfer, or move farther upstream to the
regulator-uniform spectral construction.

\clearpage
\section*{Second Volume, Version V80: Companion Audit and the Finite-Cluster Pole-Divisibility Gate}
\addcontentsline{toc}{section}{Second Volume, Version V80: Companion Audit and the Finite-Cluster Pole-Divisibility Gate}
\label{sec:v2v80}

\subsection{Mandatory companion audit}
\label{sec:v2v80-audit}

Before extending V79, the newest active companion notes in the shared
folder were read without modification:
\begin{center}
\small
\begin{tabular}{p{0.48\linewidth}r p{0.32\linewidth}}
\toprule
file&bytes&SHA--256\\
\midrule
\texttt{Renewal\_Yang\_Mills\_Mass\_Gap\_Research\_Notes\_V64.tex}
&896171&
\texttt{568479b7f8f22755276ad53fc38fb62a13ae4869de39dc4ad63a06018a7baad5}
\\
\texttt{Renewal\_NS\_Stability\_Research\_Notes\_V31.tex}
&887537&
\texttt{f1121b62238ad5491bb7cf03635ea9934942edf573ed8faaa9ee79e1d38a6696}
\\
\bottomrule
\end{tabular}
\end{center}
The NS file is unchanged from the V75 audit.  The frozen first Yang--Mills
volume remains a separate archive and was not treated as the active
continuation.

\begin{assessment}[Yang--Mills transfer]
\label{ass:v2v80-YM}
Yang--Mills V64 closes seven fully thermal orbit types by separated
estimates, but proves sharp one-power regressions for five remaining
branches.  Those branches lie in three assembled quotient-cluster
cumulants; the first target is the \(2+2\) composite covariance, which
must be pinched before its four pair-partition letters are separated.

For the SM crossing block, the transferable proof rule is to keep the
complete finite-cluster numerator assembled before applying an operator
norm.  The Yang--Mills scalar powers and Wilson-clock estimates are not
imported.
\end{assessment}

\begin{assessment}[Navier--Stokes transfer]
\label{ass:v2v80-NS}
Navier--Stokes V31 obtains an exact dyadic probability mark and a
multiplicity-free formation ladder, but removal of the mark is precisely
the missing critical first moment.  Progress requires correlation with
the actual recent-entry set and signed nonlinear transfer.

For V79 this warns that a sharp sublevel estimate for
\(\det B_{\rm cl}\) need not control any inverse moment.  The finite-type
crossing theorem records where the singularity lies; a numerator
cancellation or a physical distributional prescription is still needed.
No Navier--Stokes inequality is imported.
\end{assessment}

The audit therefore selects the algebraic numerator gate rather than
another occupation estimate.

\subsection{Finite type alone does not control the inverse}
\label{sec:v2v80-insufficiency}

\begin{proposition}[Sharp finite type can coexist with a divergent inverse]
\label{prop:v2v80-finite-type-insufficient}
Let \(B(t)=(t)\) on \((-1,1)\).  Then
\[
 \det B(t)=t,
 \qquad
 \left|\{t:|\det B(t)|<g\}\right|=2g
 \quad(0<g<1),
\]
so the strongest simple finite-type sublevel estimate holds.  Nevertheless
\begin{equation}
 \int_{-1}^1\|B(t)^{-1}\|\,\mathrm dt
 =
 \int_{-1}^1\frac{\mathrm dt}{|t|}
 =
 \infty.
\label{eq:v2v80-scalar-divergence}
\end{equation}
The same example occurs inside any finite cluster by taking
\(B(t)=\operatorname{diag}(t,1,\ldots,1)\).
\end{proposition}

\begin{proof}
The sublevel set is \((-g,g)\).  The inverse norm is \(|t|^{-1}\), whose
positive integral diverges logarithmically.
\end{proof}

\begin{firewall}[Occupation is not cancellation]
\label{fw:v2v80-occupation}
Corollary~\ref{cor:v2v79-finite-type} can make the small-eigenvalue set
quantitatively short, but a simple pole already saturates the borderline
between a short set and a divergent absolute first moment.  No refinement
of the same scalar sublevel estimate removes
\eqref{eq:v2v80-scalar-divergence}.
\end{firewall}

\subsection{The adjugate numerator}
\label{sec:v2v80-adjugate}

Let \(E\) be a finite-dimensional Hermitian space of dimension \(q\), and
let \(X,Y\) be Hilbert spaces.  On an interval about \(t_0\), consider
\[
 B(t)\in\operatorname{Herm}(E),
 \qquad
 R(t)\in\mathcal L(X,E),
 \qquad
 L(t)\in\mathcal L(E,Y).
\]
Where \(B(t)\) is invertible, define the sandwiched cluster amplitude
\begin{equation}
 \mathcal A(t)=L(t)B(t)^{-1}R(t).
\label{eq:v2v80-amplitude}
\end{equation}
Set
\begin{align}
 d(t)&=\det B(t),
 \label{eq:v2v80-d}\\
 N(t)&=L(t)\operatorname{adj}(B(t))R(t).
 \label{eq:v2v80-N}
\end{align}
Then
\begin{equation}
 \mathcal A(t)=\frac{N(t)}{d(t)}
\label{eq:v2v80-quotient}
\end{equation}
away from the zero set of \(d\).  The numerator \(N\) is formed only
after all terms contributing to \(L\), \(B\), and \(R\) have been
assembled.

\begin{theorem}[Exact pole-divisibility criterion]
\label{thm:v2v80-divisibility}
Suppose \(B,L,R\) are \(C^r\), and
\begin{equation}
 d(t)=(t-t_0)^r e(t),
 \qquad
 |e(t)|\ge e_0>0
\label{eq:v2v80-det-order}
\end{equation}
near \(t_0\).  If
\begin{equation}
 N^{(j)}(t_0)=0,
 \qquad
 0\le j\le r-1,
\label{eq:v2v80-jet-cancellation}
\end{equation}
as identities in \(\mathcal L(X,Y)\), then
\(\mathcal A(t)\) has a continuous extension through \(t_0\).  More
precisely,
\begin{equation}
 \sup_{|t-t_0|<\varepsilon}\|\mathcal A(t)\|
 \le
 \frac1{r!e_0}
 \sup_{|t-t_0|<\varepsilon}\|N^{(r)}(t)\|.
\label{eq:v2v80-extension-bound}
\end{equation}

Conversely, suppose the first nonzero numerator jet has order
\(\ell<r\):
\begin{equation}
 N^{(j)}(t_0)=0\ (j<\ell),
 \qquad
 N^{(\ell)}(t_0)\ne0.
\label{eq:v2v80-first-jet}
\end{equation}
Then, after the interval is reduced,
\begin{equation}
 \|\mathcal A(t)\|
 \ge c|t-t_0|^{\ell-r}
\label{eq:v2v80-lower-pole}
\end{equation}
on each punctured side.  In particular,
\begin{equation}
 \int_{t_0-\varepsilon}^{t_0+\varepsilon}
 \|\mathcal A(t)\|\,\mathrm dt=\infty.
\label{eq:v2v80-L1-divergence}
\end{equation}
\end{theorem}

\begin{proof}
Taylor's formula with integral remainder and
\eqref{eq:v2v80-jet-cancellation} give
\[
 N(t)
 =
 \frac{(t-t_0)^r}{(r-1)!}
 \int_0^1
 (1-\theta)^{r-1}
 N^{(r)}(t_0+\theta(t-t_0))\,\mathrm d\theta.
\]
Thus \(N(t)=(t-t_0)^r\widetilde N(t)\), where
\(\widetilde N\) is continuous and
\[
 \|\widetilde N(t)\|
 \le
 \frac1{r!}
 \sup_{|s-t_0|<\varepsilon}\|N^{(r)}(s)\|.
\]
Cancel the common factor in
\eqref{eq:v2v80-quotient} and use
\(|e|\ge e_0\) to obtain the continuous extension and
\eqref{eq:v2v80-extension-bound}.

Under \eqref{eq:v2v80-first-jet},
\[
 N(t)
 =
 \frac{(t-t_0)^\ell}{\ell!}N^{(\ell)}(t_0)
 +o(|t-t_0|^\ell)
\]
in operator norm.  Hence
\(\|N(t)\|\ge c_1|t-t_0|^\ell\) near \(t_0\).
Continuity of \(e\) also gives
\(|d(t)|\le c_2|t-t_0|^r\).  This proves
\eqref{eq:v2v80-lower-pole}.  Since \(r-\ell\) is a positive integer, the
lower bound is at least a simple pole, whose absolute integral diverges.
\end{proof}

\begin{assessment}[The required Ward order is exact]
\label{ass:v2v80-Ward-order}
If the determinant has a zero of order \(r\), cancellation of fewer than
\(r\) numerator jets cannot yield an absolutely integrable operator norm.
Cancellation through order \(r-1\) yields a bounded extension.  There is
no missing intermediate integer order for an \(L^1\) estimate.
\end{assessment}

\subsection{Simple crossing and the physical matrix element}
\label{sec:v2v80-simple}

\begin{corollary}[Simple eigenvalue crossing]
\label{cor:v2v80-simple}
Suppose \(B(t_0)\) has a one-dimensional kernel with orthogonal
projection \(\Pi_0\), all other eigenvalues are nonzero, and the crossing
eigenvalue is simple:
\[
 \lambda(t_0)=0,
 \qquad
 \lambda'(t_0)\ne0.
\]
Then \(d\) has a simple zero and
\begin{equation}
 \operatorname{adj}(B(t_0))
 =
 c_0\Pi_0,
 \qquad
 c_0=\prod_{\mu\ne0}\mu\ne0,
\label{eq:v2v80-adj-simple}
\end{equation}
where the product is over the nonzero eigenvalues of \(B(t_0)\), with
multiplicity.  Consequently
\begin{equation}
 N(t_0)=0
 \quad\Longleftrightarrow\quad
 L(t_0)\Pi_0R(t_0)=0.
\label{eq:v2v80-simple-Ward}
\end{equation}
If this matrix element vanishes, \(\mathcal A\) extends continuously.  If
it does not vanish, \(\mathcal A\) has a nonzero simple pole and its
operator norm is not locally integrable.
\end{corollary}

\begin{proof}
Diagonalize the Hermitian matrix \(B(t_0)\).  The adjugate is zero on
every nonzero eigendirection and equals the product of the nonzero
eigenvalues on the kernel direction, proving
\eqref{eq:v2v80-adj-simple}.  A transverse simple eigenvalue gives a
simple determinant zero.  Apply
Theorem~\ref{thm:v2v80-divisibility} with \(r=1\).
\end{proof}

\begin{remark}[Higher corank]
\label{rem:v2v80-higher-corank}
If \(\dim\ker B(t_0)\ge2\), then
\(\operatorname{adj}(B(t_0))=0\) automatically.  This does not prove
cancellation: the determinant also vanishes to higher order, and the
higher numerator jets in \eqref{eq:v2v80-jet-cancellation} become the
actual conditions.  The adjugate formulation retains exactly those
conditions without choosing individual eigenvalue branches.
\end{remark}

\subsection{Why the complete Ward sum must remain assembled}
\label{sec:v2v80-assembled}

Suppose the low-cluster numerator has a finite decomposition
\begin{equation}
 N(t)=\sum_{\nu=1}^J N_\nu(t)
\label{eq:v2v80-sum}
\end{equation}
coming from diagrams, BRST words, or boundary sectors.

\begin{proposition}[Cancellation may exist only after assembly]
\label{prop:v2v80-assembly}
If
\begin{equation}
 \sum_{\nu=1}^J N_\nu^{(j)}(t_0)=0,
 \qquad 0\le j\le r-1,
\label{eq:v2v80-summed-jets}
\end{equation}
then the total amplitude \(N/d\) has the bounded extension of
Theorem~\ref{thm:v2v80-divisibility}.  No corresponding bound follows
for
\begin{equation}
 \sum_{\nu=1}^J\frac{\|N_\nu(t)\|}{|d(t)|}
\label{eq:v2v80-separated-majorant}
\end{equation}
unless each numerator separately has the full \(r\)-jet divisibility.
\end{proposition}

\begin{proof}
Equation \eqref{eq:v2v80-summed-jets} is precisely
\eqref{eq:v2v80-jet-cancellation} for the sum.  The second statement is
sharp already for
\[
 d(t)=t,
 \qquad
 N_1(t)=I,
 \qquad
 N_2(t)=-I+tC.
\]
The total quotient equals \(C\), whereas the separated positive majorant
is at least \(2/|t|-O(1)\).
\end{proof}

\begin{assessment}[Companion-guided proof order]
\label{ass:v2v80-proof-order}
The exact analogue of the YM V64 instruction is now concrete: form the
complete BRST/Ward cluster numerator, prove its adjugate-jet identities,
and only then take the operator norm.  The NS V31 lesson rules out
replacing these identities by a crossing occupation estimate of the
same critical inverse-moment size.
\end{assessment}

\subsection{Strong factorization and exact sectors}
\label{sec:v2v80-factorization}

\begin{proposition}[Kinetic factorization removes the cluster inverse]
\label{prop:v2v80-factorization}
Suppose, near \(t_0\), either
\begin{equation}
 R(t)=B(t)\widetilde R(t)
\label{eq:v2v80-right-factor}
\end{equation}
or
\begin{equation}
 L(t)=\widetilde L(t)B(t)
\label{eq:v2v80-left-factor}
\end{equation}
with continuous \(\widetilde R\) or \(\widetilde L\).  Then
\(\mathcal A(t)\) extends continuously through every singular point of
\(B(t)\).  Away from the singular set its value is respectively
\[
 L(t)\widetilde R(t)
 \quad\hbox{or}\quad
 \widetilde L(t)R(t).
\]
\end{proposition}

\begin{proof}
Cancel \(B^{-1}B\) or \(BB^{-1}\) where \(B\) is invertible.  The
remaining expression is continuous and therefore defines a continuous
extension of the amplitude from the invertible locus.
\end{proof}

This is the finite-dimensional form of the desired exact-sector
mechanism: a BRST-exact or longitudinal source should carry a kinetic
factor before the cluster inverse is estimated.  V66--V68 supply the
relative-complex and Hodge framework in which such a factorization must
be proved for the complete low packet; V80 does not assume that the
physical SM numerator already has it.

\begin{firewall}[Physical massless poles are not Ward errors]
\label{fw:v2v80-physical-pole}
A transverse photon or another physical massless excitation can have a
nonzero matrix element in
\eqref{eq:v2v80-simple-Ward}.  Its pole is physical and should not be
cancelled by a BRST identity.  Such a sector requires a Lorentzian
boundary value, causal distribution, and inclusive infrared observable,
rather than a uniform Euclidean inverse norm.  The divisibility criterion
applies to sectors whose Ward or constraint structure predicts
decoupling.
\end{firewall}

\begin{firewall}[Principal value is a separate theorem]
\label{fw:v2v80-PV}
The signed principal value of \(1/(t-t_0)\) can exist although its
absolute integral diverges.  A principal-value, retarded, or Feynman
prescription cannot be substituted into the positive operator-norm
estimates of V70--V79 without a separate distributional composition
theorem.
\end{firewall}

\subsection{V80 conclusion and next acquisition}
\label{sec:v2v80-conclusion}

The mandatory audit has been completed.  It redirects the massless
crossing problem from sublevel occupation to assembled numerator
algebra.  V80 proves that finite type of the cluster determinant alone
does not control the inverse, and gives a necessary-and-sufficient jet
order for bounded cancellation of a sandwiched finite-cluster pole.
For a simple crossing this is exactly the vanishing of the source-to-
observable matrix element through the crossing eigenspace.

V81 will insert the relative BRST/Hodge decomposition of V66--V68 into
the V79 cluster block.  Its target is an exact factorization for the
longitudinal and BRST-exact cluster sources, while retaining the physical
harmonic massless block for a later causal infrared construction.

\clearpage
\section*{Second Volume, Version V81: BRST--Hodge Factorization of the Crossing-Cluster Numerator}
\addcontentsline{toc}{section}{Second Volume, Version V81: BRST--Hodge Factorization of the Crossing-Cluster Numerator}
\label{sec:v2v81}

\subsection{Purpose}
\label{sec:v2v81-purpose}

V80 reduced bounded pole cancellation to adjugate-numerator jets.  This
version proves those identities for the part of a cluster source which is
BRST exact.  The proof is finite-dimensional and occurs after the V79
Kato transport, so the cluster space and its Hilbert BRST differential
are fixed.

The conclusion is deliberately sectorial.  Exact and longitudinal
sources decouple from physical observables before the cluster determinant
is divided out.  A harmonic source is not removed and may carry a
physical massless pole.

\subsection{A fixed Hilbert BRST cluster}
\label{sec:v2v81-complex}

Let
\[
 E^\bullet=\bigoplus_gE^g
\]
be a finite-dimensional graded Hilbert space with differential
\[
 s:E^g\longrightarrow E^{g+1},
 \qquad s^2=0.
\]
Let
\begin{equation}
 \Delta_s=ss^*+s^*s,
 \qquad
 \Pi=\mathbf1_{\{0\}}(\Delta_s)
\label{eq:v2v81-Hodge}
\end{equation}
and assume the quantitative Hodge gap
\begin{equation}
 \Delta_s\ge\gamma^2(I-\Pi)
\label{eq:v2v81-Hodge-gap}
\end{equation}
for one \(\gamma>0\).  As in V68, define
\begin{equation}
 G_s=\Delta_s^{-1}(I-\Pi),
 \qquad
 h=s^*G_s.
\label{eq:v2v81-h}
\end{equation}
Then
\begin{equation}
 sh+hs=I-\Pi,
 \qquad
 \|h\|\le\gamma^{-1}.
\label{eq:v2v81-contraction}
\end{equation}

Let \(B(t)\) be a \(C^r\) family of degree-zero Hermitian operators on
\(E^\bullet\).  The decisive compatibility hypothesis is
\begin{equation}
 [B(t),s]=[B(t),s^*]=0.
\label{eq:v2v81-BRST-invariance}
\end{equation}
It follows that \(B(t)\) commutes with
\(\Delta_s,\Pi,G_s\), and \(h\).

Let \(X,Y\) be Hilbert spaces.  A source
\(R(t)\in\mathcal L(X,E^g)\) is BRST closed when
\begin{equation}
 sR(t)=0.
\label{eq:v2v81-source-closed}
\end{equation}
A physical output
\(L(t)\in\mathcal L(E^g,Y)\) annihilates BRST boundaries when
\begin{equation}
 L(t)s=0
 \quad\hbox{on }E^{g-1}.
\label{eq:v2v81-output-closed}
\end{equation}

\begin{lemma}[The adjugate is a BRST chain map]
\label{lem:v2v81-adj-chain}
Under \eqref{eq:v2v81-BRST-invariance},
\begin{equation}
 [\operatorname{adj}(B(t)),s]
 =
 [\operatorname{adj}(B(t)),s^*]
 =0.
\label{eq:v2v81-adj-chain}
\end{equation}
The same operator also commutes with \(\Pi\) and \(h\).
\end{lemma}

\begin{proof}
In finite dimension the adjugate is a polynomial in \(B\), with scalar
coefficients given by the characteristic polynomial.  Hence every
operator commuting with \(B\) commutes with \(\operatorname{adj}(B)\).
Apply this first to \(s,s^*\), and then to the operators constructed from
them by spectral calculus.
\end{proof}

\subsection{Exact cancellation before division}
\label{sec:v2v81-exact}

\begin{theorem}[BRST-exact cluster sources have zero adjugate numerator]
\label{thm:v2v81-exact}
Suppose
\eqref{eq:v2v81-BRST-invariance} and
\eqref{eq:v2v81-output-closed} hold.  If
\begin{equation}
 R_{\rm ex}(t)=sC(t)
\label{eq:v2v81-exact-source}
\end{equation}
for a \(C^r\) family
\(C(t)\in\mathcal L(X,E^{g-1})\), then
\begin{equation}
 L(t)\operatorname{adj}(B(t))R_{\rm ex}(t)=0
\label{eq:v2v81-exact-numerator}
\end{equation}
for every \(t\).  Consequently all parameter jets of this complete
numerator vanish.  Wherever \(B(t)\) is invertible,
\begin{equation}
 L(t)B(t)^{-1}R_{\rm ex}(t)=0,
\label{eq:v2v81-exact-amplitude}
\end{equation}
and zero is a continuous extension through every singular point.
\end{theorem}

\begin{proof}
Lemma~\ref{lem:v2v81-adj-chain} gives
\[
 L\operatorname{adj}(B)sC
 =
 Ls\operatorname{adj}(B)C
 =
 0.
\]
This is the pointwise identity
\eqref{eq:v2v81-exact-numerator}; differentiating the zero function gives
all jet identities.  On the invertible locus,
\[
 B^{-1}s=sB^{-1}
\]
by \eqref{eq:v2v81-BRST-invariance}, so the same calculation proves
\eqref{eq:v2v81-exact-amplitude}.  The identically zero function supplies
the extension.
\end{proof}

\begin{corollary}[Every determinant order is cancelled]
\label{cor:v2v81-all-orders}
Let \(t_0\) be a cluster crossing at which
\[
 \det B(t)=(t-t_0)^m e(t),
 \qquad e(t_0)\ne0.
\]
The exact-sector numerator in
\eqref{eq:v2v81-exact-numerator} satisfies all \(m\) jet conditions of
Theorem~\ref{thm:v2v80-divisibility}, independently of the multiplicity
or tangency order of the crossing.
\end{corollary}

\begin{proof}
The numerator vanishes identically, so every derivative vanishes.
\end{proof}

This is stronger than estimating an inverse: the singular quotient is
never formed in the exact sector.

\subsection{Closed sources reduce to their harmonic component}
\label{sec:v2v81-harmonic}

\begin{theorem}[Hodge reduction of the complete cluster numerator]
\label{thm:v2v81-harmonic}
Assume
\eqref{eq:v2v81-Hodge-gap}--\eqref{eq:v2v81-output-closed} and let \(R\)
be BRST closed.  Then
\begin{align}
 R
 &=
 \Pi R+s\,hR,
 \label{eq:v2v81-source-split}\\
 L\operatorname{adj}(B)R
 &=
 L\operatorname{adj}(B)\Pi R.
 \label{eq:v2v81-numerator-reduction}
\end{align}
On the invertible locus,
\begin{equation}
 LB^{-1}R
 =
 LB^{-1}\Pi R.
\label{eq:v2v81-amplitude-reduction}
\end{equation}
Moreover,
\begin{equation}
 \|hR\|\le\gamma^{-1}\|R\|.
\label{eq:v2v81-primitive-bound}
\end{equation}
\end{theorem}

\begin{proof}
Apply the contracting identity
\eqref{eq:v2v81-contraction} to the range of \(R\):
\[
 R=\Pi R+shR+hsR=\Pi R+shR,
\]
because \(sR=0\).  Insert this split into the numerator.  The exact term
vanishes by
Theorem~\ref{thm:v2v81-exact}, proving
\eqref{eq:v2v81-numerator-reduction}.  The same proof with \(B^{-1}\)
gives \eqref{eq:v2v81-amplitude-reduction} where the inverse exists.
Finally use the V68 bound \(\|h\|\le\gamma^{-1}\).
\end{proof}

\begin{corollary}[Purely longitudinal cluster sources are harmless]
\label{cor:v2v81-longitudinal}
Under the hypotheses of
Theorem~\ref{thm:v2v81-harmonic}, if
\begin{equation}
 \Pi R(t)=0,
\label{eq:v2v81-no-harmonic-source}
\end{equation}
then the complete physical cluster numerator is identically zero and the
sandwiched amplitude extends as zero through every crossing.
\end{corollary}

\begin{proof}
Equation \eqref{eq:v2v81-source-split} makes \(R\) exact, and
Theorem~\ref{thm:v2v81-exact} applies.
\end{proof}

\begin{assessment}[The only surviving pole is cohomological]
\label{ass:v2v81-surviving}
For a BRST-closed source and a physical output, every exact component is
removed before taking norms.  The pole-divisibility problem of V80
therefore descends to the finite harmonic space
\(\Pi E^g\).  A nonzero pole there represents a physical cohomology
channel or an uncancelled anomaly; it is not a longitudinal inverse
failure.
\end{assessment}

\subsection{Block form of the invariant kinetic matrix}
\label{sec:v2v81-block}

The finite Hodge decomposition is
\begin{equation}
 E^g
 =
 \operatorname{im}s
 \oplus\Pi E^g
 \oplus\operatorname{im}s^*.
\label{eq:v2v81-Hodge-decomposition}
\end{equation}

\begin{proposition}[Hodge reducing decomposition]
\label{prop:v2v81-Hodge-block}
Under \eqref{eq:v2v81-BRST-invariance}, all three summands in
\eqref{eq:v2v81-Hodge-decomposition} reduce \(B(t)\).  Hence
\begin{equation}
 B(t)
 =
 B_{\rm ex}(t)
 \oplus B_{\rm har}(t)
 \oplus B_{\rm coex}(t).
\label{eq:v2v81-B-blocks}
\end{equation}
For physical amplitudes with closed sources, only
\(B_{\rm har}\) can contribute a nonzero singular cluster numerator.
\end{proposition}

\begin{proof}
If \(v=su\), then \(Bv=Bsu=sBu\) lies in
\(\operatorname{im}s\).  The same argument with \(s^*\) preserves
\(\operatorname{im}s^*\).  Commutation with \(\Pi\) preserves the harmonic
space.  Since \(B\) is Hermitian, the invariant orthogonal summands are
reducing.  The last assertion is
Theorem~\ref{thm:v2v81-harmonic}.
\end{proof}

\begin{remark}[Determinant bookkeeping]
\label{rem:v2v81-det}
The full cluster determinant factors as
\[
 \det B
 =
 \det B_{\rm ex}\,
 \det B_{\rm har}\,
 \det B_{\rm coex}.
\]
Zeros in the exact or coexact blocks can make the unsandwiched inverse
singular, but Theorem~\ref{thm:v2v81-exact} removes them from a physical
closed-source amplitude.  Applying the V80 jet test to the full
determinant without first making this Hodge reduction can therefore
demand cancellations that are already exact algebraic zeros.
\end{remark}

\subsection{Sharp regression when the kinetic matrix is not a chain map}
\label{sec:v2v81-regression}

\begin{counterexample}[An exact source can acquire a physical pole]
\label{sep:v2v81-noninvariant}
Let \(E^{-1}=\mathbb C u\) and
\(E^0=\mathbb C v\oplus\mathbb C w\), with
\[
 su=v,
 \qquad
 sv=sw=0.
\]
Let \(L(v)=0\), \(L(w)=1\), so \(Ls=0\), and take the exact source
\(R(1)=v\).  For a fixed \(a\ne0\), define on \(E^0\)
\begin{equation}
 B_0(t)=
 \begin{pmatrix}
 1&a\\
 a&a^2+t
 \end{pmatrix}.
\label{eq:v2v81-bad-B}
\end{equation}
Then \(\det B_0(t)=t\), and for \(t\ne0\),
\begin{equation}
 L B_0(t)^{-1}R
 =
 -\frac a t.
\label{eq:v2v81-bad-pole}
\end{equation}
Thus a physical output and an exact source do not suffice if the kinetic
matrix mixes the exact vector \(v\) with the physical vector \(w\).
\end{counterexample}

\begin{proof}
The inverse is
\[
 B_0(t)^{-1}
 =
 \frac1t
 \begin{pmatrix}
 a^2+t&-a\\
 -a&1
 \end{pmatrix}.
\]
Apply it to \(v\) and then apply \(L\).  No choice of the operator on
\(E^{-1}\) makes \(B\) commute with \(s\), because \(B_0(t)v\) has a
nonzero \(w\)-component.
\end{proof}

\begin{firewall}[BRST invariance is an analytic identity]
\label{fw:v2v81-invariance}
The relation \([B,s]=0\) must hold for the complete regulated cluster
kinetic operator after boundary conditions, gauge fixing, and
counterterms are included.  Classical gauge invariance alone does not
prove it for a regulator.  A nonzero commutator is a Ward defect whose
local part must be removed by allowed counterterms; a surviving
cohomology class is an anomaly obstruction.
\end{firewall}

\subsection{Parameter-dependent complexes}
\label{sec:v2v81-moving}

\begin{firewall}[A moving Hodge differential creates connection terms]
\label{fw:v2v81-moving}
V81 uses one fixed differential \(s\), adjoint \(s^*\), and harmonic
projection \(\Pi\) in the Kato-transported cluster frame.  If these
objects depend on \(t\), ordinary derivatives of
\(R=\Pi R+shR\) contain derivatives of all three operators.  The
pointwise cancellation still holds when
\[
 [B(t),s(t)]=0,\qquad L(t)s(t)=0,\qquad R_{\rm ex}(t)=s(t)C(t),
\]
but uniform differentiated estimates require a BRST-compatible
connection controlling \(s',\Pi'\), and \(h'\).  They do not follow from
the spectral Kato connection alone.
\end{firewall}

\begin{firewall}[Hodge-gap dependence remains]
\label{fw:v2v81-Hodge-gap}
If an exact representation \(R=sC\) is supplied directly, the algebraic
zero in \eqref{eq:v2v81-exact-numerator} does not require a bound on
\(C\).  If exactness is obtained from the closed-source decomposition,
the primitive \(hR\) carries the factor \(\gamma^{-1}\).
A regulator-uniform representative construction therefore still needs
the V68 Hodge gap.
\end{firewall}

\subsection{V81 conclusion and next acquisition}
\label{sec:v2v81-conclusion}

V81 proves the V80 pole-divisibility identities for every BRST-exact
cluster source, at every crossing order, provided the transported
cluster kinetic matrix is a BRST chain map and the output annihilates
boundaries.  Every closed source reduces exactly to its harmonic
component.  A two-dimensional regression shows that failure of the
chain-map identity can turn an exact source into a physical simple pole.

The remaining low singularity is therefore the harmonic matrix
\(B_{\rm har}(t)\).  V82 will separate protected physical massless
channels from removable constraint channels inside that harmonic block
and formulate the causal boundary-value estimate needed when a physical
pole must be retained rather than cancelled.

\clearpage
\section*{Second Volume, Version V82: Causal Boundary Values for a Physical Harmonic Crossing}
\addcontentsline{toc}{section}{Second Volume, Version V82: Causal Boundary Values for a Physical Harmonic Crossing}
\label{sec:v2v82}

\subsection{From uniform inversion to a distribution}
\label{sec:v2v82-purpose}

V81 removes exact and longitudinal cluster sources before inversion.  A
nonzero harmonic source can instead represent a physical massless
excitation.  Its pole must be retained.  The relevant object is then not
an \(L^1\) family of operator norms but a boundary-value distribution.

This version proves the local finite-dimensional statement for one
transverse energy crossing.  It supplies the exact principal-value and
on-shell delta terms and a test-function bound.  It also shows that a
tangential zero does not admit the same unrenormalized boundary value.

\subsection{Simple harmonic eigenvalue}
\label{sec:v2v82-simple}

Let \(J\subset\mathbb R\) be an interval and
\(B(t)\in\operatorname{Herm}(E)\) a \(C^2\) family on a finite-dimensional
Hermitian space.  Suppose that, near \(t_0\), one eigenvalue
\(\lambda(t)\) is simple and
\begin{equation}
 \lambda(t_0)=0,
 \qquad
 |\lambda'(t)|\ge c_0>0.
\label{eq:v2v82-transverse}
\end{equation}
Let \(\Pi(t)\) be its rank-one spectral projection and
\(Q(t)=I-\Pi(t)\).  Assume
\begin{equation}
 \operatorname{spec}\bigl(B(t)|_{Q(t)E}\bigr)
 \cap[-\delta,\delta]=\varnothing.
\label{eq:v2v82-other-gap}
\end{equation}
Define the regular reduced inverse
\begin{equation}
 G_Q(t)
 =
 Q(t)\bigl(B(t)|_{Q(t)E}\bigr)^{-1}Q(t).
\label{eq:v2v82-GQ}
\end{equation}

For Hilbert spaces \(X,Y\), let
\[
 R(t)\in\mathcal L(X,E),
 \qquad
 L(t)\in\mathcal L(E,Y)
\]
be \(C^1\), and put
\begin{equation}
 A(t)=L(t)\Pi(t)R(t),
 \qquad
 A_Q(t)=L(t)G_Q(t)R(t).
\label{eq:v2v82-A}
\end{equation}
The second term is continuous and uniformly bounded locally.

For \(\varepsilon>0\), define
\begin{equation}
 \mathcal T_\varepsilon^\mp(t)
 =
 L(t)\bigl(B(t)\mp i\varepsilon I\bigr)^{-1}R(t).
\label{eq:v2v82-Teps}
\end{equation}

\begin{theorem}[Operator-valued Sokhotski--Plemelj formula]
\label{thm:v2v82-boundary}
In
\(\mathcal D'(J;\mathcal L(X,Y))\),
\begin{equation}
 \mathcal T_\varepsilon^\mp
 \longrightarrow
 \operatorname{pv}\frac{A(t)}{\lambda(t)}
 \ \mathbin{\pm}\
 i\pi\frac{A(t_0)}{|\lambda'(t_0)|}\delta_{t_0}
 +A_Q(t)
\label{eq:v2v82-boundary}
\end{equation}
as \(\varepsilon\downarrow0\), after the neighbourhood is chosen so that
\(t_0\) is the only zero of \(\lambda\).  The principal value is defined
by the spectral coordinate \(x=\lambda(t)\).

For every compact \(K\) in this neighbourhood there is a constant
\(C_K\) such that the limiting distribution \(\mathcal T^\mp\) obeys
\begin{equation}
 \bigl\|\langle\mathcal T^\mp,\varphi\rangle\bigr\|
 \le
 C_K\|\varphi\|_{C^1}
 \qquad
 \bigl(\varphi\in C_c^1(K)\bigr).
\label{eq:v2v82-test-bound}
\end{equation}
Thus the retained physical pole is a distribution of order at most one.
\end{theorem}

\begin{proof}
The spectral decomposition is exact:
\begin{equation}
 (B(t)\mp i\varepsilon I)^{-1}
 =
 \frac{\Pi(t)}{\lambda(t)\mp i\varepsilon}
 +
 Q(t)\bigl(B(t)|_{Q(t)E}\mp i\varepsilon\bigr)^{-1}Q(t).
\label{eq:v2v82-resolvent-split}
\end{equation}
The second term converges uniformly to \(G_Q(t)\) by
\eqref{eq:v2v82-other-gap}.

By \eqref{eq:v2v82-transverse}, \(x=\lambda(t)\) is a \(C^1\)
coordinate.  In this coordinate, pairing the first term with a test
function becomes
\[
 \int
 \frac{
 A(t(x))\varphi(t(x))
 }{x\mp i\varepsilon}
 \frac{\mathrm dx}{|\lambda'(t(x))|}.
\]
The scalar boundary-value identity
\[
 \frac1{x\mp i0}
 =
 \operatorname{pv}\frac1x
 \mathbin{\pm}i\pi\delta_0
\]
then gives
\eqref{eq:v2v82-boundary}.  This identity follows directly by writing
\[
 \frac1{x\mp i\varepsilon}
 =
 \frac{x}{x^2+\varepsilon^2}
 \mathbin{\pm}
 i\frac{\varepsilon}{x^2+\varepsilon^2},
\]
using odd cancellation for the real part and the approximate identity
\(\varepsilon/(x^2+\varepsilon^2)\to\pi\delta_0\) for the imaginary part.

For the bound, write
\[
 F(x)
 =
 \frac{A(t(x))\varphi(t(x))}
 {|\lambda'(t(x))|}.
\]
On a symmetric \(x\)-interval containing the support,
\[
 \operatorname{pv}\int\frac{F(x)}x\,\mathrm dx
 =
 \int\frac{F(x)-F(0)}x\,\mathrm dx,
\]
after extending \(F\) by zero inside a slightly larger fixed chart if
necessary.  The mean-value theorem bounds this integral by a constant
times \(\|F\|_{C^1}\).  The delta term is bounded by \(\|F(0)\|\), and
the regular complement by \(\|\varphi\|_{L^1}\).  The \(C^1\) coefficient
bounds and \(c_0>0\) give \eqref{eq:v2v82-test-bound}.
\end{proof}

\begin{remark}[Coordinate-invariant delta weight]
\label{rem:v2v82-delta}
The formula
\[
 \delta(\lambda(t))
 =
 \frac{\delta_{t_0}}{|\lambda'(t_0)|}
\]
explains the Jacobian in
\eqref{eq:v2v82-boundary}.  It is the invariant one-dimensional spectral
density of the crossing, not a choice of eigenvector phase.
\end{remark}

\subsection{Consistency with the V80 cancellation criterion}
\label{sec:v2v82-V80}

\begin{corollary}[A Ward zero removes both boundary-value terms]
\label{cor:v2v82-Ward}
If
\begin{equation}
 A(t_0)=L(t_0)\Pi(t_0)R(t_0)=0,
\label{eq:v2v82-Ward-zero}
\end{equation}
then the delta term in
\eqref{eq:v2v82-boundary} vanishes and
\(A(t)/\lambda(t)\) extends continuously through \(t_0\).
The distributional formula therefore reduces to the ordinary continuous
extension of Corollary~\ref{cor:v2v80-simple}.
\end{corollary}

\begin{proof}
Since \(A\) is \(C^1\),
\(A(t)=(t-t_0)\widetilde A(t)\) with continuous
\(\widetilde A\).  Transversality gives
\(\lambda(t)=(t-t_0)\widetilde\lambda(t)\) with
\(|\widetilde\lambda|\ge c_0\).  Cancel the common factor.  The delta
coefficient is zero by hypothesis.
\end{proof}

\begin{assessment}[Cancellation and retention are compatible]
\label{ass:v2v82-compatible}
V80--V81 and V82 are two branches of one formula.  A Ward-exact matrix
element has \(A(t_0)=0\) and becomes an ordinary bounded amplitude.  A
physical harmonic matrix element can have \(A(t_0)\ne0\); its principal
value and on-shell delta measure are retained.  No uniform inverse norm is
asserted in the second case.
\end{assessment}

\subsection{Positive on-shell spectral weight}
\label{sec:v2v82-positive}

Take \(Y=X\) and \(L(t)=R(t)^*\).  Then
\begin{equation}
 A(t)=R(t)^*\Pi(t)R(t)\ge0.
\label{eq:v2v82-positive-A}
\end{equation}

\begin{corollary}[Positive imaginary boundary measure]
\label{cor:v2v82-positive}
For the minus-sign denominator \(B-i0\), the imaginary singular measure is
\begin{equation}
 \operatorname{Im}\mathcal T^-
 =
 \pi
 \frac{R(t_0)^*\Pi(t_0)R(t_0)}
 {|\lambda'(t_0)|}
 \delta_{t_0}
\label{eq:v2v82-positive-measure}
\end{equation}
up to the continuous imaginary part of any nonsymmetric external
coefficients.  In the symmetric case displayed here,
\eqref{eq:v2v82-positive-measure} is positive semidefinite.
\end{corollary}

\begin{proof}
Use the plus sign multiplying \(i\pi\) in
\eqref{eq:v2v82-boundary} for \(B-i0\), and
\eqref{eq:v2v82-positive-A}.
\end{proof}

This is the finite-cluster spectral weight that a physical massless mode
contributes.  Positivity is asserted only after the BRST/Hodge reduction
of V81 has removed indefinite exact and ghost directions.

\subsection{Differentiated boundary values}
\label{sec:v2v82-derivatives}

\begin{proposition}[External differentiation in distribution space]
\label{prop:v2v82-derivatives}
Assume \(B,L,R\) are \(C^{m+1}\) and the simple crossing hypotheses hold
on the support chart.  Then
\begin{equation}
 \partial_t^m\mathcal T_\varepsilon^\mp
 \longrightarrow
 \partial_t^m\mathcal T^\mp
\label{eq:v2v82-derivative-limit}
\end{equation}
in distributions for every \(0\le m\le5\).  The limiting derivative has
distribution order at most \(m+1\):
\begin{equation}
 \bigl\|
 \langle\partial_t^m\mathcal T^\mp,\varphi\rangle
 \bigr\|
 \le
 C_{K,m}\|\varphi\|_{C^{m+1}}.
\label{eq:v2v82-derivative-bound}
\end{equation}
\end{proposition}

\begin{proof}
Differentiation is continuous on \(\mathcal D'\), so
\eqref{eq:v2v82-derivative-limit} follows from
Theorem~\ref{thm:v2v82-boundary}.  By the definition of distributional
derivative,
\[
 \langle\partial_t^m\mathcal T^\mp,\varphi\rangle
 =
 (-1)^m
 \langle\mathcal T^\mp,\partial_t^m\varphi\rangle.
\]
Apply \eqref{eq:v2v82-test-bound} to
\(\partial_t^m\varphi\).
\end{proof}

\begin{firewall}[Distributional derivatives are not V71 operator norms]
\label{fw:v2v82-not-V71}
The bound \eqref{eq:v2v82-derivative-bound} controls testing in the energy
parameter.  It is not the symmetric Hilbert-scale operator norm of V71.
A complete physical estimate must combine the V79 external-complement
bounds with this distributional cluster term using a causal composition
and wavefront-set theorem.
\end{firewall}

\subsection{Sharp failure at a tangential zero}
\label{sec:v2v82-tangent}

\begin{counterexample}[The unrenormalized \(t^2-i0\) limit diverges]
\label{sep:v2v82-tangent}
Let \(B(t)=(t^2)\).  For a nonnegative test function \(\varphi\) equal to
one near zero,
\begin{equation}
 \operatorname{Im}
 \int
 \frac{\varphi(t)}{t^2-i\varepsilon}\,\mathrm dt
 =
 \int
 \frac{\varepsilon\varphi(t)}
 {t^4+\varepsilon^2}\,\mathrm dt
 \longrightarrow+\infty.
\label{eq:v2v82-tangent-divergence}
\end{equation}
Hence \((t^2-i\varepsilon)^{-1}\) has no limit in
\(\mathcal D'(\mathbb R)\) without an additional subtraction or
renormalization.
\end{counterexample}

\begin{proof}
Restrict to \(|t|<a\), where \(\varphi=1\), and put
\(t=\varepsilon^{1/2}y\).  The integral becomes
\[
 \varepsilon^{-1/2}
 \int_{|y|<a/\sqrt\varepsilon}
 \frac{\mathrm dy}{y^4+1},
\]
which diverges like a positive constant times
\(\varepsilon^{-1/2}\).
\end{proof}

\begin{firewall}[Finite type does not imply a canonical causal pullback]
\label{fw:v2v82-finite-type}
A nonzero higher derivative of \(\det B\) controls the size of a crossing
set, but a tangential eigenvalue can make the naive \(i\varepsilon\)
boundary value diverge.  One must either resolve the crossing in a
different spectral variable, supply a local subtraction fixed by
renormalization conditions, or prove that the assembled numerator has
the V80 divisibility required to cancel the tangency.
\end{firewall}

\subsection{Scope of the causal statement}
\label{sec:v2v82-scope}

\begin{firewall}[Local energy theorem only]
\label{fw:v2v82-local}
V82 is a one-parameter finite-cluster boundary-value theorem.  It does not
prove analytic continuation from the Euclidean metric cone, global
hyperbolic propagation, microlocal spectrum condition, compatibility of
products of several singular propagators, or inclusive soft-photon
finiteness.  Those require spacetime wavefront and observable estimates.
\end{firewall}

\begin{firewall}[The sign must come from the physical prescription]
\label{fw:v2v82-sign}
Both boundary values in
\eqref{eq:v2v82-boundary} exist.  Retarded, advanced, and Feynman
propagators select signs and energy orientations through the Lorentzian
field equation.  The Euclidean finite-cluster algebra alone does not
choose among them.
\end{firewall}

\subsection{V82 conclusion and next acquisition}
\label{sec:v2v82-conclusion}

V82 constructs the physical simple-crossing cluster propagator as an
operator-valued boundary-value distribution.  Its singular part is the
sum of a principal value and a positive on-shell delta measure after
physical BRST projection.  A Ward zero recovers the ordinary bounded
extension of V80, while a tangential crossing provides a sharp
unrenormalized failure.

V83 will formulate a wavefront-compatible composition theorem for this
finite-cluster distribution with the uniformly controlled external
resolvent and smooth interaction vertices.  The theorem must identify
which products are defined before any loop or nonlinear causal packet is
claimed.
\clearpage
\section*{Second Volume, Version V83: Wavefront-Compatible Composition of Physical Cluster Poles}
\addcontentsline{toc}{section}{Second Volume, Version V83: Wavefront-Compatible Composition of Physical Cluster Poles}
\label{sec:v2v83}

\subsection{Purpose}
\label{sec:v2v83-purpose}

V82 constructs one physical harmonic crossing as a boundary-value
distribution.  A nonlinear or perturbative packet composes this
distribution with vertices, external resolvents, and possibly other
cluster poles.  Such products cannot be justified by multiplying
operator norms.

This version identifies the local wavefront directions of the two
boundary values.  Smooth scale-bounded factors do not change them.
Several poles with the same causal orientation can be multiplied, while
opposite boundary values at the same crossing fail the unrenormalized
product test.  This is the first composition rule needed before a
causal low-cluster packet is inserted into the V69--V79 estimates.

\subsection{One-sided wavefront of the scalar boundary value}
\label{sec:v2v83-scalar}

Use the Fourier convention
\[
 \widehat f(\xi)=\int_{\mathbb R}e^{-it\xi}f(t)\,\mathrm dt.
\]
Near a transverse zero \(t_0\), let
\[
 u^\mp(t)=\frac1{\lambda(t)\mp i0},
 \qquad
 \lambda(t_0)=0,
 \qquad
 \lambda'(t_0)\ne0.
\]

\begin{lemma}[Oriented conormal directions]
\label{lem:v2v83-WF}
The wavefront sets satisfy
\begin{align}
 \operatorname{WF}(u^-)
 &\subset
 \{(t_0,\xi):\xi\lambda'(t_0)<0\},
 \label{eq:v2v83-WF-minus}\\
 \operatorname{WF}(u^+)
 &\subset
 \{(t_0,\xi):\xi\lambda'(t_0)>0\}.
 \label{eq:v2v83-WF-plus}
\end{align}
If the pole coefficient is nonzero, the corresponding inclusion is an
equality.
\end{lemma}

\begin{proof}
First take \(\lambda(t)=t\).  The V82 boundary formula gives
\[
 \frac1{t-i0}
 =
 \operatorname{pv}\frac1t+i\pi\delta_0,
 \qquad
 \frac1{t+i0}
 =
 \operatorname{pv}\frac1t-i\pi\delta_0.
\]
With the stated Fourier convention,
\[
 \mathcal F\left(\operatorname{pv}\frac1t\right)
 =
 -i\pi\operatorname{sgn}\xi.
\]
Therefore the Fourier transform of \((t-i0)^{-1}\) is supported at high
frequency in the negative half-line, while that of
\((t+i0)^{-1}\) is supported at high frequency in the positive
half-line.  Localizing by a smooth cutoff changes the Fourier transforms
only by rapidly decreasing errors in the excluded cone.  This proves the
two one-sided wavefront statements at zero.

The map \(t\mapsto x=\lambda(t)\) is a local diffeomorphism.  Pullback of
a covector gives
\(\xi_t=\lambda'(t_0)\xi_x\), which yields
\eqref{eq:v2v83-WF-minus}--\eqref{eq:v2v83-WF-plus}.
A nonzero scalar coefficient cannot remove the singular Fourier
half-line, proving equality in that case.
\end{proof}

\begin{remark}[Why V82's coarse distribution order was insufficient]
\label{rem:v2v83-oriented}
The full conormal set
\(\{(t_0,\xi):\xi\ne0\}\) would forget the causal orientation.  The
one-sided refinement is what permits powers of one boundary value:
two covectors in the same open half-line never sum to zero.
\end{remark}

\subsection{Operator-valued cluster distribution}
\label{sec:v2v83-operator}

Let \(\mathcal T^\mp\) be the V82 limit:
\[
 \mathcal T^\mp
 =
 A(t)u^\mp(t)+A_Q(t),
\]
where \(A\) is \(C^1\) and \(A_Q\) is continuous; for wavefront statements
below assume the coefficient families are smooth.  Since the cluster
space is finite-dimensional, wavefront regularity can be tested in its
matrix entries.

\begin{proposition}[Smooth operator factors preserve the causal cone]
\label{prop:v2v83-smooth-factors}
Let \(V_0(t)\) and \(V_1(t)\) be smooth operator families with compatible
domains and codomains.  Then
\begin{equation}
 V_0\mathcal T^\mp V_1
\label{eq:v2v83-sandwich}
\end{equation}
is a well-defined operator-valued distribution and
\begin{equation}
 \operatorname{WF}(V_0\mathcal T^\mp V_1)
 \subset
 \operatorname{WF}(\mathcal T^\mp).
\label{eq:v2v83-smooth-WF}
\end{equation}
The statement remains valid when \(V_0,V_1\) act between the fixed
Hilbert-scale spaces of V79 with uniform smooth operator bounds.
\end{proposition}

\begin{proof}
Multiplication of a distribution by a smooth scalar function is always
defined and does not enlarge its wavefront set.  Expand the
finite-cluster factor in matrix entries and apply this fact to every
entry.  Composition with the smooth external operators is a finite sum
of such products.  The Hilbert-scale version is the same argument in
the Banach spaces of bounded operators; smooth bounded multiplication is
continuous there.
\end{proof}

\begin{corollary}[The V79 external inverse is composition-safe]
\label{cor:v2v83-external}
Let \(\widehat G_{\rm ext}(t)\) be the uniformly scale-bounded external
inverse from Theorem~\ref{thm:v2v79-external}.  Assume in addition the
smooth parameter dependence required by classical wavefront calculus,
and let the interaction vertices be smooth on the same scale window.
Any packet containing one
factor \(\mathcal T^\mp\), finitely many smooth vertices, and finitely
many factors \(\widehat G_{\rm ext}\) is a well-defined
operator-valued distribution.  Its singular wavefront is contained in
the oriented cone of
Lemma~\ref{lem:v2v83-WF}.
\end{corollary}

\begin{proof}
The external inverse and vertices are smooth operator multipliers.
Apply Proposition~\ref{prop:v2v83-smooth-factors} successively.
\end{proof}

\subsection{Products of equally oriented poles}
\label{sec:v2v83-same-sign}

We use the standard distribution product criterion: \(uv\) is defined if
there is no point \(t\) and covectors
\((t,\xi)\in\operatorname{WF}(u)\),
\((t,\eta)\in\operatorname{WF}(v)\) with
\(\xi+\eta=0\).

\begin{theorem}[Same-sign cluster products are defined]
\label{thm:v2v83-same-sign}
Let \(u_1^\mp,\ldots,u_N^\mp\) be simple transverse boundary values whose
zeros coincide at \(t_0\) and whose oriented derivatives
\(\lambda_j'(t_0)\) have the same sign.  Then
\begin{equation}
 u_1^\mp\cdots u_N^\mp
\label{eq:v2v83-same-product}
\end{equation}
is a well-defined distribution.  The same holds for compatible
operator-valued cluster factors separated by smooth vertices and V79
external inverses.

In the common-coordinate scalar case,
\begin{equation}
 \left(\frac1{x\mp i0}\right)^N
 =
 \frac{(-1)^{N-1}}{(N-1)!}
 \partial_x^{N-1}\left(\frac1{x\mp i0}\right).
\label{eq:v2v83-power}
\end{equation}
\end{theorem}

\begin{proof}
By Lemma~\ref{lem:v2v83-WF}, all wavefront covectors of the factors lie in
the same open half-line.  No nonempty finite sum of those covectors is
zero.  The distribution product criterion therefore applies
successively and proves
\eqref{eq:v2v83-same-product}.  Smooth operator coefficients do not add
wavefront directions by
Proposition~\ref{prop:v2v83-smooth-factors}.

For \(\varepsilon>0\), ordinary differentiation gives
\[
 \partial_x^{N-1}(x\mp i\varepsilon)^{-1}
 =
 (-1)^{N-1}(N-1)!(x\mp i\varepsilon)^{-N}.
\]
Pass to the boundary value in distributions to obtain
\eqref{eq:v2v83-power}.
\end{proof}

\begin{assessment}[Repeated causal poles increase distribution order]
\label{ass:v2v83-order}
The product exists, but
\eqref{eq:v2v83-power} shows that an \(N\)-fold pole is a distribution of
order at most \(N\).  Existence of the product is therefore only the
first gate; the vertex and test-function regularity must pay the
increased derivative order.
\end{assessment}

\subsection{Opposite orientations at one crossing}
\label{sec:v2v83-opposite}

\begin{counterexample}[The opposite-sign product has no unrenormalized limit]
\label{sep:v2v83-opposite}
For \(\varepsilon>0\),
\begin{equation}
 \frac1{x-i\varepsilon}
 \frac1{x+i\varepsilon}
 =
 \frac1{x^2+\varepsilon^2}.
\label{eq:v2v83-opposite-product}
\end{equation}
If \(\varphi\ge0\) and \(\varphi=1\) near zero, then
\begin{equation}
 \int
 \frac{\varphi(x)}{x^2+\varepsilon^2}\,\mathrm dx
 \ge
 \frac c\varepsilon
 \longrightarrow\infty.
\label{eq:v2v83-opposite-divergence}
\end{equation}
Thus the common-\(\varepsilon\) product of opposite boundary values has
no distributional limit without subtraction.
\end{counterexample}

\begin{proof}
On \(|x|<a\), put \(x=\varepsilon y\):
\[
 \int_{-a}^a\frac{\mathrm dx}{x^2+\varepsilon^2}
 =
 \varepsilon^{-1}
 \int_{-a/\varepsilon}^{a/\varepsilon}
 \frac{\mathrm dy}{1+y^2}.
\]
The last integral is bounded below by a positive constant for small
\(\varepsilon\), proving
\eqref{eq:v2v83-opposite-divergence}.
\end{proof}

The wavefront criterion predicts the failure: the two boundary values
carry opposite covectors at the same point, and these can sum to zero.

\begin{firewall}[Failure of the criterion is not a declaration of impossibility]
\label{fw:v2v83-renormalization}
When the wavefront product criterion fails, a renormalized extension may
still exist after local subtractions.  Such an extension is not
canonical; its local ambiguity must be classified by scaling degree,
symmetry, BRST cohomology, and renormalization conditions.  V83 proves
only that the naive product cannot be inserted into the continuum packet.
\end{firewall}

\subsection{Ward factors lower the singular order}
\label{sec:v2v83-Ward}

\begin{proposition}[A transverse numerator cancels one boundary factor]
\label{prop:v2v83-Ward-factor}
Let \(a\) be smooth and suppose
\[
 a(t_0)=0.
\]
Then, locally,
\begin{equation}
 a(t)\frac1{\lambda(t)\mp i0}
 =
 \frac{a(t)}{\lambda(t)}
\label{eq:v2v83-Ward-factor}
\end{equation}
is an ordinary smooth function if \(a,\lambda\) are smooth and both have
simple zeros with \(\lambda'(t_0)\ne0\).  More generally, a numerator
divisible by \(\lambda^r\) cancels \(r\) equally oriented boundary
factors.
\end{proposition}

\begin{proof}
Hadamard's lemma gives
\(a(t)=(t-t_0)\widetilde a(t)\) and
\(\lambda(t)=(t-t_0)\widetilde\lambda(t)\), with
\(\widetilde\lambda(t_0)\ne0\).  Their quotient is smooth.  Multiplication
of the V82 boundary formula by \(a\) kills the delta coefficient and
turns the principal value into this quotient.  Iterate the argument for
higher powers.
\end{proof}

\begin{assessment}[Order of operations]
\label{ass:v2v83-order-of-operations}
The V81 BRST cancellation and V80 adjugate divisibility must be applied
before the wavefront product is formed.  They can remove a pole or lower
its order.  Taking absolute norms or multiplying opposite prescriptions
first creates divergences which the later Ward sum cannot repair inside
the same estimate.
\end{assessment}

\subsection{Several crossing hypersurfaces}
\label{sec:v2v83-hypersurfaces}

For a parameter \(k\in\mathbb R^d\), a transverse eigenvalue equation
\(\lambda_j(k)=0\) defines a smooth hypersurface with conormal
\(\xi_j=c_j\,\mathrm d\lambda_j(k)\).

\begin{proposition}[Local multi-parameter product criterion]
\label{prop:v2v83-multiparameter}
A product of finitely many cluster boundary values is defined near \(k\)
provided there is no choice of allowed signed conormals
\(\xi_j\) from their wavefront sets, not all zero, such that
\begin{equation}
 \xi_1+\cdots+\xi_N=0.
\label{eq:v2v83-conormal-sum}
\end{equation}
Smooth V79 external inverses and vertices do not affect this condition.
\end{proposition}

\begin{proof}
This is the distribution product criterion applied to the pullbacks of
the one-dimensional boundary values by the submersions \(\lambda_j\).
The pullback wavefronts are the corresponding oriented conormals.
Proposition~\ref{prop:v2v83-smooth-factors} handles the smooth operator
factors.
\end{proof}

\begin{firewall}[Loop integration requires a pushforward theorem]
\label{fw:v2v83-loop}
The existence of a product in external and internal momentum variables
does not by itself justify integrating out loop momentum.  Pushforward
requires proper support on the relevant projection and exclusion of
wavefront covectors normal to its fibres, or an explicit renormalized
extension.  Those hypotheses have not yet been proved for the full SM
graph family.
\end{firewall}

\begin{firewall}[Global causality is still separate]
\label{fw:v2v83-global}
The sign cones in V83 are local energy-frequency statements.  A
Lorentzian spacetime propagator must also satisfy the appropriate
spacetime characteristic relation, support property, and microlocal
spectrum condition.  The finite cluster calculation does not construct
a global retarded or Feynman parametrix on the Einstein background.
\end{firewall}

\subsection{V83 conclusion and next acquisition}
\label{sec:v2v83-conclusion}

V83 proves that a physical cluster pole composes freely with the smooth
vertices and uniformly controlled external inverse of V79.  Products of
poles with compatible causal orientations are defined and equal higher
boundary-value derivatives.  Opposite orientations at the same crossing
have a sharp \(1/\varepsilon\) regression and require renormalization.
Ward factors must be applied before this product test.

V84 will address the next operation: pushforward over a finite-dimensional
internal energy variable.  It will give a sufficient nonstationary
criterion for the causal cluster product to be integrable as a
distributional pushforward and isolate threshold pinches where that
criterion fails.

\clearpage
\section*{Second Volume, Version V84: Nonpinched Pushforward of Causal Cluster Products}
\addcontentsline{toc}{section}{Second Volume, Version V84: Nonpinched Pushforward of Causal Cluster Products}
\label{sec:v2v84}

\subsection{Purpose and a clarification}
\label{sec:v2v84-purpose}

V83 determined when products of local cluster boundary values exist.
The next operation integrates an internal energy variable.  Proper
support is sufficient to define this pushforward as a distribution.
A further wavefront condition is needed to prove that the pushforward is
smooth in the external parameters.  V84 states that condition
quantitatively and identifies its failure as a causal pinch.

Let
\[
 \pi:K\times\mathbb R_\omega\longrightarrow K,
 \qquad
 \pi(k,\omega)=k,
\]
where \(K\subset\mathbb R^d\) is compact.  All amplitudes below have
support in one fixed compact \(\omega\)-interval, so \(\pi\) is proper on
their support.

\subsection{A finite causal product}
\label{sec:v2v84-product}

Let
\[
 \lambda_j\in C^\infty(K\times\mathbb R;\mathbb R),
 \qquad
 1\le j\le N,
\]
and choose signs \(\sigma_j\in\{+1,-1\}\).  Write
\begin{equation}
 u_j
 =
 \frac1{\lambda_j+i0\,\sigma_j}.
\label{eq:v2v84-factor}
\end{equation}
At a point where \(\lambda_j=0\), the allowed conormal covectors have the
form
\begin{equation}
 c_j\,\mathrm d\lambda_j,
 \qquad
 c_j\sigma_j>0,
\label{eq:v2v84-oriented-conormal}
\end{equation}
with the Fourier convention of V83.  Thus the sign of \(c_j\) records the
chosen boundary value.

Let \(a(k,\omega)\) be a smooth, compactly supported operator-valued
coefficient.  Suppose the V83 product criterion holds:
\begin{equation}
 \sum_{j\in S}c_j\,\mathrm d\lambda_j\ne0
\label{eq:v2v84-product-condition}
\end{equation}
for every nonempty active set
\(S\subset\{j:\lambda_j=0\}\) and every allowed choice
\eqref{eq:v2v84-oriented-conormal}.  Then
\begin{equation}
 u=a\prod_{j=1}^Nu_j
\label{eq:v2v84-u}
\end{equation}
is a well-defined compactly supported operator-valued distribution.

\begin{lemma}[Wavefront of the causal product]
\label{lem:v2v84-product-WF}
The wavefront set of \(u\) is contained in covectors
\begin{equation}
 \left(
 k,\omega;\
 \sum_{j\in S}c_j\,\partial_k\lambda_j,\
 \sum_{j\in S}c_j\,\partial_\omega\lambda_j
 \right),
\label{eq:v2v84-product-WF}
\end{equation}
where \(S\) is a nonempty active set and
\(c_j\sigma_j>0\).
\end{lemma}

\begin{proof}
The pullback of the one-dimensional boundary value by the submersion
\(\lambda_j\) has wavefront contained in its oriented conormal
\(c_j\,\mathrm d\lambda_j\).  The wavefront product theorem adds
covectors of factors singular at the same point.  Multiplication by the
smooth coefficient \(a\) does not enlarge the result.
\end{proof}

\subsection{The vertical nonpinching condition}
\label{sec:v2v84-nonpinch}

Assume there is a constant \(c_{\rm np}>0\) such that, at every common
active point and for all allowed conic coefficients,
\begin{equation}
 \left|
 \sum_{j\in S}
 c_j\,\partial_\omega\lambda_j(k,\omega)
 \right|
 \ge
 c_{\rm np}\sum_{j\in S}|c_j|.
\label{eq:v2v84-nonpinch}
\end{equation}
Because the coefficients are conic, this is a uniform angular
separation from the horizontal cotangent bundle.

\begin{theorem}[Smooth nonpinched pushforward]
\label{thm:v2v84-pushforward}
Under
\eqref{eq:v2v84-product-condition} and
\eqref{eq:v2v84-nonpinch}, the proper pushforward
\begin{equation}
 \mathcal I(k)
 =
 \pi_*u
 =
 \int_{\mathbb R}
 a(k,\omega)
 \prod_{j=1}^N
 \frac{\mathrm d\omega}
 {\lambda_j(k,\omega)+i0\,\sigma_j}
\label{eq:v2v84-pushforward}
\end{equation}
is a smooth operator-valued function of \(k\).

For fixed compact supports, a fixed number of factors, uniform phase
seminorms, and one uniform \(c_{\rm np}\), its \(C^m\) seminorm is bounded
by finitely many coefficient seminorms:
\begin{equation}
 \|\mathcal I\|_{C^m(K)}
 \le
 C_m
 \|a\|_{C^{M_m}(K\times\mathbb R)}
\label{eq:v2v84-seminorm}
\end{equation}
for some finite \(M_m\).
\end{theorem}

\begin{proof}
Proper support defines \(\pi_*u\) by
\[
 \langle\pi_*u,\varphi\rangle
 =
 \langle u,\varphi\circ\pi\rangle.
\]
The pushforward wavefront theorem gives
\[
 \operatorname{WF}(\pi_*u)
 \subset
 \left\{
 (k,\xi):
 \text{for some }\omega,\
 (k,\omega;\xi,0)\in\operatorname{WF}(u)
 \right\}.
\]
By Lemma~\ref{lem:v2v84-product-WF}, a covector on the right would
satisfy
\[
 \sum_{j\in S}
 c_j\,\partial_\omega\lambda_j=0.
\]
This contradicts \eqref{eq:v2v84-nonpinch}.  Therefore
\(\operatorname{WF}(\pi_*u)=\varnothing\), which is equivalent to
smoothness.

The distribution product and proper pushforward maps are continuous on
sets whose wavefront cones remain in fixed closed cones satisfying the
displayed separation.  Compactness and the uniform lower bound
\(c_{\rm np}\) place the present family in such a set.  Continuity gives
\eqref{eq:v2v84-seminorm}; only finitely many seminorms enter each fixed
output derivative order.
\end{proof}

\begin{assessment}[What the nonpinching inequality says]
\label{ass:v2v84-nonpinch}
The vertical component of every allowed singular covector is nonzero.
Consequently no singular direction survives integration as a covector in
the external variables.  This is stronger than existence of the
integral as a distribution and weaker than pointwise absolute
integrability.
\end{assessment}

\subsection{Elementary one-pole model}
\label{sec:v2v84-one-pole}

\begin{proposition}[Direct change-of-variable proof for one pole]
\label{prop:v2v84-one-pole}
Let
\[
 \mathcal I(k)
 =
 \int
 \frac{a(k,\omega)}
 {\lambda(k,\omega)-i0}\,\mathrm d\omega,
\]
where \(a\) is smooth and compactly supported and
\begin{equation}
 |\partial_\omega\lambda|\ge c_0>0
\label{eq:v2v84-one-transverse}
\end{equation}
on the zero set meeting the support.  Then \(\mathcal I\) is smooth in
\(k\).
\end{proposition}

\begin{proof}
Use \(x=\lambda(k,\omega)\) near every zero.  A finite partition of unity
reduces the singular part to
\[
 \int
 \frac{b(k,x)}{x-i0}\,\mathrm dx
 =
 \operatorname{pv}\int\frac{b(k,x)}x\,\mathrm dx
 +i\pi b(k,0),
\]
with \(b\) smooth and compactly supported.  Subtract \(b(k,0)\) inside a
fixed symmetric neighbourhood of zero.  The remaining quotient is
smooth and integrable with all \(k\)-derivatives.  Terms away from the
zero set are ordinary smooth integrals.
\end{proof}

This direct calculation is the \(N=1\) instance of
Theorem~\ref{thm:v2v84-pushforward}.

\subsection{A useful same-orientation criterion}
\label{sec:v2v84-same}

\begin{corollary}[Parallel energy slopes cannot pinch]
\label{cor:v2v84-same-slopes}
Suppose all active factors have the same boundary sign and there is one
\(\tau\in\{+1,-1\}\) such that
\begin{equation}
 \tau\sigma_j\partial_\omega\lambda_j
 \ge c_0>0
\label{eq:v2v84-parallel-slopes}
\end{equation}
on their zero sets.  Then
\eqref{eq:v2v84-nonpinch} holds with \(c_{\rm np}=c_0\), so the compactly
supported internal-energy pushforward is smooth.
\end{corollary}

\begin{proof}
Write \(c_j=\sigma_jd_j\) with \(d_j>0\).  Then
\(\tau c_j\partial_\omega\lambda_j
=d_j\tau\sigma_j\partial_\omega\lambda_j\ge c_0d_j\).
Thus the vertical terms all have the same sign after multiplication by
\(\tau\), and their sum cannot cancel.
\end{proof}

For example, factors
\[
 (\omega-E_j(k)-i0)^{-1}
\]
all have the same energy slope and orientation.  Their coincident poles
are legitimate higher boundary values, and their compactly supported
energy pushforward is smooth.

\subsection{Regulator-uniform decay survives a nonpinched pushforward}
\label{sec:v2v84-uniform}

Let \(a_n\) take values in bounded maps between fixed Hilbert-scale
spaces and suppose, for all seminorms required in
\eqref{eq:v2v84-seminorm},
\begin{equation}
 \|a_n\|_{C^{M_m}}
 \le
 C\Lambda_n^{-\beta}.
\label{eq:v2v84-coefficient-decay}
\end{equation}

\begin{corollary}[Transfer of the ultraviolet rate]
\label{cor:v2v84-rate}
If the support, phase seminorms, product cones, and nonpinching constant
are uniform in \(n\), then
\begin{equation}
 \|\pi_*u_n\|_{C^m(K)}
 \le
 C_m\Lambda_n^{-\beta}.
\label{eq:v2v84-rate}
\end{equation}
Thus a V70--V79 external-resolvent decay rate is not lost merely by a
finite nonpinched internal-energy integration.
\end{corollary}

\begin{proof}
Apply \eqref{eq:v2v84-seminorm} to \(a_n\) and then use
\eqref{eq:v2v84-coefficient-decay}.  Uniformity of all geometric inputs
makes \(C_m\) independent of \(n\).
\end{proof}

\begin{firewall}[Uniform support and cone separation are substantive]
\label{fw:v2v84-uniform}
If the internal support grows with the regulator, if a phase derivative
degenerates, or if \(c_{\rm np}\to0\), the constants in
\eqref{eq:v2v84-seminorm} can diverge.  Pointwise nonpinching at each
finite regulator is not enough for continuum removal.
\end{firewall}

\subsection{Sharp pinch regression}
\label{sec:v2v84-pinch}

\begin{counterexample}[Opposite poles destroy a smooth pushforward]
\label{sep:v2v84-pinch}
Let \(\chi\ge0\) be smooth, compactly supported, and equal to one near
zero.  Define
\begin{equation}
 \mathcal I_\varepsilon(k)
 =
 \int
 \frac{\chi(\omega)}
 {(\omega-i\varepsilon)(\omega-k+i\varepsilon)}
 \,\mathrm d\omega.
\label{eq:v2v84-pinch-family}
\end{equation}
At \(k=0\),
\begin{equation}
 \mathcal I_\varepsilon(0)
 =
 \int
 \frac{\chi(\omega)}
 {\omega^2+\varepsilon^2}\,\mathrm d\omega
 \ge
 \frac c\varepsilon.
\label{eq:v2v84-pinch-divergence}
\end{equation}
Hence there is no locally uniform, smooth pushforward through the
threshold \(k=0\).  Any limit in the external variable can only be a
singular distribution and requires its own boundary-value analysis.
\end{counterexample}

\begin{proof}
This is the scaling calculation of
Counterexample~\ref{sep:v2v83-opposite} integrated in the internal
variable.  The two allowed conormals have opposite vertical components,
so they can sum to a horizontal covector; at \(k=0\) they also violate the
product criterion at the common pole.
\end{proof}

\begin{assessment}[Pinch versus ordinary on-shell weight]
\label{ass:v2v84-pinch}
A single transverse pole produces the finite on-shell delta measure of
V82.  A pinch traps the contour between opposite prescriptions and
creates a divergent product before pushforward.  The latter requires
Ward cancellation, a width/resummation mechanism, or a local
renormalized extension; it is not another copy of the single-particle
spectral measure.
\end{assessment}

\subsection{Several internal variables}
\label{sec:v2v84-several}

For a submersion
\[
 \pi:(k,\ell)\longmapsto k,
 \qquad
 \ell\in\mathbb R^p,
\]
the vertical component of an allowed conormal sum is
\[
 \sum_{j\in S}c_j\,\partial_\ell\lambda_j
 \in(\mathbb R^p)^*.
\]

\begin{proposition}[Multi-energy nonpinching criterion]
\label{prop:v2v84-multi}
If this vertical sum is nonzero for every nontrivial allowed conormal
combination and is uniformly separated from zero after conic
normalization, then the properly supported pushforward is smooth and
obeys the analogue of
\eqref{eq:v2v84-seminorm}.
\end{proposition}

\begin{proof}
Repeat the proof of
Theorem~\ref{thm:v2v84-pushforward}.  A covector contributing to the
pushforward wavefront must annihilate every vertical tangent, so its
vertical component must be zero.  The hypothesis excludes precisely
that possibility.
\end{proof}

\begin{firewall}[This is not yet the full Landau analysis]
\label{fw:v2v84-Landau}
The conormal equation detects a local necessary geometry for threshold
pinches.  A full graph theorem must also include momentum-conservation
constraints, boundary faces of parameter space, ultraviolet
renormalized extensions, and the signs fixed by the physical propagator.
V84 does not classify all Standard Model Landau singularities.
\end{firewall}

\subsection{V84 conclusion and next acquisition}
\label{sec:v2v84-conclusion}

V84 proves that a properly supported causal cluster product has a smooth
internal-energy pushforward when every allowed singular covector has a
uniformly nonzero vertical component.  Under uniform cone separation,
the V70--V79 regulator decay passes through the integration unchanged.
Opposite prescriptions yield a sharp threshold-pinch regression.

V85 is the next mandatory Yang--Mills/Navier--Stokes audit.  It will use
the newest companion results to decide whether to attack the pinched
finite cluster by assembled Ward cancellation, introduce a controlled
width, or move upstream to the regulator-uniform microlocal construction.

\clearpage
\section*{Second Volume, Version V85: Companion Audit and Exact Pinch-to-External-Pole Transfer}
\addcontentsline{toc}{section}{Second Volume, Version V85: Companion Audit and Exact Pinch-to-External-Pole Transfer}
\label{sec:v2v85}

\subsection{Mandatory companion audit}
\label{sec:v2v85-audit}

The newest companion continuations in the shared folder were inspected
without modification:
\begin{center}
\small
\begin{tabular}{p{0.48\linewidth}r p{0.32\linewidth}}
\toprule
file&bytes&SHA--256\\
\midrule
\texttt{Renewal\_Yang\_Mills\_Mass\_Gap\_Research\_Notes\_V68.tex}
&954053&
\texttt{ccc5d377dbf1d461e214df771be416fb484d419350076c2af4690ad220bb004e}
\\
\texttt{Renewal\_NS\_Stability\_Research\_Notes\_V31.tex}
&887537&
\texttt{f1121b62238ad5491bb7cf03635ea9934942edf573ed8faaa9ee79e1d38a6696}
\\
\bottomrule
\end{tabular}
\end{center}
Yang--Mills V67 also remains in the shared folder as the exact predecessor
of V68.  It was not deleted.  The NS note is unchanged from the V80
audit.

\begin{assessment}[Yang--Mills V68 transfer]
\label{ass:v2v85-YM}
Yang--Mills V68 proves that private coordinates factor as one exact
multiplier of an entire first-connectivity family before any triangle
inequality or output split.  Its remaining hot \(2+2\) obstruction must
retain a joint packet containing both incident influences and the
physical output ancestry of the same shared endpoint.

For the SM threshold, the corresponding proof order is to retain both
oppositely prescribed denominators, their common numerator, and the
internal-energy integration as one packet.  Estimating either pole
separately loses the transfer identity below.  No Wilson-product estimate
is imported.
\end{assessment}

\begin{assessment}[Navier--Stokes V31 transfer]
\label{ass:v2v85-NS}
Navier--Stokes V31 remains at a correlation gate: its exact dyadic mark
and formation ladder recover the full critical cost when the active
source and exit structure are separated.  This again rules out a scalar
size estimate of the pinch set as progress.  The numerator value at the
actual common pole must remain correlated with the two causal factors.
\end{assessment}

The audit therefore selects exact assembled pushforward before either a
width insertion or a positive denominator bound.

\subsection{The universal two-pole integral}
\label{sec:v2v85-universal}

For \(\varepsilon>0\), define
\begin{equation}
 I_\varepsilon(k)
 =
 \int_{\mathbb R}
 \frac{\mathrm d\omega}
 {(\omega-i\varepsilon)(\omega-k+i\varepsilon)}.
\label{eq:v2v85-Ieps}
\end{equation}
The integrand decays quadratically, so the integral is absolutely
convergent.

\begin{lemma}[Exact pinch transfer at positive width]
\label{lem:v2v85-exact}
For every real \(k\),
\begin{equation}
 I_\varepsilon(k)
 =
 -\frac{2\pi i}{k-2i\varepsilon}.
\label{eq:v2v85-exact}
\end{equation}
Consequently,
\begin{equation}
 I_\varepsilon
 \longrightarrow
 -\frac{2\pi i}{k-i0}
 =
 -2\pi i\,\operatorname{pv}\frac1k
 +2\pi^2\delta_0
\label{eq:v2v85-boundary}
\end{equation}
in \(\mathcal D'(\mathbb R_k)\).
\end{lemma}

\begin{proof}
Close the \(\omega\)-contour in the upper half-plane.  The integrand is
\(O(|\omega|^{-2})\), and the only enclosed pole is
\(\omega=i\varepsilon\).  Its residue is
\[
 \frac1{2i\varepsilon-k}.
\]
The residue theorem gives
\[
 I_\varepsilon(k)
 =
 \frac{2\pi i}{2i\varepsilon-k}
 =
 -\frac{2\pi i}{k-2i\varepsilon}.
\]
The scalar boundary-value identity from V82 proves
\eqref{eq:v2v85-boundary}.  Replacing
\(2\varepsilon\) by \(\varepsilon\) does not change the distributional
boundary value.
\end{proof}

\begin{assessment}[The pinch is transferred, not erased]
\label{ass:v2v85-transfer}
At \(k=0\), \(I_\varepsilon(0)=\pi/\varepsilon\), so no smooth limit
exists.  After the complete internal-energy pushforward, however, the
family has the well-defined external distribution
\eqref{eq:v2v85-boundary}.  The trapped pair has become one external
causal pole with a fixed residue.
\end{assessment}

\subsection{Cutoffs change only the smooth remainder}
\label{sec:v2v85-cutoff}

Let \(\chi\in C_c^\infty(\mathbb R)\) equal one on
\([-2a,2a]\), and restrict \(|k|<a\).  Put
\begin{equation}
 I_{\varepsilon,\chi}(k)
 =
 \int
 \frac{\chi(\omega)\,\mathrm d\omega}
 {(\omega-i\varepsilon)(\omega-k+i\varepsilon)}.
\label{eq:v2v85-cutoff}
\end{equation}

\begin{proposition}[Universal local singular part]
\label{prop:v2v85-cutoff}
There is a smooth function \(S_\chi(k)\) on \((-a,a)\) such that
\begin{equation}
 I_{\varepsilon,\chi}
 \longrightarrow
 -\frac{2\pi i}{k-i0}+S_\chi(k)
\label{eq:v2v85-cutoff-limit}
\end{equation}
in distributions.
\end{proposition}

\begin{proof}
Subtract \eqref{eq:v2v85-cutoff} from
\eqref{eq:v2v85-Ieps}.  The difference has coefficient
\(1-\chi(\omega)\), whose support stays a positive distance from both
\(\omega=0\) and \(\omega=k\) for \(|k|<a\).  Every \(k\)-derivative of
its integrand is dominated by an integrable function uniformly for small
\(\varepsilon\).  Dominated differentiation gives a smooth limit
\(J_\chi(k)\).  Thus
\(S_\chi=-J_\chi\), and
Lemma~\ref{lem:v2v85-exact} supplies the singular part.
\end{proof}

\subsection{A smooth assembled numerator}
\label{sec:v2v85-numerator}

Let \(X,Y\) be Hilbert spaces and let
\[
 A\in
 C_c^\infty\bigl(
 (-a,a)_k\times\mathbb R_\omega;\
 \mathcal L(X,Y)
 \bigr).
\]
Define
\begin{equation}
 \mathcal P_\varepsilon[A](k)
 =
 \int
 \frac{A(k,\omega)}
 {(\omega-i\varepsilon)(\omega-k+i\varepsilon)}
 \,\mathrm d\omega.
\label{eq:v2v85-P}
\end{equation}

\begin{theorem}[Pinch residue is evaluation of the complete numerator]
\label{thm:v2v85-residue}
There exists
\(S_A\in C^\infty((-a,a);\mathcal L(X,Y))\) such that
\begin{equation}
 \mathcal P_\varepsilon[A]
 \longrightarrow
 -\frac{2\pi i\,A(0,0)}{k-i0}
 +S_A(k)
\label{eq:v2v85-residue}
\end{equation}
in operator-valued distributions on \((-a,a)\).
In particular,
\begin{equation}
 A(0,0)=0
 \quad\Longrightarrow\quad
 \lim_{\varepsilon\downarrow0}
 \mathcal P_\varepsilon[A]
 \in
 C^\infty((-a,a);\mathcal L(X,Y)).
\label{eq:v2v85-Ward-zero}
\end{equation}
\end{theorem}

\begin{proof}
Choose \(\chi\) as above and write, by Hadamard's lemma in the variables
\((k,\omega)\),
\begin{equation}
 A(k,\omega)
 =
 A(0,0)\chi(\omega)
 +kB(k,\omega)+\omega C(k,\omega)+A_{\rm far}(k,\omega),
\label{eq:v2v85-Hadamard}
\end{equation}
where \(B,C,A_{\rm far}\) are smooth and compactly supported and
\(A_{\rm far}\) is supported away from the common pole
\((k,\omega)=(0,0)\).  The far term has a smooth pushforward by ordinary
dominated differentiation.

The constant term is
\(A(0,0)I_{\varepsilon,\chi}\), so
Proposition~\ref{prop:v2v85-cutoff} gives the displayed singular part.
For the \(kB\) term use the exact identity
\begin{equation}
 \frac{k}
 {(\omega-i\varepsilon)(\omega-k+i\varepsilon)}
 =
 \frac1{\omega-k+i\varepsilon}
 -
 \frac1{\omega-i\varepsilon}
 +
 \frac{2i\varepsilon}
 {(\omega-i\varepsilon)(\omega-k+i\varepsilon)}.
\label{eq:v2v85-k-cancel}
\end{equation}
For the \(\omega C\) term use
\begin{equation}
 \frac{\omega}
 {(\omega-i\varepsilon)(\omega-k+i\varepsilon)}
 =
 \frac1{\omega-k+i\varepsilon}
 +
 \frac{i\varepsilon}
 {(\omega-i\varepsilon)(\omega-k+i\varepsilon)}.
\label{eq:v2v85-omega-cancel}
\end{equation}
The single-pole pushforwards have smooth limits by
Proposition~\ref{prop:v2v84-one-pole}.  The terms with an explicit
factor \(\varepsilon\) tend to zero after subtraction of their universal
pinch part, and their possible singular part is
\(\varepsilon(k-2i\varepsilon)^{-1}\), which tends to zero in
\(\mathcal D'\).  Smooth coefficient localization and a finite partition
of unity preserve this conclusion.  Hence every term other than the
constant numerator contributes to the smooth function \(S_A\).
\end{proof}

\begin{remark}[Regularity version]
\label{rem:v2v85-regularity}
If \(A\) has only \(C^m\) regularity, the same argument gives a finite
\(C^{m-2}\) remainder after the required one-pole differentiations,
provided \(m\ge2\).  The smooth statement is used here to avoid endpoint
bookkeeping.
\end{remark}

\subsection{Assembly before the pinch residue}
\label{sec:v2v85-assembly}

Suppose the numerator is a finite Ward or diagrammatic sum
\begin{equation}
 A=\sum_{\nu=1}^J A_\nu.
\label{eq:v2v85-sum}
\end{equation}

\begin{corollary}[Only the assembled value controls the external pole]
\label{cor:v2v85-assembled}
The singular part of the complete pushforward is
\begin{equation}
 -\frac{2\pi i}{k-i0}
 \sum_{\nu=1}^J A_\nu(0,0).
\label{eq:v2v85-summed-residue}
\end{equation}
Thus the external pole cancels if
\begin{equation}
 \sum_{\nu=1}^J A_\nu(0,0)=0,
\label{eq:v2v85-summed-Ward}
\end{equation}
even when no individual \(A_\nu(0,0)\) vanishes.
\end{corollary}

\begin{proof}
Apply Theorem~\ref{thm:v2v85-residue} to the sum and use linearity of
evaluation and pushforward.
\end{proof}

\begin{counterexample}[Termwise norms destroy residue cancellation]
\label{sep:v2v85-termwise}
Let
\[
 A_1(0,0)=T,
 \qquad
 A_2(0,0)=-T
\]
with \(T\ne0\).  The complete pinch pushforward is smooth, but the
positive termwise majorant assigns the nonintegrable external size
\begin{equation}
 \frac{4\pi\|T\|}{|k|}
\label{eq:v2v85-termwise-majorant}
\end{equation}
to the two residues.
\end{counterexample}

\begin{proof}
The assembled residue vanishes by
\eqref{eq:v2v85-summed-residue}.  Taking the norms of the two individual
coefficients before summing replaces \(T-T\) by \(2\|T\|\), leaving the
displayed magnitude of the two factors
\(-2\pi i(k-i0)^{-1}\).
\end{proof}

\begin{assessment}[Exact analogue of the companion proof order]
\label{ass:v2v85-proof-order}
The YM V68 private multiplier belongs to the whole first-connectivity
family.  The SM pinch residue likewise belongs to the complete numerator
evaluated at the joint on-shell point.  Neither can be assigned
consistently to separately normed letters.  The NS V31 warning appears
as the \(|k|^{-1}\) majorant: measuring the threshold without the
correlated numerator recovers the critical divergence.
\end{assessment}

\subsection{BRST and physical alternatives}
\label{sec:v2v85-BRST}

\begin{corollary}[BRST-exact pinch sources are smooth after pushforward]
\label{cor:v2v85-BRST}
If the complete on-shell numerator is in the exact sector covered by
Theorem~\ref{thm:v2v81-exact}, then
\[
 A(0,0)=0,
\]
and the pinched internal-energy pushforward is smooth.
\end{corollary}

\begin{proof}
V81 makes the physical exact-sector numerator identically zero before
the denominators are divided.  Apply
\eqref{eq:v2v85-Ward-zero}.
\end{proof}

If \(A(0,0)\ne0\) in the harmonic physical sector, the pushforward has the
external pole in
\eqref{eq:v2v85-residue}.  That pole is retained with the V82 boundary
prescription and is not estimated by a positive inverse norm.

\begin{firewall}[A width is not inserted by hand]
\label{fw:v2v85-width}
The parameter \(\varepsilon\) in V85 specifies a boundary value.  It is
not a physical decay width.  Replacing it by an interaction-dependent
finite width changes the spectral problem and requires a self-energy
resummation, positivity, and Ward-compatibility theorem.
\end{firewall}

\begin{firewall}[The local normal form must be proved]
\label{fw:v2v85-normal-form}
The denominators in
\eqref{eq:v2v85-P} are the transverse local normal form of two pinching
simple eigenvalues.  A general SM graph must first be reduced to this
form by smooth spectral coordinates with uniform Jacobians and a
separate treatment of multiple or tangential crossings.  V85 does not
assert that every threshold is simple.
\end{firewall}

\begin{firewall}[No loop-family conclusion yet]
\label{fw:v2v85-loop}
V85 treats one internal energy and one opposite pair.  It does not prove
uniform summability over graphs, compatibility of simultaneous pinches,
renormalized spatial-momentum integration, or the microlocal spectrum
condition on a curved Lorentzian background.
\end{firewall}

\subsection{V85 conclusion and next acquisition}
\label{sec:v2v85-conclusion}

The mandatory companion audit has been completed.  V85 proves that the
canonical opposite-pole pinch, kept as one packet, pushes forward to a
single external boundary-value pole with universal coefficient
\(-2\pi i\).  For a smooth operator numerator, the residue is exactly its
assembled value at the joint on-shell point.  A Ward zero there makes the
pushforward smooth; termwise norms destroy this cancellation.

V86 will replace the scalar normal form by a finite Hermitian crossing
pair.  It will express the pinch residue through the two spectral
projections and crossing Jacobians, so the BRST/Hodge test can be applied
directly to the matrix-valued on-shell carrier.
\clearpage
\section*{Second Volume, Version V86: Matrix Pinch Residue from Two Hermitian Crossing Blocks}
\addcontentsline{toc}{section}{Second Volume, Version V86: Matrix Pinch Residue from Two Hermitian Crossing Blocks}
\label{sec:v2v86}

\subsection{Purpose}
\label{sec:v2v86-purpose}

V85 computed the universal scalar pinch.  A Standard Model cluster
packet contains finite Hermitian kinetic blocks rather than scalar
denominators.  This version proves the reduction to the scalar normal
form for two simple eigenvalues, including every complementary spectral
term.

The residue is the complete interaction numerator compressed between
the two on-shell eigenprojections.  Its invariant delta coefficient is
divided by the absolute Jacobian of the two dispersion equations.  This
is the matrix carrier on which the V81 BRST/Hodge cancellation must be
tested.

\subsection{Two simple Hermitian crossings}
\label{sec:v2v86-setup}

Let \(E_1,E_2,X,Y\) be finite-dimensional Hermitian spaces.  On a
neighbourhood \(\mathcal U\) of
\((k_0,\omega_0)=(0,0)\), let
\[
 B_j(k,\omega)\in\operatorname{Herm}(E_j),
 \qquad j=1,2,
\]
be smooth.  Suppose \(B_j\) has one simple smooth eigenvalue
\(\lambda_j(k,\omega)\) with rank-one spectral projection
\(\Pi_j(k,\omega)\), and that
\begin{align}
 \lambda_1(0,0)&=\lambda_2(0,0)=0,
 \label{eq:v2v86-common-zero}\\
 \partial_\omega\lambda_j(k,\omega)&\ge c_\omega>0
 \label{eq:v2v86-positive-slopes}
\end{align}
near the common zero.  Let \(Q_j=I-\Pi_j\), and assume the complementary
spectra obey
\begin{equation}
 \operatorname{spec}
 \bigl(B_j|_{Q_jE_j}\bigr)
 \cap[-\delta,\delta]
 =
 \varnothing.
\label{eq:v2v86-complement-gap}
\end{equation}

The implicit-function theorem supplies smooth roots
\begin{equation}
 \lambda_j(k,r_j(k))=0,
 \qquad
 r_j(0)=0.
\label{eq:v2v86-roots}
\end{equation}
Define their separation
\begin{equation}
 g(k)=r_2(k)-r_1(k).
\label{eq:v2v86-g}
\end{equation}
Assume the external crossing is transverse:
\begin{equation}
 g'(0)\ne0.
\label{eq:v2v86-g-transverse}
\end{equation}

Let
\begin{align*}
 R(k,\omega)&\in\mathcal L(X,E_2),\\
 V(k,\omega)&\in\mathcal L(E_2,E_1),\\
 L(k,\omega)&\in\mathcal L(E_1,Y)
\end{align*}
be smooth and supported in a sufficiently small subset of
\(\mathcal U\).  For \(\varepsilon>0\), define
\begin{equation}
 \mathcal M_\varepsilon(k)
 =
 \int_{\mathbb R}
 L(B_1-i\varepsilon I)^{-1}
 V(B_2+i\varepsilon I)^{-1}R
 \,\mathrm d\omega.
\label{eq:v2v86-Meps}
\end{equation}
All coefficients in the integrand are evaluated at \((k,\omega)\).

\subsection{Exact spectral splitting}
\label{sec:v2v86-splitting}

Define
\begin{equation}
 G_{j,\varepsilon}^{\mp}
 =
 Q_j
 \bigl(B_j|_{Q_jE_j}\mp i\varepsilon I\bigr)^{-1}
 Q_j.
\label{eq:v2v86-Gjeps}
\end{equation}

\begin{lemma}[Singular line plus uniformly regular complement]
\label{lem:v2v86-resolvent-split}
One has
\begin{align}
 (B_1-i\varepsilon I)^{-1}
 &=
 \frac{\Pi_1}{\lambda_1-i\varepsilon}
 +G_{1,\varepsilon}^{-},
 \label{eq:v2v86-B1-split}\\
 (B_2+i\varepsilon I)^{-1}
 &=
 \frac{\Pi_2}{\lambda_2+i\varepsilon}
 +G_{2,\varepsilon}^{+}.
 \label{eq:v2v86-B2-split}
\end{align}
The complementary families converge smoothly to
\begin{equation}
 G_j
 =
 Q_j(B_j|_{Q_jE_j})^{-1}Q_j
\label{eq:v2v86-Gj}
\end{equation}
with all derivatives on a smaller compact neighbourhood.
\end{lemma}

\begin{proof}
Each \(\Pi_j\) is a reducing spectral projection, so inversion is the
orthogonal direct sum of the scalar inverse on its range and the inverse
on \(Q_jE_j\).  The gap
\eqref{eq:v2v86-complement-gap} bounds the latter inverse and all its
parameter derivatives through the ordinary differentiated inverse
formula.  Letting \(\varepsilon\) tend to zero gives smooth convergence.
\end{proof}

Expanding \eqref{eq:v2v86-Meps} with
\eqref{eq:v2v86-B1-split}--\eqref{eq:v2v86-B2-split} produces four
terms.  Three contain at most one singular eigenvalue denominator.

\begin{proposition}[Every nonpinched spectral term is smooth]
\label{prop:v2v86-three-smooth}
The terms
\begin{align*}
 \int LG_{1,\varepsilon}^{-}VG_{2,\varepsilon}^{+}R\,\mathrm d\omega,
 \qquad
 \int
 \frac{L\Pi_1VG_{2,\varepsilon}^{+}R}
 {\lambda_1-i\varepsilon}\,\mathrm d\omega,
 \qquad
 \int
 \frac{LG_{1,\varepsilon}^{-}V\Pi_2R}
 {\lambda_2+i\varepsilon}\,\mathrm d\omega
\end{align*}
converge smoothly as functions of \(k\).
\end{proposition}

\begin{proof}
The first integrand converges smoothly and has fixed compact support.
For each remaining term,
\(\partial_\omega\lambda_j\ge c_\omega\), so the one-pole change of
variables in Proposition~\ref{prop:v2v84-one-pole} applies.  The
complementary resolvents and all other coefficients are smooth.
\end{proof}

The only possible pinch is therefore
\begin{equation}
 \mathcal P_\varepsilon(k)
 =
 \int
 \frac{A(k,\omega)}
 {(\lambda_1(k,\omega)-i\varepsilon)
  (\lambda_2(k,\omega)+i\varepsilon)}
 \,\mathrm d\omega,
\label{eq:v2v86-Peps}
\end{equation}
where the complete projected numerator is
\begin{equation}
 A(k,\omega)
 =
 L(k,\omega)\Pi_1(k,\omega)
 V(k,\omega)\Pi_2(k,\omega)R(k,\omega).
\label{eq:v2v86-A}
\end{equation}

\subsection{Straightening the two dispersion roots}
\label{sec:v2v86-straightening}

Hadamard's lemma relative to the roots gives
\begin{equation}
 \lambda_j(k,\omega)
 =
 a_j(k,\omega)\bigl(\omega-r_j(k)\bigr),
\label{eq:v2v86-factor}
\end{equation}
where
\begin{equation}
 a_j(k,\omega)
 =
 \int_0^1
 \partial_\omega\lambda_j
 \bigl(k,r_j(k)+\theta(\omega-r_j(k))\bigr)
 \,\mathrm d\theta
 \ge c_\omega.
\label{eq:v2v86-aj}
\end{equation}

\begin{lemma}[Positive rescaling preserves the boundary sign]
\label{lem:v2v86-rescaling}
If \(a>0\) is smooth, then
\begin{equation}
 \frac1{a(t)x(t)\mp i0}
 =
 \frac1{a(t)}
 \frac1{x(t)\mp i0}
\label{eq:v2v86-rescaling}
\end{equation}
as distributions wherever \(x\) is a submersion.
\end{lemma}

\begin{proof}
Use the scalar principal-value and delta formula:
\[
 \operatorname{pv}\frac1{ax}
 =
 a^{-1}\operatorname{pv}\frac1x,
 \qquad
 \delta(ax)=a^{-1}\delta(x)
\]
for \(a>0\).  The two terms have the same boundary sign, proving the
identity.
\end{proof}

Translate the internal coordinate by
\[
 x=\omega-r_1(k).
\]
The second root is then at \(x=g(k)\).  Lemma
\ref{lem:v2v86-rescaling} absorbs the positive factors \(a_1,a_2\) into
the numerator.  Its on-shell value is
\begin{equation}
 C_0
 =
 \frac{
 L_0\Pi_{1,0}V_0\Pi_{2,0}R_0
 }{
 (\partial_\omega\lambda_1)_0
 (\partial_\omega\lambda_2)_0
 },
\label{eq:v2v86-C0}
\end{equation}
where the subscript \(0\) means evaluation at \((0,0)\).

\subsection{The matrix pinch theorem}
\label{sec:v2v86-main}

\begin{theorem}[Hermitian two-cluster pinch residue]
\label{thm:v2v86-pinch}
Under
\eqref{eq:v2v86-common-zero}--\eqref{eq:v2v86-g-transverse},
there is a smooth
\(S(k)\in\mathcal L(X,Y)\) such that
\begin{equation}
 \mathcal M_\varepsilon
 \longrightarrow
 -\frac{2\pi i\,C_0}{g(k)-i0}
 +S(k)
\label{eq:v2v86-main}
\end{equation}
in operator-valued distributions near \(k=0\).
Equivalently, the singular part is
\begin{equation}
 -2\pi i\,C_0\operatorname{pv}\frac1{g(k)}
 +
 2\pi^2C_0\,\delta(g(k)).
\label{eq:v2v86-PV-delta}
\end{equation}
\end{theorem}

\begin{proof}
Lemma~\ref{lem:v2v86-resolvent-split} and
Proposition~\ref{prop:v2v86-three-smooth} reduce the problem to
\eqref{eq:v2v86-Peps}.  Equations
\eqref{eq:v2v86-factor}--\eqref{eq:v2v86-rescaling}, followed by
\(x=\omega-r_1(k)\), put its boundary-value limit in the form
\[
 \int
 \frac{\widetilde A(k,x)}
 {(x-i0)(x-g(k)+i0)}
 \,\mathrm dx,
\]
with \(\widetilde A(0,0)=C_0\).  Since
\(g'(0)\ne0\), use \(g\) as the local external coordinate and apply
Theorem~\ref{thm:v2v85-residue}.  That theorem gives the singular term
\(-2\pi iC_0/(g-i0)\); every change of cutoff, variable coefficient, and
complementary spectral term contributes to the smooth remainder.
Expanding the boundary value gives
\eqref{eq:v2v86-PV-delta}.
\end{proof}

\subsection{Invariant crossing Jacobian}
\label{sec:v2v86-Jacobian}

Define
\begin{equation}
 J_0
 =
 \det
 \begin{pmatrix}
  \partial_k\lambda_1&\partial_\omega\lambda_1\\
  \partial_k\lambda_2&\partial_\omega\lambda_2
 \end{pmatrix}_{(0,0)}.
\label{eq:v2v86-J}
\end{equation}

\begin{lemma}[Root separation and joint Jacobian]
\label{lem:v2v86-J}
One has
\begin{equation}
 g'(0)
 =
 \frac{J_0}
 {(\partial_\omega\lambda_1)_0
  (\partial_\omega\lambda_2)_0}.
\label{eq:v2v86-g-J}
\end{equation}
Thus \eqref{eq:v2v86-g-transverse} is equivalent to \(J_0\ne0\).
\end{lemma}

\begin{proof}
Differentiate
\(\lambda_j(k,r_j(k))=0\):
\[
 r_j'(0)
 =
 -\frac{(\partial_k\lambda_j)_0}
 {(\partial_\omega\lambda_j)_0}.
\]
Subtract \(r_1'\) from \(r_2'\) and put the two fractions over a common
denominator.  The numerator is exactly \(J_0\).
\end{proof}

\begin{corollary}[Invariant on-shell delta carrier]
\label{cor:v2v86-delta}
The delta part of the matrix pinch is
\begin{equation}
 2\pi^2
 \frac{
 L_0\Pi_{1,0}V_0\Pi_{2,0}R_0
 }{|J_0|}
 \,\delta_0(k).
\label{eq:v2v86-invariant-delta}
\end{equation}
\end{corollary}

\begin{proof}
Use
\[
 \delta(g(k))
 =
 \frac{\delta_0(k)}{|g'(0)|}
\]
in \eqref{eq:v2v86-PV-delta}, then substitute
\eqref{eq:v2v86-C0} and
\eqref{eq:v2v86-g-J}.  The positive energy slopes cancel against the
absolute Jacobian factor.
\end{proof}

\begin{assessment}[Geometric content of the residue]
\label{ass:v2v86-geometric}
The residue has four independent pieces: the two on-shell spectral
projections, the interaction vertex between them, the external
source/output maps, and the crossing Jacobian.  Complementary eigenmodes
do not enter the singular coefficient.  This is the finite-dimensional
matrix version of cutting both internal lines on shell.
\end{assessment}

\subsection{Ward cancellation at the matrix carrier}
\label{sec:v2v86-Ward}

\begin{corollary}[Necessary on-shell matrix tested by BRST]
\label{cor:v2v86-Ward}
If the complete projected numerator satisfies
\begin{equation}
 L_0\Pi_{1,0}V_0\Pi_{2,0}R_0=0,
\label{eq:v2v86-Ward-zero}
\end{equation}
then the pinched pushforward
\(\mathcal M_\varepsilon\) has a smooth limit near \(k=0\).
This holds, in particular, when the assembled carrier is in the
BRST-exact sector of Theorem~\ref{thm:v2v81-exact}.
\end{corollary}

\begin{proof}
The condition makes \(C_0=0\).  Apply
Theorem~\ref{thm:v2v86-pinch} and the smooth-remainder conclusion in
Theorem~\ref{thm:v2v85-residue}.  V81 supplies the stated sufficient
algebraic zero for an exact carrier.
\end{proof}

\begin{proposition}[Positive cut carrier]
\label{prop:v2v86-positive}
Suppose \(X=Y\) and the complete on-shell matrix can be written
\begin{equation}
 L_0\Pi_{1,0}V_0\Pi_{2,0}R_0=T^*T
\label{eq:v2v86-positive-factor}
\end{equation}
for some \(T\).  Then the delta coefficient
\eqref{eq:v2v86-invariant-delta} is positive semidefinite.
\end{proposition}

\begin{proof}
The scalar \(2\pi^2/|J_0|\) is positive and \(T^*T\ge0\).
\end{proof}

\begin{firewall}[Positivity requires the complete cut identity]
\label{fw:v2v86-positivity}
Hermiticity of \(B_1,B_2\) alone does not imply
\eqref{eq:v2v86-positive-factor}.  That factorization is an optical or
unitarity statement about the complete physical intermediate-state sum.
Ghost, longitudinal, and individual diagram contributions need not be
positive separately.
\end{firewall}

\subsection{Sharp failures of the simple-pair theorem}
\label{sec:v2v86-failures}

\begin{firewall}[Opposite energy slopes]
\label{fw:v2v86-slopes}
The sign normalization
\eqref{eq:v2v86-positive-slopes} places the first pole above and the
second below the real \(\omega\)-axis.  If the energy slopes have opposite
signs, the same written \(i0\) prescriptions put the poles on the same
side and there is no local pinch of this type.  One must determine the
orientation from the signed slopes rather than from the displayed
denominator signs alone.
\end{firewall}

\begin{firewall}[Multiple eigenvalues]
\label{fw:v2v86-multiple}
If either crossing eigenvalue is multiple, rank-one projections and a
single scalar phase no longer suffice.  The complete degenerate cluster
matrix must be retained, and its simultaneous normal form can carry
several residues or higher-order poles.  V86 does not diagonalize a
variable-multiplicity crossing.
\end{firewall}

\begin{firewall}[Tangential joint zero]
\label{fw:v2v86-tangent}
If \(J_0=0\), then \(g'(0)=0\).  The external boundary value can be
tangential and the V82 regression applies: the naive limit may diverge
or require local subtraction.  Finite type of \(g\) is not a replacement
for the nonzero Jacobian in V86.
\end{firewall}

\begin{counterexample}[Projection norms miss an assembled zero]
\label{sep:v2v86-projection-norms}
Let the on-shell projected numerator be a sum
\[
 L_0\Pi_{1,0}V_0\Pi_{2,0}R_0
 =
 T+(-T)=0
\]
with \(T\ne0\).  The matrix pinch is smooth by
Corollary~\ref{cor:v2v86-Ward}, whereas estimating the two contributions
separately produces two copies of the external
\(|g|^{-1}\) majorant.
\end{counterexample}

\begin{proof}
The exact residue is linear in the complete projected numerator.
The triangle inequality replaces its zero by \(2\|T\|\).
\end{proof}

\subsection{V86 conclusion and next acquisition}
\label{sec:v2v86-conclusion}

V86 proves the finite Hermitian version of the V85 pinch theorem.  Only
the two simple on-shell eigenprojections contribute to the singular
coefficient.  The external delta carrier is the complete projected
interaction numerator divided by the absolute joint crossing Jacobian.
Every complementary spectral term is smooth.

The next unresolved step is uniformity.  V87 will compare regulated and
continuum crossing data.  It will give sufficient differentiated
norm-resolvent and transversality hypotheses under which the
eigenprojections, roots, Jacobians, and matrix pinch residues converge
without spectral pollution.
\clearpage
\section*{Second Volume, Version V87: Regulator-Stable Hermitian Pinch Data}
\addcontentsline{toc}{section}{Second Volume, Version V87: Regulator-Stable Hermitian Pinch Data}
\label{sec:v2v87}

\subsection{Purpose and upstream input}
\label{sec:v2v87-purpose}

V86 identifies the matrix pinch residue at one regulator.  A continuum
claim requires the crossing eigenvalues, spectral projections, root
separation, Jacobian, and causal boundary distribution to converge
together.

V39 already proves transfer of isolated Riesz clusters from uniform
norm-resolvent convergence and proves \(C^M\) convergence of compressed
cluster matrices under differentiated contour-resolvent hypotheses.
V87 starts from that finite-matrix conclusion and proves stability of
all additional data used by V86.  A final corollary states the V39
operator hypotheses which supply the matrix assumptions.

\subsection{Convergent compressed kinetic blocks}
\label{sec:v2v87-setup}

Let \(\mathcal U\) be a compact rectangle about \((0,0)\) in the
\((k,\omega)\)-plane.  In compatible Kato frames, let
\[
 B_{j,n},B_j:
 \mathcal U\longrightarrow\operatorname{Herm}(E_j),
 \qquad j=1,2,
\]
be \(C^m\), where \(m\ge3\), and assume
\begin{equation}
 B_{j,n}\longrightarrow B_j
 \quad\hbox{in }C^m(\mathcal U).
\label{eq:v2v87-B-convergence}
\end{equation}
Suppose \(B_j\) has the simple isolated eigenvalue
\(\lambda_j\) and projection \(\Pi_j\) of V86, with one uniform
complementary gap \(\delta>0\).  Assume
\begin{align}
 \lambda_1(0,0)&=\lambda_2(0,0)=0,
 \label{eq:v2v87-limit-zero}\\
 \partial_\omega\lambda_j&\ge2c_\omega>0,
 \label{eq:v2v87-limit-slopes}\\
 g'(0)&\ne0,
 \qquad
 g=r_2-r_1.
 \label{eq:v2v87-limit-transverse}
\end{align}

Let the coefficient families
\[
 L_n,\ V_n,\ R_n
\]
converge to \(L,V,R\) in \(C^m\), with one fixed compact support inside
\(\mathcal U\).

\subsection{Spectral data convergence}
\label{sec:v2v87-spectral}

\begin{lemma}[Simple eigenprojections and eigenvalues converge]
\label{lem:v2v87-spectral}
After \(\mathcal U\) is reduced and \(n\) is sufficiently large, each
\(B_{j,n}\) has exactly one rank-one eigenprojection
\(\Pi_{j,n}\) inside the continuum Riesz contour.  Its eigenvalue
\(\lambda_{j,n}\) can be labelled so that
\begin{align}
 \Pi_{j,n}&\longrightarrow\Pi_j
 &&\text{in }C^m(\mathcal U),
 \label{eq:v2v87-P-convergence}\\
 \lambda_{j,n}&\longrightarrow\lambda_j
 &&\text{in }C^m(\mathcal U).
 \label{eq:v2v87-lambda-convergence}
\end{align}
The complementary spectral gaps are at least \(\delta/2\), and
\begin{equation}
 \partial_\omega\lambda_{j,n}\ge c_\omega.
\label{eq:v2v87-reg-slopes}
\end{equation}
\end{lemma}

\begin{proof}
Choose fixed small contours enclosing the two continuum simple
eigenvalues and no complementary spectrum throughout the reduced
rectangle.  Matrix convergence and the resolvent identity give
\(C^m\) convergence of the contour resolvents.  Integrating them gives
\eqref{eq:v2v87-P-convergence}.  Once
\(\|\Pi_{j,n}-\Pi_j\|<1\), the projections have the same rank one.

The first spectral moment has the contour formula
\begin{equation}
 \lambda_{j,n}\Pi_{j,n}
 =
 \frac1{2\pi i}
 \int_{\Gamma_j}
 z(z-B_{j,n})^{-1}\,\mathrm dz.
\label{eq:v2v87-moment}
\end{equation}
It converges in \(C^m\) to
\(\lambda_j\Pi_j\).  Taking the finite-dimensional trace and using
\(\operatorname{Tr}\Pi_{j,n}=1\) proves
\eqref{eq:v2v87-lambda-convergence}.  The gap and slope statements
follow from uniform convergence and
\eqref{eq:v2v87-limit-slopes}.
\end{proof}

\subsection{Uniform roots and the shifted joint crossing}
\label{sec:v2v87-roots}

\begin{theorem}[Regulated crossing point converges]
\label{thm:v2v87-crossing}
For all sufficiently large \(n\), there are unique \(C^m\) root functions
\(r_{j,n}(k)\) near zero such that
\begin{equation}
 \lambda_{j,n}(k,r_{j,n}(k))=0.
\label{eq:v2v87-reg-roots}
\end{equation}
They satisfy
\begin{equation}
 r_{j,n}\longrightarrow r_j
 \quad\hbox{in }C^m.
\label{eq:v2v87-root-convergence}
\end{equation}
Let
\[
 g_n=r_{2,n}-r_{1,n}.
\]
There is a unique \(k_n\) near zero such that
\begin{equation}
 g_n(k_n)=0,
\label{eq:v2v87-kn}
\end{equation}
and
\begin{equation}
 k_n\longrightarrow0,
 \qquad
 \omega_n:=r_{1,n}(k_n)=r_{2,n}(k_n)\longrightarrow0.
\label{eq:v2v87-crossing-convergence}
\end{equation}
Moreover,
\begin{equation}
 g_n'\longrightarrow g'
 \quad\hbox{uniformly},
 \qquad
 |g_n'(k_n)|\ge\frac12|g'(0)|.
\label{eq:v2v87-gprime}
\end{equation}
\end{theorem}

\begin{proof}
The uniform slope bound
\eqref{eq:v2v87-reg-slopes} gives a uniform implicit-function
neighbourhood for \(\lambda_{j,n}=0\).  Uniform convergence gives
\(r_{j,n}\to r_j\).  Differentiating the implicit equation expresses
each derivative of \(r_{j,n}\) through derivatives of
\(\lambda_{j,n}\), lower root derivatives, and a power of
\((\partial_\omega\lambda_{j,n})^{-1}\).  Induction proves
\eqref{eq:v2v87-root-convergence} through order \(m\).

Hence \(g_n\to g\) in \(C^m\).  Since \(g(0)=0\) and \(g'(0)\ne0\), the
uniform implicit-function theorem gives one zero \(k_n\) of \(g_n\),
with \(k_n\to0\) and the derivative lower bound in
\eqref{eq:v2v87-gprime}.  Root convergence then gives
\(\omega_n\to0\).
\end{proof}

\subsection{Jacobian and on-shell carrier}
\label{sec:v2v87-carrier}

At the regulated crossing define
\begin{align}
 J_n
 &=
 \det
 \begin{pmatrix}
  \partial_k\lambda_{1,n}&\partial_\omega\lambda_{1,n}\\
  \partial_k\lambda_{2,n}&\partial_\omega\lambda_{2,n}
 \end{pmatrix}_{(k_n,\omega_n)},
 \label{eq:v2v87-Jn}\\
 A_n
 &=
 \bigl(
 L_n\Pi_{1,n}V_n\Pi_{2,n}R_n
 \bigr)(k_n,\omega_n).
 \label{eq:v2v87-An}
\end{align}
Let \(J_0,A_0\) be the continuum quantities of V86.

\begin{proposition}[Residue coefficients converge]
\label{prop:v2v87-residue}
One has
\begin{equation}
 J_n\longrightarrow J_0\ne0,
 \qquad
 A_n\longrightarrow A_0.
\label{eq:v2v87-data-convergence}
\end{equation}
In particular, for large \(n\),
\begin{equation}
 |J_n|\ge\frac12|J_0|,
\label{eq:v2v87-J-lower}
\end{equation}
and the delta carriers converge:
\begin{equation}
 2\pi^2\frac{A_n}{|J_n|}\delta_{k_n}
 \longrightarrow
 2\pi^2\frac{A_0}{|J_0|}\delta_0
\label{eq:v2v87-delta-convergence}
\end{equation}
in operator-valued distributions.
\end{proposition}

\begin{proof}
The derivative convergence in
Lemma~\ref{lem:v2v87-spectral} and the point convergence
\eqref{eq:v2v87-crossing-convergence} give \(J_n\to J_0\).
The projection and coefficient convergence give \(A_n\to A_0\).
The lower bound follows.  For a test function \(\varphi\),
the left side of
\eqref{eq:v2v87-delta-convergence} evaluates to
\[
 2\pi^2\frac{A_n}{|J_n|}\varphi(k_n),
\]
which converges to the stated continuum value.
\end{proof}

\subsection{Convergence of the external boundary pole}
\label{sec:v2v87-boundary}

Define
\begin{equation}
 C_n
 =
 \frac{A_n}{
 (\partial_\omega\lambda_{1,n})(k_n,\omega_n)
 (\partial_\omega\lambda_{2,n})(k_n,\omega_n)},
\label{eq:v2v87-Cn}
\end{equation}
and define \(C_0\) analogously.

\begin{lemma}[Stable pullback of a simple boundary value]
\label{lem:v2v87-pullback}
If \(g_n\to g\) in \(C^2\), each has one zero \(k_n\to0\), and their
derivatives are uniformly bounded away from zero, then
\begin{equation}
 \frac1{g_n(k)-i0}
 \longrightarrow
 \frac1{g(k)-i0}
\label{eq:v2v87-pullback}
\end{equation}
in \(\mathcal D'\) on a common neighbourhood.
\end{lemma}

\begin{proof}
Let \(\kappa_n\) be the local inverse of \(g_n\).  The uniform inverse
function theorem gives
\(\kappa_n\to\kappa\) in \(C^1\).  For a test function \(\varphi\),
change variables \(x=g_n(k)\):
\[
 \left\langle
 \frac1{g_n-i0},\varphi
 \right\rangle
 =
 \left\langle
 \frac1{x-i0},
 \frac{\varphi(\kappa_n(x))}
 {|g_n'(\kappa_n(x))|}
 \right\rangle.
\]
The transformed test functions converge in \(C^1\) with one common
compact support after a cutoff is fixed.  The distribution
\((x-i0)^{-1}\) is continuous on that test-function topology, proving
the claim.
\end{proof}

\begin{theorem}[Regulator convergence of the matrix pinch]
\label{thm:v2v87-pinch-convergence}
Let \(\mathcal M_n\) be the boundary-value limit of the regulated matrix
packet in V86 and \(\mathcal M\) its continuum counterpart.  Then
\begin{equation}
 \mathcal M_n
 =
 -\frac{2\pi i\,C_n}{g_n-i0}+S_n,
\label{eq:v2v87-Mn}
\end{equation}
where, after reducing the neighbourhood,
\begin{equation}
 C_n\longrightarrow C_0,
 \qquad
 S_n\longrightarrow S
 \quad\hbox{in }C^{m-2}.
\label{eq:v2v87-S-convergence}
\end{equation}
Consequently,
\begin{equation}
 \mathcal M_n\longrightarrow\mathcal M
\label{eq:v2v87-M-convergence}
\end{equation}
in operator-valued distributions.
\end{theorem}

\begin{proof}
Apply Theorem~\ref{thm:v2v86-pinch} at the shifted crossing
\((k_n,\omega_n)\).  The coefficient convergence follows from
\eqref{eq:v2v87-data-convergence},
\eqref{eq:v2v87-reg-slopes}, and evaluation at the convergent crossing
points.  Lemma~\ref{lem:v2v87-pullback} gives convergence of the singular
boundary term.

Every component of the smooth remainder in the V86 proof is obtained by
a fixed Riesz compression, a uniformly gapped inverse, a smooth change
of variables, or a one-pole pushforward.  These operations are continuous
under the \(C^m\) convergence and uniform support assumptions.  Two
derivatives are reserved for the changes of variables and
principal-value subtraction, giving the sufficient
\(C^{m-2}\) statement in
\eqref{eq:v2v87-S-convergence}.
\end{proof}

\subsection{Operator-level sufficient hypotheses}
\label{sec:v2v87-operator}

\begin{corollary}[Differentiated norm-resolvent input suffices]
\label{cor:v2v87-operator}
Suppose the regulated and continuum kinetic operators act on identified
Hilbert spaces and satisfy the V39 uniform norm-resolvent hypotheses on
fixed contours, with derivatives in \((k,\omega)\) through order \(m\).
Assume their compressed first spectral moments and the coefficient
families converge through the same order.  Then compatible Kato frames
may be chosen so that
\eqref{eq:v2v87-B-convergence} holds, and all conclusions of V87 follow.
\end{corollary}

\begin{proof}
Apply Proposition~\ref{prop:v2v39-operator-jets} in the two parameter
variables to both isolated clusters.  It gives \(C^m\) convergence of
the Kato-compressed matrices.  The remaining hypotheses and conclusions
are exactly those of
Lemma~\ref{lem:v2v87-spectral} through
Theorem~\ref{thm:v2v87-pinch-convergence}.
\end{proof}

\begin{firewall}[Identification maps are part of the theorem]
\label{fw:v2v87-identification}
If regulator spaces differ from the continuum Hilbert space, the
resolvents, projections, coefficient maps, and Sobolev scales must be
compared through specified quasi-unitary identification maps.  Pointwise
convergence of matrix entries in unrelated regulator bases has no
invariant meaning.
\end{firewall}

\subsection{Sharp spectral-pollution separator}
\label{sec:v2v87-pollution}

\begin{counterexample}[Strong resolvent convergence loses a crossing mode]
\label{sep:v2v87-strong}
On \(\ell^2(\mathbb N)\), let
\[
 De_j=j e_j.
\]
For \(k\) in a compact interval, define
\[
 D_n(k)e_j=
 \begin{cases}
  k e_n,&j=n,\\
  j e_j,&j\ne n.
 \end{cases}
\]
Each \(D_n(k)\) and \(D\) is self-adjoint with compact resolvent.  For
every fixed \(k\),
\begin{equation}
 D_n(k)\longrightarrow D
 \quad\hbox{in strong resolvent sense}.
\label{eq:v2v87-strong}
\end{equation}
Nevertheless, every \(D_n(k)\) has the rank-one crossing eigenvalue
\(k\) with projection
\[
 P_n=|e_n\rangle\langle e_n|,
\]
whereas \(D\) has no zero eigenvalue.  The projections have norm one and
do not converge in operator norm to a continuum zero-mode projection.
\end{counterexample}

\begin{proof}
For nonreal \(z\), the difference of the two resolvents acts only on
\(e_n\):
\[
 \bigl((D_n(k)-z)^{-1}-(D-z)^{-1}\bigr)x
 =
 \left(
 \frac1{k-z}-\frac1{n-z}
 \right)x_n e_n.
\]
The coefficient is uniformly bounded for fixed nonreal \(z\), while
\(x_n\to0\) for every \(x\in\ell^2\).  This proves strong resolvent
convergence.  The remaining statements follow from the displayed
diagonal spectra.
\end{proof}

\begin{firewall}[Strong convergence is too weak for the residue]
\label{fw:v2v87-strong}
A crossing eigenvector can escape to arbitrarily high regulator
coordinates while the resolvent converges strongly on every fixed
vector.  Norm-resolvent control on the physical contour, stable rank,
and compatible coefficient convergence are therefore essential for a
regulator-uniform pinch carrier.
\end{firewall}

\subsection{V87 conclusion and next acquisition}
\label{sec:v2v87-conclusion}

V87 proves regulator stability of the complete V86 data.  Differentiated
norm-resolvent convergence transfers the simple eigenprojections and
eigenvalues; uniform implicit-function estimates transfer the two roots
and their shifted common crossing; the joint Jacobian stays nonzero; and
both the principal-value and delta parts of the matrix pinch converge.
A compact-resolvent diagonal example proves that strong resolvent
convergence alone permits a spurious crossing mode.

V88 will add the BRST complex to this regulator limit.  Its target is a
commuting diagram between the regulated Kato identifications, the BRST
differentials, and the on-shell projections, sufficient to pass the V81
exact-sector cancellation and the physical harmonic residue to the
continuum.
\clearpage
\section*{Second Volume, Version V88: Regulator-Stable BRST Cancellation of the Matrix Pinch}
\addcontentsline{toc}{section}{Second Volume, Version V88: Regulator-Stable BRST Cancellation of the Matrix Pinch}
\label{sec:v2v88}

\subsection{Purpose}
\label{sec:v2v88-purpose}

V87 transfers the analytic crossing data to the continuum.  That result
does not by itself transfer the V81 exact-sector cancellation: the
regulated spectral projections, interaction vertex, source, and output
must form a compatible BRST diagram.

V88 gives sufficient exact and approximate diagram conditions.  Exact
chain maps make the V86 on-shell carrier vanish before division by the
joint Jacobian.  Quantified Ward defects give an explicit residue bound.
A uniform BRST Hodge gap transfers the harmonic projections and
contracting homotopies, so a closed source has a regulator-stable split
into its physical harmonic and cancelling exact parts.

\subsection{Two regulated cluster complexes}
\label{sec:v2v88-complexes}

For \(j=1,2\), let
\[
 E_{j,n}^\bullet
\]
be finite-dimensional graded Hilbert spaces with differentials
\begin{equation}
 s_{j,n}:E_{j,n}^g\longrightarrow E_{j,n}^{g+1},
 \qquad
 s_{j,n}^2=0.
\label{eq:v2v88-sn}
\end{equation}
After specified unitary identifications, regard all \(E_{j,n}\) as one
fixed \(E_j\), and assume
\begin{equation}
 s_{j,n}\longrightarrow s_j,
 \qquad
 s_{j,n}^*\longrightarrow s_j^*
\label{eq:v2v88-s-convergence}
\end{equation}
in operator norm.

Define
\begin{align}
 \Delta_{j,n}
 &=
 s_{j,n}s_{j,n}^*+s_{j,n}^*s_{j,n},
 \label{eq:v2v88-Delta-n}\\
 H_{j,n}
 &=
 \mathbf1_{\{0\}}(\Delta_{j,n}),
 \label{eq:v2v88-Hn}\\
 G_{j,n}^{\rm B}
 &=
 \Delta_{j,n}^{-1}(I-H_{j,n}),
 \label{eq:v2v88-GB}\\
 h_{j,n}
 &=
 s_{j,n}^*G_{j,n}^{\rm B}.
 \label{eq:v2v88-hn}
\end{align}
Here \(H_{j,n}\) denotes the BRST harmonic projection and is distinct
from the kinetic blocks of V79.

Assume the uniform Hodge gap
\begin{equation}
 \Delta_{j,n}
 \ge
 \gamma^2(I-H_{j,n})
\label{eq:v2v88-Hodge-gap}
\end{equation}
for one \(\gamma>0\), uniformly in \(j,n\).

\begin{theorem}[Convergence of the BRST Hodge data]
\label{thm:v2v88-Hodge-convergence}
The limiting differential satisfies \(s_j^2=0\), and its Laplacian
\(\Delta_j\) obeys the same nonzero spectral lower bound.  Moreover,
\begin{align}
 H_{j,n}&\longrightarrow H_j,
 \label{eq:v2v88-H-convergence}\\
 G_{j,n}^{\rm B}&\longrightarrow G_j^{\rm B},
 \label{eq:v2v88-GB-convergence}\\
 h_{j,n}&\longrightarrow h_j
 \label{eq:v2v88-h-convergence}
\end{align}
in operator norm, and
\begin{equation}
 s_jh_j+h_js_j=I-H_j,
 \qquad
 \|h_j\|\le\gamma^{-1}.
\label{eq:v2v88-limit-contraction}
\end{equation}
\end{theorem}

\begin{proof}
Norm convergence and \(s_{j,n}^2=0\) give \(s_j^2=0\).
Equation \eqref{eq:v2v88-s-convergence} gives
\(\Delta_{j,n}\to\Delta_j\) in norm.  The uniform gap separates zero
from the rest of every spectrum.  A fixed circle of radius less than
\(\gamma^2\) therefore gives the Riesz formula for all \(H_{j,n}\) and
\(H_j\), and the resolvent identity proves
\eqref{eq:v2v88-H-convergence}.

On the orthogonal complements, the function \(x\mapsto x^{-1}\) is
continuous on \([\gamma^2,\infty)\).  Finite-dimensional functional
calculus gives
\eqref{eq:v2v88-GB-convergence}; multiplication by
\(s_{j,n}^*\) gives
\eqref{eq:v2v88-h-convergence}.  The contracting identity and norm bound
pass from Theorem~\ref{thm:v2v68-contract} to the limit.
\end{proof}

\subsection{Kinetic spectral projections as chain maps}
\label{sec:v2v88-spectral-chain}

Let \(B_{j,n}\) be the regulated Hermitian kinetic blocks of V87 and let
\(\Pi_{j,n}\) denote their simple on-shell spectral projections at the
regulated joint crossing.

\begin{lemma}[Exact kinetic invariance descends to the Riesz projection]
\label{lem:v2v88-P-chain}
If
\begin{equation}
 [B_{j,n},s_{j,n}]=0,
\label{eq:v2v88-B-chain}
\end{equation}
then
\begin{equation}
 [\Pi_{j,n},s_{j,n}]=0.
\label{eq:v2v88-P-chain}
\end{equation}
\end{lemma}

\begin{proof}
For \(z\) on the Riesz contour,
\[
 [(z-B_{j,n})^{-1},s_{j,n}]
 =
 (z-B_{j,n})^{-1}
 [B_{j,n},s_{j,n}]
 (z-B_{j,n})^{-1}
 =
 0.
\]
Integrate this identity around the contour defining \(\Pi_{j,n}\).
\end{proof}

\begin{lemma}[Quantitative projection Ward defect]
\label{lem:v2v88-P-defect}
Suppose the Riesz contour has length \(\ell_j\) and the resolvent norm on
it is at most \(r_j\).  Then
\begin{equation}
 \|[\Pi_{j,n},s_{j,n}]\|
 \le
 \frac{\ell_jr_j^2}{2\pi}
 \|[B_{j,n},s_{j,n}]\|.
\label{eq:v2v88-P-defect}
\end{equation}
\end{lemma}

\begin{proof}
Use the same resolvent commutator identity without setting its
right-hand side to zero, take norms, and integrate around the contour.
\end{proof}

Thus differentiated norm-resolvent convergence must be accompanied by a
Ward-compatible cluster contour; analytic rank stability alone does not
make the spectral subspace a BRST subcomplex.

\subsection{The exact two-line chain diagram}
\label{sec:v2v88-diagram}

At the regulated joint crossing, let
\begin{align*}
 C_n&\in\mathcal L(X,E_{2,n}^{g-1}),\\
 \mathcal R_n&\in\mathcal L(X,E_{2,n}^{g}),\\
 V_n&\in\mathcal L(E_{2,n}^{g},E_{1,n}^{g}),\\
 L_n&\in\mathcal L(E_{1,n}^{g},Y).
\end{align*}
Assume the source entering the second kinetic line is
\begin{equation}
 \mathcal R_n=s_{2,n}C_n.
\label{eq:v2v88-exact-source}
\end{equation}
The exact chain conditions are
\begin{align}
 L_ns_{1,n}&=0,
 \label{eq:v2v88-physical-output}\\
 V_ns_{2,n}&=s_{1,n}V_n,
 \label{eq:v2v88-vertex-chain}\\
 [B_{j,n},s_{j,n}]&=0.
 \label{eq:v2v88-two-kinetic-chain}
\end{align}

\begin{theorem}[Exact regulated pinch cancellation]
\label{thm:v2v88-exact-cancellation}
Under
\eqref{eq:v2v88-exact-source}--\eqref{eq:v2v88-two-kinetic-chain},
the complete V86 on-shell carrier vanishes:
\begin{equation}
 A_n
 =
 L_n\Pi_{1,n}V_n\Pi_{2,n}\mathcal R_n
 =
 0.
\label{eq:v2v88-An-zero}
\end{equation}
Consequently both the principal-value and delta residues of the regulated
matrix pinch vanish, independently of the crossing Jacobian.
\end{theorem}

\begin{proof}
Lemma~\ref{lem:v2v88-P-chain} lets the differentials pass through the
two spectral projections.  Therefore
\begin{align*}
 L_n\Pi_{1,n}V_n\Pi_{2,n}s_{2,n}C_n
 &=
 L_n\Pi_{1,n}V_ns_{2,n}\Pi_{2,n}C_n\\
 &=
 L_n\Pi_{1,n}s_{1,n}V_n\Pi_{2,n}C_n\\
 &=
 L_ns_{1,n}\Pi_{1,n}V_n\Pi_{2,n}C_n
 =0.
\end{align*}
The V86 residue is linear in \(A_n\), proving the last statement.
\end{proof}

\begin{corollary}[Continuum exact cancellation]
\label{cor:v2v88-limit-zero}
Assume the V87 convergence hypotheses and
\eqref{eq:v2v88-s-convergence}.  If the exact chain identities hold for
every \(n\) and the coefficient families converge, then the limiting
on-shell carrier is zero and the continuum pinched pushforward is smooth
in the exact sector.
\end{corollary}

\begin{proof}
Equation \eqref{eq:v2v88-An-zero} gives \(A_n=0\).  V87 gives
\(A_n\to A_0\), hence \(A_0=0\).  Apply
Corollary~\ref{cor:v2v86-Ward}.
\end{proof}

\subsection{A quantitative approximate diagram}
\label{sec:v2v88-approximate}

Define the defects
\begin{align}
 \rho_n
 &=
 \mathcal R_n-s_{2,n}C_n,
 \label{eq:v2v88-rho}\\
 e_{L,n}
 &=
 L_ns_{1,n},
 \label{eq:v2v88-eL}\\
 E_{V,n}
 &=
 V_ns_{2,n}-s_{1,n}V_n,
 \label{eq:v2v88-EV}\\
 E_{j,n}
 &=
 [B_{j,n},s_{j,n}].
 \label{eq:v2v88-Ej}
\end{align}

\begin{theorem}[Explicit residue-defect bound]
\label{thm:v2v88-defect}
Let
\[
 p_{j,n}
 =
 \|[\Pi_{j,n},s_{j,n}]\|.
\]
Then
\begin{align}
 \|A_n\|
 \le{}&
 \|e_{L,n}\|\,\|\Pi_{1,n}\|\,\|V_n\|\,
 \|\Pi_{2,n}\|\,\|C_n\|
 \nonumber\\
 &+
 \|L_n\|\,p_{1,n}\,\|V_n\|\,
 \|\Pi_{2,n}\|\,\|C_n\|
 \nonumber\\
 &+
 \|L_n\|\,\|\Pi_{1,n}\|\,\|E_{V,n}\|\,
 \|\Pi_{2,n}\|\,\|C_n\|
 \nonumber\\
 &+
 \|L_n\|\,\|\Pi_{1,n}\|\,\|V_n\|\,
 p_{2,n}\,\|C_n\|
 \nonumber\\
 &+
 \|L_n\|\,\|\Pi_{1,n}\|\,\|V_n\|\,
 \|\Pi_{2,n}\|\,\|\rho_n\|.
\label{eq:v2v88-defect-bound}
\end{align}
Using Lemma~\ref{lem:v2v88-P-defect}, the two \(p_{j,n}\) terms are
bounded by the kinetic Ward defects \(E_{j,n}\) times the squared contour
resolvent constants.

If all coefficient norms and \(\|C_n\|\) are uniform and
\begin{equation}
 \|e_{L,n}\|+\|E_{V,n}\|+\|E_{1,n}\|
 +\|E_{2,n}\|+\|\rho_n\|
 \longrightarrow0,
\label{eq:v2v88-defects-zero}
\end{equation}
then \(A_n\to0\).  Under the V87 Jacobian lower bound, the complete pinch
residue tends to zero in the test-function topology.
\end{theorem}

\begin{proof}
Insert
\(\mathcal R_n=s_{2,n}C_n+\rho_n\).  Commute \(s_{2,n}\) successively
to the left:
\begin{align*}
 \Pi_{2,n}s_{2,n}
 &=
 s_{2,n}\Pi_{2,n}
 +[\Pi_{2,n},s_{2,n}],\\
 V_ns_{2,n}
 &=
 s_{1,n}V_n+E_{V,n},\\
 \Pi_{1,n}s_{1,n}
 &=
 s_{1,n}\Pi_{1,n}
 +[\Pi_{1,n},s_{1,n}].
\end{align*}
The term in which the differential reaches the output is
\(e_{L,n}\Pi_{1,n}V_n\Pi_{2,n}C_n\).  The three commutations produce the
next three defect terms, and \(\rho_n\) produces the last.  Take norms
only after this exact decomposition to obtain
\eqref{eq:v2v88-defect-bound}.

Lemma~\ref{lem:v2v88-P-defect} gives the kinetic estimates.  Under
\eqref{eq:v2v88-defects-zero}, the right side tends to zero.
Proposition~\ref{prop:v2v87-residue} bounds the inverse Jacobians and
V87 controls the principal-value pullbacks, so the residue tends to zero
as a distribution.
\end{proof}

\begin{assessment}[Rate required for the continuum residue]
\label{ass:v2v88-rate}
Distributional cancellation requires only \(A_n\to0\) when the Jacobian
stays uniformly nonzero.  A stronger uniform pointwise bound across the
finite-regulator crossing would require exact cancellation at each
\(n\), because every nonzero residue still multiplies an unbounded
boundary-value kernel.  The topology of the desired conclusion must
therefore be stated.
\end{assessment}

\subsection{Closed sources and stable harmonic residues}
\label{sec:v2v88-harmonic}

Suppose
\begin{equation}
 s_{2,n}\mathcal R_n=0.
\label{eq:v2v88-closed-source}
\end{equation}
The Hodge identity gives
\begin{equation}
 \mathcal R_n
 =
 H_{2,n}\mathcal R_n
 +
 s_{2,n}h_{2,n}\mathcal R_n.
\label{eq:v2v88-Hodge-split}
\end{equation}

\begin{theorem}[Only the regulator-stable harmonic carrier survives]
\label{thm:v2v88-harmonic}
Under the exact chain conditions
\eqref{eq:v2v88-physical-output}--\eqref{eq:v2v88-two-kinetic-chain},
\begin{equation}
 L_n\Pi_{1,n}V_n\Pi_{2,n}\mathcal R_n
 =
 L_n\Pi_{1,n}V_n\Pi_{2,n}
 H_{2,n}\mathcal R_n.
\label{eq:v2v88-harmonic-carrier}
\end{equation}
If \(\mathcal R_n\to\mathcal R\), then
\begin{equation}
 H_{2,n}\mathcal R_n
 \longrightarrow
 H_2\mathcal R.
\label{eq:v2v88-harmonic-source-convergence}
\end{equation}
Consequently, under V87, the physical delta carrier converges to
\begin{equation}
 2\pi^2
 \frac{
 L_0\Pi_{1,0}V_0\Pi_{2,0}H_2\mathcal R
 }{|J_0|}
 \delta_0.
\label{eq:v2v88-harmonic-delta}
\end{equation}
\end{theorem}

\begin{proof}
Apply Theorem~\ref{thm:v2v88-exact-cancellation} to the exact second
term in \eqref{eq:v2v88-Hodge-split}.  This proves
\eqref{eq:v2v88-harmonic-carrier}.  The convergence
\eqref{eq:v2v88-H-convergence} and source convergence give
\eqref{eq:v2v88-harmonic-source-convergence}.  Combine this with the
coefficient, projection, and Jacobian convergence of V87.
\end{proof}

\begin{assessment}[Two independent gaps]
\label{ass:v2v88-two-gaps}
The kinetic complementary gap of V87 stabilizes the on-shell spectral
projections and Jacobian.  The BRST Hodge gap of
\eqref{eq:v2v88-Hodge-gap} stabilizes the exact/harmonic decomposition.
Neither gap implies the other, and both constants remain in a regulator-
uniform construction.
\end{assessment}

\subsection{Failure modes}
\label{sec:v2v88-failures}

\begin{firewall}[An analytic spectral projection need not be a chain map]
\label{fw:v2v88-analytic-only}
Norm-resolvent convergence can produce a perfectly stable rank-one
projection while \([B_n,s_n]\ne0\).  Lemma
\ref{lem:v2v88-P-defect} shows that the resulting projection Ward defect
is controlled only after the kinetic commutator is controlled on the
Riesz contour.
\end{firewall}

\begin{firewall}[A convergent quotient is not a convergent primitive]
\label{fw:v2v88-primitive}
Even if every exact source has zero physical class, convergence of the
representatives \(h_{j,n}\mathcal R_n\) uses the uniform Hodge gap.
The collapsing-gap counterexample in V68 still applies.
\end{firewall}

\begin{counterexample}[A vertex defect reintroduces the physical carrier]
\label{sep:v2v88-vertex}
Let \(s_2u=v\), \(s_1\widetilde u=\widetilde v\), and let the physical
output annihilate \(\widetilde v\) but not a harmonic vector \(w\).
If
\[
 Vu=\widetilde u,
 \qquad
 Vv=\widetilde v+\eta w,
 \qquad \eta\ne0,
\]
then
\[
 (Vs_2-s_1V)u=\eta w,
\]
and an exact input can produce a nonzero on-shell matrix element of size
\(|\eta|\).  Thus separate invariance of the two kinetic blocks does not
replace the vertex chain identity.
\end{counterexample}

\begin{proof}
Apply the displayed maps to \(u\).  The \(\widetilde v\) component is a
boundary and is killed by the output, while the harmonic component
survives.
\end{proof}

\begin{firewall}[An anomaly is not a small analytic remainder]
\label{fw:v2v88-anomaly}
If the limiting Ward defect represents a nonzero local BRST cohomology
class, no choice of convergent Kato frame removes it.  One must classify
the defect and cancel it by an allowed counterterm when possible.
V88 gives the analytic consequence of a vanishing defect; it does not
prove anomaly cancellation for the complete curved Standard Model
regulator.
\end{firewall}

\subsection{V88 conclusion and next acquisition}
\label{sec:v2v88-conclusion}

V88 supplies the commuting BRST diagram missing from V87.  Exact kinetic,
vertex, source, and output chain identities annihilate the complete
matrix pinch carrier at every regulator.  Approximate identities yield an
explicit residue bound, with kinetic Ward defects amplified by the
squared Riesz-resolvent constants.  A uniform Hodge gap transfers the
harmonic projection and contracting homotopy, leaving only the physical
harmonic delta carrier in the continuum.

V89 will move upstream to the Ward defect itself.  It will formulate a
local relative-cohomology theorem which separates removable regulator
defects from a genuine anomaly class and records the uniform counterterm
bounds required to preserve the V87 spectral convergence.

\clearpage
\section*{Second Volume, Version V89: Relative Ward Counterterms Compatible with Spectral Convergence}
\addcontentsline{toc}{section}{Second Volume, Version V89: Relative Ward Counterterms Compatible with Spectral Convergence}
\label{sec:v2v89}

\subsection{Purpose and relation to V66}
\label{sec:v2v89-purpose}

V66 classified marginal fifth-derivative words modulo BRST boundaries and
contact-local terms.  V88 now produces a different but related object:
the finite tuple of kinetic, vertex, source, and output Ward defects which
controls the on-shell pinch residue.

V89 formulates the local counterterm problem for this tuple.  The new
point is quantitative and analytic.  Vanishing of the relative
cohomology class must be accompanied by a uniformly bounded right
inverse for the counterterm variation.  Only then does a small defect
produce a small counterterm, and only a small or convergent counterterm
preserves the norm-resolvent limit used by V87.

\subsection{Finite relative defect complex}
\label{sec:v2v89-complex}

At regulator \(n\), let
\begin{equation}
 \mathcal C_n^0
 \xrightarrow{\ \delta_n\ }
 \mathcal C_n^1
 \xrightarrow{\ d_n\ }
 \mathcal C_n^2,
 \qquad
 d_n\delta_n=0,
\label{eq:v2v89-complex}
\end{equation}
be a finite-dimensional Hilbert complex.  Elements of
\(\mathcal C_n^0\) are allowed local counterterms.  Elements of
\(\mathcal C_n^1\) are Ward-defect tuples, including the five components
in \eqref{eq:v2v88-rho}--\eqref{eq:v2v88-Ej}.  The equation
\begin{equation}
 d_nw_n=0
\label{eq:v2v89-consistency}
\end{equation}
is the linearized Wess--Zumino consistency condition.

Let
\[
 \mathcal L_n^1\subset\ker d_n
\]
be the subspace of permitted contact-local defects which vanish under
the physical on-shell evaluation
\begin{equation}
 \mathcal E_{{\rm os},n}:
 \ker d_n\longrightarrow\mathcal O_n,
 \qquad
 \mathcal L_n^1\subset\ker\mathcal E_{{\rm os},n}.
\label{eq:v2v89-local-null}
\end{equation}
Let
\[
 q_n:\ker d_n\longrightarrow
 \overline{\mathcal C}_n^1
 :=
 \ker d_n/\mathcal L_n^1
\]
be the quotient map, and define
\begin{equation}
 \overline\delta_n=q_n\delta_n.
\label{eq:v2v89-deltabar}
\end{equation}
The relative obstruction space is
\begin{equation}
 H_{{\rm rel},n}^1
 =
 \overline{\mathcal C}_n^1/
 \operatorname{im}\overline\delta_n.
\label{eq:v2v89-Hrel}
\end{equation}

\begin{theorem}[Exact relative removability criterion]
\label{thm:v2v89-removability}
A consistent defect \(w_n\) can be removed from the physical on-shell
carrier by an allowed counterterm if and only if
\begin{equation}
 [q_nw_n]=0
 \quad\hbox{in }H_{{\rm rel},n}^1.
\label{eq:v2v89-class-zero}
\end{equation}
More explicitly, this condition is equivalent to the existence of
\(c_n\in\mathcal C_n^0\) and
\(\ell_n\in\mathcal L_n^1\) such that
\begin{equation}
 w_n+\delta_nc_n=\ell_n.
\label{eq:v2v89-counterterm-equation}
\end{equation}
For such a correction,
\begin{equation}
 \mathcal E_{{\rm os},n}
 (w_n+\delta_nc_n)=0.
\label{eq:v2v89-onshell-zero}
\end{equation}
\end{theorem}

\begin{proof}
The class vanishes exactly when
\(q_nw_n\in\operatorname{im}\overline\delta_n\).  Thus there is a
\(c_n\) with
\[
 q_n(w_n+\delta_nc_n)=0,
\]
which is equivalent to
\(w_n+\delta_nc_n\in\mathcal L_n^1\).  This proves
\eqref{eq:v2v89-counterterm-equation}; the reverse implication is
immediate.  Apply
\eqref{eq:v2v89-local-null} to obtain
\eqref{eq:v2v89-onshell-zero}.
\end{proof}

\begin{firewall}[A nonzero relative class is an obstruction]
\label{fw:v2v89-anomaly}
If \([q_nw_n]\ne0\), no counterterm in the declared allowed space solves
\eqref{eq:v2v89-counterterm-equation}.  Enlarging
\(\mathcal C_n^0\) by a symmetry-breaking or nonlocal operator changes
the theory and is not a proof of anomaly removal.
\end{firewall}

\subsection{Quantitative counterterm selection}
\label{sec:v2v89-quantitative}

Give the quotient
\(\overline{\mathcal C}_n^1\) its Hilbert quotient norm.  Let
\(\overline\delta_n^\dagger\) be the Moore--Penrose inverse of
\(\overline\delta_n\) on its image, and let
\begin{equation}
 \sigma_n
 =
 \min\{
 \sigma>0:\sigma\text{ is a nonzero singular value of }
 \overline\delta_n
 \}.
\label{eq:v2v89-sigma}
\end{equation}
When the image is zero, the statements below are read only for the zero
defect class.

\begin{theorem}[Least-norm relative counterterm]
\label{thm:v2v89-least}
Assume \eqref{eq:v2v89-class-zero}.  The choice
\begin{equation}
 c_n
 =
 -\overline\delta_n^\dagger q_nw_n
\label{eq:v2v89-cn}
\end{equation}
has least norm among counterterms orthogonal to
\(\ker\overline\delta_n\) which solve the quotient equation, and
\begin{equation}
 \|c_n\|
 \le
 \sigma_n^{-1}\|q_nw_n\|.
\label{eq:v2v89-counterterm-bound}
\end{equation}
If
\begin{equation}
 \sigma_n\ge\sigma_0>0,
 \qquad
 \|q_nw_n\|\le\eta_n,
 \qquad
 \eta_n\longrightarrow0,
\label{eq:v2v89-uniform-gap}
\end{equation}
then \(c_n\to0\) with
\begin{equation}
 \|c_n\|\le\sigma_0^{-1}\eta_n.
\label{eq:v2v89-cn-small}
\end{equation}
\end{theorem}

\begin{proof}
The Moore--Penrose inverse maps
\(\operatorname{im}\overline\delta_n\) to
\((\ker\overline\delta_n)^\perp\) and gives the unique least-norm
solution there.  Its operator norm is the reciprocal of the smallest
nonzero singular value.  This proves
\eqref{eq:v2v89-counterterm-bound}; the uniform statement follows from
\eqref{eq:v2v89-uniform-gap}.
\end{proof}

\begin{assessment}[Cohomological zero and a uniform primitive are different]
\label{ass:v2v89-two-levels}
Equation \eqref{eq:v2v89-class-zero} is the algebraic condition.
The lower bound on \(\sigma_n\) is the quantitative condition.  The
first gives a counterterm at each regulator; the second prevents its
coefficient from diverging or remaining macroscopic while the defect
vanishes.
\end{assessment}

\begin{counterexample}[Vanishing exact defects with divergent counterterms]
\label{sep:v2v89-collapse}
Let
\[
 \mathcal C_n^0=\mathbb C u,
 \qquad
 \mathcal C_n^1=\mathbb C v,
 \qquad
 \delta_nu=\varepsilon_nv,
\]
where \(\varepsilon_n\downarrow0\), and take no local-null subspace.
Set
\[
 w_n=\sqrt{\varepsilon_n}\,v.
\]
Then \(w_n\to0\) and its relative class is zero for every \(n\), but the
unique counterterm solving
\(w_n+\delta_nc_n=0\) is
\begin{equation}
 c_n=-\varepsilon_n^{-1/2}u,
\label{eq:v2v89-divergent-cn}
\end{equation}
whose norm diverges.
\end{counterexample}

\begin{proof}
The equation
\(\sqrt{\varepsilon_n}+\varepsilon_n c=0\) gives
\(c=-\varepsilon_n^{-1/2}\).  Here
\(\sigma_n=\varepsilon_n\), so
\eqref{eq:v2v89-counterterm-bound} is sharp.
\end{proof}

\subsection{Insertion into the kinetic and vertex data}
\label{sec:v2v89-insertion}

Let
\begin{equation}
 \mathfrak I_n:
 \mathcal C_n^0\longrightarrow\mathfrak A_n
\label{eq:v2v89-insertion}
\end{equation}
insert a counterterm into the tuple of kinetic blocks, vertices, sources,
and outputs.  Assume, on the V87 parameter rectangle and for
\(0\le|\alpha|\le m\),
\begin{equation}
 \|\partial^\alpha\mathfrak I_n(c)\|_{\mathfrak A_n}
 \le
 C_{\rm ins}\|c\|.
\label{eq:v2v89-insertion-bound}
\end{equation}
For a kinetic component write
\[
 T_n=\mathfrak I_{B,n}(c_n),
 \qquad
 \widetilde B_n=B_n+T_n.
\]

\begin{lemma}[Small relative perturbations preserve a contour]
\label{lem:v2v89-contour}
Let \(\Gamma\) be a compact contour on which
\[
 \sup_{n,z\in\Gamma}
 \|(B_n-z)^{-1}\|_{\mathcal H^0\to\mathcal H^1}
 \le C_R.
\]
If
\begin{equation}
 \|T_n\|_{\mathcal H^1\to\mathcal H^0}C_R<1,
\label{eq:v2v89-Neumann}
\end{equation}
then \(\Gamma\subset\rho(\widetilde B_n)\) and
\begin{equation}
 (\widetilde B_n-z)^{-1}
 =
 (B_n-z)^{-1}
 \bigl[I+T_n(B_n-z)^{-1}\bigr]^{-1}.
\label{eq:v2v89-corrected-resolvent}
\end{equation}
Moreover,
\begin{equation}
 \|(\widetilde B_n-z)^{-1}-(B_n-z)^{-1}\|
 \le
 \frac{C_R^2\|T_n\|}
 {1-C_R\|T_n\|}.
\label{eq:v2v89-resolvent-difference}
\end{equation}
\end{lemma}

\begin{proof}
Factor
\[
 \widetilde B_n-z
 =
 \bigl[I+T_n(B_n-z)^{-1}\bigr](B_n-z).
\]
The bracket is invertible by the Neumann series under
\eqref{eq:v2v89-Neumann}, proving
\eqref{eq:v2v89-corrected-resolvent}.  Subtract the unperturbed inverse
and use the resolvent identity to obtain
\eqref{eq:v2v89-resolvent-difference}.
\end{proof}

\begin{theorem}[Counterterm-compatible spectral convergence]
\label{thm:v2v89-spectral}
Assume the V87 differentiated norm-resolvent hypotheses for \(B_n\), the
insertion estimate \eqref{eq:v2v89-insertion-bound}, and the uniform
counterterm bound
\eqref{eq:v2v89-cn-small}.  Then
\begin{equation}
 \partial^\alpha
 \left[
 (\widetilde B_n-z)^{-1}-(B_n-z)^{-1}
 \right]
 \longrightarrow0
\label{eq:v2v89-differentiated-resolvent}
\end{equation}
uniformly on the physical Riesz contours for
\(0\le|\alpha|\le m\).  Hence the corrected family has the same
continuum kinetic limit and satisfies all V87 conclusions for
sufficiently large \(n\).

More generally, if the inserted corrections converge in the declared
operator topology to a finite \(T\), the corrected family converges to
the correspondingly renormalized continuum operator \(B+T\), provided
the limiting contours and transversality remain separated.
\end{theorem}

\begin{proof}
Equation \eqref{eq:v2v89-cn-small} and
\eqref{eq:v2v89-insertion-bound} give
\(T_n\to0\) with all parameter derivatives through order \(m\).
Lemma~\ref{lem:v2v89-contour} applies uniformly for large \(n\).
Differentiate
\eqref{eq:v2v89-corrected-resolvent}.  Every difference term contains
at least one derivative of \(T_n\), while all resolvent and Neumann
factors are uniformly bounded.  The finite Leibniz sums therefore tend
to zero, proving
\eqref{eq:v2v89-differentiated-resolvent}.  Apply the V39 Riesz transfer
and V87 crossing theorem.

For a nonzero limiting correction, apply the same resolvent identity to
\(T_n-T\) about the corrected limit \(B+T\).  The explicit contour and
transversality conditions are needed because a finite correction can
move or merge spectral clusters.
\end{proof}

\begin{firewall}[A small coefficient needs the correct operator topology]
\label{fw:v2v89-topology}
The Euclidean norm of the finite counterterm coefficient does not by
itself control an unbounded kinetic insertion.  The map
\(\mathfrak I_n\) must satisfy
\eqref{eq:v2v89-insertion-bound} from the graph domain to the base
space, with differentiated constants uniform in the regulator.
\end{firewall}

\subsection{Effect on the V88 pinch residue}
\label{sec:v2v89-residue}

Let \(\widetilde w_n=w_n+\delta_nc_n\) be the corrected defect.

\begin{corollary}[Relative cancellation restores the on-shell diagram]
\label{cor:v2v89-residue}
Under
Theorem~\ref{thm:v2v89-removability},
\[
 \widetilde w_n\in\mathcal L_n^1
 \subset\ker\mathcal E_{{\rm os},n}.
\]
If the physical on-shell evaluation is the defect combination appearing
in Theorem~\ref{thm:v2v88-defect}, the corrected V86 pinch residue
vanishes.  Under the hypotheses of
Theorem~\ref{thm:v2v89-spectral}, this cancellation and the V87
continuum pinch limit are compatible.
\end{corollary}

\begin{proof}
The first statement is
\eqref{eq:v2v89-counterterm-equation} and
\eqref{eq:v2v89-local-null}.  The V88 defect decomposition is exact
before norms, so zero physical evaluation kills the residue carrier.
Theorem~\ref{thm:v2v89-spectral} preserves the spectral projections and
crossing data needed to pass to the continuum.
\end{proof}

\begin{assessment}[The order of construction]
\label{ass:v2v89-order}
The valid order is:
\begin{enumerate}
 \item form the complete consistent Ward defect;
 \item compute its relative cohomology class;
 \item choose a bounded least-norm counterterm when the class vanishes;
 \item verify its graph-operator insertion norm;
 \item only then invoke norm-resolvent and pinch-residue convergence.
\end{enumerate}
Selecting diagramwise counterterms before the complete defect is formed
can destroy both the cohomological cancellation and the spectral
estimate.
\end{assessment}

\subsection{What remains model-specific}
\label{sec:v2v89-model}

\begin{firewall}[Standard Model anomaly cancellation is still an input]
\label{fw:v2v89-SM-anomaly}
V89 does not compute the curved Standard Model anomaly polynomial or
prove cancellation among the actual fermion representations.  It proves
the consequence of a zero relative class in the declared regulator
complex.  The representation-theoretic trace identities, boundary
contributions, global anomalies, and gravitational mixed terms must be
verified in that complex.
\end{firewall}

\begin{firewall}[Local-null terms must really vanish on shell]
\label{fw:v2v89-local-null}
A polynomial contact term belongs to \(\mathcal L_n^1\) here only if the
physical on-shell evaluation kills it in the V86 carrier.  Locality alone
does not imply this; a local kinetic counterterm can move the crossing
or change its residue and must instead be included in the renormalized
spectral data.
\end{firewall}

\begin{firewall}[Finite-dimensional closure is regulator-local]
\label{fw:v2v89-finite}
The defect spaces are finite at each regulator.  Uniform bounds require
stable dimensions, compatible bases, and a lower singular-value bound
for \(\overline\delta_n\).  An increasing family of local monomials can
lose this bound even when every finite problem is exact.
\end{firewall}

\subsection{V89 conclusion and next acquisition}
\label{sec:v2v89-conclusion}

V89 separates algebraic anomaly removal from quantitative counterterm
control.  A Ward defect is removable precisely when its relative class
vanishes.  The Moore--Penrose counterterm is uniformly small only when
the counterterm variation has a regulator-uniform nonzero singular
value.  Under a uniform graph-operator insertion estimate, such a
counterterm preserves the differentiated norm-resolvent convergence and
therefore the V87 matrix pinch limit.

V90 is the next mandatory companion audit.  It will inspect the newest
Yang--Mills and Navier--Stokes notes before deciding whether the next SM
step should compute the actual anomaly-class interface, strengthen the
uniform local-complex bound, or return upstream to the Einstein-coupled
boundary regulator.
\clearpage
\section*{Second Volume, Version V90: Companion Audit and Stable Projection onto the Anomaly Cokernel}
\addcontentsline{toc}{section}{Second Volume, Version V90: Companion Audit and Stable Projection onto the Anomaly Cokernel}
\label{sec:v2v90}

\subsection{Mandatory companion audit}
\label{sec:v2v90-audit}

The current companion inventory was inspected without modification:
\begin{center}
\small
\begin{tabular}{p{0.48\linewidth}r p{0.32\linewidth}}
\toprule
file&bytes&SHA--256\\
\midrule
\texttt{Renewal\_Yang\_Mills\_Mass\_Gap\_Research\_Notes\_V70.tex}
&999512&
\texttt{19e6931c1a9a39d0b640ab668b807ab1083781a3540744c523d69ca2eb000871}
\\
\texttt{Renewal\_Yang\_Mills\_Mass\_Gap\_Research\_Notes\_V71.tex}
&999512&
\texttt{19e6931c1a9a39d0b640ab668b807ab1083781a3540744c523d69ca2eb000871}
\\
\texttt{Renewal\_NS\_Stability\_Research\_Notes\_V31.tex}
&887537&
\texttt{f1121b62238ad5491bb7cf03635ea9934942edf573ed8faaa9ee79e1d38a6696}
\\
\bottomrule
\end{tabular}
\end{center}
The two Yang--Mills files are byte-identical.  Both internally end with
the compact V70 record; hence V71 contains no additional mathematical
version.  They were left untouched.  The NS file is unchanged.

\begin{assessment}[Yang--Mills V70 transfer]
\label{ass:v2v90-YM}
Yang--Mills V70 first resolves the complete physical bank into protected
and gapped output sectors, then derives one one-use transfer certificate,
and only afterward studies the finite routing problem.  It expressly
forbids expanding the spectral split into separately normed words.

For the V89 Ward complex, the corresponding operation is the orthogonal
split of the complete consistent defect into the range of the
counterterm variation and its anomaly cokernel.  Counterterms are chosen
only on the first projection.  No Yang--Mills bank estimate or Casimir
currency is imported.
\end{assessment}

\begin{assessment}[Navier--Stokes V31 transfer]
\label{ass:v2v90-NS}
Navier--Stokes V31 still shows that a formally available mark can cost
the complete critical stock when its source correlation is discarded.
For the local Ward complex this reinforces the need to control the
smallest nonzero singular value of the complete counterterm map.  Solving
each finite problem without that control can hide a divergent
counterterm.
\end{assessment}

The audit selects a stable range/cokernel theorem before any
coefficientwise anomaly calculation.

\subsection{The removable and anomalous projections}
\label{sec:v2v90-projections}

After quotienting by the local-null space of V89, let
\[
 \delta_n:X_n\longrightarrow Y_n
\]
be the regulated counterterm variation.  Use specified unitary
identifications to regard \(X_n=X\) and \(Y_n=Y\), and assume
\begin{equation}
 \delta_n\longrightarrow\delta
 \quad\hbox{in operator norm}.
\label{eq:v2v90-delta-convergence}
\end{equation}
Assume there is \(\sigma_0>0\) such that every singular value of every
\(\delta_n\) is either zero or at least \(\sigma_0\):
\begin{equation}
 \operatorname{spec}(\delta_n\delta_n^*)
 \subset
 \{0\}\cup[\sigma_0^2,\infty).
\label{eq:v2v90-singular-gap}
\end{equation}

Let
\begin{align}
 P_n&=\delta_n\delta_n^\dagger,
 \label{eq:v2v90-Pn}\\
 A_n&=I-P_n.
 \label{eq:v2v90-An}
\end{align}
Thus \(P_n\) is the orthogonal projection onto
\(\operatorname{im}\delta_n\), while \(A_n\) projects onto
\[
 \ker\delta_n^*
 =
 (\operatorname{im}\delta_n)^\perp,
\]
the finite anomaly cokernel.

\begin{theorem}[Stable removable/anomalous split]
\label{thm:v2v90-stable-split}
Under
\eqref{eq:v2v90-delta-convergence}--\eqref{eq:v2v90-singular-gap},
the limit satisfies
\begin{equation}
 \operatorname{spec}(\delta\delta^*)
 \subset
 \{0\}\cup[\sigma_0^2,\infty),
\label{eq:v2v90-limit-gap}
\end{equation}
and
\begin{align}
 \delta_n^\dagger&\longrightarrow\delta^\dagger,
 \label{eq:v2v90-pseudoinverse-convergence}\\
 P_n&\longrightarrow P:=\delta\delta^\dagger,
 \label{eq:v2v90-P-convergence}\\
 A_n&\longrightarrow A:=I-P
 \label{eq:v2v90-A-convergence}
\end{align}
in operator norm.  In particular, the removable rank and anomaly
dimension are constant for all sufficiently large \(n\).
\end{theorem}

\begin{proof}
Equation \eqref{eq:v2v90-delta-convergence} gives
\[
 \delta_n\delta_n^*\longrightarrow\delta\delta^*
\]
in norm.  Choose a circle about zero of radius smaller than
\(\sigma_0^2\).  The Riesz projections onto the zero eigenspaces converge
in norm by the resolvent identity.  Their complements are \(P_n\), so
\eqref{eq:v2v90-P-convergence}--\eqref{eq:v2v90-A-convergence} follow,
and projections at distance below one have the same rank.

On the common spectral set
\(\{0\}\cup[\sigma_0^2,\infty)\), the function equal to zero at zero and
to \(x^{-1}\) on the positive component is continuous.  Functional
calculus gives
\[
 (\delta_n\delta_n^*)^\dagger
 \longrightarrow
 (\delta\delta^*)^\dagger.
\]
Use
\[
 \delta_n^\dagger
 =
 \delta_n^*(\delta_n\delta_n^*)^\dagger
\]
to obtain
\eqref{eq:v2v90-pseudoinverse-convergence}.  The same argument gives the
limit gap.
\end{proof}

\subsection{Canonical complete-defect reduction}
\label{sec:v2v90-reduction}

Let
\[
 y_n=q_nw_n\in Y
\]
be the complete consistent defect after the V89 local-null quotient.
Define
\begin{align}
 y_{n,\rm rem}&=P_ny_n,
 \label{eq:v2v90-yrem}\\
 \mathfrak a_n&=A_ny_n.
 \label{eq:v2v90-anomaly}
\end{align}

\begin{theorem}[Canonical counterterm leaves exactly the anomaly carrier]
\label{thm:v2v90-canonical}
The least-norm counterterm
\begin{equation}
 c_n=-\delta_n^\dagger y_n
\label{eq:v2v90-cn}
\end{equation}
satisfies
\begin{equation}
 y_n+\delta_nc_n
 =
 \mathfrak a_n.
\label{eq:v2v90-remainder}
\end{equation}
Thus \(y_{n,\rm rem}\) is removed and no component of
\(\mathfrak a_n\) is changed by an allowed counterterm.  Moreover,
\begin{equation}
 \|c_n\|\le\sigma_0^{-1}\|y_{n,\rm rem}\|.
\label{eq:v2v90-cn-bound}
\end{equation}
\end{theorem}

\begin{proof}
The Moore--Penrose identities give
\[
 \delta_n\delta_n^\dagger=P_n.
\]
Therefore
\[
 y_n+\delta_nc_n
 =
 (I-P_n)y_n=A_ny_n.
\]
The two projections are orthogonal, so the counterterm acts only on the
range component.  The norm bound follows from
\(\|\delta_n^\dagger\|\le\sigma_0^{-1}\).
\end{proof}

\begin{assessment}[The anomaly is a projected vector, not a failed solve]
\label{ass:v2v90-anomaly-vector}
Once the local-null quotient and allowed counterterm space are fixed,
\(\mathfrak a_n=A_ny_n\) is the unique orthogonal obstruction.  Changing
the counterterm by an element of \(\ker\delta_n\) does not alter it.
Estimating diagrammatic pieces of \(y_n\) before applying \(A_n\) can
destroy cancellations inside the anomaly coefficient.
\end{assessment}

\subsection{Continuum anomaly matching}
\label{sec:v2v90-matching}

\begin{theorem}[Stable anomaly and counterterm limits]
\label{thm:v2v90-matching}
Suppose
\begin{equation}
 y_n\longrightarrow y.
\label{eq:v2v90-y-convergence}
\end{equation}
Then
\begin{align}
 \mathfrak a_n=A_ny_n
 &\longrightarrow
 \mathfrak a:=Ay,
 \label{eq:v2v90-anomaly-convergence}\\
 c_n=-\delta_n^\dagger y_n
 &\longrightarrow
 c:=-\delta^\dagger y.
 \label{eq:v2v90-counterterm-convergence}
\end{align}
If \(\mathfrak a=0\), then the regulated anomaly carrier tends to zero.
If \(\mathfrak a\ne0\), no sequence of uniformly bounded allowed
counterterms can remove its limiting cokernel component.
\end{theorem}

\begin{proof}
Combine
\eqref{eq:v2v90-A-convergence},
\eqref{eq:v2v90-pseudoinverse-convergence}, and
\eqref{eq:v2v90-y-convergence}.  For the last statement, every corrected
defect has \(A_n\)-projection \(\mathfrak a_n\), because
\(A_n\delta_n=0\).  Passing to the limit gives the nonzero vector
\(\mathfrak a\).
\end{proof}

Let
\[
 \mathcal E_{{\rm os},n}:Y\longrightarrow\mathcal O
\]
be the physical on-shell evaluation after the local quotient, and assume
\begin{equation}
 \mathcal E_{{\rm os},n}
 \longrightarrow
 \mathcal E_{\rm os}
\label{eq:v2v90-E-convergence}
\end{equation}
in norm.

\begin{corollary}[The V88 residue sees only the anomaly projection]
\label{cor:v2v90-residue}
After the canonical counterterm,
\begin{equation}
 \mathcal E_{{\rm os},n}
 (y_n+\delta_nc_n)
 =
 \mathcal E_{{\rm os},n}\mathfrak a_n
 \longrightarrow
 \mathcal E_{\rm os}\mathfrak a.
\label{eq:v2v90-residue}
\end{equation}
Under the V87 Jacobian lower bound, this is exactly the coefficient
controlling the uncancelled matrix pinch distribution.
\end{corollary}

\begin{proof}
Use
\eqref{eq:v2v90-remainder},
\eqref{eq:v2v90-anomaly-convergence}, and
\eqref{eq:v2v90-E-convergence}.  V88 expresses the residue as the
on-shell evaluation of the complete corrected Ward defect.
\end{proof}

\subsection{A transported anomaly basis}
\label{sec:v2v90-basis}

\begin{proposition}[Finite stable obstruction coefficients]
\label{prop:v2v90-basis}
Let \(r=\operatorname{rank}A\).  For all large \(n\), there are
orthonormal bases
\[
 a_{1,n},\ldots,a_{r,n}
 \quad\hbox{of }\operatorname{Ran}A_n
\]
and
\[
 a_1,\ldots,a_r
 \quad\hbox{of }\operatorname{Ran}A
\]
such that
\begin{equation}
 a_{b,n}\longrightarrow a_b.
\label{eq:v2v90-basis-convergence}
\end{equation}
Writing
\begin{equation}
 \mathfrak a_n
 =
 \sum_{b=1}^r\alpha_{b,n}a_{b,n},
\label{eq:v2v90-coefficients}
\end{equation}
one has
\begin{equation}
 \alpha_{b,n}\longrightarrow
 \alpha_b:=\langle a_b,y\rangle.
\label{eq:v2v90-coefficient-convergence}
\end{equation}
\end{proposition}

\begin{proof}
Norm convergence of the equal-rank projections \(A_n\to A\) permits the
Kato unitary between their ranges.  Transport a fixed orthonormal basis
of \(\operatorname{Ran}A\) by its inverse.  This gives
\eqref{eq:v2v90-basis-convergence}.  Then
\[
 \alpha_{b,n}
 =
 \langle a_{b,n},y_n\rangle
 \longrightarrow
 \langle a_b,y\rangle.
\]
\end{proof}

This basis is used only after the complete defect has been projected.
It is not a license to assign independent counterterms to diagrammatic
letters.

\subsection{Sharp rank-loss separator}
\label{sec:v2v90-rank-loss}

\begin{counterexample}[Every regulator is anomaly-free but the limit is not]
\label{sep:v2v90-rank-loss}
Let
\[
 X=Y=\mathbb C^2,
 \qquad
 \delta_n=
 \begin{pmatrix}
 1&0\\
 0&\varepsilon_n
 \end{pmatrix},
 \qquad
 \varepsilon_n\downarrow0.
\]
For every \(n\), \(\delta_n\) is surjective, so \(A_n=0\) and every
defect is removable.  The limit is
\[
 \delta=
 \begin{pmatrix}
 1&0\\
 0&0
 \end{pmatrix},
 \qquad
 A=
 \begin{pmatrix}
 0&0\\
 0&1
 \end{pmatrix}.
\]
For the fixed defect \(y=e_2\), the regulated counterterm is
\begin{equation}
 c_n=-\varepsilon_n^{-1}e_2,
\label{eq:v2v90-divergent-counterterm}
\end{equation}
while the continuum anomaly is
\begin{equation}
 Ay=e_2\ne0.
\label{eq:v2v90-limit-anomaly}
\end{equation}
\end{counterexample}

\begin{proof}
All statements follow by diagonal matrix calculation.  The smallest
singular value is \(\varepsilon_n\), so the uniform gap
\eqref{eq:v2v90-singular-gap} fails.
\end{proof}

\begin{firewall}[Finite-regulator exactness does not prove anomaly cancellation]
\label{fw:v2v90-finite-exact}
The counterexample proves that surjectivity of every finite counterterm
map is insufficient.  A counterterm direction can collapse in the
continuum, turning an exactly removable regulator defect into a genuine
limiting cokernel class.
\end{firewall}

\subsection{Compatibility with the spectral construction}
\label{sec:v2v90-spectral}

\begin{corollary}[Stable range splitting preserves V89 control]
\label{cor:v2v90-spectral}
Assume the counterterm insertion map satisfies
\eqref{eq:v2v89-insertion-bound}.  If
\(y_{n,\rm rem}\to0\), then the canonical counterterms satisfy
\(c_n\to0\) and preserve the V87 differentiated norm-resolvent limit by
Theorem~\ref{thm:v2v89-spectral}.  The anomaly component
\(\mathfrak a_n\) is not inserted into the kinetic operator; it remains
as the explicit physical obstruction coefficient in
\eqref{eq:v2v90-residue}.
\end{corollary}

\begin{proof}
Use
\eqref{eq:v2v90-cn-bound} and
Theorem~\ref{thm:v2v89-spectral}.  Orthogonality of the split prevents
the anomaly component from entering the chosen counterterm.
\end{proof}

\begin{firewall}[A finite renormalization changes the matched continuum data]
\label{fw:v2v90-finite-renormalization}
If \(c_n\to c\ne0\), spectral convergence can still hold after the
continuum kinetic and vertex data are corrected by the matching finite
counterterm.  It is then a different renormalization condition, and the
crossing locations and Jacobians must be recomputed for that corrected
limit.
\end{firewall}

\subsection{V90 conclusion and next acquisition}
\label{sec:v2v90-conclusion}

The mandatory audit has been completed.  V90 proves a regulator-stable
orthogonal decomposition of the complete Ward defect into its removable
counterterm-range component and its anomaly cokernel component.  Under a
uniform nonzero singular-value gap, the projections, least-norm
counterterms, anomaly basis, and on-shell residue coefficients converge.
A two-dimensional example proves that anomaly-free finite regulators can
converge to an anomalous theory when that gap collapses.

The next step is representation-theoretic rather than another abstract
inverse theorem.  V91 will compute the perturbative local gauge and mixed
gravitational anomaly coefficients of one Standard Model generation in
a fixed left-handed convention, and identify exactly which global and
boundary anomalies remain outside that calculation.
\clearpage
\section*{Second Volume, Version V91: Local Anomaly Coefficients of One Standard Model Generation}
\addcontentsline{toc}{section}{Second Volume, Version V91: Local Anomaly Coefficients of One Standard Model Generation}
\label{sec:v2v91}

\subsection{Purpose and convention}
\label{sec:v2v91-purpose}

V90 isolates the continuum anomaly as a vector in the cokernel of the
allowed counterterm variation.  The next step is to compute the
coefficients of the standard four-dimensional perturbative anomaly
basis for the actual Standard Model fermion representation.

Use only left-handed Weyl fields.  A physical right-handed fermion is
replaced by its left-handed charge conjugate, which conjugates the
nonabelian representation and reverses hypercharge.  For one generation,
\begin{equation}
\begin{array}{c|c|c|c}
\text{field}&SU(3)_c&SU(2)_L&Y\\
\hline
Q&\mathbf3&\mathbf2&\frac16\\
u^c&\overline{\mathbf3}&\mathbf1&-\frac23\\
d^c&\overline{\mathbf3}&\mathbf1&\frac13\\
L&\mathbf1&\mathbf2&-\frac12\\
e^c&\mathbf1&\mathbf1&1
\end{array}
\label{eq:v2v91-field-table}
\end{equation}
An optional sterile \(\nu^c\) has
\((\mathbf1,\mathbf1)_0\) and changes none of the calculations.

Normalize the quadratic indices by
\begin{equation}
 T(\mathbf3)=T(\overline{\mathbf3})
 =
 T(\mathbf2)=\frac12,
\label{eq:v2v91-quadratic-index}
\end{equation}
and the cubic \(SU(3)\) indices by
\begin{equation}
 A(\mathbf3)=1,
 \qquad
 A(\overline{\mathbf3})=-1.
\label{eq:v2v91-cubic-index}
\end{equation}
Only vanishing is relevant, so an overall normalization of every anomaly
coefficient is immaterial.

\subsection{Complete perturbative coefficient list}
\label{sec:v2v91-list}

For a product representation
\((R_3,R_2)_Y\), write \(d_i=\dim R_i\).  The potentially nonzero local
gauge and mixed gauge--gravitational coefficients in four dimensions are
\begin{align}
 \mathcal A_{333}
 &=
 \sum_f d_2(f)A(R_3(f)),
 \label{eq:v2v91-A333}\\
 \mathcal A_{331}
 &=
 \sum_f d_2(f)Y_fT(R_3(f)),
 \label{eq:v2v91-A331}\\
 \mathcal A_{221}
 &=
 \sum_f d_3(f)Y_fT(R_2(f)),
 \label{eq:v2v91-A221}\\
 \mathcal A_{111}
 &=
 \sum_f d_3(f)d_2(f)Y_f^3,
 \label{eq:v2v91-A111}\\
 \mathcal A_{{\rm grav}{\rm grav}1}
 &=
 \sum_f d_3(f)d_2(f)Y_f.
 \label{eq:v2v91-Agg1}
\end{align}
A term is omitted when the indicated nonabelian representation is
trivial.

\begin{lemma}[Completeness of the local list]
\label{lem:v2v91-complete-list}
For the Lie algebra
\[
 \mathfrak{su}(3)\oplus\mathfrak{su}(2)\oplus\mathfrak u(1),
\]
every perturbative four-dimensional gauge-anomaly coefficient is either
one of
\eqref{eq:v2v91-A333}--\eqref{eq:v2v91-A111} or vanishes by tracelessness
or absence of a symmetric cubic invariant.  The only mixed local
gravitational coefficient is
\eqref{eq:v2v91-Agg1}.
\end{lemma}

\begin{proof}
A left-handed triangle coefficient is the symmetric trace
\[
 \operatorname{Tr}_{R_f}
 T^a\{T^b,T^c\}.
\]
For three \(SU(3)\) generators this gives the cubic index
\(\mathcal A_{333}\).  The Lie algebra \(\mathfrak{su}(2)\) has no
nonzero invariant symmetric tensor of rank three, so there is no local
\(SU(2)^3\) coefficient.  Two generators from one simple factor and one
hypercharge give
\(\mathcal A_{331}\) or \(\mathcal A_{221}\).  Three hypercharges give
\(\mathcal A_{111}\).

A trace with one generator from a simple factor is zero because every
finite-dimensional representation of a semisimple Lie algebra has
traceless generators.  This kills \(U(1)^2SU(3)\),
\(U(1)^2SU(2)\), every term with one generator from each simple factor,
and the mixed gravitational coefficient of a simple generator.
The gravitational trace can therefore retain only the hypercharge
generator, giving \eqref{eq:v2v91-Agg1}.
\end{proof}

\subsection{Nonabelian cubic anomaly}
\label{sec:v2v91-SU3}

\begin{proposition}[\(SU(3)_c^3\) cancellation]
\label{prop:v2v91-SU3}
For the representation
\eqref{eq:v2v91-field-table},
\begin{equation}
 \mathcal A_{333}
 =
 2A(\mathbf3)
 +A(\overline{\mathbf3})
 +A(\overline{\mathbf3})
 =
 2-1-1=0.
\label{eq:v2v91-SU3-cancel}
\end{equation}
\end{proposition}

\begin{proof}
The quark doublet \(Q\) contains two left-handed color triplets.
The fields \(u^c,d^c\) are one color antitriplet each.  Leptons are color
singlets.  Insert their cubic indices.
\end{proof}

\begin{remark}[Vectorlike color after electroweak multiplicity]
\label{rem:v2v91-vectorlike-color}
The cancellation is the chiral, left-handed form of the vectorlike color
content: the two triplets inside \(Q\) are paired in cubic index with
\(u^c\) and \(d^c\).
\end{remark}

\subsection{Mixed nonabelian--hypercharge anomalies}
\label{sec:v2v91-mixed}

\begin{proposition}[\(SU(3)_c^2U(1)_Y\) cancellation]
\label{prop:v2v91-SU3SU3Y}
One has
\begin{align}
 \mathcal A_{331}
 &=
 2\left(\frac16\right)\frac12
 +\left(-\frac23\right)\frac12
 +\left(\frac13\right)\frac12
 \nonumber\\
 &=
 \frac16-\frac13+\frac16=0.
\label{eq:v2v91-SU3SU3Y-cancel}
\end{align}
\end{proposition}

\begin{proof}
The factor two in the \(Q\) term is its weak multiplicity.  The quadratic
index has the same sign for a representation and its conjugate; the
hypercharge already carries the chirality-conjugation sign.
\end{proof}

\begin{proposition}[\(SU(2)_L^2U(1)_Y\) cancellation]
\label{prop:v2v91-SU2SU2Y}
One has
\begin{align}
 \mathcal A_{221}
 &=
 3\left(\frac16\right)\frac12
 +\left(-\frac12\right)\frac12
 \nonumber\\
 &=
 \frac14-\frac14=0.
\label{eq:v2v91-SU2SU2Y-cancel}
\end{align}
\end{proposition}

\begin{proof}
Only \(Q\) and \(L\) are weak doublets.  The quark term has color
multiplicity three.
\end{proof}

\subsection{Pure hypercharge and mixed gravitational coefficients}
\label{sec:v2v91-abelian}

\begin{proposition}[\(U(1)_Y^3\) cancellation]
\label{prop:v2v91-U1}
The cubic hypercharge coefficient is
\begin{align}
 \mathcal A_{111}
 &=
 6\left(\frac16\right)^3
 +3\left(-\frac23\right)^3
 +3\left(\frac13\right)^3
 +2\left(-\frac12\right)^3
 +1
 \nonumber\\
 &=
 \frac1{36}-\frac89+\frac19-\frac14+1
 \nonumber\\
 &=
 \frac{1-32+4-9+36}{36}
 =0.
\label{eq:v2v91-U1-cancel}
\end{align}
\end{proposition}

\begin{proof}
The multiplicities are respectively
\(3\cdot2,3,3,2,1\), the dimensions of the nonabelian representations.
The final numerator vanishes.
\end{proof}

\begin{proposition}[Mixed gravitational--hypercharge cancellation]
\label{prop:v2v91-gravY}
The coefficient of the local
\(\operatorname{grav}^2U(1)_Y\) anomaly is
\begin{align}
 \mathcal A_{{\rm grav}{\rm grav}1}
 &=
 6\left(\frac16\right)
 +3\left(-\frac23\right)
 +3\left(\frac13\right)
 +2\left(-\frac12\right)
 +1
 \nonumber\\
 &=1-2+1-1+1=0.
\label{eq:v2v91-gravY-cancel}
\end{align}
\end{proposition}

\begin{proof}
The gravitational vertices act as the identity on each internal gauge
multiplet, leaving the trace of hypercharge with the same representation
multiplicities as above.
\end{proof}

\subsection{Integrated local theorem}
\label{sec:v2v91-local-theorem}

\begin{theorem}[One Standard Model generation is perturbatively gauge-anomaly free]
\label{thm:v2v91-local}
For the left-handed representation in
\eqref{eq:v2v91-field-table},
\begin{equation}
 \mathcal A_{333}
 =
 \mathcal A_{331}
 =
 \mathcal A_{221}
 =
 \mathcal A_{111}
 =
 \mathcal A_{{\rm grav}{\rm grav}1}
 =
 0.
\label{eq:v2v91-all-zero}
\end{equation}
The local \(SU(2)^3\) coefficient and all remaining mixed coefficients
vanish identically for the structural reasons in
Lemma~\ref{lem:v2v91-complete-list}.  Therefore the perturbative local
gauge anomaly and the mixed gauge--gravitational anomaly cancel within
each generation.
\end{theorem}

\begin{proof}
Combine
Propositions~\ref{prop:v2v91-SU3},
\ref{prop:v2v91-SU3SU3Y},
\ref{prop:v2v91-SU2SU2Y},
\ref{prop:v2v91-U1}, and
\ref{prop:v2v91-gravY} with
Lemma~\ref{lem:v2v91-complete-list}.
\end{proof}

\begin{corollary}[Three generations and a sterile neutrino]
\label{cor:v2v91-three}
Repeating the representation three times multiplies every coefficient
in \eqref{eq:v2v91-all-zero} by three and hence preserves cancellation.
Adding any number of gauge-singlet left-handed fields of hypercharge zero
also preserves it.
\end{corollary}

\begin{proof}
Triangle coefficients add over fermion species.  A
\((\mathbf1,\mathbf1)_0\) field contributes zero to every displayed
trace.
\end{proof}

\subsection{Ordinary global \(SU(2)\) parity test}
\label{sec:v2v91-Witten}

The four-dimensional \(SU(2)\) mod-two anomaly theorem states that a
theory containing an odd number of left-handed fundamental doublets is
inconsistent.  This is a global condition and is not detected by the
local symmetric trace.

\begin{proposition}[Even doublet count per generation]
\label{prop:v2v91-Witten}
One Standard Model generation contains
\begin{equation}
 3\ \text{quark doublets}
 +
 1\ \text{lepton doublet}
 =
 4
\label{eq:v2v91-four-doublets}
\end{equation}
left-handed \(SU(2)\) doublets.  Hence the ordinary mod-two
fundamental-doublet obstruction vanishes.  Three generations contain
twelve doublets and also pass this parity test.
\end{proposition}

\begin{proof}
Each color copy of \(Q\) is a weak doublet, producing three.  The lepton
field \(L\) produces one.  The remaining fermions are weak singlets.
Four and twelve are even.
\end{proof}

\begin{firewall}[This parity count is not a full global-anomaly theorem]
\label{fw:v2v91-global}
The calculation verifies the ordinary \(SU(2)\) fundamental-doublet
mod-two condition.  It does not compute the relevant spin-bordism group
for every global form
\[
 (SU(3)\times SU(2)\times U(1))/\Gamma,
\]
nor determinant holonomy on the allowed Einstein backgrounds or
boundary mapping tori.  Those global choices can change the anomaly
classification.
\end{firewall}

\subsection{Insertion into the V90 obstruction projection}
\label{sec:v2v91-V90}

Let \(\mathcal A_{\rm loc}\) be the subspace of the continuum V90 anomaly
cokernel spanned by the standard perturbative local gauge and mixed
gauge--gravitational descent classes, and let \(A_{\rm loc}\) denote the
orthogonal projection onto this subspace.

\begin{corollary}[Vanishing of the perturbative local SM obstruction]
\label{cor:v2v91-obstruction}
If the physical on-shell defect is represented in
\(\mathcal A_{\rm loc}\) with coefficients given by the chiral fermion
traces
\eqref{eq:v2v91-A333}--\eqref{eq:v2v91-Agg1}, then its V90 anomaly
projection vanishes for the Standard Model representation:
\begin{equation}
 A_{\rm loc}y_{\rm SM}=0.
\label{eq:v2v91-local-obstruction-zero}
\end{equation}
Under the uniform singular-gap and insertion hypotheses of V89--V90,
the corresponding removable regulator defect can be corrected without
destroying the V87 pinch limit.
\end{corollary}

\begin{proof}
Theorem~\ref{thm:v2v91-local} makes every coefficient in the stated
basis zero.  Apply
Theorem~\ref{thm:v2v90-canonical} to the remaining counterterm-range
component and
Theorem~\ref{thm:v2v89-spectral} to its insertion.
\end{proof}

\begin{firewall}[The cohomology-basis identification is an input]
\label{fw:v2v91-basis}
The trace calculation proves vanishing once the regulator's local
consistent anomaly has been identified with the standard descent basis.
It does not prove that the complete boundary regulator has no additional
local-null mismatch, boundary class, or nonlocal defect.  That
identification must be established by a locality and covariance theorem
for the chosen regulator.
\end{firewall}

\subsection{Anomalies which are not removed by this table}
\label{sec:v2v91-scope}

\begin{firewall}[Trace anomaly]
\label{fw:v2v91-trace}
The curved-space trace or Weyl anomaly is generally nonzero in the
Standard Model.  It is not a gauge-consistency anomaly and is not
measured by
\eqref{eq:v2v91-Agg1}.  Its curvature counterterms and stress-tensor
Ward identity remain part of the Einstein backreaction problem.
\end{firewall}

\begin{firewall}[Global currents]
\label{fw:v2v91-global-currents}
Anomalies of ungauged baryon or lepton currents, including the
electroweak violation of \(B+L\), do not contradict
Theorem~\ref{thm:v2v91-local}.  The theorem concerns consistency of the
gauged Standard Model symmetries.
\end{firewall}

\begin{firewall}[Boundary and determinant phase]
\label{fw:v2v91-boundary}
Local bulk triangle cancellation does not fix the phase of the chiral
fermion determinant on a manifold with boundary.  Eta invariants,
boundary conditions, anomaly inflow, and large gauge transformations
remain to be matched.
\end{firewall}

\begin{firewall}[Pure gravitational anomaly]
\label{fw:v2v91-pure-gravity}
There is no perturbative pure diffeomorphism anomaly for ordinary chiral
spin-\(\tfrac12\) matter in four spacetime dimensions.  This does not
remove the mixed hypercharge calculation, the trace anomaly, or possible
global gravitational phases.
\end{firewall}

\subsection{V91 conclusion and next acquisition}
\label{sec:v2v91-conclusion}

V91 computes every standard perturbative local gauge and mixed
gauge--gravitational anomaly coefficient for one Standard Model
generation in a single left-handed convention.  All vanish generation
by generation.  The ordinary global \(SU(2)\) doublet parity also passes.
This removes the representation coefficient in the local bulk part of
the V90 anomaly cokernel, subject to identifying the regulator defect
with the standard descent basis.

V92 will address that identification on a manifold with boundary.  It
will formulate a relative bulk--boundary Ward complex, prove the exact
sequence separating bulk descent, boundary inflow, and determinant
phase, and state the analytic estimates required for a regulator-uniform
eta contribution.

\clearpage
\section*{Second Volume, Version V92: Relative Bulk--Boundary Ward Complex and Eta Stability}
\addcontentsline{toc}{section}{Second Volume, Version V92: Relative Bulk--Boundary Ward Complex and Eta Stability}
\label{sec:v2v92}

\subsection{Purpose}
\label{sec:v2v92-purpose}

V91 cancels the standard perturbative local bulk anomaly coefficients.
A manifold with boundary has additional data: the restriction of a bulk
descent cocycle, a possible boundary inflow functional, and the phase of
the chiral determinant.  These are not independent defects.

V92 packages them in a mapping-cone complex.  Its long exact sequence
states precisely when a bulk anomaly admits boundary compensation and
which pure boundary classes remain.  A second theorem gives sufficient
uniform heat-kernel hypotheses for convergence of the renormalized eta
invariant.  No claim is made that the required boundary estimates already
hold for the full Einstein-coupled regulator.

\subsection{Bulk and boundary complexes}
\label{sec:v2v92-complexes}

Let
\[
 (\mathcal B^\bullet,d_{\mathcal B})
 \quad\hbox{and}\quad
 (\mathcal E^\bullet,d_{\mathcal E})
\]
be cochain complexes describing local bulk and boundary Ward
functionals.  Let
\begin{equation}
 r:\mathcal B^\bullet\longrightarrow\mathcal E^\bullet
\label{eq:v2v92-restriction}
\end{equation}
be a chain map:
\begin{equation}
 r\,d_{\mathcal B}=d_{\mathcal E}r.
\label{eq:v2v92-chain-map}
\end{equation}

Define the shifted mapping cone
\begin{equation}
 \mathcal C_{\rm rel}^p
 =
 \mathcal B^p\oplus\mathcal E^{p-1}
\label{eq:v2v92-cone}
\end{equation}
with differential
\begin{equation}
 D(b,e)
 =
 \bigl(
 d_{\mathcal B}b,\
 rb-d_{\mathcal E}e
 \bigr).
\label{eq:v2v92-D}
\end{equation}

\begin{lemma}[Relative differential]
\label{lem:v2v92-D-square}
One has
\begin{equation}
 D^2=0.
\label{eq:v2v92-D-square}
\end{equation}
A relative cocycle of degree \(p\) is therefore a pair satisfying
\begin{equation}
 d_{\mathcal B}b=0,
 \qquad
 rb=d_{\mathcal E}e.
\label{eq:v2v92-relative-cycle}
\end{equation}
\end{lemma}

\begin{proof}
Apply \(D\) twice:
\[
 D^2(b,e)
 =
 \bigl(
 d_{\mathcal B}^2b,\
 r\,d_{\mathcal B}b-d_{\mathcal E}rb
 +d_{\mathcal E}^2e
 \bigr)=0
\]
by \eqref{eq:v2v92-chain-map}.  The cocycle equations are the two
components of \(D(b,e)=0\).
\end{proof}

For a degree-one anomaly, \(b\) is the consistent bulk Ward cocycle and
\(e\) is a degree-zero boundary functional whose variation cancels the
boundary restriction of \(b\).

\subsection{The exact bulk--boundary sequence}
\label{sec:v2v92-exact-sequence}

Define
\begin{align}
 \iota:
 H^{p-1}(\mathcal E)
 &\longrightarrow H^p(\mathcal C_{\rm rel}),
 &
 [e]&\longmapsto[(0,e)],
 \label{eq:v2v92-iota}\\
 j:
 H^p(\mathcal C_{\rm rel})
 &\longrightarrow H^p(\mathcal B),
 &
 [(b,e)]&\longmapsto[b].
 \label{eq:v2v92-j}
\end{align}

\begin{theorem}[Long exact relative Ward sequence]
\label{thm:v2v92-long-exact}
The maps above fit into the long exact sequence
\begin{equation}
 \cdots\longrightarrow
 H^{p-1}(\mathcal B)
 \xrightarrow{\,r^*\,}
 H^{p-1}(\mathcal E)
 \xrightarrow{\,\iota\,}
 H^p(\mathcal C_{\rm rel})
 \xrightarrow{\,j\,}
 H^p(\mathcal B)
 \xrightarrow{\,r^*\,}
 H^p(\mathcal E)
 \longrightarrow\cdots.
\label{eq:v2v92-long-exact}
\end{equation}
In particular:
\begin{enumerate}
 \item a bulk class \([b]\) has a relative bulk--boundary lift if and
 only if \(r^*[b]=0\);
 \item a relative class with zero bulk component is a boundary class in
 \(H^{p-1}(\mathcal E)\), modulo restrictions of bulk classes.
\end{enumerate}
\end{theorem}

\begin{proof}
First, \(\iota\) is well defined: if
\(d_{\mathcal E}e=0\), then \(D(0,e)=0\), while
\(e=d_{\mathcal E}f\) gives
\[
 (0,e)=-D(0,f).
\]
The map \(j\) is well defined because a relative boundary
\(D(c,f)\) has bulk component \(d_{\mathcal B}c\).

Exactness at \(H^{p-1}(\mathcal E)\) follows because
\[
 \iota(r^*[c])
 =
 [(0,rc)]
 =
 [D(c,0)]
 =
 0
\]
for a closed \(c\), and the converse is obtained by reading the boundary
equation for a relative primitive.

For a relative cocycle, \(rb=d_{\mathcal E}e\), so
\(r^*j[(b,e)]=0\).  Conversely, if \(b\) is closed and
\(rb=d_{\mathcal E}e\), then \((b,e)\) is a relative cocycle mapping to
\([b]\).  This proves exactness at \(H^p(\mathcal B)\).

Finally, suppose \(j[(b,e)]=0\), so
\(b=d_{\mathcal B}c\).  Subtract the relative boundary
\[
 D(c,0)=(d_{\mathcal B}c,rc).
\]
The remaining representative is
\((0,e-rc)\), and the relative cocycle equation shows that
\(e-rc\) is boundary-closed.  Hence the class lies in
\(\operatorname{im}\iota\).  These verifications repeat in every degree.
\end{proof}

\begin{corollary}[Bulk cancellation does not remove pure boundary phases]
\label{cor:v2v92-boundary-class}
If the V91 bulk class vanishes, the remaining relative obstruction lies
in
\begin{equation}
 H^0(\mathcal E)/
 r^*H^0(\mathcal B)
\label{eq:v2v92-pure-boundary}
\end{equation}
at anomaly degree one.  Thus a determinant phase or boundary inflow
class can remain even when every local bulk triangle coefficient is zero.
\end{corollary}

\begin{proof}
Set \(p=1\) in
\eqref{eq:v2v92-long-exact} and use exactness at
\(H^1(\mathcal C_{\rm rel})\).
\end{proof}

\subsection{Relative counterterms}
\label{sec:v2v92-counterterms}

A degree-zero relative counterterm is a pair
\[
 (c,f)\in\mathcal B^0\oplus\mathcal E^{-1}.
\]
It changes the anomaly pair by
\begin{equation}
 (b,e)\longmapsto
 (b,e)+D(c,f)
 =
 \bigl(
 b+d_{\mathcal B}c,\
 e+rc-d_{\mathcal E}f
 \bigr).
\label{eq:v2v92-counterterm}
\end{equation}

\begin{proposition}[Counterterms preserve the relative class]
\label{prop:v2v92-counterterm}
Two bulk--boundary anomaly pairs differ by an allowed relative
counterterm exactly when they represent the same class in
\(H^1(\mathcal C_{\rm rel})\).
\end{proposition}

\begin{proof}
This is the definition of cohomology for the mapping-cone differential.
\end{proof}

\begin{assessment}[Where inflow sits]
\label{ass:v2v92-inflow}
The equation \(rb=d_{\mathcal E}e\) is the algebraic inflow condition.
It does not require either term to vanish separately.  Taking norms of
the bulk restriction and boundary variation before forming their
difference would destroy the relative cocycle.
\end{assessment}

\subsection{Quantitative relative Hodge contraction}
\label{sec:v2v92-Hodge}

At regulator \(n\), give
\(\mathcal C_{{\rm rel},n}^\bullet\) a finite-dimensional Hilbert
structure.  Let \(D_n^*\) be the adjoint and define
\begin{equation}
 \Delta_{{\rm rel},n}
 =
 D_nD_n^*+D_n^*D_n.
\label{eq:v2v92-relative-Laplacian}
\end{equation}
Let \(\Pi_{{\rm rel},n}\) project onto its kernel.  Assume
\begin{equation}
 \Delta_{{\rm rel},n}
 \ge
 \gamma_{\rm rel}^2(I-\Pi_{{\rm rel},n})
\label{eq:v2v92-relative-gap}
\end{equation}
uniformly.

\begin{theorem}[Uniform relative primitive]
\label{thm:v2v92-relative-Hodge}
Define
\begin{equation}
 h_{{\rm rel},n}
 =
 D_n^*
 \Delta_{{\rm rel},n}^{-1}
 (I-\Pi_{{\rm rel},n}).
\label{eq:v2v92-relative-h}
\end{equation}
Then
\begin{align}
 D_nh_{{\rm rel},n}
 +h_{{\rm rel},n}D_n
 &=
 I-\Pi_{{\rm rel},n},
 \label{eq:v2v92-relative-contraction}\\
 \|h_{{\rm rel},n}\|
 &\le
 \gamma_{\rm rel}^{-1}.
 \label{eq:v2v92-relative-h-bound}
\end{align}
For every relative cocycle \(z_n\),
\begin{equation}
 z_n
 =
 \Pi_{{\rm rel},n}z_n
 +
 D_nh_{{\rm rel},n}z_n.
\label{eq:v2v92-relative-split}
\end{equation}
\end{theorem}

\begin{proof}
Apply the Hodge-contraction proof of
Theorem~\ref{thm:v2v68-contract} to the differential \(D_n\).
The mapping-cone identity \(D_n^2=0\) supplies its only algebraic input,
and \eqref{eq:v2v92-relative-gap} supplies the norm bound.
\end{proof}

\begin{corollary}[Regulator convergence of relative classes]
\label{cor:v2v92-relative-convergence}
If, after fixed identifications,
\[
 D_n\to D,
 \qquad
 D_n^*\to D^*
\]
in norm and
\eqref{eq:v2v92-relative-gap} holds, then the relative harmonic
projections and contracting homotopies converge in norm.  Hence the
bulk--boundary obstruction coefficients converge in a transported
finite harmonic basis.
\end{corollary}

\begin{proof}
Repeat the spectral-projection argument in
Theorem~\ref{thm:v2v88-Hodge-convergence} for
\(\Delta_{{\rm rel},n}\).
\end{proof}

\subsection{Boundary eta data}
\label{sec:v2v92-eta}

Let \(A_n\) be self-adjoint tangential boundary Dirac operators with
compact resolvent, after compatible identification of their boundary
Hilbert spaces.  Put
\begin{equation}
 \Theta_n(t)
 =
 \operatorname{Tr}
 \bigl(A_ne^{-tA_n^2}\bigr).
\label{eq:v2v92-Theta}
\end{equation}
The formal eta function is the Mellin transform
\begin{equation}
 \eta_n(s)
 =
 \frac1{\Gamma((s+1)/2)}
 \int_0^\infty
 t^{(s-1)/2}\Theta_n(t)\,\mathrm dt,
\label{eq:v2v92-eta-Mellin}
\end{equation}
initially where the integral converges and then continued to \(s=0\).

To state a regulator-uniform theorem, assume:
\begin{enumerate}
 \item there is a boundary gap
 \begin{equation}
  \operatorname{spec}A_n\cap(-\delta_\partial,\delta_\partial)
  =\varnothing;
 \label{eq:v2v92-boundary-gap}
 \end{equation}
 \item for \(0<t\le1\),
 \begin{equation}
  \Theta_n(t)
  =
  S_n(t)+R_n(t),
  \qquad
  S_n(t)=\sum_{\ell=1}^Nc_{\ell,n}t^{\alpha_\ell},
 \label{eq:v2v92-heat-split}
 \end{equation}
 with fixed exponents \(\alpha_\ell\),
 \(c_{\ell,n}\to c_\ell\), and
 \begin{equation}
  |R_n(t)|\le Ct^{-1/2+\rho}
  \quad\text{for some }\rho>0;
 \label{eq:v2v92-small-remainder}
 \end{equation}
 \item \(\Theta_n\to\Theta\) and \(R_n\to R\) pointwise for \(t>0\),
 with convergence dominated by the displayed small-time bound and
 \begin{equation}
  |\Theta_n(t)|\le Ce^{-\delta_\partial^2t}
  \qquad(t\ge1).
 \label{eq:v2v92-large-tail}
 \end{equation}
\end{enumerate}
The finite list \(S_n\) is the local heat subtraction fixed by the
renormalization prescription.

\begin{theorem}[Convergence of the renormalized eta invariant]
\label{thm:v2v92-eta}
Under
\eqref{eq:v2v92-boundary-gap}--\eqref{eq:v2v92-large-tail}, the finite
part at \(s=0\) of
\eqref{eq:v2v92-eta-Mellin} exists in the chosen common subtraction
scheme, and
\begin{equation}
 \eta_{{\rm ren},n}(0)
 \longrightarrow
 \eta_{\rm ren}(0).
\label{eq:v2v92-eta-convergence}
\end{equation}
Consequently, in the chosen common trivialization, every scalar phase`r`nof the form $\exp(i\kappa\eta_{{\rm ren},n}(0))$ with fixed`r`n$\kappa\in\mathbb R$ converges.
\end{theorem}

\begin{proof}
Split the Mellin integral at \(t=1\).  On \((0,1]\), subtract
\(S_n\).  At \(s=0\),
\[
 t^{-1/2}R_n(t)
 \le
 Ct^{-1+\rho},
\]
which is integrable.  Dominated convergence therefore handles the
remainder integral.  Each subtraction monomial contributes
\[
 c_{\ell,n}
 \int_0^1
 t^{\alpha_\ell+(s-1)/2}\,\mathrm dt
 =
 \frac{c_{\ell,n}}
 {\alpha_\ell+(s+1)/2}
\]
by meromorphic continuation.  Taking the finite part at \(s=0\) is a
fixed continuous linear operation on the finite coefficient vector
\((c_{\ell,n})\), so these terms converge.

The large-time integral converges uniformly by
\eqref{eq:v2v92-large-tail} and dominated convergence.  Finally,
\(\Gamma((s+1)/2)^{-1}\) is regular at \(s=0\).  Combining the three
pieces proves \eqref{eq:v2v92-eta-convergence}.  Continuity of the chosen
exponential gives the phase statement.
\end{proof}

\begin{assessment}[What the heat assumptions encode]
\label{ass:v2v92-heat}
The small-time coefficients are local boundary geometry and gauge data;
their subtraction is part of renormalization.  The large-time bound is
infrared and uses the boundary spectral gap.  Mid-time convergence is
the genuine regulator comparison.  None of these three regimes can be
inferred from the other two.
\end{assessment}

\subsection{Sharp eta discontinuity at a boundary crossing}
\label{sec:v2v92-eta-jump}

\begin{counterexample}[Norm convergence without an eta limit]
\label{sep:v2v92-eta-jump}
On \(\mathbb C\), let
\[
 A_\varepsilon=(\varepsilon).
\]
For \(\varepsilon\ne0\),
\begin{equation}
 \eta_{A_\varepsilon}(0)
 =
 \operatorname{sgn}\varepsilon.
\label{eq:v2v92-sign-eta}
\end{equation}
Although \(A_\varepsilon\to0\) in operator norm, the eta invariant has
different limits from the two sides and is not continuous at zero.
\end{counterexample}

\begin{proof}
The eta series of a one-dimensional nonzero operator is
\[
 \operatorname{sgn}(\varepsilon)|\varepsilon|^{-s},
\]
whose value at \(s=0\) is
\(\operatorname{sgn}\varepsilon\).
\end{proof}

\begin{firewall}[Spectral flow must be recorded]
\label{fw:v2v92-spectral-flow}
If the boundary gap closes, eta can jump by spectral flow.  One may still
obtain a continuous determinant-line section after combining eta with
kernel orientation and inflow data, but the scalar eta invariant alone
does not satisfy
\eqref{eq:v2v92-eta-convergence}.
\end{firewall}

\subsection{Connection to the V91 cancellation}
\label{sec:v2v92-V91}

\begin{theorem}[What local bulk cancellation leaves]
\label{thm:v2v92-after-V91}
Assume the regulator's perturbative bulk class is identified with the
standard descent basis of V91.  Then its bulk component vanishes for the
Standard Model representation.  The remaining anomaly problem consists
of:
\begin{enumerate}
 \item the pure boundary quotient
 \(H^0(\mathcal E)/r^*H^0(\mathcal B)\);
 \item global determinant-line holonomy;
 \item convergence of the eta/inflow representative under the analytic
 hypotheses of Theorem~\ref{thm:v2v92-eta}.
\end{enumerate}
No additional perturbative bulk triangle coefficient appears.
\end{theorem}

\begin{proof}
Theorem~\ref{thm:v2v91-local} kills the bulk descent coefficients.
Corollary~\ref{cor:v2v92-boundary-class} identifies the remaining
relative class, while
Theorem~\ref{thm:v2v92-eta} gives the stated sufficient analytic
continuum criterion.
\end{proof}

\begin{firewall}[Strong ellipticity and domain control]
\label{fw:v2v92-ellipticity}
The heat expansion and trace bounds above require a self-adjoint,
strongly elliptic boundary problem with a controlled domain, or a
separately proved replacement.  Formal boundary conditions which fail
strong ellipticity cannot be inserted into
Theorem~\ref{thm:v2v92-eta}.
\end{firewall}

\begin{firewall}[Eta is determinant-line data]
\label{fw:v2v92-line}
For a family of chiral Dirac operators, the exponentiated eta invariant
is naturally determinant-line data rather than a globally preferred
complex number.  V92's scalar convergence statement is made only after
a common local trivialization and subtraction convention have been
chosen.
\end{firewall}

\subsection{V92 conclusion and next acquisition}
\label{sec:v2v92-conclusion}

V92 constructs the relative bulk--boundary Ward complex and proves its
long exact sequence.  Bulk cancellation lifts to a full relative
cancellation only when the boundary restriction is exact, and pure
boundary determinant-phase classes can remain.  A uniform relative
Hodge gap controls representatives.  Uniform small-, middle-, and
large-time heat estimates give convergence of the renormalized eta
invariant, while a one-level crossing proves the necessity of recording
spectral flow.

V93 will couple this boundary construction to the Einstein equation.
It will derive the diffeomorphism Ward identity for the renormalized
matter effective action, separate the nonzero trace anomaly from
covariant conservation, and state the operator estimates required for
the resulting stress tensor to enter the V32 causal semiclassical
evolution.

\clearpage
\section*{Second Volume, Version V93: Relative Diffeomorphism Ward Identity and the Einstein Source Packet}
\addcontentsline{toc}{section}{Second Volume, Version V93: Relative Diffeomorphism Ward Identity and the Einstein Source Packet}
\label{sec:v2v93}

\subsection{Purpose and non-duplication}
\label{sec:v2v93-purpose}

V26 proved that covariant conservation is necessary for harmonic Einstein
constraint propagation.  V27 derived the retarded metric response and its
linearized conservation identity, V32 built the causal Neumann mechanism,
and V53 reduced the strong source modulus to local and memory operator
bounds.  Those results assumed the relevant Ward identities.

V93 supplies the missing implication from the anomaly analysis.  It derives
the bulk and boundary diffeomorphism Ward identities from a trivial relative
anomaly class, proves that a Weyl anomaly is compatible with covariant
conservation, and identifies the correct Ward-completed polarization packet
on a non-vacuum background.  The new point is that the metric derivative of
a conserved source is not separately transverse when the background source
is nonzero: a connection-variation term must remain attached to it.

\subsection{Effective-action conventions}
\label{sec:v2v93-conventions}

Let \(M_T\) be an oriented Lorentzian slab with timelike boundary
\(\mathcal T\).  Write \(\gamma\) for the induced metric, \(n\) for the
outward unit normal, and \(D\) for the Levi--Civita connection of \(\gamma\).
Let \(A\) be a background gauge connection and \(\phi\) collect background
matter fields.  For a renormalized effective action \(\Gamma_n\), define its
bulk variations by
\begin{equation}
 \delta\Gamma_n^{\rm bulk}
 =
 \int_{M_T}\!\sqrt{|g|}
 \left(
  \frac12 T_n^{\mu\nu}\delta g_{\mu\nu}
  +J_{n,A}^{\mu}\delta A^A_\mu
  +E_{n,i}\delta\phi^i
 \right),
\label{eq:v2v93-bulk-variation}
\end{equation}
and its intrinsic boundary variations by
\begin{equation}
 \delta\Gamma_n^{\partial}
 =
 \int_{\mathcal T}\!\sqrt{|\gamma|}
 \left(
  \frac12\tau_n^{ab}\delta\gamma_{ab}
  +j_{n,A}^{a}\delta a^A_a
  +e_{n,i}\delta\varphi^i
 \right).
\label{eq:v2v93-boundary-variation}
\end{equation}
The repeated use of \(a\) as a boundary tensor index and as a gauge index is
avoided below by writing gauge indices as \(A,B\).

Compose an infinitesimal diffeomorphism generated by \(X\) with the gauge
transformation of parameter \(-\iota_XA\).  Its gauge-covariant action is
\begin{align}
 \widehat\delta_X g_{\mu\nu}
 &=2\nabla_{(\mu}X_{\nu)},
 &
 \widehat\delta_X A^A_\mu
 &=X^\rho F^A_{\rho\mu},
 &
 \widehat\delta_X\phi^i
 &=X^\rho D_\rho\phi^i.
\label{eq:v2v93-covariant-diffeomorphism}
\end{align}
This replacement of the ordinary Lie derivative is legitimate only after
the gauge Ward identity has been restored.  V91 supplies the perturbative
bulk cancellation; V92 shows that the boundary component is an additional
relative class.

\subsection{From the relative class to local Ward identities}
\label{sec:v2v93-relative-Ward}

\begin{theorem}[Bulk and tangential boundary Ward identities]
\label{thm:v2v93-relative-Ward}
Assume the relative diffeomorphism anomaly of \(\Gamma_n\) is exact in the
V92 mapping cone and has been removed by an allowed relative counterterm.
Assume also that the gauge Ward identity used in
\eqref{eq:v2v93-covariant-diffeomorphism} has been restored.  Then invariance
under every smooth \(X\) tangent to \(\mathcal T\) implies, distributionally,
\begin{equation}
 \nabla_\mu T_n^{\mu}{}_{\nu}
 =
 F^A_{\nu\mu}J_{n,A}^{\mu}
 +E_{n,i}D_\nu\phi^i
 \qquad\hbox{in }M_T^\circ,
\label{eq:v2v93-bulk-Ward}
\end{equation}
and
\begin{equation}
 D_a\tau_n^{a}{}_{b}
 =
 n_\mu T_n^{\mu}{}_{b}
 +f^A_{ba}j_{n,A}^{a}
 +e_{n,i}D_b\varphi^i
 \qquad\hbox{on }\mathcal T.
\label{eq:v2v93-boundary-Ward}
\end{equation}
Here \(f\) is the curvature of the induced boundary connection and
\(T^{\mu}{}_{b}=T^{\mu\nu}\gamma_{\nu b}\).
\end{theorem}

\begin{proof}
For \(X\) compactly supported in the interior, insert
\eqref{eq:v2v93-covariant-diffeomorphism} into
\eqref{eq:v2v93-bulk-variation}.  Symmetry of \(T_n\) and integration by
parts give
\[
 0
 =
 \int_{M_T}\sqrt{|g|}\,X^\nu
 \left(
  -\nabla_\mu T_n^{\mu}{}_{\nu}
  +F^A_{\nu\mu}J_{n,A}^{\mu}
  +E_{n,i}D_\nu\phi^i
 \right).
\]
The fundamental lemma for distributions proves
\eqref{eq:v2v93-bulk-Ward}.

Now allow arbitrary tangent \(X\) at \(\mathcal T\).  The bulk metric
integration by parts contributes
\[
 \int_{\mathcal T}\sqrt{|\gamma|}\,
 n_\mu T_n^{\mu}{}_{b}X^b.
\]
The boundary metric term is
\[
 \int_{\mathcal T}\sqrt{|\gamma|}\,
 \tau_n^{ab}D_aX_b
 =
 -\int_{\mathcal T}\sqrt{|\gamma|}\,
 (D_a\tau_n^{a}{}_{b})X^b,
\]
where corner terms vanish by fixing \(X\) near the initial and final
corners.  The boundary gauge and matter variations contribute respectively
\[
 \int_{\mathcal T}\sqrt{|\gamma|}\,
 X^b f^A_{ba}j_{n,A}^{a},
 \qquad
 \int_{\mathcal T}\sqrt{|\gamma|}\,
 X^b e_{n,i}D_b\varphi^i.
\]
The already proved bulk identity removes the interior integral.  Since the
tangential trace of \(X\) is arbitrary, the remaining boundary coefficient
vanishes, which is \eqref{eq:v2v93-boundary-Ward}.
\end{proof}

\begin{corollary}[When the metric source is conserved]
\label{cor:v2v93-conserved-source}
If all gauge and matter variables appearing in \(\Gamma_n\) are dynamical
and satisfy their effective equations, then
\begin{equation}
 \nabla_\mu T_n^{\mu}{}_{\nu}=0.
\label{eq:v2v93-conserved-source}
\end{equation}
If \(A\) is instead an external background, the metric stress alone has the
Lorentz-force divergence \(F^A_{\nu\mu}J_{n,A}^{\mu}\); it is not a conserved
Einstein source until the stress of the gauge background is included.
\end{corollary}

\begin{proof}
For dynamical variables, stationarity gives \(J_{n,A}^{\mu}=0\) and
\(E_{n,i}=0\) for the complete effective action.  Apply
\eqref{eq:v2v93-bulk-Ward}.  The external-field statement is the same
identity without imposing the gauge equation.
\end{proof}

\begin{firewall}[The unresolved V92 boundary class is an Einstein obstruction]
\label{fw:v2v93-boundary-obstruction}
The V91 bulk cancellation alone does not imply
\eqref{eq:v2v93-boundary-Ward}.  A nonzero class in
\(H^0(\mathcal E)/r^*H^0(\mathcal B)\), determinant-line holonomy, or eta
spectral flow leaves a boundary Ward defect.  Such a defect must be cancelled
by boundary degrees of freedom or included as physical boundary flux before
the harmonic constraint theorem of V26 can be used on a slab with boundary.
\end{firewall}

\subsection{Weyl anomaly versus diffeomorphism conservation}
\label{sec:v2v93-Weyl}

Let \(\mathcal D_X\) denote the gauge-covariant diffeomorphism generator and
\(\mathcal W_\sigma\) the Weyl generator with the assigned Weyl weights of
all background fields.  Suppose
\begin{equation}
 \mathcal D_X\Gamma_n=0,
 \qquad
 \mathcal W_\sigma\Gamma_n
 =
 \int_{M_T}\sqrt{|g|}\,\sigma\mathcal A_{W,n}
 +
 \int_{\mathcal T}\sqrt{|\gamma|}\,\sigma\mathcal A_{W,n}^{\partial}.
\label{eq:v2v93-Weyl-anomaly}
\end{equation}
In four bulk dimensions, the on-shell local bulk density has the schematic
renormalized form
\begin{equation}
 \mathcal A_{W,n}
 =
 c_n C_{\mu\nu\rho\sigma}C^{\mu\nu\rho\sigma}
 -a_n E_4
 +b_n\Box R
 +\sum_I\beta_{I,n}\mathcal O_{I,n}
 +\mathcal A_{\rm mass,n},
\label{eq:v2v93-Weyl-basis}
\end{equation}
with convention-dependent normalizations.  The \(\Box R\) coefficient can
be shifted by a finite \(R^2\) counterterm.  Equation
\eqref{eq:v2v93-Weyl-basis} records the allowed structure; it does not assert
values of the interacting Standard Model coefficients.

\begin{theorem}[Trace anomaly does not violate the Bianchi identity]
\label{thm:v2v93-trace-conservation}
Under \eqref{eq:v2v93-Weyl-anomaly}, the trace Ward identity may be anomalous
while the diffeomorphism identity
\eqref{eq:v2v93-bulk-Ward} remains exact.  On the complete matter equations,
\begin{equation}
 T_n^{\mu}{}_{\mu}=\mathcal A_{W,n},
 \qquad
 \nabla_\mu T_n^{\mu}{}_{\nu}=0
\label{eq:v2v93-two-Wards}
\end{equation}
are compatible statements.
\end{theorem}

\begin{proof}
The first equation is the coefficient of an arbitrary Weyl parameter
\(\sigma\) in \(\mathcal W_\sigma\Gamma_n\), after the matter equations and
explicit weight terms have been imposed.  The second is the coefficient of
an arbitrary vector field in \(\mathcal D_X\Gamma_n\), proved in
Theorem~\ref{thm:v2v93-relative-Ward}.  These are different variations of
\(\Gamma_n\).  Their consistency condition is
\begin{equation}
 [\mathcal D_X,\mathcal W_\sigma]\Gamma_n
 =
 \mathcal W_{\mathcal L_X\sigma}\Gamma_n,
\label{eq:v2v93-WZ}
\end{equation}
which holds because \(\mathcal A_{W,n}\) and the measure form a covariant
scalar density, with the boundary density transforming intrinsically.
Thus a nonzero scalar trace does not produce a vector-valued divergence.
\end{proof}

\begin{assessment}[Where the trace anomaly enters V32]
\label{ass:v2v93-trace-V32}
The metric variation of the anomaly functional contributes to the local
contact operator and, after integration of the anomaly, possibly to a
causal nonlocal scalar response.  It therefore changes the trace-sector
polarization and the constants in the V32 norm estimate.  It does not relax
the conservation constraint and cannot be discarded merely because the TT
sector is being estimated.
\end{assessment}

\subsection{The Ward-completed linear response on a sourced background}
\label{sec:v2v93-linear-packet}

Let
\[
 \mathcal B_gS:=\nabla_g^\mu S_{\mu\nu}
\]
be the divergence operator on symmetric source tensors.  Let
\(\Theta_n(g)\) be the complete renormalized Standard Model mean source and
assume
\begin{equation}
 \mathcal B_g\Theta_n(g)=0.
\label{eq:v2v93-source-Ward}
\end{equation}
Define
\begin{equation}
 P_{n,g}h:=D\Theta_n[g](h),
 \qquad
 C_{n,g}h:=(D\mathcal B)_g(h)\Theta_n(g).
\label{eq:v2v93-P-C}
\end{equation}

\begin{proposition}[Affine linearized Ward identity]
\label{prop:v2v93-affine-Ward}
For every admissible metric variation \(h\),
\begin{equation}
 \mathcal B_gP_{n,g}h+C_{n,g}h=0.
\label{eq:v2v93-affine-Ward}
\end{equation}
Consequently \(P_{n,g}\) is separately transverse only when
\(C_{n,g}=0\), as happens after the background mean source has been absorbed
to zero in the flat-vacuum normalization.
\end{proposition}

\begin{proof}
Differentiate \eqref{eq:v2v93-source-Ward} in the direction \(h\).  The
product rule gives precisely \eqref{eq:v2v93-affine-Ward}.  If
\(\Theta_n(g)=0\), then the definition in \eqref{eq:v2v93-P-C} gives
\(C_{n,g}=0\).  For a general nonzero source no such conclusion follows.
\end{proof}

\begin{counterexample}[The connection term cannot be dropped algebraically]
\label{sep:v2v93-connection-term}
Let \(X=Y=Z=\mathbb R\), let \(\mathcal B=1\), \(P=1\), and \(C=-1\).
Then
\[
 \mathcal BP+C=0,
 \qquad
 \mathcal BP=1\ne0.
\]
Thus the differentiated Ward identity can hold exactly although the response
operator is not transverse by itself.  Estimating or projecting \(P\) before
attaching \(C\) changes the constraint equation.
\end{counterexample}

Let the covariant gravitational tensor, including whichever fixed local
curvature terms are retained, be \(\mathscr G(g)\), and suppose
\begin{equation}
 \mathcal B_g\mathscr G(g)=0,
 \qquad
 \mathscr G(g)=\kappa\Theta_n(g).
\label{eq:v2v93-background-Einstein}
\end{equation}
Write \(L_gh=D\mathscr G[g](h)\).

\begin{theorem}[Ward closure of the coupled linearized Einstein operator]
\label{thm:v2v93-coupled-linear-Ward}
At a background satisfying \eqref{eq:v2v93-background-Einstein},
\begin{equation}
 \mathcal B_g\bigl(L_g-\kappa P_{n,g}\bigr)h=0
 \qquad\hbox{for every }h.
\label{eq:v2v93-coupled-transverse}
\end{equation}
Hence an inhomogeneous linearized equation
\begin{equation}
 (L_g-\kappa P_{n,g})h=\kappa S
\label{eq:v2v93-linear-Einstein}
\end{equation}
is Bianchi compatible only if \(\mathcal B_gS=0\), together with the
appropriate boundary flux identity.
\end{theorem}

\begin{proof}
Differentiate the geometric identity
\(\mathcal B_g\mathscr G(g)=0\):
\[
 \mathcal B_gL_gh
 +(D\mathcal B)_g(h)\mathscr G(g)=0.
\]
Use the background equation in
\eqref{eq:v2v93-background-Einstein} to identify the second term with
\(\kappa C_{n,g}h\).  Proposition~\ref{prop:v2v93-affine-Ward} gives
\(\mathcal B_gP_{n,g}h=-C_{n,g}h\).  Subtraction yields
\eqref{eq:v2v93-coupled-transverse}.  Applying \(\mathcal B_g\) to
\eqref{eq:v2v93-linear-Einstein} proves the compatibility condition.
\end{proof}

\subsection{A closed Ward graph for the V32 iteration}
\label{sec:v2v93-Ward-graph}

Let \(X_T,Y_T,Z_T\) be Banach spaces on the slab such that
\(P_{n,g}:X_T\to Y_T\),
\(C_{n,g}:X_T\to Z_T\), and
\(\mathcal B_g:Y_T\to Z_T\) are bounded.  Define the closed Ward graph
\begin{equation}
 \mathfrak Y_g
 =
 \{(S,z)\in Y_T\times Z_T:\mathcal B_gS+z=0\}
\label{eq:v2v93-Ward-graph}
\end{equation}
and the completed response
\begin{equation}
 \mathfrak P_{n,g}h
 =
 (P_{n,g}h,C_{n,g}h).
\label{eq:v2v93-completed-response}
\end{equation}

\begin{lemma}[The completed response lands in the Ward graph]
\label{lem:v2v93-Ward-graph}
Under \eqref{eq:v2v93-affine-Ward},
\begin{equation}
 \mathfrak P_{n,g}:X_T\longrightarrow\mathfrak Y_g,
 \qquad
 \|\mathfrak P_{n,g}\|
 \le \|P_{n,g}\|+\|C_{n,g}\|
\label{eq:v2v93-Ward-graph-bound}
\end{equation}
for the sum norm on \(Y_T\times Z_T\).
\end{lemma}

\begin{proof}
The graph condition is exactly
\eqref{eq:v2v93-affine-Ward}.  The norm bound is the triangle inequality.
Closedness follows because \((S,z)\mapsto\mathcal B_gS+z\) is continuous.
\end{proof}

\begin{theorem}[Ward-compatible causal Neumann criterion]
\label{thm:v2v93-Ward-Neumann}
Assume a retarded augmented Einstein solver
\begin{equation}
 \widehat E_g:\mathfrak Y_g\longrightarrow X_T
\label{eq:v2v93-augmented-solver}
\end{equation}
has been constructed with the initial and boundary constraints built into
its domain.  Put
\begin{equation}
 K_{n,g}=\kappa\widehat E_g\mathfrak P_{n,g}.
\label{eq:v2v93-K}
\end{equation}
If
\begin{equation}
 \|K_{n,g}\|_{X_T\to X_T}\le\eta<1
\label{eq:v2v93-small}
\end{equation}
uniformly on the chosen background class, then for every conserved external
source \(S\), represented by \((S,0)\in\mathfrak Y_g\), the equation
\begin{equation}
 h+K_{n,g}h=\kappa\widehat E_g(S,0)
\label{eq:v2v93-augmented-equation}
\end{equation}
has a unique causal solution.  Its order-\(N\) truncation obeys
\begin{equation}
 \|h-h^{[N]}\|_{X_T}
 \le
 \frac{\eta^{N+1}}{1-\eta}
 \kappa\|\widehat E_g\|\,\|(S,0)\|_{\mathfrak Y_g}.
\label{eq:v2v93-augmented-error}
\end{equation}
Every response insertion remains in the Ward graph.
\end{theorem}

\begin{proof}
Lemma~\ref{lem:v2v93-Ward-graph} makes every composition in
\eqref{eq:v2v93-K} well defined.  Apply the geometric-series proof of
Theorems~\ref{thm:v2v32-Neumann} and~\ref{thm:v2v32-error} to
\(K_{n,g}\).  Retardedness is preserved under composition, finite sums, and
norm limits.  The range assertion follows at each insertion from
\(\mathfrak P_{n,g}X_T\subset\mathfrak Y_g\).
\end{proof}

\begin{firewall}[A flat TT projector is insufficient on a sourced geometry]
\label{fw:v2v93-flat-projector}
On a nonzero curved background, replacing
\(\mathfrak P_{n,g}=(P_{n,g},C_{n,g})\) by a flat transverse projection of
\(P_{n,g}\) discards the second term in
\eqref{eq:v2v93-affine-Ward}.  The resulting Neumann iterates need not
propagate the harmonic constraints even if their TT multiplier norm is less
than one.
\end{firewall}

\subsection{Regulator passage of the affine Ward identity}
\label{sec:v2v93-regulator}

\begin{proposition}[Norm-stable Ward packet]
\label{prop:v2v93-regulator-Ward}
Let
\(\mathcal B_n\to\mathcal B\),
\(P_n\to P\), and \(C_n\to C\) in the corresponding operator norms, with
\(\sup_n\|\mathcal B_n\|<\infty\).  If
\begin{equation}
 \|\mathcal B_nP_n+C_n\|_{X_T\to Z_T}
 \le\varepsilon_n,
 \qquad
 \varepsilon_n\longrightarrow0,
\label{eq:v2v93-Ward-defect}
\end{equation}
then
\begin{equation}
 \mathcal BP+C=0.
\label{eq:v2v93-limit-Ward}
\end{equation}
More precisely,
\begin{align}
 \|\mathcal BP+C\|
 \le{}&
 \|\mathcal B-\mathcal B_n\|\,\|P\|
 +\|\mathcal B_n\|\,\|P-P_n\|
 +\|C-C_n\|
 +\varepsilon_n.
\label{eq:v2v93-limit-Ward-rate}
\end{align}
\end{proposition}

\begin{proof}
Insert and subtract \(\mathcal B_nP\), \(\mathcal B_nP_n\), and \(C_n\),
then use the triangle inequality and
\eqref{eq:v2v93-Ward-defect}.  Every term on the right of
\eqref{eq:v2v93-limit-Ward-rate} tends to zero.
\end{proof}

\begin{assessment}[What V88 and V89 now control]
\label{ass:v2v93-V88-V89}
The V88 regulator-stable BRST diagram and V89 relative counterterm complex
supply a route to the defect bound
\eqref{eq:v2v93-Ward-defect}.  V93 shows the additional data that must be
carried into gravity: convergence of the divergence operator under metric
identification and convergence of the connection term \(C_n\).  Cohomology
alone does not provide either operator norm.
\end{assessment}

\subsection{V93 conclusion and next acquisition}
\label{sec:v2v93-conclusion}

V93 derives the local bulk and tangential boundary Ward identities from a
trivial relative anomaly class.  It proves that the nonzero trace anomaly is
compatible with diffeomorphism conservation.  On a sourced curved
background, the correct linear response is the pair
\((P_{n,g}h,C_{n,g}h)\), not \(P_{n,g}h\) alone.  This pair forms a closed
Ward graph, is stable under regulator convergence, and can be inserted into
the V32 causal Neumann construction once an augmented retarded Einstein
solver and its uniform norm bound are proved.

V94 will compute \((D\mathcal B)_g(h)\Theta(g)\) explicitly, prove its
Sobolev mapping estimate on a bounded-geometry slab, and determine whether
this connection block is lower order relative to the four-derivative stress
polarization.  That estimate is the next input needed for
\eqref{eq:v2v93-small}.

\clearpage
\section*{Second Volume, Version V94: Explicit Connection-Variation Source and Sobolev Order}
\addcontentsline{toc}{section}{Second Volume, Version V94: Explicit Connection-Variation Source and Sobolev Order}
\label{sec:v2v94}

\subsection{Purpose}
\label{sec:v2v94-purpose}

V93 showed that the linearized Ward identity on a sourced background is
\[
 \mathcal B_gP_gh+C_gh=0,
 \qquad
 C_gh=(D\mathcal B)_g(h)\Theta(g),
\]
and that \(C_g\) must remain attached to the stress polarization.  V94
computes \(C_g\) exactly for a symmetric covariant source, proves its
Sobolev mapping and regulator-stability estimates, and checks the formula
on a cosmological source.  The result places \(C_g\) two spatial Sobolev
orders above the four-derivative stress response in the V53 scale, provided
the background source has the stated regularity.  Lower order does not by
itself mean small; a retarded short-slab estimate is still required.

\subsection{Variation of the covariant divergence}
\label{sec:v2v94-formula}

Fix a metric \(g\) and a symmetric covariant tensor \(S_{\mu\nu}\).  Use the
common covariant-component chart in which \(S_{\mu\nu}\) is held fixed while
the divergence operator is varied.  Thus
\begin{equation}
 (\mathcal B_gS)_\nu
 =g^{\mu\rho}\nabla_\rho S_{\mu\nu}.
\label{eq:v2v94-divergence}
\end{equation}
For \(g_\varepsilon=g+\varepsilon h\), put
\begin{equation}
 (C_{g,S}h)_\nu
 =
 \left.\frac{\mathrm d}{\mathrm d\varepsilon}\right|_{\varepsilon=0}
 (\mathcal B_{g_\varepsilon}S)_\nu.
\label{eq:v2v94-C-definition}
\end{equation}

\begin{theorem}[Exact connection-variation formula]
\label{thm:v2v94-C-formula}
Let \(h_{\mu\nu}\) and \(S_{\mu\nu}\) be symmetric.  Then
\begin{align}
 (C_{g,S}h)_\nu
 ={}&
 -h^{\mu\rho}\nabla_\rho S_{\mu\nu}
 -S_{\lambda\nu}
 \left(
  \nabla_\mu h^{\lambda\mu}
  -\frac12\nabla^\lambda\operatorname{tr}_g h
 \right)
 \nonumber\\
 &
 -\frac12 S^{\rho\lambda}\nabla_\nu h_{\rho\lambda}.
\label{eq:v2v94-C-formula}
\end{align}
In particular, \(C_{g,S}\) contains at most one derivative of \(h\) and one
derivative of the fixed coefficient \(S\).
\end{theorem}

\begin{proof}
The elementary metric variations are
\begin{align}
 \delta g^{\mu\rho}&=-h^{\mu\rho},
\label{eq:v2v94-inverse-variation}\\
 \delta\Gamma^\lambda_{\rho\mu}
 &=
 \frac12
 \left(
  \nabla_\rho h^\lambda{}_{\mu}
  +\nabla_\mu h^\lambda{}_{\rho}
  -\nabla^\lambda h_{\rho\mu}
 \right).
\label{eq:v2v94-Christoffel-variation}
\end{align}
Holding the covariant components of \(S\) fixed gives
\begin{align}
 \delta(\mathcal B_gS)_\nu
 ={}&
 -h^{\mu\rho}\nabla_\rho S_{\mu\nu}
 -g^{\mu\rho}\delta\Gamma^\lambda_{\rho\mu}S_{\lambda\nu}
 -g^{\mu\rho}\delta\Gamma^\lambda_{\rho\nu}S_{\mu\lambda}.
\label{eq:v2v94-unsimplified}
\end{align}
Contracting \eqref{eq:v2v94-Christoffel-variation} in its two lower
indices gives
\begin{equation}
 g^{\mu\rho}\delta\Gamma^\lambda_{\rho\mu}
 =
 \nabla_\mu h^{\lambda\mu}
 -\frac12\nabla^\lambda\operatorname{tr}_g h.
\label{eq:v2v94-first-contraction}
\end{equation}
For the other contraction, symmetry of \(S\) gives
\begin{align}
 g^{\mu\rho}S_{\mu\lambda}
 \delta\Gamma^\lambda_{\rho\nu}
 &=
 \frac12S^{\rho\lambda}
 \left(
  \nabla_\rho h_{\lambda\nu}
  +\nabla_\nu h_{\lambda\rho}
  -\nabla_\lambda h_{\rho\nu}
 \right)
 \nonumber\\
 &=
 \frac12S^{\rho\lambda}\nabla_\nu h_{\rho\lambda}.
\label{eq:v2v94-second-contraction}
\end{align}
Indeed, the first and third terms cancel after interchanging
\(\rho\) and \(\lambda\).  Substitution into
\eqref{eq:v2v94-unsimplified} proves
\eqref{eq:v2v94-C-formula}.
\end{proof}

\begin{remark}[Chart dependence is accounted for in the response]
\label{rem:v2v94-chart}
If the common tensor chart transports contravariant or mixed components
instead, \(C_{g,S}\) changes by an algebraic \(hS\) term.  The derivative
\(P_g=D\Theta[g]\) changes by the opposite chart term, so the invariant
combination \(\mathcal B_gP_g+(D\mathcal B)_g\Theta(g)\) is unchanged.
Equation \eqref{eq:v2v94-C-formula} fixes the covariant-component convention
used in the estimates below.
\end{remark}

\subsection{Cosmological-source sign test}
\label{sec:v2v94-cosmological}

\begin{proposition}[Covariantly constant source need not give \(C=0\)]
\label{prop:v2v94-cosmological}
Let \(S_{\mu\nu}=\lambda g_{\mu\nu}\) with constant \(\lambda\), while the
covariant components of \(S\) are held fixed in
\eqref{eq:v2v94-C-definition}.  Then
\begin{equation}
 (C_{g,\lambda g}h)_\nu
 =
 -\lambda\nabla_\mu h^\mu{}_{\nu}.
\label{eq:v2v94-cosmological-C}
\end{equation}
If the source map is \(\Theta(g)=\lambda g\), then
\begin{equation}
 P_gh=\lambda h,
 \qquad
 \mathcal B_gP_gh+C_{g,\lambda g}h=0.
\label{eq:v2v94-cosmological-Ward}
\end{equation}
\end{proposition}

\begin{proof}
In \eqref{eq:v2v94-C-formula}, \(\nabla S=0\).  The two trace-gradient
terms cancel, leaving
\(-\lambda\nabla_\mu h^\mu{}_{\nu}\), which is
\eqref{eq:v2v94-cosmological-C}.  Since
\(D(\lambda g)[h]=\lambda h\), its divergence is the opposite term.
\end{proof}

\begin{corollary}[Moving vacuum energy preserves the Ward packet]
\label{cor:v2v94-vacuum-transfer}
Moving \(\kappa\lambda g\) from the matter source to the cosmological term
on the geometric side moves \(\kappa\lambda h\) from \(P_g\) to the
linearized geometric operator.  Equation
\eqref{eq:v2v94-cosmological-Ward} proves that the coupled transverse
operator of Theorem~\ref{thm:v2v93-coupled-linear-Ward} is unchanged.
\end{corollary}

\begin{proof}
Both the background equation and its metric derivative change by the same
term with opposite placement.  The divergence of the transferred response
is exactly cancelled by \eqref{eq:v2v94-cosmological-C}.
\end{proof}

\subsection{Slice Sobolev estimate}
\label{sec:v2v94-Sobolev}

Let the spatial slices have dimension \(d\), and assume a bounded-geometry
atlas in which the metric, inverse metric, and connection coefficients have
uniform Sobolev multiplier bounds.  For an integer
\begin{equation}
 s>\frac d2+2,
\label{eq:v2v94-s}
\end{equation}
put
\begin{align}
 \|h\|_{X_s}
 &=
 \|h\|_{C(I;H^s_x)}
 +\|\partial_t h\|_{C(I;H^{s-1}_x)},
\label{eq:v2v94-Xs}\\
 \|S\|_{\mathcal S_{s-1}}
 &=
 \|S\|_{L^\infty(I;H^{s-1}_x)}
 +\|\nabla S\|_{L^\infty(I;H^{s-2}_x)}.
\label{eq:v2v94-Ss}
\end{align}
The second summand is displayed to record the differentiated coefficient;
on a fixed bounded-geometry slab it is controlled by the first plus the
corresponding time regularity.

\begin{theorem}[Connection block loses at most two spatial derivatives]
\label{thm:v2v94-Sobolev}
Under \eqref{eq:v2v94-s}, there is a constant
\(K_{s,g}\), uniform on a bounded-geometry metric ball, such that
\begin{equation}
 \|C_{g,S}h\|_{L^1(I;H^{s-2}_x)}
 \le
 T K_{s,g}
 \|S\|_{\mathcal S_{s-1}}
 \|h\|_{X_s}.
\label{eq:v2v94-Sobolev}
\end{equation}
If instead \(S\) is uniformly in \(H^s_x\) with one spatial covariant
derivative in \(H^{s-1}_x\), then
\begin{equation}
 \|C_{g,S}h\|_{L^1(I;H^{s-1}_x)}
 \le
 T K'_{s,g}
 \|S\|_{\mathcal S_s}
 \|h\|_{X_s}.
\label{eq:v2v94-Sobolev-strong}
\end{equation}
\end{theorem}

\begin{proof}
Because \(s-1>d/2+1\), the standard Sobolev multiplication estimates on
each bounded-geometry chart give
\begin{align*}
 \|h\,\nabla S\|_{H^{s-2}_x}
 &\le K\|h\|_{H^s_x}\|\nabla S\|_{H^{s-2}_x},\\
 \|S\,\nabla h\|_{H^{s-2}_x}
 &\le K\|S\|_{H^{s-1}_x}\|h\|_{H^s_x}.
\end{align*}
The inverse metrics and connection coefficients appearing in
\eqref{eq:v2v94-C-formula} are uniform multipliers on the metric ball.
Summing over finitely many charts and integrating the pointwise bound in
time proves \eqref{eq:v2v94-Sobolev}.  With one more derivative on \(S\),
the same multiplication theorem places both products in
\(H^{s-1}_x\), proving \eqref{eq:v2v94-Sobolev-strong}.
\end{proof}

\begin{corollary}[Relative order in the V53 source scale]
\label{cor:v2v94-relative-order}
In the V53 scale \(P_g:X_s\to Y_{s-4}\), the connection block factors as
\begin{equation}
 X_s
 \xrightarrow{\ C_{g,S}\ }
 L^1(I;H^{s-2}_x)
 \hookrightarrow
 Y_{s-4}
\label{eq:v2v94-factorization}
\end{equation}
whenever \(Y_{s-4}=L^1(I;H^{s-4}_x)\), or whenever that space is
continuously dominated by the target used in V53.  Thus \(C_{g,S}\) is two
Sobolev orders lower than the worst four-derivative stress response.
\end{corollary}

\begin{proof}
Apply Theorem~\ref{thm:v2v94-Sobolev} and the continuous embedding
\(H^{s-2}_x\hookrightarrow H^{s-4}_x\).
\end{proof}

\begin{firewall}[Lower order is not automatic smallness]
\label{fw:v2v94-lower-order}
The factor \(T\) in \eqref{eq:v2v94-Sobolev} is useful only if the augmented
retarded solver is estimated in compatible time-integrated norms.  Large
background stress, degeneration of the bounded-geometry constants, or a
boundary trace loss can defeat the V93 smallness condition even though
\(C_{g,S}\) has lower differential order.
\end{firewall}

\subsection{Regulator stability}
\label{sec:v2v94-regulator}

Let \((g_n,S_n)\) lie in a common bounded-geometry ball, and use fixed
background identifications of their tensor bundles and Sobolev spaces.

\begin{theorem}[Operator-norm continuity of the connection block]
\label{thm:v2v94-regulator}
Suppose
\begin{equation}
 \|g_n-g\|_{X_s}
 +\|S_n-S\|_{\mathcal S_{s-1}}
 \longrightarrow0,
\label{eq:v2v94-data-convergence}
\end{equation}
and the metrics retain a common nondegeneracy margin.  Then
\begin{equation}
 \|C_{g_n,S_n}-C_{g,S}\|_{
 \mathcal L(X_s,L^1(I;H^{s-2}_x))}
 \longrightarrow0.
\label{eq:v2v94-C-convergence}
\end{equation}
On a sufficiently small common ball there is a constant \(K\) for which
\begin{align}
 \|C_{g_n,S_n}-C_{g,S}\|
 \le
 KT
 \left(
  \|g_n-g\|_{X_s}
  +\|S_n-S\|_{\mathcal S_{s-1}}
 \right).
\label{eq:v2v94-C-rate}
\end{align}
\end{theorem}

\begin{proof}
Formula \eqref{eq:v2v94-C-formula} is a finite sum of contractions of
\(g^{-1}\), the connection of \(g\), \(h\), \(\nabla h\), \(S\), and
\(\nabla S\).  On a uniformly nondegenerate \(H^s\) metric ball, inversion
and the Christoffel map are locally Lipschitz in the corresponding Sobolev
multiplier norms.  Subtract the two formulas term by term.  In each
difference, telescope one factor at a time and apply the product estimates
used in Theorem~\ref{thm:v2v94-Sobolev}.  This gives
\eqref{eq:v2v94-C-rate}, hence \eqref{eq:v2v94-C-convergence}.
\end{proof}

\begin{corollary}[Closure of the V93 regulator hypothesis]
\label{cor:v2v94-V93-regulator}
If, in addition,
\(\mathcal B_{g_n}\to\mathcal B_g\) and
\(P_{n,g_n}\to P_g\) in the V93 operator norms and the finite-regulator Ward
defects tend to zero, then the limiting affine Ward identity
\begin{equation}
 \mathcal B_gP_g+C_{g,\Theta(g)}=0
\label{eq:v2v94-limit-affine-Ward}
\end{equation}
holds as an operator \(X_s\to L^1(I;H^{s-2}_x)\), and hence also in the
weaker V53 source-divergence topology.
\end{corollary}

\begin{proof}
Theorem~\ref{thm:v2v94-regulator} supplies the convergence of \(C_n\)
required in Proposition~\ref{prop:v2v93-regulator-Ward}.  Apply that
proposition and then the continuous Sobolev embedding.
\end{proof}

\subsection{Boundary connection block}
\label{sec:v2v94-boundary}

For the intrinsic boundary divergence \((\mathcal B_\gamma^\partial
\tau)_b=D^a\tau_{ab}\), the same calculation with
\((g,\nabla,h,S)\) replaced by \((\gamma,D,k,\tau)\) gives
\begin{align}
 ((D\mathcal B^\partial)_\gamma(k)\tau)_b
 ={}&
 -k^{ac}D_c\tau_{ab}
 -\tau_{db}
 \left(
  D_ak^{da}-\frac12D^d\operatorname{tr}_\gamma k
 \right)
 -\frac12\tau^{cd}D_bk_{cd}.
\label{eq:v2v94-boundary-C}
\end{align}

\begin{proposition}[Intrinsic boundary estimate]
\label{prop:v2v94-boundary-estimate}
On a uniformly timelike bounded-geometry boundary, the analogues of
\eqref{eq:v2v94-Sobolev} and \eqref{eq:v2v94-C-rate} hold in intrinsic
boundary Sobolev spaces.  To differentiate the full boundary Ward identity
\eqref{eq:v2v93-boundary-Ward}, these terms must be combined with the
variations of \(n_\mu T^\mu{}_b\), the induced connection, and the boundary
matter forces before taking norms.
\end{proposition}

\begin{proof}
Equation \eqref{eq:v2v94-boundary-C} has the same product structure as
\eqref{eq:v2v94-C-formula}; the slice multiplication proof applies in
boundary charts.  Differentiating the right side of
\eqref{eq:v2v93-boundary-Ward} produces the listed normal, projection, and
matter-response terms by the product rule.
\end{proof}

\begin{firewall}[A boundary trace estimate is still missing]
\label{fw:v2v94-boundary-trace}
Bulk control in \(H^s_x\) yields boundary traces only with the usual
half-derivative cost.  V94 does not prove that the complete normal-flux
response and eta-sector boundary stress fit the intrinsic spaces of
Proposition~\ref{prop:v2v94-boundary-estimate}.  That remains a separate
boundary-domain estimate.
\end{firewall}

\subsection{V94 conclusion and next acquisition}
\label{sec:v2v94-conclusion}

V94 computes the connection-variation source exactly.  With a background
source in \(H^{s-1}\), it maps \(X_s\) into
\(L^1H^{s-2}\), two orders above the worst V53 stress-response target
\(L^1H^{s-4}\).  It is locally Lipschitz under joint regulator convergence
of the metric and source.  The cosmological-source calculation proves both
the sign and the invariance of the Ward packet when vacuum energy is moved
between the matter and gravitational sides.

V95 is a mandatory companion audit.  After reading the current Yang--Mills
and Navier--Stokes notes, it will combine any transferable packet lesson
with V94 and formulate the augmented retarded energy estimate needed to
turn lower differential order into an explicit contribution to the V93
Neumann constant.
\clearpage
\section*{Second Volume, Version V95: Companion Audit and Constraint-Propagating Reduced Einstein Iteration}
\addcontentsline{toc}{section}{Second Volume, Version V95: Companion Audit and Constraint-Propagating Reduced Einstein Iteration}
\label{sec:v2v95}

\subsection{Mandatory companion audit}
\label{sec:v2v95-audit}

Before continuing the Einstein estimate, the current companion files in the
shared \texttt{latex\_notes} directory were read without modification.
Their audit fingerprints are:
\begin{center}
\begin{tabular}{lll}
 programme & bytes & SHA--256\\
\hline
Yang--Mills V76 & (1095484) &
\texttt{f5da5bc0e3d9a0174f7e94aa0e212ee9}\\
& &\texttt{205539423539b36d11db910e9c7b72ed}\\
Navier--Stokes V31 & (887537) &
\texttt{f1121b62238ad5491bb7cf03635ea993}\\
& &\texttt{4942edf573ed8faaa9ee79e1d38a6696}
\end{tabular}
\end{center}
The frozen Yang--Mills first volume and the frozen SM first volume were left
untouched.

The transferable Yang--Mills lesson is the exact two-endpoint rule of V76:
each binary minimum pays its own endpoint, and a paired product may not be
enlarged and charged twice before the full source packet is formed.  Its
large nonprefix residual is split into disjoint stocks before order
optimization.  The transferable Navier--Stokes lesson remains the sharp V31
formation ladder: a heat or Duhamel factor cannot create a missing scale
moment; every conversion from a marked or weak source to an unmarked strong
quantity must display its cost.

\begin{assessment}[Type-correct transfer to the Einstein Ward packet]
\label{ass:v2v95-audit-transfer}
The SM connection term \(C_gh\) and response divergence
\(\mathcal B_gP_gh\) are the two entries of one exact Ward packet.  They must
be combined before norms are taken.  Once the geometric Ward identity is
included, the same \(C_gh\) occurs with the opposite coefficient and cancels.
It is therefore incorrect to add \(\|C_g\|\) automatically to the V32
Neumann constant.  The V94 bound is used to pass this cancellation through
the regulator and to control a genuine mismatch; it is not a second charge
on an exact packet.
\end{assessment}

\subsection{Exact coupled Ward cancellation}
\label{sec:v2v95-cancellation}

Let \(X\) be a metric-perturbation space, \(Y\) a symmetric-source space,
and \(Z\) a divergence-defect space.  Fix bounded operators
\begin{equation}
 L:X\to Y,
 \qquad
 P:X\to Y,
 \qquad
 \mathcal B:Y\to Z,
 \qquad
 C:X\to Z.
\label{eq:v2v95-operators}
\end{equation}
Here \(L=D\mathscr G[g]\), \(P=D\Theta[g]\), and
\(C=(D\mathcal B)_g(\,cdot\,)\Theta(g)\) at a background satisfying
\(\mathscr G(g)=\kappa\Theta(g)\).  Allow finite Ward defects
\begin{align}
 d_G&:=\mathcal BL+\kappa C,
\label{eq:v2v95-dG}\\
 d_M&:=\mathcal BP+C.
\label{eq:v2v95-dM}
\end{align}

\begin{theorem}[Connection cancellation before norms]
\label{thm:v2v95-C-cancellation}
The complete linearized Einstein--matter operator
\begin{equation}
 A:=L-\kappa P
\label{eq:v2v95-A}
\end{equation}
satisfies the exact identity
\begin{equation}
 \mathcal BA=d_G-\kappa d_M.
\label{eq:v2v95-complete-defect}
\end{equation}
Consequently, if \(d_G=d_M=0\), then \(\mathcal BA=0\) with no smallness
assumption on \(C\).  In general,
\begin{equation}
 \|\mathcal BA\|_{\mathcal L(X,Z)}
 \le
 \|d_G\|+\kappa\|d_M\|.
\label{eq:v2v95-defect-bound}
\end{equation}
\end{theorem}

\begin{proof}
Use \eqref{eq:v2v95-dG}--\eqref{eq:v2v95-dM}:
\begin{align*}
 \mathcal B(L-\kappa P)
 &=
 (d_G-\kappa C)
 -\kappa(d_M-C)\\
 &=d_G-\kappa d_M.
\end{align*}
The two copies of \(C\) cancel before any norm is taken.  The operator-norm
bound is the triangle inequality applied only after this cancellation.
\end{proof}

\begin{corollary}[Regulator-uniform coupled conservation]
\label{cor:v2v95-regulator-cancellation}
If
\begin{equation}
 \sup_{g\in\mathcal K}
 \bigl(\|d_{G,n,g}\|+\kappa\|d_{M,n,g}\|\bigr)
 \le\epsilon_n,
 \qquad
 \epsilon_n\longrightarrow0,
\label{eq:v2v95-defects-vanish}
\end{equation}
then the complete cutoff operators obey
\begin{equation}
 \sup_{g\in\mathcal K}
 \|\mathcal B_{n,g}(L_{n,g}-\kappa P_{n,g})\|
 \le\epsilon_n.
\label{eq:v2v95-complete-defect-rate}
\end{equation}
If the operators converge in the V93--V94 norms, the limiting coupled
operator is exactly divergence free.
\end{corollary}

\begin{proof}
Apply Theorem~\ref{thm:v2v95-C-cancellation} at each regulator.  The limit
follows from Proposition~\ref{prop:v2v93-regulator-Ward} and
Theorem~\ref{thm:v2v94-regulator}, together with the analogous geometric
identity.
\end{proof}

\begin{firewall}[Separate transverse projection changes the theorem]
\label{fw:v2v95-separate-projection}
Replacing \(P\) by a separately transverse projection before forming
\(L-\kappa P\) generally removes the \(-C\) entry from
\eqref{eq:v2v95-dM}.  The cancellation in
\eqref{eq:v2v95-complete-defect} is then unavailable.  A projected response
is admissible only if the projector is constructed on the full affine Ward
packet or a new defect estimate is proved.
\end{firewall}

\subsection{Harmonic reduction and the constraint wave equation}
\label{sec:v2v95-harmonic}

Let \(H:X\to U\) be the linearized harmonic gauge-constraint operator.
Suppose the reduced geometric operator has the form
\begin{equation}
 L_R=L+QH,
\label{eq:v2v95-reduced-L}
\end{equation}
where \(Q:U\to Y\) is the gauge-restoring differential operator.  Assume
there is a normally hyperbolic constraint operator
\begin{equation}
 W:=\mathcal BQ:U\to Z
\label{eq:v2v95-W}
\end{equation}
with uniqueness for zero Cauchy and admissible homogeneous boundary data.

\begin{theorem}[Constraint equation with Ward defects]
\label{thm:v2v95-constraint-equation}
If \(h\) solves the reduced linearized equation
\begin{equation}
 (L_R-\kappa P)h=\kappa S,
\label{eq:v2v95-reduced-equation}
\end{equation}
then
\begin{equation}
 WHh
 =
 \kappa\mathcal BS
 -(d_G-\kappa d_M)h.
\label{eq:v2v95-constraint-forcing}
\end{equation}
In particular, exact geometric and matter Ward identities, a conserved
external source, and zero constraint data imply
\begin{equation}
 Hh=0.
\label{eq:v2v95-H-zero}
\end{equation}
The reduced solution therefore solves the unreduced linearized Einstein
system.
\end{theorem}

\begin{proof}
Apply \(\mathcal B\) to \eqref{eq:v2v95-reduced-equation}.  By
\eqref{eq:v2v95-reduced-L},
\[
 \mathcal B(L_R-\kappa P)h
 =
 \mathcal B(L-\kappa P)h+\mathcal BQHh.
\]
Insert \eqref{eq:v2v95-complete-defect} and
\eqref{eq:v2v95-W}; rearrangement gives
\eqref{eq:v2v95-constraint-forcing}.  Under the exact hypotheses the right
side vanishes.  Uniqueness for \(W\) with zero data proves
\eqref{eq:v2v95-H-zero}.  Then \(QHh=0\), so reduced and unreduced equations
coincide.
\end{proof}

Assume a retarded energy estimate for the constraint system,
\begin{equation}
 \|u\|_U
 \le C_W\|Wu\|_Z
\label{eq:v2v95-W-energy}
\end{equation}
for zero constraint data and the chosen homogeneous boundary condition.

\begin{corollary}[Quantitative constraint defect]
\label{cor:v2v95-constraint-defect}
Under \eqref{eq:v2v95-W-energy}, every solution of
\eqref{eq:v2v95-reduced-equation} satisfies
\begin{equation}
 \|Hh\|_U
 \le
 C_W
 \left[
  \kappa\|\mathcal BS\|_Z
  +(\|d_G\|+\kappa\|d_M\|)\|h\|_X
 \right].
\label{eq:v2v95-constraint-bound}
\end{equation}
\end{corollary}

\begin{proof}
Apply \eqref{eq:v2v95-W-energy} to
\eqref{eq:v2v95-constraint-forcing} and use
\eqref{eq:v2v95-defect-bound}.
\end{proof}

\subsection{The causal Neumann solution and its constraints}
\label{sec:v2v95-Neumann}

Assume \(L_R\) has a retarded inverse
\begin{equation}
 E_R:Y\to X,
 \qquad
 L_RE_R=I
\label{eq:v2v95-ER}
\end{equation}
on the selected source domain.  Define
\begin{equation}
 K:=-\kappa E_RP.
\label{eq:v2v95-K}
\end{equation}
Then \eqref{eq:v2v95-reduced-equation} is equivalent to
\begin{equation}
 h+Kh=\kappa E_RS.
\label{eq:v2v95-Neumann-equation}
\end{equation}

\begin{theorem}[Constraint-propagating V32 criterion]
\label{thm:v2v95-Neumann}
Assume
\begin{equation}
 \|K\|_{X\to X}\le\eta<1.
\label{eq:v2v95-eta}
\end{equation}
Then \eqref{eq:v2v95-Neumann-equation} has the unique causal solution
\begin{equation}
 h=\sum_{j=0}^{\infty}(-K)^j\kappa E_RS,
 \qquad
 \|h\|_X
 \le
 \frac{\kappa\|E_R\|}{1-\eta}\|S\|_Y.
\label{eq:v2v95-Neumann-solution}
\end{equation}
Its harmonic-constraint defect obeys
\begin{align}
 \|Hh\|_U
 \le C_W\biggl[
 &\kappa\|\mathcal BS\|_Z\\
 &+
 \frac{\kappa\|E_R\|}{1-\eta}
 (\|d_G\|+\kappa\|d_M\|)\|S\|_Y
 \biggr].
\label{eq:v2v95-Neumann-constraint}
\end{align}
Thus \(C\) is not an extra summand in \(\eta\); it appears only through a
non-cancelling Ward defect.
\end{theorem}

\begin{proof}
The Neumann series and solution bound are
Theorem~\ref{thm:v2v32-Neumann} applied to \(K\).  Substitute that bound for
\(\|h\|_X\) in \eqref{eq:v2v95-constraint-bound}.  The resulting expression
is \eqref{eq:v2v95-Neumann-constraint}.
\end{proof}

\begin{assessment}[Meaning of the V94 lower-order estimate]
\label{ass:v2v95-C-role}
The estimate for \(C\) in V94 has three precise uses: it proves convergence
of the affine Ward packet, bounds \(d_G\) and \(d_M\) when the two identities
are only approximate, and controls chart or boundary errors.  When both Ward
identities are exact, adding \(\kappa\|E_R\|\|C\|\) to \(\eta\) would charge
the same connection variation twice and weaken the theorem without need.
\end{assessment}

\subsection{Constraint error of a finite order reduction}
\label{sec:v2v95-truncation}

Let
\begin{equation}
 h^{[N]}
 =
 \sum_{j=0}^{N}(-K)^j\kappa E_RS
\label{eq:v2v95-hN}
\end{equation}
and define the reduced differential residual
\begin{equation}
 R_N
 =
 (L_R-\kappa P)h^{[N]}-\kappa S.
\label{eq:v2v95-RN}
\end{equation}

\begin{proposition}[Geometric constraint decay for truncations]
\label{prop:v2v95-truncation}
Suppose \(L_R:X\to Y\), \(\mathcal B:Y\to Z\), and \(H:X\to U\) are
bounded on the truncation domain.  Under exact Ward identities and
\(\mathcal BS=0\),
\begin{align}
 \|Hh^{[N]}\|_U
 &\le C_W\|\mathcal B\|\,\|R_N\|_Y,
\label{eq:v2v95-HN-residual}\\
 \|R_N\|_Y
 &\le
 \kappa\|L_R\|\,\|E_R\|\,
 \eta^{N+1}\|S\|_Y.
\label{eq:v2v95-RN-rate}
\end{align}
Consequently
\begin{equation}
 \|Hh^{[N]}\|_U
 \le
 C_W\|\mathcal B\|\,\|L_R\|
 \kappa\|E_R\|\eta^{N+1}\|S\|_Y.
\label{eq:v2v95-HN-rate}
\end{equation}
Finite Neumann iterates need not satisfy the harmonic constraint exactly,
but their violation decays with the same geometric factor as the equation
residual.
\end{proposition}

\begin{proof}
Apply \(\mathcal B\) to \eqref{eq:v2v95-RN}.  The exact version of
\eqref{eq:v2v95-complete-defect} and conservation of \(S\) give
\[
 \mathcal BR_N=WHh^{[N]}.
\]
The constraint energy estimate proves
\eqref{eq:v2v95-HN-residual}.  The normalized residual in
\eqref{eq:v2v95-Neumann-equation} is, by the telescoping identity of V32,
\[
 r_N=(-1)^NK^{N+1}\kappa E_RS.
\]
Since \(L_Rr_N=R_N\), boundedness and
\eqref{eq:v2v95-eta} give \eqref{eq:v2v95-RN-rate}.  Combine the two
bounds.
\end{proof}

\begin{firewall}[Partial iterates are not separately physical metrics]
\label{fw:v2v95-partial-iterates}
The individual terms of the Neumann series and its finite truncations need
not be constraint-satisfying Einstein perturbations.  The exact infinite
solution is physical by Theorem~\ref{thm:v2v95-constraint-equation}; a finite
order reduction is physical only up to the explicit constraint error
\eqref{eq:v2v95-HN-rate}, unless a separate constraint correction is
applied.
\end{firewall}

\subsection{Boundary Ward forcing}
\label{sec:v2v95-boundary}

On a slab with timelike boundary, let \(d_\partial\) denote the defect in the
linearization of the complete boundary identity
\eqref{eq:v2v93-boundary-Ward}, including the intrinsic connection block
\eqref{eq:v2v94-boundary-C}, the normal-flux variation, and boundary matter
forces.  Assume the constraint energy estimate has the form
\begin{equation}
 \|u\|_U
 \le
 C_W
 \left(
  \|Wu\|_Z+|\mathcal T_Hu\|_{Z_\partial}
 \right),
\label{eq:v2v95-boundary-energy}
\end{equation}
where \(\mathcal T_Hu\) is the incoming constraint trace.

\begin{corollary}[Bulk and boundary defect ledger]
\label{cor:v2v95-boundary-ledger}
If the boundary condition converts the differentiated boundary Ward identity
into
\begin{equation}
 \|\mathcal T_HHh\|_{Z_\partial}
 \le
 \|d_\partial h\|_{Z_\partial}
 +\kappa\|b_\partial\|_{Z_\partial},
\label{eq:v2v95-boundary-defect}
\end{equation}
then the right side of \eqref{eq:v2v95-Neumann-constraint} acquires the
additional term
\begin{equation}
 C_W
 \left(
  \|d_\partial\|\,
  \frac{\kappa\|E_R\|}{1-\eta}\|S\|_Y
  +\kappa\|b_\partial\|_{Z_\partial}
 \right).
\label{eq:v2v95-boundary-ledger}
\end{equation}
Thus bulk anomaly cancellation cannot absorb a boundary Ward defect.
\end{corollary}

\begin{proof}
Use \eqref{eq:v2v95-boundary-defect} in
\eqref{eq:v2v95-boundary-energy}, then apply the solution bound in
\eqref{eq:v2v95-Neumann-solution}.
\end{proof}

\begin{firewall}[The boundary estimate is conditional]
\label{fw:v2v95-boundary-conditional}
V95 does not construct a strongly elliptic Lorentzian boundary condition for
the complete chiral Standard Model, nor prove the incoming constraint trace
estimate \eqref{eq:v2v95-boundary-defect}.  The eta convergence theorem of
V92 and the Sobolev trace warning of V94 remain independent inputs.
\end{firewall}

\subsection{V95 conclusion and next acquisition}
\label{sec:v2v95-conclusion}

The mandatory companion audit has been completed.  V95 proves that the
connection-variation source cancels exactly between the geometric and matter
Ward identities before norms are taken.  Only geometric, matter, source, and
boundary Ward defects force the harmonic constraint.  The V32 Neumann
solution therefore uses the original response constant
\(\eta=\kappa\|E_RP\|\), while its constraint defect has the explicit bound
\eqref{eq:v2v95-Neumann-constraint}.  Finite order reductions have the
geometric constraint error \eqref{eq:v2v95-HN-rate}.

V96 will attack the remaining conditional estimate in this chain: a
maximally dissipative energy theorem for the harmonic constraint wave system
with a timelike boundary, expressing its incoming trace in terms of the
linearized relative Ward defect and normal stress flux.  If that boundary
condition cannot be made compatible with the chiral Dirac domain of V92,
the programme will return upstream to the boundary-domain construction
rather than assuming the trace estimate.
\clearpage
\section*{Second Volume, Version V96: Maximally Dissipative Constraint Energy and Relative Boundary Defects}
\addcontentsline{toc}{section}{Second Volume, Version V96: Maximally Dissipative Constraint Energy and Relative Boundary Defects}
\label{sec:v2v96}

\subsection{Purpose}
\label{sec:v2v96-purpose}

V95 reduced harmonic constraint propagation to an energy estimate for a
normally hyperbolic constraint operator with timelike boundary.  V96 proves
the required energy-level theorem for a uniformly timelike bounded-geometry
slab.  The incoming characteristic datum is allowed to be a bounded image of
the differentiated relative boundary Ward defect.  Substitution gives an
explicit bulk--boundary constraint bound for the V95 Neumann solution.

The theorem concerns the bosonic harmonic constraint system.  Compatibility
with the chiral Dirac boundary domain and with determinant-line data remains
a separate domain problem; no Euclidean strong-ellipticity statement is
identified with Lorentzian maximal dissipation.

\subsection{Constraint wave system and characteristic fields}
\label{sec:v2v96-system}

Let \(M_T=\bigcup_{0\le t\le T}\Sigma_t\) be a Lorentzian slab with smooth
timelike boundary \(\mathcal T\).  Let \(N\) be the future unit normal to
\(\Sigma_t\), and let \(n\) be the outward unit spacelike normal to
\(\partial\Sigma_t\) within \(\Sigma_t\).  Let \(u\) be a section of a
finite-rank bundle with a positive fibre metric and compatible connection.
Consider a normally hyperbolic operator
\begin{equation}
 Wu
 =
 \nabla_N^2u-\nabla^i\nabla_i u
 +A^\alpha\nabla_\alpha u+Bu
 =f,
\label{eq:v2v96-W}
\end{equation}
where lapse and shift factors are absorbed into uniformly bounded
coefficients.  Assume the principal energy is uniformly equivalent to
\begin{equation}
 E(t)
 =
 \frac12
 \left(
  \|\nabla_Nu(t)\|_{L^2(\Sigma_t)}^2
  +\|\nabla_xu(t)\|_{L^2(\Sigma_t)}^2
  +\|u(t)\|_{L^2(\Sigma_t)}^2
 \right).
\label{eq:v2v96-energy}
\end{equation}
Define boundary characteristic fields by
\begin{equation}
 w_{\rm in}=\nabla_Nu+\nabla_nu,
 \qquad
 w_{\rm out}=\nabla_Nu-\nabla_nu.
\label{eq:v2v96-characteristics}
\end{equation}
With these conventions the principal outward energy flux is
\begin{equation}
 \langle\nabla_Nu,\nabla_nu\rangle
 =
 \frac14
 \left(
  |w_{\rm in}|^2-|w_{\rm out}|^2
 \right).
\label{eq:v2v96-flux}
\end{equation}

Impose the boundary condition
\begin{equation}
 w_{\rm in}=\mathsf K w_{\rm out}+g_\partial,
 \qquad
 \|\mathsf K\|\le q<1
\label{eq:v2v96-boundary-condition}
\end{equation}
pointwise in the boundary fibre, with a fixed contraction margin.

\begin{lemma}[Strict boundary flux inequality]
\label{lem:v2v96-flux}
Under \eqref{eq:v2v96-boundary-condition},
\begin{equation}
 \frac14
 \left(
  |w_{\rm in}|^2-|w_{\rm out}|^2
 \right)
 \le
 -\frac{1-q^2}{8}|w_{\rm out}|^2
 +\frac{1+q^2}{4(1-q^2)}|g_\partial|^2.
\label{eq:v2v96-flux-bound}
\end{equation}
\end{lemma}

\begin{proof}
For \(q>0\), use
\[
 |\mathsf K v+g|^2
 \le
 (1+\epsilon)q^2|v|^2
 +(1+\epsilon^{-1})|g|^2
\]
with
\(\epsilon=(1-q^2)/(2q^2)\).  Then
\((1+\epsilon)q^2=(1+q^2)/2\) and
\(1+\epsilon^{-1}=(1+q^2)/(1-q^2)\).  Substitute
\(v=w_{\rm out}\) in \eqref{eq:v2v96-flux}.  For \(q=0\), the same bound
follows directly and the displayed right-side constant equals \(1/4\).
\end{proof}

\subsection{Energy theorem}
\label{sec:v2v96-energy-theorem}

Assume the deformation tensors of \(N\), the coefficients of \(W\), and the
slice/boundary geometry are bounded by a common constant \(M\), and that the
energy equivalence constants stay uniformly positive.

\begin{theorem}[Energy-level maximally dissipative estimate]
\label{thm:v2v96-energy}
Let \(u\) be a sufficiently regular solution of
\eqref{eq:v2v96-W}--\eqref{eq:v2v96-boundary-condition}.  There are constants
\(C_0,C_1\), depending only on the bounded-geometry and contraction margins,
such that
\begin{align}
 &\sup_{0\le t\le T}E(t)^{1/2}
 +\left(
   \int_{\mathcal T}|w_{\rm out}|^2\,\mathrm d\sigma
  \right)^{1/2}
 \nonumber\\
 &\qquad\le
 C_0e^{C_1T}
 \left[
  E(0)^{1/2}
  +\|f\|_{L^1(I;L^2_x)}
  +\|g_\partial\|_{L^2(\mathcal T)}
 \right].
\label{eq:v2v96-energy-estimate}
\end{align}
The estimate extends by density to finite-energy weak solutions whenever the
corresponding maximally dissipative problem is constructed.
\end{theorem}

\begin{proof}
Multiply \eqref{eq:v2v96-W} by \(\nabla_Nu\), integrate on \(\Sigma_t\),
and integrate the spatial principal part by parts.  The time dependence of
the volume form and metric, together with the lower-order terms, contributes
at most \(C E(t)\).  The forcing term is bounded by
\[
 \|f(t)\|_{L^2_x}\,\|\nabla_Nu(t)\|_{L^2_x}
 \le C\|f(t)\|_{L^2_x}E(t)^{1/2}.
\]
The boundary contribution is the integral of
\eqref{eq:v2v96-flux}; Lemma~\ref{lem:v2v96-flux} gives
\begin{align}
 E(t)
 +c_q\int_0^t\!\|w_{\rm out}(s)\|_{L^2(\partial\Sigma_s)}^2\,\mathrm ds
 \le{}&
 E(0)+C\int_0^t E(s)\,\mathrm ds
 \nonumber\\
 &+C\int_0^t\|f(s)\|_{L^2_x}E(s)^{1/2}\,\mathrm ds
 +C_q\|g_\partial\|_{L^2(\mathcal T_t)}^2,
\label{eq:v2v96-energy-intermediate}
\end{align}
where \(c_q=(1-q^2)/8\) up to the uniform energy-equivalence factor.
Apply the integral form of Gronwall's inequality to obtain
\[
 \sup_{r\le t}E(r)^{1/2}
 \le
 C_0e^{C_1t}
 \left(E(0)^{1/2}
 +\|f\|_{L^1([0,t];L^2_x)}
 +\|g_\partial\|_{L^2(\mathcal T_t)}\right).
\]
Insert this bound back into
\eqref{eq:v2v96-energy-intermediate} to control the outgoing trace.  This
proves \eqref{eq:v2v96-energy-estimate}.  Approximation by smooth compatible
data and lower semicontinuity give the finite-energy extension.
\end{proof}

\begin{corollary}[Uniqueness and causality]
\label{cor:v2v96-uniqueness}
If \(E(0)=0\), \(f=0\), and \(g_\partial=0\), then \(u=0\).  The solution on
\(M_t\) depends only on initial data, bulk forcing, and boundary input in the
causal past of \(M_t\).
\end{corollary}

\begin{proof}
The estimate makes \(E(t)=0\) for every \(t\).  The causal statement follows
by applying uniqueness to the difference of two solutions on the common
past domain.
\end{proof}

\begin{firewall}[The strict margin has a job]
\label{fw:v2v96-q-margin}
At \(q=1\), the boundary may be energy conserving, but
\eqref{eq:v2v96-energy-estimate} no longer controls the outgoing trace with a
positive coefficient.  For \(q>1\), the homogeneous boundary flux can be
positive.  A uniform trace estimate therefore requires either the strict
margin used here or a different coercive boundary mechanism proved from the
full system.
\end{firewall}

\subsection{Higher regularity without hidden normal derivatives}
\label{sec:v2v96-higher}

For \(m\ge1\), commute tangential derivatives of order at most \(m\) through
\eqref{eq:v2v96-W} and the boundary condition.  Normal derivatives are then
recovered from the equation and the characteristic relation.

\begin{proposition}[Conditional commuted estimate]
\label{prop:v2v96-higher}
Assume:
\begin{enumerate}
 \item the coefficients and geometry have bounded derivatives through order
 \(m+1\);
 \item the differentiated boundary conditions satisfy the usual corner
 compatibility conditions through order \(m\);
 \item commutators of tangential derivatives with \(\mathsf K\) are bounded
 in the corresponding boundary Sobolev norms; and
 \item the equation and boundary relation recover the missing normal
 derivatives without loss.
\end{enumerate}
Then
\begin{align}
 \|u\|_{U_m}
 \le
 C_me^{C_mT}
 \left(
  \|u[0]\|_{H^{m+1}\times H^m}
  +\|f\|_{L^1H^m}
  +\|g_\partial\|_{L^2H^m(\mathcal T)}
 \right),
\label{eq:v2v96-higher-estimate}
\end{align}
where \(U_m\) contains the slice
\(H^{m+1}\times H^m\) energy and the outgoing boundary trace through order
\(m\).
\end{proposition}

\begin{proof}
Apply Theorem~\ref{thm:v2v96-energy} to every tangential derivative.  The
commutators are lower order and are absorbed inductively by Gronwall.  The
last hypothesis converts normal derivatives to tangential, temporal, and
forcing terms, so no unrecorded trace derivative is used.  Summing the
commuted estimates proves \eqref{eq:v2v96-higher-estimate}.
\end{proof}

\begin{firewall}[No automatic half-derivative repair]
\label{fw:v2v96-half-derivative}
The last hypothesis of Proposition~\ref{prop:v2v96-higher} is a boundary
regularity statement, not a formal consequence of bulk hyperbolicity.  The
half-derivative trace issue identified in V94 cannot be paid by writing more
commuted energy identities.  This is the boundary version of the NS V31
warning that smoothing does not create a missing scale moment.
\end{firewall}

\subsection{Insertion of the relative Ward defect}
\label{sec:v2v96-relative-insertion}

Return to the harmonic constraint \(u=Hh\).  In V95 notation its bulk
forcing is
\begin{equation}
 f_H
 =
 \kappa\mathcal BS-(d_G-\kappa d_M)h.
\label{eq:v2v96-fH}
\end{equation}
Let the differentiated boundary Ward residual be
\begin{align}
 \mathfrak d_\partial(h)
 :={}&
 \delta_h(D_a\tau^a{}_b)
 -\delta_h(n_\mu T^\mu{}_b)
 \nonumber\\
 &-\delta_h(f^A_{ba}j_A^a)
 -\delta_h(e_iD_b\varphi^i),
\label{eq:v2v96-boundary-Ward-residual}
\end{align}
where the first term includes
\eqref{eq:v2v94-boundary-C}.  Let \(b_\partial\) be an external physical
boundary flux.  Assume the chosen gravitational boundary condition has the
characteristic form
\begin{equation}
 w_{\rm in}(Hh)
 =
 \mathsf K w_{\rm out}(Hh)
 +R_\partial(\mathfrak d_\partial(h)+\kappa b_\partial),
\label{eq:v2v96-relative-boundary-condition}
\end{equation}
where
\begin{equation}
 \|R_\partial v\|_{L^2(\mathcal T)}
 \le C_\partial\|v\|_{Z_\partial}.
\label{eq:v2v96-R-bound}
\end{equation}

\begin{theorem}[Explicit bulk--boundary constraint ledger]
\label{thm:v2v96-ledger}
Let \(h\) be the V95 Neumann solution, and put
\begin{equation}
 H_*=
 \frac{\kappa\|E_R\|}{1-\eta}\|S\|_Y.
\label{eq:v2v96-Hstar}
\end{equation}
Assume zero initial harmonic-constraint data and
\begin{align}
 \|d_G\|&\le\delta_G,
 &\|d_M\|&\le\delta_M,
 &\|\mathfrak d_\partial(h)\|_{Z_\partial}
 &\le\delta_\partial\|h\|_X.
\label{eq:v2v96-defect-sizes}
\end{align}
Then
\begin{align}
 \|Hh\|_{U_0}
 \le
 C_0e^{C_1T}\bigl[&
  \kappa\|\mathcal BS\|_{L^1L^2}
  +(\delta_G+\kappa\delta_M)H_*
 \nonumber\\
 &+C_\partial\delta_\partial H_*
  +\kappa C_\partial\|b_\partial\|_{Z_\partial}
 \bigr].
\label{eq:v2v96-ledger}
\end{align}
For exact bulk and boundary Ward identities and conserved source data, the
right side vanishes and \(Hh=0\).
\end{theorem}

\begin{proof}
Use \eqref{eq:v2v96-fH} to estimate the bulk forcing in
Theorem~\ref{thm:v2v96-energy}.  Equations
\eqref{eq:v2v96-relative-boundary-condition} and
\eqref{eq:v2v96-R-bound} bound the incoming datum.  Finally use the V95
Neumann estimate \(\|h\|_X\le H_*\).  This gives
\eqref{eq:v2v96-ledger}.  Exactness makes both the bulk forcing and incoming
boundary datum zero, so Corollary~\ref{cor:v2v96-uniqueness} applies.
\end{proof}

\begin{corollary}[Uniform regulator constraint convergence]
\label{cor:v2v96-regulator}
Suppose the geometry and contraction margins in
Theorem~\ref{thm:v2v96-energy} are uniform in \(n\), the Neumann constants
satisfy one common \(\eta<1\), the source ranges are uniformly bounded, and
\begin{equation}
 \|\mathcal B_nS_n\|
 +\delta_{G,n}+\delta_{M,n}+\delta_{\partial,n}
 +\|b_{\partial,n}\|
 \longrightarrow0.
\label{eq:v2v96-regulator-defects}
\end{equation}
Then
\begin{equation}
 \|H_nh_n\|_{U_0}\longrightarrow0.
\label{eq:v2v96-constraint-convergence}
\end{equation}
If the limiting reduced solutions converge in \(X\), the reduced operators
converge on their graph domains, and the constraint and trace maps are closed
under that convergence, then the limit satisfies the unreduced Einstein
equation and homogeneous constraint boundary condition.
\end{corollary}

\begin{proof}
All coefficients multiplying the defect quantities in
\eqref{eq:v2v96-ledger} are uniform, so its right side tends to zero.  Closed
operator convergence of the reduced equation and continuity of the
constraint trace identify the limit.
\end{proof}

\subsection{What remains at the chiral boundary}
\label{sec:v2v96-chiral-boundary}

\begin{programme}[Boundary-domain compatibility gate]
\label{prog:v2v96-boundary-domain}
To apply Theorem~\ref{thm:v2v96-ledger} to the complete Standard Model,
prove on one common boundary domain:
\begin{enumerate}
 \item self-adjointness and the V92 uniform gap/heat bounds for the tangential
 chiral Dirac family;
 \item gauge and diffeomorphism covariance of that domain under the relative
 counterterms;
 \item a bosonic gravitational boundary condition of the form
 \eqref{eq:v2v96-relative-boundary-condition} with \(q<1\);
 \item the trace map \eqref{eq:v2v96-R-bound}, including the eta-sector and
 normal stress response; and
 \item uniform preservation of all four properties under regulator removal.
\end{enumerate}
\end{programme}

\begin{firewall}[Euclidean and Lorentzian boundary tests are distinct]
\label{fw:v2v96-Euclidean-Lorentzian}
Strong ellipticity of the Euclidean Dirac boundary problem supplies heat
asymptotics for V92.  Maximal dissipation of the Lorentzian harmonic
constraint problem supplies the energy estimate in V96.  Neither property
implies the other.  The full boundary theorem needs a domain on which both
sector-specific tests and the relative Ward identity are simultaneously
valid.
\end{firewall}

\subsection{V96 conclusion and next acquisition}
\label{sec:v2v96-conclusion}

V96 proves the energy-level timelike-boundary estimate required by V95.  A
strict characteristic contraction controls the outgoing constraint trace,
and the relative boundary Ward residual becomes the incoming datum.  The
resulting estimate \eqref{eq:v2v96-ledger} separates conserved source error,
bulk geometric and matter Ward defects, boundary Ward defect, and physical
boundary flux without double charging the connection block.

V97 will go upstream to the common boundary-domain problem.  It will
formulate a graph-norm criterion under which the tangential Dirac projector,
relative BRST differential, metric variation, and Lorentzian boundary trace
share one stable domain.  A finite-dimensional counterexample will test why
separate domain convergence does not imply convergence of their
intersection.
\clearpage
\section*{Second Volume, Version V97: Common Boundary Graph, Uniform Angle, and Domain Transport}
\addcontentsline{toc}{section}{Second Volume, Version V97: Common Boundary Graph, Uniform Angle, and Domain Transport}
\label{sec:v2v97}

\subsection{Purpose}
\label{sec:v2v97-purpose}

V92 requires a stable self-adjoint tangential Dirac domain and uniform heat
data.  V93 requires relative BRST and diffeomorphism invariance of the same
boundary theory.  V96 requires a maximally dissipative incoming harmonic
constraint condition.  V97 packages these requirements into one bounded
constraint map on a common graph space.

The central theorem is an elementary but decisive uniform-angle result.  If
the stacked boundary map is uniformly surjective, its kernel has an
orthogonal projection which converges in norm and varies differentiably with
the metric.  Kato transport then identifies the varying common domains.
Separate convergence of the component domains is not enough; a two-dimensional
counterexample exhibits a jumping intersection when their relative angle
collapses.

\subsection{The stacked boundary constraint}
\label{sec:v2v97-stacked}

Let \(X\) be a Hilbert graph space for the complete linearized bulk field
multiplet on a fixed slab, after the regulator-dependent bundles have been
identified.  Its norm includes the bulk operator graph norms needed for the
Dirac, gauge, ghost, Higgs, and harmonic constraint sectors.  Let
\begin{equation}
 Y
 =Y_{\rm D}\oplus Y_{\rm BRST}\oplus Y_{\rm H}
\label{eq:v2v97-Y}
\end{equation}
be the boundary data space.  At regulator \(n\), write
\begin{equation}
 F_n=
 \begin{pmatrix}
  F_{{\rm D},n}\\
  F_{{\rm BRST},n}\\
  F_{{\rm H},n}
 \end{pmatrix}:X\longrightarrow Y.
\label{eq:v2v97-Fn}
\end{equation}
The entries have the following intended meanings.  In the Dirac entry, the`r`nselected spectral range is first identified unitarily with the fixed space`r`n\(Y_{\rm D}\); this range identification is included in the notation`r`n\(P_{\partial,n}\):
\begin{align}
 F_{{\rm D},n}u
 &=P_{\partial,n}\gamma_{{\rm D},n}u,
\label{eq:v2v97-Dirac-entry}\\
 F_{{\rm BRST},n}u
 &=G_{\partial,n}\gamma_{{\rm BRST},n}u,
\label{eq:v2v97-BRST-entry}\\
 F_{{\rm H},n}u
 &=w_{{\rm in},n}(u)-\mathsf K_nw_{{\rm out},n}(u).
\label{eq:v2v97-H-entry}
\end{align}
Here \(P_{\partial,n}\) is the chosen tangential spectral projector,
\(G_{\partial,n}\) encodes homogeneous relative gauge/ghost conditions, and
\eqref{eq:v2v97-H-entry} is the homogeneous part of the V96 characteristic
condition.  Inhomogeneous physical boundary data are treated by an affine
translate after the homogeneous domain is controlled.

Define the common homogeneous domain
\begin{equation}
 \mathcal D_n=\ker F_n\subset X.
\label{eq:v2v97-Dn}
\end{equation}

\begin{definition}[Uniform common-domain angle]
\label{def:v2v97-angle}
The family \(F_n\) has angle margin \(\gamma>0\) if
\begin{equation}
 F_nF_n^*\ge\gamma^2 I_Y
\label{eq:v2v97-angle}
\end{equation}
as quadratic forms on \(Y\), uniformly in \(n\).
\end{definition}

The condition implies that \(F_n\) is onto and that the Moore--Penrose right
inverse
\begin{equation}
 R_n=F_n^*(F_nF_n^*)^{-1}
\label{eq:v2v97-Rn}
\end{equation}
has norm at most \(\gamma^{-1}\).

\subsection{Stable common-domain projection}
\label{sec:v2v97-projection}

\begin{theorem}[Norm convergence of the common boundary domain]
\label{thm:v2v97-domain-projection}
Assume
\begin{equation}
 F_n\longrightarrow F
 \quad\hbox{in }\mathcal L(X,Y)
\label{eq:v2v97-F-convergence}
\end{equation}
and the uniform angle condition \eqref{eq:v2v97-angle}.  Then
\(FF^*\ge\gamma^2I\), and the orthogonal projections
\begin{equation}
 \Pi_n
 =I_X-F_n^*(F_nF_n^*)^{-1}F_n
\label{eq:v2v97-Pin}
\end{equation}
onto \(\mathcal D_n\) satisfy
\begin{equation}
 \Pi_n\longrightarrow
 \Pi:=I_X-F^*(FF^*)^{-1}F
 \quad\hbox{in }\mathcal L(X).
\label{eq:v2v97-Pi-convergence}
\end{equation}
In particular:
\begin{enumerate}
 \item if \(u\in\ker F\), then \(u_n:=\Pi_nu\in\mathcal D_n\) and
 \(u_n\to u\) in \(X\);
 \item if \(u_n\in\mathcal D_n\) and \(u_n\to u\) in \(X\), then
 \(u\in\ker F\).
\end{enumerate}
\end{theorem}

\begin{proof}
Norm convergence of \(F_n\) implies norm convergence of
\(F_nF_n^*\) to \(FF^*\).  Taking limits in
\(\langle F_nF_n^*y,y\rangle\ge\gamma^2\|y\|^2\) gives the same lower
bound for \(FF^*\).  The resolvent identity yields
\begin{align*}
 &(F_nF_n^*)^{-1}-(FF^*)^{-1}\\
 &\qquad=
 (F_nF_n^*)^{-1}(FF^*-F_nF_n^*)(FF^*)^{-1},
\end{align*}
so the inverses converge in norm and are uniformly bounded by
\(\gamma^{-2}\).  Substitution in \eqref{eq:v2v97-Pin} proves
\eqref{eq:v2v97-Pi-convergence}.

The formula \eqref{eq:v2v97-Pin} is the orthogonal projection onto
\(\ker F_n\): it is self-adjoint and idempotent, its range lies in the
kernel because \(F_nF_n^*(F_nF_n^*)^{-1}=I\), and it fixes every kernel
vector.  Item 1 follows from projection convergence and \(\Pi u=u\).
For item 2,
\[
 Fu=\lim_{n\to\infty}F_nu_n=0.
\]
\end{proof}

\begin{corollary}[Convergence of restricted bulk operators]
\label{cor:v2v97-restricted-operators}
Let \(K_n,K:X\to H\) be the graph realizations of bulk operators and assume
\(K_n\to K\) in \(\mathcal L(X,H)\).  Then
\begin{equation}
 K_n\Pi_n\longrightarrow K\Pi
 \quad\hbox{in }\mathcal L(X,H).
\label{eq:v2v97-restricted-convergence}
\end{equation}
Thus the restrictions to \(\mathcal D_n\), transported by their common-domain
projections, converge in graph norm.
\end{corollary}

\begin{proof}
Use
\[
 K_n\Pi_n-K\Pi
 =(K_n-K)\Pi_n+K(\Pi_n-\Pi)
\]
and the uniform bound \(\|\Pi_n\|=1\).
\end{proof}

\subsection{Why separate domain convergence fails}
\label{sec:v2v97-counterexample}

\begin{counterexample}[Jumping intersection under an angle collapse]
\label{sep:v2v97-angle-collapse}
Let \(X=\mathbb R^2\) and define
\begin{equation}
 F_n^{(1)}(x,y)=x,
 \qquad
 F_n^{(2)}(x,y)=x+n^{-1}y.
\label{eq:v2v97-two-maps}
\end{equation}
Each map is onto, converges in norm, and has a right inverse of norm at most
one.  Nevertheless,
\begin{equation}
 \ker F_n^{(1)}\cap\ker F_n^{(2)}=\{0\}
 \quad(n<\infty),
\label{eq:v2v97-finite-intersection}
\end{equation}
whereas
\begin{equation}
 \ker F^{(1)}\cap\ker F^{(2)}
 =\{(0,y):y\in\mathbb R\}.
\label{eq:v2v97-limit-intersection}
\end{equation}
\end{counterexample}

\begin{proof}
The two finite equations are \(x=0\) and \(x+n^{-1}y=0\), hence
\(x=y=0\).  In the limit both equations are \(x=0\).  The stacked map
\[
 \mathbf F_n(x,y)=(x,x+n^{-1}y)
\]
has determinant \(n^{-1}\), so the norm of its inverse diverges at least
linearly in \(n\).  Thus the component maps are individually stable while
the common angle condition fails exactly where the intersection jumps.
\end{proof}

\begin{firewall}[Component gaps do not imply a common angle]
\label{fw:v2v97-component-gaps}
The V92 Dirac spectral gap, the V88 BRST Hodge gap, and the V96 dissipative
margin control three different component problems.  They do not imply
\eqref{eq:v2v97-angle} for the stacked map.  Cross-sector boundary relations
can become asymptotically dependent just as in
Counterexample~\ref{sep:v2v97-angle-collapse}.
\end{firewall}

\subsection{Input from the tangential Dirac gap}
\label{sec:v2v97-Dirac-gap}

Let \(A_{\partial,n}\) be self-adjoint tangential Dirac operators and let
\(P_{\partial,n}\) be a Riesz projector associated with a spectral cluster
separated by a contour \(\Gamma\).  The V87 norm-resolvent theorem gives
\begin{equation}
 \|P_{\partial,n}-P_\partial\|
 \le
 \frac{|\Gamma|}{2\pi}
 \sup_{z\in\Gamma}
 \|(A_{\partial,n}-z)^{-1}-(A_\partial-z)^{-1}\|.
\label{eq:v2v97-Riesz-rate}
\end{equation}

\begin{proposition}[Convergence of the Dirac boundary entry]
\label{prop:v2v97-Dirac-entry}
If \(P_{\partial,n}\to P_\partial\) in operator norm and
\(\gamma_{{\rm D},n}\to\gamma_{\rm D}\) in
\(\mathcal L(X,Y_{\rm D})\), then
\begin{equation}
 F_{{\rm D},n}longrightarrow
 F_{\rm D}=P_\partial\gamma_{\rm D}
 \quad\hbox{in }\mathcal L(X,Y_{\rm D}).
\label{eq:v2v97-Dirac-entry-convergence}
\end{equation}
\end{proposition}

\begin{proof}
Write
\[
 P_{\partial,n}\gamma_{{\rm D},n}
 -P_\partial\gamma_{\rm D}
 =
 (P_{\partial,n}-P_\partial)\gamma_{{\rm D},n}
 +P_\partial(\gamma_{{\rm D},n}-\gamma_{\rm D})
\]
and take norms.
\end{proof}

\begin{firewall}[Projection convergence is not heat convergence]
\label{fw:v2v97-projection-not-heat}
Equation \eqref{eq:v2v97-Riesz-rate} stabilizes the boundary domain near a
spectral cluster.  It does not supply the uniform small-time heat subtraction
or large-time eta bound of V92.  Those remain additional analytic estimates
on the restricted Dirac family.
\end{firewall}

\subsection{Preservation of self-adjoint boundary geometry}
\label{sec:v2v97-selfadjoint}

Let \(H_\partial\) be the fermionic boundary Hilbert space.  Let
\(J_n\) be the bounded skew-unitary Green-form operator, so that
\begin{equation}
 \omega_n(x,y)=\langle J_nx,y\rangle,
 \qquad
 J_n^*=-J_n,
 \qquad
 J_n^*J_n=I.
\label{eq:v2v97-Green-form}
\end{equation}
Let \(Q_n\) be the orthogonal projector onto the allowed fermionic boundary
subspace \(L_n\).

\begin{proposition}[Closedness of the Lagrangian condition]
\label{prop:v2v97-Lagrangian}
Assume \(J_n\to J\) and \(Q_n\to Q\) in operator norm and
\begin{equation}
 J_nQ_nJ_n^{-1}=I-Q_n
\label{eq:v2v97-Lagrangian-n}
\end{equation}
for every \(n\).  Then
\begin{equation}
 JQJ^{-1}=I-Q.
\label{eq:v2v97-Lagrangian-limit}
\end{equation}
Thus the limiting subspace is Lagrangian for the limiting Green form, and
the corresponding symmetric Dirac boundary relation remains self-adjoint
whenever the standard Green-form characterization applies.
\end{proposition}

\begin{proof}
Skew unitarity is closed under norm convergence, so \(J^{-1}\) exists and
\(J_n^{-1}\to J^{-1}\).  Take the norm limit in
\eqref{eq:v2v97-Lagrangian-n}.  The identity
\eqref{eq:v2v97-Lagrangian-limit} says
\(JL=(L)^\perp\), which is the maximal isotropic condition for
\(\omega(x,y)=\langle Jx,y\rangle\).
\end{proof}

\subsection{Metric differentiability and Kato transport}
\label{sec:v2v97-transport}

Let \(r\mapsto F(r)\in\mathcal L(X,Y)\) be \(C^1\), with
\begin{equation}
 \|F(r)\|\le M,
 \qquad
 F(r)F(r)^*\ge\gamma^2I
\label{eq:v2v97-path-margin}
\end{equation}
on an interval.  Put
\begin{equation}
 \Pi(r)=I-F(r)^*(F(r)F(r)^*)^{-1}F(r).
\label{eq:v2v97-path-Pi}
\end{equation}

\begin{lemma}[Derivative bound for the domain projector]
\label{lem:v2v97-Pi-derivative}
The path \(\Pi(r)\) is \(C^1\), and
\begin{equation}
 \|\dot\Pi(r)\|
 \le
 \left(
  2M\gamma^{-2}+2M^3\gamma^{-4}
 \right)
 \|\dot F(r)\|.
\label{eq:v2v97-Pi-derivative}
\end{equation}
\end{lemma}

\begin{proof}
Let \(G=FF^*\).  Differentiate
\(\Pi=I-F^*G^{-1}F\) and use
\(dot G^{-1}=-G^{-1}\dot G G^{-1}\):
\begin{align*}
 \dot\Pi={}&
 -\dot F^*G^{-1}F-F^*G^{-1}\dot F\\
 &+F^*G^{-1}(\dot F F^*+F\dot F^*)G^{-1}F.
\end{align*}
Now use \(\|G^{-1}\|\le\gamma^{-2}\),
\(\|F\|\le M\), and
\(\|\dot G\|\le2M\|\dot F\|\).
\end{proof}

\begin{theorem}[Unitary transport of the common domain]
\label{thm:v2v97-Kato}
Let \(U(r)\) solve
\begin{equation}
 \dot U(r)=[\dot\Pi(r),\Pi(r)]U(r),
 \qquad
 U(0)=I.
\label{eq:v2v97-Kato-ODE}
\end{equation}
Then \(U(r)\) is unitary and
\begin{equation}
 \Pi(r)U(r)=U(r)\Pi(0).
\label{eq:v2v97-intertwining}
\end{equation}
Consequently
\begin{equation}
 U(r)\mathcal D(0)=\mathcal D(r),
 \qquad
 \|\dot U(r)\|
 \le2\|\dot\Pi(r)\|.
\label{eq:v2v97-domain-transport}
\end{equation}
\end{theorem}

\begin{proof}
For an orthogonal projection, differentiation of \(\Pi^2=\Pi\) gives
\(\Pi\dot\Pi\Pi=0\).  The commutator
\([\dot\Pi,\Pi]\) is skew-adjoint, so the ODE preserves \(U^*U=I\).
A direct calculation using \(\Pi\dot\Pi\Pi=0\) shows
\[
 \frac{\mathrm d}{\mathrm dr}
 \bigl(U(r)^*\Pi(r)U(r)\bigr)=0.
\]
Hence \(U^*\Pi U=\Pi(0)\), which is
\eqref{eq:v2v97-intertwining}.  The range identity follows, and the norm
bound uses \(\|[\dot\Pi,\Pi]\|\le2\|\dot\Pi\|\).
\end{proof}

\begin{corollary}[Controlled chart derivative]
\label{cor:v2v97-chart-derivative}
Under \eqref{eq:v2v97-path-margin}, differentiation after transporting all
fields to \(\mathcal D(0)\) adds a connection term with norm at most
\begin{equation}
 4\left(M\gamma^{-2}+M^3\gamma^{-4}\right)
 \|\dot F(r)\|.
\label{eq:v2v97-chart-cost}
\end{equation}
This is a concrete domain-chart contribution to the V53 relative-Cauchy and
stress-response derivatives.
\end{corollary}

\begin{proof}
The transport connection is \(U^{-1}\dot U\).  Apply
\eqref{eq:v2v97-domain-transport} and
\eqref{eq:v2v97-Pi-derivative}.
\end{proof}

\subsection{BRST invariance of the common domain}
\label{sec:v2v97-BRST}

Let \(Q_n:X\to X\) be the finite-regulator BRST differential in the graph
realization.  The common domain is invariant precisely when
\begin{equation}
 F_nQ_n\Pi_n=0.
\label{eq:v2v97-BRST-domain}
\end{equation}

\begin{proposition}[Regulator-stable domain invariance]
\label{prop:v2v97-BRST-limit}
Suppose \(Q_n\to Q\), \(F_n\to F\), and the hypotheses of
Theorem~\ref{thm:v2v97-domain-projection} hold.  If
\begin{equation}
 \|F_nQ_n\Pi_n\|\longrightarrow0,
\label{eq:v2v97-BRST-defect}
\end{equation}
then
\begin{equation}
 FQ\Pi=0.
\label{eq:v2v97-BRST-limit}
\end{equation}
Thus \(Q(\ker F)\subseteq\ker F\).
\end{proposition}

\begin{proof}
Every factor converges in operator norm, so
\(F_nQ_n\Pi_n\to FQ\Pi\).  The assumed defect convergence makes the limit
zero.  Since \(\Pi\) projects onto \(\ker F\), the last statement follows.
\end{proof}

\begin{firewall}[Nilpotence alone does not preserve a boundary domain]
\label{fw:v2v97-nilpotence}
The identity \(Q_n^2=0\) does not imply
\eqref{eq:v2v97-BRST-domain}.  Boundary invariance is a mixed statement
about \(Q_n\) and the stacked domain projector.  It must be proved before the
relative anomaly complex is restricted to boundary-admissible fields.
\end{firewall}

\subsection{Dissipative margin inside the common graph}
\label{sec:v2v97-dissipative}

\begin{proposition}[Closedness of the characteristic contraction]
\label{prop:v2v97-dissipative-limit}
If \(\mathsf K_n\to\mathsf K\) in boundary operator norm and
\begin{equation}
 \|\mathsf K_n\|\le q<1
\label{eq:v2v97-uniform-q}
\end{equation}
uniformly, then \(\|\mathsf K\|\le q<1\).  Hence the V96 boundary flux
constant remains positive in the limit.
\end{proposition}

\begin{proof}
The operator norm is continuous under norm convergence.  Apply the V96 flux
lemma with the same \(q\).
\end{proof}

\begin{assessment}[What the common angle adds]
\label{ass:v2v97-common-angle}
The Dirac Riesz gap stabilizes \(P_{\partial,n}\), BRST estimates stabilize
the algebraic differential, and the characteristic margin stabilizes the
constraint flux.  The additional condition
\(F_nF_n^*\ge\gamma^2I\) prevents these stable components from becoming
mutually dependent when imposed together.  It is the precise analytic
replacement for the phrase ``choose compatible boundary conditions.''
\end{assessment}

\subsection{V97 conclusion and next acquisition}
\label{sec:v2v97-conclusion}

V97 constructs a common boundary graph for the tangential Dirac, relative
BRST, and harmonic constraint conditions.  A uniform inf--sup angle gives a
norm-convergent orthogonal projector onto the joint domain, graph convergence
of restricted bulk operators, and a differentiable Kato transport under
metric variation.  Self-adjoint Green-form geometry, BRST domain invariance,
and strict characteristic dissipation all pass to the limit under their
stated norm hypotheses.  A finite-dimensional example proves that separate
domain convergence cannot replace the stacked angle.

The remaining acquisition is now concrete.  V98 will derive a sufficient
block-Schur criterion for the common angle in
\eqref{eq:v2v97-angle}.  It will quantify how much cross-sector boundary
coupling can be tolerated relative to the separate Dirac, BRST, and harmonic
right-inverse margins, and give a sharp two-block failure at equality.

\clearpage
\section*{Second Volume, Version V98: Block--Schur Criterion for the Common Boundary Angle}
\addcontentsline{toc}{section}{Second Volume, Version V98: Block--Schur Criterion for the Common Boundary Angle}
\label{sec:v2v98}

\subsection{Purpose}
\label{sec:v2v98-purpose}

V97 reduced stability of the joint Dirac--BRST--harmonic boundary domain to
the uniform inequality
\(FF^*\ge\gamma^2I\).  V98 turns that condition into block estimates.  Each
component boundary map supplies a diagonal surjectivity margin; normalized
cross-sector Gram operators measure the angle loss.  If their scalar row
sum is below one, the stacked map has an explicit common margin.  A
rotating two-row example proves sharp degeneration at equality.

\subsection{Gram normalization}
\label{sec:v2v98-Gram}

Let
\begin{equation}
 F=
 \begin{pmatrix}F_1\\ \vdots\\ F_m\end{pmatrix}:
 X\longrightarrow Y_1\oplus\cdots\oplus Y_m=:Y
\label{eq:v2v98-F}
\end{equation}
be bounded between Hilbert spaces.  Put
\begin{equation}
 G=FF^*,
 \qquad
 G_{ij}=F_iF_j^*:Y_j\to Y_i.
\label{eq:v2v98-G}
\end{equation}
Assume the component margins
\begin{equation}
 G_{ii}\ge\gamma_i^2I_{Y_i},
 \qquad
 \gamma_i>0.
\label{eq:v2v98-component-margin}
\end{equation}
Define normalized cross couplings
\begin{equation}
 T_{ij}
 =G_{ii}^{-1/2}G_{ij}G_{jj}^{-1/2},
 \qquad
 \rho_{ij}=\|T_{ij}\|,
 \qquad i\ne j.
\label{eq:v2v98-Tij}
\end{equation}
Since \(G_{ji}=G_{ij}^*\), one has \(\rho_{ji}=\rho_{ij}\).

Let \(R\) be the real symmetric \(m\)-by-\(m\) matrix with zero diagonal
and off-diagonal entries \(R_{ij}=\rho_{ij}\).

\begin{theorem}[Normalized block--Gram angle criterion]
\label{thm:v2v98-Gram-criterion}
If
\begin{equation}
 \|R\|_{\ell^2_m\to\ell^2_m}\le\rho_*<1,
\label{eq:v2v98-rho}
\end{equation}
then
\begin{equation}
 FF^*
 \ge
 (1-\rho_*)
 \left(\min_{1\le i\le m}\gamma_i^2\right)I_Y.
\label{eq:v2v98-common-margin}
\end{equation}
Thus the common-domain angle in V97 may be chosen as
\begin{equation}
 \gamma_{
m com}
 =
 \sqrt{1-\rho_*}\,min_i\gamma_i.
\label{eq:v2v98-gamma-common}
\end{equation}
\end{theorem}

\begin{proof}
Let
\(D=\operatorname{diag}(G_{11},\ldots,G_{mm})\) and write
\begin{equation}
 G=D^{1/2}(I+T)D^{1/2},
\label{eq:v2v98-normalized-Gram}
\end{equation}
where \(T_{ii}=0\) and the off-diagonal blocks are
\eqref{eq:v2v98-Tij}.  For \(z=(z_1,\ldots,z_m)\), put
\(a_i=\|z_i\|\).  Then
\begin{align*}
 |\langle Tz,z\rangle|
 &\le
 \sum_{i\ne j}\rho_{ij}a_ia_j
 =\langle Ra,a\rangle
 \le\rho_*\|a\|_{\ell^2}^2
 =\rho_*\|z\|_Y^2.
\end{align*}
Since \(T\) is self-adjoint, this proves \(I+T\ge(1-\rho_*)I\).
Also \(D\ge(\min_i\gamma_i^2)I\) by
\eqref{eq:v2v98-component-margin}.  Insert both inequalities in
\eqref{eq:v2v98-normalized-Gram} to obtain
\eqref{eq:v2v98-common-margin}.
\end{proof}

\begin{corollary}[Row-sum test]
\label{cor:v2v98-row-sum}
It is sufficient that
\begin{equation}
 \max_i\sum_{j\ne i}\rho_{ij}le r_*<1.
\label{eq:v2v98-row-sum}
\end{equation}
Then \eqref{eq:v2v98-common-margin} holds with \(\rho_*=r_*\).
\end{corollary}

\begin{proof}
For a real symmetric matrix,
\(\|R\|_{2\to2}\le\sqrt{\|R\|_1\|R\|_\infty}=\|R\|_\infty\).
Apply Theorem~\ref{thm:v2v98-Gram-criterion}.
\end{proof}

\subsection{The three-sector criterion}
\label{sec:v2v98-three}

For the V97 decomposition, use indices
\({\rm D},{\rm B},{\rm H}\) for Dirac, relative BRST, and harmonic
constraint boundary data.  Put
\begin{align}
 \rho_{{\rm DB}}
 &=
 \|G_{\rm DD}^{-1/2}G_{\rm DB}G_{\rm BB}^{-1/2}\|,
\label{eq:v2v98-rho-DB}\\
 \rho_{{\rm DH}}
 &=
 \|G_{\rm DD}^{-1/2}G_{\rm DH}G_{\rm HH}^{-1/2}\|,
\label{eq:v2v98-rho-DH}\\
 \rho_{{\rm BH}}
 &=
 \|G_{\rm BB}^{-1/2}G_{\rm BH}G_{\rm HH}^{-1/2}\|.
\label{eq:v2v98-rho-BH}
\end{align}

\begin{corollary}[Explicit Dirac--BRST--harmonic margin]
\label{cor:v2v98-three-margin}
If
\begin{equation}
 r_*:=
 \max\{
  \rho_{{\rm DB}}+\rho_{{\rm DH}},
  \rho_{{\rm DB}}+\rho_{{\rm BH}},
  \rho_{{\rm DH}}+\rho_{{\rm BH}}
 \}<1,
\label{eq:v2v98-three-row-sum}
\end{equation}
then the V97 stacked map satisfies
\begin{equation}
 FF^*\ge
 (1-r_*)
 \min\{\gamma_{\rm D}^2,
       \gamma_{\rm B}^2,
       \gamma_{\rm H}^2\}I.
\label{eq:v2v98-three-angle}
\end{equation}
\end{corollary}

\begin{proof}
The three quantities in \eqref{eq:v2v98-three-row-sum} are the row sums of
the scalar comparison matrix.  Apply
Corollary~\ref{cor:v2v98-row-sum}.
\end{proof}

\begin{assessment}[Orthogonal sectors close immediately]
\label{ass:v2v98-orthogonal}
If the three boundary maps act on mutually orthogonal summands of the bulk
graph space, then every off-diagonal Gram block vanishes.  Hence
\(r_*=0\) and
\(\gamma_{\rm com}=\min\gamma_i\).  The cross couplings arise only from
shared boundary variables, relative Ward matching, gauge fixing, or
metric-dependent identifications.  Those are the blocks that require an
estimate.
\end{assessment}

\subsection{Exact two-block Schur estimate}
\label{sec:v2v98-two-block}

\begin{theorem}[Two-block normalized Schur bound]
\label{thm:v2v98-two-block}
Let
\begin{equation}
 G=
 \begin{pmatrix}
  A&C\\ C^*&D
 \end{pmatrix},
 \qquad
 A\ge a^2I,
 \qquad
 D\ge d^2I,
\label{eq:v2v98-two-Gram}
\end{equation}
and suppose
\begin{equation}
 \|A^{-1/2}CD^{-1/2}\|\le\rho<1.
\label{eq:v2v98-two-rho}
\end{equation}
Then
\begin{equation}
 G\ge(1-\rho)\min\{a^2,d^2\}I.
\label{eq:v2v98-two-margin}
\end{equation}
Equivalently, the Schur complements satisfy
\begin{align}
 A-CD^{-1}C^*&\ge(1-\rho^2)A,
\label{eq:v2v98-Schur-A}\\
 D-C^*A^{-1}C&\ge(1-\rho^2)D.
\label{eq:v2v98-Schur-D}
\end{align}
\end{theorem}

\begin{proof}
Factor
\[
 G=
 \begin{pmatrix}A^{1/2}&0\\0&D^{1/2}\end{pmatrix}
 \begin{pmatrix}I&T\\T^*&I\end{pmatrix}
 \begin{pmatrix}A^{1/2}&0\\0&D^{1/2}\end{pmatrix},
 \qquad
 T=A^{-1/2}CD^{-1/2}.
\]
The middle matrix is bounded below by \((1-\|T\|)I\), proving
\eqref{eq:v2v98-two-margin}.  Also
\[
 CD^{-1}C^*
 =A^{1/2}TT^*A^{1/2}
 \le\rho^2A,
\]
and the other Schur inequality is analogous.
\end{proof}

\subsection{Sharp failure at equality}
\label{sec:v2v98-sharp}

\begin{counterexample}[Parallel boundary rows]
\label{sep:v2v98-parallel}
Let \(X=\mathbb R^2\), \(Y_1=Y_2=\mathbb R\), and
\begin{equation}
 F_{1,\theta}(x,y)=x,
 \qquad
 F_{2,\theta}(x,y)=\cos\theta\,x+\sin\theta\,y.
\label{eq:v2v98-rotating-rows}
\end{equation}
Each component has margin one.  The stacked Gram matrix is
\begin{equation}
 G_\theta=
 \begin{pmatrix}
  1&\cos\theta\\
  \cos\theta&1
 \end{pmatrix}
\label{eq:v2v98-Gtheta}
\end{equation}
with eigenvalues \(1\pm|\cos\theta|\).  Hence
\begin{equation}
 \gamma_{\rm com}(\theta)
 =\sqrt{1-|\cos\theta|}
 \sim\frac{|\theta|}{\sqrt2}
 \quad(\theta\to0).
\label{eq:v2v98-sharp-rate}
\end{equation}
At \(\theta=0\), the normalized cross coupling equals one, the stacked map
is not onto, and the common kernel jumps from \(\{0\}\) to the vertical
axis.
\end{counterexample}

\begin{proof}
The component row vectors have unit norm and inner product
\(\cos\theta\), giving \eqref{eq:v2v98-Gtheta}.  Diagonalization gives the
eigenvalues and \eqref{eq:v2v98-sharp-rate}.  At zero the two rows coincide.
\end{proof}

\begin{firewall}[Separate unit margins permit a zero common margin]
\label{fw:v2v98-separate-unit}
Counterexample~\ref{sep:v2v98-parallel} has the best possible component
margins for every \(\theta\).  The common inverse nevertheless diverges as
\(|\theta|^{-1}\).  No estimate depending only on
\(\gamma_{\rm D},\gamma_{\rm B},\gamma_{\rm H}\) can replace a cross-angle
bound.
\end{firewall}

\subsection{Small coupling around a decoupled domain}
\label{sec:v2v98-perturbation}

\begin{theorem}[One-step perturbation criterion]
\label{thm:v2v98-perturbation}
Let \(F_0:X\to Y\) satisfy
\begin{equation}
 F_0F_0^*\ge\gamma_0^2I.
\label{eq:v2v98-F0-margin}
\end{equation}
If
\begin{equation}
 F=F_0+C,
 \qquad
 \|C\|\le\delta<\gamma_0,
\label{eq:v2v98-small-C}
\end{equation}
then
\begin{equation}
 FF^*\ge(\gamma_0-\delta)^2I.
\label{eq:v2v98-perturbed-margin}
\end{equation}
The right inverse and common-domain projector therefore obey
\begin{equation}
 \|F^*(FF^*)^{-1}\|
 \le(\gamma_0-\delta)^{-1}.
\label{eq:v2v98-right-inverse-bound}
\end{equation}
\end{theorem}

\begin{proof}
For every \(y\in Y\),
\[
 \|F^*y\|
 \ge\|F_0^*y\|-\|C^*y\|
 \ge(\gamma_0-\delta)\|y\|.
\]
Squaring proves \eqref{eq:v2v98-perturbed-margin}.  The right-inverse bound
is the singular-value estimate used in V97.
\end{proof}

\begin{corollary}[Tolerance for relative boundary counterterms]
\label{cor:v2v98-counterterm-tolerance}
Suppose a reference common domain has angle \(\gamma_0\), and the change of
boundary restriction, relative counterterm, metric identification, and
characteristic operator produces a total stacked perturbation
\(C_n\) with
\begin{equation}
 \sup_n\|C_n\|\le\delta<\gamma_0.
\label{eq:v2v98-counterterm-size}
\end{equation}
Then every corrected common domain has the uniform angle
\(\gamma_0-\delta\).  Individual correction norms may be summed only after
the exact stacked perturbation has been formed; cancellation inside \(C_n\)
should be retained.
\end{corollary}

\begin{proof}
Apply Theorem~\ref{thm:v2v98-perturbation} to each \(F_n=F_0+C_n\).
\end{proof}

\subsection{Consequences for metric-domain transport}
\label{sec:v2v98-transport}

Let \(F(r)\) be the metric path of V97.  Suppose the component margins are
bounded below by \(\gamma_*>0\), the normalized scalar coupling matrices
satisfy \(\|R(r)\|\le\rho_*<1\), and \(\|F(r)\|\le M\).

\begin{corollary}[Uniform Kato connection bound from block data]
\label{cor:v2v98-Kato-bound}
The common angle satisfies
\begin{equation}
 \gamma_{\rm com}\ge\sqrt{1-\rho_*}\,\gamma_*.
\label{eq:v2v98-path-gamma}
\end{equation}
Consequently the V97 domain projector obeys
\begin{align}
 \|\dot\Pi(r)\|
 \le
 \biggl[
 &\frac{2M}{(1-\rho_*)\gamma_*^2}
 +\frac{2M^3}{(1-\rho_*)^2\gamma_*^4}
 \biggr]
 \|\dot F(r)\|.
\label{eq:v2v98-Kato-bound}
\end{align}
Thus the common-domain chart cost diverges at most quadratically in the
inverse normalized angle margin in this general estimate.
\end{corollary}

\begin{proof}
Use Theorem~\ref{thm:v2v98-Gram-criterion} in the derivative estimate
\eqref{eq:v2v97-Pi-derivative}.
\end{proof}

\begin{assessment}[The cross-angle is now the upstream boundary datum]
\label{ass:v2v98-upstream}
The V92 spectral gap, V88 BRST Hodge gap, and V96 characteristic contraction
are component margins.  V98 shows exactly what is still missing to combine
them: bounds for the three normalized cross Grams
\eqref{eq:v2v98-rho-DB}--\eqref{eq:v2v98-rho-BH}.  If the boundary
construction is block orthogonal they vanish.  If relative Ward matching
couples the sectors, their row sums must remain below one or a sharper Schur
estimate must be proved.
\end{assessment}

\begin{firewall}[A common angle does not prove all boundary analysis]
\label{fw:v2v98-angle-scope}
The common angle gives stable kernels, right inverses, and graph transport.
It does not prove the V92 heat coefficients, determinant-line triviality,
the V96 commuted normal-derivative recovery, or the interacting stress
response bound.  Those estimates live on the stable domain supplied here.
\end{firewall}

\subsection{V98 conclusion and next acquisition}
\label{sec:v2v98-conclusion}

V98 gives two quantitative routes to the V97 common-domain angle: a
normalized block--Gram row-sum criterion and a small perturbation theorem
around a decoupled boundary domain.  The two-row model proves sharp loss when
the normalized coupling reaches one.  Substitution into the V97 Kato bound
makes the metric-domain chart cost explicit.

V99 will inspect the actual field-sector block structure.  It will separate
orthogonal fermion, ghost/gauge, and harmonic variables from the shared
metric boundary data, triangularize the relative Ward matching map, and
reduce the three cross Grams to the genuinely shared metric/normal-flux
block.  If such a triangularization requires an unproved trace right
inverse, that missing inverse will be isolated instead of being hidden in
the common angle.
\clearpage
\section*{Second Volume, Version V99: Private Boundary Lifts and the Shared-Flux Cokernel}
\addcontentsline{toc}{section}{Second Volume, Version V99: Private Boundary Lifts and the Shared-Flux Cokernel}
\label{sec:v2v99}

\subsection{Purpose}
\label{sec:v2v99-purpose}

V98 bounded the common boundary angle by estimating normalized cross Grams.
That estimate allows destructive off-diagonal interference and is therefore
conservative.  In the field-sector graph, many cross terms arise from one
shared metric or normal-flux variable.  Kept as a complete block, their Gram
contribution is positive and cannot reduce a diagonal margin supplied by
independent sector lifts.

V99 proves this exact private-lift theorem.  If the private Dirac, BRST, and
harmonic trace maps are jointly onto with a uniform right inverse, arbitrary
bounded shared metric coupling preserves the common angle.  If the private
map is not onto, all remaining compatibility is compressed to its cokernel;
a single shared-flux right inverse there is necessary and sufficient for
surjectivity and yields an explicit angle.  This isolates the missing
boundary datum without taking norms of individual cross terms.

\subsection{Private and shared graph variables}
\label{sec:v2v99-decomposition}

Let
\begin{equation}
 X
 =X_{\rm D}\oplus X_{\rm B}\oplus X_{\rm H}\oplus Z
 =:X_{\rm priv}\oplus Z
\label{eq:v2v99-X}
\end{equation}
be an orthogonal graph-space decomposition.  The first three summands are
private fermionic, gauge/ghost, and harmonic boundary lift variables.  The
space \(Z\) contains shared metric, induced-geometry, and normal-flux
variables.  Let
\begin{equation}
 Y=Y_{\rm D}\oplus Y_{\rm B}\oplus Y_{\rm H}.
\label{eq:v2v99-Y}
\end{equation}
Consider boundary constraints of the form
\begin{align}
 F_{\rm D}(x,z)&=D_{\rm D}x_{\rm D}+A_{\rm D}z,
\label{eq:v2v99-FD}\\
 F_{\rm B}(x,z)&=D_{\rm B}x_{\rm B}+A_{\rm B}z,
\label{eq:v2v99-FB}\\
 F_{\rm H}(x,z)&=D_{\rm H}x_{\rm H}+A_{\rm H}z.
\label{eq:v2v99-FH}
\end{align}
Equivalently,
\begin{equation}
 F=[D,A]:X_{\rm priv}\oplus Z\longrightarrow Y,
\label{eq:v2v99-FDA}
\end{equation}
where
\begin{equation}
 D=\operatorname{diag}(D_{\rm D},D_{\rm B},D_{\rm H}),
 \qquad
 A=\begin{pmatrix}A_{\rm D}\\A_{\rm B}\\A_{\rm H}\end{pmatrix}.
\label{eq:v2v99-D-A}
\end{equation}
The decomposition is a structural hypothesis to be verified in the complete
boundary theory.  It does not assert that BRST and metric variables decouple;
their coupling is retained in \(A\).

\subsection{The complete shared packet is positive}
\label{sec:v2v99-positive-packet}

\begin{theorem}[Private lifts protect the common angle]
\label{thm:v2v99-private-lifts}
Assume
\begin{equation}
 DD^*\ge\gamma_{
m priv}^2I_Y
\label{eq:v2v99-private-margin}
\end{equation}
for some \(\gamma_{\rm priv}>0\).  Then, for every bounded shared coupling
\(A:Z\to Y\),
\begin{equation}
 FF^*
 =DD^*+AA^*
 \ge\gamma_{
m priv}^2I_Y.
\label{eq:v2v99-full-margin}
\end{equation}
Thus the V97 common-domain angle is at least \(\gamma_{\rm priv}\), with no
smallness condition on \(A\).
\end{theorem}

\begin{proof}
The adjoint of \(F=[D,A]\) is
\(F^*y=(D^*y,A^*y)\).  Orthogonality of
\(X_{\rm priv}\oplus Z\) gives
\[
 FF^*=DD^*+AA^*.
\]
The second summand is positive.  Apply
\eqref{eq:v2v99-private-margin}.
\end{proof}

\begin{corollary}[Explicit private right inverse]
\label{cor:v2v99-private-right-inverse}
Under \eqref{eq:v2v99-private-margin},
\begin{equation}
 R_Fy
 =
 \bigl(D^*(DD^*)^{-1}y,0\bigr)
\label{eq:v2v99-private-right-inverse}
\end{equation}
is a right inverse of \(F\) and
\begin{equation}
 \|R_F\|\le\gamma_{\rm priv}^{-1}.
\label{eq:v2v99-private-right-bound}
\end{equation}
\end{corollary}

\begin{proof}
The shared component is zero, so
\(FR_F=DD^*(DD^*)^{-1}=I_Y\).  The norm is bounded by the inverse smallest
singular value of \(D\).
\end{proof}

\begin{assessment}[Why the V98 row-sum can overcharge]
\label{ass:v2v99-row-sum}
Expanding \(AA^*\) into its sector blocks produces off-diagonal terms
\(A_iA_j^*\).  Estimating their absolute values separately can make the V98
comparison row sum exceed one.  Theorem~\ref{thm:v2v99-private-lifts} keeps
the complete shared packet \(AA^*\ge0\), so those entries cannot reduce the
private margin.  This is the boundary-domain analogue of the V95 Ward
cancellation and the Yang--Mills V76 complete-packet rule.
\end{assessment}

\subsection{What a private lift means in each sector}
\label{sec:v2v99-sector-meaning}

\begin{proposition}[Sufficient sectorwise lift criterion]
\label{prop:v2v99-sector-lifts}
Suppose each
\(D_i:X_i\to Y_i\),
\(i\in\{{\rm D},{\rm B},{\rm H}\}\), has a right inverse \(R_i\) with
\begin{equation}
 \|R_i\|\le\gamma_i^{-1}.
\label{eq:v2v99-sector-right-inverse}
\end{equation}
Then \eqref{eq:v2v99-private-margin} holds with
\begin{equation}
 \gamma_{\rm priv}=\min_i\gamma_i.
\label{eq:v2v99-private-gamma}
\end{equation}
\end{proposition}

\begin{proof}
For \(y=(y_i)\in Y\),
\[
 \|D^*y\|^2
 =\sum_i\|D_i^*y_i\|^2
 \ge\sum_i\gamma_i^2\|y_i\|^2
 \ge(\min_i\gamma_i^2)\|y\|^2.
\]
This is equivalent to \eqref{eq:v2v99-private-margin}.
\end{proof}

For the intended boundary theory, the three right inverses have distinct
content:
\begin{enumerate}
 \item \(R_{\rm D}\) is a graph-norm Poisson or collar lift into the selected
 tangential Dirac boundary range;
 \item \(R_{\rm B}\) lifts relative gauge/ghost boundary data while
 preserving the chosen BRST graph regularity; and
 \item \(R_{\rm H}\) lifts the incoming harmonic characteristic field in
 the V96 energy domain.
\end{enumerate}
Their uniform bounds are trace/domain estimates.  Spectral or Hodge gaps
alone do not construct these lifts.

\begin{firewall}[A spectral projector is not a bulk trace right inverse]
\label{fw:v2v99-projector-lift}
Norm convergence of the tangential Riesz projector controls which boundary
modes are selected.  It does not prove that every selected boundary datum
has a bulk extension with uniformly bounded Dirac graph norm.  The latter is
exactly the \(R_{\rm D}\) hypothesis in
Proposition~\ref{prop:v2v99-sector-lifts}.
\end{firewall}

\subsection{Compression to the missing private cokernel}
\label{sec:v2v99-cokernel}

The private map need not be onto after gauge quotienting, chirality, corner
compatibility, or removal of redundant boundary equations.  Assume only that
\(\operatorname{Ran}D\) is closed.  Let
\begin{equation}
 Q:Y\to Y_0:=\ker D^*
\label{eq:v2v99-Q}
\end{equation}
be the orthogonal cokernel projector, and let
\(Y_1=(I-Q)Y=\operatorname{Ran}D\).  Assume
\begin{equation}
 DD^*\ge\gamma_D^2(I-Q)
\label{eq:v2v99-D-closed-range}
\end{equation}
on \(Y_1\).  Define the compressed shared map
\begin{equation}
 A_0:=QA:Z\longrightarrow Y_0.
\label{eq:v2v99-A0}
\end{equation}

\begin{theorem}[Cokernel reduction of the common-domain problem]
\label{thm:v2v99-cokernel}
The stacked map \(F=[D,A]\) is onto \(Y\) if and only if
\(A_0=QA\) is onto \(Y_0\).  If \(A_0\) has a right inverse
\(R_0:Y_0\to Z\) satisfying
\begin{equation}
 \|R_0\|\le\gamma_0^{-1},
\label{eq:v2v99-R0}
\end{equation}
then \(F\) has a right inverse of norm at most
\begin{equation}
 C_F
 =
 \gamma_D^{-1}
 +\gamma_0^{-1}
 \left(1+\gamma_D^{-1}\|A\|\right).
\label{eq:v2v99-CF}
\end{equation}
Consequently
\begin{equation}
 FF^*\ge C_F^{-2}I_Y.
\label{eq:v2v99-cokernel-angle}
\end{equation}
\end{theorem}

\begin{proof}
If \(F\) is onto and \(y_0\in Y_0\), choose \((x,z)\) with
\(Dx+Az=y_0\).  Applying \(Q\) gives \(A_0z=y_0\), so \(A_0\) is onto.

Conversely, let \(y=y_1+y_0\) with \(y_0=Qy\).  Set
\begin{equation}
 z=R_0y_0,
 \qquad
 r=y_1-(I-Q)Az.
\label{eq:v2v99-triangular-step}
\end{equation}
Then \(r\in Y_1=\operatorname{Ran}D\).  Let
\begin{equation}
 R_D=D^*(DD^*|_{Y_1})^{-1}:Y_1\to X_{\rm priv}
\label{eq:v2v99-RD}
\end{equation}
and put \(x=R_Dr\).  Since \(DR_Dr=r\) and \(QAz=A_0z=y_0\),
\[
 Dx+Az=r+(I-Q)Az+QAz=y.
\]
Thus \((x,z)\) is a right inverse construction.  Moreover,
\begin{align*}
 \|z\|&\le\gamma_0^{-1}\|y_0\|,\\
 \|x\|&\le
 \gamma_D^{-1}
 \left(\|y_1\|+\|A\|\gamma_0^{-1}\|y_0\|\right).
\end{align*}
Using the sum norm first and then its domination of the Hilbert norm gives
the safe bound \eqref{eq:v2v99-CF}.  A surjective Hilbert-space operator with
a right inverse of norm \(C_F\) has smallest adjoint singular value at least
\(C_F^{-1}\), proving \eqref{eq:v2v99-cokernel-angle}.
\end{proof}

\begin{corollary}[Only shared metric/flux data need pay the cokernel]
\label{cor:v2v99-only-cokernel}
All components of \(A\) taking values in \(Y_1=\operatorname{Ran}D\) can be
removed by the private correction \(x=R_Dr\).  The only genuinely new
common-domain condition is the right inverse for
\begin{equation}
 QA:Z\to\ker D^*.
\label{eq:v2v99-shared-target}
\end{equation}
\end{corollary}

\begin{proof}
This is the triangular construction
\eqref{eq:v2v99-triangular-step}; only \(QAz\) is used to reach the private
cokernel.
\end{proof}

\subsection{Sharp failure of the shared cokernel lift}
\label{sec:v2v99-failure}

\begin{counterexample}[Vanishing shared-flux singular value]
\label{sep:v2v99-cokernel-failure}
Let \(X_{\rm priv}=Z=\mathbb R\), \(Y=\mathbb R^2\), and
\begin{equation}
 Dx=(x,0),
 \qquad
 A_nz=(0,n^{-1}z).
\label{eq:v2v99-diagonal-example}
\end{equation}
The private range is the first coordinate with margin one, while the
compressed shared map on the cokernel has margin \(n^{-1}\).  The full map is
\begin{equation}
 F_n(x,z)=(x,n^{-1}z)
\label{eq:v2v99-Fn-example}
\end{equation}
with common angle \(n^{-1}\).  It is onto for every finite \(n\) but loses
surjectivity in the limit.
\end{counterexample}

\begin{proof}
Here \(Q(y_1,y_2)=(0,y_2)\), so
\(QA_nz=(0,n^{-1}z)\).  The singular values of \(F_n\) are \(1\) and
\(n^{-1}\).  At the limit the second boundary equation has no lift.
\end{proof}

\begin{firewall}[Finite-regulator solvability is insufficient]
\label{fw:v2v99-finite-solvability}
Counterexample~\ref{sep:v2v99-cokernel-failure} has a common boundary domain
and an onto stacked map at every finite regulator.  The right-inverse norm
still diverges.  Uniform trace/right-inverse control is required before
regulator removal; solvability at each cutoff does not supply it.
\end{firewall}

\subsection{Metric transport with arbitrary shared coupling}
\label{sec:v2v99-metric}

\begin{proposition}[Uniform domain transport from private lifts]
\label{prop:v2v99-metric}
Let \(r\mapsto D(r),A(r)\) be \(C^1\), and assume
\begin{equation}
 D(r)D(r)^*\ge\gamma_{
m priv}^2I
\label{eq:v2v99-path-private}
\end{equation}
uniformly.  Then the common angle of
\(F(r)=[D(r),A(r)]\) is at least \(\gamma_{
m priv}\), regardless of
\(\|A(r)\|\).  The V97 Kato projector derivative obeys
\begin{align}
 \|\dot\Pi(r)\|
 \le
 \left(
 2M\gamma_{
m priv}^{-2}
 +2M^3\gamma_{
m priv}^{-4}
 \right)
 \bigl(\|\dot D(r)\|+\|\dot A(r)\|\bigr),
\label{eq:v2v99-metric-bound}
\end{align}
where \(M\ge\sup_r\|[D(r),A(r)]\|\).
\end{proposition}

\begin{proof}
The angle is Theorem~\ref{thm:v2v99-private-lifts}.  Since
\(\|\dot F\|\le\|\dot D\|+\|\dot A\|\), insert the protected angle in
Lemma~\ref{lem:v2v97-Pi-derivative}.
\end{proof}

\begin{assessment}[Large coupling affects variation, not solvability]
\label{ass:v2v99-large-A}
A large shared block \(A\) cannot destroy the angle when private lifts are
onto, but it enlarges \(M\), \(\dot F\), and hence the metric-domain chart
cost.  Static domain solvability and differentiable source control are
therefore separate estimates even in the protected private-lift regime.
\end{assessment}

\subsection{Application ledger for the Standard Model boundary}
\label{sec:v2v99-SM-ledger}

\begin{programme}[Private-lift/cokernel audit]
\label{prog:v2v99-audit}
For the intended regulator family:
\begin{enumerate}
 \item decompose the graph space into private fermion, relative gauge/ghost,
 harmonic characteristic, and shared metric/normal-flux variables;
 \item construct a collar Poisson right inverse for the selected chiral
 Dirac trace with a regulator-uniform graph norm;
 \item construct a relative BRST boundary lift after quotienting redundant
 gauge equations;
 \item construct the incoming harmonic characteristic lift in the V96
 energy domain;
 \item compute the cokernel \(\ker D_n^*\) left by any failures of items
 2--4; and
 \item prove a uniform right inverse for the compressed shared map
 \(Q_nA_n\), including normal stress and eta-sector boundary data.
\end{enumerate}
If items 2--4 are all onto with uniform bounds, item 6 is vacuous and the
shared block cannot close the common angle.
\end{programme}

\begin{firewall}[The eta sector is not declared a private lift]
\label{fw:v2v99-eta-lift}
The eta invariant is determinant-line phase data, not an arbitrary local
boundary value.  It may enter the affine relative Ward datum or the shared
cokernel, but V99 does not treat it as a freely prescribable private trace.
Doing so would erase the spectral-flow and global-holonomy gates of V92.
\end{firewall}

\subsection{V99 conclusion and V100 acquisition}
\label{sec:v2v99-conclusion}

V99 improves the generic V98 angle estimate by preserving the exact shared
Gram packet.  Uniform private lifts give the full common angle with no
smallness condition on metric or normal-flux coupling.  When private
surjectivity fails, the complete boundary-domain problem reduces to the one
compressed map \(QA:Z\to\ker D^*\), with the explicit right-inverse bound
\eqref{eq:v2v99-CF}.  A two-coordinate example proves that finite-cutoff
solvability does not prevent collapse of this shared lift.

V100 will close this volume by proving a collar right inverse for the model
Dirac graph trace under a uniform tangential spectral gap and bounded
product-geometry perturbation.  It will state exactly which estimate survives
for the interacting curved boundary and which BRST/shared-cokernel lifts
remain.  In accordance with the volume rule, the completed V100 file will be
frozen with suffix \texttt{\_f2}; the following continuation will restart at
V1 with a context and major-results summary.

\clearpage
\section*{Second Volume, Version V100: Collar Right Inverse for the Chiral Dirac Trace}
\addcontentsline{toc}{section}{Second Volume, Version V100: Collar Right Inverse for the Chiral Dirac Trace}
\label{sec:v2v100}

\subsection{Purpose}
\label{sec:v2v100-purpose}

V99 identified a concrete boundary-domain acquisition: a uniformly bounded
bulk graph lift for every selected tangential Dirac boundary datum.  V100
constructs that lift on a product collar and proves stability under uniformly
bounded first-order coefficient perturbations.  The construction separates
three inputs which had previously been conflated: the boundary spectral
projector selects the data, a Poisson extension lifts them, and a uniform
bulk coefficient bound controls the Dirac graph norm.

This closes the Dirac private-lift entry in the model bounded-geometry
setting.  It does not construct the relative BRST lift or the shared
normal-flux cokernel lift, and it does not derive the V92 heat coefficients
from the trace extension.

\subsection{Reference collar and Sobolev scale}
\label{sec:v2v100-collar}

Let \(\Sigma\) be a closed boundary manifold and let
\(A_0\) be a self-adjoint first-order elliptic operator on a Hermitian bundle
\(E_\Sigma\to\Sigma\).  Put
\begin{equation}
 B=(I+A_0^2)^{1/2}\ge I
\label{eq:v2v100-B}
\end{equation}
and define the spectral Sobolev norm
\begin{equation}
 \|y\|_{H_B^s}=\|B^sy\|_{L^2(\Sigma)}.
\label{eq:v2v100-HBs}
\end{equation}
Ellipticity makes this norm equivalent to the usual \(H^s\) norm, with
constants fixed by the reference geometry.

On the collar
\(\mathcal C_\ell=[0,\ell)_r\times\Sigma\), choose
\(\chi\in C_c^\infty([0,\ell))\) with \(\chi(0)=1\).  Define
\begin{equation}
 (\mathcal Ey)(r)
 =\chi(r)e^{-rB}y.
\label{eq:v2v100-Poisson-extension}
\end{equation}

\begin{lemma}[Poisson collar estimate]
\label{lem:v2v100-Poisson}
The map \(\mathcal E\) is a right inverse of the boundary trace and obeys
\begin{equation}
 \|\mathcal Ey\|_{H^1_B(\mathcal C_\ell)}
 \le C_\chi\|y\|_{H_B^{1/2}},
\label{eq:v2v100-Poisson-bound}
\end{equation}
where
\begin{equation}
 \|u\|_{H^1_B}^2
 =
 \int_0^\ell
 \left(
  \|u(r)\|^2+|Bu(r)\|^2+|\partial_ru(r)\|^2
 \right)\,\mathrm dr.
\label{eq:v2v100-H1B}
\end{equation}
The constant depends only on a finite norm of \(\chi\), not on \(y\).
\end{lemma}

\begin{proof}
The trace at \(r=0\) is \(y\).  By the spectral theorem, for every
nonnegative Borel function of \(B\), integration can be performed on each
spectral value \(b\ge1\).  Without the cutoff,
\begin{align}
 \int_0^\infty\|Be^{-rB}y\|^2\,\mathrm dr
 &=\frac12\|B^{1/2}y\|^2,
\label{eq:v2v100-B-integral}\\
 \int_0^\infty\|\partial_re^{-rB}y\|^2\,\mathrm dr
 &=\frac12\|B^{1/2}y\|^2,
\label{eq:v2v100-r-integral}\\
 \int_0^\infty\|e^{-rB}y\|^2\,\mathrm dr
 &=\frac12\|B^{-1/2}y\|^2
 \le\frac12\|B^{1/2}y\|^2.
\label{eq:v2v100-L2-integral}
\end{align}
Multiplication by \(\chi\) only decreases the terms without
\(\partial_r\), up to \(\|\chi\|_\infty\).  For the normal derivative,
\[
 \partial_r(\chi e^{-rB}y)
 =\chi'e^{-rB}y-\chi Be^{-rB}y.
\]
Use \eqref{eq:v2v100-B-integral}--\eqref{eq:v2v100-L2-integral} and the
triangle inequality.  This proves \eqref{eq:v2v100-Poisson-bound}.
\end{proof}

\subsection{Uniform Dirac graph control}
\label{sec:v2v100-graph}

Let \(\mathscr D_n\) be first-order Dirac-type operators on the collar,
after all bundles are identified with the reference bundle.  Assume the
uniform coefficient estimate
\begin{equation}
 \|\mathscr D_nu\|_{L^2(\mathcal C_\ell)}
 \le C_D
 \left(
  \|\partial_ru\|_{L^2}
  +\|Bu\|_{L^2}
  +\|u\|_{L^2}
 \right)
\label{eq:v2v100-coefficient-bound}
\end{equation}
for every smooth collar-supported \(u\), with \(C_D\) independent of \(n\).
This follows from uniform bounded geometry and uniform equivalence of the
first-order tangential Sobolev norms; it is stated explicitly because the
spectral gap alone does not imply it.

Define the graph norm
\begin{equation}
 \|u\|_{X_{{\rm D},n}}^2
 =
 \|u\|_{L^2}^2+|\mathscr D_nu\|_{L^2}^2.
\label{eq:v2v100-Dirac-graph}
\end{equation}

\begin{theorem}[Uniform collar graph lift]
\label{thm:v2v100-graph-lift}
Under \eqref{eq:v2v100-coefficient-bound},
\begin{equation}
 \|\mathcal Ey\|_{X_{{\rm D},n}}
 \le C_{\rm lift}\|y\|_{H_B^{1/2}}
\label{eq:v2v100-graph-lift}
\end{equation}
uniformly in \(n\), where \(C_{\rm lift}\) depends only on
\(C_D\), the reference norm equivalence, and \(\chi\).
\end{theorem}

\begin{proof}
Apply \eqref{eq:v2v100-coefficient-bound} to \(u=\mathcal Ey\), then use
Lemma~\ref{lem:v2v100-Poisson}.  The \(L^2\) part of the graph norm is
already one of the terms in \eqref{eq:v2v100-H1B}.
\end{proof}

\begin{remark}[No inverse elliptic estimate is used]
\label{rem:v2v100-no-Garding}
The theorem needs only an upper bound for \(\mathscr D_n\mathcal Ey\).
It does not claim that the unrestricted Dirac graph norm controls the full
\(H^1\) norm without a boundary condition.  The extension is constructed in
\(H^1\) first and then estimated in the graph norm.
\end{remark}

\subsection{Selected chiral spectral range}
\label{sec:v2v100-selected-range}

Let \(P_n\) be the tangential spectral projector defining the forbidden or
prescribed chiral boundary component.  Let
\begin{equation}
 U_n:\operatorname{Ran}P_n\longrightarrow Y_{\rm D}
\label{eq:v2v100-Un}
\end{equation}
identify its range with one fixed Hilbert boundary data space.  Assume
\begin{align}
 \|P_n\|_{\mathcal L(H_B^{1/2})}&\le C_P,
\label{eq:v2v100-Pn-Hhalf}\\
 \|U_n^{-1}y\|_{H_B^{1/2}}&\le C_U\|y\|_{Y_{\rm D}}
\label{eq:v2v100-Un-bound}
\end{align}
uniformly.  A uniform tangential gap and norm-resolvent comparison stabilize
the Riesz ranges; the \(H^{1/2}\) estimate is an additional graph-scale
version of that comparison.

Define
\begin{equation}
 F_{{\rm D},n}u
 =U_nP_n\gamma u,
\label{eq:v2v100-FDn}
\end{equation}
where \(\gamma\) is the trace at \(r=0\).

\begin{theorem}[Uniform right inverse for the selected Dirac trace]
\label{thm:v2v100-Dirac-right-inverse}
The operator
\begin{equation}
 R_{{\rm D},n}y
 =\mathcal E(U_n^{-1}y)
\label{eq:v2v100-RDn}
\end{equation}
is a right inverse of \(F_{{\rm D},n}\) and satisfies
\begin{equation}
 \|R_{{\rm D},n}\|_{
 \mathcal L(Y_{\rm D},X_{{\rm D},n})}
 \le C_{\rm lift}C_U.
\label{eq:v2v100-RDn-bound}
\end{equation}
Consequently the Dirac private margin in V99 obeys
\begin{equation}
 \gamma_{\rm D}
 \ge(C_{\rm lift}C_U)^{-1}.
\label{eq:v2v100-gammaD}
\end{equation}
\end{theorem}

\begin{proof}
For \(y\in Y_{\rm D}\), one has
\(U_n^{-1}y\in\operatorname{Ran}P_n\).  Since
\(\gamma\mathcal E=I\),
\[
 F_{{\rm D},n}R_{{\rm D},n}y
 =U_nP_nU_n^{-1}y=y.
\]
The norm bound is Theorem~\ref{thm:v2v100-graph-lift} followed by
\eqref{eq:v2v100-Un-bound}.  The smallest adjoint singular value is at least
the reciprocal of a right-inverse norm, giving \eqref{eq:v2v100-gammaD}.
\end{proof}

\begin{corollary}[Private-lift protection of shared metric coupling]
\label{cor:v2v100-private-protection}
If the relative BRST and harmonic private lifts also have margins
\(\gamma_{\rm B},\gamma_{\rm H}>0\), then the full common-domain angle is at
least
\begin{equation}
 \min\{(C_{\rm lift}C_U)^{-1},
       \gamma_{\rm B},
       \gamma_{\rm H}\},
\label{eq:v2v100-common-angle}
\end{equation}
independently of the size of the bounded shared metric/normal-flux block.
\end{corollary}

\begin{proof}
Apply Proposition~\ref{prop:v2v99-sector-lifts} and
Theorem~\ref{thm:v2v99-private-lifts}.
\end{proof}

\subsection{Metric and regulator stability of the lift}
\label{sec:v2v100-stability}

\begin{proposition}[Convergence in the fixed collar chart]
\label{prop:v2v100-lift-convergence}
Suppose
\begin{equation}
 U_n^{-1}\longrightarrow U^{-1}
 \quad\hbox{in }\mathcal L(Y_{\rm D},H_B^{1/2})
\label{eq:v2v100-U-convergence}
\end{equation}
and \(\mathcal E\) is the fixed reference extension
\eqref{eq:v2v100-Poisson-extension}.  Then
\begin{equation}
 R_{{\rm D},n}\longrightarrow R_{\rm D}:=\mathcal EU^{-1}
 \quad\hbox{in }\mathcal L(Y_{\rm D},H_B^1(\mathcal C_\ell)).
\label{eq:v2v100-R-convergence}
\end{equation}
If also \(\mathscr D_n\to\mathscr D\) in
\(\mathcal L(H_B^1,L^2)\), the convergence holds in the identified Dirac
graph realizations.
\end{proposition}

\begin{proof}
The first assertion follows by composing
\eqref{eq:v2v100-U-convergence} with the bounded map \(\mathcal E\).  For
the graph term, write
\[
 \mathscr D_nR_{{\rm D},n}-\mathscr DR_{\rm D}
 =
 (\mathscr D_n-\mathscr D)R_{{\rm D},n}
 +\mathscr D(R_{{\rm D},n}-R_{\rm D})
\]
and take norms.
\end{proof}

\begin{proposition}[Differentiable range transport]
\label{prop:v2v100-lift-derivative}
Let \(r\mapsto U(r)^{-1}\) be \(C^1\) as a map
\(Y_{\rm D}\to H_B^{1/2}\).  Then
\begin{equation}
 R_{\rm D}(r)=\mathcal EU(r)^{-1}
\label{eq:v2v100-R-path}
\end{equation}
is \(C^1\) and
\begin{equation}
 \|\dot R_{\rm D}(r)\|_{
 \mathcal L(Y_{\rm D},H_B^1)}
 \le C_\chi\|\partial_rU(r)^{-1}\|_{
 \mathcal L(Y_{\rm D},H_B^{1/2})}.
\label{eq:v2v100-R-derivative}
\end{equation}
This lift derivative is one component of the common-domain chart cost in
V97.
\end{proposition}

\begin{proof}
Differentiate \eqref{eq:v2v100-R-path}; \(\mathcal E\) is fixed and bounded
by Lemma~\ref{lem:v2v100-Poisson}.
\end{proof}

\subsection{Why the tangential gap is not enough}
\label{sec:v2v100-counterexample}

\begin{counterexample}[Loss of the normal coefficient margin]
\label{sep:v2v100-normal-loss}
Let
\begin{equation}
 X_n=\{u\in H^1(0,1):u(1)=0\},
 \qquad
 \|u\|_{X_n}^2=\|u\|_{L^2}^2+n^2\|u'\|_{L^2}^2,
\label{eq:v2v100-Xn-example}
\end{equation}
and let \(\gamma u=u(0)\).  Every right inverse
\(R_n:\mathbb R\to X_n\) of \(\gamma\) satisfies
\begin{equation}
 \|R_n\|\ge n.
\label{eq:v2v100-right-inverse-loss}
\end{equation}
The boundary projector is the identity and has a perfect fixed gap; the
failure is entirely in the unbounded normal principal coefficient.
\end{counterexample}

\begin{proof}
If \(u(0)=1\) and \(u(1)=0\), the fundamental theorem of calculus and
Cauchy--Schwarz give
\[
 1=\left|\int_0^1u'(r)\,\mathrm dr\right|
 \le\|u'\|_{L^2}.
\]
Hence \(\|u\|_{X_n}\ge n\).  Apply this to \(u=R_n1\).
\end{proof}

\begin{firewall}[Three independent Dirac boundary estimates]
\label{fw:v2v100-three-estimates}
The tangential Riesz gap controls the selected range.  The collar coefficient
bound controls the bulk graph lift.  Strong ellipticity and heat estimates
control the eta invariant.  None of these three statements implies the
other two, although all are required on the same common domain.
\end{firewall}

\subsection{Second-volume result ledger}
\label{sec:v2v100-ledger}

The final fifteen versions of this volume establish the following chain:
\begin{enumerate}
 \item V86--V90 isolate regulator-stable spectral clusters, relative
 counterterms, and stable removable/anomalous cohomology splittings, with
 counterexamples when norm-resolvent or gap hypotheses are weakened.
 \item V91 verifies every perturbative local one-generation Standard Model
 gauge and mixed hypercharge--gravity anomaly coefficient and the ordinary
 \(SU(2)\) parity test.
 \item V92 constructs the relative bulk--boundary Ward complex and a
 uniform heat-subtraction criterion for eta convergence.
 \item V93--V96 derive the relative diffeomorphism Ward identities, separate
 the Weyl anomaly from conservation, compute the connection-variation
 source, prove its exact cancellation in the coupled Einstein operator, and
 obtain the maximally dissipative harmonic-constraint estimate.
 \item V97--V99 construct the stable common boundary graph, prove the
 common-angle and Kato transport theorems, and reduce any failure of private
 boundary lifting to a shared metric/normal-flux cokernel.
 \item V100 supplies the model uniform Dirac private lift and an explicit
 lower bound for its contribution to the common angle.
\end{enumerate}

\begin{firewall}[Boundary and full-continuum claim boundary]
\label{fw:v2v100-claim-boundary}
The remaining boundary acquisitions include a regulator-uniform relative
BRST lift, the compressed shared normal-flux/eta cokernel lift, the commuted
V96 normal-derivative estimate, the V92 heat bounds on the interacting
domain, and determinant-line triviality or controlled holonomy.  Beyond the
boundary, the full interacting stress-response norm and the global
all-frequency Einstein stability problem remain open.  V100 does not claim
the full Standard Model--Einstein continuum.
\end{firewall}

\subsection{V100 conclusion and volume rollover}
\label{sec:v2v100-conclusion}

V100 proves a uniform Poisson collar right inverse from the selected
\(H^{1/2}\) chiral boundary range into the bulk Dirac graph.  Combined with
V99, this protects the common boundary angle against arbitrary bounded
shared metric coupling once the BRST and harmonic private lifts are also
constructed.  The normal-coefficient counterexample proves why a tangential
spectral gap alone cannot supply this result.

This completes the second note volume at V100.  The file is to be frozen with
suffix \texttt{\_f2}.  The next active note restarts at V1 with a concise
context and major-results summary, then attacks the relative BRST private
lift and the shared normal-flux cokernel without reproducing the proofs
already retained in the frozen volumes.
\end{document}
